% GENERATED by bin/curate/gen_tutorials_guide.py -- DO NOT EDIT BY HAND.
% Assembled from each case's controlDict description + header comment.

\subsection{adsorption}
\grouptheory{adsorption}
\clearpage
\tutentry{isotherm01\_eval\_13x\_co2}{tutorials/props/adsorption/isotherm01\_eval\_13x\_co2}
\begin{tutscheme}
\begin{tikzpicture}[>={Stealth[length=2mm]}, every node/.style={font=\scriptsize}]
  \node[draw, rounded corners=2pt, align=center, minimum width=3.1cm, minimum height=0.8cm, fill=black!4] at (0.0,-0.0) {\textbf{eval\_13x\_co2}\\[-1pt]{\tiny isothermEval}};
\end{tikzpicture}
\end{tutscheme}
\tutwhat{Adsorption isotherm bench: CO2/zeolite13X q(T,p) grid + Henry/saturation/pin/unit-invariance gates}
\begin{tutnarrative}
One isothermEval operation: the curated CO2/zeolite13X langmuir record on a 3-T x 4-p grid (K, Pa -- grid values are RAW canonical SI), with the unit-invariance twin (gate d) declared case-locally in constant/. Temperatures match the Jaschik 2020 dataset the FIT tutorial (isotherm02\_fit\_jaschik) regresses -- eval here, fit there.

\end{tutnarrative}
\begin{tutgolden}
\goldrow{diag}{eval\_13x\_co2}{T0}{293.15}
\goldrow{diag}{eval\_13x\_co2}{T1}{313.15}
\goldrow{diag}{eval\_13x\_co2}{T2}{333.15}
\goldrow{diag}{eval\_13x\_co2}{b\_T0}{12.2661169044}
\goldrow{diag}{eval\_13x\_co2}{b\_T1}{3.77235323402}
\goldrow{diag}{eval\_13x\_co2}{b\_T2}{1.33660410787}
\goldrow{diag}{eval\_13x\_co2}{henry\_T0}{0.000613305845222}
\goldrow{diag}{eval\_13x\_co2}{henry\_T1}{0.000188617661701}
\goldrow{diag}{eval\_13x\_co2}{henry\_T2}{6.68302053933e-05}
\goldrow{diag}{eval\_13x\_co2}{qsat}{5}
\goldrow{diag}{eval\_13x\_co2}{tRef}{298.15}
\goldrow{diag}{eval\_13x\_co2}{gate\_henry}{1}
\goldrow{diag}{eval\_13x\_co2}{gate\_henry\_maxdev}{1.22661154271e-07}
\goldrow{diag}{eval\_13x\_co2}{gate\_saturation}{1}
\goldmore{9}
\end{tutgolden}

\clearpage
\tutentry{isotherm02\_fit\_jaschik}{tutorials/props/adsorption/isotherm02\_fit\_jaschik}
\begin{tutscheme}
\begin{tikzpicture}[>={Stealth[length=2mm]}, every node/.style={font=\scriptsize}]
  \node[draw, rounded corners=2pt, align=center, minimum width=3.1cm, minimum height=0.8cm, fill=black!4] at (0.0,-0.0) {\textbf{fit\_jaschik}\\[-1pt]{\tiny fitParameters}};
  \node[draw, rounded corners=2pt, align=center, minimum width=3.1cm, minimum height=0.8cm, fill=black!4] at (3.6,-0.0) {\textbf{recovery\_multiT}\\[-1pt]{\tiny fitParameters}};
  \node[draw, rounded corners=2pt, align=center, minimum width=3.1cm, minimum height=0.8cm, fill=black!4] at (7.2,-0.0) {\textbf{gate\_singleT}\\[-1pt]{\tiny fitParameters}};
  \node[draw, rounded corners=2pt, align=center, minimum width=3.1cm, minimum height=0.8cm, fill=black!4] at (0.0,-1.3) {\textbf{gate\_henry}\\[-1pt]{\tiny fitParameters}};
\end{tikzpicture}
\end{tutscheme}
\tutwhat{Isotherm regression: CO2/13X Jaschik 2020 (CC BY) multi-T Langmuir fit + recovery/single-T/Henry identifiability gates}
\begin{tutnarrative}
Four fitParameters(kind isotherm) operations: the REAL Jaschik fit plus the three identifiability gates on zero-noise synthetics generated from a KNOWN surface (q\_max 4.6 mol/kg, b\_ref 3.0e-4 1/Pa @ 298.15 K, dH\_ads -38000 J/mol) -- so every gate has an exact truth to recover or to honestly refuse. Inline synthetic data is canonical SI (K, Pa, mol/kg); the vendored Jaschik dataset declares its own units (bar, mol/kg) and is converted at the data boundary.

\end{tutnarrative}
\begin{tutgolden}
\goldrow{diag}{fit\_jaschik}{converged}{1.00000e+00}
\goldrow{diag}{fit\_jaschik}{identifiable}{1.00000e+00}
\goldrow{diag}{fit\_jaschik}{n\_data}{1.32000e+02}
\goldrow{diag}{fit\_jaschik}{n\_params}{3.00000e+00}
\goldrow{diag}{fit\_jaschik}{dof}{1.29000e+02}
\goldrow{diag}{fit\_jaschik}{single\_T}{0.00000e+00}
\goldrow{diag}{fit\_jaschik}{henry\_regime}{0.00000e+00}
\goldrow{diag}{fit\_jaschik}{endothermic\_suspicious}{0.00000e+00}
\goldrow{diag}{fit\_jaschik}{fit.q\_max.value}{4.31448e+00}
\goldrow{diag}{fit\_jaschik}{fit.q\_max.stderr}{5.17246e-02}
\goldrow{diag}{fit\_jaschik}{fit.b\_ref.value}{4.06758e-04}
\goldrow{diag}{fit\_jaschik}{fit.b\_ref.stderr}{2.90989e-05}
\goldrow{diag}{fit\_jaschik}{fit.dH\_ads.value}{-3.73835e+04}
\goldrow{diag}{fit\_jaschik}{fit.dH\_ads.stderr}{1.90453e+03}
\goldmore{69}
\end{tutgolden}

\clearpage
\tutentry{isotherm03\_fit\_refusals}{tutorials/props/adsorption/isotherm03\_fit\_refusals}
\begin{tutscheme}
\begin{tikzpicture}[>={Stealth[length=2mm]}, every node/.style={font=\scriptsize}]
  \node[draw, rounded corners=2pt, align=center, minimum width=3.1cm, minimum height=0.8cm, fill=black!4] at (0.0,-0.0) {\textbf{gate\_henry\_nonconverged}\\[-1pt]{\tiny fitParameters}};
  \node[draw, rounded corners=2pt, align=center, minimum width=3.1cm, minimum height=0.8cm, fill=black!4] at (3.6,-0.0) {\textbf{gate\_nonconverged}\\[-1pt]{\tiny fitParameters}};
  \node[draw, rounded corners=2pt, align=center, minimum width=3.1cm, minimum height=0.8cm, fill=black!4] at (7.2,-0.0) {\textbf{gate\_langfull\_nonconverged}\\[-1pt]{\tiny fitParameters}};
\end{tikzpicture}
\end{tutscheme}
\tutwhat{Fit refusal gates: non-converged fits refuse their verdicts by name}
\begin{tutnarrative}
Two ADVERSARIAL fitParameters(kind isotherm) operations -- see the controlDict header. Synthetic data generated EXACTLY from the known surface (q\_max 4.6 mol/kg, b\_298 3.0e-4 1/Pa, dH\_ads -38000 J/mol, tRef 298.15 K) -- zero noise, so the ONLY failure mode under test is the strangled optimiser, never the data.

\end{tutnarrative}

\subsection{compare}
\grouptheory{compare}
\clearpage
\tutentry{aceticWater01\_txy\_association}{tutorials/props/compare/aceticWater01\_txy\_association}
\begin{tutscheme}
\begin{tikzpicture}[>={Stealth[length=2mm]}, every node/.style={font=\scriptsize}]
  \node[draw, rounded corners=2pt, align=center, minimum width=3.1cm, minimum height=0.8cm, fill=black!4] at (0.0,-0.0) {\textbf{bp\_aceticAcid}\\[-1pt]{\tiny propertyPoint}};
  \node[draw, rounded corners=2pt, align=center, minimum width=3.1cm, minimum height=0.8cm, fill=black!4] at (3.6,-0.0) {\textbf{bp\_water}\\[-1pt]{\tiny propertyPoint}};
  \node[draw, rounded corners=2pt, align=center, minimum width=3.1cm, minimum height=0.8cm, fill=black!4] at (7.2,-0.0) {\textbf{txy\_calvar2005}\\[-1pt]{\tiny propertyScan1D}};
\end{tikzpicture}
\end{tutscheme}
\tutwhat{Acetic acid/water T-x-y at 1 atm on UNIFAC with the ASSOCIATED ideal-gas vapour (monomer + dimer), against Calvar, Dominguez \& Tojo (2005)}
\begin{tutnarrative}
WHAT A DIMERISED VAPOUR DOES TO A T-x-y DIAGRAM. Acetic acid's vapour is largely dimers near its boiling point, so the acid is LESS volatile than its vapour pressure alone suggests: the dimers take up partial pressure that the monomer -- the species in equilibrium with the liquid -- does not. This case evaluates the bubble temperature and the vapour composition at the 24 liquid compositions Calvar, Dominguez \& Tojo measured at 101.325 kPa (constant/experimental/calvar2005\_txy.csv), plus the two pure boiling points, on UNIFAC (predictive, no fitted pair) and the associated ideal-gas vapour. The gate check\_reference\_rung arm (k) holds the pure limit and the latent heat at Tb and prints the comparison with these measurements; the numbers are recorded in docs/design/one-record-two-vapours-acetic-acid.md section 7. Nothing here is tuned.

WHAT IS NOT IN IT: no regressed acetic-acid/water pair (the catalogue holds none); the measured points were read from the NIST ThermoML archive copy of the article and NOT checked against the article itself.

\end{tutnarrative}
\begin{tutgolden}
\goldrow{diag}{bp\_aceticAcid}{Cp\_ig}{80.6641982498}
\goldrow{diag}{bp\_aceticAcid}{H\_R}{-22841.7769643}
\goldrow{diag}{bp\_aceticAcid}{H\_ig}{-425398.226035}
\goldrow{diag}{bp\_aceticAcid}{H\_real}{-448240.002999}
\goldrow{diag}{bp\_aceticAcid}{P}{101325}
\goldrow{diag}{bp\_aceticAcid}{S\_R}{-51.4737821304}
\goldrow{diag}{bp\_aceticAcid}{S\_ig}{302.310579052}
\goldrow{diag}{bp\_aceticAcid}{S\_real}{250.836796922}
\goldrow{diag}{bp\_aceticAcid}{T}{391}
\goldrow{diag}{bp\_aceticAcid}{T\_bubble}{391.008522453}
\goldrow{diag}{bp\_aceticAcid}{Z}{1}
\goldrow{diag}{bp\_aceticAcid}{gamma}{1.11492042846}
\goldrow{diag}{bp\_aceticAcid}{v\_molar}{0.0320844301371}
\goldrow{diag}{bp\_water}{Cp\_ig}{34.2425437543}
\goldmore{15}
\end{tutgolden}

\clearpage
\tutentry{acetone01\_ipa\_water\_azeotrope}{tutorials/props/compare/acetone01\_ipa\_water\_azeotrope}
\begin{tutscheme}
\begin{tikzpicture}[>={Stealth[length=2mm]}, every node/.style={font=\scriptsize}]
  \node[draw, rounded corners=2pt, align=center, minimum width=3.1cm, minimum height=0.8cm, fill=black!4] at (0.0,-0.0) {\textbf{bp\_ipa}\\[-1pt]{\tiny propertyPoint}};
  \node[draw, rounded corners=2pt, align=center, minimum width=3.1cm, minimum height=0.8cm, fill=black!4] at (3.6,-0.0) {\textbf{bp\_water}\\[-1pt]{\tiny propertyPoint}};
  \node[draw, rounded corners=2pt, align=center, minimum width=3.1cm, minimum height=0.8cm, fill=black!4] at (7.2,-0.0) {\textbf{azeo\_luyben\_x}\\[-1pt]{\tiny propertyPoint}};
  \node[draw, rounded corners=2pt, align=center, minimum width=3.1cm, minimum height=0.8cm, fill=black!4] at (0.0,-1.3) {\textbf{txy\_unifac}\\[-1pt]{\tiny propertyScan1D}};
\end{tikzpicture}
\end{tutscheme}
\tutwhat{IPA/water T-x-y at 1 atm from PREDICTIVE UNIFAC on a lake-estimated isopropanol, held against Luyben's published azeotrope}
\begin{tutnarrative}
THE FIRST RUNNABLE SLICE of the acetone-from-2-propanol reference case (docs/design/acetone-ipa-reference-case.md). It answers ONE question, and the question is chosen because it can be answered honestly today:

how close does a component Choupo has never curated, priced by a model that saw no data for it, land on a published azeotrope?

WHAT IS BEING TESTED, and it is a stack of three things at once -- said up front, because a single number that agrees can hide two errors that cancel:

1. ISOPROPANOL IS A LAKE ESTIMATE. It is NOT in data/standards/. The record in this case's constant/components/ came from data/groupEstimative/, where every NUMBER is a Joback / Lee-Kesler estimate from the molecular structure -- including Tb, Tc, Pc, omega and hence the whole Ambrose-Walton vapour pressure. Its estimated Tb is 359.98 K; the measured normal boiling point Luyben quotes is 355.4 K. The pure component is already 4.6 K out before any mixture model is asked anything. 2. UNIFAC IS PREDICTIVE. No pair was fitted, because Choupo ships no UNIQUAC or NRTL pair for IPA-water and inventing one would convert `unsourced` into `falsely sourced`. gamma comes from the groups. 3. WATER IS CURATED. It is the one leg of the three that rests on reviewed data, which is what makes the comparison legible at all.

THE COMPARISON POINTS: W. L. Luyben, "Design and Control of the Acetone Process via Dehydrogenation of 2-Propanol", Ind. Eng. Chem. Res. 50 (2011) 1206-1218, DOI 10.1021/ie901923a. He states the IPA/water azeotrope as 67.32 mol \% IPA at 353.4 K and 1 atm, and normal boiling points IPA 355.4 K, water 373 K. Those numbers come out of his UNIQUAC model, fitted to VLE data neither this case nor its components have ever seen -- so for the predictive route here they are an INDEPENDENT target.

A NOTE ON WHAT KIND OF ROW PINS THIS, because it is not an `anchor`. Luyben's azeotrope is his UNIQUAC model's value, not a measurement, and the AAD computed here is derived from that number and this model TOGETHER -- so nobody published it and no primary can be quoted for it. It is pinned as a self-recorded `aad` row: what the suite holds is that the agreement has not MOVED. Anchoring properly needs a bubble temperature published as an operation diagnostic, which is a separate gap -- see the note at the end.

WHAT THIS CASE IS NOT. It is not a reproduction of Luyben's flowsheet, and it makes no claim about the acetone process. The full reproduction needs three curated UNIQUAC pairs that do not exist; the design record's section 4 lists them and section 5 says which case category that would be. This is the bottom rung: the thermodynamics, alone, against three published points.

A GAP FOUND BUILDING THIS, and left visible rather than worked around. The natural shape for this case is three `propertyPoint` ops asking for `T\_bubble` at the three compositions, because a propertyPoint publishes into the operation diagnostics where a golden row can read it. It does not work: `PropertyPoint` never reads its own `properties ( ... )` list -- it emits a fixed set and discards what the case asked for. Both the class header ("properties (all); // or pick a subset") and the published JSON schema (which names `T\_bubble` among the accepted keywords) say otherwise. That is a parsed-and-dropped field, the shape this project has paid for three times, and it is why the comparison below rides the overlay instead.

\end{tutnarrative}
\begin{tutgolden}
\goldrow{diag}{bp\_ipa}{Cp\_ig}{99.085675}
\goldrow{diag}{bp\_ipa}{H\_R}{0}
\goldrow{diag}{bp\_ipa}{H\_ig}{-257924.667357}
\goldrow{diag}{bp\_ipa}{H\_real}{-257924.667357}
\goldrow{diag}{bp\_ipa}{P}{101325}
\goldrow{diag}{bp\_ipa}{S\_R}{0}
\goldrow{diag}{bp\_ipa}{S\_ig}{329.037676662}
\goldrow{diag}{bp\_ipa}{S\_real}{329.037676662}
\goldrow{diag}{bp\_ipa}{T}{350}
\goldrow{diag}{bp\_ipa}{T\_bubble}{359.907520874}
\goldrow{diag}{bp\_ipa}{Z}{1}
\goldrow{diag}{bp\_ipa}{gamma}{1.09159801219}
\goldrow{diag}{bp\_ipa}{v\_molar}{0.0287200781278}
\goldrow{diag}{bp\_water}{Cp\_ig}{34.055209}
\goldmore{30}
\end{tutgolden}

\clearpage
\tutentry{acetone04\_acetone\_water\_vle}{tutorials/props/compare/acetone04\_acetone\_water\_vle}
\begin{tutscheme}
\begin{tikzpicture}[>={Stealth[length=2mm]}, every node/.style={font=\scriptsize}]
  \node[draw, rounded corners=2pt, align=center, minimum width=3.1cm, minimum height=0.8cm, fill=black!4] at (0.0,-0.0) {\textbf{bp\_acetone}\\[-1pt]{\tiny propertyPoint}};
  \node[draw, rounded corners=2pt, align=center, minimum width=3.1cm, minimum height=0.8cm, fill=black!4] at (3.6,-0.0) {\textbf{bp\_water}\\[-1pt]{\tiny propertyPoint}};
  \node[draw, rounded corners=2pt, align=center, minimum width=3.1cm, minimum height=0.8cm, fill=black!4] at (7.2,-0.0) {\textbf{dilute\_acetone}\\[-1pt]{\tiny propertyPoint}};
  \node[draw, rounded corners=2pt, align=center, minimum width=3.1cm, minimum height=0.8cm, fill=black!4] at (0.0,-1.3) {\textbf{txy\_top\_fine}\\[-1pt]{\tiny propertyScan1D}};
  \node[draw, rounded corners=2pt, align=center, minimum width=3.1cm, minimum height=0.8cm, fill=black!4] at (3.6,-1.3) {\textbf{txy\_unifac}\\[-1pt]{\tiny propertyScan1D}};
\end{tikzpicture}
\end{tutscheme}
\tutwhat{Acetone/water T-x-y at 1 atm on PREDICTIVE UNIFAC -- both components curated, so this isolates the mixture model}
\begin{tutnarrative}
MEASURE THE THERMODYNAMICS BEFORE BUILDING THE UNIT THAT RESTS ON IT.

Luyben's absorber recovers acetone from the reactor's hydrogen-rich offgas by washing it with water, and it is the next unit of the flowsheet to build (docs/design/acetone-ipa-reference-case.md). What that unit costs, and whether it can work at all, is set by ONE thing: how volatile acetone is out of a dilute aqueous solution. So that is measured here first -- the same order that made acetone01 useful before acetone03 was written.

WHAT MAKES THIS A CLEANER TEST THAN acetone01. Both components are CURATED: acetone and water are in data/standards/, neither comes from the estimate lake. acetone01 had to measure a mixture model and a lake-estimated component at once and separate them afterwards; here there is nothing to separate. Whatever deviation appears is UNIFAC, plus whatever the pure vapour pressures carry -- and the two pure boiling points below say how much that second part is.

WHAT LUYBEN GIVES, and it is less than for the IPA binary: acetone's normal boiling point 329.4 K, and the qualitative statement that acetone/water has NO azeotrope but a pinch at the high-acetone end. A qualitative claim is still falsifiable -- a model that invents an azeotrope here would be wrong in a way no AAD would show.

NOTE THE WATER LEG IS KNOWN TO BE 0.70 K LOW before it is measured again here: Choupo's curated water record puts the normal boiling point at 372.4536 K against the definition's 373.15 K, because its Antoine set is being used at the very top of its own declared window (see acetone01\_ipa\_water\_azeotrope/CLAUDE.md). It is repeated in this case rather than assumed, because a number quoted from another case is a number with two homes.

\end{tutnarrative}
\begin{tutgolden}
\goldrow{diag}{bp\_acetone}{Cp\_ig}{79.38622606}
\goldrow{diag}{bp\_acetone}{H\_R}{0}
\goldrow{diag}{bp\_acetone}{H\_ig}{-214666.297086}
\goldrow{diag}{bp\_acetone}{H\_real}{-214666.297086}
\goldrow{diag}{bp\_acetone}{P}{101325}
\goldrow{diag}{bp\_acetone}{S\_R}{0}
\goldrow{diag}{bp\_acetone}{S\_ig}{302.940865079}
\goldrow{diag}{bp\_acetone}{S\_real}{302.940865079}
\goldrow{diag}{bp\_acetone}{T}{330}
\goldrow{diag}{bp\_acetone}{T\_bubble}{329.417905068}
\goldrow{diag}{bp\_acetone}{Z}{1}
\goldrow{diag}{bp\_acetone}{gamma}{1.11698686251}
\goldrow{diag}{bp\_acetone}{v\_molar}{0.0270789308062}
\goldrow{diag}{bp\_water}{Cp\_ig}{34.217711312}
\goldmore{33}
\end{tutgolden}

\clearpage
\tutentry{compare\_activity\_etoh\_water}{tutorials/props/compare/compare\_activity\_etoh\_water}
\begin{tutscheme}
\begin{tikzpicture}[>={Stealth[length=2mm]}, every node/.style={font=\scriptsize}]
  \node[draw, rounded corners=2pt, align=center, minimum width=3.1cm, minimum height=0.8cm, fill=black!4] at (0.0,-0.0) {\textbf{g\_ideal}\\[-1pt]{\tiny propertyScan1D}};
  \node[draw, rounded corners=2pt, align=center, minimum width=3.1cm, minimum height=0.8cm, fill=black!4] at (3.6,-0.0) {\textbf{g\_nrtl}\\[-1pt]{\tiny propertyScan1D}};
  \node[draw, rounded corners=2pt, align=center, minimum width=3.1cm, minimum height=0.8cm, fill=black!4] at (7.2,-0.0) {\textbf{g\_wilson}\\[-1pt]{\tiny propertyScan1D}};
\end{tikzpicture}
\end{tutscheme}
\tutwhat{Compare 3 activity models (ideal / NRTL / Wilson) for $\gamma$\_ethanol in equimolar mixture with water vs T}
\begin{tutnarrative}
COMPARE 3 activity-coefficient models for $\gamma$\_ethanol in an equimolar ethanol/water mixture, T = 280 K to 360 K, P = 1 bar (as all three ops below declare; this line said 1 atm, which is 1.01325 bar).

ideal : $\gamma$ = 1 always (Raoult's law assumption). NRTL : strong non-ideality, $\gamma$\_ethanol > 1 (azeotrope behaviour). Wilson : similar to NRTL for this binary, slightly different curve.

Pedagogical point: a student studying ethanol/water VLE will see that "ideal" is a bad approximation here ($\gamma$\_ethanol reaches 1.267 under NRTL and 1.377 under Wilson at 280 K, the coldest and most non-ideal end of this scan, against exactly 1 for ideal), and that NRTL/Wilson give similar non-ideal corrections. The choice between the two is then decided by which fits the project's experimental data better (workflow Mode B, see tutorials/props/README.md).

Each op overrides the global activityModel block via `thermo \{\}`. The three CSVs (gamma\_ideal.csv, gamma\_nrtl.csv, gamma\_wilson.csv) are auto-overlaid by the GUI.

\end{tutnarrative}
\tutreadme{An equimolar ethanol/water liquid, held at 1 bar, scanned from 280 K to 360 K in 41 points, and asked for exactly one property: gamma\_ethanol. The same scan is run three times --- g\_ideal, g\_nrtl, g\_wilson --- and the only thing that changes between them is the activity model. There is no unit operation here on purpose: this is the thermodynamics that every unit in the rest of the path will be standing on.}
\tutreadme{Each op writes its own curve (gamma\_ideal.csv, gamma\_nrtl.csv, gamma\_wilson.csv), and the case's system/propsDict says the three are overlaid on one set of axes in the app. The golden pins the SHAPE of the result rather than the numbers: 41 points, 0 failures, 0 empty curves for each of the three ops.}
\tutreadme{1. Ideal is a model, and this case declares it. The case's constant/thermoPhysPropDict says activityModel ideal; with a header that spells out why --- $\ast$Raoult DECLARED, not defaulted$\ast$. The g\_ideal curve is flat at 1 across the whole 80 K span not because something was missing but because that is the answer the declared model gives. An absence and a choice would draw the same line; only the dict tells you which one you are looking at. 2. A thermo \{\} block overrides the models for one op only. All three ops declare the identical vary and state blocks; each carries its own thermo \{ equilibrium \{ liquid \{ activityModel ... \} \} \}. The components stay global --- it is the MODEL that is swapped, never the mixture --- which is the same per-unit override a flowsheet uses to run one unit in a different thermodynamic world. 3. The dict owns the parameters, in the open. NRTL is declared with a\_ij -0.8009, b\_ij 246.18, a\_ji 3.4578, b\_ji -586.0809, alpha 0.30; Wilson with A\_ij 1157.95 and A\_ji 4081.65 in J/mol. Two models, two different parameter shapes, both written where you can read them. Those four NRTL numbers are worth remembering: they reappear verbatim in flash02\_ethanol\_water and again in column03\_azeotrope\_mesh, which is how one thermodynamic decision travels through a curriculum. 4. A property scan is a case like any other. application choupoProps; in controlDict, the work declared as a list of operations in system/propsDict, the same dict grammar, the same .cho marker, the same golden file. Nothing about "just properties" is a lesser kind of run.}
\tutreadme{Change composition in all three ops to ethanol 0.05; water 0.95; and run again: the ideal curve does not move at all, and the other two do. Activity is a function of composition, and that is precisely the information the ideal model does not carry. Then open flash02\_ethanol\_water, where this same mixture and these same NRTL parameters arrive in a unit operation and change its answer.}
\begin{tutgolden}
\goldrow{diag}{g\_ideal}{n\_emptyCurves}{0}
\goldrow{diag}{g\_ideal}{n\_failures}{0}
\goldrow{diag}{g\_ideal}{n\_points}{41}
\goldrow{diag}{g\_nrtl}{n\_emptyCurves}{0}
\goldrow{diag}{g\_nrtl}{n\_failures}{0}
\goldrow{diag}{g\_nrtl}{n\_points}{41}
\goldrow{diag}{g\_wilson}{n\_emptyCurves}{0}
\goldrow{diag}{g\_wilson}{n\_failures}{0}
\goldrow{diag}{g\_wilson}{n\_points}{41}
\end{tutgolden}

\clearpage
\tutentry{compare\_eos\_co2}{tutorials/props/compare/compare\_eos\_co2}
\begin{tutscheme}
\begin{tikzpicture}[>={Stealth[length=2mm]}, every node/.style={font=\scriptsize}]
  \node[draw, rounded corners=2pt, align=center, minimum width=3.1cm, minimum height=0.8cm, fill=black!4] at (0.0,-0.0) {\textbf{z\_ideal}\\[-1pt]{\tiny propertyScan1D}};
  \node[draw, rounded corners=2pt, align=center, minimum width=3.1cm, minimum height=0.8cm, fill=black!4] at (3.6,-0.0) {\textbf{z\_srk}\\[-1pt]{\tiny propertyScan1D}};
  \node[draw, rounded corners=2pt, align=center, minimum width=3.1cm, minimum height=0.8cm, fill=black!4] at (7.2,-0.0) {\textbf{z\_pr}\\[-1pt]{\tiny propertyScan1D}};
\end{tikzpicture}
\end{tutscheme}
\tutwhat{Compare 3 EoS (idealGas / SRK / PR) for Z of CO2 near critical (T=310 K, P from 1 to 100 bar)}
\begin{tutnarrative}
COMPARE 3 equations of state for the same property at the same conditions. This is what choupoProps is $\ast$for$\ast$: helping a student or researcher CHOOSE the model that best describes the system at hand --- not just look at pretty Z curves of one EoS at a time.

System: pure CO2 at T = 310 K (just above its critical temperature Tc = 304.13 K), pressure swept from 1 bar to 100 bar. Near the critical pressure (Pc = 7.38 MPa) the compressibility plunges as the fluid turns liquid-like; the three curves read (at P = 73.7 bar, the swept point nearest Pc):

idealGas: Z = 1 always (textbook reference; ignores attraction). SRK : Z = 0.538 at Pc. PR : Z = 0.509 at Pc --- systematically lower than SRK (more attraction), and essentially on the experimental value. NIST : the $\ast$experimental$\ast$ Z near the critical point is \textasciitilde{}0.49.

The lesson: at and just above Tc, BOTH cubics track CO2 well --- PR sits almost exactly on the \textasciitilde{}0.49 datum and SRK is only \textasciitilde{}0.05 high. This is a regime where cubic EoS are at their best, not their worst. Both Z curves keep falling past Pc to a shallow minimum (SRK \textasciitilde{}0.30, PR \textasciitilde{}0.275 near 98 bar) before rising again; the SRK-vs-PR gap is the diagnostic a student should read off the overlay.

Each operation overrides the global thermoPhysPropDict with a different EoS via the `thermo \{\}` per-op block. The three CSVs are independent, so the GUI overlays them on the same Plotly canvas (same X axis "P", same Y "Z"). The result is one figure with three curves, ready for the project report.

\end{tutnarrative}
\begin{tutgolden}
\goldrow{diag}{z\_ideal}{n\_emptyCurves}{0}
\goldrow{diag}{z\_ideal}{n\_failures}{0}
\goldrow{diag}{z\_ideal}{n\_points}{50}
\goldrow{diag}{z\_srk}{n\_emptyCurves}{0}
\goldrow{diag}{z\_srk}{n\_failures}{0}
\goldrow{diag}{z\_srk}{n\_points}{50}
\goldrow{diag}{z\_pr}{n\_emptyCurves}{0}
\goldrow{diag}{z\_pr}{n\_failures}{0}
\goldrow{diag}{z\_pr}{n\_points}{50}
\end{tutgolden}

\clearpage
\tutentry{compare\_kinetics\_order}{tutorials/props/compare/compare\_kinetics\_order}
\begin{tutscheme}
\begin{tikzpicture}[>={Stealth[length=2mm]}, every node/.style={font=\scriptsize}]
  \node[draw, rounded corners=2pt, align=center, minimum width=3.1cm, minimum height=0.8cm, fill=black!4] at (0.0,-0.0) {\textbf{rate\_order1}\\[-1pt]{\tiny kinetics1D}};
  \node[draw, rounded corners=2pt, align=center, minimum width=3.1cm, minimum height=0.8cm, fill=black!4] at (3.6,-0.0) {\textbf{rate\_order2}\\[-1pt]{\tiny kinetics1D}};
\end{tikzpicture}
\end{tutscheme}
\tutwhat{Rate-law order from DATA: 1st- vs 2nd-order fits overlaid on measured c(t) --- the data picks the order}
\begin{tutnarrative}
The "see, then decide" of CHEMICAL KINETICS: fit a 1st- and a 2nd-order rate law to the SAME measured c(t) and overlay both on the data. No `k` is given, so each op REGRESSES its own k from the points (provenance = fitted). The data here are second-order -- the 1st-order curve misses the tail, the 2nd lands on the points. The student picks the order from the evidence.

\end{tutnarrative}
\begin{tutgolden}
\goldrow{diag}{rate\_order1}{R2}{0.973397}
\goldrow{diag}{rate\_order1}{chi2}{0.0105136}
\goldrow{diag}{rate\_order1}{fitted}{1}
\goldrow{diag}{rate\_order1}{k}{0.00065871}
\goldrow{diag}{rate\_order1}{max\_abs\_resid}{0.0544611}
\goldrow{diag}{rate\_order1}{n\_data}{8}
\goldrow{diag}{rate\_order1}{n\_points}{61}
\goldrow{diag}{rate\_order1}{order}{1}
\goldrow{diag}{rate\_order1}{rms}{0.0362518}
\goldrow{diag}{rate\_order2}{R2}{0.999421}
\goldrow{diag}{rate\_order2}{chi2}{0.000228631}
\goldrow{diag}{rate\_order2}{fitted}{1}
\goldrow{diag}{rate\_order2}{k}{0.000992909}
\goldrow{diag}{rate\_order2}{max\_abs\_resid}{0.00829088}
\goldmore{8}
\end{tutgolden}

\clearpage
\tutentry{compare\_psat\_ambrose\_walton}{tutorials/props/compare/compare\_psat\_ambrose\_walton}
\begin{tutscheme}
\begin{tikzpicture}[>={Stealth[length=2mm]}, every node/.style={font=\scriptsize}]
  \node[draw, rounded corners=2pt, align=center, minimum width=3.1cm, minimum height=0.8cm, fill=black!4] at (0.0,-0.0) {\textbf{psat\_aw\_vs\_antoine}\\[-1pt]{\tiny propertyScan1D}};
\end{tikzpicture}
\end{tutscheme}
\tutwhat{Psat: Ambrose-Walton corresponding-states PREDICTION vs fitted Antoine DATA, for benzene (non-polar, good) and ethanol (associating, degrades) --- see the error, then decide}
\begin{tutnarrative}
Ambrose-Walton corresponding states (Ambrose \& Walton 1989; Poling et al. 5th ed., Eq. 7-2.7) predicts vapour pressure from Tc, Pc and omega ALONE -- exactly the constants a Joback + Lee-Kesler estimate produces. So an estimated component becomes flashable with NO measured Psat data. But it is an ESTIMATE, and the glass-box rule is: SEE the error before you trust it.

This sweeps T = 270 $\to$ 360 K and tabulates four curves: Psat\_benzene -- fitted Antoine (catalogue DATA) \} non-polar pair Psat\_benzeneAW -- Ambrose-Walton PREDICTION \} Psat\_ethanol -- fitted Antoine (catalogue DATA) \} associating pair Psat\_ethanolAW -- Ambrose-Walton PREDICTION \}

Open `psat\_aw\_vs\_antoine.csv` (GUI plot, or gnuplot). What you SEE: Near the normal boiling point (\textasciitilde{}350-360 K) BOTH agree to \textasciitilde{}1 \% -- corresponding states is anchored there and is excellent. Towards the cold end (270 K) the curves separate, and MORE for the associating ethanol than for non-polar benzene: benzene AW vs Antoine: +3.3 \% average, +8.9 \% at 270 K ethanol AW vs Antoine: +5.2 \% average, -16 \% at 270 K So corresponding states is a solid FREE first predictor (it makes an estimated component flashable with no data), best near Tb and degrading at low reduced temperature -- the degradation worse for hydrogen-bonding species. The lesson: prediction gets you started; for an associating species, or far below Tb, fit Antoine to data (choupoProps vaporPressureFit) before any design number.

\end{tutnarrative}
\begin{tutgolden}
\goldrow{diag}{psat\_aw\_vs\_antoine}{n\_emptyCurves}{0}
\goldrow{diag}{psat\_aw\_vs\_antoine}{n\_failures}{0}
\goldrow{diag}{psat\_aw\_vs\_antoine}{n\_points}{46}
\end{tutgolden}

\clearpage
\tutentry{compare\_transport\_water}{tutorials/props/compare/compare\_transport\_water}
\begin{tutscheme}
\begin{tikzpicture}[>={Stealth[length=2mm]}, every node/.style={font=\scriptsize}]
  \node[draw, rounded corners=2pt, align=center, minimum width=3.1cm, minimum height=0.8cm, fill=black!4] at (0.0,-0.0) {\textbf{mu\_iapws}\\[-1pt]{\tiny propertyScan1D}};
  \node[draw, rounded corners=2pt, align=center, minimum width=3.1cm, minimum height=0.8cm, fill=black!4] at (3.6,-0.0) {\textbf{mu\_generic}\\[-1pt]{\tiny propertyScan1D}};
  \node[draw, rounded corners=2pt, align=center, minimum width=3.1cm, minimum height=0.8cm, fill=black!4] at (7.2,-0.0) {\textbf{k\_iapws}\\[-1pt]{\tiny propertyScan1D}};
  \node[draw, rounded corners=2pt, align=center, minimum width=3.1cm, minimum height=0.8cm, fill=black!4] at (0.0,-1.3) {\textbf{k\_generic}\\[-1pt]{\tiny propertyScan1D}};
\end{tikzpicture}
\end{tutscheme}
\tutwhat{Compare water transport: IAPWS (IF97 R12-08/R15-11) vs generic correlations (Andrade / SatoRiedel), 280-370 K}
\begin{tutnarrative}
COMPARE water transport properties computed TWO WAYS on the SAME water, over the saturated-liquid window T = 280-370 K:

(a) IAPWS reference -- water flagged `pureFluids \{ water \{ method IF97; \} \}` so the transport accessors route to the matching IAPWS releases: viscosity $\to$ R12-08 (industrial) thermal conductivity $\to$ R15-11 (background term) evaluated at the saturation state (the legacy accessors carry no P, so p = p\_sat(T) -- saturated liquid). These are THE international standard reference values (verified to 1e-8 by IAPWSTransport::verify(), which fires on every run).

(b) generic correlations -- the same property KEY, but with a generic transport block and NO pureFluids flag, so the accessors fall to: viscosity $\to$ Andrade ln(mu) = A + B/T thermal conductivity $\to$ SatoRiedel fitted from the per-component .dat (Tc/Tb + the Andrade pair).

Because BOTH ops emit the same property key (`viscosity\_liquid`, `thermal\_conductivity\_liquid`) the two curves overlay directly -- the student SEES where the cheap generic correlation deviates from the IAPWS reference.

THE LESSON (the actual computed gaps): VISCOSITY -- Andrade tracks IAPWS within a few percent over the window (\textasciitilde{}7 \% at 280 K, \textasciitilde{}2-4 \% mid-range, \textasciitilde{}3 \% at 370 K). The 2-parameter ln(mu)=A+B/T cannot follow the full curvature, but a few \% is fine for a Reynolds number; it is NOT good enough for a precision steam duty. CONDUCTIVITY -- SatoRiedel is a corresponding-states estimate and it FAILS BADLY for water: it underpredicts IAPWS by \textasciitilde{}50-60 \% across the window. Water is the textbook SatoRiedel failure (strong hydrogen bonding, anomalous k); this is exactly why "SEE before you commit" matters -- a blind generic pick would have halved the heat-transfer coefficient. For water conductivity use IAPWS R15-11, full stop. Generic correlations are CHEAP and GENERAL (any compound with the handful of .dat parameters); IAPWS is EXACT but water-only. For a water-dominated process use the IF97 + IAPWS route; for a screening sweep over many fluids the generic correlations are the right tool. SEE the gap, then decide -- do not pick blindly.

The GUI renders the CSVs via the existing compare machinery (overlay of the $\ast$\_iapws / $\ast$\_generic curves vs T) -- no new Explorer kind.

\end{tutnarrative}
\begin{tutgolden}
\goldrow{diag}{mu\_iapws}{n\_failures}{0}
\goldrow{diag}{mu\_iapws}{n\_points}{46}
\goldrow{diag}{mu\_generic}{n\_failures}{0}
\goldrow{diag}{mu\_generic}{n\_points}{46}
\goldrow{diag}{k\_iapws}{n\_failures}{0}
\goldrow{diag}{k\_iapws}{n\_points}{46}
\goldrow{diag}{k\_generic}{n\_failures}{0}
\goldrow{diag}{k\_generic}{n\_points}{46}
\goldrow{diag}{mu\_iapws}{n\_emptyCurves}{0}
\goldrow{diag}{mu\_generic}{n\_emptyCurves}{0}
\goldrow{diag}{k\_iapws}{n\_emptyCurves}{0}
\goldrow{diag}{k\_generic}{n\_emptyCurves}{0}
\end{tutgolden}

\clearpage
\tutentry{compare\_visc\_liquid\_ammonia}{tutorials/props/compare/compare\_visc\_liquid\_ammonia}
\begin{tutscheme}
\begin{tikzpicture}[>={Stealth[length=2mm]}, every node/.style={font=\scriptsize}]
  \node[draw, rounded corners=2pt, align=center, minimum width=3.1cm, minimum height=0.8cm, fill=black!4] at (0.0,-0.0) {\textbf{mu\_andrade}\\[-1pt]{\tiny propertyScan1D}};
\end{tikzpicture}
\end{tutscheme}
\tutwhat{Liquid-viscosity Andrade for ammonia, 235-300 K (cross-check mu@239.85K \textasciitilde{} 0.255 mPa.s)}
\begin{tutnarrative}
Liquid-viscosity Andrade correlation for pure ammonia (NH3), 235-300 K. This case is SEALED: it runs its own copy, constant/components/NH3.dat, whose Andrade pair was fitted to the Engineering Toolbox table -- a no-grant source (CLAUDE.md 10) -- and, measured against the reference correlation (Monogenidou, Assael \& Huber, J. Phys. Chem. Ref. Data 47, 023102 (2018), Table 6), is +10 \% at 240 K and +65 \% at 300 K. Since 2026-10-11 (DEV.md 4c C65 item c) the CATALOGUE record data/standards/components/NH3.dat carries a fit to that Table 6 (AAD 0.51 \% over 200-320 K; 0.253 mPa.s at the normal boiling point, against Table 6's 0.2510 at 240 K). This copy keeps the old pair until the case is re-imported, and check\_source\_licence pins it until then; the golden pins only the scan's point counts, so the re-import moves no golden row. Note this case scans a SINGLE method (Andrade) -- it is a single-correlation curve, not a multi-method comparison. Ammonia is liquid only under pressure in this T window; viscosity\_liquid is a pure-T correlation, so the state P is immaterial to the value.

\end{tutnarrative}
\begin{tutgolden}
\goldrow{diag}{mu\_andrade}{n\_emptyCurves}{0}
\goldrow{diag}{mu\_andrade}{n\_failures}{0}
\goldrow{diag}{mu\_andrade}{n\_points}{27}
\end{tutgolden}

\clearpage
\tutentry{compare\_visc\_liquid\_water}{tutorials/props/compare/compare\_visc\_liquid\_water}
\begin{tutscheme}
\begin{tikzpicture}[>={Stealth[length=2mm]}, every node/.style={font=\scriptsize}]
  \node[draw, rounded corners=2pt, align=center, minimum width=3.1cm, minimum height=0.8cm, fill=black!4] at (0.0,-0.0) {\textbf{mu\_andrade}\\[-1pt]{\tiny propertyScan1D}};
  \node[draw, rounded corners=2pt, align=center, minimum width=3.1cm, minimum height=0.8cm, fill=black!4] at (3.6,-0.0) {\textbf{mu\_vogel}\\[-1pt]{\tiny propertyScan1D}};
\end{tikzpicture}
\end{tutscheme}
\tutwhat{Compare 2 liquid viscosity correlations (Andrade vs Vogel) for water, 280-360 K}
\begin{tutnarrative}
COMPARE 2 liquid-viscosity correlations for pure water, T = 280-360 K.

Andrade: ln($\mu$) = A + B/T --- 2-parameter, simplest fit, OK over a narrow T window. Vogel : ln($\mu$) = A + B/(T + C) --- 3-parameter, captures the curvature better, recommended for wider T ranges.

Pedagogical point: for a project using water transport in a non-trivial T window (e.g. 25-90 $^\circ$C for a heated reactor), the student should know WHICH correlation he is committing to. Both correlations give viscosity FALLING with temperature (the physical trend): \textasciitilde{}1.3-1.4 mPa$\cdot$s near 280 K down to \textasciitilde{}0.32 mPa$\cdot$s at 360 K. This case shows the two models disagree MOST at the cold end (\textasciitilde{}6 \% at 280 K), are within \textasciitilde{}2 \% at room T, and agree BEST at the hot end (\textasciitilde{}1 \% at 360 K) --- the 2-parameter Andrade fit drifts where it is extrapolated coolest. Validate against the textbook value: $\mu$\_water @ 298 K $\approx$ 8.9 $\times$ 10?? Pa$\cdot$s --- Vogel reads 8.93 $\times$ 10?? (essentially on it), Andrade 9.09 $\times$ 10?? (\textasciitilde{}2 \% high).

\end{tutnarrative}
\begin{tutgolden}
\goldrow{diag}{mu\_andrade}{n\_emptyCurves}{0}
\goldrow{diag}{mu\_andrade}{n\_failures}{0}
\goldrow{diag}{mu\_andrade}{n\_points}{41}
\goldrow{diag}{mu\_vogel}{n\_emptyCurves}{0}
\goldrow{diag}{mu\_vogel}{n\_failures}{0}
\goldrow{diag}{mu\_vogel}{n\_points}{41}
\end{tutgolden}

\clearpage
\tutentry{compare\_vle\_etoh\_water}{tutorials/props/compare/compare\_vle\_etoh\_water}
\begin{tutscheme}
\begin{tikzpicture}[>={Stealth[length=2mm]}, every node/.style={font=\scriptsize}]
  \node[draw, rounded corners=2pt, align=center, minimum width=3.1cm, minimum height=0.8cm, fill=black!4] at (0.0,-0.0) {\textbf{txy\_ideal}\\[-1pt]{\tiny propertyScan1D}};
  \node[draw, rounded corners=2pt, align=center, minimum width=3.1cm, minimum height=0.8cm, fill=black!4] at (3.6,-0.0) {\textbf{txy\_nrtl}\\[-1pt]{\tiny propertyScan1D}};
  \node[draw, rounded corners=2pt, align=center, minimum width=3.1cm, minimum height=0.8cm, fill=black!4] at (7.2,-0.0) {\textbf{txy\_unifac}\\[-1pt]{\tiny propertyScan1D}};
  \node[draw, rounded corners=2pt, align=center, minimum width=3.1cm, minimum height=0.8cm, fill=black!4] at (0.0,-1.3) {\textbf{consistency\_etoh\_water}\\[-1pt]{\tiny vleConsistency}};
  \node[draw, rounded corners=2pt, align=center, minimum width=3.1cm, minimum height=0.8cm, fill=black!4] at (3.6,-1.3) {\textbf{txy\_wilson}\\[-1pt]{\tiny propertyScan1D}};
\end{tikzpicture}
\end{tutscheme}
\tutwhat{Ethanol/water T-x-y at 1 atm: ideal vs NRTL model curves OVERLAID with measured VLE data --- see the data, pick the model}
\begin{tutnarrative}
The canonical "SEE, then DECIDE" props case: ethanol/water T-x-y at 1 atm, with two MODEL curves and the MEASURED data on ONE plot.

Piece 1 -- two scan operations, identical but for the activity model: txy\_ideal Raoult (gamma = 1) $\to$ a smooth lens, NO azeotrope txy\_nrtl NRTL $\to$ the real minimum-boiling azeotrope Each emits the (x, T\_bubble, y\_eq) trio $\to$ the T-x-y / x-y / T-x / T-y toggle (x-y is the McCabe-Thiele diagram).

Piece 2 -- the `experimental \{\}` overlay: lists the two ops as `models`, and points `dataset` at REAL measured points (constant/experiments/...). The GUI draws both model curves + the lab points together; the student SEES that the ideal model misses the azeotrope and NRTL captures it, and picks NRTL from evidence -- not from a recommended-by badge.

\end{tutnarrative}
\begin{tutgolden}
\goldrow{diag}{txy\_ideal}{n\_emptyCurves}{0}
\goldrow{diag}{txy\_ideal}{n\_failures}{0}
\goldrow{diag}{txy\_ideal}{n\_points}{51}
\goldrow{diag}{txy\_nrtl}{n\_emptyCurves}{0}
\goldrow{diag}{txy\_nrtl}{n\_failures}{0}
\goldrow{diag}{txy\_nrtl}{n\_points}{51}
\goldrow{diag}{txy\_unifac}{n\_emptyCurves}{0}
\goldrow{diag}{txy\_unifac}{n\_failures}{0}
\goldrow{diag}{txy\_unifac}{n\_points}{51}
\goldrow{diag}{consistency\_etoh\_water}{gamma1\_inf}{4.68506}
\goldrow{diag}{consistency\_etoh\_water}{gamma2\_inf}{2.27627}
\goldrow{diag}{consistency\_etoh\_water}{gd\_max\_resid}{0.0964023}
\goldrow{diag}{consistency\_etoh\_water}{gd\_mean\_resid}{0.05208}
\goldrow{diag}{consistency\_etoh\_water}{herington\_D}{12.0482}
\goldmore{14}
\end{tutgolden}

\clearpage
\tutentry{evidence01\_methanol\_water\_points}{tutorials/props/compare/evidence01\_methanol\_water\_points}
\begin{tutscheme}
\begin{tikzpicture}[>={Stealth[length=2mm]}, every node/.style={font=\scriptsize}]
  \node[draw, rounded corners=2pt, align=center, minimum width=3.1cm, minimum height=0.8cm, fill=black!4] at (0.0,-0.0) {\textbf{txy\_nrtl\_points}\\[-1pt]{\tiny propertyScan1D}};
  \node[draw, rounded corners=2pt, align=center, minimum width=3.1cm, minimum height=0.8cm, fill=black!4] at (3.6,-0.0) {\textbf{txy\_unifac\_points}\\[-1pt]{\tiny propertyScan1D}};
  \node[draw, rounded corners=2pt, align=center, minimum width=3.1cm, minimum height=0.8cm, fill=black!4] at (7.2,-0.0) {\textbf{txy\_ideal\_points}\\[-1pt]{\tiny propertyScan1D}};
\end{tikzpicture}
\end{tutscheme}
\tutwhat{Methanol/water bubble T evaluated AT the 13 measured x of a held-out study (doi 10.1021/je200341c): NRTL, UNIFAC and ideal, residual = engine minus measurement at the measured point}
\begin{tutnarrative}
A MODEL EVALUATED AT THE MEASURED CONDITIONS, NOT ON A GRID (DEV.md 4c C117).

THE QUESTION. A measured dataset is a list of CONDITIONS somebody set and the value they read there. To compare a model with it, the honest residual is model(x\_measured) - measurement -- the engine asked at the measured x, never a value read off a smooth curve between two grid neighbours. The `experimental \{\}` overlay's AAD interpolates the model CSV at the measured x; here the CSV HOLDS the measured x, so that interpolation is exact.

THE GRAMMAR. `vary \{ variable x[methanol]; points ( ... ); \}` -- the 13 interior liquid mole fractions of the held-out study, in its own order, instead of `from / to / n`. Both in one `vary \{\}` is refused by name. The values are canonical SI (a mole fraction here; K or Pa for T or P).

THE EVIDENCE. constant/experiments/methanol-water-bubble-je200341c.dat: Yang, Ma \& Yin, J. Chem. Eng. Data (2011), doi 10.1021/je200341c, 13 interior points at 101.32 kPa, the HELD-OUT set of the curate04 regression (tutorials/props/curation/curate04\_nrtl\_methanol\_water), kept here as this case's own copy -- a self-contained tutorial keeps its own copy. Its `reviewStatus transcribedNotCheckedAgainstArticle` is announced by the run.

THE THREE CANDIDATES. The catalogue NRTL pair (regressed by Choupo in curate04 on ANOTHER study, so this set is a validation of it), UNIFAC (predictive, from the components' groups) and ideal mixing. What the curate04 case found holds here to the digit: UNIFAC is closer than the fitted pair on this set, and the fitted pair is `validated` -- inside its declared band, never better than every alternative.

WHAT IS NOT CLAIMED. Nothing is fitted here. Part of every residual is the catalogue water Antoine record's own 0.70 K at 101.32 kPa, which no pair can remove. This is the engine witness of the Property Explorer's Evidence view, which synthesises exactly these operations (one choupoProps run per candidate) and only subtracts.

\end{tutnarrative}
\begin{tutgolden}
\goldrow{diag}{txy\_nrtl\_points}{n\_emptyCurves}{0}
\goldrow{diag}{txy\_nrtl\_points}{n\_failures}{0}
\goldrow{diag}{txy\_nrtl\_points}{n\_points}{13}
\goldrow{diag}{txy\_unifac\_points}{n\_emptyCurves}{0}
\goldrow{diag}{txy\_unifac\_points}{n\_failures}{0}
\goldrow{diag}{txy\_unifac\_points}{n\_points}{13}
\goldrow{diag}{txy\_ideal\_points}{n\_emptyCurves}{0}
\goldrow{diag}{txy\_ideal\_points}{n\_failures}{0}
\goldrow{diag}{txy\_ideal\_points}{n\_points}{13}
\goldrow{aad}{heldout\_je200341c}{txy\_nrtl\_points.T\_bubble.aadAbs}{0.534750769231}
\goldrow{aad}{heldout\_je200341c}{txy\_nrtl\_points.T\_bubble.aadRelPct}{0.150061454094}
\goldrow{aad}{heldout\_je200341c}{txy\_nrtl\_points.y\_eq\_methanol.aadAbs}{0.0125426852538}
\goldrow{aad}{heldout\_je200341c}{txy\_unifac\_points.T\_bubble.aadAbs}{0.366152846154}
\goldrow{aad}{heldout\_je200341c}{txy\_unifac\_points.T\_bubble.aadRelPct}{0.10093737486}
\goldmore{4}
\end{tutgolden}

\clearpage
\tutentry{extractionShortcut01\_perry\_example15\_1}{tutorials/props/compare/extractionShortcut01\_perry\_example15\_1}
\begin{tutscheme}
\begin{tikzpicture}[>={Stealth[length=2mm]}, every node/.style={font=\scriptsize}]
  \node[draw, rounded corners=2pt, align=center, minimum width=3.1cm, minimum height=0.8cm, fill=black!4] at (0.0,-0.0) {\textbf{ex15\_1\_caseA}\\[-1pt]{\tiny extractionShortcut}};
  \node[draw, rounded corners=2pt, align=center, minimum width=3.1cm, minimum height=0.8cm, fill=black!4] at (3.6,-0.0) {\textbf{ex15\_1\_E\_from\_slopes}\\[-1pt]{\tiny extractionShortcut}};
  \node[draw, rounded corners=2pt, align=center, minimum width=3.1cm, minimum height=0.8cm, fill=black!4] at (7.2,-0.0) {\textbf{ex15\_2\_caseB}\\[-1pt]{\tiny extractionShortcut}};
\end{tikzpicture}
\end{tutscheme}
\tutwhat{Extraction shortcut (Perry Eq. 15-43) on Examples 15-1 to 15-3}
\begin{tutnarrative}
WHAT IT IS. The Kremser-Souders-Brown shortcut for a dilute liquid-liquid extraction cascade, run on the numbers of two worked examples in Perry's Chemical Engineers' Handbook, 9th ed., Section 15, p. 15-48:

Example 15-1 (case A, immiscible solvents). 100 kg/h of 20 wt\% acetic acid in water, extracted by 200 kg/h of recycle MIBK, raffinate down to 1 wt\% acid. Perry reads the ratios X' = 0.25, 0.01 (feed, raffinate) and Y' = 0.001 (solvent), the equilibrium slope K' = 0.656 at the dilute end, and the extraction factor E = sqrt(m1' mr') S'/F' = 0.739 x 199.8/80 = 1.85, and prints N = 4.3 theoretical stages from Eq. (15-43). Example 15-2 (case B, partly miscible solvents). The same problem with the mutual solubilities of water and MIBK accounted for by the pseudo-solute concentrations of Eqs. (15-90) and (15-91): X\_f\textasciicircum{}B = 0.251, Y\_s\textasciicircum{}B = 0.0016, N = 4.5. Example 15-3. The transfer units of Example 15-2: N\_or = 6.0.

WHAT IT IS NOT. No property package is read -- the op evaluates a formula on declared numbers, and the package in constant/ is only what the props driver requires of a case. The answers are held to the PRINTED numbers, to the digits Perry prints (the `anchor` rows of `expected`); that is a check of Choupo's transcription of the equations, not a validation of the shortcut against an extractor. Treybal, Mass-Transfer Operations 3rd ed., Eq. (10.38) is the same cascade written as a fraction taken, and the `fractionTaken` / `fractionRequired` pair printed beside N is the agreement of the two books' forms.

The `extractor` unit prints the same shortcut beside its rigorous stage count (KPI N\_shortcut), as a CHECK on a straight equilibrium line, never a design.

Source: Perry's Chemical Engineers' Handbook, 9th ed., Examples 15-1 to 15-3, read from the page image (p. 15-48); Treybal 3rd ed., Eq. (10.38).

\end{tutnarrative}
\begin{tutgolden}
\goldrow{diag}{ex15\_1\_caseA}{E}{1.85}
\goldrow{diag}{ex15\_1\_caseA}{N}{4.3}
\goldrow{diag}{ex15\_1\_caseA}{N\_or}{5.74543214891}
\goldrow{diag}{ex15\_1\_caseA}{fractionRequired}{0.965889570552}
\goldrow{diag}{ex15\_1\_caseA}{fractionTaken}{0.965889570552}
\goldrow{diag}{ex15\_1\_E\_from\_slopes}{E}{1.85}
\goldrow{diag}{ex15\_1\_E\_from\_slopes}{N}{4.3}
\goldrow{diag}{ex15\_1\_E\_from\_slopes}{N\_or}{5.75445508328}
\goldrow{diag}{ex15\_1\_E\_from\_slopes}{fractionRequired}{0.965889570552}
\goldrow{diag}{ex15\_1\_E\_from\_slopes}{fractionTaken}{0.965889570552}
\goldrow{diag}{ex15\_2\_caseB}{E}{1.85}
\goldrow{diag}{ex15\_2\_caseB}{N}{4.5}
\goldrow{diag}{ex15\_2\_caseB}{N\_caseB}{4.48053341968}
\goldrow{diag}{ex15\_2\_caseB}{N\_or}{6.0}
\goldmore{4}
\end{tutgolden}

\clearpage
\tutentry{lle01\_acetic\_mibk\_perry}{tutorials/props/compare/lle01\_acetic\_mibk\_perry}
\begin{tutscheme}
\begin{tikzpicture}[>={Stealth[length=2mm]}, every node/.style={font=\scriptsize}]
  \node[draw, rounded corners=2pt, align=center, minimum width=3.1cm, minimum height=0.8cm, fill=black!4] at (0.0,-0.0) {\textbf{tie\_unifac}\\[-1pt]{\tiny lleTieLines}};
\end{tikzpicture}
\end{tutscheme}
\tutwhat{LLE validation: water + acetic acid + MIBK, 25 C, UNIFAC against Perry Table 15-8}
\begin{tutnarrative}
WHAT IT IS. The first VALIDATION READING of the liquid-liquid flash that the `extractor` runs at every stage: the catalogue's UNIFAC, with no parameter touched, laid beside the six measured tie lines of water + acetic acid + methyl isobutyl ketone (MIBK) at 25 C, 1 atm (constant/experiments: Sherwood, Evans \& Longcor 1939, as printed in Perry's Chemical Engineers' Handbook 9th ed., Table 15-8, p. 15-48).

HOW IT READS. `lleTieLines` feeds the flash along a straight line in the triangle (acid mass fraction 0.005 $\to$ 0.40, the rest water : MIBK 1 : 1), and writes the conjugate pair (raffinate, extract) of every feed that splits, in MASS fractions. The overlay interpolates the model's extract composition at each MEASURED raffinate acid fraction (the abscissa) and reports the average absolute deviation per column. Nothing is fitted: this is a reading, pinned by `aad` rows, not a calibration.

WHAT IT FOUND (mass fractions, AAD over the six tie lines):

raffinate water 0.036 extract acetic acid 0.031 raffinate MIBK 0.036 extract water 0.043 extract MIBK 0.056

UNIFAC DOES split the system and has the right shape of the binodal, but it MISSES the measurements by 3 to 6 mass percent -- a miss, and the verdict word the gate holds is `misses` (band declared here, a priori and for every column: AAD <= 0.01 `agrees`, <= 0.03 `approximate`, above `misses`). The misses have a direction: the model puts too much acetic acid in the MIBK (extract acid 0.130 against the measured 0.089 at the raffinate's 0.117, and 0.222 against 0.173 at 0.205), dissolves about twice the measured water in the dilute extract (0.088 against 0.054 at 0.117), and keeps the two liquids apart up to a raffinate acid fraction of about 0.395, past the last measured tie line (0.346, where the measured extract is still 0.31 water). This is the corpus UNIFAC, the VLE parameter table (CLAUDE.md sec. 6: the same table keeps water / 1-butanol miscible); a curated LLE parameter set, or NRTL pairs fitted on a held-out design, is the remedy and is NOT done here. Consequence for a student: a countercurrent cascade on this package (tutorials/steady/absorption/extract03\_acetic\_acid\_mibk\_perry) takes more acid out per stage than the real system does.

NOT IN THE READING. The all-binary row of Table 15-8 (water / MIBK mutual solubility) is not carried; the data are transcribed from the table's page image and NOT checked against the 1939 article; the MIBK vapour pressure advisory in the log is about a property the LLE never uses.

\end{tutnarrative}
\begin{tutgolden}
\goldrow{diag}{tie\_unifac}{n\_feeds}{80}
\goldrow{diag}{tie\_unifac}{n\_refused}{0}
\goldrow{diag}{tie\_unifac}{n\_single}{0}
\goldrow{diag}{tie\_unifac}{n\_split}{80}
\goldrow{diag}{tie\_unifac}{wR\_solute\_max}{0.395303060021}
\goldrow{aad}{perry\_table15\_8}{tie\_unifac.wR\_water.aadAbs}{0.0362145055063}
\goldrow{aad}{perry\_table15\_8}{tie\_unifac.wR\_methylIsobutylKetone.aadAbs}{0.0361978386946}
\goldrow{aad}{perry\_table15\_8}{tie\_unifac.wE\_aceticAcid.aadAbs}{0.0310105107019}
\goldrow{aad}{perry\_table15\_8}{tie\_unifac.wE\_water.aadAbs}{0.0432948894869}
\goldrow{aad}{perry\_table15\_8}{tie\_unifac.wE\_methylIsobutylKetone.aadAbs}{0.0561117991127}
\end{tutgolden}

\subsection{curation}
\grouptheory{curation}
\clearpage
\tutentry{curate01\_water\_evidence\_partition}{tutorials/props/curation/curate01\_water\_evidence\_partition}
\begin{tutscheme}
\begin{tikzpicture}[>={Stealth[length=2mm]}, every node/.style={font=\scriptsize}]
  \node[draw, rounded corners=2pt, align=center, minimum width=3.1cm, minimum height=0.8cm, fill=black!4] at (0.0,-0.0) {\textbf{psat\_water\_validated}\\[-1pt]{\tiny vaporPressureFit}};
  \node[draw, rounded corners=2pt, align=center, minimum width=3.1cm, minimum height=0.8cm, fill=black!4] at (3.6,-0.0) {\textbf{psat\_water\_refused}\\[-1pt]{\tiny vaporPressureFit}};
\end{tikzpicture}
\end{tutscheme}
\tutwhat{The FIT/HELD-OUT partition contract: the same property fitted twice, differing only in whether evidence was held out -- one VALIDATED, one VALIDATION REFUSED}
\begin{tutnarrative}
THE FIT / HELD-OUT VALIDATION CONTRACT, and its refusal, side by side.

TWO OPERATIONS. SAME COMPONENT. SAME PROPERTY. SAME TWO DATASETS. The ONLY thing that differs is whether any evidence was held out -- which is exactly the variable the contract governs, and isolating it is why this witness does not use two different properties.

psat\_water\_validated lowT `role fit` + highT `role validation` psat\_water\_refused lowT `role fit` + highT `role fit`

The first can report a held-out residual because the fitter never received the 325-370 K points. The second cannot, and says so by name instead of presenting its in-sample R2 as though it were an external check:

VALIDATION REFUSED: No independent experimental evidence remains after fitting.

THE ANCHOR IS ENFORCED BY CONSTRUCTION. The partition is parsed and frozen before any fitting code runs, and the fitter is handed only the `role fit` points. It is not forbidden to move a held-out point into the fit -- it is never given one. A rule can be broken by the next edit; a data-flow fact cannot. `role` has NO default, because a defaulted role would make every dataset a fitting dataset and bring the in-sample problem back invisibly.

WHAT THIS CASE DOES NOT SHOW. Both datasets are generated from the SAME curated Antoine correlation (see either dataset header), so the held-out residual is arithmetic rather than evidence about the physics. This is a STRUCTURAL witness -- no number in it is validated against measurement.

Contract: src/propertyOps/EvidencePartition.H

\end{tutnarrative}
\begin{tutgolden}
\goldrow{diag}{psat\_water\_validated}{A}{5.40221}
\goldrow{diag}{psat\_water\_validated}{B}{1838.68}
\goldrow{diag}{psat\_water\_validated}{C}{-31.737}
\goldrow{diag}{psat\_water\_validated}{R2}{1}
\goldrow{diag}{psat\_water\_validated}{aad\_heldout\_pct}{5.95532e-08}
\goldrow{diag}{psat\_water\_validated}{dHvap\_kJ\_per\_mol}{43.3787}
\goldrow{diag}{psat\_water\_validated}{n\_heldout}{10}
\goldrow{diag}{psat\_water\_validated}{n\_points}{10}
\goldrow{diag}{psat\_water\_validated}{rms\_heldout\_log10P}{2.76024e-10}
\goldrow{diag}{psat\_water\_refused}{A}{5.40221}
\goldrow{diag}{psat\_water\_refused}{B}{1838.68}
\goldrow{diag}{psat\_water\_refused}{C}{-31.737}
\goldrow{diag}{psat\_water\_refused}{R2}{1}
\goldrow{diag}{psat\_water\_refused}{dHvap\_kJ\_per\_mol}{42.1162}
\goldmore{4}
\end{tutgolden}

\clearpage
\tutentry{curate02\_vle\_heldout\_ethanol\_water}{tutorials/props/curation/curate02\_vle\_heldout\_ethanol\_water}
\begin{tutscheme}
\begin{tikzpicture}[>={Stealth[length=2mm]}, every node/.style={font=\scriptsize}]
  \node[draw, rounded corners=2pt, align=center, minimum width=3.1cm, minimum height=0.8cm, fill=black!4] at (0.0,-0.0) {\textbf{nrtl\_etoh\_water\_heldout}\\[-1pt]{\tiny fitParameters}};
\end{tikzpicture}
\end{tutscheme}
\tutwhat{A binary NRTL pair regressed on eight MEASURED bubble temperatures and scored on three it never saw -- the first held-out validation in the corpus against real evidence}
\begin{tutnarrative}
curate02 -- THE FIT / HELD-OUT CONTRACT ON A BINARY PHASE EQUILIBRIUM.

The NRTL ethanol-water pair is regressed against eight measured bubble temperatures at 101.325 kPa, and then asked to predict THREE it never saw.

WHY THIS CASE EXISTS BESIDE curate01. curate01 proves the machinery: its two halves come from the same Antoine equation, so its held-out residual is round-off and its acceptance band is a numerical criterion, which its own `acceptance.origin` says out loud. Here the evidence is MEASURED, so the held-out residual is a statement about the MODEL -- the first place in the corpus where a fitted binary pair is scored on real evidence it was never handed.

THE SPLIT IS DECLARED AND ARITHMETIC. Every third point of the eleven-point series (x1 = 0.20, 0.50, 0.80) is held out, spread across the composition range rather than cut contiguously, because a model can interpolate well and extrapolate badly and a contiguous cut tests only one of those. The rule was fixed before anything was fitted.

THE ACCEPTANCE BAND IS DECLARED BEFORE THE FIT, WITH ITS ORIGIN, AND IS NOT ADJUSTED AFTERWARDS. Whatever verdict falls out of it is the result. A band chosen after seeing the residual would make the whole contract decorative.

Read the run's four sections in order: FIT QUALITY is what the model does on what it saw; HELD-OUT EVIDENCE is what it did on what it did not.

TWO THINGS THIS CASE DOES NOT SHOW, BOTH VISIBLE IN ITS OWN GOLDEN.

(1) THE PARAMETERS ARE NOT IDENTIFIABLE. `identifiable 0`, `cond\_JtJ \textasciitilde{}4e11` and pairwise correlations of 0.9999 between each a/b pair: eight points at ONE pressure cannot separate the entropic and enthalpic halves of an NRTL parameter, because they only ever appear together over a narrow T span. The model PREDICTS well and the four numbers are individually meaningless -- a validated prediction is not a determined parameter set, and a student must be able to read both facts off the same run. This is why the verdict is scoped `binaryVLE.T\_bubble` in the dossier and not to the pair itself.

(2) THE HELD-OUT MARGIN IS NOT LARGE. 0.0737 \% against a 0.1 \% band is a pass by a factor of 1.4, not by an order of magnitude. The band was declared before the fit and is not adjusted; the margin is reported rather than described as comfortable.

The run also announces an extrapolation: ethanol's vapour pressure is evaluated at 372 K, outside its declared Trange (273-369). Announced, not judged -- it rides the caveat surface like every other extrapolation.

WHAT THE STUDENT TAKES AWAY: `ethanol-water.proposal.dat`.

The console says all of the above and then scrolls. The file is what survives -- it is moved into `constant/parameters/NRTL/ethanol-water.dat` and simulated with, months later, by someone who never saw this run. So it carries the whole qualification, not the numbers alone: the two datasets and the role each was frozen in, the fingerprint of that declaration, the held-out AAD over the three withheld points, the 0.1 \% band with the reason it was set at 0.1 \%, the verdict, and -- in the header, above everything -- the fact that these four coefficients are NOT individually identifiable.

Both facts of (1) are readable off that one file: `verdict validated` beside `identifiable false`. A validated prediction is not a determined parameter set, and the record says so in the artefact rather than only in a run log.

\end{tutnarrative}
\begin{tutgolden}
\goldrow{diag}{nrtl\_etoh\_water\_heldout}{aad\_heldout\_pct}{0.0736714746429}
\goldrow{diag}{nrtl\_etoh\_water\_heldout}{chi2}{0.809699405251}
\goldrow{diag}{nrtl\_etoh\_water\_heldout}{chi2\_in\_sample}{0.809699405251}
\goldrow{diag}{nrtl\_etoh\_water\_heldout}{chi2\_reduced}{0.202424851313}
\goldrow{diag}{nrtl\_etoh\_water\_heldout}{cond\_JtJ}{394331596992}
\goldrow{diag}{nrtl\_etoh\_water\_heldout}{converged}{1}
\goldrow{diag}{nrtl\_etoh\_water\_heldout}{corr\_computed}{1}
\goldrow{diag}{nrtl\_etoh\_water\_heldout}{dof}{4}
\goldrow{diag}{nrtl\_etoh\_water\_heldout}{identifiable}{0}
\goldrow{diag}{nrtl\_etoh\_water\_heldout}{invertible}{1}
\goldrow{diag}{nrtl\_etoh\_water\_heldout}{iter}{29}
\goldrow{diag}{nrtl\_etoh\_water\_heldout}{max\_abs\_corr}{0.999936039821}
\goldrow{diag}{nrtl\_etoh\_water\_heldout}{max\_abs\_resid}{0.582462818864}
\goldrow{diag}{nrtl\_etoh\_water\_heldout}{n\_data}{8}
\goldmore{9}
\end{tutgolden}

\clearpage
\tutentry{curate03\_thermoml\_fixture\_bubble}{tutorials/props/curation/curate03\_thermoml\_fixture\_bubble}
\begin{tutscheme}
\begin{tikzpicture}[>={Stealth[length=2mm]}, every node/.style={font=\scriptsize}]
  \node[draw, rounded corners=2pt, align=center, minimum width=3.1cm, minimum height=0.8cm, fill=black!4] at (0.0,-0.0) {\textbf{thermoml\_fixture\_bubble}\\[-1pt]{\tiny fitParameters}};
\end{tikzpicture}
\end{tutscheme}
\tutwhat{ThermoML $\to$ extract $\to$ partition $\to$ fit $\to$ held-out $\to$ dossier, on a SYNTHETIC structural fixture: the plumbing proven end to end, and no physical claim anywhere}
\begin{tutnarrative}
curate03 -- THE THERMOML PATH, END TO END, ON A STRUCTURAL FIXTURE.

ThermoML .xml $\to$ extract-vle $\to$ EvidencePartition $\to$ FIT only $\to$ T\_bubble regression $\to$ held-out $\to$ dossier

EVERY NUMBER HERE IS INVENTED. Both datasets were extracted by `bin/choupo-thermoml extract-vle` from synthetic fixtures under tutorials/props/curation/fixtures/thermoml/, and each carries `provenance \{ source synthetic; \}` -- so the run ANNOUNCES that its evidence was generated, and the dossier records it per dataset. A reader cannot arrive at the verdict without having been told what the numbers are.

WHAT THIS CASE PROVES: that a ThermoML file crosses into the evidence contract without losing its identity -- archive file, DOI, system, pressure and checksum all survive into the dossier -- and that the partition, the regression and the verdict work on evidence that arrived through the parser rather than being typed into a case.

WHAT IT DOES NOT PROVE: anything about ethanol, water, or NRTL. The acceptance band below is a PLUMBING tolerance and says so. The scientific witness on measured evidence is curate02, next door.

\end{tutnarrative}
\begin{tutgolden}
\goldrow{diag}{thermoml\_fixture\_bubble}{aad\_heldout\_pct}{0.0528782856338}
\goldrow{diag}{thermoml\_fixture\_bubble}{chi2}{0.669244065892}
\goldrow{diag}{thermoml\_fixture\_bubble}{chi2\_in\_sample}{0.669244065892}
\goldrow{diag}{thermoml\_fixture\_bubble}{chi2\_reduced}{0.223081355297}
\goldrow{diag}{thermoml\_fixture\_bubble}{cond\_JtJ}{272780104374}
\goldrow{diag}{thermoml\_fixture\_bubble}{converged}{1}
\goldrow{diag}{thermoml\_fixture\_bubble}{corr\_computed}{1}
\goldrow{diag}{thermoml\_fixture\_bubble}{dof}{3}
\goldrow{diag}{thermoml\_fixture\_bubble}{identifiable}{0}
\goldrow{diag}{thermoml\_fixture\_bubble}{invertible}{1}
\goldrow{diag}{thermoml\_fixture\_bubble}{iter}{42}
\goldrow{diag}{thermoml\_fixture\_bubble}{max\_abs\_corr}{0.999960515815}
\goldrow{diag}{thermoml\_fixture\_bubble}{max\_abs\_resid}{0.656169896008}
\goldrow{diag}{thermoml\_fixture\_bubble}{n\_data}{7}
\goldmore{9}
\end{tutgolden}

\clearpage
\tutentry{curate04\_nrtl\_methanol\_water}{tutorials/props/curation/curate04\_nrtl\_methanol\_water}
\begin{tutscheme}
\begin{tikzpicture}[>={Stealth[length=2mm]}, every node/.style={font=\scriptsize}]
  \node[draw, rounded corners=2pt, align=center, minimum width=3.1cm, minimum height=0.8cm, fill=black!4] at (0.0,-0.0) {\textbf{fit\_nrtl}\\[-1pt]{\tiny fitParameters}};
  \node[draw, rounded corners=2pt, align=center, minimum width=3.1cm, minimum height=0.8cm, fill=black!4] at (3.6,-0.0) {\textbf{fit\_nrtl\_start2}\\[-1pt]{\tiny fitParameters}};
  \node[draw, rounded corners=2pt, align=center, minimum width=3.1cm, minimum height=0.8cm, fill=black!4] at (7.2,-0.0) {\textbf{fit\_nrtl\_ab}\\[-1pt]{\tiny fitParameters}};
  \node[draw, rounded corners=2pt, align=center, minimum width=3.1cm, minimum height=0.8cm, fill=black!4] at (0.0,-1.3) {\textbf{txy\_ideal}\\[-1pt]{\tiny propertyScan1D}};
  \node[draw, rounded corners=2pt, align=center, minimum width=3.1cm, minimum height=0.8cm, fill=black!4] at (3.6,-1.3) {\textbf{txy\_unifac}\\[-1pt]{\tiny propertyScan1D}};
  \node[draw, rounded corners=2pt, align=center, minimum width=3.1cm, minimum height=0.8cm, fill=black!4] at (7.2,-1.3) {\textbf{txy\_nrtl}\\[-1pt]{\tiny propertyScan1D}};
\end{tikzpicture}
\end{tutscheme}
\tutwhat{NRTL methanol-water regressed to measured bubble temperatures (ThermoML locates, the article sources), held out against an independent study, drawn against both beside ideal and UNIFAC}
\begin{tutnarrative}
AN NRTL PAIR FROM PUBLISHED MEASUREMENTS, END TO END.

WHY THIS SYSTEM. the commonest aqueous alcohol separation in a first distillation course; the fit set spans six isobars (15-95 kPa), so it can ask whether a SECOND temperature coefficient is identified.

THE CHAIN, EACH STEP A COMMAND (DEV.md C16; record docs/design/binary-pairs-from-open-measurements.md): 1. bin/choupo-thermoml sync the NIST/TRC ThermoML Archive (ThermoML.v2020-09-30.tgz, doi 10.18434/mds2-2422, sha256-verified) 2. bin/choupo-thermoml identity CAS $\to$ InChIKey for the catalogue, two resolvers that must agree 3. bin/choupo-thermoml vle-index every binary VLE block, searchable 4. bin/choupo-thermoml vle-search methanol water --geometry bubbleT 5. bin/choupo-thermoml extract-vle <file> --block N --interior-only $\to$ the two files under constant/experiments/ 6. choupoProps (this case) fitParameters $\to$ a proposal record 7. promotion (a human act) $\to$ data/standards/parameters/NRTL/methanol-water.dat 8. bin/choupo-import the seal mirrors that record here

THE EVIDENCE. FIT: doi 10.1016/j.jct.2009.11.020 (54 interior points, several isobars). HELD OUT: doi 10.1021/je200341c (13 interior points, 101.32 kPa) -- a different article, so the held-out AAD is a validation, not a slice of the fit. Each file names its article, its archive file and block, the sha256 of the archive file, the two compounds' CAS and InChIKey, and the per-point expanded uncertainty; each says `reviewStatus transcribedNotCheckedAgainstArticle` because nobody has read the numbers back against the paper yet, and the run announces it. Pure-component endpoints are set aside (they measure the vapour-pressure record, not the pair) and kept in each file under `selection`.

THE FIT. NRTL with a\_ij = a\_ji = 0 and alpha = 0.30 fixed: TWO parameters (b\_ij, b\_ji), because bubble temperatures over the ranges below do not identify more (fit\_nrtl\_ab in the methanol + water case asks that question directly). Each point is weighted by 1/U, U its OWN stated expanded uncertainty (`weighting uncertainty;`), so a point measured to 0.1 K is not outvoted by one measured to 1 K. The held-out AAD is plain kelvin, unweighted.

WHAT THE RUN FOUND (pinned in `expected`; measured, not argued). fit\_nrtl: b\_ij = 165.0 K, b\_ji = 41.8 K, in-sample rms 0.50 K over 54 points on six isobars. fit\_nrtl\_start2, from (800, -400), lands on the same pair to 0.01 K -- one minimum, not a start point's choice. HELD OUT (101.32 kPa, another laboratory): AAD 0.53 K (0.150 \%) against a band of 0.33 \% declared before the fit $\to$ `validated`, and the pair is the catalogue record NRTL/methanol-water.dat. fit\_nrtl\_ab frees a\_ij and a\_ji too: the run says `NOT individually identifiable` (a and b trade off, max|corr| 0.996) and its held-out AAD is WORSE (0.64 K). Six isobars over 15-95 kPa do not pay for a second temperature coefficient -- which is why the record carries two. THE OVERLAY SAYS SOMETHING THE VERDICT DOES NOT: on the held-out set UNIFAC (0.36 K) is closer than the fitted pair (0.53 K). The pair was regressed below atmospheric pressure and is validated, not best; `validated` means inside the declared band, never better than every alternative. Ideal mixing misses by 2.8 K. Part of every residual here is not the pair's: the catalogue water Antoine record boils 0.70 K low at 101.32 kPa (measured at this set's own pure endpoint), and no pair can remove that.

WHAT IS NOT CLAIMED. The pair is correct at these pressures and in this frame (ideal-gas vapour, the catalogue's Antoine records); nothing here tests it at high pressure, against LLE, or with another vapour model.

\end{tutnarrative}
\begin{tutgolden}
\goldrow{diag}{fit\_nrtl}{aad\_heldout\_K}{0.534750672125}
\goldrow{diag}{fit\_nrtl}{aad\_heldout\_pct}{0.150061425821}
\goldrow{diag}{fit\_nrtl}{chi2}{315.937277131}
\goldrow{diag}{fit\_nrtl}{chi2\_in\_sample}{315.937277131}
\goldrow{diag}{fit\_nrtl}{chi2\_reduced}{6.07571686791}
\goldrow{diag}{fit\_nrtl}{cond\_JtJ}{116.915709849}
\goldrow{diag}{fit\_nrtl}{converged}{1}
\goldrow{diag}{fit\_nrtl}{corr\_computed}{1}
\goldrow{diag}{fit\_nrtl}{dof}{52}
\goldrow{diag}{fit\_nrtl}{identifiable}{1}
\goldrow{diag}{fit\_nrtl}{invertible}{1}
\goldrow{diag}{fit\_nrtl}{iter}{7}
\goldrow{diag}{fit\_nrtl}{max\_abs\_corr}{0.995750654438}
\goldrow{diag}{fit\_nrtl}{max\_abs\_resid}{2.0802446823}
\goldmore{76}
\end{tutgolden}

\clearpage
\tutentry{curate05\_nrtl\_acetone\_methanol}{tutorials/props/curation/curate05\_nrtl\_acetone\_methanol}
\begin{tutscheme}
\begin{tikzpicture}[>={Stealth[length=2mm]}, every node/.style={font=\scriptsize}]
  \node[draw, rounded corners=2pt, align=center, minimum width=3.1cm, minimum height=0.8cm, fill=black!4] at (0.0,-0.0) {\textbf{fit\_nrtl}\\[-1pt]{\tiny fitParameters}};
  \node[draw, rounded corners=2pt, align=center, minimum width=3.1cm, minimum height=0.8cm, fill=black!4] at (3.6,-0.0) {\textbf{fit\_nrtl\_start2}\\[-1pt]{\tiny fitParameters}};
  \node[draw, rounded corners=2pt, align=center, minimum width=3.1cm, minimum height=0.8cm, fill=black!4] at (7.2,-0.0) {\textbf{txy\_ideal}\\[-1pt]{\tiny propertyScan1D}};
  \node[draw, rounded corners=2pt, align=center, minimum width=3.1cm, minimum height=0.8cm, fill=black!4] at (0.0,-1.3) {\textbf{txy\_unifac}\\[-1pt]{\tiny propertyScan1D}};
  \node[draw, rounded corners=2pt, align=center, minimum width=3.1cm, minimum height=0.8cm, fill=black!4] at (3.6,-1.3) {\textbf{txy\_nrtl}\\[-1pt]{\tiny propertyScan1D}};
  \node[draw, rounded corners=2pt, align=center, minimum width=3.1cm, minimum height=0.8cm, fill=black!4] at (7.2,-1.3) {\textbf{txy\_ideal\_fitP}\\[-1pt]{\tiny propertyScan1D}};
  \node[draw, rounded corners=2pt, align=center, minimum width=3.1cm, minimum height=0.8cm, fill=black!4] at (0.0,-2.6) {\textbf{txy\_unifac\_fitP}\\[-1pt]{\tiny propertyScan1D}};
  \node[draw, rounded corners=2pt, align=center, minimum width=3.1cm, minimum height=0.8cm, fill=black!4] at (3.6,-2.6) {\textbf{txy\_nrtl\_fitP}\\[-1pt]{\tiny propertyScan1D}};
\end{tikzpicture}
\end{tutscheme}
\tutwhat{NRTL acetone-methanol regressed to measured bubble temperatures (ThermoML locates, the article sources), held out against an independent study, drawn against both beside ideal and UNIFAC}
\begin{tutnarrative}
AN NRTL PAIR FROM PUBLISHED MEASUREMENTS, END TO END.

WHY THIS SYSTEM. a minimum-boiling azeotrope between two components a student meets in the acetone and methanol plants of this corpus.

THE CHAIN, EACH STEP A COMMAND (DEV.md C16; record docs/design/binary-pairs-from-open-measurements.md): 1. bin/choupo-thermoml sync the NIST/TRC ThermoML Archive (ThermoML.v2020-09-30.tgz, doi 10.18434/mds2-2422, sha256-verified) 2. bin/choupo-thermoml identity CAS $\to$ InChIKey for the catalogue, two resolvers that must agree 3. bin/choupo-thermoml vle-index every binary VLE block, searchable 4. bin/choupo-thermoml vle-search acetone methanol --geometry bubbleT 5. bin/choupo-thermoml extract-vle <file> --block N --interior-only $\to$ the two files under constant/experiments/ 6. choupoProps (this case) fitParameters $\to$ a proposal record 7. promotion NOT TAKEN: the verdict is notValidated, so the record stays in THIS case (constant/parameters/NRTL/acetone-methanol.dat) 8. bin/choupo-import the seal keeps that authored record as the case's own and mirrors the rest

THE EVIDENCE. FIT: doi 10.1021/je5007373 (32 interior points, 101.3 kPa). HELD OUT: doi 10.1021/je0603170 (30 interior points, 100 kPa) -- a different article, so the held-out AAD is a validation, not a slice of the fit. Each file names its article, its archive file and block, the sha256 of the archive file, the two compounds' CAS and InChIKey, and the per-point expanded uncertainty; each says `reviewStatus transcribedNotCheckedAgainstArticle` because nobody has read the numbers back against the paper yet, and the run announces it. Pure-component endpoints are set aside (they measure the vapour-pressure record, not the pair) and kept in each file under `selection`.

THE FIT. NRTL with a\_ij = a\_ji = 0 and alpha = 0.30 fixed: TWO parameters (b\_ij, b\_ji), because bubble temperatures over the ranges below do not identify more (fit\_nrtl\_ab in the methanol + water case asks that question directly). Each point is weighted by 1/U, U its OWN stated expanded uncertainty (`weighting uncertainty;`), so a point measured to 0.1 K is not outvoted by one measured to 1 K. The held-out AAD is plain kelvin, unweighted.

WHAT THE RUN FOUND (pinned in `expected`; measured, not argued). fit\_nrtl: b\_ij = 369.8 K, b\_ji = -92.8 K, in-sample rms 0.13 K; fit\_nrtl\_start2 lands on the same pair. HELD OUT (100 kPa, another laboratory): AAD 0.32 K (0.096 \%) against a band of 0.047 \% declared before the fit $\to$ `notValidated`. THE PAIR IS NOT PROMOTED: its record lives in this case ONLY (constant/parameters/NRTL/acetone-methanol.dat, reviewStatus interim, verdict notValidated), so the scans can still draw it. WHY, as far as the run can see: the two studies disagree with each other. UNIFAC, which saw neither, reproduces the held-out study to 0.04 K and misses the fit study by 0.34 K; the fitted pair does the mirror image (0.10 K on its own study, 0.32 K on the other). Two measurements of one mixture, 0.3 K apart where each states about 0.1-0.35 K: a finding about the EVIDENCE, for a curator to settle by reading both articles. Swapping the roles now would be choosing the partition after seeing the residual, which is the tuning the declared partition exists to prevent.

WHAT IS NOT CLAIMED. The pair is correct at these pressures and in this frame (ideal-gas vapour, the catalogue's Antoine records); nothing here tests it at high pressure, against LLE, or with another vapour model.

\end{tutnarrative}
\begin{tutgolden}
\goldrow{diag}{fit\_nrtl}{aad\_heldout\_K}{0.315969383935}
\goldrow{diag}{fit\_nrtl}{aad\_heldout\_pct}{0.0960285506587}
\goldrow{diag}{fit\_nrtl}{chi2}{4.45236740196}
\goldrow{diag}{fit\_nrtl}{chi2\_in\_sample}{4.45236740196}
\goldrow{diag}{fit\_nrtl}{chi2\_reduced}{0.148412246732}
\goldrow{diag}{fit\_nrtl}{cond\_JtJ}{95.8483654535}
\goldrow{diag}{fit\_nrtl}{converged}{1}
\goldrow{diag}{fit\_nrtl}{corr\_computed}{1}
\goldrow{diag}{fit\_nrtl}{dof}{30}
\goldrow{diag}{fit\_nrtl}{identifiable}{1}
\goldrow{diag}{fit\_nrtl}{invertible}{1}
\goldrow{diag}{fit\_nrtl}{iter}{11}
\goldrow{diag}{fit\_nrtl}{max\_abs\_corr}{0.994823894221}
\goldrow{diag}{fit\_nrtl}{max\_abs\_resid}{0.242779577201}
\goldmore{70}
\end{tutgolden}

\clearpage
\tutentry{curate06\_nrtl\_isopropanol\_water}{tutorials/props/curation/curate06\_nrtl\_isopropanol\_water}
\begin{tutscheme}
\begin{tikzpicture}[>={Stealth[length=2mm]}, every node/.style={font=\scriptsize}]
  \node[draw, rounded corners=2pt, align=center, minimum width=3.1cm, minimum height=0.8cm, fill=black!4] at (0.0,-0.0) {\textbf{fit\_nrtl}\\[-1pt]{\tiny fitParameters}};
  \node[draw, rounded corners=2pt, align=center, minimum width=3.1cm, minimum height=0.8cm, fill=black!4] at (3.6,-0.0) {\textbf{fit\_nrtl\_start2}\\[-1pt]{\tiny fitParameters}};
  \node[draw, rounded corners=2pt, align=center, minimum width=3.1cm, minimum height=0.8cm, fill=black!4] at (7.2,-0.0) {\textbf{txy\_ideal}\\[-1pt]{\tiny propertyScan1D}};
  \node[draw, rounded corners=2pt, align=center, minimum width=3.1cm, minimum height=0.8cm, fill=black!4] at (0.0,-1.3) {\textbf{txy\_unifac}\\[-1pt]{\tiny propertyScan1D}};
  \node[draw, rounded corners=2pt, align=center, minimum width=3.1cm, minimum height=0.8cm, fill=black!4] at (3.6,-1.3) {\textbf{txy\_nrtl}\\[-1pt]{\tiny propertyScan1D}};
  \node[draw, rounded corners=2pt, align=center, minimum width=3.1cm, minimum height=0.8cm, fill=black!4] at (7.2,-1.3) {\textbf{txy\_ideal\_fitP}\\[-1pt]{\tiny propertyScan1D}};
  \node[draw, rounded corners=2pt, align=center, minimum width=3.1cm, minimum height=0.8cm, fill=black!4] at (0.0,-2.6) {\textbf{txy\_unifac\_fitP}\\[-1pt]{\tiny propertyScan1D}};
  \node[draw, rounded corners=2pt, align=center, minimum width=3.1cm, minimum height=0.8cm, fill=black!4] at (3.6,-2.6) {\textbf{txy\_nrtl\_fitP}\\[-1pt]{\tiny propertyScan1D}};
\end{tikzpicture}
\end{tutscheme}
\tutwhat{NRTL isopropanol-water regressed to measured bubble temperatures (ThermoML locates, the article sources), held out against an independent study, drawn against both beside ideal and UNIFAC}
\begin{tutnarrative}
AN NRTL PAIR FROM PUBLISHED MEASUREMENTS, END TO END.

WHY THIS SYSTEM. the azeotrope the acetone-from-isopropanol tutorials have to break; the fit set's stated uncertainty varies tenfold along x, which is what the uncertainty weighting is for.

THE CHAIN, EACH STEP A COMMAND (DEV.md C16; record docs/design/binary-pairs-from-open-measurements.md): 1. bin/choupo-thermoml sync the NIST/TRC ThermoML Archive (ThermoML.v2020-09-30.tgz, doi 10.18434/mds2-2422, sha256-verified) 2. bin/choupo-thermoml identity CAS $\to$ InChIKey for the catalogue, two resolvers that must agree 3. bin/choupo-thermoml vle-index every binary VLE block, searchable 4. bin/choupo-thermoml vle-search isopropanol water --geometry bubbleT 5. bin/choupo-thermoml extract-vle <file> --block N --interior-only $\to$ the two files under constant/experiments/ 6. choupoProps (this case) fitParameters $\to$ a proposal record 7. promotion (a human act) $\to$ data/standards/parameters/NRTL/isopropanol-water.dat 8. bin/choupo-import the seal mirrors that record here

THE EVIDENCE. FIT: doi 10.1016/j.jct.2017.02.005 (29 interior points, 100 kPa). HELD OUT: doi 10.1021/acs.jced.7b00523 (16 interior points, 101.058 kPa) -- a different article, so the held-out AAD is a validation, not a slice of the fit. Each file names its article, its archive file and block, the sha256 of the archive file, the two compounds' CAS and InChIKey, and the per-point expanded uncertainty; each says `reviewStatus transcribedNotCheckedAgainstArticle` because nobody has read the numbers back against the paper yet, and the run announces it. Pure-component endpoints are set aside (they measure the vapour-pressure record, not the pair) and kept in each file under `selection`.

THE FIT. NRTL with a\_ij = a\_ji = 0 and alpha = 0.30 fixed: TWO parameters (b\_ij, b\_ji), because bubble temperatures over the ranges below do not identify more (fit\_nrtl\_ab in the methanol + water case asks that question directly). Each point is weighted by 1/U, U its OWN stated expanded uncertainty (`weighting uncertainty;`), so a point measured to 0.1 K is not outvoted by one measured to 1 K. The held-out AAD is plain kelvin, unweighted.

WHAT THE RUN FOUND (pinned in `expected`; measured, not argued). fit\_nrtl: b\_ij = -69.6 K, b\_ji = 968.2 K, in-sample rms 0.52 K; fit\_nrtl\_start2 lands on the same pair. The fit set's stated uncertainty runs from 0.13 to 1.35 K along x, so the 1/U weighting changes which points the pair is held to. HELD OUT (101.058 kPa, another laboratory): AAD 0.83 K (0.232 \%) against a band of 0.63 \% $\to$ `validated`; the pair is the catalogue record NRTL/isopropanol-water.dat. Overlay on the held-out set: NRTL 0.83 K, UNIFAC 1.27 K, ideal 8.2 K. The held-out study's own stated uncertainty is 1.6-1.7 K, and the catalogue water record boils 0.52 K low at its pressure.

WHAT IS NOT CLAIMED. The pair is correct at these pressures and in this frame (ideal-gas vapour, the catalogue's Antoine records); nothing here tests it at high pressure, against LLE, or with another vapour model.

\end{tutnarrative}
\begin{tutgolden}
\goldrow{diag}{fit\_nrtl}{aad\_heldout\_K}{0.84147821571}
\goldrow{diag}{fit\_nrtl}{aad\_heldout\_pct}{0.234200546209}
\goldrow{diag}{fit\_nrtl}{chi2}{90.1257883215}
\goldrow{diag}{fit\_nrtl}{chi2\_in\_sample}{90.1257883215}
\goldrow{diag}{fit\_nrtl}{chi2\_reduced}{3.33799216005}
\goldrow{diag}{fit\_nrtl}{cond\_JtJ}{39.3738330039}
\goldrow{diag}{fit\_nrtl}{converged}{1}
\goldrow{diag}{fit\_nrtl}{corr\_computed}{1}
\goldrow{diag}{fit\_nrtl}{dof}{27}
\goldrow{diag}{fit\_nrtl}{identifiable}{1}
\goldrow{diag}{fit\_nrtl}{invertible}{1}
\goldrow{diag}{fit\_nrtl}{iter}{21}
\goldrow{diag}{fit\_nrtl}{max\_abs\_corr}{0.89819126525}
\goldrow{diag}{fit\_nrtl}{max\_abs\_resid}{1.65678230014}
\goldmore{70}
\end{tutgolden}

\clearpage
\tutentry{curate07\_nrtl\_ethylacetate\_ethanol}{tutorials/props/curation/curate07\_nrtl\_ethylacetate\_ethanol}
\begin{tutscheme}
\begin{tikzpicture}[>={Stealth[length=2mm]}, every node/.style={font=\scriptsize}]
  \node[draw, rounded corners=2pt, align=center, minimum width=3.1cm, minimum height=0.8cm, fill=black!4] at (0.0,-0.0) {\textbf{fit\_nrtl}\\[-1pt]{\tiny fitParameters}};
  \node[draw, rounded corners=2pt, align=center, minimum width=3.1cm, minimum height=0.8cm, fill=black!4] at (3.6,-0.0) {\textbf{fit\_nrtl\_start2}\\[-1pt]{\tiny fitParameters}};
  \node[draw, rounded corners=2pt, align=center, minimum width=3.1cm, minimum height=0.8cm, fill=black!4] at (7.2,-0.0) {\textbf{txy\_ideal}\\[-1pt]{\tiny propertyScan1D}};
  \node[draw, rounded corners=2pt, align=center, minimum width=3.1cm, minimum height=0.8cm, fill=black!4] at (0.0,-1.3) {\textbf{txy\_unifac}\\[-1pt]{\tiny propertyScan1D}};
  \node[draw, rounded corners=2pt, align=center, minimum width=3.1cm, minimum height=0.8cm, fill=black!4] at (3.6,-1.3) {\textbf{txy\_nrtl}\\[-1pt]{\tiny propertyScan1D}};
  \node[draw, rounded corners=2pt, align=center, minimum width=3.1cm, minimum height=0.8cm, fill=black!4] at (7.2,-1.3) {\textbf{txy\_ideal\_fitP}\\[-1pt]{\tiny propertyScan1D}};
  \node[draw, rounded corners=2pt, align=center, minimum width=3.1cm, minimum height=0.8cm, fill=black!4] at (0.0,-2.6) {\textbf{txy\_unifac\_fitP}\\[-1pt]{\tiny propertyScan1D}};
  \node[draw, rounded corners=2pt, align=center, minimum width=3.1cm, minimum height=0.8cm, fill=black!4] at (3.6,-2.6) {\textbf{txy\_nrtl\_fitP}\\[-1pt]{\tiny propertyScan1D}};
\end{tikzpicture}
\end{tutscheme}
\tutwhat{NRTL ethylAcetate-ethanol regressed to measured bubble temperatures (ThermoML locates, the article sources), held out against an independent study, drawn against both beside ideal and UNIFAC}
\begin{tutnarrative}
AN NRTL PAIR FROM PUBLISHED MEASUREMENTS, END TO END.

WHY THIS SYSTEM. the esterification product and its reactant, a minimum-boiling azeotrope the esterification tutorials run through.

THE CHAIN, EACH STEP A COMMAND (DEV.md C16; record docs/design/binary-pairs-from-open-measurements.md): 1. bin/choupo-thermoml sync the NIST/TRC ThermoML Archive (ThermoML.v2020-09-30.tgz, doi 10.18434/mds2-2422, sha256-verified) 2. bin/choupo-thermoml identity CAS $\to$ InChIKey for the catalogue, two resolvers that must agree 3. bin/choupo-thermoml vle-index every binary VLE block, searchable 4. bin/choupo-thermoml vle-search ethylAcetate ethanol --geometry bubbleT 5. bin/choupo-thermoml extract-vle <file> --block N --interior-only $\to$ the two files under constant/experiments/ 6. choupoProps (this case) fitParameters $\to$ a proposal record 7. promotion (a human act) $\to$ data/standards/parameters/NRTL/ethanol-ethylAcetate.dat 8. bin/choupo-import the seal mirrors that record here

THE EVIDENCE. FIT: doi 10.1016/j.fluid.2005.07.010 (24 interior points, 101.325 kPa). HELD OUT: doi 10.1021/je700322p (28 interior points, 100 kPa) -- a different article, so the held-out AAD is a validation, not a slice of the fit. Each file names its article, its archive file and block, the sha256 of the archive file, the two compounds' CAS and InChIKey, and the per-point expanded uncertainty; each says `reviewStatus transcribedNotCheckedAgainstArticle` because nobody has read the numbers back against the paper yet, and the run announces it. Pure-component endpoints are set aside (they measure the vapour-pressure record, not the pair) and kept in each file under `selection`.

THE FIT. NRTL with a\_ij = a\_ji = 0 and alpha = 0.30 fixed: TWO parameters (b\_ij, b\_ji), because bubble temperatures over the ranges below do not identify more (fit\_nrtl\_ab in the methanol + water case asks that question directly). Each point is weighted by 1/U, U its OWN stated expanded uncertainty (`weighting uncertainty;`), so a point measured to 0.1 K is not outvoted by one measured to 1 K. The held-out AAD is plain kelvin, unweighted.

WHAT THE RUN FOUND (pinned in `expected`; measured, not argued). fit\_nrtl: b\_ij = 95.4 K, b\_ji = 214.8 K, in-sample rms 0.077 K; fit\_nrtl\_start2 lands on the same pair. HELD OUT (100 kPa, another laboratory): AAD 0.11 K (0.032 \%) against a band of 0.052 \% $\to$ `validated`; the pair is the catalogue record NRTL/ethanol-ethylAcetate.dat. Overlay on the held-out set: NRTL 0.11 K, UNIFAC 0.69 K, ideal 4.1 K.

WHAT IS NOT CLAIMED. The pair is correct at these pressures and in this frame (ideal-gas vapour, the catalogue's Antoine records); nothing here tests it at high pressure, against LLE, or with another vapour model.

\end{tutnarrative}
\begin{tutgolden}
\goldrow{diag}{fit\_nrtl}{aad\_heldout\_K}{0.11098102751}
\goldrow{diag}{fit\_nrtl}{aad\_heldout\_pct}{0.0321107150642}
\goldrow{diag}{fit\_nrtl}{chi2}{4.34313933595}
\goldrow{diag}{fit\_nrtl}{chi2\_in\_sample}{4.34313933595}
\goldrow{diag}{fit\_nrtl}{chi2\_reduced}{0.197415424361}
\goldrow{diag}{fit\_nrtl}{cond\_JtJ}{151.806801216}
\goldrow{diag}{fit\_nrtl}{converged}{1}
\goldrow{diag}{fit\_nrtl}{corr\_computed}{1}
\goldrow{diag}{fit\_nrtl}{dof}{22}
\goldrow{diag}{fit\_nrtl}{identifiable}{1}
\goldrow{diag}{fit\_nrtl}{invertible}{1}
\goldrow{diag}{fit\_nrtl}{iter}{6}
\goldrow{diag}{fit\_nrtl}{max\_abs\_corr}{0.995865280749}
\goldrow{diag}{fit\_nrtl}{max\_abs\_resid}{0.194183542276}
\goldmore{70}
\end{tutgolden}

\clearpage
\tutentry{curate08\_nrtl\_methylacetate\_methanol}{tutorials/props/curation/curate08\_nrtl\_methylacetate\_methanol}
\begin{tutscheme}
\begin{tikzpicture}[>={Stealth[length=2mm]}, every node/.style={font=\scriptsize}]
  \node[draw, rounded corners=2pt, align=center, minimum width=3.1cm, minimum height=0.8cm, fill=black!4] at (0.0,-0.0) {\textbf{fit\_nrtl}\\[-1pt]{\tiny fitParameters}};
  \node[draw, rounded corners=2pt, align=center, minimum width=3.1cm, minimum height=0.8cm, fill=black!4] at (3.6,-0.0) {\textbf{fit\_nrtl\_start2}\\[-1pt]{\tiny fitParameters}};
  \node[draw, rounded corners=2pt, align=center, minimum width=3.1cm, minimum height=0.8cm, fill=black!4] at (7.2,-0.0) {\textbf{txy\_ideal}\\[-1pt]{\tiny propertyScan1D}};
  \node[draw, rounded corners=2pt, align=center, minimum width=3.1cm, minimum height=0.8cm, fill=black!4] at (0.0,-1.3) {\textbf{txy\_unifac}\\[-1pt]{\tiny propertyScan1D}};
  \node[draw, rounded corners=2pt, align=center, minimum width=3.1cm, minimum height=0.8cm, fill=black!4] at (3.6,-1.3) {\textbf{txy\_nrtl}\\[-1pt]{\tiny propertyScan1D}};
  \node[draw, rounded corners=2pt, align=center, minimum width=3.1cm, minimum height=0.8cm, fill=black!4] at (7.2,-1.3) {\textbf{txy\_ideal\_fitP}\\[-1pt]{\tiny propertyScan1D}};
  \node[draw, rounded corners=2pt, align=center, minimum width=3.1cm, minimum height=0.8cm, fill=black!4] at (0.0,-2.6) {\textbf{txy\_unifac\_fitP}\\[-1pt]{\tiny propertyScan1D}};
  \node[draw, rounded corners=2pt, align=center, minimum width=3.1cm, minimum height=0.8cm, fill=black!4] at (3.6,-2.6) {\textbf{txy\_nrtl\_fitP}\\[-1pt]{\tiny propertyScan1D}};
\end{tikzpicture}
\end{tutscheme}
\tutwhat{NRTL methylAcetate-methanol regressed to measured bubble temperatures (ThermoML locates, the article sources), held out against an independent study, drawn against both beside ideal and UNIFAC}
\begin{tutnarrative}
AN NRTL PAIR FROM PUBLISHED MEASUREMENTS, END TO END.

WHY THIS SYSTEM. the methyl acetate reactive-distillation system, where the catalogue already carries a UNIQUAC pair (Popken 2000) -- an NRTL pair fitted to independent measurements puts a second model beside it.

THE CHAIN, EACH STEP A COMMAND (DEV.md C16; record docs/design/binary-pairs-from-open-measurements.md): 1. bin/choupo-thermoml sync the NIST/TRC ThermoML Archive (ThermoML.v2020-09-30.tgz, doi 10.18434/mds2-2422, sha256-verified) 2. bin/choupo-thermoml identity CAS $\to$ InChIKey for the catalogue, two resolvers that must agree 3. bin/choupo-thermoml vle-index every binary VLE block, searchable 4. bin/choupo-thermoml vle-search methylAcetate methanol --geometry bubbleT 5. bin/choupo-thermoml extract-vle <file> --block N --interior-only $\to$ the two files under constant/experiments/ 6. choupoProps (this case) fitParameters $\to$ a proposal record 7. promotion (a human act) $\to$ data/standards/parameters/NRTL/methanol-methylAcetate.dat 8. bin/choupo-import the seal mirrors that record here

THE EVIDENCE. FIT: doi 10.1021/je600518s (27 interior points, 100 kPa). HELD OUT: doi 10.1016/j.fluid.2015.04.018 (19 interior points, 101.3 kPa) -- a different article, so the held-out AAD is a validation, not a slice of the fit. Each file names its article, its archive file and block, the sha256 of the archive file, the two compounds' CAS and InChIKey, and the per-point expanded uncertainty; each says `reviewStatus transcribedNotCheckedAgainstArticle` because nobody has read the numbers back against the paper yet, and the run announces it. Pure-component endpoints are set aside (they measure the vapour-pressure record, not the pair) and kept in each file under `selection`.

THE FIT. NRTL with a\_ij = a\_ji = 0 and alpha = 0.30 fixed: TWO parameters (b\_ij, b\_ji), because bubble temperatures over the ranges below do not identify more (fit\_nrtl\_ab in the methanol + water case asks that question directly). Each point is weighted by 1/U, U its OWN stated expanded uncertainty (`weighting uncertainty;`), so a point measured to 0.1 K is not outvoted by one measured to 1 K. The held-out AAD is plain kelvin, unweighted.

WHAT THE RUN FOUND (pinned in `expected`; measured, not argued). fit\_nrtl: b\_ij = 182.9 K, b\_ji = 185.6 K, in-sample rms 0.029 K; fit\_nrtl\_start2 lands on the same pair. HELD OUT (101.3 kPa, another laboratory): AAD 0.15 K (0.045 \%) against a band of 0.089 \% $\to$ `validated`; the pair is the catalogue record NRTL/methanol-methylAcetate.dat, beside the UNIQUAC record the catalogue already carries for the same system. Overlay on the held-out set: NRTL 0.15 K, UNIFAC 0.22 K, ideal 4.4 K. A third study was set aside BEFORE fitting: the archive's j.fluid.2015.12.029 block repeats every (x, P, T) row twice, and `extract-vle` refuses a duplicated observation rather than keep one copy on a guess.

WHAT IS NOT CLAIMED. The pair is correct at these pressures and in this frame (ideal-gas vapour, the catalogue's Antoine records); nothing here tests it at high pressure, against LLE, or with another vapour model.

\end{tutnarrative}
\begin{tutgolden}
\goldrow{diag}{fit\_nrtl}{aad\_heldout\_K}{0.150139370506}
\goldrow{diag}{fit\_nrtl}{aad\_heldout\_pct}{0.0454327520711}
\goldrow{diag}{fit\_nrtl}{chi2}{0.7380738708}
\goldrow{diag}{fit\_nrtl}{chi2\_in\_sample}{0.7380738708}
\goldrow{diag}{fit\_nrtl}{chi2\_reduced}{0.029522954832}
\goldrow{diag}{fit\_nrtl}{cond\_JtJ}{64.7492559018}
\goldrow{diag}{fit\_nrtl}{converged}{1}
\goldrow{diag}{fit\_nrtl}{corr\_computed}{1}
\goldrow{diag}{fit\_nrtl}{dof}{25}
\goldrow{diag}{fit\_nrtl}{identifiable}{1}
\goldrow{diag}{fit\_nrtl}{invertible}{1}
\goldrow{diag}{fit\_nrtl}{iter}{5}
\goldrow{diag}{fit\_nrtl}{max\_abs\_corr}{0.991685096395}
\goldrow{diag}{fit\_nrtl}{max\_abs\_resid}{0.0690603448025}
\goldmore{70}
\end{tutgolden}

\clearpage
\tutentry{curate09\_nrtl\_propanol\_water}{tutorials/props/curation/curate09\_nrtl\_propanol\_water}
\begin{tutscheme}
\begin{tikzpicture}[>={Stealth[length=2mm]}, every node/.style={font=\scriptsize}]
  \node[draw, rounded corners=2pt, align=center, minimum width=3.1cm, minimum height=0.8cm, fill=black!4] at (0.0,-0.0) {\textbf{fit\_nrtl}\\[-1pt]{\tiny fitParameters}};
  \node[draw, rounded corners=2pt, align=center, minimum width=3.1cm, minimum height=0.8cm, fill=black!4] at (3.6,-0.0) {\textbf{fit\_nrtl\_start2}\\[-1pt]{\tiny fitParameters}};
  \node[draw, rounded corners=2pt, align=center, minimum width=3.1cm, minimum height=0.8cm, fill=black!4] at (7.2,-0.0) {\textbf{txy\_ideal}\\[-1pt]{\tiny propertyScan1D}};
  \node[draw, rounded corners=2pt, align=center, minimum width=3.1cm, minimum height=0.8cm, fill=black!4] at (0.0,-1.3) {\textbf{txy\_unifac}\\[-1pt]{\tiny propertyScan1D}};
  \node[draw, rounded corners=2pt, align=center, minimum width=3.1cm, minimum height=0.8cm, fill=black!4] at (3.6,-1.3) {\textbf{txy\_nrtl}\\[-1pt]{\tiny propertyScan1D}};
  \node[draw, rounded corners=2pt, align=center, minimum width=3.1cm, minimum height=0.8cm, fill=black!4] at (7.2,-1.3) {\textbf{txy\_ideal\_fitP}\\[-1pt]{\tiny propertyScan1D}};
  \node[draw, rounded corners=2pt, align=center, minimum width=3.1cm, minimum height=0.8cm, fill=black!4] at (0.0,-2.6) {\textbf{txy\_unifac\_fitP}\\[-1pt]{\tiny propertyScan1D}};
  \node[draw, rounded corners=2pt, align=center, minimum width=3.1cm, minimum height=0.8cm, fill=black!4] at (3.6,-2.6) {\textbf{txy\_nrtl\_fitP}\\[-1pt]{\tiny propertyScan1D}};
\end{tikzpicture}
\end{tutscheme}
\tutwhat{NRTL 1Propanol-water regressed to measured bubble temperatures (ThermoML locates, the article sources), held out against an independent study, drawn against both beside ideal and UNIFAC}
\begin{tutnarrative}
AN NRTL PAIR FROM PUBLISHED MEASUREMENTS, END TO END.

WHY THIS SYSTEM. the third water + alcohol pair of the catalogue after methanol and isopropanol, and the one with the homologue's minimum-boiling azeotrope (both studies bottom out near 360.5 K): with methanol + water and isopropanol + water it completes the C1-C3 series a first distillation course compares. The archive holds five low-pressure studies of it, the most of any candidate.

THE CHAIN, EACH STEP A COMMAND (DEV.md C16; record docs/design/binary-pairs-from-open-measurements.md): 1. bin/choupo-thermoml sync the NIST/TRC ThermoML Archive (ThermoML.v2020-09-30.tgz, doi 10.18434/mds2-2422, sha256-verified) 2. bin/choupo-thermoml identity CAS $\to$ InChIKey for the catalogue, two resolvers that must agree 3. bin/choupo-thermoml vle-index every binary VLE block, searchable 4. bin/choupo-thermoml vle-search 1Propanol water --geometry bubbleT 5. bin/choupo-thermoml extract-vle <file> --block N --interior-only $\to$ the two files under constant/experiments/ 6. choupoProps (this case) fitParameters $\to$ a proposal record 7. promotion (a human act) $\to$ data/standards/parameters/NRTL/1Propanol-water.dat ONLY if the held-out verdict is `validated` 8. bin/choupo-import the seal mirrors that record here

THE EVIDENCE. FIT: doi 10.1021/je800436u (26 interior points, 100 kPa). HELD OUT: doi 10.1016/j.jct.2017.06.014 (19 interior points, 101.3 kPa). Chosen BEFORE any fit by the slice-2 partition rule: the FIT set is the low-pressure bubble-temperature study with the most interior points; the HELD-OUT set is the next largest that shares NO author with it (the archive carries no affiliations, so a shared author is the dependence the rule can check). Each file names its article, its archive file and block, the sha256 of the archive file, the two compounds' CAS and InChIKey, and the per-point expanded uncertainty; each says `reviewStatus transcribedNotCheckedAgainstArticle` because nobody has read the numbers back against the paper yet, and the run announces it. Pure-component endpoints are set aside (they measure the vapour-pressure record, not the pair) and kept in each file under `selection`.

THE FIT. NRTL with a\_ij = a\_ji = 0 and alpha = 0.30 fixed: TWO parameters (b\_ij, b\_ji), because bubble temperatures over the ranges below do not identify more (fit\_nrtl\_ab in the methanol + water case asks that question directly). Each point is weighted by 1/U, U its OWN stated expanded uncertainty (`weighting uncertainty;`), so a point measured to 0.1 K is not outvoted by one measured to 1 K. The held-out AAD is plain kelvin, unweighted.

WHAT THE RUN FOUND (pinned in `expected`; measured, not argued). fit\_nrtl: b\_ij = -25.1 K, b\_ji = 979.8 K, in-sample rms 0.52 K; fit\_nrtl\_start2 lands on the same pair. The in-sample reduced chi2 is 8.9: two parameters cannot follow the fit set within its own stated U (0.09-0.89 K) -- recorded, not hidden by a third parameter. HELD OUT (101.3 kPa, no author in common): AAD 0.21 K (0.057 \%) against a band of 0.26 \% $\to$ `validated`; the pair is the catalogue record NRTL/1Propanol-water.dat. Overlay on the held-out set: NRTL 0.21 K, UNIFAC 0.36 K, ideal 8.9 K. PHASE STABILITY, checked beside the verdict because the verdict cannot see it: one liquid across the validity span (360.5-368.6 K), but BELOW about 355.5 K the pair splits a liquid that is miscible in all proportions. The record says so, and the validity block is the guard.

WHAT IS NOT CLAIMED. The pair is correct at these pressures and in this frame (ideal-gas vapour, the catalogue's Antoine records); nothing here tests it at high pressure, against LLE, or with another vapour model.

\end{tutnarrative}
\begin{tutgolden}
\goldrow{diag}{fit\_nrtl}{aad\_heldout\_K}{0.205309666797}
\goldrow{diag}{fit\_nrtl}{aad\_heldout\_pct}{0.0565382893231}
\goldrow{diag}{fit\_nrtl}{chi2}{213.201326383}
\goldrow{diag}{fit\_nrtl}{chi2\_in\_sample}{213.201326383}
\goldrow{diag}{fit\_nrtl}{chi2\_reduced}{8.8833885993}
\goldrow{diag}{fit\_nrtl}{cond\_JtJ}{45.7035038408}
\goldrow{diag}{fit\_nrtl}{converged}{1}
\goldrow{diag}{fit\_nrtl}{corr\_computed}{1}
\goldrow{diag}{fit\_nrtl}{dof}{24}
\goldrow{diag}{fit\_nrtl}{identifiable}{1}
\goldrow{diag}{fit\_nrtl}{invertible}{1}
\goldrow{diag}{fit\_nrtl}{iter}{12}
\goldrow{diag}{fit\_nrtl}{max\_abs\_corr}{0.919864189194}
\goldrow{diag}{fit\_nrtl}{max\_abs\_resid}{1.46662361676}
\goldmore{70}
\end{tutgolden}

\clearpage
\tutentry{curate100\_nrtl\_acetone\_methanol}{tutorials/props/curation/curate100\_nrtl\_acetone\_methanol}
\begin{tutscheme}
\begin{tikzpicture}[>={Stealth[length=2mm]}, every node/.style={font=\scriptsize}]
  \node[draw, rounded corners=2pt, align=center, minimum width=3.1cm, minimum height=0.8cm, fill=black!4] at (0.0,-0.0) {\textbf{fit\_nrtl}\\[-1pt]{\tiny fitParameters}};
  \node[draw, rounded corners=2pt, align=center, minimum width=3.1cm, minimum height=0.8cm, fill=black!4] at (3.6,-0.0) {\textbf{fit\_nrtl\_start2}\\[-1pt]{\tiny fitParameters}};
  \node[draw, rounded corners=2pt, align=center, minimum width=3.1cm, minimum height=0.8cm, fill=black!4] at (7.2,-0.0) {\textbf{txy\_ideal}\\[-1pt]{\tiny propertyScan1D}};
  \node[draw, rounded corners=2pt, align=center, minimum width=3.1cm, minimum height=0.8cm, fill=black!4] at (0.0,-1.3) {\textbf{txy\_unifac}\\[-1pt]{\tiny propertyScan1D}};
  \node[draw, rounded corners=2pt, align=center, minimum width=3.1cm, minimum height=0.8cm, fill=black!4] at (3.6,-1.3) {\textbf{txy\_nrtl}\\[-1pt]{\tiny propertyScan1D}};
\end{tikzpicture}
\end{tutscheme}
\tutwhat{NRTL acetone-methanol regressed by the curation pipeline to every laboratory's measured bubble temperatures (ThermoML locates, the articles source), accepted at max |dT| < 1 K}
\begin{tutnarrative}
AN NRTL PAIR FROM EVERY LABORATORY'S MEASUREMENTS.

WRITTEN BY bin/choupo-curate pair acetone methanol (DEV.md 4c C284, the bulk run; record docs/design/binary-pairs-from-open-measurements.md section 11).

THE DECLARATION -- WRITTEN AND COMMITTED BEFORE ANY FIT RAN (several declarations of the bulk run share one commit).

IDENTITY. acetone (CAS 67-64-1, InChIKey CSCPPACGZOOCGX-UHFFFAOYSA-N) + methanol (CAS 67-56-1, InChIKey OKKJLVBELUTLKV-UHFFFAOYSA-N), joined to the archive through identity.json -- never by name. CLASS (C284): both records are volatile molecular liquids (no role nonvolatile/solute, no dissociatesTo, no charge, no solid, a vapour-pressure model), and every measured temperature lies below both critical temperatures.

NO MEASURED SERIES IN THIS TREE. constant/archiveDatasets names each dataset (DOI, NIST URL, article sha256, extraction selector, extract sha256); `bin/choupo-curate fetch` writes the points into constant/archiveExtracts/, which is gitignored.

STUDIES: every bubble-temperature block of the pair at or below 110 kPa with a declared uncertainty (7 block(s)). FITTED, ALL OF THEM (7): 10.1021/je5007373 block 0 32 pts 101.3 kPa 327.96-337.05 K U 0.34-0.35 K 10.1016/j.fluid.2006.09.007 block 2 30 pts 100 kPa 328.15-336.95 K U 0.22-0.23 K 10.1021/je0603170 block 0 30 pts 100 kPa 328.13-336.90 K U 0.09-0.11 K 10.1016/j.fluid.2017.03.014 block 0 19 pts 101.3 kPa 328.15-335.55 K U 0.09-0.11 K 10.1016/j.fluid.2017.11.035 block 2 19 pts 101 kPa 328.20-336.10 K U 0.09-0.1 K 10.1016/j.fluid.2014.10.011 block 1 17 pts 101.3 kPa 328.38-336.25 K U 0.09-0.1 K 10.1021/je049688b block 2 9 pts 101 kPa 328.55-334.85 K U 0.25-0.58 K SET ASIDE BY THE EXTRACTION RULES, BY NAME: (none)

THE RULE (DEV.md 4c C283 amended, C284). Every laboratory above is fitted; none is held out (the C16 partition and its band are retired as the acceptance test). NRTL, a\_ij = a\_ji = 0, alpha = 0.30 fixed: two parameters (b\_ij, b\_ji); each point weighted by 1/U, U its own stated expanded uncertainty. Objective twoStage (least squares, then minimax; T\_bubble absolute, in K). OUTLIERS: after a least-squares stage over every point, a point with |T\_model - T| > 3 K is discarded BY NAME (`discard` on its evidence entry, listed below) and the fit repeated without it. ACCEPTED when max |T\_model - T| over the KEPT points is below 1 K. C283's 5 \% relative clause judges composition; this fit has no composition residual, so only the temperature clause applies. A pair that splits ONE liquid inside its validity span is not promoted (arm f). Two starts: (0, 0) and (800, -400).

OUTLIERS DISCARDED AFTER THE LEAST-SQUARES STAGE: @@OUTLIERS-BEGIN@@ none: every one of the 156 points lies within 3 K after the least-squares stage. @@OUTLIERS-END@@

@@RESULT-BEGIN@@ WHAT THE RUN FOUND (pinned in `expected`; measured, not argued). fit\_nrtl: b\_ij = 103.6 K, b\_ji = 112.5 K over 156 kept points of 7 laboratories (0 discarded): max |dT| 0.622 K, mean 0.163 K against the acceptance 1 K $\to$ accepted. fit\_nrtl\_start2 lands on the same pair (103.6, 112.5). PROMOTED on 2026-11-10 into data/standards/parameters/NRTL/acetone-methanol.dat as reviewStatus autoPromoted, humanReviewed false, because: max |dT| 0.622 K < 1 K over the 156 kept points, one liquid across the validity span @@RESULT-END@@

WHAT IS NOT CLAIMED. The pair is correct at these pressures and in this frame (ideal-gas vapour, the catalogue's own vapour-pressure records); the measurements were transcribed by the archive and have NOT been read back against the articles (the CHECKLIST that `fetch` writes is the aid).

\end{tutnarrative}
\begin{tutgolden}
\goldrow{diag}{fit\_nrtl}{bubble\_aad\_fit}{0.163269428411}
\goldrow{diag}{fit\_nrtl}{bubble\_max\_fit}{0.622096208455}
\goldrow{diag}{fit\_nrtl}{chi2}{408.267446898}
\goldrow{diag}{fit\_nrtl}{chi2\_in\_sample}{408.267446898}
\goldrow{diag}{fit\_nrtl}{chi2\_reduced}{2.65108731752}
\goldrow{diag}{fit\_nrtl}{cond\_JtJ}{191.68150614}
\goldrow{diag}{fit\_nrtl}{converged}{1}
\goldrow{diag}{fit\_nrtl}{corr\_computed}{1}
\goldrow{diag}{fit\_nrtl}{dof}{154}
\goldrow{diag}{fit\_nrtl}{identifiable}{1}
\goldrow{diag}{fit\_nrtl}{invertible}{1}
\goldrow{diag}{fit\_nrtl}{iter}{6}
\goldrow{diag}{fit\_nrtl}{max\_abs\_corr}{0.997227444503}
\goldrow{diag}{fit\_nrtl}{max\_abs\_resid}{0.622096208455}
\goldmore{98}
\end{tutgolden}

\clearpage
\tutentry{curate101\_nrtl\_ethanol\_224Trimethylpentane}{tutorials/props/curation/curate101\_nrtl\_ethanol\_224Trimethylpentane}
\begin{tutscheme}
\begin{tikzpicture}[>={Stealth[length=2mm]}, every node/.style={font=\scriptsize}]
  \node[draw, rounded corners=2pt, align=center, minimum width=3.1cm, minimum height=0.8cm, fill=black!4] at (0.0,-0.0) {\textbf{fit\_nrtl}\\[-1pt]{\tiny fitParameters}};
  \node[draw, rounded corners=2pt, align=center, minimum width=3.1cm, minimum height=0.8cm, fill=black!4] at (3.6,-0.0) {\textbf{fit\_nrtl\_start2}\\[-1pt]{\tiny fitParameters}};
  \node[draw, rounded corners=2pt, align=center, minimum width=3.1cm, minimum height=0.8cm, fill=black!4] at (7.2,-0.0) {\textbf{txy\_ideal}\\[-1pt]{\tiny propertyScan1D}};
  \node[draw, rounded corners=2pt, align=center, minimum width=3.1cm, minimum height=0.8cm, fill=black!4] at (0.0,-1.3) {\textbf{txy\_unifac}\\[-1pt]{\tiny propertyScan1D}};
  \node[draw, rounded corners=2pt, align=center, minimum width=3.1cm, minimum height=0.8cm, fill=black!4] at (3.6,-1.3) {\textbf{txy\_nrtl}\\[-1pt]{\tiny propertyScan1D}};
\end{tikzpicture}
\end{tutscheme}
\tutwhat{NRTL ethanol-224Trimethylpentane regressed by the curation pipeline to every laboratory's measured bubble temperatures (ThermoML locates, the articles source), accepted at max |dT| < 1 K}
\begin{tutnarrative}
AN NRTL PAIR FROM EVERY LABORATORY'S MEASUREMENTS.

WRITTEN BY bin/choupo-curate pair ethanol 224Trimethylpentane (DEV.md 4c C284, the bulk run; record docs/design/binary-pairs-from-open-measurements.md section 11).

THE DECLARATION -- WRITTEN AND COMMITTED BEFORE ANY FIT RAN (several declarations of the bulk run share one commit).

IDENTITY. ethanol (CAS 64-17-5, InChIKey LFQSCWFLJHTTHZ-UHFFFAOYSA-N) + 224Trimethylpentane (CAS 540-84-1, InChIKey NHTMVDHEPJAVLT-UHFFFAOYSA-N), joined to the archive through identity.json -- never by name. CLASS (C284): both records are volatile molecular liquids (no role nonvolatile/solute, no dissociatesTo, no charge, no solid, a vapour-pressure model), and every measured temperature lies below both critical temperatures.

NO MEASURED SERIES IN THIS TREE. constant/archiveDatasets names each dataset (DOI, NIST URL, article sha256, extraction selector, extract sha256); `bin/choupo-curate fetch` writes the points into constant/archiveExtracts/, which is gitignored.

STUDIES: every bubble-temperature block of the pair at or below 110 kPa with a declared uncertainty (5 block(s)). FITTED, ALL OF THEM (5): 10.1021/je300847v block 6 22 pts 101.3 kPa 344.57-347.65 K U 0.32-2.3 K 10.1021/acs.jced.8b00680 block 12 21 pts 101.3 kPa 344.33-355.43 K U 0.43-2.41 K 10.1016/j.fluid.2005.01.007 block 13 19 pts 101.32 kPa 344.44-362.55 K U 0.16-1.93 K 10.1016/j.fluid.2007.06.005 block 10 19 pts 101.3 kPa 344.44-361.27 K U 0.13-1.9 K 10.1016/j.jct.2005.04.006 block 25 6 pts 95.8 kPa 343.95-352.55 K U 1.01-1.01 K SET ASIDE BY THE EXTRACTION RULES, BY NAME: (none)

THE RULE (DEV.md 4c C283 amended, C284). Every laboratory above is fitted; none is held out (the C16 partition and its band are retired as the acceptance test). NRTL, a\_ij = a\_ji = 0, alpha = 0.30 fixed: two parameters (b\_ij, b\_ji); each point weighted by 1/U, U its own stated expanded uncertainty. Objective twoStage (least squares, then minimax; T\_bubble absolute, in K). OUTLIERS: after a least-squares stage over every point, a point with |T\_model - T| > 3 K is discarded BY NAME (`discard` on its evidence entry, listed below) and the fit repeated without it. ACCEPTED when max |T\_model - T| over the KEPT points is below 1 K. C283's 5 \% relative clause judges composition; this fit has no composition residual, so only the temperature clause applies. A pair that splits ONE liquid inside its validity span is not promoted (arm f). Two starts: (0, 0) and (800, -400).

OUTLIERS DISCARDED AFTER THE LEAST-SQUARES STAGE: @@OUTLIERS-BEGIN@@ 2 point(s) of 87: doi 10.1016/j.jct.2005.04.006 row 1: x 0.1186, P 95.800 kPa, T 352.55 K, dT -5.88 K; doi 10.1016/j.jct.2005.04.006 row 2: x 0.2876, P 95.800 kPa, T 346.55 K, dT -3.31 K @@OUTLIERS-END@@

@@RESULT-BEGIN@@ WHAT THE RUN FOUND (pinned in `expected`; measured, not argued). fit\_nrtl: b\_ij = 358.3 K, b\_ji = 555.8 K over 85 kept points of 5 laboratories (2 discarded): max |dT| 1.882 K, mean 0.516 K against the acceptance 1 K $\to$ notAccepted. fit\_nrtl\_start2 lands on the same pair (358.3, 555.8). NOT PROMOTABLE: max |dT| 1.882 K over the 85 kept points is not below 1 K @@RESULT-END@@

WHAT IS NOT CLAIMED. The pair is correct at these pressures and in this frame (ideal-gas vapour, the catalogue's own vapour-pressure records); the measurements were transcribed by the archive and have NOT been read back against the articles (the CHECKLIST that `fetch` writes is the aid).

\end{tutnarrative}
\begin{tutgolden}
\goldrow{diag}{fit\_nrtl}{bubble\_aad\_fit}{0.515998186857}
\goldrow{diag}{fit\_nrtl}{bubble\_max\_fit}{1.88160296803}
\goldrow{diag}{fit\_nrtl}{chi2}{352.63557199}
\goldrow{diag}{fit\_nrtl}{chi2\_in\_sample}{352.63557199}
\goldrow{diag}{fit\_nrtl}{chi2\_reduced}{4.24862134928}
\goldrow{diag}{fit\_nrtl}{cond\_JtJ}{24.1692021985}
\goldrow{diag}{fit\_nrtl}{converged}{1}
\goldrow{diag}{fit\_nrtl}{corr\_computed}{1}
\goldrow{diag}{fit\_nrtl}{dof}{83}
\goldrow{diag}{fit\_nrtl}{identifiable}{1}
\goldrow{diag}{fit\_nrtl}{invertible}{1}
\goldrow{diag}{fit\_nrtl}{iter}{5}
\goldrow{diag}{fit\_nrtl}{max\_abs\_corr}{0.961597410325}
\goldrow{diag}{fit\_nrtl}{max\_abs\_resid}{1.88160296803}
\goldmore{86}
\end{tutgolden}

\clearpage
\tutentry{curate102\_nrtl\_ethanol\_dimethylCarbonate}{tutorials/props/curation/curate102\_nrtl\_ethanol\_dimethylCarbonate}
\begin{tutscheme}
\begin{tikzpicture}[>={Stealth[length=2mm]}, every node/.style={font=\scriptsize}]
  \node[draw, rounded corners=2pt, align=center, minimum width=3.1cm, minimum height=0.8cm, fill=black!4] at (0.0,-0.0) {\textbf{fit\_nrtl}\\[-1pt]{\tiny fitParameters}};
  \node[draw, rounded corners=2pt, align=center, minimum width=3.1cm, minimum height=0.8cm, fill=black!4] at (3.6,-0.0) {\textbf{fit\_nrtl\_start2}\\[-1pt]{\tiny fitParameters}};
  \node[draw, rounded corners=2pt, align=center, minimum width=3.1cm, minimum height=0.8cm, fill=black!4] at (7.2,-0.0) {\textbf{txy\_ideal}\\[-1pt]{\tiny propertyScan1D}};
  \node[draw, rounded corners=2pt, align=center, minimum width=3.1cm, minimum height=0.8cm, fill=black!4] at (0.0,-1.3) {\textbf{txy\_nrtl}\\[-1pt]{\tiny propertyScan1D}};
\end{tikzpicture}
\end{tutscheme}
\tutwhat{NRTL ethanol-dimethylCarbonate regressed by the curation pipeline to every laboratory's measured bubble temperatures (ThermoML locates, the articles source), accepted at max |dT| < 1 K}
\begin{tutnarrative}
AN NRTL PAIR FROM EVERY LABORATORY'S MEASUREMENTS.

WRITTEN BY bin/choupo-curate pair ethanol dimethylCarbonate (DEV.md 4c C284, the bulk run; record docs/design/binary-pairs-from-open-measurements.md section 11).

THE DECLARATION -- WRITTEN AND COMMITTED BEFORE ANY FIT RAN (several declarations of the bulk run share one commit).

IDENTITY. ethanol (CAS 64-17-5, InChIKey LFQSCWFLJHTTHZ-UHFFFAOYSA-N) + dimethylCarbonate (CAS 616-38-6, InChIKey IEJIGPNLZYLLBP-UHFFFAOYSA-N), joined to the archive through identity.json -- never by name. CLASS (C284): both records are volatile molecular liquids (no role nonvolatile/solute, no dissociatesTo, no charge, no solid, a vapour-pressure model), and every measured temperature lies below both critical temperatures.

NO MEASURED SERIES IN THIS TREE. constant/archiveDatasets names each dataset (DOI, NIST URL, article sha256, extraction selector, extract sha256); `bin/choupo-curate fetch` writes the points into constant/archiveExtracts/, which is gitignored.

STUDIES: every bubble-temperature block of the pair at or below 110 kPa with a declared uncertainty (4 block(s)). FITTED, ALL OF THEM (4): 10.1021/je0601309 block 9 95 pts 5 isobars 40-93.32 kPa 324.96-356.51 K U 0.19-0.29 K 10.1021/acs.jced.9b00414 block 5 20 pts 101.3 kPa 348.69-362.31 K U 0.26-0.51 K 10.1016/j.fluid.2018.05.024 block 0 19 pts 101.3 kPa 348.15-360.25 K U 0.22-0.24 K 10.1021/je030131q block 23 9 pts 95.8 kPa 346.65-348.55 K U 0.32-0.32 K SET ASIDE BY THE EXTRACTION RULES, BY NAME: (none)

THE RULE (DEV.md 4c C283 amended, C284). Every laboratory above is fitted; none is held out (the C16 partition and its band are retired as the acceptance test). NRTL, a\_ij = a\_ji = 0, alpha = 0.30 fixed: two parameters (b\_ij, b\_ji); each point weighted by 1/U, U its own stated expanded uncertainty. Objective twoStage (least squares, then minimax; T\_bubble absolute, in K). OUTLIERS: after a least-squares stage over every point, a point with |T\_model - T| > 3 K is discarded BY NAME (`discard` on its evidence entry, listed below) and the fit repeated without it. ACCEPTED when max |T\_model - T| over the KEPT points is below 1 K. C283's 5 \% relative clause judges composition; this fit has no composition residual, so only the temperature clause applies. A pair that splits ONE liquid inside its validity span is not promoted (arm f). Two starts: (0, 0) and (800, -400).

OUTLIERS DISCARDED AFTER THE LEAST-SQUARES STAGE: @@OUTLIERS-BEGIN@@ none: every one of the 143 points lies within 3 K after the least-squares stage. @@OUTLIERS-END@@

@@RESULT-BEGIN@@ WHAT THE RUN FOUND (pinned in `expected`; measured, not argued). fit\_nrtl: b\_ij = 507.9 K, b\_ji = 12.1 K over 143 kept points of 4 laboratories (0 discarded): max |dT| 2.600 K, mean 0.544 K against the acceptance 1 K $\to$ notAccepted. fit\_nrtl\_start2 lands on A DIFFERENT pair (507.6, 12.3). NOT PROMOTABLE: max |dT| 2.600 K over the 143 kept points is not below 1 K @@RESULT-END@@

WHAT IS NOT CLAIMED. The pair is correct at these pressures and in this frame (ideal-gas vapour, the catalogue's own vapour-pressure records); the measurements were transcribed by the archive and have NOT been read back against the articles (the CHECKLIST that `fetch` writes is the aid).

\end{tutnarrative}
\begin{tutgolden}
\goldrow{diag}{fit\_nrtl}{bubble\_aad\_fit}{0.544392906236}
\goldrow{diag}{fit\_nrtl}{bubble\_max\_fit}{2.59982030708}
\goldrow{diag}{fit\_nrtl}{chi2}{1196.93830259}
\goldrow{diag}{fit\_nrtl}{chi2\_in\_sample}{1196.93830259}
\goldrow{diag}{fit\_nrtl}{chi2\_reduced}{8.48892413181}
\goldrow{diag}{fit\_nrtl}{cond\_JtJ}{27.0679142243}
\goldrow{diag}{fit\_nrtl}{converged}{1}
\goldrow{diag}{fit\_nrtl}{corr\_computed}{1}
\goldrow{diag}{fit\_nrtl}{dof}{141}
\goldrow{diag}{fit\_nrtl}{identifiable}{1}
\goldrow{diag}{fit\_nrtl}{invertible}{1}
\goldrow{diag}{fit\_nrtl}{iter}{5}
\goldrow{diag}{fit\_nrtl}{max\_abs\_corr}{0.982024505789}
\goldrow{diag}{fit\_nrtl}{max\_abs\_resid}{2.59982030708}
\goldmore{50}
\end{tutgolden}

\clearpage
\tutentry{curate103\_nrtl\_diisopropylEther\_isopropanol}{tutorials/props/curation/curate103\_nrtl\_diisopropylEther\_isopropanol}
\begin{tutscheme}
\begin{tikzpicture}[>={Stealth[length=2mm]}, every node/.style={font=\scriptsize}]
  \node[draw, rounded corners=2pt, align=center, minimum width=3.1cm, minimum height=0.8cm, fill=black!4] at (0.0,-0.0) {\textbf{fit\_nrtl}\\[-1pt]{\tiny fitParameters}};
  \node[draw, rounded corners=2pt, align=center, minimum width=3.1cm, minimum height=0.8cm, fill=black!4] at (3.6,-0.0) {\textbf{fit\_nrtl\_start2}\\[-1pt]{\tiny fitParameters}};
  \node[draw, rounded corners=2pt, align=center, minimum width=3.1cm, minimum height=0.8cm, fill=black!4] at (7.2,-0.0) {\textbf{txy\_ideal}\\[-1pt]{\tiny propertyScan1D}};
  \node[draw, rounded corners=2pt, align=center, minimum width=3.1cm, minimum height=0.8cm, fill=black!4] at (0.0,-1.3) {\textbf{txy\_nrtl}\\[-1pt]{\tiny propertyScan1D}};
\end{tikzpicture}
\end{tutscheme}
\tutwhat{NRTL diisopropylEther-isopropanol regressed by the curation pipeline to every laboratory's measured bubble temperatures (ThermoML locates, the articles source), accepted at max |dT| < 1 K}
\begin{tutnarrative}
AN NRTL PAIR FROM EVERY LABORATORY'S MEASUREMENTS.

WRITTEN BY bin/choupo-curate pair diisopropylEther isopropanol (DEV.md 4c C284, the bulk run; record docs/design/binary-pairs-from-open-measurements.md section 11).

THE DECLARATION -- WRITTEN AND COMMITTED BEFORE ANY FIT RAN (several declarations of the bulk run share one commit).

IDENTITY. diisopropylEther (CAS 108-20-3, InChIKey ZAFNJMIOTHYJRJ-UHFFFAOYSA-N) + isopropanol (CAS 67-63-0, InChIKey KFZMGEQAYNKOFK-UHFFFAOYSA-N), joined to the archive through identity.json -- never by name. CLASS (C284): both records are volatile molecular liquids (no role nonvolatile/solute, no dissociatesTo, no charge, no solid, a vapour-pressure model), and every measured temperature lies below both critical temperatures.

NO MEASURED SERIES IN THIS TREE. constant/archiveDatasets names each dataset (DOI, NIST URL, article sha256, extraction selector, extract sha256); `bin/choupo-curate fetch` writes the points into constant/archiveExtracts/, which is gitignored.

STUDIES: every bubble-temperature block of the pair at or below 110 kPa with a declared uncertainty (4 block(s)). FITTED, ALL OF THEM (4): 10.1016/j.fluid.2014.09.014 block 8 67 pts 4 isobars 53.33-94.79 kPa 320.37-345.39 K U 0.31-1.63 K 10.1016/j.fluid.2007.09.014 block 11 37 pts 2 isobars 30-101.3 kPa 307.18-352.60 K U 0.14-0.25 K 10.1021/je300994y block 9 19 pts 101.3 kPa 339.13-352.94 K U 0.1-0.61 K 10.1016/j.jct.2018.11.013 block 12 13 pts 101.3 kPa 339.28-352.73 K U 0.72-1.14 K SET ASIDE BY THE EXTRACTION RULES, BY NAME: (none)

THE RULE (DEV.md 4c C283 amended, C284). Every laboratory above is fitted; none is held out (the C16 partition and its band are retired as the acceptance test). NRTL, a\_ij = a\_ji = 0, alpha = 0.30 fixed: two parameters (b\_ij, b\_ji); each point weighted by 1/U, U its own stated expanded uncertainty. Objective twoStage (least squares, then minimax; T\_bubble absolute, in K). OUTLIERS: after a least-squares stage over every point, a point with |T\_model - T| > 3 K is discarded BY NAME (`discard` on its evidence entry, listed below) and the fit repeated without it. ACCEPTED when max |T\_model - T| over the KEPT points is below 1 K. C283's 5 \% relative clause judges composition; this fit has no composition residual, so only the temperature clause applies. A pair that splits ONE liquid inside its validity span is not promoted (arm f). Two starts: (0, 0) and (800, -400).

OUTLIERS DISCARDED AFTER THE LEAST-SQUARES STAGE: @@OUTLIERS-BEGIN@@ 10 point(s) of 136: doi 10.1016/j.fluid.2014.09.014 row 3: x 0.0510, P 79.990 kPa, T 341.61 K, dT +3.61 K; doi 10.1016/j.fluid.2014.09.014 row 4: x 0.0510, P 94.790 kPa, T 345.39 K, dT +4.13 K; doi 10.1016/j.fluid.2014.09.014 row 5: x 0.0720, P 53.330 kPa, T 330.55 K, dT +3.37 K; doi 10.1016/j.fluid.2014.09.014 row 6: x 0.0720, P 66.670 kPa, T 334.24 K, dT +5.06 K; doi 10.1016/j.fluid.2014.09.014 row 7: x 0.0720, P 94.790 kPa, T 343.55 K, dT +4.64 K; doi 10.1016/j.fluid.2014.09.014 row 8: x 0.1530, P 53.330 kPa, T 326.50 K, dT +3.15 K; doi 10.1016/j.fluid.2014.09.014 row 10: x 0.1530, P 79.990 kPa, T 335.32 K, dT +4.54 K; doi 10.1016/j.fluid.2014.09.014 row 11: x 0.1530, P 94.790 kPa, T 340.38 K, dT +3.95 K; doi 10.1016/j.fluid.2014.09.014 row 12: x 0.1940, P 79.990 kPa, T 335.17 K, dT +3.28 K; doi 10.1016/j.fluid.2014.09.014 row 13: x 0.1940, P 94.790 kPa, T 339.23 K, dT +3.72 K @@OUTLIERS-END@@

@@RESULT-BEGIN@@ WHAT THE RUN FOUND (pinned in `expected`; measured, not argued). fit\_nrtl: b\_ij = 250.5 K, b\_ji = 207.6 K over 126 kept points of 4 laboratories (10 discarded): max |dT| 2.337 K, mean 0.718 K against the acceptance 1 K $\to$ notAccepted. fit\_nrtl\_start2 lands on the same pair (250.5, 207.6). NOT PROMOTABLE: max |dT| 2.337 K over the 126 kept points is not below 1 K @@RESULT-END@@

WHAT IS NOT CLAIMED. The pair is correct at these pressures and in this frame (ideal-gas vapour, the catalogue's own vapour-pressure records); the measurements were transcribed by the archive and have NOT been read back against the articles (the CHECKLIST that `fetch` writes is the aid).

\end{tutnarrative}
\begin{tutgolden}
\goldrow{diag}{fit\_nrtl}{bubble\_aad\_fit}{0.717526141651}
\goldrow{diag}{fit\_nrtl}{bubble\_max\_fit}{2.33719153375}
\goldrow{diag}{fit\_nrtl}{chi2}{934.643731478}
\goldrow{diag}{fit\_nrtl}{chi2\_in\_sample}{934.643731478}
\goldrow{diag}{fit\_nrtl}{chi2\_reduced}{7.5374494474}
\goldrow{diag}{fit\_nrtl}{cond\_JtJ}{37.5874000088}
\goldrow{diag}{fit\_nrtl}{converged}{1}
\goldrow{diag}{fit\_nrtl}{corr\_computed}{1}
\goldrow{diag}{fit\_nrtl}{dof}{124}
\goldrow{diag}{fit\_nrtl}{identifiable}{1}
\goldrow{diag}{fit\_nrtl}{invertible}{1}
\goldrow{diag}{fit\_nrtl}{iter}{5}
\goldrow{diag}{fit\_nrtl}{max\_abs\_corr}{0.986240281782}
\goldrow{diag}{fit\_nrtl}{max\_abs\_resid}{2.33719153375}
\goldmore{50}
\end{tutgolden}

\clearpage
\tutentry{curate104\_nrtl\_acetonitrile\_methanol}{tutorials/props/curation/curate104\_nrtl\_acetonitrile\_methanol}
\begin{tutscheme}
\begin{tikzpicture}[>={Stealth[length=2mm]}, every node/.style={font=\scriptsize}]
  \node[draw, rounded corners=2pt, align=center, minimum width=3.1cm, minimum height=0.8cm, fill=black!4] at (0.0,-0.0) {\textbf{fit\_nrtl}\\[-1pt]{\tiny fitParameters}};
  \node[draw, rounded corners=2pt, align=center, minimum width=3.1cm, minimum height=0.8cm, fill=black!4] at (3.6,-0.0) {\textbf{fit\_nrtl\_start2}\\[-1pt]{\tiny fitParameters}};
  \node[draw, rounded corners=2pt, align=center, minimum width=3.1cm, minimum height=0.8cm, fill=black!4] at (7.2,-0.0) {\textbf{txy\_ideal}\\[-1pt]{\tiny propertyScan1D}};
  \node[draw, rounded corners=2pt, align=center, minimum width=3.1cm, minimum height=0.8cm, fill=black!4] at (0.0,-1.3) {\textbf{txy\_nrtl}\\[-1pt]{\tiny propertyScan1D}};
\end{tikzpicture}
\end{tutscheme}
\tutwhat{NRTL acetonitrile-methanol regressed by the curation pipeline to every laboratory's measured bubble temperatures (ThermoML locates, the articles source), accepted at max |dT| < 1 K}
\begin{tutnarrative}
AN NRTL PAIR FROM EVERY LABORATORY'S MEASUREMENTS.

WRITTEN BY bin/choupo-curate pair acetonitrile methanol (DEV.md 4c C284, the bulk run; record docs/design/binary-pairs-from-open-measurements.md section 11).

THE DECLARATION -- WRITTEN AND COMMITTED BEFORE ANY FIT RAN (several declarations of the bulk run share one commit).

IDENTITY. acetonitrile (CAS 75-05-8, InChIKey WEVYAHXRMPXWCK-UHFFFAOYSA-N) + methanol (CAS 67-56-1, InChIKey OKKJLVBELUTLKV-UHFFFAOYSA-N), joined to the archive through identity.json -- never by name. CLASS (C284): both records are volatile molecular liquids (no role nonvolatile/solute, no dissociatesTo, no charge, no solid, a vapour-pressure model), and every measured temperature lies below both critical temperatures.

NO MEASURED SERIES IN THIS TREE. constant/archiveDatasets names each dataset (DOI, NIST URL, article sha256, extraction selector, extract sha256); `bin/choupo-curate fetch` writes the points into constant/archiveExtracts/, which is gitignored.

STUDIES: every bubble-temperature block of the pair at or below 110 kPa with a declared uncertainty (4 block(s)). FITTED, ALL OF THEM (4): 10.1016/j.fluid.2018.08.009 block 0 14 pts 101.3 kPa 337.24-347.91 K U 0.13-0.19 K 10.1021/acs.jced.5b00796 block 2 14 pts 101.32 kPa 336.55-347.85 K U 0.22-0.31 K (reports x of the other member: written as x[acetonitrile] = 1 - x by --x-of) 10.1016/j.jct.2016.05.011 block 6 13 pts 101.3 kPa 336.57-351.02 K U 0.31-0.42 K (reports x of the other member: written as x[acetonitrile] = 1 - x by --x-of) 10.1016/j.fluid.2012.12.003 block 7 10 pts 101.3 kPa 336.96-350.19 K U 0.08-0.1 K (reports x of the other member: written as x[acetonitrile] = 1 - x by --x-of) SET ASIDE BY THE EXTRACTION RULES, BY NAME: (none)

THE RULE (DEV.md 4c C283 amended, C284). Every laboratory above is fitted; none is held out (the C16 partition and its band are retired as the acceptance test). NRTL, a\_ij = a\_ji = 0, alpha = 0.30 fixed: two parameters (b\_ij, b\_ji); each point weighted by 1/U, U its own stated expanded uncertainty. Objective twoStage (least squares, then minimax; T\_bubble absolute, in K). OUTLIERS: after a least-squares stage over every point, a point with |T\_model - T| > 3 K is discarded BY NAME (`discard` on its evidence entry, listed below) and the fit repeated without it. ACCEPTED when max |T\_model - T| over the KEPT points is below 1 K. C283's 5 \% relative clause judges composition; this fit has no composition residual, so only the temperature clause applies. A pair that splits ONE liquid inside its validity span is not promoted (arm f). Two starts: (0, 0) and (800, -400).

OUTLIERS DISCARDED AFTER THE LEAST-SQUARES STAGE: @@OUTLIERS-BEGIN@@ none: every one of the 51 points lies within 3 K after the least-squares stage. @@OUTLIERS-END@@

@@RESULT-BEGIN@@ WHAT THE RUN FOUND (pinned in `expected`; measured, not argued). fit\_nrtl: b\_ij = 351.5 K, b\_ji = -20.1 K over 51 kept points of 4 laboratories (0 discarded): max |dT| 0.991 K, mean 0.452 K against the acceptance 1 K $\to$ accepted. fit\_nrtl\_start2 lands on the same pair (351.5, -20.1). PROMOTED on 2026-11-10 into data/standards/parameters/NRTL/acetonitrile-methanol.dat as reviewStatus autoPromoted, humanReviewed false, because: max |dT| 0.991 K < 1 K over the 51 kept points, one liquid across the validity span @@RESULT-END@@

WHAT IS NOT CLAIMED. The pair is correct at these pressures and in this frame (ideal-gas vapour, the catalogue's own vapour-pressure records); the measurements were transcribed by the archive and have NOT been read back against the articles (the CHECKLIST that `fetch` writes is the aid).

\end{tutnarrative}
\begin{tutgolden}
\goldrow{diag}{fit\_nrtl}{bubble\_aad\_fit}{0.452337135612}
\goldrow{diag}{fit\_nrtl}{bubble\_max\_fit}{0.991013327368}
\goldrow{diag}{fit\_nrtl}{chi2}{481.716129653}
\goldrow{diag}{fit\_nrtl}{chi2\_in\_sample}{481.716129653}
\goldrow{diag}{fit\_nrtl}{chi2\_reduced}{9.8309414215}
\goldrow{diag}{fit\_nrtl}{cond\_JtJ}{49.2124268801}
\goldrow{diag}{fit\_nrtl}{converged}{1}
\goldrow{diag}{fit\_nrtl}{corr\_computed}{1}
\goldrow{diag}{fit\_nrtl}{dof}{49}
\goldrow{diag}{fit\_nrtl}{identifiable}{1}
\goldrow{diag}{fit\_nrtl}{invertible}{1}
\goldrow{diag}{fit\_nrtl}{iter}{4}
\goldrow{diag}{fit\_nrtl}{max\_abs\_corr}{0.989992227888}
\goldrow{diag}{fit\_nrtl}{max\_abs\_resid}{0.991013327368}
\goldmore{74}
\end{tutgolden}

\clearpage
\tutentry{curate105\_nrtl\_water\_triethyleneGlycol}{tutorials/props/curation/curate105\_nrtl\_water\_triethyleneGlycol}
\begin{tutscheme}
\begin{tikzpicture}[>={Stealth[length=2mm]}, every node/.style={font=\scriptsize}]
  \node[draw, rounded corners=2pt, align=center, minimum width=3.1cm, minimum height=0.8cm, fill=black!4] at (0.0,-0.0) {\textbf{fit\_nrtl}\\[-1pt]{\tiny fitParameters}};
  \node[draw, rounded corners=2pt, align=center, minimum width=3.1cm, minimum height=0.8cm, fill=black!4] at (3.6,-0.0) {\textbf{fit\_nrtl\_start2}\\[-1pt]{\tiny fitParameters}};
  \node[draw, rounded corners=2pt, align=center, minimum width=3.1cm, minimum height=0.8cm, fill=black!4] at (7.2,-0.0) {\textbf{txy\_ideal}\\[-1pt]{\tiny propertyScan1D}};
  \node[draw, rounded corners=2pt, align=center, minimum width=3.1cm, minimum height=0.8cm, fill=black!4] at (0.0,-1.3) {\textbf{txy\_nrtl}\\[-1pt]{\tiny propertyScan1D}};
\end{tikzpicture}
\end{tutscheme}
\tutwhat{NRTL water-triethyleneGlycol regressed by the curation pipeline to every laboratory's measured bubble temperatures (ThermoML locates, the articles source), accepted at max |dT| < 1 K}
\begin{tutnarrative}
AN NRTL PAIR FROM EVERY LABORATORY'S MEASUREMENTS.

WRITTEN BY bin/choupo-curate pair water triethyleneGlycol (DEV.md 4c C284, the bulk run; record docs/design/binary-pairs-from-open-measurements.md section 11).

THE DECLARATION -- WRITTEN AND COMMITTED BEFORE ANY FIT RAN (several declarations of the bulk run share one commit).

IDENTITY. water (CAS 7732-18-5, InChIKey XLYOFNOQVPJJNP-UHFFFAOYSA-N) + triethyleneGlycol (CAS 112-27-6, InChIKey ZIBGPFATKBEMQZ-UHFFFAOYSA-N), joined to the archive through identity.json -- never by name. CLASS (C284): both records are volatile molecular liquids (no role nonvolatile/solute, no dissociatesTo, no charge, no solid, a vapour-pressure model), and every measured temperature lies below both critical temperatures.

NO MEASURED SERIES IN THIS TREE. constant/archiveDatasets names each dataset (DOI, NIST URL, article sha256, extraction selector, extract sha256); `bin/choupo-curate fetch` writes the points into constant/archiveExtracts/, which is gitignored.

STUDIES: every bubble-temperature block of the pair at or below 110 kPa with a declared uncertainty (3 block(s)). FITTED, ALL OF THEM (3): 10.1016/j.fluid.2015.07.030 block 6 52 pts 4 isobars 53.33-95.99 kPa 357.46-542.62 K U 0.12-0.39 K 10.1021/acs.jced.7b00945 block 5 51 pts 3 isobars 50-100 kPa 355.04-420.64 K U 2.03-2.13 K 10.1021/je800675u block 0 18 pts 85 kPa 371.55-421.25 K U 2.7-2.7 K SET ASIDE BY THE EXTRACTION RULES, BY NAME: (none)

THE RULE (DEV.md 4c C283 amended, C284). Every laboratory above is fitted; none is held out (the C16 partition and its band are retired as the acceptance test). NRTL, a\_ij = a\_ji = 0, alpha = 0.30 fixed: two parameters (b\_ij, b\_ji); each point weighted by 1/U, U its own stated expanded uncertainty. Objective twoStage (least squares, then minimax; T\_bubble absolute, in K). OUTLIERS: after a least-squares stage over every point, a point with |T\_model - T| > 3 K is discarded BY NAME (`discard` on its evidence entry, listed below) and the fit repeated without it. ACCEPTED when max |T\_model - T| over the KEPT points is below 1 K. C283's 5 \% relative clause judges composition; this fit has no composition residual, so only the temperature clause applies. A pair that splits ONE liquid inside its validity span is not promoted (arm f). Two starts: (0, 0) and (800, -400).

OUTLIERS DISCARDED AFTER THE LEAST-SQUARES STAGE: @@OUTLIERS-BEGIN@@ 25 point(s) of 121: doi 10.1016/j.fluid.2015.07.030 row 1: x 0.0500, P 53.330 kPa, T 511.31 K, dT -5.30 K; doi 10.1016/j.fluid.2015.07.030 row 2: x 0.0500, P 66.660 kPa, T 518.45 K, dT -4.13 K; doi 10.1016/j.fluid.2015.07.030 row 3: x 0.0500, P 79.990 kPa, T 529.77 K, dT -8.41 K; doi 10.1016/j.fluid.2015.07.030 row 4: x 0.0500, P 95.990 kPa, T 542.62 K, dT -13.97 K; doi 10.1016/j.fluid.2015.07.030 row 6: x 0.0990, P 66.660 kPa, T 492.68 K, dT -3.64 K; doi 10.1016/j.fluid.2015.07.030 row 8: x 0.0990, P 95.990 kPa, T 511.26 K, dT -8.27 K; doi 10.1016/j.fluid.2015.07.030 row 11: x 0.1510, P 79.990 kPa, T 468.35 K, dT +4.59 K; doi 10.1016/j.fluid.2015.07.030 row 15: x 0.2001, P 79.990 kPa, T 449.00 K, dT +5.88 K; doi 10.1016/j.fluid.2015.07.030 row 16: x 0.2001, P 95.990 kPa, T 456.67 K, dT +4.88 K; doi 10.1016/j.fluid.2015.07.030 row 19: x 0.3003, P 79.990 kPa, T 421.06 K, dT +5.86 K; doi 10.1016/j.fluid.2015.07.030 row 20: x 0.3003, P 95.990 kPa, T 426.62 K, dT +6.55 K; doi 10.1016/j.fluid.2015.07.030 row 23: x 0.4011, P 79.990 kPa, T 403.47 K, dT +4.13 K; doi 10.1016/j.fluid.2015.07.030 row 24: x 0.4011, P 95.990 kPa, T 408.98 K, dT +4.52 K; doi 10.1021/acs.jced.7b00945 row 1: x 0.3390, P 70.000 kPa, T 410.80 K, dT +3.49 K; doi 10.1021/acs.jced.7b00945 row 3: x 0.3670, P 50.000 kPa, T 395.52 K, dT +3.16 K; doi 10.1021/je800675u row 1: x 0.1170, P 85.000 kPa, T 421.25 K, dT +68.55 K; doi 10.1021/je800675u row 2: x 0.1400, P 85.000 kPa, T 417.15 K, dT +62.59 K; doi 10.1021/je800675u row 3: x 0.1610, P 85.000 kPa, T 416.25 K, dT +55.00 K; doi 10.1021/je800675u row 4: x 0.1760, P 85.000 kPa, T 414.75 K, dT +50.82 K; doi 10.1021/je800675u row 5: x 0.1820, P 85.000 kPa, T 413.65 K, dT +49.74 K; doi 10.1021/je800675u row 6: x 0.2420, P 85.000 kPa, T 413.25 K, dT +30.78 K; doi 10.1021/je800675u row 7: x 0.2510, P 85.000 kPa, T 412.55 K, dT +28.93 K; doi 10.1021/je800675u row 8: x 0.3250, P 85.000 kPa, T 403.95 K, dT +19.58 K; doi 10.1021/je800675u row 9: x 0.3730, P 85.000 kPa, T 400.85 K, dT +13.41 K; doi 10.1021/je800675u row 10: x 0.4100, P 85.000 kPa, T 397.95 K, dT +10.20 K @@OUTLIERS-END@@

@@RESULT-BEGIN@@ WHAT THE RUN FOUND (pinned in `expected`; measured, not argued). fit\_nrtl: b\_ij = -928.6 K, b\_ji = 609.0 K over 96 kept points of 3 laboratories (25 discarded): max |dT| 3.833 K, mean 1.069 K against the acceptance 1 K $\to$ notAccepted. fit\_nrtl\_start2 lands on A DIFFERENT pair (3000.0, -1057.2). NOT PROMOTABLE: max |dT| 3.833 K over the 96 kept points is not below 1 K @@RESULT-END@@

WHAT IS NOT CLAIMED. The pair is correct at these pressures and in this frame (ideal-gas vapour, the catalogue's own vapour-pressure records); the measurements were transcribed by the archive and have NOT been read back against the articles (the CHECKLIST that `fetch` writes is the aid).

\end{tutnarrative}
\begin{tutgolden}
\goldrow{diag}{fit\_nrtl}{bubble\_aad\_fit}{1.06906291836}
\goldrow{diag}{fit\_nrtl}{bubble\_max\_fit}{3.8325257483}
\goldrow{diag}{fit\_nrtl}{chi2}{2608.16411374}
\goldrow{diag}{fit\_nrtl}{chi2\_in\_sample}{2608.16411374}
\goldrow{diag}{fit\_nrtl}{chi2\_reduced}{27.746426742}
\goldrow{diag}{fit\_nrtl}{cond\_JtJ}{43479.0876858}
\goldrow{diag}{fit\_nrtl}{converged}{1}
\goldrow{diag}{fit\_nrtl}{corr\_computed}{1}
\goldrow{diag}{fit\_nrtl}{dof}{94}
\goldrow{diag}{fit\_nrtl}{identifiable}{0}
\goldrow{diag}{fit\_nrtl}{invertible}{1}
\goldrow{diag}{fit\_nrtl}{iter}{18}
\goldrow{diag}{fit\_nrtl}{max\_abs\_corr}{0.99992202318}
\goldrow{diag}{fit\_nrtl}{max\_abs\_resid}{3.8325257483}
\goldmore{50}
\end{tutgolden}

\clearpage
\tutentry{curate106\_nrtl\_2Butanol\_dimethylCarbonate}{tutorials/props/curation/curate106\_nrtl\_2Butanol\_dimethylCarbonate}
\begin{tutscheme}
\begin{tikzpicture}[>={Stealth[length=2mm]}, every node/.style={font=\scriptsize}]
  \node[draw, rounded corners=2pt, align=center, minimum width=3.1cm, minimum height=0.8cm, fill=black!4] at (0.0,-0.0) {\textbf{fit\_nrtl}\\[-1pt]{\tiny fitParameters}};
  \node[draw, rounded corners=2pt, align=center, minimum width=3.1cm, minimum height=0.8cm, fill=black!4] at (3.6,-0.0) {\textbf{fit\_nrtl\_start2}\\[-1pt]{\tiny fitParameters}};
  \node[draw, rounded corners=2pt, align=center, minimum width=3.1cm, minimum height=0.8cm, fill=black!4] at (7.2,-0.0) {\textbf{txy\_ideal}\\[-1pt]{\tiny propertyScan1D}};
  \node[draw, rounded corners=2pt, align=center, minimum width=3.1cm, minimum height=0.8cm, fill=black!4] at (0.0,-1.3) {\textbf{txy\_nrtl}\\[-1pt]{\tiny propertyScan1D}};
\end{tikzpicture}
\end{tutscheme}
\tutwhat{NRTL 2Butanol-dimethylCarbonate regressed by the curation pipeline to every laboratory's measured bubble temperatures (ThermoML locates, the articles source), accepted at max |dT| < 1 K}
\begin{tutnarrative}
AN NRTL PAIR FROM EVERY LABORATORY'S MEASUREMENTS.

WRITTEN BY bin/choupo-curate pair 2Butanol dimethylCarbonate (DEV.md 4c C284, the bulk run; record docs/design/binary-pairs-from-open-measurements.md section 11).

THE DECLARATION -- WRITTEN AND COMMITTED BEFORE ANY FIT RAN (several declarations of the bulk run share one commit).

IDENTITY. 2Butanol (CAS 78-92-2, InChIKey BTANRVKWQNVYAZ-UHFFFAOYSA-N) + dimethylCarbonate (CAS 616-38-6, InChIKey IEJIGPNLZYLLBP-UHFFFAOYSA-N), joined to the archive through identity.json -- never by name. CLASS (C284): both records are volatile molecular liquids (no role nonvolatile/solute, no dissociatesTo, no charge, no solid, a vapour-pressure model), and every measured temperature lies below both critical temperatures.

NO MEASURED SERIES IN THIS TREE. constant/archiveDatasets names each dataset (DOI, NIST URL, article sha256, extraction selector, extract sha256); `bin/choupo-curate fetch` writes the points into constant/archiveExtracts/, which is gitignored.

STUDIES: every bubble-temperature block of the pair at or below 110 kPa with a declared uncertainty (3 block(s)). FITTED, ALL OF THEM (3): 10.1016/j.jct.2011.08.006 block 10 70 pts 5 isobars 40-93.32 kPa 335.83-366.43 K U 0.12-5.4 K 10.1016/j.jct.2004.09.009 block 18 32 pts 101.3 kPa 361.93-371.33 K U 0.34-0.35 K (reports x of the other member: written as x[2Butanol] = 1 - x by --x-of) 10.1021/je030131q block 28 9 pts 95.8 kPa 360.25-365.75 K U 0.32-0.32 K (reports x of the other member: written as x[2Butanol] = 1 - x by --x-of) SET ASIDE BY THE EXTRACTION RULES, BY NAME: (none)

THE RULE (DEV.md 4c C283 amended, C284). Every laboratory above is fitted; none is held out (the C16 partition and its band are retired as the acceptance test). NRTL, a\_ij = a\_ji = 0, alpha = 0.30 fixed: two parameters (b\_ij, b\_ji); each point weighted by 1/U, U its own stated expanded uncertainty. Objective twoStage (least squares, then minimax; T\_bubble absolute, in K). OUTLIERS: after a least-squares stage over every point, a point with |T\_model - T| > 3 K is discarded BY NAME (`discard` on its evidence entry, listed below) and the fit repeated without it. ACCEPTED when max |T\_model - T| over the KEPT points is below 1 K. C283's 5 \% relative clause judges composition; this fit has no composition residual, so only the temperature clause applies. A pair that splits ONE liquid inside its validity span is not promoted (arm f). Two starts: (0, 0) and (800, -400).

OUTLIERS DISCARDED AFTER THE LEAST-SQUARES STAGE: @@OUTLIERS-BEGIN@@ none: every one of the 111 points lies within 3 K after the least-squares stage. @@OUTLIERS-END@@

@@RESULT-BEGIN@@ WHAT THE RUN FOUND (pinned in `expected`; measured, not argued). fit\_nrtl: b\_ij = 169.8 K, b\_ji = 144.1 K over 111 kept points of 3 laboratories (0 discarded): max |dT| 1.377 K, mean 0.368 K against the acceptance 1 K $\to$ notAccepted. fit\_nrtl\_start2 lands on the same pair (169.8, 144.1). NOT PROMOTABLE: max |dT| 1.377 K over the 111 kept points is not below 1 K @@RESULT-END@@

WHAT IS NOT CLAIMED. The pair is correct at these pressures and in this frame (ideal-gas vapour, the catalogue's own vapour-pressure records); the measurements were transcribed by the archive and have NOT been read back against the articles (the CHECKLIST that `fetch` writes is the aid).

\end{tutnarrative}
\begin{tutgolden}
\goldrow{diag}{fit\_nrtl}{bubble\_aad\_fit}{0.368170648763}
\goldrow{diag}{fit\_nrtl}{bubble\_max\_fit}{1.37696244795}
\goldrow{diag}{fit\_nrtl}{chi2}{471.776033354}
\goldrow{diag}{fit\_nrtl}{chi2\_in\_sample}{471.776033354}
\goldrow{diag}{fit\_nrtl}{chi2\_reduced}{4.32822048948}
\goldrow{diag}{fit\_nrtl}{cond\_JtJ}{93.7245999322}
\goldrow{diag}{fit\_nrtl}{converged}{1}
\goldrow{diag}{fit\_nrtl}{corr\_computed}{1}
\goldrow{diag}{fit\_nrtl}{dof}{109}
\goldrow{diag}{fit\_nrtl}{identifiable}{1}
\goldrow{diag}{fit\_nrtl}{invertible}{1}
\goldrow{diag}{fit\_nrtl}{iter}{4}
\goldrow{diag}{fit\_nrtl}{max\_abs\_corr}{0.994707535031}
\goldrow{diag}{fit\_nrtl}{max\_abs\_resid}{1.37696244795}
\goldmore{50}
\end{tutgolden}

\clearpage
\tutentry{curate107\_nrtl\_methanol\_isopropanol}{tutorials/props/curation/curate107\_nrtl\_methanol\_isopropanol}
\begin{tutscheme}
\begin{tikzpicture}[>={Stealth[length=2mm]}, every node/.style={font=\scriptsize}]
  \node[draw, rounded corners=2pt, align=center, minimum width=3.1cm, minimum height=0.8cm, fill=black!4] at (0.0,-0.0) {\textbf{fit\_nrtl}\\[-1pt]{\tiny fitParameters}};
  \node[draw, rounded corners=2pt, align=center, minimum width=3.1cm, minimum height=0.8cm, fill=black!4] at (3.6,-0.0) {\textbf{fit\_nrtl\_start2}\\[-1pt]{\tiny fitParameters}};
  \node[draw, rounded corners=2pt, align=center, minimum width=3.1cm, minimum height=0.8cm, fill=black!4] at (7.2,-0.0) {\textbf{txy\_ideal}\\[-1pt]{\tiny propertyScan1D}};
  \node[draw, rounded corners=2pt, align=center, minimum width=3.1cm, minimum height=0.8cm, fill=black!4] at (0.0,-1.3) {\textbf{txy\_unifac}\\[-1pt]{\tiny propertyScan1D}};
  \node[draw, rounded corners=2pt, align=center, minimum width=3.1cm, minimum height=0.8cm, fill=black!4] at (3.6,-1.3) {\textbf{txy\_nrtl}\\[-1pt]{\tiny propertyScan1D}};
\end{tikzpicture}
\end{tutscheme}
\tutwhat{NRTL methanol-isopropanol regressed by the curation pipeline to every laboratory's measured bubble temperatures (ThermoML locates, the articles source), accepted at max |dT| < 1 K}
\begin{tutnarrative}
AN NRTL PAIR FROM EVERY LABORATORY'S MEASUREMENTS.

WRITTEN BY bin/choupo-curate pair methanol isopropanol (DEV.md 4c C284, the bulk run; record docs/design/binary-pairs-from-open-measurements.md section 11).

THE DECLARATION -- WRITTEN AND COMMITTED BEFORE ANY FIT RAN (several declarations of the bulk run share one commit).

IDENTITY. methanol (CAS 67-56-1, InChIKey OKKJLVBELUTLKV-UHFFFAOYSA-N) + isopropanol (CAS 67-63-0, InChIKey KFZMGEQAYNKOFK-UHFFFAOYSA-N), joined to the archive through identity.json -- never by name. CLASS (C284): both records are volatile molecular liquids (no role nonvolatile/solute, no dissociatesTo, no charge, no solid, a vapour-pressure model), and every measured temperature lies below both critical temperatures.

NO MEASURED SERIES IN THIS TREE. constant/archiveDatasets names each dataset (DOI, NIST URL, article sha256, extraction selector, extract sha256); `bin/choupo-curate fetch` writes the points into constant/archiveExtracts/, which is gitignored.

STUDIES: every bubble-temperature block of the pair at or below 110 kPa with a declared uncertainty (3 block(s)). FITTED, ALL OF THEM (3): 10.1016/j.fluid.2015.07.023 block 6 77 pts 4 isobars 53.33-94.79 kPa 322.45-353.74 K U 0.16-0.23 K 10.1021/je060476j block 6 18 pts 101.3 kPa 339.20-354.22 K U 0.14-0.15 K 10.1021/je201349k block 3 16 pts 101 kPa 339.07-354.69 K U 0.18-0.19 K SET ASIDE BY THE EXTRACTION RULES, BY NAME: (none)

THE RULE (DEV.md 4c C283 amended, C284). Every laboratory above is fitted; none is held out (the C16 partition and its band are retired as the acceptance test). NRTL, a\_ij = a\_ji = 0, alpha = 0.30 fixed: two parameters (b\_ij, b\_ji); each point weighted by 1/U, U its own stated expanded uncertainty. Objective twoStage (least squares, then minimax; T\_bubble absolute, in K). OUTLIERS: after a least-squares stage over every point, a point with |T\_model - T| > 3 K is discarded BY NAME (`discard` on its evidence entry, listed below) and the fit repeated without it. ACCEPTED when max |T\_model - T| over the KEPT points is below 1 K. C283's 5 \% relative clause judges composition; this fit has no composition residual, so only the temperature clause applies. A pair that splits ONE liquid inside its validity span is not promoted (arm f). Two starts: (0, 0) and (800, -400).

OUTLIERS DISCARDED AFTER THE LEAST-SQUARES STAGE: @@OUTLIERS-BEGIN@@ none: every one of the 111 points lies within 3 K after the least-squares stage. @@OUTLIERS-END@@

@@RESULT-BEGIN@@ WHAT THE RUN FOUND (pinned in `expected`; measured, not argued). fit\_nrtl: b\_ij = 82.7 K, b\_ji = -60.7 K over 111 kept points of 3 laboratories (0 discarded): max |dT| 1.736 K, mean 0.631 K against the acceptance 1 K $\to$ notAccepted. fit\_nrtl\_start2 lands on A DIFFERENT pair (85.7, -63.3). NOT PROMOTABLE: max |dT| 1.736 K over the 111 kept points is not below 1 K @@RESULT-END@@

WHAT IS NOT CLAIMED. The pair is correct at these pressures and in this frame (ideal-gas vapour, the catalogue's own vapour-pressure records); the measurements were transcribed by the archive and have NOT been read back against the articles (the CHECKLIST that `fetch` writes is the aid).

\end{tutnarrative}
\begin{tutgolden}
\goldrow{diag}{fit\_nrtl}{bubble\_aad\_fit}{0.631094376454}
\goldrow{diag}{fit\_nrtl}{bubble\_max\_fit}{1.73640726967}
\goldrow{diag}{fit\_nrtl}{chi2}{2132.77523556}
\goldrow{diag}{fit\_nrtl}{chi2\_in\_sample}{2132.77523556}
\goldrow{diag}{fit\_nrtl}{chi2\_reduced}{19.5667452804}
\goldrow{diag}{fit\_nrtl}{cond\_JtJ}{14792.0108691}
\goldrow{diag}{fit\_nrtl}{converged}{1}
\goldrow{diag}{fit\_nrtl}{corr\_computed}{1}
\goldrow{diag}{fit\_nrtl}{dof}{109}
\goldrow{diag}{fit\_nrtl}{identifiable}{0}
\goldrow{diag}{fit\_nrtl}{invertible}{1}
\goldrow{diag}{fit\_nrtl}{iter}{11}
\goldrow{diag}{fit\_nrtl}{max\_abs\_corr}{0.999966199683}
\goldrow{diag}{fit\_nrtl}{max\_abs\_resid}{1.73640726967}
\goldmore{53}
\end{tutgolden}

\clearpage
\tutentry{curate108\_nrtl\_1Pentanol\_nNonane}{tutorials/props/curation/curate108\_nrtl\_1Pentanol\_nNonane}
\begin{tutscheme}
\begin{tikzpicture}[>={Stealth[length=2mm]}, every node/.style={font=\scriptsize}]
  \node[draw, rounded corners=2pt, align=center, minimum width=3.1cm, minimum height=0.8cm, fill=black!4] at (0.0,-0.0) {\textbf{fit\_nrtl}\\[-1pt]{\tiny fitParameters}};
  \node[draw, rounded corners=2pt, align=center, minimum width=3.1cm, minimum height=0.8cm, fill=black!4] at (3.6,-0.0) {\textbf{fit\_nrtl\_start2}\\[-1pt]{\tiny fitParameters}};
  \node[draw, rounded corners=2pt, align=center, minimum width=3.1cm, minimum height=0.8cm, fill=black!4] at (7.2,-0.0) {\textbf{txy\_ideal}\\[-1pt]{\tiny propertyScan1D}};
  \node[draw, rounded corners=2pt, align=center, minimum width=3.1cm, minimum height=0.8cm, fill=black!4] at (0.0,-1.3) {\textbf{txy\_unifac}\\[-1pt]{\tiny propertyScan1D}};
  \node[draw, rounded corners=2pt, align=center, minimum width=3.1cm, minimum height=0.8cm, fill=black!4] at (3.6,-1.3) {\textbf{txy\_nrtl}\\[-1pt]{\tiny propertyScan1D}};
\end{tikzpicture}
\end{tutscheme}
\tutwhat{NRTL 1Pentanol-nNonane regressed by the curation pipeline to every laboratory's measured bubble temperatures (ThermoML locates, the articles source), accepted at max |dT| < 1 K}
\begin{tutnarrative}
AN NRTL PAIR FROM EVERY LABORATORY'S MEASUREMENTS.

WRITTEN BY bin/choupo-curate pair 1Pentanol nNonane (DEV.md 4c C284, the bulk run; record docs/design/binary-pairs-from-open-measurements.md section 11).

THE DECLARATION -- WRITTEN AND COMMITTED BEFORE ANY FIT RAN (several declarations of the bulk run share one commit).

IDENTITY. 1Pentanol (CAS 71-41-0, InChIKey AMQJEAYHLZJPGS-UHFFFAOYSA-N) + nNonane (CAS 111-84-2, InChIKey BKIMMITUMNQMOS-UHFFFAOYSA-N), joined to the archive through identity.json -- never by name. CLASS (C284): both records are volatile molecular liquids (no role nonvolatile/solute, no dissociatesTo, no charge, no solid, a vapour-pressure model), and every measured temperature lies below both critical temperatures.

NO MEASURED SERIES IN THIS TREE. constant/archiveDatasets names each dataset (DOI, NIST URL, article sha256, extraction selector, extract sha256); `bin/choupo-curate fetch` writes the points into constant/archiveExtracts/, which is gitignored.

STUDIES: every bubble-temperature block of the pair at or below 110 kPa with a declared uncertainty (3 block(s)). FITTED, ALL OF THEM (3): 10.1021/acs.jced.8b00680 block 14 44 pts 40 kPa 377.93-390.17 K U 0.55-3.21 K 10.1021/je0502905 block 12 22 pts 3 isobars 26.66-79.99 kPa 368.52-406.20 K U 0.25-0.29 K 10.1021/je060237w block 13 9 pts 101.32 kPa 405.95-412.69 K U 0.13-0.15 K SET ASIDE BY THE EXTRACTION RULES, BY NAME: (none)

THE RULE (DEV.md 4c C283 amended, C284). Every laboratory above is fitted; none is held out (the C16 partition and its band are retired as the acceptance test). NRTL, a\_ij = a\_ji = 0, alpha = 0.30 fixed: two parameters (b\_ij, b\_ji); each point weighted by 1/U, U its own stated expanded uncertainty. Objective twoStage (least squares, then minimax; T\_bubble absolute, in K). OUTLIERS: after a least-squares stage over every point, a point with |T\_model - T| > 3 K is discarded BY NAME (`discard` on its evidence entry, listed below) and the fit repeated without it. ACCEPTED when max |T\_model - T| over the KEPT points is below 1 K. C283's 5 \% relative clause judges composition; this fit has no composition residual, so only the temperature clause applies. A pair that splits ONE liquid inside its validity span is not promoted (arm f). Two starts: (0, 0) and (800, -400).

OUTLIERS DISCARDED AFTER THE LEAST-SQUARES STAGE: @@OUTLIERS-BEGIN@@ none: every one of the 75 points lies within 3 K after the least-squares stage. @@OUTLIERS-END@@

@@RESULT-BEGIN@@ WHAT THE RUN FOUND (pinned in `expected`; measured, not argued). fit\_nrtl: b\_ij = 261.4 K, b\_ji = 282.7 K over 75 kept points of 3 laboratories (0 discarded): max |dT| 2.122 K, mean 0.755 K against the acceptance 1 K $\to$ notAccepted. fit\_nrtl\_start2 lands on the same pair (261.4, 282.7). NOT PROMOTABLE: max |dT| 2.122 K over the 75 kept points is not below 1 K @@RESULT-END@@

WHAT IS NOT CLAIMED. The pair is correct at these pressures and in this frame (ideal-gas vapour, the catalogue's own vapour-pressure records); the measurements were transcribed by the archive and have NOT been read back against the articles (the CHECKLIST that `fetch` writes is the aid).

\end{tutnarrative}
\begin{tutgolden}
\goldrow{diag}{fit\_nrtl}{bubble\_aad\_fit}{0.754543982262}
\goldrow{diag}{fit\_nrtl}{bubble\_max\_fit}{2.12159936043}
\goldrow{diag}{fit\_nrtl}{chi2}{202.222113729}
\goldrow{diag}{fit\_nrtl}{chi2\_in\_sample}{202.222113729}
\goldrow{diag}{fit\_nrtl}{chi2\_reduced}{2.77016594149}
\goldrow{diag}{fit\_nrtl}{cond\_JtJ}{47.5272782049}
\goldrow{diag}{fit\_nrtl}{converged}{1}
\goldrow{diag}{fit\_nrtl}{corr\_computed}{1}
\goldrow{diag}{fit\_nrtl}{dof}{73}
\goldrow{diag}{fit\_nrtl}{identifiable}{1}
\goldrow{diag}{fit\_nrtl}{invertible}{1}
\goldrow{diag}{fit\_nrtl}{iter}{3}
\goldrow{diag}{fit\_nrtl}{max\_abs\_corr}{0.989082638251}
\goldrow{diag}{fit\_nrtl}{max\_abs\_resid}{2.12159936043}
\goldmore{59}
\end{tutgolden}

\clearpage
\tutentry{curate109\_nrtl\_water\_ethyleneGlycol}{tutorials/props/curation/curate109\_nrtl\_water\_ethyleneGlycol}
\begin{tutscheme}
\begin{tikzpicture}[>={Stealth[length=2mm]}, every node/.style={font=\scriptsize}]
  \node[draw, rounded corners=2pt, align=center, minimum width=3.1cm, minimum height=0.8cm, fill=black!4] at (0.0,-0.0) {\textbf{fit\_nrtl}\\[-1pt]{\tiny fitParameters}};
  \node[draw, rounded corners=2pt, align=center, minimum width=3.1cm, minimum height=0.8cm, fill=black!4] at (3.6,-0.0) {\textbf{fit\_nrtl\_start2}\\[-1pt]{\tiny fitParameters}};
  \node[draw, rounded corners=2pt, align=center, minimum width=3.1cm, minimum height=0.8cm, fill=black!4] at (7.2,-0.0) {\textbf{txy\_ideal}\\[-1pt]{\tiny propertyScan1D}};
  \node[draw, rounded corners=2pt, align=center, minimum width=3.1cm, minimum height=0.8cm, fill=black!4] at (0.0,-1.3) {\textbf{txy\_nrtl}\\[-1pt]{\tiny propertyScan1D}};
\end{tikzpicture}
\end{tutscheme}
\tutwhat{NRTL water-ethyleneGlycol regressed by the curation pipeline to every laboratory's measured bubble temperatures (ThermoML locates, the articles source), accepted at max |dT| < 1 K}
\begin{tutnarrative}
AN NRTL PAIR FROM EVERY LABORATORY'S MEASUREMENTS.

WRITTEN BY bin/choupo-curate pair water ethyleneGlycol (DEV.md 4c C284, the bulk run; record docs/design/binary-pairs-from-open-measurements.md section 11).

THE DECLARATION -- WRITTEN AND COMMITTED BEFORE ANY FIT RAN (several declarations of the bulk run share one commit).

IDENTITY. water (CAS 7732-18-5, InChIKey XLYOFNOQVPJJNP-UHFFFAOYSA-N) + ethyleneGlycol (CAS 107-21-1, InChIKey LYCAIKOWRPUZTN-UHFFFAOYSA-N), joined to the archive through identity.json -- never by name. CLASS (C284): both records are volatile molecular liquids (no role nonvolatile/solute, no dissociatesTo, no charge, no solid, a vapour-pressure model), and every measured temperature lies below both critical temperatures.

NO MEASURED SERIES IN THIS TREE. constant/archiveDatasets names each dataset (DOI, NIST URL, article sha256, extraction selector, extract sha256); `bin/choupo-curate fetch` writes the points into constant/archiveExtracts/, which is gitignored.

STUDIES: every bubble-temperature block of the pair at or below 110 kPa with a declared uncertainty (3 block(s)). FITTED, ALL OF THEM (3): 10.1021/acs.jced.7b00945 block 1 49 pts 3 isobars 50-100 kPa 353.64-410.95 K U 2.03-2.11 K 10.1021/je2008704 block 5 19 pts 101.3 kPa 376.05-462.57 K U 0.46-4.49 K 10.1021/acs.jced.6b00264 block 9 6 pts 101.3 kPa 388.93-424.82 K U 0.52-0.74 K SET ASIDE BY THE EXTRACTION RULES, BY NAME: (none)

THE RULE (DEV.md 4c C283 amended, C284). Every laboratory above is fitted; none is held out (the C16 partition and its band are retired as the acceptance test). NRTL, a\_ij = a\_ji = 0, alpha = 0.30 fixed: two parameters (b\_ij, b\_ji); each point weighted by 1/U, U its own stated expanded uncertainty. Objective twoStage (least squares, then minimax; T\_bubble absolute, in K). OUTLIERS: after a least-squares stage over every point, a point with |T\_model - T| > 3 K is discarded BY NAME (`discard` on its evidence entry, listed below) and the fit repeated without it. ACCEPTED when max |T\_model - T| over the KEPT points is below 1 K. C283's 5 \% relative clause judges composition; this fit has no composition residual, so only the temperature clause applies. A pair that splits ONE liquid inside its validity span is not promoted (arm f). Two starts: (0, 0) and (800, -400).

OUTLIERS DISCARDED AFTER THE LEAST-SQUARES STAGE: @@OUTLIERS-BEGIN@@ 1 point(s) of 74: doi 10.1021/acs.jced.7b00945 row 3: x 0.3170, P 100.000 kPa, T 410.95 K, dT -3.01 K @@OUTLIERS-END@@

@@RESULT-BEGIN@@ WHAT THE RUN FOUND (pinned in `expected`; measured, not argued). fit\_nrtl: b\_ij = 350.0 K, b\_ji = -358.9 K over 73 kept points of 3 laboratories (1 discarded): max |dT| 2.584 K, mean 0.714 K against the acceptance 1 K $\to$ notAccepted. fit\_nrtl\_start2 lands on the same pair (350.0, -358.9). NOT PROMOTABLE: max |dT| 2.584 K over the 73 kept points is not below 1 K @@RESULT-END@@

WHAT IS NOT CLAIMED. The pair is correct at these pressures and in this frame (ideal-gas vapour, the catalogue's own vapour-pressure records); the measurements were transcribed by the archive and have NOT been read back against the articles (the CHECKLIST that `fetch` writes is the aid).

\end{tutnarrative}
\begin{tutgolden}
\goldrow{diag}{fit\_nrtl}{bubble\_aad\_fit}{0.713674285229}
\goldrow{diag}{fit\_nrtl}{bubble\_max\_fit}{2.58420587805}
\goldrow{diag}{fit\_nrtl}{chi2}{23.0891958565}
\goldrow{diag}{fit\_nrtl}{chi2\_in\_sample}{23.0891958565}
\goldrow{diag}{fit\_nrtl}{chi2\_reduced}{0.325199941641}
\goldrow{diag}{fit\_nrtl}{cond\_JtJ}{15503.5656057}
\goldrow{diag}{fit\_nrtl}{converged}{1}
\goldrow{diag}{fit\_nrtl}{corr\_computed}{1}
\goldrow{diag}{fit\_nrtl}{dof}{71}
\goldrow{diag}{fit\_nrtl}{identifiable}{0}
\goldrow{diag}{fit\_nrtl}{invertible}{1}
\goldrow{diag}{fit\_nrtl}{iter}{15}
\goldrow{diag}{fit\_nrtl}{max\_abs\_corr}{0.999967750914}
\goldrow{diag}{fit\_nrtl}{max\_abs\_resid}{2.58420587805}
\goldmore{50}
\end{tutgolden}

\clearpage
\tutentry{curate10\_nrtl\_acetonitrile\_water}{tutorials/props/curation/curate10\_nrtl\_acetonitrile\_water}
\begin{tutscheme}
\begin{tikzpicture}[>={Stealth[length=2mm]}, every node/.style={font=\scriptsize}]
  \node[draw, rounded corners=2pt, align=center, minimum width=3.1cm, minimum height=0.8cm, fill=black!4] at (0.0,-0.0) {\textbf{fit\_nrtl}\\[-1pt]{\tiny fitParameters}};
  \node[draw, rounded corners=2pt, align=center, minimum width=3.1cm, minimum height=0.8cm, fill=black!4] at (3.6,-0.0) {\textbf{fit\_nrtl\_start2}\\[-1pt]{\tiny fitParameters}};
  \node[draw, rounded corners=2pt, align=center, minimum width=3.1cm, minimum height=0.8cm, fill=black!4] at (7.2,-0.0) {\textbf{txy\_ideal}\\[-1pt]{\tiny propertyScan1D}};
  \node[draw, rounded corners=2pt, align=center, minimum width=3.1cm, minimum height=0.8cm, fill=black!4] at (0.0,-1.3) {\textbf{txy\_nrtl}\\[-1pt]{\tiny propertyScan1D}};
  \node[draw, rounded corners=2pt, align=center, minimum width=3.1cm, minimum height=0.8cm, fill=black!4] at (3.6,-1.3) {\textbf{txy\_ideal\_fitP}\\[-1pt]{\tiny propertyScan1D}};
  \node[draw, rounded corners=2pt, align=center, minimum width=3.1cm, minimum height=0.8cm, fill=black!4] at (7.2,-1.3) {\textbf{txy\_nrtl\_fitP}\\[-1pt]{\tiny propertyScan1D}};
\end{tikzpicture}
\end{tutscheme}
\tutwhat{NRTL acetonitrile-water regressed to measured bubble temperatures (ThermoML locates, the article sources), held out against an independent study, drawn against both beside ideal mixing}
\begin{tutnarrative}
AN NRTL PAIR FROM PUBLISHED MEASUREMENTS, END TO END.

WHY THIS SYSTEM. a solvent every analytical laboratory recovers, and the textbook system for PRESSURE-SWING distillation: a minimum-boiling azeotrope near 349.6 K at 101.3 kPa (both studies). The catalogue's acetonitrile record carries no UNIFAC group decomposition, so the overlay draws the pair beside ideal mixing only -- the predictive comparison is a gap of the COMPONENT record, said here rather than filled by a decomposition this case would have to invent.

THE CHAIN, EACH STEP A COMMAND (DEV.md C16; record docs/design/binary-pairs-from-open-measurements.md): 1. bin/choupo-thermoml sync the NIST/TRC ThermoML Archive (ThermoML.v2020-09-30.tgz, doi 10.18434/mds2-2422, sha256-verified) 2. bin/choupo-thermoml identity CAS $\to$ InChIKey for the catalogue, two resolvers that must agree 3. bin/choupo-thermoml vle-index every binary VLE block, searchable 4. bin/choupo-thermoml vle-search acetonitrile water --geometry bubbleT 5. bin/choupo-thermoml extract-vle <file> --block N --interior-only $\to$ the two files under constant/experiments/ 6. choupoProps (this case) fitParameters $\to$ a proposal record 7. promotion (a human act) $\to$ data/standards/parameters/NRTL/acetonitrile-water.dat ONLY if the held-out verdict is `validated` 8. bin/choupo-import the seal mirrors that record here

THE EVIDENCE. FIT: doi 10.1021/je400531a (19 interior points, 101.3 kPa). HELD OUT: doi 10.1016/j.jct.2019.07.006 (12 interior points, 101.3 kPa). Chosen BEFORE any fit by the slice-2 partition rule: the FIT set is the low-pressure bubble-temperature study with the most interior points; the HELD-OUT set is the next largest that shares NO author with it (the archive carries no affiliations, so a shared author is the dependence the rule can check). Each file names its article, its archive file and block, the sha256 of the archive file, the two compounds' CAS and InChIKey, and the per-point expanded uncertainty; each says `reviewStatus transcribedNotCheckedAgainstArticle` because nobody has read the numbers back against the paper yet, and the run announces it. Pure-component endpoints are set aside (they measure the vapour-pressure record, not the pair) and kept in each file under `selection`.

THE FIT. NRTL with a\_ij = a\_ji = 0 and alpha = 0.30 fixed: TWO parameters (b\_ij, b\_ji), because bubble temperatures over the ranges below do not identify more (fit\_nrtl\_ab in the methanol + water case asks that question directly). Each point is weighted by 1/U, U its OWN stated expanded uncertainty (`weighting uncertainty;`), so a point measured to 0.1 K is not outvoted by one measured to 1 K. The held-out AAD is plain kelvin, unweighted.

WHAT THE RUN FOUND (pinned in `expected`; measured, not argued). fit\_nrtl: b\_ij = 218.3 K, b\_ji = 616.3 K, in-sample rms 0.52 K; fit\_nrtl\_start2 lands on the same pair. HELD OUT (101.3 kPa, no author in common): AAD 0.50 K (0.137 \%) against a band of 0.72 \% $\to$ `validated`; the pair is the catalogue record NRTL/acetonitrile-water.dat. The band is wide because the held-out study states U95 = 1.69-1.85 K on every point: a held-out set can only validate to the precision it was measured to. Overlay on the held-out set: NRTL 0.50 K, ideal 10.7 K (no UNIFAC: see WHY THIS SYSTEM). PHASE STABILITY, checked beside the verdict because the verdict cannot see it: one liquid across the validity span (350.0-367.6 K), but BELOW about 323.6 K the pair splits a liquid that is miscible at room temperature. The record says so, and the validity block is the guard.

WHAT IS NOT CLAIMED. The pair is correct at these pressures and in this frame (ideal-gas vapour, the catalogue's Antoine records); nothing here tests it at high pressure, against LLE, or with another vapour model.

\end{tutnarrative}
\begin{tutgolden}
\goldrow{diag}{fit\_nrtl}{aad\_heldout\_K}{0.495564366362}
\goldrow{diag}{fit\_nrtl}{aad\_heldout\_pct}{0.136733099834}
\goldrow{diag}{fit\_nrtl}{chi2}{27.5888548619}
\goldrow{diag}{fit\_nrtl}{chi2\_in\_sample}{27.5888548619}
\goldrow{diag}{fit\_nrtl}{chi2\_reduced}{1.6228738154}
\goldrow{diag}{fit\_nrtl}{cond\_JtJ}{15.0626260021}
\goldrow{diag}{fit\_nrtl}{converged}{1}
\goldrow{diag}{fit\_nrtl}{corr\_computed}{1}
\goldrow{diag}{fit\_nrtl}{dof}{17}
\goldrow{diag}{fit\_nrtl}{identifiable}{1}
\goldrow{diag}{fit\_nrtl}{invertible}{1}
\goldrow{diag}{fit\_nrtl}{iter}{6}
\goldrow{diag}{fit\_nrtl}{max\_abs\_corr}{0.940389969107}
\goldrow{diag}{fit\_nrtl}{max\_abs\_resid}{1.15866160924}
\goldmore{58}
\end{tutgolden}

\clearpage
\tutentry{curate110\_nrtl\_2Methyl2Propanol\_water}{tutorials/props/curation/curate110\_nrtl\_2Methyl2Propanol\_water}
\begin{tutscheme}
\begin{tikzpicture}[>={Stealth[length=2mm]}, every node/.style={font=\scriptsize}]
  \node[draw, rounded corners=2pt, align=center, minimum width=3.1cm, minimum height=0.8cm, fill=black!4] at (0.0,-0.0) {\textbf{fit\_nrtl}\\[-1pt]{\tiny fitParameters}};
  \node[draw, rounded corners=2pt, align=center, minimum width=3.1cm, minimum height=0.8cm, fill=black!4] at (3.6,-0.0) {\textbf{fit\_nrtl\_start2}\\[-1pt]{\tiny fitParameters}};
  \node[draw, rounded corners=2pt, align=center, minimum width=3.1cm, minimum height=0.8cm, fill=black!4] at (7.2,-0.0) {\textbf{txy\_ideal}\\[-1pt]{\tiny propertyScan1D}};
  \node[draw, rounded corners=2pt, align=center, minimum width=3.1cm, minimum height=0.8cm, fill=black!4] at (0.0,-1.3) {\textbf{txy\_unifac}\\[-1pt]{\tiny propertyScan1D}};
  \node[draw, rounded corners=2pt, align=center, minimum width=3.1cm, minimum height=0.8cm, fill=black!4] at (3.6,-1.3) {\textbf{txy\_nrtl}\\[-1pt]{\tiny propertyScan1D}};
\end{tikzpicture}
\end{tutscheme}
\tutwhat{NRTL 2Methyl2Propanol-water regressed by the curation pipeline to every laboratory's measured bubble temperatures (ThermoML locates, the articles source), accepted at max |dT| < 1 K}
\begin{tutnarrative}
AN NRTL PAIR FROM EVERY LABORATORY'S MEASUREMENTS.

WRITTEN BY bin/choupo-curate pair 2Methyl2Propanol water (DEV.md 4c C284, the bulk run; record docs/design/binary-pairs-from-open-measurements.md section 11).

THE DECLARATION -- WRITTEN AND COMMITTED BEFORE ANY FIT RAN (several declarations of the bulk run share one commit).

IDENTITY. 2Methyl2Propanol (CAS 75-65-0, InChIKey DKGAVHZHDRPRBM-UHFFFAOYSA-N) + water (CAS 7732-18-5, InChIKey XLYOFNOQVPJJNP-UHFFFAOYSA-N), joined to the archive through identity.json -- never by name. CLASS (C284): both records are volatile molecular liquids (no role nonvolatile/solute, no dissociatesTo, no charge, no solid, a vapour-pressure model), and every measured temperature lies below both critical temperatures.

NO MEASURED SERIES IN THIS TREE. constant/archiveDatasets names each dataset (DOI, NIST URL, article sha256, extraction selector, extract sha256); `bin/choupo-curate fetch` writes the points into constant/archiveExtracts/, which is gitignored.

STUDIES: every bubble-temperature block of the pair at or below 110 kPa with a declared uncertainty (3 block(s)). FITTED, ALL OF THEM (3): 10.1016/j.jct.2019.07.008 block 2 29 pts 100 kPa 352.60-371.10 K U 0.4-0.4 K 10.1021/je501098j block 2 18 pts 101.3 kPa 352.69-356.41 K U 0.33-0.43 K 10.1016/j.jct.2016.03.013 block 2 16 pts 101.3 kPa 352.60-361.70 K U 0.6-0.7 K SET ASIDE BY THE EXTRACTION RULES, BY NAME: (none)

THE RULE (DEV.md 4c C283 amended, C284). Every laboratory above is fitted; none is held out (the C16 partition and its band are retired as the acceptance test). NRTL, a\_ij = a\_ji = 0, alpha = 0.30 fixed: two parameters (b\_ij, b\_ji); each point weighted by 1/U, U its own stated expanded uncertainty. Objective twoStage (least squares, then minimax; T\_bubble absolute, in K). OUTLIERS: after a least-squares stage over every point, a point with |T\_model - T| > 3 K is discarded BY NAME (`discard` on its evidence entry, listed below) and the fit repeated without it. ACCEPTED when max |T\_model - T| over the KEPT points is below 1 K. C283's 5 \% relative clause judges composition; this fit has no composition residual, so only the temperature clause applies. A pair that splits ONE liquid inside its validity span is not promoted (arm f). Two starts: (0, 0) and (800, -400).

OUTLIERS DISCARDED AFTER THE LEAST-SQUARES STAGE: @@OUTLIERS-BEGIN@@ none: every one of the 63 points lies within 3 K after the least-squares stage. @@OUTLIERS-END@@

@@RESULT-BEGIN@@ WHAT THE RUN FOUND (pinned in `expected`; measured, not argued). fit\_nrtl: b\_ij = -121.0 K, b\_ji = 1271.4 K over 63 kept points of 3 laboratories (0 discarded): max |dT| 2.191 K, mean 1.053 K against the acceptance 1 K $\to$ notAccepted. fit\_nrtl\_start2 lands on the same pair (-121.0, 1271.4). NOT PROMOTABLE: max |dT| 2.191 K over the 63 kept points is not below 1 K @@RESULT-END@@

WHAT IS NOT CLAIMED. The pair is correct at these pressures and in this frame (ideal-gas vapour, the catalogue's own vapour-pressure records); the measurements were transcribed by the archive and have NOT been read back against the articles (the CHECKLIST that `fetch` writes is the aid).

\end{tutnarrative}
\begin{tutgolden}
\goldrow{diag}{fit\_nrtl}{bubble\_aad\_fit}{1.05292959532}
\goldrow{diag}{fit\_nrtl}{bubble\_max\_fit}{2.19080265371}
\goldrow{diag}{fit\_nrtl}{chi2}{473.799723936}
\goldrow{diag}{fit\_nrtl}{chi2\_in\_sample}{473.799723936}
\goldrow{diag}{fit\_nrtl}{chi2\_reduced}{7.76720858912}
\goldrow{diag}{fit\_nrtl}{cond\_JtJ}{20.1539910313}
\goldrow{diag}{fit\_nrtl}{converged}{1}
\goldrow{diag}{fit\_nrtl}{corr\_computed}{1}
\goldrow{diag}{fit\_nrtl}{dof}{61}
\goldrow{diag}{fit\_nrtl}{identifiable}{1}
\goldrow{diag}{fit\_nrtl}{invertible}{1}
\goldrow{diag}{fit\_nrtl}{iter}{11}
\goldrow{diag}{fit\_nrtl}{max\_abs\_corr}{0.797913929801}
\goldrow{diag}{fit\_nrtl}{max\_abs\_resid}{2.19080265371}
\goldmore{62}
\end{tutgolden}

\clearpage
\tutentry{curate111\_nrtl\_acetonitrile\_1Propanol}{tutorials/props/curation/curate111\_nrtl\_acetonitrile\_1Propanol}
\begin{tutscheme}
\begin{tikzpicture}[>={Stealth[length=2mm]}, every node/.style={font=\scriptsize}]
  \node[draw, rounded corners=2pt, align=center, minimum width=3.1cm, minimum height=0.8cm, fill=black!4] at (0.0,-0.0) {\textbf{fit\_nrtl}\\[-1pt]{\tiny fitParameters}};
  \node[draw, rounded corners=2pt, align=center, minimum width=3.1cm, minimum height=0.8cm, fill=black!4] at (3.6,-0.0) {\textbf{fit\_nrtl\_start2}\\[-1pt]{\tiny fitParameters}};
  \node[draw, rounded corners=2pt, align=center, minimum width=3.1cm, minimum height=0.8cm, fill=black!4] at (7.2,-0.0) {\textbf{txy\_ideal}\\[-1pt]{\tiny propertyScan1D}};
  \node[draw, rounded corners=2pt, align=center, minimum width=3.1cm, minimum height=0.8cm, fill=black!4] at (0.0,-1.3) {\textbf{txy\_nrtl}\\[-1pt]{\tiny propertyScan1D}};
\end{tikzpicture}
\end{tutscheme}
\tutwhat{NRTL acetonitrile-1Propanol regressed by the curation pipeline to every laboratory's measured bubble temperatures (ThermoML locates, the articles source), accepted at max |dT| < 1 K}
\begin{tutnarrative}
AN NRTL PAIR FROM EVERY LABORATORY'S MEASUREMENTS.

WRITTEN BY bin/choupo-curate pair acetonitrile 1Propanol (DEV.md 4c C284, the bulk run; record docs/design/binary-pairs-from-open-measurements.md section 11).

THE DECLARATION -- WRITTEN AND COMMITTED BEFORE ANY FIT RAN (several declarations of the bulk run share one commit).

IDENTITY. acetonitrile (CAS 75-05-8, InChIKey WEVYAHXRMPXWCK-UHFFFAOYSA-N) + 1Propanol (CAS 71-23-8, InChIKey BDERNNFJNOPAEC-UHFFFAOYSA-N), joined to the archive through identity.json -- never by name. CLASS (C284): both records are volatile molecular liquids (no role nonvolatile/solute, no dissociatesTo, no charge, no solid, a vapour-pressure model), and every measured temperature lies below both critical temperatures.

NO MEASURED SERIES IN THIS TREE. constant/archiveDatasets names each dataset (DOI, NIST URL, article sha256, extraction selector, extract sha256); `bin/choupo-curate fetch` writes the points into constant/archiveExtracts/, which is gitignored.

STUDIES: every bubble-temperature block of the pair at or below 110 kPa with a declared uncertainty (3 block(s)). FITTED, ALL OF THEM (3): 10.1021/je100033s block 9 20 pts 101.3 kPa 352.92-363.98 K U 0.17-0.27 K 10.1016/j.jct.2013.08.011 block 9 19 pts 101.3 kPa 352.91-363.97 K U 0.1-0.53 K 10.1021/acs.jced.9b00015 block 2 19 pts 101.3 kPa 353.16-366.05 K U 0.38-0.58 K SET ASIDE BY THE EXTRACTION RULES, BY NAME: (none)

THE RULE (DEV.md 4c C283 amended, C284). Every laboratory above is fitted; none is held out (the C16 partition and its band are retired as the acceptance test). NRTL, a\_ij = a\_ji = 0, alpha = 0.30 fixed: two parameters (b\_ij, b\_ji); each point weighted by 1/U, U its own stated expanded uncertainty. Objective twoStage (least squares, then minimax; T\_bubble absolute, in K). OUTLIERS: after a least-squares stage over every point, a point with |T\_model - T| > 3 K is discarded BY NAME (`discard` on its evidence entry, listed below) and the fit repeated without it. ACCEPTED when max |T\_model - T| over the KEPT points is below 1 K. C283's 5 \% relative clause judges composition; this fit has no composition residual, so only the temperature clause applies. A pair that splits ONE liquid inside its validity span is not promoted (arm f). Two starts: (0, 0) and (800, -400).

OUTLIERS DISCARDED AFTER THE LEAST-SQUARES STAGE: @@OUTLIERS-BEGIN@@ none: every one of the 58 points lies within 3 K after the least-squares stage. @@OUTLIERS-END@@

@@RESULT-BEGIN@@ WHAT THE RUN FOUND (pinned in `expected`; measured, not argued). fit\_nrtl: b\_ij = 191.7 K, b\_ji = 249.3 K over 58 kept points of 3 laboratories (0 discarded): max |dT| 1.320 K, mean 0.283 K against the acceptance 1 K $\to$ notAccepted. fit\_nrtl\_start2 lands on the same pair (191.7, 249.3). NOT PROMOTABLE: max |dT| 1.320 K over the 58 kept points is not below 1 K @@RESULT-END@@

WHAT IS NOT CLAIMED. The pair is correct at these pressures and in this frame (ideal-gas vapour, the catalogue's own vapour-pressure records); the measurements were transcribed by the archive and have NOT been read back against the articles (the CHECKLIST that `fetch` writes is the aid).

\end{tutnarrative}
\begin{tutgolden}
\goldrow{diag}{fit\_nrtl}{bubble\_aad\_fit}{0.283328743772}
\goldrow{diag}{fit\_nrtl}{bubble\_max\_fit}{1.32039408006}
\goldrow{diag}{fit\_nrtl}{chi2}{223.424839131}
\goldrow{diag}{fit\_nrtl}{chi2\_in\_sample}{223.424839131}
\goldrow{diag}{fit\_nrtl}{chi2\_reduced}{3.9897292702}
\goldrow{diag}{fit\_nrtl}{cond\_JtJ}{65.2371995695}
\goldrow{diag}{fit\_nrtl}{converged}{1}
\goldrow{diag}{fit\_nrtl}{corr\_computed}{1}
\goldrow{diag}{fit\_nrtl}{dof}{56}
\goldrow{diag}{fit\_nrtl}{identifiable}{1}
\goldrow{diag}{fit\_nrtl}{invertible}{1}
\goldrow{diag}{fit\_nrtl}{iter}{5}
\goldrow{diag}{fit\_nrtl}{max\_abs\_corr}{0.989959746746}
\goldrow{diag}{fit\_nrtl}{max\_abs\_resid}{1.32039408006}
\goldmore{68}
\end{tutgolden}

\clearpage
\tutentry{curate112\_nrtl\_ethanol\_methylEthylKetone}{tutorials/props/curation/curate112\_nrtl\_ethanol\_methylEthylKetone}
\begin{tutscheme}
\begin{tikzpicture}[>={Stealth[length=2mm]}, every node/.style={font=\scriptsize}]
  \node[draw, rounded corners=2pt, align=center, minimum width=3.1cm, minimum height=0.8cm, fill=black!4] at (0.0,-0.0) {\textbf{fit\_nrtl}\\[-1pt]{\tiny fitParameters}};
  \node[draw, rounded corners=2pt, align=center, minimum width=3.1cm, minimum height=0.8cm, fill=black!4] at (3.6,-0.0) {\textbf{fit\_nrtl\_start2}\\[-1pt]{\tiny fitParameters}};
  \node[draw, rounded corners=2pt, align=center, minimum width=3.1cm, minimum height=0.8cm, fill=black!4] at (7.2,-0.0) {\textbf{txy\_ideal}\\[-1pt]{\tiny propertyScan1D}};
  \node[draw, rounded corners=2pt, align=center, minimum width=3.1cm, minimum height=0.8cm, fill=black!4] at (0.0,-1.3) {\textbf{txy\_unifac}\\[-1pt]{\tiny propertyScan1D}};
  \node[draw, rounded corners=2pt, align=center, minimum width=3.1cm, minimum height=0.8cm, fill=black!4] at (3.6,-1.3) {\textbf{txy\_nrtl}\\[-1pt]{\tiny propertyScan1D}};
\end{tikzpicture}
\end{tutscheme}
\tutwhat{NRTL ethanol-methylEthylKetone regressed by the curation pipeline to every laboratory's measured bubble temperatures (ThermoML locates, the articles source), accepted at max |dT| < 1 K}
\begin{tutnarrative}
AN NRTL PAIR FROM EVERY LABORATORY'S MEASUREMENTS.

WRITTEN BY bin/choupo-curate pair ethanol methylEthylKetone (DEV.md 4c C284, the bulk run; record docs/design/binary-pairs-from-open-measurements.md section 11).

THE DECLARATION -- WRITTEN AND COMMITTED BEFORE ANY FIT RAN (several declarations of the bulk run share one commit).

IDENTITY. ethanol (CAS 64-17-5, InChIKey LFQSCWFLJHTTHZ-UHFFFAOYSA-N) + methylEthylKetone (CAS 78-93-3, InChIKey ZWEHNKRNPOVVGH-UHFFFAOYSA-N), joined to the archive through identity.json -- never by name. CLASS (C284): both records are volatile molecular liquids (no role nonvolatile/solute, no dissociatesTo, no charge, no solid, a vapour-pressure model), and every measured temperature lies below both critical temperatures.

NO MEASURED SERIES IN THIS TREE. constant/archiveDatasets names each dataset (DOI, NIST URL, article sha256, extraction selector, extract sha256); `bin/choupo-curate fetch` writes the points into constant/archiveExtracts/, which is gitignored.

STUDIES: every bubble-temperature block of the pair at or below 110 kPa with a declared uncertainty (3 block(s)). FITTED, ALL OF THEM (3): 10.1016/j.fluid.2007.06.005 block 6 19 pts 101.3 kPa 347.48-351.69 K U 0.13-0.2 K 10.1016/j.fluid.2017.10.001 block 2 17 pts 101.3 kPa 347.50-351.89 K U 0.09-0.1 K (reports x of the other member: written as x[ethanol] = 1 - x by --x-of) 10.1021/acs.jced.7b00783 block 2 17 pts 101.3 kPa 347.50-351.89 K U 0.09-0.1 K (reports x of the other member: written as x[ethanol] = 1 - x by --x-of) SET ASIDE BY THE EXTRACTION RULES, BY NAME: (none)

THE RULE (DEV.md 4c C283 amended, C284). Every laboratory above is fitted; none is held out (the C16 partition and its band are retired as the acceptance test). NRTL, a\_ij = a\_ji = 0, alpha = 0.30 fixed: two parameters (b\_ij, b\_ji); each point weighted by 1/U, U its own stated expanded uncertainty. Objective twoStage (least squares, then minimax; T\_bubble absolute, in K). OUTLIERS: after a least-squares stage over every point, a point with |T\_model - T| > 3 K is discarded BY NAME (`discard` on its evidence entry, listed below) and the fit repeated without it. ACCEPTED when max |T\_model - T| over the KEPT points is below 1 K. C283's 5 \% relative clause judges composition; this fit has no composition residual, so only the temperature clause applies. A pair that splits ONE liquid inside its validity span is not promoted (arm f). Two starts: (0, 0) and (800, -400).

OUTLIERS DISCARDED AFTER THE LEAST-SQUARES STAGE: @@OUTLIERS-BEGIN@@ none: every one of the 53 points lies within 3 K after the least-squares stage. @@OUTLIERS-END@@

@@RESULT-BEGIN@@ WHAT THE RUN FOUND (pinned in `expected`; measured, not argued). fit\_nrtl: b\_ij = 159.6 K, b\_ji = 82.5 K over 53 kept points of 3 laboratories (0 discarded): max |dT| 0.155 K, mean 0.056 K against the acceptance 1 K $\to$ accepted. fit\_nrtl\_start2 lands on the same pair (159.6, 82.5). PROMOTED on 2026-11-10 into data/standards/parameters/NRTL/ethanol-methylEthylKetone.dat as reviewStatus autoPromoted, humanReviewed false, because: max |dT| 0.155 K < 1 K over the 53 kept points, one liquid across the validity span @@RESULT-END@@

WHAT IS NOT CLAIMED. The pair is correct at these pressures and in this frame (ideal-gas vapour, the catalogue's own vapour-pressure records); the measurements were transcribed by the archive and have NOT been read back against the articles (the CHECKLIST that `fetch` writes is the aid).

\end{tutnarrative}
\begin{tutgolden}
\goldrow{diag}{fit\_nrtl}{bubble\_aad\_fit}{0.0557687915762}
\goldrow{diag}{fit\_nrtl}{bubble\_max\_fit}{0.154628815367}
\goldrow{diag}{fit\_nrtl}{chi2}{20.6144530088}
\goldrow{diag}{fit\_nrtl}{chi2\_in\_sample}{20.6144530088}
\goldrow{diag}{fit\_nrtl}{chi2\_reduced}{0.404204960957}
\goldrow{diag}{fit\_nrtl}{cond\_JtJ}{143.316630099}
\goldrow{diag}{fit\_nrtl}{converged}{1}
\goldrow{diag}{fit\_nrtl}{corr\_computed}{1}
\goldrow{diag}{fit\_nrtl}{dof}{51}
\goldrow{diag}{fit\_nrtl}{identifiable}{1}
\goldrow{diag}{fit\_nrtl}{invertible}{1}
\goldrow{diag}{fit\_nrtl}{iter}{5}
\goldrow{diag}{fit\_nrtl}{max\_abs\_corr}{0.996529373194}
\goldrow{diag}{fit\_nrtl}{max\_abs\_resid}{0.154628815367}
\goldmore{80}
\end{tutgolden}

\clearpage
\tutentry{curate113\_nrtl\_isopropylAcetate\_isopropanol}{tutorials/props/curation/curate113\_nrtl\_isopropylAcetate\_isopropanol}
\begin{tutscheme}
\begin{tikzpicture}[>={Stealth[length=2mm]}, every node/.style={font=\scriptsize}]
  \node[draw, rounded corners=2pt, align=center, minimum width=3.1cm, minimum height=0.8cm, fill=black!4] at (0.0,-0.0) {\textbf{fit\_nrtl}\\[-1pt]{\tiny fitParameters}};
  \node[draw, rounded corners=2pt, align=center, minimum width=3.1cm, minimum height=0.8cm, fill=black!4] at (3.6,-0.0) {\textbf{fit\_nrtl\_start2}\\[-1pt]{\tiny fitParameters}};
  \node[draw, rounded corners=2pt, align=center, minimum width=3.1cm, minimum height=0.8cm, fill=black!4] at (7.2,-0.0) {\textbf{txy\_ideal}\\[-1pt]{\tiny propertyScan1D}};
  \node[draw, rounded corners=2pt, align=center, minimum width=3.1cm, minimum height=0.8cm, fill=black!4] at (0.0,-1.3) {\textbf{txy\_unifac}\\[-1pt]{\tiny propertyScan1D}};
  \node[draw, rounded corners=2pt, align=center, minimum width=3.1cm, minimum height=0.8cm, fill=black!4] at (3.6,-1.3) {\textbf{txy\_nrtl}\\[-1pt]{\tiny propertyScan1D}};
\end{tikzpicture}
\end{tutscheme}
\tutwhat{NRTL isopropylAcetate-isopropanol regressed by the curation pipeline to every laboratory's measured bubble temperatures (ThermoML locates, the articles source), accepted at max |dT| < 1 K}
\begin{tutnarrative}
AN NRTL PAIR FROM EVERY LABORATORY'S MEASUREMENTS.

WRITTEN BY bin/choupo-curate pair isopropylAcetate isopropanol (DEV.md 4c C284, the bulk run; record docs/design/binary-pairs-from-open-measurements.md section 11).

THE DECLARATION -- WRITTEN AND COMMITTED BEFORE ANY FIT RAN (several declarations of the bulk run share one commit).

IDENTITY. isopropylAcetate (CAS 108-21-4, InChIKey JMMWKPVZQRWMSS-UHFFFAOYSA-N) + isopropanol (CAS 67-63-0, InChIKey KFZMGEQAYNKOFK-UHFFFAOYSA-N), joined to the archive through identity.json -- never by name. CLASS (C284): both records are volatile molecular liquids (no role nonvolatile/solute, no dissociatesTo, no charge, no solid, a vapour-pressure model), and every measured temperature lies below both critical temperatures.

NO MEASURED SERIES IN THIS TREE. constant/archiveDatasets names each dataset (DOI, NIST URL, article sha256, extraction selector, extract sha256); `bin/choupo-curate fetch` writes the points into constant/archiveExtracts/, which is gitignored.

STUDIES: every bubble-temperature block of the pair at or below 110 kPa with a declared uncertainty (3 block(s)). FITTED, ALL OF THEM (3): 10.1021/acs.jced.5b00360 block 5 29 pts 101.32 kPa 353.51-361.31 K U 0.63-0.73 K 10.1016/j.fluid.2017.09.010 block 10 14 pts 101 kPa 353.60-359.70 K U 0.2-0.3 K (reports x of the other member: written as x[isopropylAcetate] = 1 - x by --x-of) 10.1016/j.fluid.2019.06.016 block 0 14 pts 101.3 kPa 353.67-360.65 K U 0.09-0.11 K SET ASIDE BY THE EXTRACTION RULES, BY NAME: (none)

THE RULE (DEV.md 4c C283 amended, C284). Every laboratory above is fitted; none is held out (the C16 partition and its band are retired as the acceptance test). NRTL, a\_ij = a\_ji = 0, alpha = 0.30 fixed: two parameters (b\_ij, b\_ji); each point weighted by 1/U, U its own stated expanded uncertainty. Objective twoStage (least squares, then minimax; T\_bubble absolute, in K). OUTLIERS: after a least-squares stage over every point, a point with |T\_model - T| > 3 K is discarded BY NAME (`discard` on its evidence entry, listed below) and the fit repeated without it. ACCEPTED when max |T\_model - T| over the KEPT points is below 1 K. C283's 5 \% relative clause judges composition; this fit has no composition residual, so only the temperature clause applies. A pair that splits ONE liquid inside its validity span is not promoted (arm f). Two starts: (0, 0) and (800, -400).

OUTLIERS DISCARDED AFTER THE LEAST-SQUARES STAGE: @@OUTLIERS-BEGIN@@ none: every one of the 57 points lies within 3 K after the least-squares stage. @@OUTLIERS-END@@

@@RESULT-BEGIN@@ WHAT THE RUN FOUND (pinned in `expected`; measured, not argued). fit\_nrtl: b\_ij = 177.7 K, b\_ji = 37.4 K over 57 kept points of 3 laboratories (0 discarded): max |dT| 0.455 K, mean 0.110 K against the acceptance 1 K $\to$ accepted. fit\_nrtl\_start2 lands on the same pair (177.7, 37.4). PROMOTED on 2026-11-10 into data/standards/parameters/NRTL/isopropanol-isopropylAcetate.dat as reviewStatus autoPromoted, humanReviewed false, because: max |dT| 0.455 K < 1 K over the 57 kept points, one liquid across the validity span @@RESULT-END@@

WHAT IS NOT CLAIMED. The pair is correct at these pressures and in this frame (ideal-gas vapour, the catalogue's own vapour-pressure records); the measurements were transcribed by the archive and have NOT been read back against the articles (the CHECKLIST that `fetch` writes is the aid).

\end{tutnarrative}
\begin{tutgolden}
\goldrow{diag}{fit\_nrtl}{bubble\_aad\_fit}{0.109832166085}
\goldrow{diag}{fit\_nrtl}{bubble\_max\_fit}{0.455162563814}
\goldrow{diag}{fit\_nrtl}{chi2}{22.4859939959}
\goldrow{diag}{fit\_nrtl}{chi2\_in\_sample}{22.4859939959}
\goldrow{diag}{fit\_nrtl}{chi2\_reduced}{0.408836254471}
\goldrow{diag}{fit\_nrtl}{cond\_JtJ}{167.943974306}
\goldrow{diag}{fit\_nrtl}{converged}{1}
\goldrow{diag}{fit\_nrtl}{corr\_computed}{1}
\goldrow{diag}{fit\_nrtl}{dof}{55}
\goldrow{diag}{fit\_nrtl}{identifiable}{1}
\goldrow{diag}{fit\_nrtl}{invertible}{1}
\goldrow{diag}{fit\_nrtl}{iter}{5}
\goldrow{diag}{fit\_nrtl}{max\_abs\_corr}{0.997036046475}
\goldrow{diag}{fit\_nrtl}{max\_abs\_resid}{0.455162563814}
\goldmore{80}
\end{tutgolden}

\clearpage
\tutentry{curate114\_nrtl\_pXylene\_nDecane}{tutorials/props/curation/curate114\_nrtl\_pXylene\_nDecane}
\begin{tutscheme}
\begin{tikzpicture}[>={Stealth[length=2mm]}, every node/.style={font=\scriptsize}]
  \node[draw, rounded corners=2pt, align=center, minimum width=3.1cm, minimum height=0.8cm, fill=black!4] at (0.0,-0.0) {\textbf{fit\_nrtl}\\[-1pt]{\tiny fitParameters}};
  \node[draw, rounded corners=2pt, align=center, minimum width=3.1cm, minimum height=0.8cm, fill=black!4] at (3.6,-0.0) {\textbf{fit\_nrtl\_start2}\\[-1pt]{\tiny fitParameters}};
  \node[draw, rounded corners=2pt, align=center, minimum width=3.1cm, minimum height=0.8cm, fill=black!4] at (7.2,-0.0) {\textbf{txy\_ideal}\\[-1pt]{\tiny propertyScan1D}};
  \node[draw, rounded corners=2pt, align=center, minimum width=3.1cm, minimum height=0.8cm, fill=black!4] at (0.0,-1.3) {\textbf{txy\_unifac}\\[-1pt]{\tiny propertyScan1D}};
  \node[draw, rounded corners=2pt, align=center, minimum width=3.1cm, minimum height=0.8cm, fill=black!4] at (3.6,-1.3) {\textbf{txy\_nrtl}\\[-1pt]{\tiny propertyScan1D}};
\end{tikzpicture}
\end{tutscheme}
\tutwhat{NRTL pXylene-nDecane regressed by the curation pipeline to every laboratory's measured bubble temperatures (ThermoML locates, the articles source), accepted at max |dT| < 1 K}
\begin{tutnarrative}
AN NRTL PAIR FROM EVERY LABORATORY'S MEASUREMENTS.

WRITTEN BY bin/choupo-curate pair pXylene nDecane (DEV.md 4c C284, the bulk run; record docs/design/binary-pairs-from-open-measurements.md section 11).

THE DECLARATION -- WRITTEN AND COMMITTED BEFORE ANY FIT RAN (several declarations of the bulk run share one commit).

IDENTITY. pXylene (CAS 106-42-3, InChIKey URLKBWYHVLBVBO-UHFFFAOYSA-N) + nDecane (CAS 124-18-5, InChIKey DIOQZVSQGTUSAI-UHFFFAOYSA-N), joined to the archive through identity.json -- never by name. CLASS (C284): both records are volatile molecular liquids (no role nonvolatile/solute, no dissociatesTo, no charge, no solid, a vapour-pressure model), and every measured temperature lies below both critical temperatures.

NO MEASURED SERIES IN THIS TREE. constant/archiveDatasets names each dataset (DOI, NIST URL, article sha256, extraction selector, extract sha256); `bin/choupo-curate fetch` writes the points into constant/archiveExtracts/, which is gitignored.

STUDIES: every bubble-temperature block of the pair at or below 110 kPa with a declared uncertainty (3 block(s)). FITTED, ALL OF THEM (3): 10.1021/je0501585 block 8 18 pts 97.3 kPa 410.10-439.60 K U 1-1 K 10.1016/j.fluid.2012.09.019 block 10 11 pts 20 kPa 360.62-389.31 K U 0.12-0.15 K (reports x of the other member: written as x[pXylene] = 1 - x by --x-of) 10.1016/j.fluid.2004.10.027 block 26 6 pts 95.8 kPa 413.05-436.35 K U 1.01-1.01 K SET ASIDE BY THE EXTRACTION RULES, BY NAME: (none)

THE RULE (DEV.md 4c C283 amended, C284). Every laboratory above is fitted; none is held out (the C16 partition and its band are retired as the acceptance test). NRTL, a\_ij = a\_ji = 0, alpha = 0.30 fixed: two parameters (b\_ij, b\_ji); each point weighted by 1/U, U its own stated expanded uncertainty. Objective twoStage (least squares, then minimax; T\_bubble absolute, in K). OUTLIERS: after a least-squares stage over every point, a point with |T\_model - T| > 3 K is discarded BY NAME (`discard` on its evidence entry, listed below) and the fit repeated without it. ACCEPTED when max |T\_model - T| over the KEPT points is below 1 K. C283's 5 \% relative clause judges composition; this fit has no composition residual, so only the temperature clause applies. A pair that splits ONE liquid inside its validity span is not promoted (arm f). Two starts: (0, 0) and (800, -400).

OUTLIERS DISCARDED AFTER THE LEAST-SQUARES STAGE: @@OUTLIERS-BEGIN@@ 7 point(s) of 35: doi 10.1021/je0501585 row 1: x 0.0338, P 97.300 kPa, T 439.60 K, dT +4.37 K; doi 10.1021/je0501585 row 2: x 0.0450, P 97.300 kPa, T 438.40 K, dT +4.97 K; doi 10.1021/je0501585 row 3: x 0.0850, P 97.300 kPa, T 435.80 K, dT +5.51 K; doi 10.1021/je0501585 row 4: x 0.1363, P 97.300 kPa, T 434.20 K, dT +4.59 K; doi 10.1021/je0501585 row 5: x 0.1500, P 97.300 kPa, T 433.70 K, dT +4.43 K; doi 10.1021/je0501585 row 6: x 0.2000, P 97.300 kPa, T 431.80 K, dT +4.04 K; doi 10.1021/je0501585 row 7: x 0.2500, P 97.300 kPa, T 430.60 K, dT +3.05 K @@OUTLIERS-END@@

@@RESULT-BEGIN@@ WHAT THE RUN FOUND (pinned in `expected`; measured, not argued). fit\_nrtl: b\_ij = 550.6 K, b\_ji = -378.9 K over 28 kept points of 3 laboratories (7 discarded): max |dT| 1.424 K, mean 0.334 K against the acceptance 1 K $\to$ notAccepted. fit\_nrtl\_start2 lands on the same pair (550.6, -378.9). NOT PROMOTABLE: max |dT| 1.424 K over the 28 kept points is not below 1 K @@RESULT-END@@

WHAT IS NOT CLAIMED. The pair is correct at these pressures and in this frame (ideal-gas vapour, the catalogue's own vapour-pressure records); the measurements were transcribed by the archive and have NOT been read back against the articles (the CHECKLIST that `fetch` writes is the aid).

\end{tutnarrative}
\begin{tutgolden}
\goldrow{diag}{fit\_nrtl}{bubble\_aad\_fit}{0.333550850372}
\goldrow{diag}{fit\_nrtl}{bubble\_max\_fit}{1.42373089764}
\goldrow{diag}{fit\_nrtl}{chi2}{17.6232918634}
\goldrow{diag}{fit\_nrtl}{chi2\_in\_sample}{17.6232918634}
\goldrow{diag}{fit\_nrtl}{chi2\_reduced}{0.677818917823}
\goldrow{diag}{fit\_nrtl}{cond\_JtJ}{115.617491724}
\goldrow{diag}{fit\_nrtl}{converged}{1}
\goldrow{diag}{fit\_nrtl}{corr\_computed}{1}
\goldrow{diag}{fit\_nrtl}{dof}{26}
\goldrow{diag}{fit\_nrtl}{identifiable}{1}
\goldrow{diag}{fit\_nrtl}{invertible}{1}
\goldrow{diag}{fit\_nrtl}{iter}{25}
\goldrow{diag}{fit\_nrtl}{max\_abs\_corr}{0.995703247485}
\goldrow{diag}{fit\_nrtl}{max\_abs\_resid}{1.42373089764}
\goldmore{62}
\end{tutgolden}

\clearpage
\tutentry{curate115\_nrtl\_ethylAcetate\_acetonitrile}{tutorials/props/curation/curate115\_nrtl\_ethylAcetate\_acetonitrile}
\begin{tutscheme}
\begin{tikzpicture}[>={Stealth[length=2mm]}, every node/.style={font=\scriptsize}]
  \node[draw, rounded corners=2pt, align=center, minimum width=3.1cm, minimum height=0.8cm, fill=black!4] at (0.0,-0.0) {\textbf{fit\_nrtl}\\[-1pt]{\tiny fitParameters}};
  \node[draw, rounded corners=2pt, align=center, minimum width=3.1cm, minimum height=0.8cm, fill=black!4] at (3.6,-0.0) {\textbf{fit\_nrtl\_start2}\\[-1pt]{\tiny fitParameters}};
  \node[draw, rounded corners=2pt, align=center, minimum width=3.1cm, minimum height=0.8cm, fill=black!4] at (7.2,-0.0) {\textbf{txy\_ideal}\\[-1pt]{\tiny propertyScan1D}};
  \node[draw, rounded corners=2pt, align=center, minimum width=3.1cm, minimum height=0.8cm, fill=black!4] at (0.0,-1.3) {\textbf{txy\_nrtl}\\[-1pt]{\tiny propertyScan1D}};
\end{tikzpicture}
\end{tutscheme}
\tutwhat{NRTL ethylAcetate-acetonitrile regressed by the curation pipeline to every laboratory's measured bubble temperatures (ThermoML locates, the articles source), accepted at max |dT| < 1 K}
\begin{tutnarrative}
AN NRTL PAIR FROM EVERY LABORATORY'S MEASUREMENTS.

WRITTEN BY bin/choupo-curate pair ethylAcetate acetonitrile (DEV.md 4c C284, the bulk run; record docs/design/binary-pairs-from-open-measurements.md section 11).

THE DECLARATION -- WRITTEN AND COMMITTED BEFORE ANY FIT RAN (several declarations of the bulk run share one commit).

IDENTITY. ethylAcetate (CAS 141-78-6, InChIKey XEKOWRVHYACXOJ-UHFFFAOYSA-N) + acetonitrile (CAS 75-05-8, InChIKey WEVYAHXRMPXWCK-UHFFFAOYSA-N), joined to the archive through identity.json -- never by name. CLASS (C284): both records are volatile molecular liquids (no role nonvolatile/solute, no dissociatesTo, no charge, no solid, a vapour-pressure model), and every measured temperature lies below both critical temperatures.

NO MEASURED SERIES IN THIS TREE. constant/archiveDatasets names each dataset (DOI, NIST URL, article sha256, extraction selector, extract sha256); `bin/choupo-curate fetch` writes the points into constant/archiveExtracts/, which is gitignored.

STUDIES: every bubble-temperature block of the pair at or below 110 kPa with a declared uncertainty (3 block(s)). FITTED, ALL OF THEM (3): 10.1016/j.fluid.2016.06.022 block 2 10 pts 101.3 kPa 348.45-352.45 K U 0.22-0.23 K 10.1021/acs.jced.5b00995 block 0 10 pts 101.3 kPa 348.45-352.45 K U 0.22-0.23 K 10.1021/je301204w block 0 10 pts 101.3 kPa 348.50-353.20 K U 0.4-0.4 K SET ASIDE BY THE EXTRACTION RULES, BY NAME: (none)

THE RULE (DEV.md 4c C283 amended, C284). Every laboratory above is fitted; none is held out (the C16 partition and its band are retired as the acceptance test). NRTL, a\_ij = a\_ji = 0, alpha = 0.30 fixed: two parameters (b\_ij, b\_ji); each point weighted by 1/U, U its own stated expanded uncertainty. Objective twoStage (least squares, then minimax; T\_bubble absolute, in K). OUTLIERS: after a least-squares stage over every point, a point with |T\_model - T| > 3 K is discarded BY NAME (`discard` on its evidence entry, listed below) and the fit repeated without it. ACCEPTED when max |T\_model - T| over the KEPT points is below 1 K. C283's 5 \% relative clause judges composition; this fit has no composition residual, so only the temperature clause applies. A pair that splits ONE liquid inside its validity span is not promoted (arm f). Two starts: (0, 0) and (800, -400).

OUTLIERS DISCARDED AFTER THE LEAST-SQUARES STAGE: @@OUTLIERS-BEGIN@@ none: every one of the 30 points lies within 3 K after the least-squares stage. @@OUTLIERS-END@@

@@RESULT-BEGIN@@ WHAT THE RUN FOUND (pinned in `expected`; measured, not argued). fit\_nrtl: b\_ij = -23.9 K, b\_ji = 215.2 K over 30 kept points of 3 laboratories (0 discarded): max |dT| 0.304 K, mean 0.152 K against the acceptance 1 K $\to$ accepted. fit\_nrtl\_start2 lands on the same pair (-23.9, 215.2). PROMOTED on 2026-11-10 into data/standards/parameters/NRTL/acetonitrile-ethylAcetate.dat as reviewStatus autoPromoted, humanReviewed false, because: max |dT| 0.304 K < 1 K over the 30 kept points, one liquid across the validity span @@RESULT-END@@

WHAT IS NOT CLAIMED. The pair is correct at these pressures and in this frame (ideal-gas vapour, the catalogue's own vapour-pressure records); the measurements were transcribed by the archive and have NOT been read back against the articles (the CHECKLIST that `fetch` writes is the aid).

\end{tutnarrative}
\begin{tutgolden}
\goldrow{diag}{fit\_nrtl}{bubble\_aad\_fit}{0.152441197848}
\goldrow{diag}{fit\_nrtl}{bubble\_max\_fit}{0.303992412903}
\goldrow{diag}{fit\_nrtl}{chi2}{13.2192251392}
\goldrow{diag}{fit\_nrtl}{chi2\_in\_sample}{13.2192251392}
\goldrow{diag}{fit\_nrtl}{chi2\_reduced}{0.472115183544}
\goldrow{diag}{fit\_nrtl}{cond\_JtJ}{431.539931707}
\goldrow{diag}{fit\_nrtl}{converged}{1}
\goldrow{diag}{fit\_nrtl}{corr\_computed}{1}
\goldrow{diag}{fit\_nrtl}{dof}{28}
\goldrow{diag}{fit\_nrtl}{identifiable}{1}
\goldrow{diag}{fit\_nrtl}{invertible}{1}
\goldrow{diag}{fit\_nrtl}{iter}{7}
\goldrow{diag}{fit\_nrtl}{max\_abs\_corr}{0.998199712083}
\goldrow{diag}{fit\_nrtl}{max\_abs\_resid}{0.303992412903}
\goldmore{66}
\end{tutgolden}

\clearpage
\tutentry{curate117\_nrtl\_nOctane\_ethanol}{tutorials/props/curation/curate117\_nrtl\_nOctane\_ethanol}
\begin{tutscheme}
\begin{tikzpicture}[>={Stealth[length=2mm]}, every node/.style={font=\scriptsize}]
  \node[draw, rounded corners=2pt, align=center, minimum width=3.1cm, minimum height=0.8cm, fill=black!4] at (0.0,-0.0) {\textbf{fit\_nrtl}\\[-1pt]{\tiny fitParameters}};
  \node[draw, rounded corners=2pt, align=center, minimum width=3.1cm, minimum height=0.8cm, fill=black!4] at (3.6,-0.0) {\textbf{fit\_nrtl\_start2}\\[-1pt]{\tiny fitParameters}};
  \node[draw, rounded corners=2pt, align=center, minimum width=3.1cm, minimum height=0.8cm, fill=black!4] at (7.2,-0.0) {\textbf{txy\_ideal}\\[-1pt]{\tiny propertyScan1D}};
  \node[draw, rounded corners=2pt, align=center, minimum width=3.1cm, minimum height=0.8cm, fill=black!4] at (0.0,-1.3) {\textbf{txy\_unifac}\\[-1pt]{\tiny propertyScan1D}};
  \node[draw, rounded corners=2pt, align=center, minimum width=3.1cm, minimum height=0.8cm, fill=black!4] at (3.6,-1.3) {\textbf{txy\_nrtl}\\[-1pt]{\tiny propertyScan1D}};
\end{tikzpicture}
\end{tutscheme}
\tutwhat{NRTL nOctane-ethanol regressed by the curation pipeline to every laboratory's measured bubble temperatures (ThermoML locates, the articles source), accepted at max |dT| < 1 K}
\begin{tutnarrative}
AN NRTL PAIR FROM EVERY LABORATORY'S MEASUREMENTS.

WRITTEN BY bin/choupo-curate pair nOctane ethanol (DEV.md 4c C284, the bulk run; record docs/design/binary-pairs-from-open-measurements.md section 11).

THE DECLARATION -- WRITTEN AND COMMITTED BEFORE ANY FIT RAN (several declarations of the bulk run share one commit).

IDENTITY. nOctane (CAS 111-65-9, InChIKey TVMXDCGIABBOFY-UHFFFAOYSA-N) + ethanol (CAS 64-17-5, InChIKey LFQSCWFLJHTTHZ-UHFFFAOYSA-N), joined to the archive through identity.json -- never by name. CLASS (C284): both records are volatile molecular liquids (no role nonvolatile/solute, no dissociatesTo, no charge, no solid, a vapour-pressure model), and every measured temperature lies below both critical temperatures.

NO MEASURED SERIES IN THIS TREE. constant/archiveDatasets names each dataset (DOI, NIST URL, article sha256, extraction selector, extract sha256); `bin/choupo-curate fetch` writes the points into constant/archiveExtracts/, which is gitignored.

STUDIES: every bubble-temperature block of the pair at or below 110 kPa with a declared uncertainty (2 block(s)). FITTED, ALL OF THEM (2): 10.1021/je400238j block 2 84 pts 4 isobars 30-101.3 kPa 321.60-366.08 K U 0.39-0.86 K 10.1016/j.fluid.2018.07.013 block 4 21 pts 101.3 kPa 348.65-390.55 K U 0.22-2.26 K (reports x of the other member: written as x[nOctane] = 1 - x by --x-of) SET ASIDE BY THE EXTRACTION RULES, BY NAME: (none)

THE RULE (DEV.md 4c C283 amended, C284). Every laboratory above is fitted; none is held out (the C16 partition and its band are retired as the acceptance test). NRTL, a\_ij = a\_ji = 0, alpha = 0.30 fixed: two parameters (b\_ij, b\_ji); each point weighted by 1/U, U its own stated expanded uncertainty. Objective twoStage (least squares, then minimax; T\_bubble absolute, in K). OUTLIERS: after a least-squares stage over every point, a point with |T\_model - T| > 3 K is discarded BY NAME (`discard` on its evidence entry, listed below) and the fit repeated without it. ACCEPTED when max |T\_model - T| over the KEPT points is below 1 K. C283's 5 \% relative clause judges composition; this fit has no composition residual, so only the temperature clause applies. A pair that splits ONE liquid inside its validity span is not promoted (arm f). Two starts: (0, 0) and (800, -400).

OUTLIERS DISCARDED AFTER THE LEAST-SQUARES STAGE: @@OUTLIERS-BEGIN@@ 6 point(s) of 105: doi 10.1021/je400238j row 82: x 0.9542, P 60.000 kPa, T 348.36 K, dT +4.14 K; doi 10.1016/j.fluid.2018.07.013 row 1: x 0.9884, P 101.300 kPa, T 390.55 K, dT -4.82 K; doi 10.1016/j.fluid.2018.07.013 row 2: x 0.9762, P 101.300 kPa, T 385.55 K, dT -8.27 K; doi 10.1016/j.fluid.2018.07.013 row 3: x 0.9692, P 101.300 kPa, T 382.15 K, dT -8.44 K; doi 10.1016/j.fluid.2018.07.013 row 4: x 0.9498, P 101.300 kPa, T 371.55 K, dT -4.90 K; doi 10.1016/j.fluid.2018.07.013 row 5: x 0.9413, P 101.300 kPa, T 368.05 K, dT -3.63 K @@OUTLIERS-END@@

@@RESULT-BEGIN@@ WHAT THE RUN FOUND (pinned in `expected`; measured, not argued). fit\_nrtl: b\_ij = 506.6 K, b\_ji = 442.1 K over 99 kept points of 2 laboratories (6 discarded): max |dT| 3.047 K, mean 0.745 K against the acceptance 1 K $\to$ notAccepted. fit\_nrtl\_start2 lands on the same pair (506.6, 442.1). NOT PROMOTABLE: max |dT| 3.047 K over the 99 kept points is not below 1 K @@RESULT-END@@

WHAT IS NOT CLAIMED. The pair is correct at these pressures and in this frame (ideal-gas vapour, the catalogue's own vapour-pressure records); the measurements were transcribed by the archive and have NOT been read back against the articles (the CHECKLIST that `fetch` writes is the aid).

\end{tutnarrative}
\begin{tutgolden}
\goldrow{diag}{fit\_nrtl}{bubble\_aad\_fit}{0.745204678956}
\goldrow{diag}{fit\_nrtl}{bubble\_max\_fit}{3.04684066943}
\goldrow{diag}{fit\_nrtl}{chi2}{647.898223366}
\goldrow{diag}{fit\_nrtl}{chi2\_in\_sample}{647.898223366}
\goldrow{diag}{fit\_nrtl}{chi2\_reduced}{6.67936312748}
\goldrow{diag}{fit\_nrtl}{cond\_JtJ}{5.62821530265}
\goldrow{diag}{fit\_nrtl}{converged}{1}
\goldrow{diag}{fit\_nrtl}{corr\_computed}{1}
\goldrow{diag}{fit\_nrtl}{dof}{97}
\goldrow{diag}{fit\_nrtl}{identifiable}{1}
\goldrow{diag}{fit\_nrtl}{invertible}{1}
\goldrow{diag}{fit\_nrtl}{iter}{5}
\goldrow{diag}{fit\_nrtl}{max\_abs\_corr}{0.899290914635}
\goldrow{diag}{fit\_nrtl}{max\_abs\_resid}{3.04684066943}
\goldmore{62}
\end{tutgolden}

\clearpage
\tutentry{curate118\_nrtl\_2Methyl2Propanol\_14Dioxane}{tutorials/props/curation/curate118\_nrtl\_2Methyl2Propanol\_14Dioxane}
\begin{tutscheme}
\begin{tikzpicture}[>={Stealth[length=2mm]}, every node/.style={font=\scriptsize}]
  \node[draw, rounded corners=2pt, align=center, minimum width=3.1cm, minimum height=0.8cm, fill=black!4] at (0.0,-0.0) {\textbf{fit\_nrtl}\\[-1pt]{\tiny fitParameters}};
  \node[draw, rounded corners=2pt, align=center, minimum width=3.1cm, minimum height=0.8cm, fill=black!4] at (3.6,-0.0) {\textbf{fit\_nrtl\_start2}\\[-1pt]{\tiny fitParameters}};
  \node[draw, rounded corners=2pt, align=center, minimum width=3.1cm, minimum height=0.8cm, fill=black!4] at (7.2,-0.0) {\textbf{txy\_ideal}\\[-1pt]{\tiny propertyScan1D}};
  \node[draw, rounded corners=2pt, align=center, minimum width=3.1cm, minimum height=0.8cm, fill=black!4] at (0.0,-1.3) {\textbf{txy\_unifac}\\[-1pt]{\tiny propertyScan1D}};
  \node[draw, rounded corners=2pt, align=center, minimum width=3.1cm, minimum height=0.8cm, fill=black!4] at (3.6,-1.3) {\textbf{txy\_nrtl}\\[-1pt]{\tiny propertyScan1D}};
\end{tikzpicture}
\end{tutscheme}
\tutwhat{NRTL 2Methyl2Propanol-14Dioxane regressed by the curation pipeline to every laboratory's measured bubble temperatures (ThermoML locates, the articles source), accepted at max |dT| < 1 K}
\begin{tutnarrative}
AN NRTL PAIR FROM EVERY LABORATORY'S MEASUREMENTS.

WRITTEN BY bin/choupo-curate pair 2Methyl2Propanol 14Dioxane (DEV.md 4c C284, the bulk run; record docs/design/binary-pairs-from-open-measurements.md section 11).

THE DECLARATION -- WRITTEN AND COMMITTED BEFORE ANY FIT RAN (several declarations of the bulk run share one commit).

IDENTITY. 2Methyl2Propanol (CAS 75-65-0, InChIKey DKGAVHZHDRPRBM-UHFFFAOYSA-N) + 14Dioxane (CAS 123-91-1, InChIKey RYHBNJHYFVUHQT-UHFFFAOYSA-N), joined to the archive through identity.json -- never by name. CLASS (C284): both records are volatile molecular liquids (no role nonvolatile/solute, no dissociatesTo, no charge, no solid, a vapour-pressure model), and every measured temperature lies below both critical temperatures.

NO MEASURED SERIES IN THIS TREE. constant/archiveDatasets names each dataset (DOI, NIST URL, article sha256, extraction selector, extract sha256); `bin/choupo-curate fetch` writes the points into constant/archiveExtracts/, which is gitignored.

STUDIES: every bubble-temperature block of the pair at or below 110 kPa with a declared uncertainty (2 block(s)). FITTED, ALL OF THEM (2): 10.1016/j.fluid.2007.11.002 block 21 76 pts 4 isobars 26.66-101.325 kPa 324.45-372.55 K U 0.34-0.55 K 10.1016/j.jct.2003.09.001 block 18 24 pts 2 isobars 40-101.3 kPa 334.21-370.84 K U 0.09-0.98 K (reports x of the other member: written as x[2Methyl2Propanol] = 1 - x by --x-of) SET ASIDE BY THE EXTRACTION RULES, BY NAME: (none)

THE RULE (DEV.md 4c C283 amended, C284). Every laboratory above is fitted; none is held out (the C16 partition and its band are retired as the acceptance test). NRTL, a\_ij = a\_ji = 0, alpha = 0.30 fixed: two parameters (b\_ij, b\_ji); each point weighted by 1/U, U its own stated expanded uncertainty. Objective twoStage (least squares, then minimax; T\_bubble absolute, in K). OUTLIERS: after a least-squares stage over every point, a point with |T\_model - T| > 3 K is discarded BY NAME (`discard` on its evidence entry, listed below) and the fit repeated without it. ACCEPTED when max |T\_model - T| over the KEPT points is below 1 K. C283's 5 \% relative clause judges composition; this fit has no composition residual, so only the temperature clause applies. A pair that splits ONE liquid inside its validity span is not promoted (arm f). Two starts: (0, 0) and (800, -400).

OUTLIERS DISCARDED AFTER THE LEAST-SQUARES STAGE: @@OUTLIERS-BEGIN@@ none: every one of the 100 points lies within 3 K after the least-squares stage. @@OUTLIERS-END@@

@@RESULT-BEGIN@@ WHAT THE RUN FOUND (pinned in `expected`; measured, not argued). fit\_nrtl: b\_ij = 299.4 K, b\_ji = -98.2 K over 100 kept points of 2 laboratories (0 discarded): max |dT| 1.987 K, mean 0.699 K against the acceptance 1 K $\to$ notAccepted. fit\_nrtl\_start2 lands on the same pair (299.4, -98.2). NOT PROMOTABLE: max |dT| 1.987 K over the 100 kept points is not below 1 K @@RESULT-END@@

WHAT IS NOT CLAIMED. The pair is correct at these pressures and in this frame (ideal-gas vapour, the catalogue's own vapour-pressure records); the measurements were transcribed by the archive and have NOT been read back against the articles (the CHECKLIST that `fetch` writes is the aid).

\end{tutnarrative}
\begin{tutgolden}
\goldrow{diag}{fit\_nrtl}{bubble\_aad\_fit}{0.698854004157}
\goldrow{diag}{fit\_nrtl}{bubble\_max\_fit}{1.98691251367}
\goldrow{diag}{fit\_nrtl}{chi2}{375.892218354}
\goldrow{diag}{fit\_nrtl}{chi2\_in\_sample}{375.892218354}
\goldrow{diag}{fit\_nrtl}{chi2\_reduced}{3.83563488116}
\goldrow{diag}{fit\_nrtl}{cond\_JtJ}{151.936081511}
\goldrow{diag}{fit\_nrtl}{converged}{1}
\goldrow{diag}{fit\_nrtl}{corr\_computed}{1}
\goldrow{diag}{fit\_nrtl}{dof}{98}
\goldrow{diag}{fit\_nrtl}{identifiable}{1}
\goldrow{diag}{fit\_nrtl}{invertible}{1}
\goldrow{diag}{fit\_nrtl}{iter}{10}
\goldrow{diag}{fit\_nrtl}{max\_abs\_corr}{0.996725298484}
\goldrow{diag}{fit\_nrtl}{max\_abs\_resid}{1.98691251367}
\goldmore{53}
\end{tutgolden}

\clearpage
\tutentry{curate119\_nrtl\_1Propanol\_dimethylCarbonate}{tutorials/props/curation/curate119\_nrtl\_1Propanol\_dimethylCarbonate}
\begin{tutscheme}
\begin{tikzpicture}[>={Stealth[length=2mm]}, every node/.style={font=\scriptsize}]
  \node[draw, rounded corners=2pt, align=center, minimum width=3.1cm, minimum height=0.8cm, fill=black!4] at (0.0,-0.0) {\textbf{fit\_nrtl}\\[-1pt]{\tiny fitParameters}};
  \node[draw, rounded corners=2pt, align=center, minimum width=3.1cm, minimum height=0.8cm, fill=black!4] at (3.6,-0.0) {\textbf{fit\_nrtl\_start2}\\[-1pt]{\tiny fitParameters}};
  \node[draw, rounded corners=2pt, align=center, minimum width=3.1cm, minimum height=0.8cm, fill=black!4] at (7.2,-0.0) {\textbf{txy\_ideal}\\[-1pt]{\tiny propertyScan1D}};
  \node[draw, rounded corners=2pt, align=center, minimum width=3.1cm, minimum height=0.8cm, fill=black!4] at (0.0,-1.3) {\textbf{txy\_nrtl}\\[-1pt]{\tiny propertyScan1D}};
\end{tikzpicture}
\end{tutscheme}
\tutwhat{NRTL 1Propanol-dimethylCarbonate regressed by the curation pipeline to every laboratory's measured bubble temperatures (ThermoML locates, the articles source), accepted at max |dT| < 1 K}
\begin{tutnarrative}
AN NRTL PAIR FROM EVERY LABORATORY'S MEASUREMENTS.

WRITTEN BY bin/choupo-curate pair 1Propanol dimethylCarbonate (DEV.md 4c C284, the bulk run; record docs/design/binary-pairs-from-open-measurements.md section 11).

THE DECLARATION -- WRITTEN AND COMMITTED BEFORE ANY FIT RAN (several declarations of the bulk run share one commit).

IDENTITY. 1Propanol (CAS 71-23-8, InChIKey BDERNNFJNOPAEC-UHFFFAOYSA-N) + dimethylCarbonate (CAS 616-38-6, InChIKey IEJIGPNLZYLLBP-UHFFFAOYSA-N), joined to the archive through identity.json -- never by name. CLASS (C284): both records are volatile molecular liquids (no role nonvolatile/solute, no dissociatesTo, no charge, no solid, a vapour-pressure model), and every measured temperature lies below both critical temperatures.

NO MEASURED SERIES IN THIS TREE. constant/archiveDatasets names each dataset (DOI, NIST URL, article sha256, extraction selector, extract sha256); `bin/choupo-curate fetch` writes the points into constant/archiveExtracts/, which is gitignored.

STUDIES: every bubble-temperature block of the pair at or below 110 kPa with a declared uncertainty (2 block(s)). FITTED, ALL OF THEM (2): 10.1016/j.jct.2011.08.006 block 3 75 pts 5 isobars 40-93.32 kPa 334.45-363.46 K U 0.11-0.16 K 10.1021/je030131q block 24 9 pts 95.8 kPa 359.65-364.05 K U 0.32-0.32 K (reports x of the other member: written as x[1Propanol] = 1 - x by --x-of) SET ASIDE BY THE EXTRACTION RULES, BY NAME: (none)

THE RULE (DEV.md 4c C283 amended, C284). Every laboratory above is fitted; none is held out (the C16 partition and its band are retired as the acceptance test). NRTL, a\_ij = a\_ji = 0, alpha = 0.30 fixed: two parameters (b\_ij, b\_ji); each point weighted by 1/U, U its own stated expanded uncertainty. Objective twoStage (least squares, then minimax; T\_bubble absolute, in K). OUTLIERS: after a least-squares stage over every point, a point with |T\_model - T| > 3 K is discarded BY NAME (`discard` on its evidence entry, listed below) and the fit repeated without it. ACCEPTED when max |T\_model - T| over the KEPT points is below 1 K. C283's 5 \% relative clause judges composition; this fit has no composition residual, so only the temperature clause applies. A pair that splits ONE liquid inside its validity span is not promoted (arm f). Two starts: (0, 0) and (800, -400).

OUTLIERS DISCARDED AFTER THE LEAST-SQUARES STAGE: @@OUTLIERS-BEGIN@@ none: every one of the 84 points lies within 3 K after the least-squares stage. @@OUTLIERS-END@@

@@RESULT-BEGIN@@ WHAT THE RUN FOUND (pinned in `expected`; measured, not argued). fit\_nrtl: b\_ij = 240.6 K, b\_ji = 163.9 K over 84 kept points of 2 laboratories (0 discarded): max |dT| 1.955 K, mean 0.377 K against the acceptance 1 K $\to$ notAccepted. fit\_nrtl\_start2 lands on the same pair (240.6, 163.9). NOT PROMOTABLE: max |dT| 1.955 K over the 84 kept points is not below 1 K @@RESULT-END@@

WHAT IS NOT CLAIMED. The pair is correct at these pressures and in this frame (ideal-gas vapour, the catalogue's own vapour-pressure records); the measurements were transcribed by the archive and have NOT been read back against the articles (the CHECKLIST that `fetch` writes is the aid).

\end{tutnarrative}
\begin{tutgolden}
\goldrow{diag}{fit\_nrtl}{bubble\_aad\_fit}{0.376865938318}
\goldrow{diag}{fit\_nrtl}{bubble\_max\_fit}{1.95457878377}
\goldrow{diag}{fit\_nrtl}{chi2}{740.741568692}
\goldrow{diag}{fit\_nrtl}{chi2\_in\_sample}{740.741568692}
\goldrow{diag}{fit\_nrtl}{chi2\_reduced}{9.03343376454}
\goldrow{diag}{fit\_nrtl}{cond\_JtJ}{61.1552804338}
\goldrow{diag}{fit\_nrtl}{converged}{1}
\goldrow{diag}{fit\_nrtl}{corr\_computed}{1}
\goldrow{diag}{fit\_nrtl}{dof}{82}
\goldrow{diag}{fit\_nrtl}{identifiable}{1}
\goldrow{diag}{fit\_nrtl}{invertible}{1}
\goldrow{diag}{fit\_nrtl}{iter}{4}
\goldrow{diag}{fit\_nrtl}{max\_abs\_corr}{0.991923012314}
\goldrow{diag}{fit\_nrtl}{max\_abs\_resid}{1.95457878377}
\goldmore{50}
\end{tutgolden}

\clearpage
\tutentry{curate11\_nrtl\_acetone\_ethanol}{tutorials/props/curation/curate11\_nrtl\_acetone\_ethanol}
\begin{tutscheme}
\begin{tikzpicture}[>={Stealth[length=2mm]}, every node/.style={font=\scriptsize}]
  \node[draw, rounded corners=2pt, align=center, minimum width=3.1cm, minimum height=0.8cm, fill=black!4] at (0.0,-0.0) {\textbf{fit\_nrtl}\\[-1pt]{\tiny fitParameters}};
  \node[draw, rounded corners=2pt, align=center, minimum width=3.1cm, minimum height=0.8cm, fill=black!4] at (3.6,-0.0) {\textbf{fit\_nrtl\_start2}\\[-1pt]{\tiny fitParameters}};
  \node[draw, rounded corners=2pt, align=center, minimum width=3.1cm, minimum height=0.8cm, fill=black!4] at (7.2,-0.0) {\textbf{txy\_ideal}\\[-1pt]{\tiny propertyScan1D}};
  \node[draw, rounded corners=2pt, align=center, minimum width=3.1cm, minimum height=0.8cm, fill=black!4] at (0.0,-1.3) {\textbf{txy\_unifac}\\[-1pt]{\tiny propertyScan1D}};
  \node[draw, rounded corners=2pt, align=center, minimum width=3.1cm, minimum height=0.8cm, fill=black!4] at (3.6,-1.3) {\textbf{txy\_nrtl}\\[-1pt]{\tiny propertyScan1D}};
  \node[draw, rounded corners=2pt, align=center, minimum width=3.1cm, minimum height=0.8cm, fill=black!4] at (7.2,-1.3) {\textbf{txy\_ideal\_fitP}\\[-1pt]{\tiny propertyScan1D}};
  \node[draw, rounded corners=2pt, align=center, minimum width=3.1cm, minimum height=0.8cm, fill=black!4] at (0.0,-2.6) {\textbf{txy\_unifac\_fitP}\\[-1pt]{\tiny propertyScan1D}};
  \node[draw, rounded corners=2pt, align=center, minimum width=3.1cm, minimum height=0.8cm, fill=black!4] at (3.6,-2.6) {\textbf{txy\_nrtl\_fitP}\\[-1pt]{\tiny propertyScan1D}};
\end{tikzpicture}
\end{tutscheme}
\tutwhat{NRTL acetone-ethanol regressed to measured bubble temperatures (ThermoML locates, the article sources), held out against an independent study, drawn against both beside ideal and UNIFAC}
\begin{tutnarrative}
AN NRTL PAIR FROM PUBLISHED MEASUREMENTS, END TO END.

WHY THIS SYSTEM. two of the commonest laboratory solvents, and a positive deviation from Raoult's law WITHOUT an azeotrope -- the contrast to every other pair of this family, showing that a non-ideal pair need not pinch a column.

THE CHAIN, EACH STEP A COMMAND (DEV.md C16; record docs/design/binary-pairs-from-open-measurements.md): 1. bin/choupo-thermoml sync the NIST/TRC ThermoML Archive (ThermoML.v2020-09-30.tgz, doi 10.18434/mds2-2422, sha256-verified) 2. bin/choupo-thermoml identity CAS $\to$ InChIKey for the catalogue, two resolvers that must agree 3. bin/choupo-thermoml vle-index every binary VLE block, searchable 4. bin/choupo-thermoml vle-search acetone ethanol --geometry bubbleT 5. bin/choupo-thermoml extract-vle <file> --block N --interior-only $\to$ the two files under constant/experiments/ 6. choupoProps (this case) fitParameters $\to$ a proposal record 7. promotion (a human act) $\to$ data/standards/parameters/NRTL/acetone-ethanol.dat ONLY if the held-out verdict is `validated` 8. bin/choupo-import the seal mirrors that record here

THE EVIDENCE. FIT: doi 10.1016/j.fluid.2005.01.007 (19 interior points, 101.32 kPa). HELD OUT: doi 10.1021/je400544h (11 interior points, 101.3 kPa) -- an article whose title names another system; the acetone + ethanol block is the one the archive files under it, and reading it back is part of the curator's check. Chosen BEFORE any fit by the slice-2 partition rule: the FIT set is the low-pressure bubble-temperature study with the most interior points; the HELD-OUT set is the next largest that shares NO author with it (the archive carries no affiliations, so a shared author is the dependence the rule can check). Each file names its article, its archive file and block, the sha256 of the archive file, the two compounds' CAS and InChIKey, and the per-point expanded uncertainty; each says `reviewStatus transcribedNotCheckedAgainstArticle` because nobody has read the numbers back against the paper yet, and the run announces it. Pure- component endpoints are set aside (they measure the vapour-pressure record, not the pair) and kept in each file under `selection`.

THE FIT. NRTL with a\_ij = a\_ji = 0 and alpha = 0.30 fixed: TWO parameters (b\_ij, b\_ji), because bubble temperatures over the ranges below do not identify more (fit\_nrtl\_ab in the methanol + water case asks that question directly). Each point is weighted by 1/U, U its OWN stated expanded uncertainty (`weighting uncertainty;`), so a point measured to 0.1 K is not outvoted by one measured to 1 K. The held-out AAD is plain kelvin, unweighted.

WHAT THE RUN FOUND (pinned in `expected`; measured, not argued). fit\_nrtl: b\_ij = 149.8 K, b\_ji = 68.1 K, in-sample rms 0.072 K; fit\_nrtl\_start2 lands on the same pair. max|corr| 0.997: the two b's are nearly one direction in the data, as for every near-symmetric pair here. HELD OUT (101.3 kPa, no author in common): AAD 0.093 K (0.028 \%) against a band of 0.14 \% $\to$ `validated`; the pair is the catalogue record NRTL/acetone-ethanol.dat. Overlay on the held-out set: NRTL 0.093 K, UNIFAC 0.79 K, ideal 3.6 K. PHASE STABILITY: one liquid at every temperature down to 200 K.

WHAT IS NOT CLAIMED. The pair is correct at these pressures and in this frame (ideal-gas vapour, the catalogue's Antoine records); nothing here tests it at high pressure, against LLE, or with another vapour model.

\end{tutnarrative}
\begin{tutgolden}
\goldrow{diag}{fit\_nrtl}{aad\_heldout\_K}{0.0928861936011}
\goldrow{diag}{fit\_nrtl}{aad\_heldout\_pct}{0.0274764440669}
\goldrow{diag}{fit\_nrtl}{chi2}{1.96410744483}
\goldrow{diag}{fit\_nrtl}{chi2\_in\_sample}{1.96410744483}
\goldrow{diag}{fit\_nrtl}{chi2\_reduced}{0.115535732049}
\goldrow{diag}{fit\_nrtl}{cond\_JtJ}{162.970032107}
\goldrow{diag}{fit\_nrtl}{converged}{1}
\goldrow{diag}{fit\_nrtl}{corr\_computed}{1}
\goldrow{diag}{fit\_nrtl}{dof}{17}
\goldrow{diag}{fit\_nrtl}{identifiable}{1}
\goldrow{diag}{fit\_nrtl}{invertible}{1}
\goldrow{diag}{fit\_nrtl}{iter}{9}
\goldrow{diag}{fit\_nrtl}{max\_abs\_corr}{0.996945998842}
\goldrow{diag}{fit\_nrtl}{max\_abs\_resid}{0.212432153701}
\goldmore{70}
\end{tutgolden}

\clearpage
\tutentry{curate120\_nrtl\_ethanol\_pXylene}{tutorials/props/curation/curate120\_nrtl\_ethanol\_pXylene}
\begin{tutscheme}
\begin{tikzpicture}[>={Stealth[length=2mm]}, every node/.style={font=\scriptsize}]
  \node[draw, rounded corners=2pt, align=center, minimum width=3.1cm, minimum height=0.8cm, fill=black!4] at (0.0,-0.0) {\textbf{fit\_nrtl}\\[-1pt]{\tiny fitParameters}};
  \node[draw, rounded corners=2pt, align=center, minimum width=3.1cm, minimum height=0.8cm, fill=black!4] at (3.6,-0.0) {\textbf{fit\_nrtl\_start2}\\[-1pt]{\tiny fitParameters}};
  \node[draw, rounded corners=2pt, align=center, minimum width=3.1cm, minimum height=0.8cm, fill=black!4] at (7.2,-0.0) {\textbf{txy\_ideal}\\[-1pt]{\tiny propertyScan1D}};
  \node[draw, rounded corners=2pt, align=center, minimum width=3.1cm, minimum height=0.8cm, fill=black!4] at (0.0,-1.3) {\textbf{txy\_unifac}\\[-1pt]{\tiny propertyScan1D}};
  \node[draw, rounded corners=2pt, align=center, minimum width=3.1cm, minimum height=0.8cm, fill=black!4] at (3.6,-1.3) {\textbf{txy\_nrtl}\\[-1pt]{\tiny propertyScan1D}};
\end{tikzpicture}
\end{tutscheme}
\tutwhat{NRTL ethanol-pXylene regressed by the curation pipeline to every laboratory's measured bubble temperatures (ThermoML locates, the articles source), accepted at max |dT| < 1 K}
\begin{tutnarrative}
AN NRTL PAIR FROM EVERY LABORATORY'S MEASUREMENTS.

WRITTEN BY bin/choupo-curate pair ethanol pXylene (DEV.md 4c C284, the bulk run; record docs/design/binary-pairs-from-open-measurements.md section 11).

THE DECLARATION -- WRITTEN AND COMMITTED BEFORE ANY FIT RAN (several declarations of the bulk run share one commit).

IDENTITY. ethanol (CAS 64-17-5, InChIKey LFQSCWFLJHTTHZ-UHFFFAOYSA-N) + pXylene (CAS 106-42-3, InChIKey URLKBWYHVLBVBO-UHFFFAOYSA-N), joined to the archive through identity.json -- never by name. CLASS (C284): both records are volatile molecular liquids (no role nonvolatile/solute, no dissociatesTo, no charge, no solid, a vapour-pressure model), and every measured temperature lies below both critical temperatures.

NO MEASURED SERIES IN THIS TREE. constant/archiveDatasets names each dataset (DOI, NIST URL, article sha256, extraction selector, extract sha256); `bin/choupo-curate fetch` writes the points into constant/archiveExtracts/, which is gitignored.

STUDIES: every bubble-temperature block of the pair at or below 110 kPa with a declared uncertainty (2 block(s)). FITTED, ALL OF THEM (2): 10.1021/je200868y block 6 61 pts 6 isobars 40-101.3 kPa 329.95-369.38 K U 0.68-0.77 K 10.1021/je0342318 block 13 21 pts 101.3 kPa 351.48-404.93 K U 0.68-1.36 K SET ASIDE BY THE EXTRACTION RULES, BY NAME: (none)

THE RULE (DEV.md 4c C283 amended, C284). Every laboratory above is fitted; none is held out (the C16 partition and its band are retired as the acceptance test). NRTL, a\_ij = a\_ji = 0, alpha = 0.30 fixed: two parameters (b\_ij, b\_ji); each point weighted by 1/U, U its own stated expanded uncertainty. Objective twoStage (least squares, then minimax; T\_bubble absolute, in K). OUTLIERS: after a least-squares stage over every point, a point with |T\_model - T| > 3 K is discarded BY NAME (`discard` on its evidence entry, listed below) and the fit repeated without it. ACCEPTED when max |T\_model - T| over the KEPT points is below 1 K. C283's 5 \% relative clause judges composition; this fit has no composition residual, so only the temperature clause applies. A pair that splits ONE liquid inside its validity span is not promoted (arm f). Two starts: (0, 0) and (800, -400).

OUTLIERS DISCARDED AFTER THE LEAST-SQUARES STAGE: @@OUTLIERS-BEGIN@@ 1 point(s) of 82: doi 10.1021/je0342318 row 2: x 0.0310, P 101.300 kPa, T 391.42 K, dT -5.11 K @@OUTLIERS-END@@

@@RESULT-BEGIN@@ WHAT THE RUN FOUND (pinned in `expected`; measured, not argued). fit\_nrtl: b\_ij = 209.5 K, b\_ji = 514.2 K over 81 kept points of 2 laboratories (1 discarded): max |dT| 3.080 K, mean 0.772 K against the acceptance 1 K $\to$ notAccepted. fit\_nrtl\_start2 lands on the same pair (209.5, 514.2). NOT PROMOTABLE: max |dT| 3.080 K over the 81 kept points is not below 1 K @@RESULT-END@@

WHAT IS NOT CLAIMED. The pair is correct at these pressures and in this frame (ideal-gas vapour, the catalogue's own vapour-pressure records); the measurements were transcribed by the archive and have NOT been read back against the articles (the CHECKLIST that `fetch` writes is the aid).

\end{tutnarrative}
\begin{tutgolden}
\goldrow{diag}{fit\_nrtl}{bubble\_aad\_fit}{0.771501146805}
\goldrow{diag}{fit\_nrtl}{bubble\_max\_fit}{3.07957791309}
\goldrow{diag}{fit\_nrtl}{chi2}{122.54657753}
\goldrow{diag}{fit\_nrtl}{chi2\_in\_sample}{122.54657753}
\goldrow{diag}{fit\_nrtl}{chi2\_reduced}{1.55122250038}
\goldrow{diag}{fit\_nrtl}{cond\_JtJ}{16.8974869544}
\goldrow{diag}{fit\_nrtl}{converged}{1}
\goldrow{diag}{fit\_nrtl}{corr\_computed}{1}
\goldrow{diag}{fit\_nrtl}{dof}{79}
\goldrow{diag}{fit\_nrtl}{identifiable}{1}
\goldrow{diag}{fit\_nrtl}{invertible}{1}
\goldrow{diag}{fit\_nrtl}{iter}{5}
\goldrow{diag}{fit\_nrtl}{max\_abs\_corr}{0.95247136628}
\goldrow{diag}{fit\_nrtl}{max\_abs\_resid}{3.07957791309}
\goldmore{62}
\end{tutgolden}

\clearpage
\tutentry{curate121\_nrtl\_cyclohexane\_cyclohexene}{tutorials/props/curation/curate121\_nrtl\_cyclohexane\_cyclohexene}
\begin{tutscheme}
\begin{tikzpicture}[>={Stealth[length=2mm]}, every node/.style={font=\scriptsize}]
  \node[draw, rounded corners=2pt, align=center, minimum width=3.1cm, minimum height=0.8cm, fill=black!4] at (0.0,-0.0) {\textbf{fit\_nrtl}\\[-1pt]{\tiny fitParameters}};
  \node[draw, rounded corners=2pt, align=center, minimum width=3.1cm, minimum height=0.8cm, fill=black!4] at (3.6,-0.0) {\textbf{fit\_nrtl\_start2}\\[-1pt]{\tiny fitParameters}};
  \node[draw, rounded corners=2pt, align=center, minimum width=3.1cm, minimum height=0.8cm, fill=black!4] at (7.2,-0.0) {\textbf{txy\_ideal}\\[-1pt]{\tiny propertyScan1D}};
  \node[draw, rounded corners=2pt, align=center, minimum width=3.1cm, minimum height=0.8cm, fill=black!4] at (0.0,-1.3) {\textbf{txy\_unifac}\\[-1pt]{\tiny propertyScan1D}};
  \node[draw, rounded corners=2pt, align=center, minimum width=3.1cm, minimum height=0.8cm, fill=black!4] at (3.6,-1.3) {\textbf{txy\_nrtl}\\[-1pt]{\tiny propertyScan1D}};
\end{tikzpicture}
\end{tutscheme}
\tutwhat{NRTL cyclohexane-cyclohexene regressed by the curation pipeline to every laboratory's measured bubble temperatures (ThermoML locates, the articles source), accepted at max |dT| < 1 K}
\begin{tutnarrative}
AN NRTL PAIR FROM EVERY LABORATORY'S MEASUREMENTS.

WRITTEN BY bin/choupo-curate pair cyclohexane cyclohexene (DEV.md 4c C284, the bulk run; record docs/design/binary-pairs-from-open-measurements.md section 11).

THE DECLARATION -- WRITTEN AND COMMITTED BEFORE ANY FIT RAN (several declarations of the bulk run share one commit).

IDENTITY. cyclohexane (CAS 110-82-7, InChIKey XDTMQSROBMDMFD-UHFFFAOYSA-N) + cyclohexene (CAS 110-83-8, InChIKey HGCIXCUEYOPUTN-UHFFFAOYSA-N), joined to the archive through identity.json -- never by name. CLASS (C284): both records are volatile molecular liquids (no role nonvolatile/solute, no dissociatesTo, no charge, no solid, a vapour-pressure model), and every measured temperature lies below both critical temperatures.

NO MEASURED SERIES IN THIS TREE. constant/archiveDatasets names each dataset (DOI, NIST URL, article sha256, extraction selector, extract sha256); `bin/choupo-curate fetch` writes the points into constant/archiveExtracts/, which is gitignored.

STUDIES: every bubble-temperature block of the pair at or below 110 kPa with a declared uncertainty (2 block(s)). FITTED, ALL OF THEM (2): 10.1016/j.fluid.2008.12.007 block 10 63 pts 3 isobars 30-101.3 kPa 318.11-355.94 K U 0.1-0.27 K 10.1021/je049902w block 10 18 pts 12 isobars 98.9-101.2 kPa 353.40-356.10 K U 1-1 K (reports x of the other member: written as x[cyclohexane] = 1 - x by --x-of) SET ASIDE BY THE EXTRACTION RULES, BY NAME: (none)

THE RULE (DEV.md 4c C283 amended, C284). Every laboratory above is fitted; none is held out (the C16 partition and its band are retired as the acceptance test). NRTL, a\_ij = a\_ji = 0, alpha = 0.30 fixed: two parameters (b\_ij, b\_ji); each point weighted by 1/U, U its own stated expanded uncertainty. Objective twoStage (least squares, then minimax; T\_bubble absolute, in K). OUTLIERS: after a least-squares stage over every point, a point with |T\_model - T| > 3 K is discarded BY NAME (`discard` on its evidence entry, listed below) and the fit repeated without it. ACCEPTED when max |T\_model - T| over the KEPT points is below 1 K. C283's 5 \% relative clause judges composition; this fit has no composition residual, so only the temperature clause applies. A pair that splits ONE liquid inside its validity span is not promoted (arm f). Two starts: (0, 0) and (800, -400).

OUTLIERS DISCARDED AFTER THE LEAST-SQUARES STAGE: @@OUTLIERS-BEGIN@@ none: every one of the 81 points lies within 3 K after the least-squares stage. @@OUTLIERS-END@@

@@RESULT-BEGIN@@ WHAT THE RUN FOUND (pinned in `expected`; measured, not argued). fit\_nrtl: b\_ij = 601.5 K, b\_ji = -370.9 K over 81 kept points of 2 laboratories (0 discarded): max |dT| 1.624 K, mean 0.248 K against the acceptance 1 K $\to$ notAccepted. fit\_nrtl\_start2 lands on the same pair (601.5, -370.9). NOT PROMOTABLE: max |dT| 1.624 K over the 81 kept points is not below 1 K @@RESULT-END@@

WHAT IS NOT CLAIMED. The pair is correct at these pressures and in this frame (ideal-gas vapour, the catalogue's own vapour-pressure records); the measurements were transcribed by the archive and have NOT been read back against the articles (the CHECKLIST that `fetch` writes is the aid).

\end{tutnarrative}
\begin{tutgolden}
\goldrow{diag}{fit\_nrtl}{bubble\_aad\_fit}{0.247841045624}
\goldrow{diag}{fit\_nrtl}{bubble\_max\_fit}{1.62425112667}
\goldrow{diag}{fit\_nrtl}{chi2}{193.248618398}
\goldrow{diag}{fit\_nrtl}{chi2\_in\_sample}{193.248618398}
\goldrow{diag}{fit\_nrtl}{chi2\_reduced}{2.44618504302}
\goldrow{diag}{fit\_nrtl}{cond\_JtJ}{52.4584254712}
\goldrow{diag}{fit\_nrtl}{converged}{1}
\goldrow{diag}{fit\_nrtl}{corr\_computed}{1}
\goldrow{diag}{fit\_nrtl}{dof}{79}
\goldrow{diag}{fit\_nrtl}{identifiable}{1}
\goldrow{diag}{fit\_nrtl}{invertible}{1}
\goldrow{diag}{fit\_nrtl}{iter}{16}
\goldrow{diag}{fit\_nrtl}{max\_abs\_corr}{0.99060278356}
\goldrow{diag}{fit\_nrtl}{max\_abs\_resid}{1.62425112667}
\goldmore{53}
\end{tutgolden}

\clearpage
\tutentry{curate122\_nrtl\_1Pentanol\_anisole}{tutorials/props/curation/curate122\_nrtl\_1Pentanol\_anisole}
\begin{tutscheme}
\begin{tikzpicture}[>={Stealth[length=2mm]}, every node/.style={font=\scriptsize}]
  \node[draw, rounded corners=2pt, align=center, minimum width=3.1cm, minimum height=0.8cm, fill=black!4] at (0.0,-0.0) {\textbf{fit\_nrtl}\\[-1pt]{\tiny fitParameters}};
  \node[draw, rounded corners=2pt, align=center, minimum width=3.1cm, minimum height=0.8cm, fill=black!4] at (3.6,-0.0) {\textbf{fit\_nrtl\_start2}\\[-1pt]{\tiny fitParameters}};
  \node[draw, rounded corners=2pt, align=center, minimum width=3.1cm, minimum height=0.8cm, fill=black!4] at (7.2,-0.0) {\textbf{txy\_ideal}\\[-1pt]{\tiny propertyScan1D}};
  \node[draw, rounded corners=2pt, align=center, minimum width=3.1cm, minimum height=0.8cm, fill=black!4] at (0.0,-1.3) {\textbf{txy\_nrtl}\\[-1pt]{\tiny propertyScan1D}};
\end{tikzpicture}
\end{tutscheme}
\tutwhat{NRTL 1Pentanol-anisole regressed by the curation pipeline to every laboratory's measured bubble temperatures (ThermoML locates, the articles source), accepted at max |dT| < 1 K}
\begin{tutnarrative}
AN NRTL PAIR FROM EVERY LABORATORY'S MEASUREMENTS.

WRITTEN BY bin/choupo-curate pair 1Pentanol anisole (DEV.md 4c C284, the bulk run; record docs/design/binary-pairs-from-open-measurements.md section 11).

THE DECLARATION -- WRITTEN AND COMMITTED BEFORE ANY FIT RAN (several declarations of the bulk run share one commit).

IDENTITY. 1Pentanol (CAS 71-41-0, InChIKey AMQJEAYHLZJPGS-UHFFFAOYSA-N) + anisole (CAS 100-66-3, InChIKey RDOXTESZEPMUJZ-UHFFFAOYSA-N), joined to the archive through identity.json -- never by name. CLASS (C284): both records are volatile molecular liquids (no role nonvolatile/solute, no dissociatesTo, no charge, no solid, a vapour-pressure model), and every measured temperature lies below both critical temperatures.

NO MEASURED SERIES IN THIS TREE. constant/archiveDatasets names each dataset (DOI, NIST URL, article sha256, extraction selector, extract sha256); `bin/choupo-curate fetch` writes the points into constant/archiveExtracts/, which is gitignored.

STUDIES: every bubble-temperature block of the pair at or below 110 kPa with a declared uncertainty (2 block(s)). FITTED, ALL OF THEM (2): 10.1021/je700594n block 14 40 pts 4 isobars 26.66-101.32 kPa 372.65-419.85 K U 0.13-0.26 K 10.1021/je034189g block 4 28 pts 101.32 kPa 410.06-422.80 K U 0.12-0.34 K SET ASIDE BY THE EXTRACTION RULES, BY NAME: (none)

THE RULE (DEV.md 4c C283 amended, C284). Every laboratory above is fitted; none is held out (the C16 partition and its band are retired as the acceptance test). NRTL, a\_ij = a\_ji = 0, alpha = 0.30 fixed: two parameters (b\_ij, b\_ji); each point weighted by 1/U, U its own stated expanded uncertainty. Objective twoStage (least squares, then minimax; T\_bubble absolute, in K). OUTLIERS: after a least-squares stage over every point, a point with |T\_model - T| > 3 K is discarded BY NAME (`discard` on its evidence entry, listed below) and the fit repeated without it. ACCEPTED when max |T\_model - T| over the KEPT points is below 1 K. C283's 5 \% relative clause judges composition; this fit has no composition residual, so only the temperature clause applies. A pair that splits ONE liquid inside its validity span is not promoted (arm f). Two starts: (0, 0) and (800, -400).

OUTLIERS DISCARDED AFTER THE LEAST-SQUARES STAGE: @@OUTLIERS-BEGIN@@ none: every one of the 68 points lies within 3 K after the least-squares stage. @@OUTLIERS-END@@

@@RESULT-BEGIN@@ WHAT THE RUN FOUND (pinned in `expected`; measured, not argued). fit\_nrtl: b\_ij = 143.6 K, b\_ji = 178.8 K over 68 kept points of 2 laboratories (0 discarded): max |dT| 0.884 K, mean 0.197 K against the acceptance 1 K $\to$ accepted. fit\_nrtl\_start2 lands on the same pair (143.6, 178.8). PROMOTED on 2026-11-10 into data/standards/parameters/NRTL/1Pentanol-anisole.dat as reviewStatus autoPromoted, humanReviewed false, because: max |dT| 0.884 K < 1 K over the 68 kept points, one liquid across the validity span @@RESULT-END@@

WHAT IS NOT CLAIMED. The pair is correct at these pressures and in this frame (ideal-gas vapour, the catalogue's own vapour-pressure records); the measurements were transcribed by the archive and have NOT been read back against the articles (the CHECKLIST that `fetch` writes is the aid).

\end{tutnarrative}
\begin{tutgolden}
\goldrow{diag}{fit\_nrtl}{bubble\_aad\_fit}{0.19731402643}
\goldrow{diag}{fit\_nrtl}{bubble\_max\_fit}{0.884027313968}
\goldrow{diag}{fit\_nrtl}{chi2}{193.935355194}
\goldrow{diag}{fit\_nrtl}{chi2\_in\_sample}{193.935355194}
\goldrow{diag}{fit\_nrtl}{chi2\_reduced}{2.93841447264}
\goldrow{diag}{fit\_nrtl}{cond\_JtJ}{154.311597704}
\goldrow{diag}{fit\_nrtl}{converged}{1}
\goldrow{diag}{fit\_nrtl}{corr\_computed}{1}
\goldrow{diag}{fit\_nrtl}{dof}{66}
\goldrow{diag}{fit\_nrtl}{identifiable}{1}
\goldrow{diag}{fit\_nrtl}{invertible}{1}
\goldrow{diag}{fit\_nrtl}{iter}{3}
\goldrow{diag}{fit\_nrtl}{max\_abs\_corr}{0.996317560615}
\goldrow{diag}{fit\_nrtl}{max\_abs\_resid}{0.884027313968}
\goldmore{56}
\end{tutgolden}

\clearpage
\tutentry{curate123\_nrtl\_ethylBenzene\_styrene}{tutorials/props/curation/curate123\_nrtl\_ethylBenzene\_styrene}
\begin{tutscheme}
\begin{tikzpicture}[>={Stealth[length=2mm]}, every node/.style={font=\scriptsize}]
  \node[draw, rounded corners=2pt, align=center, minimum width=3.1cm, minimum height=0.8cm, fill=black!4] at (0.0,-0.0) {\textbf{fit\_nrtl}\\[-1pt]{\tiny fitParameters}};
  \node[draw, rounded corners=2pt, align=center, minimum width=3.1cm, minimum height=0.8cm, fill=black!4] at (3.6,-0.0) {\textbf{fit\_nrtl\_start2}\\[-1pt]{\tiny fitParameters}};
  \node[draw, rounded corners=2pt, align=center, minimum width=3.1cm, minimum height=0.8cm, fill=black!4] at (7.2,-0.0) {\textbf{txy\_ideal}\\[-1pt]{\tiny propertyScan1D}};
  \node[draw, rounded corners=2pt, align=center, minimum width=3.1cm, minimum height=0.8cm, fill=black!4] at (0.0,-1.3) {\textbf{txy\_nrtl}\\[-1pt]{\tiny propertyScan1D}};
\end{tikzpicture}
\end{tutscheme}
\tutwhat{NRTL ethylBenzene-styrene regressed by the curation pipeline to every laboratory's measured bubble temperatures (ThermoML locates, the articles source), accepted at max |dT| < 1 K}
\begin{tutnarrative}
AN NRTL PAIR FROM EVERY LABORATORY'S MEASUREMENTS.

WRITTEN BY bin/choupo-curate pair ethylBenzene styrene (DEV.md 4c C284, the bulk run; record docs/design/binary-pairs-from-open-measurements.md section 11).

THE DECLARATION -- WRITTEN AND COMMITTED BEFORE ANY FIT RAN (several declarations of the bulk run share one commit).

IDENTITY. ethylBenzene (CAS 100-41-4, InChIKey YNQLUTRBYVCPMQ-UHFFFAOYSA-N) + styrene (CAS 100-42-5, InChIKey PPBRXRYQALVLMV-UHFFFAOYSA-N), joined to the archive through identity.json -- never by name. CLASS (C284): both records are volatile molecular liquids (no role nonvolatile/solute, no dissociatesTo, no charge, no solid, a vapour-pressure model), and every measured temperature lies below both critical temperatures.

NO MEASURED SERIES IN THIS TREE. constant/archiveDatasets names each dataset (DOI, NIST URL, article sha256, extraction selector, extract sha256); `bin/choupo-curate fetch` writes the points into constant/archiveExtracts/, which is gitignored.

STUDIES: every bubble-temperature block of the pair at or below 110 kPa with a declared uncertainty (2 block(s)). FITTED, ALL OF THEM (2): 10.1021/je200098d block 2 36 pts 3 isobars 5-20 kPa 324.50-366.00 K U 0.3-0.5 K 10.1021/je050523s block 15 35 pts 2 isobars 5-15 kPa 324.38-357.84 K U 0.27-0.29 K SET ASIDE BY THE EXTRACTION RULES, BY NAME: (none)

THE RULE (DEV.md 4c C283 amended, C284). Every laboratory above is fitted; none is held out (the C16 partition and its band are retired as the acceptance test). NRTL, a\_ij = a\_ji = 0, alpha = 0.30 fixed: two parameters (b\_ij, b\_ji); each point weighted by 1/U, U its own stated expanded uncertainty. Objective twoStage (least squares, then minimax; T\_bubble absolute, in K). OUTLIERS: after a least-squares stage over every point, a point with |T\_model - T| > 3 K is discarded BY NAME (`discard` on its evidence entry, listed below) and the fit repeated without it. ACCEPTED when max |T\_model - T| over the KEPT points is below 1 K. C283's 5 \% relative clause judges composition; this fit has no composition residual, so only the temperature clause applies. A pair that splits ONE liquid inside its validity span is not promoted (arm f). Two starts: (0, 0) and (800, -400).

OUTLIERS DISCARDED AFTER THE LEAST-SQUARES STAGE: @@OUTLIERS-BEGIN@@ none: every one of the 71 points lies within 3 K after the least-squares stage. @@OUTLIERS-END@@

@@RESULT-BEGIN@@ WHAT THE RUN FOUND (pinned in `expected`; measured, not argued). fit\_nrtl: b\_ij = -292.4 K, b\_ji = 258.3 K over 71 kept points of 2 laboratories (0 discarded): max |dT| 1.919 K, mean 1.075 K against the acceptance 1 K $\to$ notAccepted. fit\_nrtl\_start2 lands on A DIFFERENT pair (1397.8, -566.7). NOT PROMOTABLE: max |dT| 1.919 K over the 71 kept points is not below 1 K @@RESULT-END@@

WHAT IS NOT CLAIMED. The pair is correct at these pressures and in this frame (ideal-gas vapour, the catalogue's own vapour-pressure records); the measurements were transcribed by the archive and have NOT been read back against the articles (the CHECKLIST that `fetch` writes is the aid).

\end{tutnarrative}
\begin{tutgolden}
\goldrow{diag}{fit\_nrtl}{bubble\_aad\_fit}{1.07507591112}
\goldrow{diag}{fit\_nrtl}{bubble\_max\_fit}{1.91922768679}
\goldrow{diag}{fit\_nrtl}{chi2}{1204.20954194}
\goldrow{diag}{fit\_nrtl}{chi2\_in\_sample}{1204.20954194}
\goldrow{diag}{fit\_nrtl}{chi2\_reduced}{17.4523122021}
\goldrow{diag}{fit\_nrtl}{cond\_JtJ}{6971.01131135}
\goldrow{diag}{fit\_nrtl}{converged}{1}
\goldrow{diag}{fit\_nrtl}{corr\_computed}{1}
\goldrow{diag}{fit\_nrtl}{dof}{69}
\goldrow{diag}{fit\_nrtl}{identifiable}{0}
\goldrow{diag}{fit\_nrtl}{invertible}{1}
\goldrow{diag}{fit\_nrtl}{iter}{11}
\goldrow{diag}{fit\_nrtl}{max\_abs\_corr}{0.99981356453}
\goldrow{diag}{fit\_nrtl}{max\_abs\_resid}{1.91922768679}
\goldmore{50}
\end{tutgolden}

\clearpage
\tutentry{curate124\_nrtl\_oCresol\_ethanol}{tutorials/props/curation/curate124\_nrtl\_oCresol\_ethanol}
\begin{tutscheme}
\begin{tikzpicture}[>={Stealth[length=2mm]}, every node/.style={font=\scriptsize}]
  \node[draw, rounded corners=2pt, align=center, minimum width=3.1cm, minimum height=0.8cm, fill=black!4] at (0.0,-0.0) {\textbf{fit\_nrtl}\\[-1pt]{\tiny fitParameters}};
  \node[draw, rounded corners=2pt, align=center, minimum width=3.1cm, minimum height=0.8cm, fill=black!4] at (3.6,-0.0) {\textbf{fit\_nrtl\_start2}\\[-1pt]{\tiny fitParameters}};
  \node[draw, rounded corners=2pt, align=center, minimum width=3.1cm, minimum height=0.8cm, fill=black!4] at (7.2,-0.0) {\textbf{txy\_ideal}\\[-1pt]{\tiny propertyScan1D}};
  \node[draw, rounded corners=2pt, align=center, minimum width=3.1cm, minimum height=0.8cm, fill=black!4] at (0.0,-1.3) {\textbf{txy\_nrtl}\\[-1pt]{\tiny propertyScan1D}};
\end{tikzpicture}
\end{tutscheme}
\tutwhat{NRTL oCresol-ethanol regressed by the curation pipeline to every laboratory's measured bubble temperatures (ThermoML locates, the articles source), accepted at max |dT| < 1 K}
\begin{tutnarrative}
AN NRTL PAIR FROM EVERY LABORATORY'S MEASUREMENTS.

WRITTEN BY bin/choupo-curate pair oCresol ethanol (DEV.md 4c C284, the bulk run; record docs/design/binary-pairs-from-open-measurements.md section 11).

THE DECLARATION -- WRITTEN AND COMMITTED BEFORE ANY FIT RAN (several declarations of the bulk run share one commit).

IDENTITY. oCresol (CAS 95-48-7, InChIKey QWVGKYWNOKOFNN-UHFFFAOYSA-N) + ethanol (CAS 64-17-5, InChIKey LFQSCWFLJHTTHZ-UHFFFAOYSA-N), joined to the archive through identity.json -- never by name. CLASS (C284): both records are volatile molecular liquids (no role nonvolatile/solute, no dissociatesTo, no charge, no solid, a vapour-pressure model), and every measured temperature lies below both critical temperatures.

NO MEASURED SERIES IN THIS TREE. constant/archiveDatasets names each dataset (DOI, NIST URL, article sha256, extraction selector, extract sha256); `bin/choupo-curate fetch` writes the points into constant/archiveExtracts/, which is gitignored.

STUDIES: every bubble-temperature block of the pair at or below 110 kPa with a declared uncertainty (2 block(s)). FITTED, ALL OF THEM (2): 10.1016/j.fluid.2017.09.008 block 16 55 pts 12 isobars 9-99.9 kPa 302.23-394.11 K U 1.17-3.31 K 10.1016/j.fluid.2006.03.018 block 34 12 pts 95.76 kPa 351.35-434.85 K U 1.01-1.01 K (reports x of the other member: written as x[oCresol] = 1 - x by --x-of) SET ASIDE BY THE EXTRACTION RULES, BY NAME: (none)

THE RULE (DEV.md 4c C283 amended, C284). Every laboratory above is fitted; none is held out (the C16 partition and its band are retired as the acceptance test). NRTL, a\_ij = a\_ji = 0, alpha = 0.30 fixed: two parameters (b\_ij, b\_ji); each point weighted by 1/U, U its own stated expanded uncertainty. Objective twoStage (least squares, then minimax; T\_bubble absolute, in K). OUTLIERS: after a least-squares stage over every point, a point with |T\_model - T| > 3 K is discarded BY NAME (`discard` on its evidence entry, listed below) and the fit repeated without it. ACCEPTED when max |T\_model - T| over the KEPT points is below 1 K. C283's 5 \% relative clause judges composition; this fit has no composition residual, so only the temperature clause applies. A pair that splits ONE liquid inside its validity span is not promoted (arm f). Two starts: (0, 0) and (800, -400).

OUTLIERS DISCARDED AFTER THE LEAST-SQUARES STAGE: @@OUTLIERS-BEGIN@@ 11 point(s) of 67: doi 10.1016/j.fluid.2017.09.008 row 39: x 0.5970, P 55.100 kPa, T 367.82 K, dT -3.15 K; doi 10.1016/j.fluid.2017.09.008 row 40: x 0.5970, P 75.000 kPa, T 376.10 K, dT -3.34 K; doi 10.1016/j.fluid.2017.09.008 row 41: x 0.5970, P 85.100 kPa, T 379.52 K, dT -3.33 K; doi 10.1016/j.fluid.2017.09.008 row 42: x 0.5970, P 98.500 kPa, T 384.56 K, dT -4.30 K; doi 10.1016/j.fluid.2017.09.008 row 47: x 0.6990, P 75.000 kPa, T 385.84 K, dT -3.45 K; doi 10.1016/j.fluid.2017.09.008 row 48: x 0.6990, P 85.100 kPa, T 389.62 K, dT -3.66 K; doi 10.1016/j.fluid.2017.09.008 row 49: x 0.6990, P 99.400 kPa, T 394.11 K, dT -3.64 K; doi 10.1016/j.fluid.2017.09.008 row 52: x 0.7910, P 35.100 kPa, T 375.90 K, dT -3.18 K; doi 10.1016/j.fluid.2017.09.008 row 53: x 0.7910, P 55.100 kPa, T 389.39 K, dT -4.52 K; doi 10.1016/j.fluid.2017.09.008 row 54: x 0.8980, P 9.000 kPa, T 359.99 K, dT -3.91 K; doi 10.1016/j.fluid.2017.09.008 row 55: x 0.8980, P 15.000 kPa, T 373.10 K, dT -4.88 K @@OUTLIERS-END@@

@@RESULT-BEGIN@@ WHAT THE RUN FOUND (pinned in `expected`; measured, not argued). fit\_nrtl: b\_ij = -693.3 K, b\_ji = 1272.4 K over 56 kept points of 2 laboratories (11 discarded): max |dT| 4.030 K, mean 1.658 K against the acceptance 1 K $\to$ notAccepted. fit\_nrtl\_start2 lands on A DIFFERENT pair (993.3, -656.6). NOT PROMOTABLE: max |dT| 4.030 K over the 56 kept points is not below 1 K @@RESULT-END@@

WHAT IS NOT CLAIMED. The pair is correct at these pressures and in this frame (ideal-gas vapour, the catalogue's own vapour-pressure records); the measurements were transcribed by the archive and have NOT been read back against the articles (the CHECKLIST that `fetch` writes is the aid).

\end{tutnarrative}
\begin{tutgolden}
\goldrow{diag}{fit\_nrtl}{bubble\_aad\_fit}{1.65813685599}
\goldrow{diag}{fit\_nrtl}{bubble\_max\_fit}{4.02955760965}
\goldrow{diag}{fit\_nrtl}{chi2}{70.5495016543}
\goldrow{diag}{fit\_nrtl}{chi2\_in\_sample}{70.5495016543}
\goldrow{diag}{fit\_nrtl}{chi2\_reduced}{1.30647225286}
\goldrow{diag}{fit\_nrtl}{cond\_JtJ}{194.771699039}
\goldrow{diag}{fit\_nrtl}{converged}{1}
\goldrow{diag}{fit\_nrtl}{corr\_computed}{1}
\goldrow{diag}{fit\_nrtl}{dof}{54}
\goldrow{diag}{fit\_nrtl}{identifiable}{1}
\goldrow{diag}{fit\_nrtl}{invertible}{1}
\goldrow{diag}{fit\_nrtl}{iter}{13}
\goldrow{diag}{fit\_nrtl}{max\_abs\_corr}{0.845869037624}
\goldrow{diag}{fit\_nrtl}{max\_abs\_resid}{4.02955760965}
\goldmore{50}
\end{tutgolden}

\clearpage
\tutentry{curate125\_nrtl\_ethanol\_2Methyl1Butanol}{tutorials/props/curation/curate125\_nrtl\_ethanol\_2Methyl1Butanol}
\begin{tutscheme}
\begin{tikzpicture}[>={Stealth[length=2mm]}, every node/.style={font=\scriptsize}]
  \node[draw, rounded corners=2pt, align=center, minimum width=3.1cm, minimum height=0.8cm, fill=black!4] at (0.0,-0.0) {\textbf{fit\_nrtl}\\[-1pt]{\tiny fitParameters}};
  \node[draw, rounded corners=2pt, align=center, minimum width=3.1cm, minimum height=0.8cm, fill=black!4] at (3.6,-0.0) {\textbf{fit\_nrtl\_start2}\\[-1pt]{\tiny fitParameters}};
  \node[draw, rounded corners=2pt, align=center, minimum width=3.1cm, minimum height=0.8cm, fill=black!4] at (7.2,-0.0) {\textbf{txy\_ideal}\\[-1pt]{\tiny propertyScan1D}};
  \node[draw, rounded corners=2pt, align=center, minimum width=3.1cm, minimum height=0.8cm, fill=black!4] at (0.0,-1.3) {\textbf{txy\_unifac}\\[-1pt]{\tiny propertyScan1D}};
  \node[draw, rounded corners=2pt, align=center, minimum width=3.1cm, minimum height=0.8cm, fill=black!4] at (3.6,-1.3) {\textbf{txy\_nrtl}\\[-1pt]{\tiny propertyScan1D}};
\end{tikzpicture}
\end{tutscheme}
\tutwhat{NRTL ethanol-2Methyl1Butanol regressed by the curation pipeline to every laboratory's measured bubble temperatures (ThermoML locates, the articles source), accepted at max |dT| < 1 K}
\begin{tutnarrative}
AN NRTL PAIR FROM EVERY LABORATORY'S MEASUREMENTS.

WRITTEN BY bin/choupo-curate pair ethanol 2Methyl1Butanol (DEV.md 4c C284, the bulk run; record docs/design/binary-pairs-from-open-measurements.md section 11).

THE DECLARATION -- WRITTEN AND COMMITTED BEFORE ANY FIT RAN (several declarations of the bulk run share one commit).

IDENTITY. ethanol (CAS 64-17-5, InChIKey LFQSCWFLJHTTHZ-UHFFFAOYSA-N) + 2Methyl1Butanol (CAS 137-32-6, InChIKey QPRQEDXDYOZYLA-UHFFFAOYSA-N), joined to the archive through identity.json -- never by name. CLASS (C284): both records are volatile molecular liquids (no role nonvolatile/solute, no dissociatesTo, no charge, no solid, a vapour-pressure model), and every measured temperature lies below both critical temperatures.

NO MEASURED SERIES IN THIS TREE. constant/archiveDatasets names each dataset (DOI, NIST URL, article sha256, extraction selector, extract sha256); `bin/choupo-curate fetch` writes the points into constant/archiveExtracts/, which is gitignored.

STUDIES: every bubble-temperature block of the pair at or below 110 kPa with a declared uncertainty (2 block(s)). FITTED, ALL OF THEM (2): 10.1021/je010272x block 4 39 pts 3 isobars 33.3-101.3 kPa 327.65-393.45 K U 0.2-0.4 K 10.1021/je0500330 block 19 20 pts 101.3 kPa 351.61-398.77 K U 0.29-1.31 K SET ASIDE BY THE EXTRACTION RULES, BY NAME: (none)

THE RULE (DEV.md 4c C283 amended, C284). Every laboratory above is fitted; none is held out (the C16 partition and its band are retired as the acceptance test). NRTL, a\_ij = a\_ji = 0, alpha = 0.30 fixed: two parameters (b\_ij, b\_ji); each point weighted by 1/U, U its own stated expanded uncertainty. Objective twoStage (least squares, then minimax; T\_bubble absolute, in K). OUTLIERS: after a least-squares stage over every point, a point with |T\_model - T| > 3 K is discarded BY NAME (`discard` on its evidence entry, listed below) and the fit repeated without it. ACCEPTED when max |T\_model - T| over the KEPT points is below 1 K. C283's 5 \% relative clause judges composition; this fit has no composition residual, so only the temperature clause applies. A pair that splits ONE liquid inside its validity span is not promoted (arm f). Two starts: (0, 0) and (800, -400).

OUTLIERS DISCARDED AFTER THE LEAST-SQUARES STAGE: @@OUTLIERS-BEGIN@@ none: every one of the 59 points lies within 3 K after the least-squares stage. @@OUTLIERS-END@@

@@RESULT-BEGIN@@ WHAT THE RUN FOUND (pinned in `expected`; measured, not argued). fit\_nrtl: b\_ij = 444.7 K, b\_ji = -300.1 K over 59 kept points of 2 laboratories (0 discarded): max |dT| 0.719 K, mean 0.264 K against the acceptance 1 K $\to$ accepted. fit\_nrtl\_start2 lands on the same pair (444.7, -300.1). PROMOTED on 2026-11-10 into data/standards/parameters/NRTL/2Methyl1Butanol-ethanol.dat as reviewStatus autoPromoted, humanReviewed false, because: max |dT| 0.719 K < 1 K over the 59 kept points, one liquid across the validity span @@RESULT-END@@

WHAT IS NOT CLAIMED. The pair is correct at these pressures and in this frame (ideal-gas vapour, the catalogue's own vapour-pressure records); the measurements were transcribed by the archive and have NOT been read back against the articles (the CHECKLIST that `fetch` writes is the aid).

\end{tutnarrative}
\begin{tutgolden}
\goldrow{diag}{fit\_nrtl}{bubble\_aad\_fit}{0.26364900524}
\goldrow{diag}{fit\_nrtl}{bubble\_max\_fit}{0.718717896523}
\goldrow{diag}{fit\_nrtl}{chi2}{44.3092818429}
\goldrow{diag}{fit\_nrtl}{chi2\_in\_sample}{44.3092818429}
\goldrow{diag}{fit\_nrtl}{chi2\_reduced}{0.777355821805}
\goldrow{diag}{fit\_nrtl}{cond\_JtJ}{186.095271219}
\goldrow{diag}{fit\_nrtl}{converged}{1}
\goldrow{diag}{fit\_nrtl}{corr\_computed}{1}
\goldrow{diag}{fit\_nrtl}{dof}{57}
\goldrow{diag}{fit\_nrtl}{identifiable}{1}
\goldrow{diag}{fit\_nrtl}{invertible}{1}
\goldrow{diag}{fit\_nrtl}{iter}{11}
\goldrow{diag}{fit\_nrtl}{max\_abs\_corr}{0.997323984209}
\goldrow{diag}{fit\_nrtl}{max\_abs\_resid}{0.718717896523}
\goldmore{62}
\end{tutgolden}

\clearpage
\tutentry{curate126\_nrtl\_nButanol\_dimethylCarbonate}{tutorials/props/curation/curate126\_nrtl\_nButanol\_dimethylCarbonate}
\begin{tutscheme}
\begin{tikzpicture}[>={Stealth[length=2mm]}, every node/.style={font=\scriptsize}]
  \node[draw, rounded corners=2pt, align=center, minimum width=3.1cm, minimum height=0.8cm, fill=black!4] at (0.0,-0.0) {\textbf{fit\_nrtl}\\[-1pt]{\tiny fitParameters}};
  \node[draw, rounded corners=2pt, align=center, minimum width=3.1cm, minimum height=0.8cm, fill=black!4] at (3.6,-0.0) {\textbf{fit\_nrtl\_start2}\\[-1pt]{\tiny fitParameters}};
  \node[draw, rounded corners=2pt, align=center, minimum width=3.1cm, minimum height=0.8cm, fill=black!4] at (7.2,-0.0) {\textbf{txy\_ideal}\\[-1pt]{\tiny propertyScan1D}};
  \node[draw, rounded corners=2pt, align=center, minimum width=3.1cm, minimum height=0.8cm, fill=black!4] at (0.0,-1.3) {\textbf{txy\_nrtl}\\[-1pt]{\tiny propertyScan1D}};
\end{tikzpicture}
\end{tutscheme}
\tutwhat{NRTL nButanol-dimethylCarbonate regressed by the curation pipeline to every laboratory's measured bubble temperatures (ThermoML locates, the articles source), accepted at max |dT| < 1 K}
\begin{tutnarrative}
AN NRTL PAIR FROM EVERY LABORATORY'S MEASUREMENTS.

WRITTEN BY bin/choupo-curate pair nButanol dimethylCarbonate (DEV.md 4c C284, the bulk run; record docs/design/binary-pairs-from-open-measurements.md section 11).

THE DECLARATION -- WRITTEN AND COMMITTED BEFORE ANY FIT RAN (several declarations of the bulk run share one commit).

IDENTITY. nButanol (CAS 71-36-3, InChIKey LRHPLDYGYMQRHN-UHFFFAOYSA-N) + dimethylCarbonate (CAS 616-38-6, InChIKey IEJIGPNLZYLLBP-UHFFFAOYSA-N), joined to the archive through identity.json -- never by name. CLASS (C284): both records are volatile molecular liquids (no role nonvolatile/solute, no dissociatesTo, no charge, no solid, a vapour-pressure model), and every measured temperature lies below both critical temperatures.

NO MEASURED SERIES IN THIS TREE. constant/archiveDatasets names each dataset (DOI, NIST URL, article sha256, extraction selector, extract sha256); `bin/choupo-curate fetch` writes the points into constant/archiveExtracts/, which is gitignored.

STUDIES: every bubble-temperature block of the pair at or below 110 kPa with a declared uncertainty (2 block(s)). FITTED, ALL OF THEM (2): 10.1016/j.jct.2011.08.006 block 7 45 pts 5 isobars 40-93.32 kPa 337.45-378.75 K U 0.12-0.22 K 10.1021/je030131q block 26 9 pts 95.8 kPa 362.05-375.95 K U 0.32-0.34 K (reports x of the other member: written as x[nButanol] = 1 - x by --x-of) SET ASIDE BY THE EXTRACTION RULES, BY NAME: (none)

THE RULE (DEV.md 4c C283 amended, C284). Every laboratory above is fitted; none is held out (the C16 partition and its band are retired as the acceptance test). NRTL, a\_ij = a\_ji = 0, alpha = 0.30 fixed: two parameters (b\_ij, b\_ji); each point weighted by 1/U, U its own stated expanded uncertainty. Objective twoStage (least squares, then minimax; T\_bubble absolute, in K). OUTLIERS: after a least-squares stage over every point, a point with |T\_model - T| > 3 K is discarded BY NAME (`discard` on its evidence entry, listed below) and the fit repeated without it. ACCEPTED when max |T\_model - T| over the KEPT points is below 1 K. C283's 5 \% relative clause judges composition; this fit has no composition residual, so only the temperature clause applies. A pair that splits ONE liquid inside its validity span is not promoted (arm f). Two starts: (0, 0) and (800, -400).

OUTLIERS DISCARDED AFTER THE LEAST-SQUARES STAGE: @@OUTLIERS-BEGIN@@ 2 point(s) of 54: doi 10.1021/je030131q row 1: x 0.9090, P 95.800 kPa, T 375.95 K, dT +4.73 K; doi 10.1021/je030131q row 2: x 0.8349, P 95.800 kPa, T 372.35 K, dT +3.70 K @@OUTLIERS-END@@

@@RESULT-BEGIN@@ WHAT THE RUN FOUND (pinned in `expected`; measured, not argued). fit\_nrtl: b\_ij = 75.9 K, b\_ji = 309.1 K over 52 kept points of 2 laboratories (2 discarded): max |dT| 2.151 K, mean 0.373 K against the acceptance 1 K $\to$ notAccepted. fit\_nrtl\_start2 lands on the same pair (75.9, 309.1). NOT PROMOTABLE: max |dT| 2.151 K over the 52 kept points is not below 1 K @@RESULT-END@@

WHAT IS NOT CLAIMED. The pair is correct at these pressures and in this frame (ideal-gas vapour, the catalogue's own vapour-pressure records); the measurements were transcribed by the archive and have NOT been read back against the articles (the CHECKLIST that `fetch` writes is the aid).

\end{tutnarrative}
\begin{tutgolden}
\goldrow{diag}{fit\_nrtl}{bubble\_aad\_fit}{0.373170603696}
\goldrow{diag}{fit\_nrtl}{bubble\_max\_fit}{2.15071595923}
\goldrow{diag}{fit\_nrtl}{chi2}{405.010951098}
\goldrow{diag}{fit\_nrtl}{chi2\_in\_sample}{405.010951098}
\goldrow{diag}{fit\_nrtl}{chi2\_reduced}{8.10021902196}
\goldrow{diag}{fit\_nrtl}{cond\_JtJ}{68.61740274}
\goldrow{diag}{fit\_nrtl}{converged}{1}
\goldrow{diag}{fit\_nrtl}{corr\_computed}{1}
\goldrow{diag}{fit\_nrtl}{dof}{50}
\goldrow{diag}{fit\_nrtl}{identifiable}{1}
\goldrow{diag}{fit\_nrtl}{invertible}{1}
\goldrow{diag}{fit\_nrtl}{iter}{4}
\goldrow{diag}{fit\_nrtl}{max\_abs\_corr}{0.989794901441}
\goldrow{diag}{fit\_nrtl}{max\_abs\_resid}{2.15071595923}
\goldmore{50}
\end{tutgolden}

\clearpage
\tutentry{curate127\_nrtl\_methylPropionate\_nHeptane}{tutorials/props/curation/curate127\_nrtl\_methylPropionate\_nHeptane}
\begin{tutscheme}
\begin{tikzpicture}[>={Stealth[length=2mm]}, every node/.style={font=\scriptsize}]
  \node[draw, rounded corners=2pt, align=center, minimum width=3.1cm, minimum height=0.8cm, fill=black!4] at (0.0,-0.0) {\textbf{fit\_nrtl}\\[-1pt]{\tiny fitParameters}};
  \node[draw, rounded corners=2pt, align=center, minimum width=3.1cm, minimum height=0.8cm, fill=black!4] at (3.6,-0.0) {\textbf{fit\_nrtl\_start2}\\[-1pt]{\tiny fitParameters}};
  \node[draw, rounded corners=2pt, align=center, minimum width=3.1cm, minimum height=0.8cm, fill=black!4] at (7.2,-0.0) {\textbf{txy\_ideal}\\[-1pt]{\tiny propertyScan1D}};
  \node[draw, rounded corners=2pt, align=center, minimum width=3.1cm, minimum height=0.8cm, fill=black!4] at (0.0,-1.3) {\textbf{txy\_nrtl}\\[-1pt]{\tiny propertyScan1D}};
\end{tikzpicture}
\end{tutscheme}
\tutwhat{NRTL methylPropionate-nHeptane regressed by the curation pipeline to every laboratory's measured bubble temperatures (ThermoML locates, the articles source), accepted at max |dT| < 1 K}
\begin{tutnarrative}
AN NRTL PAIR FROM EVERY LABORATORY'S MEASUREMENTS.

WRITTEN BY bin/choupo-curate pair methylPropionate nHeptane (DEV.md 4c C284, the bulk run; record docs/design/binary-pairs-from-open-measurements.md section 11).

THE DECLARATION -- WRITTEN AND COMMITTED BEFORE ANY FIT RAN (several declarations of the bulk run share one commit).

IDENTITY. methylPropionate (CAS 554-12-1, InChIKey RJUFJBKOKNCXHH-UHFFFAOYSA-N) + nHeptane (CAS 142-82-5, InChIKey IMNFDUFMRHMDMM-UHFFFAOYSA-N), joined to the archive through identity.json -- never by name. CLASS (C284): both records are volatile molecular liquids (no role nonvolatile/solute, no dissociatesTo, no charge, no solid, a vapour-pressure model), and every measured temperature lies below both critical temperatures.

NO MEASURED SERIES IN THIS TREE. constant/archiveDatasets names each dataset (DOI, NIST URL, article sha256, extraction selector, extract sha256); `bin/choupo-curate fetch` writes the points into constant/archiveExtracts/, which is gitignored.

STUDIES: every bubble-temperature block of the pair at or below 110 kPa with a declared uncertainty (2 block(s)). FITTED, ALL OF THEM (2): 10.1021/je4009415 block 20 39 pts 101.32 kPa 351.86-369.65 K U 0.11-0.39 K 10.1021/je030117d block 32 25 pts 101.325 kPa 350.54-370.26 K U 0.32-0.39 K SET ASIDE BY THE EXTRACTION RULES, BY NAME: (none)

THE RULE (DEV.md 4c C283 amended, C284). Every laboratory above is fitted; none is held out (the C16 partition and its band are retired as the acceptance test). NRTL, a\_ij = a\_ji = 0, alpha = 0.30 fixed: two parameters (b\_ij, b\_ji); each point weighted by 1/U, U its own stated expanded uncertainty. Objective twoStage (least squares, then minimax; T\_bubble absolute, in K). OUTLIERS: after a least-squares stage over every point, a point with |T\_model - T| > 3 K is discarded BY NAME (`discard` on its evidence entry, listed below) and the fit repeated without it. ACCEPTED when max |T\_model - T| over the KEPT points is below 1 K. C283's 5 \% relative clause judges composition; this fit has no composition residual, so only the temperature clause applies. A pair that splits ONE liquid inside its validity span is not promoted (arm f). Two starts: (0, 0) and (800, -400).

OUTLIERS DISCARDED AFTER THE LEAST-SQUARES STAGE: @@OUTLIERS-BEGIN@@ none: every one of the 64 points lies within 3 K after the least-squares stage. @@OUTLIERS-END@@

@@RESULT-BEGIN@@ WHAT THE RUN FOUND (pinned in `expected`; measured, not argued). fit\_nrtl: b\_ij = 381.2 K, b\_ji = -0.5 K over 64 kept points of 2 laboratories (0 discarded): max |dT| 1.849 K, mean 0.509 K against the acceptance 1 K $\to$ notAccepted. fit\_nrtl\_start2 lands on A DIFFERENT pair (378.0, 1.4). NOT PROMOTABLE: max |dT| 1.849 K over the 64 kept points is not below 1 K @@RESULT-END@@

WHAT IS NOT CLAIMED. The pair is correct at these pressures and in this frame (ideal-gas vapour, the catalogue's own vapour-pressure records); the measurements were transcribed by the archive and have NOT been read back against the articles (the CHECKLIST that `fetch` writes is the aid).

\end{tutnarrative}
\begin{tutgolden}
\goldrow{diag}{fit\_nrtl}{bubble\_aad\_fit}{0.508737899267}
\goldrow{diag}{fit\_nrtl}{bubble\_max\_fit}{1.84932175126}
\goldrow{diag}{fit\_nrtl}{chi2}{450.359059388}
\goldrow{diag}{fit\_nrtl}{chi2\_in\_sample}{450.359059388}
\goldrow{diag}{fit\_nrtl}{chi2\_reduced}{7.26385579659}
\goldrow{diag}{fit\_nrtl}{cond\_JtJ}{87.485290604}
\goldrow{diag}{fit\_nrtl}{converged}{1}
\goldrow{diag}{fit\_nrtl}{corr\_computed}{1}
\goldrow{diag}{fit\_nrtl}{dof}{62}
\goldrow{diag}{fit\_nrtl}{identifiable}{1}
\goldrow{diag}{fit\_nrtl}{invertible}{1}
\goldrow{diag}{fit\_nrtl}{iter}{5}
\goldrow{diag}{fit\_nrtl}{max\_abs\_corr}{0.994333287487}
\goldrow{diag}{fit\_nrtl}{max\_abs\_resid}{1.84932175126}
\goldmore{62}
\end{tutgolden}

\clearpage
\tutentry{curate128\_nrtl\_2Methyl1Propanol\_isobutylAcetate}{tutorials/props/curation/curate128\_nrtl\_2Methyl1Propanol\_isobutylAcetate}
\begin{tutscheme}
\begin{tikzpicture}[>={Stealth[length=2mm]}, every node/.style={font=\scriptsize}]
  \node[draw, rounded corners=2pt, align=center, minimum width=3.1cm, minimum height=0.8cm, fill=black!4] at (0.0,-0.0) {\textbf{fit\_nrtl}\\[-1pt]{\tiny fitParameters}};
  \node[draw, rounded corners=2pt, align=center, minimum width=3.1cm, minimum height=0.8cm, fill=black!4] at (3.6,-0.0) {\textbf{fit\_nrtl\_start2}\\[-1pt]{\tiny fitParameters}};
  \node[draw, rounded corners=2pt, align=center, minimum width=3.1cm, minimum height=0.8cm, fill=black!4] at (7.2,-0.0) {\textbf{txy\_ideal}\\[-1pt]{\tiny propertyScan1D}};
  \node[draw, rounded corners=2pt, align=center, minimum width=3.1cm, minimum height=0.8cm, fill=black!4] at (0.0,-1.3) {\textbf{txy\_unifac}\\[-1pt]{\tiny propertyScan1D}};
  \node[draw, rounded corners=2pt, align=center, minimum width=3.1cm, minimum height=0.8cm, fill=black!4] at (3.6,-1.3) {\textbf{txy\_nrtl}\\[-1pt]{\tiny propertyScan1D}};
\end{tikzpicture}
\end{tutscheme}
\tutwhat{NRTL 2Methyl1Propanol-isobutylAcetate regressed by the curation pipeline to every laboratory's measured bubble temperatures (ThermoML locates, the articles source), accepted at max |dT| < 1 K}
\begin{tutnarrative}
AN NRTL PAIR FROM EVERY LABORATORY'S MEASUREMENTS.

WRITTEN BY bin/choupo-curate pair 2Methyl1Propanol isobutylAcetate (DEV.md 4c C284, the bulk run; record docs/design/binary-pairs-from-open-measurements.md section 11).

THE DECLARATION -- WRITTEN AND COMMITTED BEFORE ANY FIT RAN (several declarations of the bulk run share one commit).

IDENTITY. 2Methyl1Propanol (CAS 78-83-1, InChIKey ZXEKIIBDNHEJCQ-UHFFFAOYSA-N) + isobutylAcetate (CAS 110-19-0, InChIKey GJRQTCIYDGXPES-UHFFFAOYSA-N), joined to the archive through identity.json -- never by name. CLASS (C284): both records are volatile molecular liquids (no role nonvolatile/solute, no dissociatesTo, no charge, no solid, a vapour-pressure model), and every measured temperature lies below both critical temperatures.

NO MEASURED SERIES IN THIS TREE. constant/archiveDatasets names each dataset (DOI, NIST URL, article sha256, extraction selector, extract sha256); `bin/choupo-curate fetch` writes the points into constant/archiveExtracts/, which is gitignored.

STUDIES: every bubble-temperature block of the pair at or below 110 kPa with a declared uncertainty (2 block(s)). FITTED, ALL OF THEM (2): 10.1016/j.fluid.2004.10.022 block 12 40 pts 2 isobars 20-100 kPa 339.13-388.48 K U 0.1-0.16 K 10.1021/acs.jced.7b00241 block 6 18 pts 101.3 kPa 380.32-388.55 K U 0.22-0.47 K SET ASIDE BY THE EXTRACTION RULES, BY NAME: (none)

THE RULE (DEV.md 4c C283 amended, C284). Every laboratory above is fitted; none is held out (the C16 partition and its band are retired as the acceptance test). NRTL, a\_ij = a\_ji = 0, alpha = 0.30 fixed: two parameters (b\_ij, b\_ji); each point weighted by 1/U, U its own stated expanded uncertainty. Objective twoStage (least squares, then minimax; T\_bubble absolute, in K). OUTLIERS: after a least-squares stage over every point, a point with |T\_model - T| > 3 K is discarded BY NAME (`discard` on its evidence entry, listed below) and the fit repeated without it. ACCEPTED when max |T\_model - T| over the KEPT points is below 1 K. C283's 5 \% relative clause judges composition; this fit has no composition residual, so only the temperature clause applies. A pair that splits ONE liquid inside its validity span is not promoted (arm f). Two starts: (0, 0) and (800, -400).

OUTLIERS DISCARDED AFTER THE LEAST-SQUARES STAGE: @@OUTLIERS-BEGIN@@ none: every one of the 58 points lies within 3 K after the least-squares stage. @@OUTLIERS-END@@

@@RESULT-BEGIN@@ WHAT THE RUN FOUND (pinned in `expected`; measured, not argued). fit\_nrtl: b\_ij = 170.4 K, b\_ji = 5.7 K over 58 kept points of 2 laboratories (0 discarded): max |dT| 1.317 K, mean 0.706 K against the acceptance 1 K $\to$ notAccepted. fit\_nrtl\_start2 lands on the same pair (170.4, 5.7). NOT PROMOTABLE: max |dT| 1.317 K over the 58 kept points is not below 1 K @@RESULT-END@@

WHAT IS NOT CLAIMED. The pair is correct at these pressures and in this frame (ideal-gas vapour, the catalogue's own vapour-pressure records); the measurements were transcribed by the archive and have NOT been read back against the articles (the CHECKLIST that `fetch` writes is the aid).

\end{tutnarrative}
\begin{tutgolden}
\goldrow{diag}{fit\_nrtl}{bubble\_aad\_fit}{0.705976245205}
\goldrow{diag}{fit\_nrtl}{bubble\_max\_fit}{1.31675076307}
\goldrow{diag}{fit\_nrtl}{chi2}{1970.65303473}
\goldrow{diag}{fit\_nrtl}{chi2\_in\_sample}{1970.65303473}
\goldrow{diag}{fit\_nrtl}{chi2\_reduced}{35.1902327631}
\goldrow{diag}{fit\_nrtl}{cond\_JtJ}{321.589422526}
\goldrow{diag}{fit\_nrtl}{converged}{1}
\goldrow{diag}{fit\_nrtl}{corr\_computed}{1}
\goldrow{diag}{fit\_nrtl}{dof}{56}
\goldrow{diag}{fit\_nrtl}{identifiable}{1}
\goldrow{diag}{fit\_nrtl}{invertible}{1}
\goldrow{diag}{fit\_nrtl}{iter}{7}
\goldrow{diag}{fit\_nrtl}{max\_abs\_corr}{0.998448839124}
\goldrow{diag}{fit\_nrtl}{max\_abs\_resid}{1.31675076307}
\goldmore{53}
\end{tutgolden}

\clearpage
\tutentry{curate129\_nrtl\_methylFormate\_nHeptane}{tutorials/props/curation/curate129\_nrtl\_methylFormate\_nHeptane}
\begin{tutscheme}
\begin{tikzpicture}[>={Stealth[length=2mm]}, every node/.style={font=\scriptsize}]
  \node[draw, rounded corners=2pt, align=center, minimum width=3.1cm, minimum height=0.8cm, fill=black!4] at (0.0,-0.0) {\textbf{fit\_nrtl}\\[-1pt]{\tiny fitParameters}};
  \node[draw, rounded corners=2pt, align=center, minimum width=3.1cm, minimum height=0.8cm, fill=black!4] at (3.6,-0.0) {\textbf{fit\_nrtl\_start2}\\[-1pt]{\tiny fitParameters}};
  \node[draw, rounded corners=2pt, align=center, minimum width=3.1cm, minimum height=0.8cm, fill=black!4] at (7.2,-0.0) {\textbf{txy\_ideal}\\[-1pt]{\tiny propertyScan1D}};
  \node[draw, rounded corners=2pt, align=center, minimum width=3.1cm, minimum height=0.8cm, fill=black!4] at (0.0,-1.3) {\textbf{txy\_nrtl}\\[-1pt]{\tiny propertyScan1D}};
\end{tikzpicture}
\end{tutscheme}
\tutwhat{NRTL methylFormate-nHeptane regressed by the curation pipeline to every laboratory's measured bubble temperatures (ThermoML locates, the articles source), accepted at max |dT| < 1 K}
\begin{tutnarrative}
AN NRTL PAIR FROM EVERY LABORATORY'S MEASUREMENTS.

WRITTEN BY bin/choupo-curate pair methylFormate nHeptane (DEV.md 4c C284, the bulk run; record docs/design/binary-pairs-from-open-measurements.md section 11).

THE DECLARATION -- WRITTEN AND COMMITTED BEFORE ANY FIT RAN (several declarations of the bulk run share one commit).

IDENTITY. methylFormate (CAS 107-31-3, InChIKey TZIHFWKZFHZASV-UHFFFAOYSA-N) + nHeptane (CAS 142-82-5, InChIKey IMNFDUFMRHMDMM-UHFFFAOYSA-N), joined to the archive through identity.json -- never by name. CLASS (C284): both records are volatile molecular liquids (no role nonvolatile/solute, no dissociatesTo, no charge, no solid, a vapour-pressure model), and every measured temperature lies below both critical temperatures.

NO MEASURED SERIES IN THIS TREE. constant/archiveDatasets names each dataset (DOI, NIST URL, article sha256, extraction selector, extract sha256); `bin/choupo-curate fetch` writes the points into constant/archiveExtracts/, which is gitignored.

STUDIES: every bubble-temperature block of the pair at or below 110 kPa with a declared uncertainty (2 block(s)). FITTED, ALL OF THEM (2): 10.1021/je030117d block 22 36 pts 101.325 kPa 304.68-363.95 K U 0.25-3.24 K 10.1016/j.fluid.2015.05.031 block 30 24 pts 101.32 kPa 304.91-367.67 K U 0.09-2.47 K SET ASIDE BY THE EXTRACTION RULES, BY NAME: (none)

THE RULE (DEV.md 4c C283 amended, C284). Every laboratory above is fitted; none is held out (the C16 partition and its band are retired as the acceptance test). NRTL, a\_ij = a\_ji = 0, alpha = 0.30 fixed: two parameters (b\_ij, b\_ji); each point weighted by 1/U, U its own stated expanded uncertainty. Objective twoStage (least squares, then minimax; T\_bubble absolute, in K). OUTLIERS: after a least-squares stage over every point, a point with |T\_model - T| > 3 K is discarded BY NAME (`discard` on its evidence entry, listed below) and the fit repeated without it. ACCEPTED when max |T\_model - T| over the KEPT points is below 1 K. C283's 5 \% relative clause judges composition; this fit has no composition residual, so only the temperature clause applies. A pair that splits ONE liquid inside its validity span is not promoted (arm f). Two starts: (0, 0) and (800, -400).

OUTLIERS DISCARDED AFTER THE LEAST-SQUARES STAGE: @@OUTLIERS-BEGIN@@ 19 point(s) of 60: doi 10.1021/je030117d row 1: x 0.0101, P 101.325 kPa, T 363.95 K, dT +6.11 K; doi 10.1021/je030117d row 2: x 0.0163, P 101.325 kPa, T 358.95 K, dT +9.80 K; doi 10.1021/je030117d row 3: x 0.0255, P 101.325 kPa, T 355.79 K, dT +10.64 K; doi 10.1021/je030117d row 4: x 0.0295, P 101.325 kPa, T 353.00 K, dT +12.30 K; doi 10.1021/je030117d row 5: x 0.0400, P 101.325 kPa, T 349.38 K, dT +12.73 K; doi 10.1021/je030117d row 6: x 0.0543, P 101.325 kPa, T 344.78 K, dT +12.64 K; doi 10.1021/je030117d row 7: x 0.0700, P 101.325 kPa, T 340.75 K, dT +11.39 K; doi 10.1021/je030117d row 8: x 0.0887, P 101.325 kPa, T 337.55 K, dT +8.51 K; doi 10.1021/je030117d row 9: x 0.1068, P 101.325 kPa, T 334.15 K, dT +6.46 K; doi 10.1021/je030117d row 10: x 0.1250, P 101.325 kPa, T 331.20 K, dT +4.50 K; doi 10.1021/je030117d row 11: x 0.1371, P 101.325 kPa, T 329.22 K, dT +3.54 K; doi 10.1016/j.fluid.2015.05.031 row 1: x 0.0060, P 101.320 kPa, T 367.67 K, dT +3.11 K; doi 10.1016/j.fluid.2015.05.031 row 2: x 0.0130, P 101.320 kPa, T 363.52 K, dT +5.96 K; doi 10.1016/j.fluid.2015.05.031 row 3: x 0.0230, P 101.320 kPa, T 358.34 K, dT +8.76 K; doi 10.1016/j.fluid.2015.05.031 row 4: x 0.0330, P 101.320 kPa, T 353.86 K, dT +10.41 K; doi 10.1016/j.fluid.2015.05.031 row 5: x 0.0420, P 101.320 kPa, T 350.79 K, dT +10.68 K; doi 10.1016/j.fluid.2015.05.031 row 6: x 0.0620, P 101.320 kPa, T 344.59 K, dT +10.24 K; doi 10.1016/j.fluid.2015.05.031 row 7: x 0.0920, P 101.320 kPa, T 337.26 K, dT +7.76 K; doi 10.1016/j.fluid.2015.05.031 row 8: x 0.1110, P 101.320 kPa, T 332.97 K, dT +6.46 K @@OUTLIERS-END@@

@@RESULT-BEGIN@@ WHAT THE RUN FOUND (pinned in `expected`; measured, not argued). fit\_nrtl: b\_ij = 710.3 K, b\_ji = 158.2 K over 41 kept points of 2 laboratories (19 discarded): max |dT| 5.372 K, mean 1.181 K against the acceptance 1 K $\to$ notAccepted. fit\_nrtl\_start2 lands on the same pair (710.3, 158.2). NOT PROMOTABLE: max |dT| 5.372 K over the 41 kept points is not below 1 K @@RESULT-END@@

WHAT IS NOT CLAIMED. The pair is correct at these pressures and in this frame (ideal-gas vapour, the catalogue's own vapour-pressure records); the measurements were transcribed by the archive and have NOT been read back against the articles (the CHECKLIST that `fetch` writes is the aid).

\end{tutnarrative}
\begin{tutgolden}
\goldrow{diag}{fit\_nrtl}{bubble\_aad\_fit}{1.18089807013}
\goldrow{diag}{fit\_nrtl}{bubble\_max\_fit}{5.37236714957}
\goldrow{diag}{fit\_nrtl}{chi2}{1113.72160623}
\goldrow{diag}{fit\_nrtl}{chi2\_in\_sample}{1113.72160623}
\goldrow{diag}{fit\_nrtl}{chi2\_reduced}{28.5569642622}
\goldrow{diag}{fit\_nrtl}{cond\_JtJ}{3.63618057036}
\goldrow{diag}{fit\_nrtl}{converged}{1}
\goldrow{diag}{fit\_nrtl}{corr\_computed}{1}
\goldrow{diag}{fit\_nrtl}{dof}{39}
\goldrow{diag}{fit\_nrtl}{identifiable}{1}
\goldrow{diag}{fit\_nrtl}{invertible}{1}
\goldrow{diag}{fit\_nrtl}{iter}{6}
\goldrow{diag}{fit\_nrtl}{max\_abs\_corr}{0.885610077189}
\goldrow{diag}{fit\_nrtl}{max\_abs\_resid}{5.37236714957}
\goldmore{62}
\end{tutgolden}

\clearpage
\tutentry{curate12\_nrtl\_ethanol\_cyclohexane}{tutorials/props/curation/curate12\_nrtl\_ethanol\_cyclohexane}
\begin{tutscheme}
\begin{tikzpicture}[>={Stealth[length=2mm]}, every node/.style={font=\scriptsize}]
  \node[draw, rounded corners=2pt, align=center, minimum width=3.1cm, minimum height=0.8cm, fill=black!4] at (0.0,-0.0) {\textbf{fit\_nrtl}\\[-1pt]{\tiny fitParameters}};
  \node[draw, rounded corners=2pt, align=center, minimum width=3.1cm, minimum height=0.8cm, fill=black!4] at (3.6,-0.0) {\textbf{fit\_nrtl\_start2}\\[-1pt]{\tiny fitParameters}};
  \node[draw, rounded corners=2pt, align=center, minimum width=3.1cm, minimum height=0.8cm, fill=black!4] at (7.2,-0.0) {\textbf{txy\_ideal}\\[-1pt]{\tiny propertyScan1D}};
  \node[draw, rounded corners=2pt, align=center, minimum width=3.1cm, minimum height=0.8cm, fill=black!4] at (0.0,-1.3) {\textbf{txy\_unifac}\\[-1pt]{\tiny propertyScan1D}};
  \node[draw, rounded corners=2pt, align=center, minimum width=3.1cm, minimum height=0.8cm, fill=black!4] at (3.6,-1.3) {\textbf{txy\_nrtl}\\[-1pt]{\tiny propertyScan1D}};
  \node[draw, rounded corners=2pt, align=center, minimum width=3.1cm, minimum height=0.8cm, fill=black!4] at (7.2,-1.3) {\textbf{txy\_ideal\_fitP}\\[-1pt]{\tiny propertyScan1D}};
  \node[draw, rounded corners=2pt, align=center, minimum width=3.1cm, minimum height=0.8cm, fill=black!4] at (0.0,-2.6) {\textbf{txy\_unifac\_fitP}\\[-1pt]{\tiny propertyScan1D}};
  \node[draw, rounded corners=2pt, align=center, minimum width=3.1cm, minimum height=0.8cm, fill=black!4] at (3.6,-2.6) {\textbf{txy\_nrtl\_fitP}\\[-1pt]{\tiny propertyScan1D}};
\end{tikzpicture}
\end{tutscheme}
\tutwhat{NRTL ethanol-cyclohexane regressed to measured bubble temperatures (ThermoML locates, the article sources), held out against an independent study, drawn against both beside ideal and UNIFAC}
\begin{tutnarrative}
AN NRTL PAIR FROM PUBLISHED MEASUREMENTS, END TO END.

WHY THIS SYSTEM. the binary under the industrial dehydration of ethanol by heterogeneous azeotropic distillation, where cyclohexane is the entrainer: an alcohol + hydrocarbon pair with a strong positive deviation and a minimum-boiling azeotrope near 338 K (both studies).

THE CHAIN, EACH STEP A COMMAND (DEV.md C16; record docs/design/binary-pairs-from-open-measurements.md): 1. bin/choupo-thermoml sync the NIST/TRC ThermoML Archive (ThermoML.v2020-09-30.tgz, doi 10.18434/mds2-2422, sha256-verified) 2. bin/choupo-thermoml identity CAS $\to$ InChIKey for the catalogue, two resolvers that must agree 3. bin/choupo-thermoml vle-index every binary VLE block, searchable 4. bin/choupo-thermoml vle-search ethanol cyclohexane --geometry bubbleT 5. bin/choupo-thermoml extract-vle <file> --block N --interior-only $\to$ the two files under constant/experiments/ 6. choupoProps (this case) fitParameters $\to$ a proposal record 7. promotion (a human act) $\to$ NOT promoted -- see WHAT THE RUN FOUND below (validated, but it splits the liquid) 8. bin/choupo-import the seal mirrors nothing for the pair (not promoted)

THE EVIDENCE. FIT: doi 10.1016/j.fluid.2018.07.013 (25 interior points, 101.3 kPa). HELD OUT: doi 10.1016/j.fluid.2006.01.023 (14 interior points, 101.325 kPa) -- an article whose title names ethanol + water with ionic liquids; the ethanol + cyclohexane block is the one the archive files under it, and reading it back is part of the curator's check. Neither study reports a pure-component bubble point. Chosen BEFORE any fit by the slice-2 partition rule: the FIT set is the low-pressure bubble- temperature study with the most interior points; the HELD-OUT set is the next largest that shares NO author with it (the archive carries no affiliations, so a shared author is the dependence the rule can check). Each file names its article, its archive file and block, the sha256 of the archive file, the two compounds' CAS and InChIKey, and the per-point expanded uncertainty; each says `reviewStatus transcribedNotCheckedAgainstArticle` because nobody has read the numbers back against the paper yet, and the run announces it. Pure-component endpoints are set aside (they measure the vapour-pressure record, not the pair) and kept in each file under `selection`.

THE FIT. NRTL with a\_ij = a\_ji = 0 and alpha = 0.30 fixed: TWO parameters (b\_ij, b\_ji), because bubble temperatures over the ranges below do not identify more (fit\_nrtl\_ab in the methanol + water case asks that question directly). Each point is weighted by 1/U, U its OWN stated expanded uncertainty (`weighting uncertainty;`), so a point measured to 0.1 K is not outvoted by one measured to 1 K. The held-out AAD is plain kelvin, unweighted.

WHAT THE RUN FOUND (pinned in `expected`; measured, not argued). fit\_nrtl: b\_ij = 354.7 K, b\_ji = 565.4 K, in-sample rms 1.30 K (reduced chi2 5.5); fit\_nrtl\_start2 lands on the same pair. HELD OUT (101.325 kPa, no author in common): AAD 0.32 K (0.094 \%) against a band of 0.41 \% $\to$ `validated`. Overlay on the held-out set: NRTL 0.33 K, UNIFAC 0.59 K, ideal 9.2 K. AND NOT PROMOTED. The same pair predicts a liquid-liquid split up to about 358.8 K -- across the WHOLE of its validity span (337.7-353.5 K), where both studies report one liquid (the archive files each as liquid + gas). The engine's VLLE flash at 340 K and 3 bar, run with this pair, splits a 43/57 feed into two liquids. A bubble-temperature scan solves for ONE liquid by construction, so the held-out verdict cannot see a second one: `validated` is necessary, not sufficient. The pair stays in this case as the authored record constant/parameters/NRTL/cyclohexane-ethanol.dat (verdict validated, reviewStatus interim); check\_regressed\_pairs arm (f) refuses it in the catalogue. The remedy -- a third parameter constrained by the phase count, or LLE data -- is a curation decision, Vitor's.

WHAT IS NOT CLAIMED. The pair is correct at these pressures and in this frame (ideal-gas vapour, the catalogue's Antoine records); nothing here tests it at high pressure, against LLE, or with another vapour model.

\end{tutnarrative}
\begin{tutgolden}
\goldrow{diag}{fit\_nrtl}{aad\_heldout\_K}{0.323443931811}
\goldrow{diag}{fit\_nrtl}{aad\_heldout\_pct}{0.093961584242}
\goldrow{diag}{fit\_nrtl}{chi2}{125.860527228}
\goldrow{diag}{fit\_nrtl}{chi2\_in\_sample}{125.860527228}
\goldrow{diag}{fit\_nrtl}{chi2\_reduced}{5.47219683598}
\goldrow{diag}{fit\_nrtl}{cond\_JtJ}{15.8520572952}
\goldrow{diag}{fit\_nrtl}{converged}{1}
\goldrow{diag}{fit\_nrtl}{corr\_computed}{1}
\goldrow{diag}{fit\_nrtl}{dof}{23}
\goldrow{diag}{fit\_nrtl}{identifiable}{1}
\goldrow{diag}{fit\_nrtl}{invertible}{1}
\goldrow{diag}{fit\_nrtl}{iter}{6}
\goldrow{diag}{fit\_nrtl}{max\_abs\_corr}{0.940983736236}
\goldrow{diag}{fit\_nrtl}{max\_abs\_resid}{2.85975730621}
\goldmore{70}
\end{tutgolden}

\clearpage
\tutentry{curate130\_nrtl\_14Dioxane\_cyclohexane}{tutorials/props/curation/curate130\_nrtl\_14Dioxane\_cyclohexane}
\begin{tutscheme}
\begin{tikzpicture}[>={Stealth[length=2mm]}, every node/.style={font=\scriptsize}]
  \node[draw, rounded corners=2pt, align=center, minimum width=3.1cm, minimum height=0.8cm, fill=black!4] at (0.0,-0.0) {\textbf{fit\_nrtl}\\[-1pt]{\tiny fitParameters}};
  \node[draw, rounded corners=2pt, align=center, minimum width=3.1cm, minimum height=0.8cm, fill=black!4] at (3.6,-0.0) {\textbf{fit\_nrtl\_start2}\\[-1pt]{\tiny fitParameters}};
  \node[draw, rounded corners=2pt, align=center, minimum width=3.1cm, minimum height=0.8cm, fill=black!4] at (7.2,-0.0) {\textbf{txy\_ideal}\\[-1pt]{\tiny propertyScan1D}};
  \node[draw, rounded corners=2pt, align=center, minimum width=3.1cm, minimum height=0.8cm, fill=black!4] at (0.0,-1.3) {\textbf{txy\_unifac}\\[-1pt]{\tiny propertyScan1D}};
  \node[draw, rounded corners=2pt, align=center, minimum width=3.1cm, minimum height=0.8cm, fill=black!4] at (3.6,-1.3) {\textbf{txy\_nrtl}\\[-1pt]{\tiny propertyScan1D}};
\end{tikzpicture}
\end{tutscheme}
\tutwhat{NRTL 14Dioxane-cyclohexane regressed by the curation pipeline to every laboratory's measured bubble temperatures (ThermoML locates, the articles source), accepted at max |dT| < 1 K}
\begin{tutnarrative}
AN NRTL PAIR FROM EVERY LABORATORY'S MEASUREMENTS.

WRITTEN BY bin/choupo-curate pair 14Dioxane cyclohexane (DEV.md 4c C284, the bulk run; record docs/design/binary-pairs-from-open-measurements.md section 11).

THE DECLARATION -- WRITTEN AND COMMITTED BEFORE ANY FIT RAN (several declarations of the bulk run share one commit).

IDENTITY. 14Dioxane (CAS 123-91-1, InChIKey RYHBNJHYFVUHQT-UHFFFAOYSA-N) + cyclohexane (CAS 110-82-7, InChIKey XDTMQSROBMDMFD-UHFFFAOYSA-N), joined to the archive through identity.json -- never by name. CLASS (C284): both records are volatile molecular liquids (no role nonvolatile/solute, no dissociatesTo, no charge, no solid, a vapour-pressure model), and every measured temperature lies below both critical temperatures.

NO MEASURED SERIES IN THIS TREE. constant/archiveDatasets names each dataset (DOI, NIST URL, article sha256, extraction selector, extract sha256); `bin/choupo-curate fetch` writes the points into constant/archiveExtracts/, which is gitignored.

STUDIES: every bubble-temperature block of the pair at or below 110 kPa with a declared uncertainty (2 block(s)). FITTED, ALL OF THEM (2): 10.1016/j.fluid.2005.09.010 block 10 42 pts 3 isobars 40-101.3 kPa 325.16-369.18 K U 0.46-0.5 K 10.1021/je020148t block 31 8 pts 94.7 kPa 351.45-360.85 K U 1.04-1.04 K (reports x of the other member: written as x[14Dioxane] = 1 - x by --x-of) SET ASIDE BY THE EXTRACTION RULES, BY NAME: (none)

THE RULE (DEV.md 4c C283 amended, C284). Every laboratory above is fitted; none is held out (the C16 partition and its band are retired as the acceptance test). NRTL, a\_ij = a\_ji = 0, alpha = 0.30 fixed: two parameters (b\_ij, b\_ji); each point weighted by 1/U, U its own stated expanded uncertainty. Objective twoStage (least squares, then minimax; T\_bubble absolute, in K). OUTLIERS: after a least-squares stage over every point, a point with |T\_model - T| > 3 K is discarded BY NAME (`discard` on its evidence entry, listed below) and the fit repeated without it. ACCEPTED when max |T\_model - T| over the KEPT points is below 1 K. C283's 5 \% relative clause judges composition; this fit has no composition residual, so only the temperature clause applies. A pair that splits ONE liquid inside its validity span is not promoted (arm f). Two starts: (0, 0) and (800, -400).

OUTLIERS DISCARDED AFTER THE LEAST-SQUARES STAGE: @@OUTLIERS-BEGIN@@ 1 point(s) of 50: doi 10.1021/je020148t row 1: x 0.9430, P 94.700 kPa, T 360.85 K, dT +4.86 K @@OUTLIERS-END@@

@@RESULT-BEGIN@@ WHAT THE RUN FOUND (pinned in `expected`; measured, not argued). fit\_nrtl: b\_ij = 306.7 K, b\_ji = 144.3 K over 49 kept points of 2 laboratories (1 discarded): max |dT| 2.016 K, mean 0.411 K against the acceptance 1 K $\to$ notAccepted. fit\_nrtl\_start2 lands on the same pair (306.7, 144.3). NOT PROMOTABLE: max |dT| 2.016 K over the 49 kept points is not below 1 K @@RESULT-END@@

WHAT IS NOT CLAIMED. The pair is correct at these pressures and in this frame (ideal-gas vapour, the catalogue's own vapour-pressure records); the measurements were transcribed by the archive and have NOT been read back against the articles (the CHECKLIST that `fetch` writes is the aid).

\end{tutnarrative}
\begin{tutgolden}
\goldrow{diag}{fit\_nrtl}{bubble\_aad\_fit}{0.411327270668}
\goldrow{diag}{fit\_nrtl}{bubble\_max\_fit}{2.01585822031}
\goldrow{diag}{fit\_nrtl}{chi2}{40.2863051603}
\goldrow{diag}{fit\_nrtl}{chi2\_in\_sample}{40.2863051603}
\goldrow{diag}{fit\_nrtl}{chi2\_reduced}{0.857155428942}
\goldrow{diag}{fit\_nrtl}{cond\_JtJ}{25.3428674946}
\goldrow{diag}{fit\_nrtl}{converged}{1}
\goldrow{diag}{fit\_nrtl}{corr\_computed}{1}
\goldrow{diag}{fit\_nrtl}{dof}{47}
\goldrow{diag}{fit\_nrtl}{identifiable}{1}
\goldrow{diag}{fit\_nrtl}{invertible}{1}
\goldrow{diag}{fit\_nrtl}{iter}{4}
\goldrow{diag}{fit\_nrtl}{max\_abs\_corr}{0.980835898849}
\goldrow{diag}{fit\_nrtl}{max\_abs\_resid}{2.01585822031}
\goldmore{53}
\end{tutgolden}

\clearpage
\tutentry{curate131\_nrtl\_thiophene\_nHeptane}{tutorials/props/curation/curate131\_nrtl\_thiophene\_nHeptane}
\begin{tutscheme}
\begin{tikzpicture}[>={Stealth[length=2mm]}, every node/.style={font=\scriptsize}]
  \node[draw, rounded corners=2pt, align=center, minimum width=3.1cm, minimum height=0.8cm, fill=black!4] at (0.0,-0.0) {\textbf{fit\_nrtl}\\[-1pt]{\tiny fitParameters}};
  \node[draw, rounded corners=2pt, align=center, minimum width=3.1cm, minimum height=0.8cm, fill=black!4] at (3.6,-0.0) {\textbf{fit\_nrtl\_start2}\\[-1pt]{\tiny fitParameters}};
  \node[draw, rounded corners=2pt, align=center, minimum width=3.1cm, minimum height=0.8cm, fill=black!4] at (7.2,-0.0) {\textbf{txy\_ideal}\\[-1pt]{\tiny propertyScan1D}};
  \node[draw, rounded corners=2pt, align=center, minimum width=3.1cm, minimum height=0.8cm, fill=black!4] at (0.0,-1.3) {\textbf{txy\_nrtl}\\[-1pt]{\tiny propertyScan1D}};
\end{tikzpicture}
\end{tutscheme}
\tutwhat{NRTL thiophene-nHeptane regressed by the curation pipeline to every laboratory's measured bubble temperatures (ThermoML locates, the articles source), accepted at max |dT| < 1 K}
\begin{tutnarrative}
AN NRTL PAIR FROM EVERY LABORATORY'S MEASUREMENTS.

WRITTEN BY bin/choupo-curate pair thiophene nHeptane (DEV.md 4c C284, the bulk run; record docs/design/binary-pairs-from-open-measurements.md section 11).

THE DECLARATION -- WRITTEN AND COMMITTED BEFORE ANY FIT RAN (several declarations of the bulk run share one commit).

IDENTITY. thiophene (CAS 110-02-1, InChIKey YTPLMLYBLZKORZ-UHFFFAOYSA-N) + nHeptane (CAS 142-82-5, InChIKey IMNFDUFMRHMDMM-UHFFFAOYSA-N), joined to the archive through identity.json -- never by name. CLASS (C284): both records are volatile molecular liquids (no role nonvolatile/solute, no dissociatesTo, no charge, no solid, a vapour-pressure model), and every measured temperature lies below both critical temperatures.

NO MEASURED SERIES IN THIS TREE. constant/archiveDatasets names each dataset (DOI, NIST URL, article sha256, extraction selector, extract sha256); `bin/choupo-curate fetch` writes the points into constant/archiveExtracts/, which is gitignored.

STUDIES: every bubble-temperature block of the pair at or below 110 kPa with a declared uncertainty (2 block(s)). FITTED, ALL OF THEM (1): 10.1016/j.fluid.2011.11.020 block 8 18 pts 101.33 kPa 356.34-368.69 K U 0.09-0.11 K SET ASIDE BY THE EXTRACTION RULES, BY NAME: - doi 10.1016/j.fluid.2013.05.019 block 9: extract-vle refused it: This extractor supports fixed-pressure binary bubble-point evidence (constant P, liquid mole fraction x1, bubble temperature T). j.fluid.2013.05.019.xml: none of this file's 1 dataset block(s) is a binary bubble-point series:

THE RULE (DEV.md 4c C283 amended, C284). Every laboratory above is fitted; none is held out (the C16 partition and its band are retired as the acceptance test). NRTL, a\_ij = a\_ji = 0, alpha = 0.30 fixed: two parameters (b\_ij, b\_ji); each point weighted by 1/U, U its own stated expanded uncertainty. Objective twoStage (least squares, then minimax; T\_bubble absolute, in K). OUTLIERS: after a least-squares stage over every point, a point with |T\_model - T| > 3 K is discarded BY NAME (`discard` on its evidence entry, listed below) and the fit repeated without it. ACCEPTED when max |T\_model - T| over the KEPT points is below 1 K. C283's 5 \% relative clause judges composition; this fit has no composition residual, so only the temperature clause applies. A pair that splits ONE liquid inside its validity span is not promoted (arm f). Two starts: (0, 0) and (800, -400).

OUTLIERS DISCARDED AFTER THE LEAST-SQUARES STAGE: @@OUTLIERS-BEGIN@@ none: every one of the 18 points lies within 3 K after the least-squares stage. @@OUTLIERS-END@@

@@RESULT-BEGIN@@ WHAT THE RUN FOUND (pinned in `expected`; measured, not argued). fit\_nrtl: b\_ij = 675.2 K, b\_ji = -252.3 K over 18 kept points of 1 laboratories (0 discarded): max |dT| 0.394 K, mean 0.239 K against the acceptance 1 K $\to$ accepted. fit\_nrtl\_start2 lands on the same pair (675.2, -252.3). PROMOTED on 2026-11-10 into data/standards/parameters/NRTL/nHeptane-thiophene.dat as reviewStatus autoPromoted, humanReviewed false, because: max |dT| 0.394 K < 1 K over the 18 kept points, one liquid across the validity span @@RESULT-END@@

WHAT IS NOT CLAIMED. The pair is correct at these pressures and in this frame (ideal-gas vapour, the catalogue's own vapour-pressure records); the measurements were transcribed by the archive and have NOT been read back against the articles (the CHECKLIST that `fetch` writes is the aid).

\end{tutnarrative}
\begin{tutgolden}
\goldrow{diag}{fit\_nrtl}{bubble\_aad\_fit}{0.239174281445}
\goldrow{diag}{fit\_nrtl}{bubble\_max\_fit}{0.394412443095}
\goldrow{diag}{fit\_nrtl}{chi2}{129.344311742}
\goldrow{diag}{fit\_nrtl}{chi2\_in\_sample}{129.344311742}
\goldrow{diag}{fit\_nrtl}{chi2\_reduced}{8.08401948387}
\goldrow{diag}{fit\_nrtl}{cond\_JtJ}{27.8288481444}
\goldrow{diag}{fit\_nrtl}{converged}{1}
\goldrow{diag}{fit\_nrtl}{corr\_computed}{1}
\goldrow{diag}{fit\_nrtl}{dof}{16}
\goldrow{diag}{fit\_nrtl}{identifiable}{1}
\goldrow{diag}{fit\_nrtl}{invertible}{1}
\goldrow{diag}{fit\_nrtl}{iter}{13}
\goldrow{diag}{fit\_nrtl}{max\_abs\_corr}{0.982503192484}
\goldrow{diag}{fit\_nrtl}{max\_abs\_resid}{0.394412443095}
\goldmore{56}
\end{tutgolden}

\clearpage
\tutentry{curate132\_nrtl\_1Propanol\_propionicAcid}{tutorials/props/curation/curate132\_nrtl\_1Propanol\_propionicAcid}
\begin{tutscheme}
\begin{tikzpicture}[>={Stealth[length=2mm]}, every node/.style={font=\scriptsize}]
  \node[draw, rounded corners=2pt, align=center, minimum width=3.1cm, minimum height=0.8cm, fill=black!4] at (0.0,-0.0) {\textbf{fit\_nrtl}\\[-1pt]{\tiny fitParameters}};
  \node[draw, rounded corners=2pt, align=center, minimum width=3.1cm, minimum height=0.8cm, fill=black!4] at (3.6,-0.0) {\textbf{fit\_nrtl\_start2}\\[-1pt]{\tiny fitParameters}};
  \node[draw, rounded corners=2pt, align=center, minimum width=3.1cm, minimum height=0.8cm, fill=black!4] at (7.2,-0.0) {\textbf{txy\_ideal}\\[-1pt]{\tiny propertyScan1D}};
  \node[draw, rounded corners=2pt, align=center, minimum width=3.1cm, minimum height=0.8cm, fill=black!4] at (0.0,-1.3) {\textbf{txy\_unifac}\\[-1pt]{\tiny propertyScan1D}};
  \node[draw, rounded corners=2pt, align=center, minimum width=3.1cm, minimum height=0.8cm, fill=black!4] at (3.6,-1.3) {\textbf{txy\_nrtl}\\[-1pt]{\tiny propertyScan1D}};
\end{tikzpicture}
\end{tutscheme}
\tutwhat{NRTL 1Propanol-propionicAcid regressed by the curation pipeline to every laboratory's measured bubble temperatures (ThermoML locates, the articles source), accepted at max |dT| < 1 K}
\begin{tutnarrative}
AN NRTL PAIR FROM EVERY LABORATORY'S MEASUREMENTS.

WRITTEN BY bin/choupo-curate pair 1Propanol propionicAcid (DEV.md 4c C284, the bulk run; record docs/design/binary-pairs-from-open-measurements.md section 11).

THE DECLARATION -- WRITTEN AND COMMITTED BEFORE ANY FIT RAN (several declarations of the bulk run share one commit).

IDENTITY. 1Propanol (CAS 71-23-8, InChIKey BDERNNFJNOPAEC-UHFFFAOYSA-N) + propionicAcid (CAS 79-09-4, InChIKey XBDQKXXYIPTUBI-UHFFFAOYSA-N), joined to the archive through identity.json -- never by name. CLASS (C284): both records are volatile molecular liquids (no role nonvolatile/solute, no dissociatesTo, no charge, no solid, a vapour-pressure model), and every measured temperature lies below both critical temperatures.

NO MEASURED SERIES IN THIS TREE. constant/archiveDatasets names each dataset (DOI, NIST URL, article sha256, extraction selector, extract sha256); `bin/choupo-curate fetch` writes the points into constant/archiveExtracts/, which is gitignored.

STUDIES: every bubble-temperature block of the pair at or below 110 kPa with a declared uncertainty (2 block(s)). FITTED, ALL OF THEM (2): 10.1021/je100912t block 11 24 pts 101.325 kPa 372.03-408.17 K U 0.22-0.29 K 10.1021/je101302p block 6 23 pts 101.3 kPa 371.27-413.55 K U 0.71-0.94 K SET ASIDE BY THE EXTRACTION RULES, BY NAME: (none)

THE RULE (DEV.md 4c C283 amended, C284). Every laboratory above is fitted; none is held out (the C16 partition and its band are retired as the acceptance test). NRTL, a\_ij = a\_ji = 0, alpha = 0.30 fixed: two parameters (b\_ij, b\_ji); each point weighted by 1/U, U its own stated expanded uncertainty. Objective twoStage (least squares, then minimax; T\_bubble absolute, in K). OUTLIERS: after a least-squares stage over every point, a point with |T\_model - T| > 3 K is discarded BY NAME (`discard` on its evidence entry, listed below) and the fit repeated without it. ACCEPTED when max |T\_model - T| over the KEPT points is below 1 K. C283's 5 \% relative clause judges composition; this fit has no composition residual, so only the temperature clause applies. A pair that splits ONE liquid inside its validity span is not promoted (arm f). Two starts: (0, 0) and (800, -400).

OUTLIERS DISCARDED AFTER THE LEAST-SQUARES STAGE: @@OUTLIERS-BEGIN@@ 5 point(s) of 47: doi 10.1021/je101302p row 5: x 0.0854, P 101.300 kPa, T 408.15 K, dT -3.08 K; doi 10.1021/je101302p row 6: x 0.1055, P 101.300 kPa, T 406.55 K, dT -3.18 K; doi 10.1021/je101302p row 7: x 0.1474, P 101.300 kPa, T 404.15 K, dT -3.95 K; doi 10.1021/je101302p row 8: x 0.2114, P 101.300 kPa, T 399.45 K, dT -3.34 K; doi 10.1021/je101302p row 9: x 0.2559, P 101.300 kPa, T 396.71 K, dT -3.04 K @@OUTLIERS-END@@

@@RESULT-BEGIN@@ WHAT THE RUN FOUND (pinned in `expected`; measured, not argued). fit\_nrtl: b\_ij = -384.8 K, b\_ji = 632.4 K over 42 kept points of 2 laboratories (5 discarded): max |dT| 2.916 K, mean 0.557 K against the acceptance 1 K $\to$ notAccepted. fit\_nrtl\_start2 lands on the same pair (-384.8, 632.4). NOT PROMOTABLE: max |dT| 2.916 K over the 42 kept points is not below 1 K @@RESULT-END@@

WHAT IS NOT CLAIMED. The pair is correct at these pressures and in this frame (ideal-gas vapour, the catalogue's own vapour-pressure records); the measurements were transcribed by the archive and have NOT been read back against the articles (the CHECKLIST that `fetch` writes is the aid).

\end{tutnarrative}
\begin{tutgolden}
\goldrow{diag}{fit\_nrtl}{bubble\_aad\_fit}{0.557423426986}
\goldrow{diag}{fit\_nrtl}{bubble\_max\_fit}{2.91580077035}
\goldrow{diag}{fit\_nrtl}{chi2}{131.377851739}
\goldrow{diag}{fit\_nrtl}{chi2\_in\_sample}{131.377851739}
\goldrow{diag}{fit\_nrtl}{chi2\_reduced}{3.28444629347}
\goldrow{diag}{fit\_nrtl}{cond\_JtJ}{187.561019549}
\goldrow{diag}{fit\_nrtl}{converged}{1}
\goldrow{diag}{fit\_nrtl}{corr\_computed}{1}
\goldrow{diag}{fit\_nrtl}{dof}{40}
\goldrow{diag}{fit\_nrtl}{identifiable}{1}
\goldrow{diag}{fit\_nrtl}{invertible}{1}
\goldrow{diag}{fit\_nrtl}{iter}{11}
\goldrow{diag}{fit\_nrtl}{max\_abs\_corr}{0.984720642179}
\goldrow{diag}{fit\_nrtl}{max\_abs\_resid}{2.91580077035}
\goldmore{68}
\end{tutgolden}

\clearpage
\tutentry{curate133\_nrtl\_methanol\_2Butanol}{tutorials/props/curation/curate133\_nrtl\_methanol\_2Butanol}
\begin{tutscheme}
\begin{tikzpicture}[>={Stealth[length=2mm]}, every node/.style={font=\scriptsize}]
  \node[draw, rounded corners=2pt, align=center, minimum width=3.1cm, minimum height=0.8cm, fill=black!4] at (0.0,-0.0) {\textbf{fit\_nrtl}\\[-1pt]{\tiny fitParameters}};
  \node[draw, rounded corners=2pt, align=center, minimum width=3.1cm, minimum height=0.8cm, fill=black!4] at (3.6,-0.0) {\textbf{fit\_nrtl\_start2}\\[-1pt]{\tiny fitParameters}};
  \node[draw, rounded corners=2pt, align=center, minimum width=3.1cm, minimum height=0.8cm, fill=black!4] at (7.2,-0.0) {\textbf{txy\_ideal}\\[-1pt]{\tiny propertyScan1D}};
  \node[draw, rounded corners=2pt, align=center, minimum width=3.1cm, minimum height=0.8cm, fill=black!4] at (0.0,-1.3) {\textbf{txy\_unifac}\\[-1pt]{\tiny propertyScan1D}};
  \node[draw, rounded corners=2pt, align=center, minimum width=3.1cm, minimum height=0.8cm, fill=black!4] at (3.6,-1.3) {\textbf{txy\_nrtl}\\[-1pt]{\tiny propertyScan1D}};
\end{tikzpicture}
\end{tutscheme}
\tutwhat{NRTL methanol-2Butanol regressed by the curation pipeline to every laboratory's measured bubble temperatures (ThermoML locates, the articles source), accepted at max |dT| < 1 K}
\begin{tutnarrative}
AN NRTL PAIR FROM EVERY LABORATORY'S MEASUREMENTS.

WRITTEN BY bin/choupo-curate pair methanol 2Butanol (DEV.md 4c C284, the bulk run; record docs/design/binary-pairs-from-open-measurements.md section 11).

THE DECLARATION -- WRITTEN AND COMMITTED BEFORE ANY FIT RAN (several declarations of the bulk run share one commit).

IDENTITY. methanol (CAS 67-56-1, InChIKey OKKJLVBELUTLKV-UHFFFAOYSA-N) + 2Butanol (CAS 78-92-2, InChIKey BTANRVKWQNVYAZ-UHFFFAOYSA-N), joined to the archive through identity.json -- never by name. CLASS (C284): both records are volatile molecular liquids (no role nonvolatile/solute, no dissociatesTo, no charge, no solid, a vapour-pressure model), and every measured temperature lies below both critical temperatures.

NO MEASURED SERIES IN THIS TREE. constant/archiveDatasets names each dataset (DOI, NIST URL, article sha256, extraction selector, extract sha256); `bin/choupo-curate fetch` writes the points into constant/archiveExtracts/, which is gitignored.

STUDIES: every bubble-temperature block of the pair at or below 110 kPa with a declared uncertainty (2 block(s)). FITTED, ALL OF THEM (2): 10.1016/j.jct.2013.12.021 block 3 20 pts 101.33 kPa 337.95-368.15 K U 0.13-0.13 K 10.1021/acs.jced.8b00499 block 17 20 pts 101.3 kPa 340.73-367.83 K U 0.29-0.65 K SET ASIDE BY THE EXTRACTION RULES, BY NAME: (none)

THE RULE (DEV.md 4c C283 amended, C284). Every laboratory above is fitted; none is held out (the C16 partition and its band are retired as the acceptance test). NRTL, a\_ij = a\_ji = 0, alpha = 0.30 fixed: two parameters (b\_ij, b\_ji); each point weighted by 1/U, U its own stated expanded uncertainty. Objective twoStage (least squares, then minimax; T\_bubble absolute, in K). OUTLIERS: after a least-squares stage over every point, a point with |T\_model - T| > 3 K is discarded BY NAME (`discard` on its evidence entry, listed below) and the fit repeated without it. ACCEPTED when max |T\_model - T| over the KEPT points is below 1 K. C283's 5 \% relative clause judges composition; this fit has no composition residual, so only the temperature clause applies. A pair that splits ONE liquid inside its validity span is not promoted (arm f). Two starts: (0, 0) and (800, -400).

OUTLIERS DISCARDED AFTER THE LEAST-SQUARES STAGE: @@OUTLIERS-BEGIN@@ none: every one of the 40 points lies within 3 K after the least-squares stage. @@OUTLIERS-END@@

@@RESULT-BEGIN@@ WHAT THE RUN FOUND (pinned in `expected`; measured, not argued). fit\_nrtl: b\_ij = 149.3 K, b\_ji = -148.0 K over 40 kept points of 2 laboratories (0 discarded): max |dT| 1.087 K, mean 0.234 K against the acceptance 1 K $\to$ notAccepted. fit\_nrtl\_start2 lands on the same pair (149.3, -148.0). NOT PROMOTABLE: max |dT| 1.087 K over the 40 kept points is not below 1 K @@RESULT-END@@

WHAT IS NOT CLAIMED. The pair is correct at these pressures and in this frame (ideal-gas vapour, the catalogue's own vapour-pressure records); the measurements were transcribed by the archive and have NOT been read back against the articles (the CHECKLIST that `fetch` writes is the aid).

\end{tutnarrative}
\begin{tutgolden}
\goldrow{diag}{fit\_nrtl}{bubble\_aad\_fit}{0.233571520387}
\goldrow{diag}{fit\_nrtl}{bubble\_max\_fit}{1.08673399433}
\goldrow{diag}{fit\_nrtl}{chi2}{48.8451424039}
\goldrow{diag}{fit\_nrtl}{chi2\_in\_sample}{48.8451424039}
\goldrow{diag}{fit\_nrtl}{chi2\_reduced}{1.28539848431}
\goldrow{diag}{fit\_nrtl}{cond\_JtJ}{11667254.6959}
\goldrow{diag}{fit\_nrtl}{converged}{1}
\goldrow{diag}{fit\_nrtl}{corr\_computed}{1}
\goldrow{diag}{fit\_nrtl}{dof}{38}
\goldrow{diag}{fit\_nrtl}{identifiable}{0}
\goldrow{diag}{fit\_nrtl}{invertible}{1}
\goldrow{diag}{fit\_nrtl}{iter}{12}
\goldrow{diag}{fit\_nrtl}{max\_abs\_corr}{0.999999957145}
\goldrow{diag}{fit\_nrtl}{max\_abs\_resid}{1.08673399433}
\goldmore{71}
\end{tutgolden}

\clearpage
\tutentry{curate134\_nrtl\_isopropanol\_224Trimethylpentane}{tutorials/props/curation/curate134\_nrtl\_isopropanol\_224Trimethylpentane}
\begin{tutscheme}
\begin{tikzpicture}[>={Stealth[length=2mm]}, every node/.style={font=\scriptsize}]
  \node[draw, rounded corners=2pt, align=center, minimum width=3.1cm, minimum height=0.8cm, fill=black!4] at (0.0,-0.0) {\textbf{fit\_nrtl}\\[-1pt]{\tiny fitParameters}};
  \node[draw, rounded corners=2pt, align=center, minimum width=3.1cm, minimum height=0.8cm, fill=black!4] at (3.6,-0.0) {\textbf{fit\_nrtl\_start2}\\[-1pt]{\tiny fitParameters}};
  \node[draw, rounded corners=2pt, align=center, minimum width=3.1cm, minimum height=0.8cm, fill=black!4] at (7.2,-0.0) {\textbf{txy\_ideal}\\[-1pt]{\tiny propertyScan1D}};
  \node[draw, rounded corners=2pt, align=center, minimum width=3.1cm, minimum height=0.8cm, fill=black!4] at (0.0,-1.3) {\textbf{txy\_unifac}\\[-1pt]{\tiny propertyScan1D}};
  \node[draw, rounded corners=2pt, align=center, minimum width=3.1cm, minimum height=0.8cm, fill=black!4] at (3.6,-1.3) {\textbf{txy\_nrtl}\\[-1pt]{\tiny propertyScan1D}};
\end{tikzpicture}
\end{tutscheme}
\tutwhat{NRTL isopropanol-224Trimethylpentane regressed by the curation pipeline to every laboratory's measured bubble temperatures (ThermoML locates, the articles source), accepted at max |dT| < 1 K}
\begin{tutnarrative}
AN NRTL PAIR FROM EVERY LABORATORY'S MEASUREMENTS.

WRITTEN BY bin/choupo-curate pair isopropanol 224Trimethylpentane (DEV.md 4c C284, the bulk run; record docs/design/binary-pairs-from-open-measurements.md section 11).

THE DECLARATION -- WRITTEN AND COMMITTED BEFORE ANY FIT RAN (several declarations of the bulk run share one commit).

IDENTITY. isopropanol (CAS 67-63-0, InChIKey KFZMGEQAYNKOFK-UHFFFAOYSA-N) + 224Trimethylpentane (CAS 540-84-1, InChIKey NHTMVDHEPJAVLT-UHFFFAOYSA-N), joined to the archive through identity.json -- never by name. CLASS (C284): both records are volatile molecular liquids (no role nonvolatile/solute, no dissociatesTo, no charge, no solid, a vapour-pressure model), and every measured temperature lies below both critical temperatures.

NO MEASURED SERIES IN THIS TREE. constant/archiveDatasets names each dataset (DOI, NIST URL, article sha256, extraction selector, extract sha256); `bin/choupo-curate fetch` writes the points into constant/archiveExtracts/, which is gitignored.

STUDIES: every bubble-temperature block of the pair at or below 110 kPa with a declared uncertainty (2 block(s)). FITTED, ALL OF THEM (2): 10.1016/j.jct.2011.10.025 block 12 19 pts 101.3 kPa 349.65-361.23 K U 0.14-1.14 K 10.1021/je900327x block 14 19 pts 101.3 kPa 349.47-360.65 K U 0.19-1.09 K SET ASIDE BY THE EXTRACTION RULES, BY NAME: (none)

THE RULE (DEV.md 4c C283 amended, C284). Every laboratory above is fitted; none is held out (the C16 partition and its band are retired as the acceptance test). NRTL, a\_ij = a\_ji = 0, alpha = 0.30 fixed: two parameters (b\_ij, b\_ji); each point weighted by 1/U, U its own stated expanded uncertainty. Objective twoStage (least squares, then minimax; T\_bubble absolute, in K). OUTLIERS: after a least-squares stage over every point, a point with |T\_model - T| > 3 K is discarded BY NAME (`discard` on its evidence entry, listed below) and the fit repeated without it. ACCEPTED when max |T\_model - T| over the KEPT points is below 1 K. C283's 5 \% relative clause judges composition; this fit has no composition residual, so only the temperature clause applies. A pair that splits ONE liquid inside its validity span is not promoted (arm f). Two starts: (0, 0) and (800, -400).

OUTLIERS DISCARDED AFTER THE LEAST-SQUARES STAGE: @@OUTLIERS-BEGIN@@ none: every one of the 38 points lies within 3 K after the least-squares stage. @@OUTLIERS-END@@

@@RESULT-BEGIN@@ WHAT THE RUN FOUND (pinned in `expected`; measured, not argued). fit\_nrtl: b\_ij = 346.3 K, b\_ji = 347.0 K over 38 kept points of 2 laboratories (0 discarded): max |dT| 0.507 K, mean 0.174 K against the acceptance 1 K $\to$ accepted. fit\_nrtl\_start2 lands on the same pair (346.3, 347.0). PROMOTED on 2026-11-10 into data/standards/parameters/NRTL/224Trimethylpentane-isopropanol.dat as reviewStatus autoPromoted, humanReviewed false, because: max |dT| 0.507 K < 1 K over the 38 kept points, one liquid across the validity span @@RESULT-END@@

WHAT IS NOT CLAIMED. The pair is correct at these pressures and in this frame (ideal-gas vapour, the catalogue's own vapour-pressure records); the measurements were transcribed by the archive and have NOT been read back against the articles (the CHECKLIST that `fetch` writes is the aid).

\end{tutnarrative}
\begin{tutgolden}
\goldrow{diag}{fit\_nrtl}{bubble\_aad\_fit}{0.174265020545}
\goldrow{diag}{fit\_nrtl}{bubble\_max\_fit}{0.506740299916}
\goldrow{diag}{fit\_nrtl}{chi2}{22.1415438127}
\goldrow{diag}{fit\_nrtl}{chi2\_in\_sample}{22.1415438127}
\goldrow{diag}{fit\_nrtl}{chi2\_reduced}{0.615042883686}
\goldrow{diag}{fit\_nrtl}{cond\_JtJ}{27.6776653195}
\goldrow{diag}{fit\_nrtl}{converged}{1}
\goldrow{diag}{fit\_nrtl}{corr\_computed}{1}
\goldrow{diag}{fit\_nrtl}{dof}{36}
\goldrow{diag}{fit\_nrtl}{identifiable}{1}
\goldrow{diag}{fit\_nrtl}{invertible}{1}
\goldrow{diag}{fit\_nrtl}{iter}{4}
\goldrow{diag}{fit\_nrtl}{max\_abs\_corr}{0.979055444365}
\goldrow{diag}{fit\_nrtl}{max\_abs\_resid}{0.506740299916}
\goldmore{71}
\end{tutgolden}

\clearpage
\tutentry{curate135\_nrtl\_2Butanol\_nNDimethylformamide}{tutorials/props/curation/curate135\_nrtl\_2Butanol\_nNDimethylformamide}
\begin{tutscheme}
\begin{tikzpicture}[>={Stealth[length=2mm]}, every node/.style={font=\scriptsize}]
  \node[draw, rounded corners=2pt, align=center, minimum width=3.1cm, minimum height=0.8cm, fill=black!4] at (0.0,-0.0) {\textbf{fit\_nrtl}\\[-1pt]{\tiny fitParameters}};
  \node[draw, rounded corners=2pt, align=center, minimum width=3.1cm, minimum height=0.8cm, fill=black!4] at (3.6,-0.0) {\textbf{fit\_nrtl\_start2}\\[-1pt]{\tiny fitParameters}};
  \node[draw, rounded corners=2pt, align=center, minimum width=3.1cm, minimum height=0.8cm, fill=black!4] at (7.2,-0.0) {\textbf{txy\_ideal}\\[-1pt]{\tiny propertyScan1D}};
  \node[draw, rounded corners=2pt, align=center, minimum width=3.1cm, minimum height=0.8cm, fill=black!4] at (0.0,-1.3) {\textbf{txy\_nrtl}\\[-1pt]{\tiny propertyScan1D}};
\end{tikzpicture}
\end{tutscheme}
\tutwhat{NRTL 2Butanol-nNDimethylformamide regressed by the curation pipeline to every laboratory's measured bubble temperatures (ThermoML locates, the articles source), accepted at max |dT| < 1 K}
\begin{tutnarrative}
AN NRTL PAIR FROM EVERY LABORATORY'S MEASUREMENTS.

WRITTEN BY bin/choupo-curate pair 2Butanol nNDimethylformamide (DEV.md 4c C284, the bulk run; record docs/design/binary-pairs-from-open-measurements.md section 11).

THE DECLARATION -- WRITTEN AND COMMITTED BEFORE ANY FIT RAN (several declarations of the bulk run share one commit).

IDENTITY. 2Butanol (CAS 78-92-2, InChIKey BTANRVKWQNVYAZ-UHFFFAOYSA-N) + nNDimethylformamide (CAS 68-12-2, InChIKey ZMXDDKWLCZADIW-UHFFFAOYSA-N), joined to the archive through identity.json -- never by name. CLASS (C284): both records are volatile molecular liquids (no role nonvolatile/solute, no dissociatesTo, no charge, no solid, a vapour-pressure model), and every measured temperature lies below both critical temperatures.

NO MEASURED SERIES IN THIS TREE. constant/archiveDatasets names each dataset (DOI, NIST URL, article sha256, extraction selector, extract sha256); `bin/choupo-curate fetch` writes the points into constant/archiveExtracts/, which is gitignored.

STUDIES: every bubble-temperature block of the pair at or below 110 kPa with a declared uncertainty (2 block(s)). FITTED, ALL OF THEM (2): 10.1016/j.fluid.2014.09.021 block 0 21 pts 101.3 kPa 374.10-424.08 K U 0.26-0.72 K 10.1016/j.fluid.2006.08.018 block 33 19 pts 95.5 kPa 372.35-415.65 K U 0.69-0.69 K SET ASIDE BY THE EXTRACTION RULES, BY NAME: (none)

THE RULE (DEV.md 4c C283 amended, C284). Every laboratory above is fitted; none is held out (the C16 partition and its band are retired as the acceptance test). NRTL, a\_ij = a\_ji = 0, alpha = 0.30 fixed: two parameters (b\_ij, b\_ji); each point weighted by 1/U, U its own stated expanded uncertainty. Objective twoStage (least squares, then minimax; T\_bubble absolute, in K). OUTLIERS: after a least-squares stage over every point, a point with |T\_model - T| > 3 K is discarded BY NAME (`discard` on its evidence entry, listed below) and the fit repeated without it. ACCEPTED when max |T\_model - T| over the KEPT points is below 1 K. C283's 5 \% relative clause judges composition; this fit has no composition residual, so only the temperature clause applies. A pair that splits ONE liquid inside its validity span is not promoted (arm f). Two starts: (0, 0) and (800, -400).

OUTLIERS DISCARDED AFTER THE LEAST-SQUARES STAGE: @@OUTLIERS-BEGIN@@ none: every one of the 40 points lies within 3 K after the least-squares stage. @@OUTLIERS-END@@

@@RESULT-BEGIN@@ WHAT THE RUN FOUND (pinned in `expected`; measured, not argued). fit\_nrtl: b\_ij = 348.6 K, b\_ji = -368.3 K over 40 kept points of 2 laboratories (0 discarded): max |dT| 2.566 K, mean 1.291 K against the acceptance 1 K $\to$ notAccepted. fit\_nrtl\_start2 lands on the same pair (348.6, -368.3). NOT PROMOTABLE: max |dT| 2.566 K over the 40 kept points is not below 1 K @@RESULT-END@@

WHAT IS NOT CLAIMED. The pair is correct at these pressures and in this frame (ideal-gas vapour, the catalogue's own vapour-pressure records); the measurements were transcribed by the archive and have NOT been read back against the articles (the CHECKLIST that `fetch` writes is the aid).

\end{tutnarrative}
\begin{tutgolden}
\goldrow{diag}{fit\_nrtl}{bubble\_aad\_fit}{1.29060108776}
\goldrow{diag}{fit\_nrtl}{bubble\_max\_fit}{2.565931175}
\goldrow{diag}{fit\_nrtl}{chi2}{255.980936084}
\goldrow{diag}{fit\_nrtl}{chi2\_in\_sample}{255.980936084}
\goldrow{diag}{fit\_nrtl}{chi2\_reduced}{6.73634042325}
\goldrow{diag}{fit\_nrtl}{cond\_JtJ}{165348.054238}
\goldrow{diag}{fit\_nrtl}{converged}{1}
\goldrow{diag}{fit\_nrtl}{corr\_computed}{1}
\goldrow{diag}{fit\_nrtl}{dof}{38}
\goldrow{diag}{fit\_nrtl}{identifiable}{0}
\goldrow{diag}{fit\_nrtl}{invertible}{1}
\goldrow{diag}{fit\_nrtl}{iter}{12}
\goldrow{diag}{fit\_nrtl}{max\_abs\_corr}{0.999996976089}
\goldrow{diag}{fit\_nrtl}{max\_abs\_resid}{2.565931175}
\goldmore{56}
\end{tutgolden}

\clearpage
\tutentry{curate136\_nrtl\_ethylBenzene\_sulfolane}{tutorials/props/curation/curate136\_nrtl\_ethylBenzene\_sulfolane}
\begin{tutscheme}
\begin{tikzpicture}[>={Stealth[length=2mm]}, every node/.style={font=\scriptsize}]
  \node[draw, rounded corners=2pt, align=center, minimum width=3.1cm, minimum height=0.8cm, fill=black!4] at (0.0,-0.0) {\textbf{fit\_nrtl}\\[-1pt]{\tiny fitParameters}};
  \node[draw, rounded corners=2pt, align=center, minimum width=3.1cm, minimum height=0.8cm, fill=black!4] at (3.6,-0.0) {\textbf{fit\_nrtl\_start2}\\[-1pt]{\tiny fitParameters}};
  \node[draw, rounded corners=2pt, align=center, minimum width=3.1cm, minimum height=0.8cm, fill=black!4] at (7.2,-0.0) {\textbf{txy\_ideal}\\[-1pt]{\tiny propertyScan1D}};
  \node[draw, rounded corners=2pt, align=center, minimum width=3.1cm, minimum height=0.8cm, fill=black!4] at (0.0,-1.3) {\textbf{txy\_nrtl}\\[-1pt]{\tiny propertyScan1D}};
\end{tikzpicture}
\end{tutscheme}
\tutwhat{NRTL ethylBenzene-sulfolane regressed by the curation pipeline to every laboratory's measured bubble temperatures (ThermoML locates, the articles source), accepted at max |dT| < 1 K}
\begin{tutnarrative}
AN NRTL PAIR FROM EVERY LABORATORY'S MEASUREMENTS.

WRITTEN BY bin/choupo-curate pair ethylBenzene sulfolane (DEV.md 4c C284, the bulk run; record docs/design/binary-pairs-from-open-measurements.md section 11).

THE DECLARATION -- WRITTEN AND COMMITTED BEFORE ANY FIT RAN (several declarations of the bulk run share one commit).

IDENTITY. ethylBenzene (CAS 100-41-4, InChIKey YNQLUTRBYVCPMQ-UHFFFAOYSA-N) + sulfolane (CAS 126-33-0, InChIKey HXJUTPCZVOIRIF-UHFFFAOYSA-N), joined to the archive through identity.json -- never by name. CLASS (C284): both records are volatile molecular liquids (no role nonvolatile/solute, no dissociatesTo, no charge, no solid, a vapour-pressure model), and every measured temperature lies below both critical temperatures.

NO MEASURED SERIES IN THIS TREE. constant/archiveDatasets names each dataset (DOI, NIST URL, article sha256, extraction selector, extract sha256); `bin/choupo-curate fetch` writes the points into constant/archiveExtracts/, which is gitignored.

STUDIES: every bubble-temperature block of the pair at or below 110 kPa with a declared uncertainty (2 block(s)). FITTED, ALL OF THEM (2): 10.1021/je200098d block 4 30 pts 3 isobars 5-20 kPa 325.00-379.40 K U 0.3-0.5 K 10.1016/j.jct.2011.06.015 block 6 12 pts 101.33 kPa 411.69-492.15 K U 0.12-8 K SET ASIDE BY THE EXTRACTION RULES, BY NAME: (none)

THE RULE (DEV.md 4c C283 amended, C284). Every laboratory above is fitted; none is held out (the C16 partition and its band are retired as the acceptance test). NRTL, a\_ij = a\_ji = 0, alpha = 0.30 fixed: two parameters (b\_ij, b\_ji); each point weighted by 1/U, U its own stated expanded uncertainty. Objective twoStage (least squares, then minimax; T\_bubble absolute, in K). OUTLIERS: after a least-squares stage over every point, a point with |T\_model - T| > 3 K is discarded BY NAME (`discard` on its evidence entry, listed below) and the fit repeated without it. ACCEPTED when max |T\_model - T| over the KEPT points is below 1 K. C283's 5 \% relative clause judges composition; this fit has no composition residual, so only the temperature clause applies. A pair that splits ONE liquid inside its validity span is not promoted (arm f). Two starts: (0, 0) and (800, -400).

OUTLIERS DISCARDED AFTER THE LEAST-SQUARES STAGE: @@OUTLIERS-BEGIN@@ 14 point(s) of 42: doi 10.1021/je200098d row 3: x 0.0920, P 20.000 kPa, T 379.40 K, dT +3.37 K; doi 10.1021/je200098d row 4: x 0.1970, P 20.000 kPa, T 361.40 K, dT +6.67 K; doi 10.1021/je200098d row 7: x 0.2570, P 10.000 kPa, T 342.20 K, dT +3.13 K; doi 10.1021/je200098d row 8: x 0.2570, P 20.000 kPa, T 360.30 K, dT +4.62 K; doi 10.1021/je200098d row 10: x 0.3470, P 20.000 kPa, T 358.40 K, dT +4.27 K; doi 10.1021/je200098d row 14: x 0.4300, P 20.000 kPa, T 358.40 K, dT +3.45 K; doi 10.1021/je200098d row 17: x 0.5150, P 20.000 kPa, T 358.30 K, dT +3.29 K; doi 10.1021/je200098d row 20: x 0.6500, P 20.000 kPa, T 358.40 K, dT +3.11 K; doi 10.1021/je200098d row 24: x 0.7410, P 20.000 kPa, T 358.20 K, dT +3.12 K; doi 10.1016/j.jct.2011.06.015 row 2: x 0.0613, P 101.330 kPa, T 473.00 K, dT -7.09 K; doi 10.1016/j.jct.2011.06.015 row 4: x 0.0949, P 101.330 kPa, T 456.60 K, dT -7.28 K; doi 10.1016/j.jct.2011.06.015 row 5: x 0.1462, P 101.330 kPa, T 446.35 K, dT -10.64 K; doi 10.1016/j.jct.2011.06.015 row 6: x 0.2184, P 101.330 kPa, T 432.07 K, dT -5.89 K; doi 10.1016/j.jct.2011.06.015 row 7: x 0.3436, P 101.330 kPa, T 425.50 K, dT -6.16 K @@OUTLIERS-END@@

@@RESULT-BEGIN@@ WHAT THE RUN FOUND (pinned in `expected`; measured, not argued). fit\_nrtl: b\_ij = 583.7 K, b\_ji = 312.7 K over 28 kept points of 2 laboratories (14 discarded): max |dT| 3.584 K, mean 1.527 K against the acceptance 1 K $\to$ notAccepted. fit\_nrtl\_start2 lands on the same pair (583.7, 312.7). NOT PROMOTABLE: max |dT| 3.584 K over the 28 kept points is not below 1 K @@RESULT-END@@

WHAT IS NOT CLAIMED. The pair is correct at these pressures and in this frame (ideal-gas vapour, the catalogue's own vapour-pressure records); the measurements were transcribed by the archive and have NOT been read back against the articles (the CHECKLIST that `fetch` writes is the aid).

\end{tutnarrative}
\begin{tutgolden}
\goldrow{diag}{fit\_nrtl}{bubble\_aad\_fit}{1.52713877479}
\goldrow{diag}{fit\_nrtl}{bubble\_max\_fit}{3.58390605665}
\goldrow{diag}{fit\_nrtl}{chi2}{1065.23860605}
\goldrow{diag}{fit\_nrtl}{chi2\_in\_sample}{1065.23860605}
\goldrow{diag}{fit\_nrtl}{chi2\_reduced}{40.9707156174}
\goldrow{diag}{fit\_nrtl}{cond\_JtJ}{2.93355845401}
\goldrow{diag}{fit\_nrtl}{converged}{1}
\goldrow{diag}{fit\_nrtl}{corr\_computed}{1}
\goldrow{diag}{fit\_nrtl}{dof}{26}
\goldrow{diag}{fit\_nrtl}{identifiable}{1}
\goldrow{diag}{fit\_nrtl}{invertible}{1}
\goldrow{diag}{fit\_nrtl}{iter}{5}
\goldrow{diag}{fit\_nrtl}{max\_abs\_corr}{0.842497299853}
\goldrow{diag}{fit\_nrtl}{max\_abs\_resid}{3.58390605665}
\goldmore{50}
\end{tutgolden}

\clearpage
\tutentry{curate137\_nrtl\_2Methyl1Propanol\_nNDimethylformamide}{tutorials/props/curation/curate137\_nrtl\_2Methyl1Propanol\_nNDimethylformamide}
\begin{tutscheme}
\begin{tikzpicture}[>={Stealth[length=2mm]}, every node/.style={font=\scriptsize}]
  \node[draw, rounded corners=2pt, align=center, minimum width=3.1cm, minimum height=0.8cm, fill=black!4] at (0.0,-0.0) {\textbf{fit\_nrtl}\\[-1pt]{\tiny fitParameters}};
  \node[draw, rounded corners=2pt, align=center, minimum width=3.1cm, minimum height=0.8cm, fill=black!4] at (3.6,-0.0) {\textbf{fit\_nrtl\_start2}\\[-1pt]{\tiny fitParameters}};
  \node[draw, rounded corners=2pt, align=center, minimum width=3.1cm, minimum height=0.8cm, fill=black!4] at (7.2,-0.0) {\textbf{txy\_ideal}\\[-1pt]{\tiny propertyScan1D}};
  \node[draw, rounded corners=2pt, align=center, minimum width=3.1cm, minimum height=0.8cm, fill=black!4] at (0.0,-1.3) {\textbf{txy\_nrtl}\\[-1pt]{\tiny propertyScan1D}};
\end{tikzpicture}
\end{tutscheme}
\tutwhat{NRTL 2Methyl1Propanol-nNDimethylformamide regressed by the curation pipeline to every laboratory's measured bubble temperatures (ThermoML locates, the articles source), accepted at max |dT| < 1 K}
\begin{tutnarrative}
AN NRTL PAIR FROM EVERY LABORATORY'S MEASUREMENTS.

WRITTEN BY bin/choupo-curate pair 2Methyl1Propanol nNDimethylformamide (DEV.md 4c C284, the bulk run; record docs/design/binary-pairs-from-open-measurements.md section 11).

THE DECLARATION -- WRITTEN AND COMMITTED BEFORE ANY FIT RAN (several declarations of the bulk run share one commit).

IDENTITY. 2Methyl1Propanol (CAS 78-83-1, InChIKey ZXEKIIBDNHEJCQ-UHFFFAOYSA-N) + nNDimethylformamide (CAS 68-12-2, InChIKey ZMXDDKWLCZADIW-UHFFFAOYSA-N), joined to the archive through identity.json -- never by name. CLASS (C284): both records are volatile molecular liquids (no role nonvolatile/solute, no dissociatesTo, no charge, no solid, a vapour-pressure model), and every measured temperature lies below both critical temperatures.

NO MEASURED SERIES IN THIS TREE. constant/archiveDatasets names each dataset (DOI, NIST URL, article sha256, extraction selector, extract sha256); `bin/choupo-curate fetch` writes the points into constant/archiveExtracts/, which is gitignored.

STUDIES: every bubble-temperature block of the pair at or below 110 kPa with a declared uncertainty (2 block(s)). FITTED, ALL OF THEM (2): 10.1016/j.fluid.2006.08.018 block 32 19 pts 95.5 kPa 380.75-417.45 K U 1.03-1.03 K 10.1016/j.fluid.2005.03.014 block 7 18 pts 101.3 kPa 382.06-421.95 K U 0.49-0.51 K SET ASIDE BY THE EXTRACTION RULES, BY NAME: (none)

THE RULE (DEV.md 4c C283 amended, C284). Every laboratory above is fitted; none is held out (the C16 partition and its band are retired as the acceptance test). NRTL, a\_ij = a\_ji = 0, alpha = 0.30 fixed: two parameters (b\_ij, b\_ji); each point weighted by 1/U, U its own stated expanded uncertainty. Objective twoStage (least squares, then minimax; T\_bubble absolute, in K). OUTLIERS: after a least-squares stage over every point, a point with |T\_model - T| > 3 K is discarded BY NAME (`discard` on its evidence entry, listed below) and the fit repeated without it. ACCEPTED when max |T\_model - T| over the KEPT points is below 1 K. C283's 5 \% relative clause judges composition; this fit has no composition residual, so only the temperature clause applies. A pair that splits ONE liquid inside its validity span is not promoted (arm f). Two starts: (0, 0) and (800, -400).

OUTLIERS DISCARDED AFTER THE LEAST-SQUARES STAGE: @@OUTLIERS-BEGIN@@ none: every one of the 37 points lies within 3 K after the least-squares stage. @@OUTLIERS-END@@

@@RESULT-BEGIN@@ WHAT THE RUN FOUND (pinned in `expected`; measured, not argued). fit\_nrtl: b\_ij = 311.5 K, b\_ji = -366.5 K over 37 kept points of 2 laboratories (0 discarded): max |dT| 2.043 K, mean 0.879 K against the acceptance 1 K $\to$ notAccepted. fit\_nrtl\_start2 lands on A DIFFERENT pair (336.3, -381.4). NOT PROMOTABLE: max |dT| 2.043 K over the 37 kept points is not below 1 K @@RESULT-END@@

WHAT IS NOT CLAIMED. The pair is correct at these pressures and in this frame (ideal-gas vapour, the catalogue's own vapour-pressure records); the measurements were transcribed by the archive and have NOT been read back against the articles (the CHECKLIST that `fetch` writes is the aid).

\end{tutnarrative}
\begin{tutgolden}
\goldrow{diag}{fit\_nrtl}{bubble\_aad\_fit}{0.879196945619}
\goldrow{diag}{fit\_nrtl}{bubble\_max\_fit}{2.04296773041}
\goldrow{diag}{fit\_nrtl}{chi2}{59.8965704687}
\goldrow{diag}{fit\_nrtl}{chi2\_in\_sample}{59.8965704687}
\goldrow{diag}{fit\_nrtl}{chi2\_reduced}{1.71133058482}
\goldrow{diag}{fit\_nrtl}{cond\_JtJ}{4737.64485331}
\goldrow{diag}{fit\_nrtl}{converged}{1}
\goldrow{diag}{fit\_nrtl}{corr\_computed}{1}
\goldrow{diag}{fit\_nrtl}{dof}{35}
\goldrow{diag}{fit\_nrtl}{identifiable}{0}
\goldrow{diag}{fit\_nrtl}{invertible}{1}
\goldrow{diag}{fit\_nrtl}{iter}{8}
\goldrow{diag}{fit\_nrtl}{max\_abs\_corr}{0.999894479034}
\goldrow{diag}{fit\_nrtl}{max\_abs\_resid}{2.04296773041}
\goldmore{54}
\end{tutgolden}

\clearpage
\tutentry{curate138\_nrtl\_ethanol\_cyclohexane}{tutorials/props/curation/curate138\_nrtl\_ethanol\_cyclohexane}
\begin{tutscheme}
\begin{tikzpicture}[>={Stealth[length=2mm]}, every node/.style={font=\scriptsize}]
  \node[draw, rounded corners=2pt, align=center, minimum width=3.1cm, minimum height=0.8cm, fill=black!4] at (0.0,-0.0) {\textbf{fit\_nrtl}\\[-1pt]{\tiny fitParameters}};
  \node[draw, rounded corners=2pt, align=center, minimum width=3.1cm, minimum height=0.8cm, fill=black!4] at (3.6,-0.0) {\textbf{fit\_nrtl\_start2}\\[-1pt]{\tiny fitParameters}};
  \node[draw, rounded corners=2pt, align=center, minimum width=3.1cm, minimum height=0.8cm, fill=black!4] at (7.2,-0.0) {\textbf{txy\_ideal}\\[-1pt]{\tiny propertyScan1D}};
  \node[draw, rounded corners=2pt, align=center, minimum width=3.1cm, minimum height=0.8cm, fill=black!4] at (0.0,-1.3) {\textbf{txy\_unifac}\\[-1pt]{\tiny propertyScan1D}};
  \node[draw, rounded corners=2pt, align=center, minimum width=3.1cm, minimum height=0.8cm, fill=black!4] at (3.6,-1.3) {\textbf{txy\_nrtl}\\[-1pt]{\tiny propertyScan1D}};
\end{tikzpicture}
\end{tutscheme}
\tutwhat{NRTL ethanol-cyclohexane regressed by the curation pipeline to every laboratory's measured bubble temperatures (ThermoML locates, the articles source), accepted at max |dT| < 1 K}
\begin{tutnarrative}
AN NRTL PAIR FROM EVERY LABORATORY'S MEASUREMENTS.

WRITTEN BY bin/choupo-curate pair ethanol cyclohexane (DEV.md 4c C284, the bulk run; record docs/design/binary-pairs-from-open-measurements.md section 11).

THE DECLARATION -- WRITTEN AND COMMITTED BEFORE ANY FIT RAN (several declarations of the bulk run share one commit).

IDENTITY. ethanol (CAS 64-17-5, InChIKey LFQSCWFLJHTTHZ-UHFFFAOYSA-N) + cyclohexane (CAS 110-82-7, InChIKey XDTMQSROBMDMFD-UHFFFAOYSA-N), joined to the archive through identity.json -- never by name. CLASS (C284): both records are volatile molecular liquids (no role nonvolatile/solute, no dissociatesTo, no charge, no solid, a vapour-pressure model), and every measured temperature lies below both critical temperatures.

NO MEASURED SERIES IN THIS TREE. constant/archiveDatasets names each dataset (DOI, NIST URL, article sha256, extraction selector, extract sha256); `bin/choupo-curate fetch` writes the points into constant/archiveExtracts/, which is gitignored.

STUDIES: every bubble-temperature block of the pair at or below 110 kPa with a declared uncertainty (2 block(s)). FITTED, ALL OF THEM (2): 10.1016/j.fluid.2018.07.013 block 0 25 pts 101.3 kPa 337.65-353.45 K U 0.22-1.71 K 10.1016/j.fluid.2006.01.023 block 0 14 pts 101.325 kPa 338.08-350.50 K U 1.36-1.4 K SET ASIDE BY THE EXTRACTION RULES, BY NAME: (none)

THE RULE (DEV.md 4c C283 amended, C284). Every laboratory above is fitted; none is held out (the C16 partition and its band are retired as the acceptance test). NRTL, a\_ij = a\_ji = 0, alpha = 0.30 fixed: two parameters (b\_ij, b\_ji); each point weighted by 1/U, U its own stated expanded uncertainty. Objective twoStage (least squares, then minimax; T\_bubble absolute, in K). OUTLIERS: after a least-squares stage over every point, a point with |T\_model - T| > 3 K is discarded BY NAME (`discard` on its evidence entry, listed below) and the fit repeated without it. ACCEPTED when max |T\_model - T| over the KEPT points is below 1 K. C283's 5 \% relative clause judges composition; this fit has no composition residual, so only the temperature clause applies. A pair that splits ONE liquid inside its validity span is not promoted (arm f). Two starts: (0, 0) and (800, -400).

OUTLIERS DISCARDED AFTER THE LEAST-SQUARES STAGE: @@OUTLIERS-BEGIN@@ none: every one of the 39 points lies within 3 K after the least-squares stage. @@OUTLIERS-END@@

@@RESULT-BEGIN@@ WHAT THE RUN FOUND (pinned in `expected`; measured, not argued). fit\_nrtl: b\_ij = 437.4 K, b\_ji = 485.2 K over 39 kept points of 2 laboratories (0 discarded): max |dT| 2.682 K, mean 0.903 K against the acceptance 1 K $\to$ notAccepted. fit\_nrtl\_start2 lands on the same pair (437.4, 485.2). NOT PROMOTABLE: max |dT| 2.682 K over the 39 kept points is not below 1 K @@RESULT-END@@

WHAT IS NOT CLAIMED. The pair is correct at these pressures and in this frame (ideal-gas vapour, the catalogue's own vapour-pressure records); the measurements were transcribed by the archive and have NOT been read back against the articles (the CHECKLIST that `fetch` writes is the aid).

\end{tutnarrative}
\begin{tutgolden}
\goldrow{diag}{fit\_nrtl}{bubble\_aad\_fit}{0.902767851878}
\goldrow{diag}{fit\_nrtl}{bubble\_max\_fit}{2.68151186137}
\goldrow{diag}{fit\_nrtl}{chi2}{186.221155823}
\goldrow{diag}{fit\_nrtl}{chi2\_in\_sample}{186.221155823}
\goldrow{diag}{fit\_nrtl}{chi2\_reduced}{5.03300421142}
\goldrow{diag}{fit\_nrtl}{cond\_JtJ}{11.2810412069}
\goldrow{diag}{fit\_nrtl}{converged}{1}
\goldrow{diag}{fit\_nrtl}{corr\_computed}{1}
\goldrow{diag}{fit\_nrtl}{dof}{37}
\goldrow{diag}{fit\_nrtl}{identifiable}{1}
\goldrow{diag}{fit\_nrtl}{invertible}{1}
\goldrow{diag}{fit\_nrtl}{iter}{5}
\goldrow{diag}{fit\_nrtl}{max\_abs\_corr}{0.940785157718}
\goldrow{diag}{fit\_nrtl}{max\_abs\_resid}{2.68151186137}
\goldmore{71}
\end{tutgolden}

\clearpage
\tutentry{curate139\_nrtl\_methanol\_methylEthylKetone}{tutorials/props/curation/curate139\_nrtl\_methanol\_methylEthylKetone}
\begin{tutscheme}
\begin{tikzpicture}[>={Stealth[length=2mm]}, every node/.style={font=\scriptsize}]
  \node[draw, rounded corners=2pt, align=center, minimum width=3.1cm, minimum height=0.8cm, fill=black!4] at (0.0,-0.0) {\textbf{fit\_nrtl}\\[-1pt]{\tiny fitParameters}};
  \node[draw, rounded corners=2pt, align=center, minimum width=3.1cm, minimum height=0.8cm, fill=black!4] at (3.6,-0.0) {\textbf{fit\_nrtl\_start2}\\[-1pt]{\tiny fitParameters}};
  \node[draw, rounded corners=2pt, align=center, minimum width=3.1cm, minimum height=0.8cm, fill=black!4] at (7.2,-0.0) {\textbf{txy\_ideal}\\[-1pt]{\tiny propertyScan1D}};
  \node[draw, rounded corners=2pt, align=center, minimum width=3.1cm, minimum height=0.8cm, fill=black!4] at (0.0,-1.3) {\textbf{txy\_unifac}\\[-1pt]{\tiny propertyScan1D}};
  \node[draw, rounded corners=2pt, align=center, minimum width=3.1cm, minimum height=0.8cm, fill=black!4] at (3.6,-1.3) {\textbf{txy\_nrtl}\\[-1pt]{\tiny propertyScan1D}};
\end{tikzpicture}
\end{tutscheme}
\tutwhat{NRTL methanol-methylEthylKetone regressed by the curation pipeline to every laboratory's measured bubble temperatures (ThermoML locates, the articles source), accepted at max |dT| < 1 K}
\begin{tutnarrative}
AN NRTL PAIR FROM EVERY LABORATORY'S MEASUREMENTS.

WRITTEN BY bin/choupo-curate pair methanol methylEthylKetone (DEV.md 4c C284, the bulk run; record docs/design/binary-pairs-from-open-measurements.md section 11).

THE DECLARATION -- WRITTEN AND COMMITTED BEFORE ANY FIT RAN (several declarations of the bulk run share one commit).

IDENTITY. methanol (CAS 67-56-1, InChIKey OKKJLVBELUTLKV-UHFFFAOYSA-N) + methylEthylKetone (CAS 78-93-3, InChIKey ZWEHNKRNPOVVGH-UHFFFAOYSA-N), joined to the archive through identity.json -- never by name. CLASS (C284): both records are volatile molecular liquids (no role nonvolatile/solute, no dissociatesTo, no charge, no solid, a vapour-pressure model), and every measured temperature lies below both critical temperatures.

NO MEASURED SERIES IN THIS TREE. constant/archiveDatasets names each dataset (DOI, NIST URL, article sha256, extraction selector, extract sha256); `bin/choupo-curate fetch` writes the points into constant/archiveExtracts/, which is gitignored.

STUDIES: every bubble-temperature block of the pair at or below 110 kPa with a declared uncertainty (2 block(s)). FITTED, ALL OF THEM (2): 10.1016/j.fluid.2016.06.048 block 0 19 pts 101.3 kPa 337.40-349.80 K U 0.2-0.3 K 10.1021/je301062g block 0 17 pts 101.32 kPa 337.40-348.20 K U 0.2-0.3 K SET ASIDE BY THE EXTRACTION RULES, BY NAME: (none)

THE RULE (DEV.md 4c C283 amended, C284). Every laboratory above is fitted; none is held out (the C16 partition and its band are retired as the acceptance test). NRTL, a\_ij = a\_ji = 0, alpha = 0.30 fixed: two parameters (b\_ij, b\_ji); each point weighted by 1/U, U its own stated expanded uncertainty. Objective twoStage (least squares, then minimax; T\_bubble absolute, in K). OUTLIERS: after a least-squares stage over every point, a point with |T\_model - T| > 3 K is discarded BY NAME (`discard` on its evidence entry, listed below) and the fit repeated without it. ACCEPTED when max |T\_model - T| over the KEPT points is below 1 K. C283's 5 \% relative clause judges composition; this fit has no composition residual, so only the temperature clause applies. A pair that splits ONE liquid inside its validity span is not promoted (arm f). Two starts: (0, 0) and (800, -400).

OUTLIERS DISCARDED AFTER THE LEAST-SQUARES STAGE: @@OUTLIERS-BEGIN@@ none: every one of the 36 points lies within 3 K after the least-squares stage. @@OUTLIERS-END@@

@@RESULT-BEGIN@@ WHAT THE RUN FOUND (pinned in `expected`; measured, not argued). fit\_nrtl: b\_ij = 243.0 K, b\_ji = 22.8 K over 36 kept points of 2 laboratories (0 discarded): max |dT| 0.659 K, mean 0.297 K against the acceptance 1 K $\to$ accepted. fit\_nrtl\_start2 lands on the same pair (243.0, 22.8). PROMOTED on 2026-11-10 into data/standards/parameters/NRTL/methanol-methylEthylKetone.dat as reviewStatus autoPromoted, humanReviewed false, because: max |dT| 0.659 K < 1 K over the 36 kept points, one liquid across the validity span @@RESULT-END@@

WHAT IS NOT CLAIMED. The pair is correct at these pressures and in this frame (ideal-gas vapour, the catalogue's own vapour-pressure records); the measurements were transcribed by the archive and have NOT been read back against the articles (the CHECKLIST that `fetch` writes is the aid).

\end{tutnarrative}
\begin{tutgolden}
\goldrow{diag}{fit\_nrtl}{bubble\_aad\_fit}{0.297190281837}
\goldrow{diag}{fit\_nrtl}{bubble\_max\_fit}{0.659208674087}
\goldrow{diag}{fit\_nrtl}{chi2}{76.9526853042}
\goldrow{diag}{fit\_nrtl}{chi2\_in\_sample}{76.9526853042}
\goldrow{diag}{fit\_nrtl}{chi2\_reduced}{2.26331427365}
\goldrow{diag}{fit\_nrtl}{cond\_JtJ}{127.085706346}
\goldrow{diag}{fit\_nrtl}{converged}{1}
\goldrow{diag}{fit\_nrtl}{corr\_computed}{1}
\goldrow{diag}{fit\_nrtl}{dof}{34}
\goldrow{diag}{fit\_nrtl}{identifiable}{1}
\goldrow{diag}{fit\_nrtl}{invertible}{1}
\goldrow{diag}{fit\_nrtl}{iter}{4}
\goldrow{diag}{fit\_nrtl}{max\_abs\_corr}{0.996088714724}
\goldrow{diag}{fit\_nrtl}{max\_abs\_resid}{0.659208674087}
\goldmore{68}
\end{tutgolden}

\clearpage
\tutentry{curate13\_synthetic\_lle\_butanol\_water}{tutorials/props/curation/curate13\_synthetic\_lle\_butanol\_water}
\begin{tutscheme}
\begin{tikzpicture}[>={Stealth[length=2mm]}, every node/.style={font=\scriptsize}]
  \node[draw, rounded corners=2pt, align=center, minimum width=3.1cm, minimum height=0.8cm, fill=black!4] at (0.0,-0.0) {\textbf{fit\_mixed}\\[-1pt]{\tiny fitParameters}};
  \node[draw, rounded corners=2pt, align=center, minimum width=3.1cm, minimum height=0.8cm, fill=black!4] at (3.6,-0.0) {\textbf{fit\_mixed\_start2}\\[-1pt]{\tiny fitParameters}};
\end{tikzpicture}
\end{tutscheme}
\tutwhat{SYNTHETIC water + 1-butanol: NRTL recovered from engine-generated bubble temperatures AND tie lines in ONE fit (kind T\_bubble+LLE), held out against other compositions and temperatures}
\begin{tutnarrative}
BUBBLE TEMPERATURES AND TIE LINES IN ONE FIT, on a pair whose parameters are KNOWN (DEV.md 4c C226 A1).

WHAT THIS IS. A SYNTHETIC witness, labelled as such everywhere it can be: the four datasets under constant/experiments/ were GENERATED by running Choupo's own engine (propertyScan1D for T\_bubble at 101.325 kPa, propertyScanBinary for the binodal) on a water + 1-butanol NRTL pair of known parameters (alpha 0.30, a\_ij = a\_ji = 0, b\_ij = 150 K, b\_ji = 1050 K, i = nButanol, j = water). No laboratory measurement is involved, so this case shows that the machinery recovers what it was given and nothing about water + 1-butanol itself. THE REAL PAIR IS STILL TO DO: it needs measured tie lines (a ThermoML LLE series, or a published table read by a curator), and the cache that would locate them lives outside this tree.

WHY ONE FIT. A partially miscible pair has two kinds of evidence about the same two numbers: where its boiling liquid sits (T\_bubble) and where its two liquids sit (the binodal). A pair regressed to the first alone is free to split the liquid wherever it likes or not at all (C16 slice 2: ethanol + cyclohexane, validated and splitting its whole span). The second is the measurement that pins the gap. `kind T\_bubble+LLE` adds one row per measured binodal composition (x\_model - x\_measured)/sigma to the bubble rows (T\_model - T\_measured)/U, so a single chi2 weighs both, with the weights the author declared -- `weighting uncertainty` for the kelvin rows, and `liquidLiquid \{ sigma \}` for the mole-fraction rows.

THE RULE C16 ALREADY HAS, REUSED. Fit on one set of datasets, test on others the fitter is never handed, with the bands declared BEFORE the fit (`maxAAD` for the bubble temperatures, `maxLLE` for the binodal), and a verdict. A held-out tie line the fitted model leaves as ONE liquid is named on stderr and is a miss, never a skip.

THE START. The inline pair of constant/thermoPhysPropDict (b\_ij 0, b\_ji 1000) splits the liquid at every tie line. A start that splits nothing (ideal mixing) is REFUSED by name: the penalty for a missing split is flat, so the search would have no gradient.

WHAT THE RUN SHOWS (pinned in `expected`). Both starts recover b\_ij and b\_ji of the generating pair to round-off, and the held-out bubble temperatures and tie lines are reproduced to round-off: the band-judged verdict is `validated`. That is arithmetic, not evidence about water + 1-butanol -- the data come from the model.

\end{tutnarrative}
\begin{tutgolden}
\goldrow{diag}{fit\_mixed}{cond\_JtJ}{29.8172762939}
\goldrow{diag}{fit\_mixed}{converged}{1}
\goldrow{diag}{fit\_mixed}{corr\_computed}{1}
\goldrow{diag}{fit\_mixed}{dof}{26}
\goldrow{diag}{fit\_mixed}{identifiable}{1}
\goldrow{diag}{fit\_mixed}{invertible}{1}
\goldrow{diag}{fit\_mixed}{iter}{5}
\goldrow{diag}{fit\_mixed}{lle\_n\_ties}{8}
\goldrow{diag}{fit\_mixed}{lle\_n\_ties\_heldout}{5}
\goldrow{diag}{fit\_mixed}{lle\_penalised\_ties}{0}
\goldrow{diag}{fit\_mixed}{lle\_penalised\_ties\_heldout}{0}
\goldrow{diag}{fit\_mixed}{max\_abs\_corr}{0.750661460302}
\goldrow{diag}{fit\_mixed}{n\_data}{28}
\goldrow{diag}{fit\_mixed}{n\_heldout}{10}
\goldmore{32}
\end{tutgolden}

\clearpage
\tutentry{curate140\_nrtl\_nButanol\_pXylene}{tutorials/props/curation/curate140\_nrtl\_nButanol\_pXylene}
\begin{tutscheme}
\begin{tikzpicture}[>={Stealth[length=2mm]}, every node/.style={font=\scriptsize}]
  \node[draw, rounded corners=2pt, align=center, minimum width=3.1cm, minimum height=0.8cm, fill=black!4] at (0.0,-0.0) {\textbf{fit\_nrtl}\\[-1pt]{\tiny fitParameters}};
  \node[draw, rounded corners=2pt, align=center, minimum width=3.1cm, minimum height=0.8cm, fill=black!4] at (3.6,-0.0) {\textbf{fit\_nrtl\_start2}\\[-1pt]{\tiny fitParameters}};
  \node[draw, rounded corners=2pt, align=center, minimum width=3.1cm, minimum height=0.8cm, fill=black!4] at (7.2,-0.0) {\textbf{txy\_ideal}\\[-1pt]{\tiny propertyScan1D}};
  \node[draw, rounded corners=2pt, align=center, minimum width=3.1cm, minimum height=0.8cm, fill=black!4] at (0.0,-1.3) {\textbf{txy\_unifac}\\[-1pt]{\tiny propertyScan1D}};
  \node[draw, rounded corners=2pt, align=center, minimum width=3.1cm, minimum height=0.8cm, fill=black!4] at (3.6,-1.3) {\textbf{txy\_nrtl}\\[-1pt]{\tiny propertyScan1D}};
\end{tikzpicture}
\end{tutscheme}
\tutwhat{NRTL nButanol-pXylene regressed by the curation pipeline to every laboratory's measured bubble temperatures (ThermoML locates, the articles source), accepted at max |dT| < 1 K}
\begin{tutnarrative}
AN NRTL PAIR FROM EVERY LABORATORY'S MEASUREMENTS.

WRITTEN BY bin/choupo-curate pair nButanol pXylene (DEV.md 4c C284, the bulk run; record docs/design/binary-pairs-from-open-measurements.md section 11).

THE DECLARATION -- WRITTEN AND COMMITTED BEFORE ANY FIT RAN (several declarations of the bulk run share one commit).

IDENTITY. nButanol (CAS 71-36-3, InChIKey LRHPLDYGYMQRHN-UHFFFAOYSA-N) + pXylene (CAS 106-42-3, InChIKey URLKBWYHVLBVBO-UHFFFAOYSA-N), joined to the archive through identity.json -- never by name. CLASS (C284): both records are volatile molecular liquids (no role nonvolatile/solute, no dissociatesTo, no charge, no solid, a vapour-pressure model), and every measured temperature lies below both critical temperatures.

NO MEASURED SERIES IN THIS TREE. constant/archiveDatasets names each dataset (DOI, NIST URL, article sha256, extraction selector, extract sha256); `bin/choupo-curate fetch` writes the points into constant/archiveExtracts/, which is gitignored.

STUDIES: every bubble-temperature block of the pair at or below 110 kPa with a declared uncertainty (2 block(s)). FITTED, ALL OF THEM (2): 10.1021/je201053v block 16 19 pts 101.33 kPa 388.94-407.14 K U 0.37-0.97 K 10.1007/s10765-006-0018-5 block 4 15 pts 97.3 kPa 389.00-406.20 K U 0.3-2.6 K SET ASIDE BY THE EXTRACTION RULES, BY NAME: (none)

THE RULE (DEV.md 4c C283 amended, C284). Every laboratory above is fitted; none is held out (the C16 partition and its band are retired as the acceptance test). NRTL, a\_ij = a\_ji = 0, alpha = 0.30 fixed: two parameters (b\_ij, b\_ji); each point weighted by 1/U, U its own stated expanded uncertainty. Objective twoStage (least squares, then minimax; T\_bubble absolute, in K). OUTLIERS: after a least-squares stage over every point, a point with |T\_model - T| > 3 K is discarded BY NAME (`discard` on its evidence entry, listed below) and the fit repeated without it. ACCEPTED when max |T\_model - T| over the KEPT points is below 1 K. C283's 5 \% relative clause judges composition; this fit has no composition residual, so only the temperature clause applies. A pair that splits ONE liquid inside its validity span is not promoted (arm f). Two starts: (0, 0) and (800, -400).

OUTLIERS DISCARDED AFTER THE LEAST-SQUARES STAGE: @@OUTLIERS-BEGIN@@ none: every one of the 34 points lies within 3 K after the least-squares stage. @@OUTLIERS-END@@

@@RESULT-BEGIN@@ WHAT THE RUN FOUND (pinned in `expected`; measured, not argued). fit\_nrtl: b\_ij = 248.2 K, b\_ji = 191.8 K over 34 kept points of 2 laboratories (0 discarded): max |dT| 2.060 K, mean 0.709 K against the acceptance 1 K $\to$ notAccepted. fit\_nrtl\_start2 lands on the same pair (248.2, 191.8). NOT PROMOTABLE: max |dT| 2.060 K over the 34 kept points is not below 1 K @@RESULT-END@@

WHAT IS NOT CLAIMED. The pair is correct at these pressures and in this frame (ideal-gas vapour, the catalogue's own vapour-pressure records); the measurements were transcribed by the archive and have NOT been read back against the articles (the CHECKLIST that `fetch` writes is the aid).

\end{tutnarrative}
\begin{tutgolden}
\goldrow{diag}{fit\_nrtl}{bubble\_aad\_fit}{0.709411923703}
\goldrow{diag}{fit\_nrtl}{bubble\_max\_fit}{2.05988300875}
\goldrow{diag}{fit\_nrtl}{chi2}{85.4276550747}
\goldrow{diag}{fit\_nrtl}{chi2\_in\_sample}{85.4276550747}
\goldrow{diag}{fit\_nrtl}{chi2\_reduced}{2.66961422108}
\goldrow{diag}{fit\_nrtl}{cond\_JtJ}{74.8106212741}
\goldrow{diag}{fit\_nrtl}{converged}{1}
\goldrow{diag}{fit\_nrtl}{corr\_computed}{1}
\goldrow{diag}{fit\_nrtl}{dof}{32}
\goldrow{diag}{fit\_nrtl}{identifiable}{1}
\goldrow{diag}{fit\_nrtl}{invertible}{1}
\goldrow{diag}{fit\_nrtl}{iter}{3}
\goldrow{diag}{fit\_nrtl}{max\_abs\_corr}{0.992699920241}
\goldrow{diag}{fit\_nrtl}{max\_abs\_resid}{2.05988300875}
\goldmore{62}
\end{tutgolden}

\clearpage
\tutentry{curate141\_nrtl\_nNDimethylformamide\_nNDimethylacetamide}{tutorials/props/curation/curate141\_nrtl\_nNDimethylformamide\_nNDimethylacetamide}
\begin{tutscheme}
\begin{tikzpicture}[>={Stealth[length=2mm]}, every node/.style={font=\scriptsize}]
  \node[draw, rounded corners=2pt, align=center, minimum width=3.1cm, minimum height=0.8cm, fill=black!4] at (0.0,-0.0) {\textbf{fit\_nrtl}\\[-1pt]{\tiny fitParameters}};
  \node[draw, rounded corners=2pt, align=center, minimum width=3.1cm, minimum height=0.8cm, fill=black!4] at (3.6,-0.0) {\textbf{fit\_nrtl\_start2}\\[-1pt]{\tiny fitParameters}};
  \node[draw, rounded corners=2pt, align=center, minimum width=3.1cm, minimum height=0.8cm, fill=black!4] at (7.2,-0.0) {\textbf{txy\_ideal}\\[-1pt]{\tiny propertyScan1D}};
  \node[draw, rounded corners=2pt, align=center, minimum width=3.1cm, minimum height=0.8cm, fill=black!4] at (0.0,-1.3) {\textbf{txy\_nrtl}\\[-1pt]{\tiny propertyScan1D}};
\end{tikzpicture}
\end{tutscheme}
\tutwhat{NRTL nNDimethylformamide-nNDimethylacetamide regressed by the curation pipeline to every laboratory's measured bubble temperatures (ThermoML locates, the articles source), accepted at max |dT| < 1 K}
\begin{tutnarrative}
AN NRTL PAIR FROM EVERY LABORATORY'S MEASUREMENTS.

WRITTEN BY bin/choupo-curate pair nNDimethylformamide nNDimethylacetamide (DEV.md 4c C284, the bulk run; record docs/design/binary-pairs-from-open-measurements.md section 11).

THE DECLARATION -- WRITTEN AND COMMITTED BEFORE ANY FIT RAN (several declarations of the bulk run share one commit).

IDENTITY. nNDimethylformamide (CAS 68-12-2, InChIKey ZMXDDKWLCZADIW-UHFFFAOYSA-N) + nNDimethylacetamide (CAS 127-19-5, InChIKey FXHOOIRPVKKKFG-UHFFFAOYSA-N), joined to the archive through identity.json -- never by name. CLASS (C284): both records are volatile molecular liquids (no role nonvolatile/solute, no dissociatesTo, no charge, no solid, a vapour-pressure model), and every measured temperature lies below both critical temperatures.

NO MEASURED SERIES IN THIS TREE. constant/archiveDatasets names each dataset (DOI, NIST URL, article sha256, extraction selector, extract sha256); `bin/choupo-curate fetch` writes the points into constant/archiveExtracts/, which is gitignored.

STUDIES: every bubble-temperature block of the pair at or below 110 kPa with a declared uncertainty (2 block(s)). FITTED, ALL OF THEM (2): 10.1016/j.fluid.2015.12.052 block 0 18 pts 101.3 kPa 424.02-437.59 K U 0.1-0.54 K 10.1016/j.fluid.2008.07.018 block 19 16 pts 95.5 kPa 415.85-421.15 K U 0.88-0.88 K (reports x of the other member: written as x[nNDimethylformamide] = 1 - x by --x-of) SET ASIDE BY THE EXTRACTION RULES, BY NAME: (none)

THE RULE (DEV.md 4c C283 amended, C284). Every laboratory above is fitted; none is held out (the C16 partition and its band are retired as the acceptance test). NRTL, a\_ij = a\_ji = 0, alpha = 0.30 fixed: two parameters (b\_ij, b\_ji); each point weighted by 1/U, U its own stated expanded uncertainty. Objective twoStage (least squares, then minimax; T\_bubble absolute, in K). OUTLIERS: after a least-squares stage over every point, a point with |T\_model - T| > 3 K is discarded BY NAME (`discard` on its evidence entry, listed below) and the fit repeated without it. ACCEPTED when max |T\_model - T| over the KEPT points is below 1 K. C283's 5 \% relative clause judges composition; this fit has no composition residual, so only the temperature clause applies. A pair that splits ONE liquid inside its validity span is not promoted (arm f). Two starts: (0, 0) and (800, -400).

OUTLIERS DISCARDED AFTER THE LEAST-SQUARES STAGE: @@OUTLIERS-BEGIN@@ 15 point(s) of 34: doi 10.1016/j.fluid.2008.07.018 row 2: x 0.8445, P 95.500 kPa, T 419.85 K, dT +3.75 K; doi 10.1016/j.fluid.2008.07.018 row 3: x 0.7835, P 95.500 kPa, T 419.15 K, dT +4.76 K; doi 10.1016/j.fluid.2008.07.018 row 4: x 0.7307, P 95.500 kPa, T 418.65 K, dT +5.58 K; doi 10.1016/j.fluid.2008.07.018 row 5: x 0.6847, P 95.500 kPa, T 418.25 K, dT +6.29 K; doi 10.1016/j.fluid.2008.07.018 row 6: x 0.6441, P 95.500 kPa, T 417.95 K, dT +6.89 K; doi 10.1016/j.fluid.2008.07.018 row 7: x 0.6071, P 95.500 kPa, T 417.75 K, dT +7.39 K; doi 10.1016/j.fluid.2008.07.018 row 8: x 0.5468, P 95.500 kPa, T 417.45 K, dT +8.22 K; doi 10.1016/j.fluid.2008.07.018 row 9: x 0.5175, P 95.500 kPa, T 417.25 K, dT +8.71 K; doi 10.1016/j.fluid.2008.07.018 row 10: x 0.4841, P 95.500 kPa, T 417.05 K, dT +9.26 K; doi 10.1016/j.fluid.2008.07.018 row 11: x 0.4457, P 95.500 kPa, T 416.95 K, dT +9.79 K; doi 10.1016/j.fluid.2008.07.018 row 12: x 0.4012, P 95.500 kPa, T 416.75 K, dT +10.54 K; doi 10.1016/j.fluid.2008.07.018 row 13: x 0.3489, P 95.500 kPa, T 416.55 K, dT +11.46 K; doi 10.1016/j.fluid.2008.07.018 row 14: x 0.2867, P 95.500 kPa, T 416.35 K, dT +12.64 K; doi 10.1016/j.fluid.2008.07.018 row 15: x 0.2114, P 95.500 kPa, T 416.15 K, dT +14.27 K; doi 10.1016/j.fluid.2008.07.018 row 16: x 0.1182, P 95.500 kPa, T 415.85 K, dT +16.83 K @@OUTLIERS-END@@

@@RESULT-BEGIN@@ WHAT THE RUN FOUND (pinned in `expected`; measured, not argued). fit\_nrtl: b\_ij = 742.2 K, b\_ji = -361.3 K over 19 kept points of 2 laboratories (15 discarded): max |dT| 1.437 K, mean 0.862 K against the acceptance 1 K $\to$ notAccepted. fit\_nrtl\_start2 lands on the same pair (742.2, -361.3). NOT PROMOTABLE: max |dT| 1.437 K over the 19 kept points is not below 1 K @@RESULT-END@@

WHAT IS NOT CLAIMED. The pair is correct at these pressures and in this frame (ideal-gas vapour, the catalogue's own vapour-pressure records); the measurements were transcribed by the archive and have NOT been read back against the articles (the CHECKLIST that `fetch` writes is the aid).

\end{tutnarrative}
\begin{tutgolden}
\goldrow{diag}{fit\_nrtl}{bubble\_aad\_fit}{0.861676942458}
\goldrow{diag}{fit\_nrtl}{bubble\_max\_fit}{1.43730671345}
\goldrow{diag}{fit\_nrtl}{chi2}{984.472247923}
\goldrow{diag}{fit\_nrtl}{chi2\_in\_sample}{984.472247923}
\goldrow{diag}{fit\_nrtl}{chi2\_reduced}{57.9101322308}
\goldrow{diag}{fit\_nrtl}{cond\_JtJ}{53.000831938}
\goldrow{diag}{fit\_nrtl}{converged}{1}
\goldrow{diag}{fit\_nrtl}{corr\_computed}{1}
\goldrow{diag}{fit\_nrtl}{dof}{17}
\goldrow{diag}{fit\_nrtl}{identifiable}{1}
\goldrow{diag}{fit\_nrtl}{invertible}{1}
\goldrow{diag}{fit\_nrtl}{iter}{11}
\goldrow{diag}{fit\_nrtl}{max\_abs\_corr}{0.990697616217}
\goldrow{diag}{fit\_nrtl}{max\_abs\_resid}{1.43730671345}
\goldmore{56}
\end{tutgolden}

\clearpage
\tutentry{curate142\_nrtl\_ethylAcetate\_methanol}{tutorials/props/curation/curate142\_nrtl\_ethylAcetate\_methanol}
\begin{tutscheme}
\begin{tikzpicture}[>={Stealth[length=2mm]}, every node/.style={font=\scriptsize}]
  \node[draw, rounded corners=2pt, align=center, minimum width=3.1cm, minimum height=0.8cm, fill=black!4] at (0.0,-0.0) {\textbf{fit\_nrtl}\\[-1pt]{\tiny fitParameters}};
  \node[draw, rounded corners=2pt, align=center, minimum width=3.1cm, minimum height=0.8cm, fill=black!4] at (3.6,-0.0) {\textbf{fit\_nrtl\_start2}\\[-1pt]{\tiny fitParameters}};
  \node[draw, rounded corners=2pt, align=center, minimum width=3.1cm, minimum height=0.8cm, fill=black!4] at (7.2,-0.0) {\textbf{txy\_ideal}\\[-1pt]{\tiny propertyScan1D}};
  \node[draw, rounded corners=2pt, align=center, minimum width=3.1cm, minimum height=0.8cm, fill=black!4] at (0.0,-1.3) {\textbf{txy\_unifac}\\[-1pt]{\tiny propertyScan1D}};
  \node[draw, rounded corners=2pt, align=center, minimum width=3.1cm, minimum height=0.8cm, fill=black!4] at (3.6,-1.3) {\textbf{txy\_nrtl}\\[-1pt]{\tiny propertyScan1D}};
\end{tikzpicture}
\end{tutscheme}
\tutwhat{NRTL ethylAcetate-methanol regressed by the curation pipeline to every laboratory's measured bubble temperatures (ThermoML locates, the articles source), accepted at max |dT| < 1 K}
\begin{tutnarrative}
AN NRTL PAIR FROM EVERY LABORATORY'S MEASUREMENTS.

WRITTEN BY bin/choupo-curate pair ethylAcetate methanol (DEV.md 4c C284, the bulk run; record docs/design/binary-pairs-from-open-measurements.md section 11).

THE DECLARATION -- WRITTEN AND COMMITTED BEFORE ANY FIT RAN (several declarations of the bulk run share one commit).

IDENTITY. ethylAcetate (CAS 141-78-6, InChIKey XEKOWRVHYACXOJ-UHFFFAOYSA-N) + methanol (CAS 67-56-1, InChIKey OKKJLVBELUTLKV-UHFFFAOYSA-N), joined to the archive through identity.json -- never by name. CLASS (C284): both records are volatile molecular liquids (no role nonvolatile/solute, no dissociatesTo, no charge, no solid, a vapour-pressure model), and every measured temperature lies below both critical temperatures.

NO MEASURED SERIES IN THIS TREE. constant/archiveDatasets names each dataset (DOI, NIST URL, article sha256, extraction selector, extract sha256); `bin/choupo-curate fetch` writes the points into constant/archiveExtracts/, which is gitignored.

STUDIES: every bubble-temperature block of the pair at or below 110 kPa with a declared uncertainty (2 block(s)). FITTED, ALL OF THEM (2): 10.1021/acs.jced.5b00546 block 2 18 pts 101.3 kPa 335.59-348.74 K U 5-5.07 K 10.1021/acs.jced.8b00424 block 2 14 pts 101.3 kPa 335.45-346.80 K U 0.09-0.16 K SET ASIDE BY THE EXTRACTION RULES, BY NAME: (none)

THE RULE (DEV.md 4c C283 amended, C284). Every laboratory above is fitted; none is held out (the C16 partition and its band are retired as the acceptance test). NRTL, a\_ij = a\_ji = 0, alpha = 0.30 fixed: two parameters (b\_ij, b\_ji); each point weighted by 1/U, U its own stated expanded uncertainty. Objective twoStage (least squares, then minimax; T\_bubble absolute, in K). OUTLIERS: after a least-squares stage over every point, a point with |T\_model - T| > 3 K is discarded BY NAME (`discard` on its evidence entry, listed below) and the fit repeated without it. ACCEPTED when max |T\_model - T| over the KEPT points is below 1 K. C283's 5 \% relative clause judges composition; this fit has no composition residual, so only the temperature clause applies. A pair that splits ONE liquid inside its validity span is not promoted (arm f). Two starts: (0, 0) and (800, -400).

OUTLIERS DISCARDED AFTER THE LEAST-SQUARES STAGE: @@OUTLIERS-BEGIN@@ none: every one of the 32 points lies within 3 K after the least-squares stage. @@OUTLIERS-END@@

@@RESULT-BEGIN@@ WHAT THE RUN FOUND (pinned in `expected`; measured, not argued). fit\_nrtl: b\_ij = 202.5 K, b\_ji = 173.6 K over 32 kept points of 2 laboratories (0 discarded): max |dT| 0.498 K, mean 0.111 K against the acceptance 1 K $\to$ accepted. fit\_nrtl\_start2 lands on the same pair (202.5, 173.6). PROMOTED on 2026-11-10 into data/standards/parameters/NRTL/ethylAcetate-methanol.dat as reviewStatus autoPromoted, humanReviewed false, because: max |dT| 0.498 K < 1 K over the 32 kept points, one liquid across the validity span @@RESULT-END@@

WHAT IS NOT CLAIMED. The pair is correct at these pressures and in this frame (ideal-gas vapour, the catalogue's own vapour-pressure records); the measurements were transcribed by the archive and have NOT been read back against the articles (the CHECKLIST that `fetch` writes is the aid).

\end{tutnarrative}
\begin{tutgolden}
\goldrow{diag}{fit\_nrtl}{bubble\_aad\_fit}{0.110537463813}
\goldrow{diag}{fit\_nrtl}{bubble\_max\_fit}{0.497643236317}
\goldrow{diag}{fit\_nrtl}{chi2}{0.523318618177}
\goldrow{diag}{fit\_nrtl}{chi2\_in\_sample}{0.523318618177}
\goldrow{diag}{fit\_nrtl}{chi2\_reduced}{0.0174439539392}
\goldrow{diag}{fit\_nrtl}{cond\_JtJ}{48.4028048604}
\goldrow{diag}{fit\_nrtl}{converged}{1}
\goldrow{diag}{fit\_nrtl}{corr\_computed}{1}
\goldrow{diag}{fit\_nrtl}{dof}{30}
\goldrow{diag}{fit\_nrtl}{identifiable}{1}
\goldrow{diag}{fit\_nrtl}{invertible}{1}
\goldrow{diag}{fit\_nrtl}{iter}{5}
\goldrow{diag}{fit\_nrtl}{max\_abs\_corr}{0.98982737602}
\goldrow{diag}{fit\_nrtl}{max\_abs\_resid}{0.497643236317}
\goldmore{71}
\end{tutgolden}

\clearpage
\tutentry{curate143\_nrtl\_ethylBenzene\_monochlorobenzene}{tutorials/props/curation/curate143\_nrtl\_ethylBenzene\_monochlorobenzene}
\begin{tutscheme}
\begin{tikzpicture}[>={Stealth[length=2mm]}, every node/.style={font=\scriptsize}]
  \node[draw, rounded corners=2pt, align=center, minimum width=3.1cm, minimum height=0.8cm, fill=black!4] at (0.0,-0.0) {\textbf{fit\_nrtl}\\[-1pt]{\tiny fitParameters}};
  \node[draw, rounded corners=2pt, align=center, minimum width=3.1cm, minimum height=0.8cm, fill=black!4] at (3.6,-0.0) {\textbf{fit\_nrtl\_start2}\\[-1pt]{\tiny fitParameters}};
  \node[draw, rounded corners=2pt, align=center, minimum width=3.1cm, minimum height=0.8cm, fill=black!4] at (7.2,-0.0) {\textbf{txy\_ideal}\\[-1pt]{\tiny propertyScan1D}};
  \node[draw, rounded corners=2pt, align=center, minimum width=3.1cm, minimum height=0.8cm, fill=black!4] at (0.0,-1.3) {\textbf{txy\_nrtl}\\[-1pt]{\tiny propertyScan1D}};
\end{tikzpicture}
\end{tutscheme}
\tutwhat{NRTL ethylBenzene-monochlorobenzene regressed by the curation pipeline to every laboratory's measured bubble temperatures (ThermoML locates, the articles source), accepted at max |dT| < 1 K}
\begin{tutnarrative}
AN NRTL PAIR FROM EVERY LABORATORY'S MEASUREMENTS.

WRITTEN BY bin/choupo-curate pair ethylBenzene monochlorobenzene (DEV.md 4c C284, the bulk run; record docs/design/binary-pairs-from-open-measurements.md section 11).

THE DECLARATION -- WRITTEN AND COMMITTED BEFORE ANY FIT RAN (several declarations of the bulk run share one commit).

IDENTITY. ethylBenzene (CAS 100-41-4, InChIKey YNQLUTRBYVCPMQ-UHFFFAOYSA-N) + monochlorobenzene (CAS 108-90-7, InChIKey MVPPADPHJFYWMZ-UHFFFAOYSA-N), joined to the archive through identity.json -- never by name. CLASS (C284): both records are volatile molecular liquids (no role nonvolatile/solute, no dissociatesTo, no charge, no solid, a vapour-pressure model), and every measured temperature lies below both critical temperatures.

NO MEASURED SERIES IN THIS TREE. constant/archiveDatasets names each dataset (DOI, NIST URL, article sha256, extraction selector, extract sha256); `bin/choupo-curate fetch` writes the points into constant/archiveExtracts/, which is gitignored.

STUDIES: every bubble-temperature block of the pair at or below 110 kPa with a declared uncertainty (2 block(s)). FITTED, ALL OF THEM (2): 10.1021/je020186c block 5 20 pts 4 isobars 26.66-101.32 kPa 362.12-408.76 K U 0.25-0.31 K 10.1016/j.jct.2005.12.009 block 30 6 pts 94.7 kPa 403.15-406.35 K U 0.69-0.69 K (reports x of the other member: written as x[ethylBenzene] = 1 - x by --x-of) SET ASIDE BY THE EXTRACTION RULES, BY NAME: (none)

THE RULE (DEV.md 4c C283 amended, C284). Every laboratory above is fitted; none is held out (the C16 partition and its band are retired as the acceptance test). NRTL, a\_ij = a\_ji = 0, alpha = 0.30 fixed: two parameters (b\_ij, b\_ji); each point weighted by 1/U, U its own stated expanded uncertainty. Objective twoStage (least squares, then minimax; T\_bubble absolute, in K). OUTLIERS: after a least-squares stage over every point, a point with |T\_model - T| > 3 K is discarded BY NAME (`discard` on its evidence entry, listed below) and the fit repeated without it. ACCEPTED when max |T\_model - T| over the KEPT points is below 1 K. C283's 5 \% relative clause judges composition; this fit has no composition residual, so only the temperature clause applies. A pair that splits ONE liquid inside its validity span is not promoted (arm f). Two starts: (0, 0) and (800, -400).

OUTLIERS DISCARDED AFTER THE LEAST-SQUARES STAGE: @@OUTLIERS-BEGIN@@ none: every one of the 26 points lies within 3 K after the least-squares stage. @@OUTLIERS-END@@

@@RESULT-BEGIN@@ WHAT THE RUN FOUND (pinned in `expected`; measured, not argued). fit\_nrtl: b\_ij = 319.8 K, b\_ji = -258.9 K over 26 kept points of 2 laboratories (0 discarded): max |dT| 0.377 K, mean 0.131 K against the acceptance 1 K $\to$ accepted. fit\_nrtl\_start2 lands on the same pair (319.8, -258.9). PROMOTED on 2026-11-10 into data/standards/parameters/NRTL/ethylBenzene-monochlorobenzene.dat as reviewStatus autoPromoted, humanReviewed false, because: max |dT| 0.377 K < 1 K over the 26 kept points, one liquid across the validity span @@RESULT-END@@

WHAT IS NOT CLAIMED. The pair is correct at these pressures and in this frame (ideal-gas vapour, the catalogue's own vapour-pressure records); the measurements were transcribed by the archive and have NOT been read back against the articles (the CHECKLIST that `fetch` writes is the aid).

\end{tutnarrative}
\begin{tutgolden}
\goldrow{diag}{fit\_nrtl}{bubble\_aad\_fit}{0.130732217619}
\goldrow{diag}{fit\_nrtl}{bubble\_max\_fit}{0.377382817192}
\goldrow{diag}{fit\_nrtl}{chi2}{3.83323491789}
\goldrow{diag}{fit\_nrtl}{chi2\_in\_sample}{3.83323491789}
\goldrow{diag}{fit\_nrtl}{chi2\_reduced}{0.159718121579}
\goldrow{diag}{fit\_nrtl}{cond\_JtJ}{806.361311335}
\goldrow{diag}{fit\_nrtl}{converged}{1}
\goldrow{diag}{fit\_nrtl}{corr\_computed}{1}
\goldrow{diag}{fit\_nrtl}{dof}{24}
\goldrow{diag}{fit\_nrtl}{identifiable}{0}
\goldrow{diag}{fit\_nrtl}{invertible}{1}
\goldrow{diag}{fit\_nrtl}{iter}{21}
\goldrow{diag}{fit\_nrtl}{max\_abs\_corr}{0.999380506702}
\goldrow{diag}{fit\_nrtl}{max\_abs\_resid}{0.377382817192}
\goldmore{50}
\end{tutgolden}

\clearpage
\tutentry{curate144\_nrtl\_2Hexanone\_1Pentanol}{tutorials/props/curation/curate144\_nrtl\_2Hexanone\_1Pentanol}
\begin{tutscheme}
\begin{tikzpicture}[>={Stealth[length=2mm]}, every node/.style={font=\scriptsize}]
  \node[draw, rounded corners=2pt, align=center, minimum width=3.1cm, minimum height=0.8cm, fill=black!4] at (0.0,-0.0) {\textbf{fit\_nrtl}\\[-1pt]{\tiny fitParameters}};
  \node[draw, rounded corners=2pt, align=center, minimum width=3.1cm, minimum height=0.8cm, fill=black!4] at (3.6,-0.0) {\textbf{fit\_nrtl\_start2}\\[-1pt]{\tiny fitParameters}};
  \node[draw, rounded corners=2pt, align=center, minimum width=3.1cm, minimum height=0.8cm, fill=black!4] at (7.2,-0.0) {\textbf{txy\_ideal}\\[-1pt]{\tiny propertyScan1D}};
  \node[draw, rounded corners=2pt, align=center, minimum width=3.1cm, minimum height=0.8cm, fill=black!4] at (0.0,-1.3) {\textbf{txy\_unifac}\\[-1pt]{\tiny propertyScan1D}};
  \node[draw, rounded corners=2pt, align=center, minimum width=3.1cm, minimum height=0.8cm, fill=black!4] at (3.6,-1.3) {\textbf{txy\_nrtl}\\[-1pt]{\tiny propertyScan1D}};
\end{tikzpicture}
\end{tutscheme}
\tutwhat{NRTL 2Hexanone-1Pentanol regressed by the curation pipeline to every laboratory's measured bubble temperatures (ThermoML locates, the articles source), accepted at max |dT| < 1 K}
\begin{tutnarrative}
AN NRTL PAIR FROM EVERY LABORATORY'S MEASUREMENTS.

WRITTEN BY bin/choupo-curate pair 2Hexanone 1Pentanol (DEV.md 4c C284, the bulk run; record docs/design/binary-pairs-from-open-measurements.md section 11).

THE DECLARATION -- WRITTEN AND COMMITTED BEFORE ANY FIT RAN (several declarations of the bulk run share one commit).

IDENTITY. 2Hexanone (CAS 591-78-6, InChIKey QQZOPKMRPOGIEB-UHFFFAOYSA-N) + 1Pentanol (CAS 71-41-0, InChIKey AMQJEAYHLZJPGS-UHFFFAOYSA-N), joined to the archive through identity.json -- never by name. CLASS (C284): both records are volatile molecular liquids (no role nonvolatile/solute, no dissociatesTo, no charge, no solid, a vapour-pressure model), and every measured temperature lies below both critical temperatures.

NO MEASURED SERIES IN THIS TREE. constant/archiveDatasets names each dataset (DOI, NIST URL, article sha256, extraction selector, extract sha256); `bin/choupo-curate fetch` writes the points into constant/archiveExtracts/, which is gitignored.

STUDIES: every bubble-temperature block of the pair at or below 110 kPa with a declared uncertainty (2 block(s)). FITTED, ALL OF THEM (2): 10.1021/je0502905 block 9 21 pts 3 isobars 26.66-79.99 kPa 360.11-400.55 K U 0.25-0.29 K 10.1021/je900058p block 10 6 pts 101.32 kPa 401.07-408.31 K U 0.19-0.19 K SET ASIDE BY THE EXTRACTION RULES, BY NAME: (none)

THE RULE (DEV.md 4c C283 amended, C284). Every laboratory above is fitted; none is held out (the C16 partition and its band are retired as the acceptance test). NRTL, a\_ij = a\_ji = 0, alpha = 0.30 fixed: two parameters (b\_ij, b\_ji); each point weighted by 1/U, U its own stated expanded uncertainty. Objective twoStage (least squares, then minimax; T\_bubble absolute, in K). OUTLIERS: after a least-squares stage over every point, a point with |T\_model - T| > 3 K is discarded BY NAME (`discard` on its evidence entry, listed below) and the fit repeated without it. ACCEPTED when max |T\_model - T| over the KEPT points is below 1 K. C283's 5 \% relative clause judges composition; this fit has no composition residual, so only the temperature clause applies. A pair that splits ONE liquid inside its validity span is not promoted (arm f). Two starts: (0, 0) and (800, -400).

OUTLIERS DISCARDED AFTER THE LEAST-SQUARES STAGE: @@OUTLIERS-BEGIN@@ none: every one of the 27 points lies within 3 K after the least-squares stage. @@OUTLIERS-END@@

@@RESULT-BEGIN@@ WHAT THE RUN FOUND (pinned in `expected`; measured, not argued). fit\_nrtl: b\_ij = 268.1 K, b\_ji = -141.6 K over 27 kept points of 2 laboratories (0 discarded): max |dT| 0.498 K, mean 0.191 K against the acceptance 1 K $\to$ accepted. fit\_nrtl\_start2 lands on the same pair (268.1, -141.6). PROMOTED on 2026-11-10 into data/standards/parameters/NRTL/1Pentanol-2Hexanone.dat as reviewStatus autoPromoted, humanReviewed false, because: max |dT| 0.498 K < 1 K over the 27 kept points, one liquid across the validity span @@RESULT-END@@

WHAT IS NOT CLAIMED. The pair is correct at these pressures and in this frame (ideal-gas vapour, the catalogue's own vapour-pressure records); the measurements were transcribed by the archive and have NOT been read back against the articles (the CHECKLIST that `fetch` writes is the aid).

\end{tutnarrative}
\begin{tutgolden}
\goldrow{diag}{fit\_nrtl}{bubble\_aad\_fit}{0.190866253233}
\goldrow{diag}{fit\_nrtl}{bubble\_max\_fit}{0.497664850908}
\goldrow{diag}{fit\_nrtl}{chi2}{22.8018123012}
\goldrow{diag}{fit\_nrtl}{chi2\_in\_sample}{22.8018123012}
\goldrow{diag}{fit\_nrtl}{chi2\_reduced}{0.912072492048}
\goldrow{diag}{fit\_nrtl}{cond\_JtJ}{582.370960517}
\goldrow{diag}{fit\_nrtl}{converged}{1}
\goldrow{diag}{fit\_nrtl}{corr\_computed}{1}
\goldrow{diag}{fit\_nrtl}{dof}{25}
\goldrow{diag}{fit\_nrtl}{identifiable}{0}
\goldrow{diag}{fit\_nrtl}{invertible}{1}
\goldrow{diag}{fit\_nrtl}{iter}{10}
\goldrow{diag}{fit\_nrtl}{max\_abs\_corr}{0.999142544814}
\goldrow{diag}{fit\_nrtl}{max\_abs\_resid}{0.497664850908}
\goldmore{53}
\end{tutgolden}

\clearpage
\tutentry{curate145\_nrtl\_ethanol\_ethyleneGlycol}{tutorials/props/curation/curate145\_nrtl\_ethanol\_ethyleneGlycol}
\begin{tutscheme}
\begin{tikzpicture}[>={Stealth[length=2mm]}, every node/.style={font=\scriptsize}]
  \node[draw, rounded corners=2pt, align=center, minimum width=3.1cm, minimum height=0.8cm, fill=black!4] at (0.0,-0.0) {\textbf{fit\_nrtl}\\[-1pt]{\tiny fitParameters}};
  \node[draw, rounded corners=2pt, align=center, minimum width=3.1cm, minimum height=0.8cm, fill=black!4] at (3.6,-0.0) {\textbf{fit\_nrtl\_start2}\\[-1pt]{\tiny fitParameters}};
  \node[draw, rounded corners=2pt, align=center, minimum width=3.1cm, minimum height=0.8cm, fill=black!4] at (7.2,-0.0) {\textbf{txy\_ideal}\\[-1pt]{\tiny propertyScan1D}};
  \node[draw, rounded corners=2pt, align=center, minimum width=3.1cm, minimum height=0.8cm, fill=black!4] at (0.0,-1.3) {\textbf{txy\_nrtl}\\[-1pt]{\tiny propertyScan1D}};
\end{tikzpicture}
\end{tutscheme}
\tutwhat{NRTL ethanol-ethyleneGlycol regressed by the curation pipeline to every laboratory's measured bubble temperatures (ThermoML locates, the articles source), accepted at max |dT| < 1 K}
\begin{tutnarrative}
AN NRTL PAIR FROM EVERY LABORATORY'S MEASUREMENTS.

WRITTEN BY bin/choupo-curate pair ethanol ethyleneGlycol (DEV.md 4c C284, the bulk run; record docs/design/binary-pairs-from-open-measurements.md section 11).

THE DECLARATION -- WRITTEN AND COMMITTED BEFORE ANY FIT RAN (several declarations of the bulk run share one commit).

IDENTITY. ethanol (CAS 64-17-5, InChIKey LFQSCWFLJHTTHZ-UHFFFAOYSA-N) + ethyleneGlycol (CAS 107-21-1, InChIKey LYCAIKOWRPUZTN-UHFFFAOYSA-N), joined to the archive through identity.json -- never by name. CLASS (C284): both records are volatile molecular liquids (no role nonvolatile/solute, no dissociatesTo, no charge, no solid, a vapour-pressure model), and every measured temperature lies below both critical temperatures.

NO MEASURED SERIES IN THIS TREE. constant/archiveDatasets names each dataset (DOI, NIST URL, article sha256, extraction selector, extract sha256); `bin/choupo-curate fetch` writes the points into constant/archiveExtracts/, which is gitignored.

STUDIES: every bubble-temperature block of the pair at or below 110 kPa with a declared uncertainty (2 block(s)). FITTED, ALL OF THEM (2): 10.1021/acs.jced.6b00384 block 16 19 pts 101 kPa 352.50-384.80 K U 0.8-0.9 K 10.1021/je2008704 block 7 15 pts 101.3 kPa 354.00-389.79 K U 0.37-1.9 K SET ASIDE BY THE EXTRACTION RULES, BY NAME: (none)

THE RULE (DEV.md 4c C283 amended, C284). Every laboratory above is fitted; none is held out (the C16 partition and its band are retired as the acceptance test). NRTL, a\_ij = a\_ji = 0, alpha = 0.30 fixed: two parameters (b\_ij, b\_ji); each point weighted by 1/U, U its own stated expanded uncertainty. Objective twoStage (least squares, then minimax; T\_bubble absolute, in K). OUTLIERS: after a least-squares stage over every point, a point with |T\_model - T| > 3 K is discarded BY NAME (`discard` on its evidence entry, listed below) and the fit repeated without it. ACCEPTED when max |T\_model - T| over the KEPT points is below 1 K. C283's 5 \% relative clause judges composition; this fit has no composition residual, so only the temperature clause applies. A pair that splits ONE liquid inside its validity span is not promoted (arm f). Two starts: (0, 0) and (800, -400).

OUTLIERS DISCARDED AFTER THE LEAST-SQUARES STAGE: @@OUTLIERS-BEGIN@@ none: every one of the 34 points lies within 3 K after the least-squares stage. @@OUTLIERS-END@@

@@RESULT-BEGIN@@ WHAT THE RUN FOUND (pinned in `expected`; measured, not argued). fit\_nrtl: b\_ij = 12.3 K, b\_ji = 221.8 K over 34 kept points of 2 laboratories (0 discarded): max |dT| 1.434 K, mean 0.257 K against the acceptance 1 K $\to$ notAccepted. fit\_nrtl\_start2 lands on the same pair (12.3, 221.8). NOT PROMOTABLE: max |dT| 1.434 K over the 34 kept points is not below 1 K @@RESULT-END@@

WHAT IS NOT CLAIMED. The pair is correct at these pressures and in this frame (ideal-gas vapour, the catalogue's own vapour-pressure records); the measurements were transcribed by the archive and have NOT been read back against the articles (the CHECKLIST that `fetch` writes is the aid).

\end{tutnarrative}
\begin{tutgolden}
\goldrow{diag}{fit\_nrtl}{bubble\_aad\_fit}{0.256604555617}
\goldrow{diag}{fit\_nrtl}{bubble\_max\_fit}{1.43354829854}
\goldrow{diag}{fit\_nrtl}{chi2}{4.35228033761}
\goldrow{diag}{fit\_nrtl}{chi2\_in\_sample}{4.35228033761}
\goldrow{diag}{fit\_nrtl}{chi2\_reduced}{0.13600876055}
\goldrow{diag}{fit\_nrtl}{cond\_JtJ}{154.125919267}
\goldrow{diag}{fit\_nrtl}{converged}{1}
\goldrow{diag}{fit\_nrtl}{corr\_computed}{1}
\goldrow{diag}{fit\_nrtl}{dof}{32}
\goldrow{diag}{fit\_nrtl}{identifiable}{1}
\goldrow{diag}{fit\_nrtl}{invertible}{1}
\goldrow{diag}{fit\_nrtl}{iter}{5}
\goldrow{diag}{fit\_nrtl}{max\_abs\_corr}{0.994480655698}
\goldrow{diag}{fit\_nrtl}{max\_abs\_resid}{1.43354829854}
\goldmore{62}
\end{tutgolden}

\clearpage
\tutentry{curate146\_nrtl\_2Methyl2Propanol\_nNDimethylformamide}{tutorials/props/curation/curate146\_nrtl\_2Methyl2Propanol\_nNDimethylformamide}
\begin{tutscheme}
\begin{tikzpicture}[>={Stealth[length=2mm]}, every node/.style={font=\scriptsize}]
  \node[draw, rounded corners=2pt, align=center, minimum width=3.1cm, minimum height=0.8cm, fill=black!4] at (0.0,-0.0) {\textbf{fit\_nrtl}\\[-1pt]{\tiny fitParameters}};
  \node[draw, rounded corners=2pt, align=center, minimum width=3.1cm, minimum height=0.8cm, fill=black!4] at (3.6,-0.0) {\textbf{fit\_nrtl\_start2}\\[-1pt]{\tiny fitParameters}};
  \node[draw, rounded corners=2pt, align=center, minimum width=3.1cm, minimum height=0.8cm, fill=black!4] at (7.2,-0.0) {\textbf{txy\_ideal}\\[-1pt]{\tiny propertyScan1D}};
  \node[draw, rounded corners=2pt, align=center, minimum width=3.1cm, minimum height=0.8cm, fill=black!4] at (0.0,-1.3) {\textbf{txy\_nrtl}\\[-1pt]{\tiny propertyScan1D}};
\end{tikzpicture}
\end{tutscheme}
\tutwhat{NRTL 2Methyl2Propanol-nNDimethylformamide regressed by the curation pipeline to every laboratory's measured bubble temperatures (ThermoML locates, the articles source), accepted at max |dT| < 1 K}
\begin{tutnarrative}
AN NRTL PAIR FROM EVERY LABORATORY'S MEASUREMENTS.

WRITTEN BY bin/choupo-curate pair 2Methyl2Propanol nNDimethylformamide (DEV.md 4c C284, the bulk run; record docs/design/binary-pairs-from-open-measurements.md section 11).

THE DECLARATION -- WRITTEN AND COMMITTED BEFORE ANY FIT RAN (several declarations of the bulk run share one commit).

IDENTITY. 2Methyl2Propanol (CAS 75-65-0, InChIKey DKGAVHZHDRPRBM-UHFFFAOYSA-N) + nNDimethylformamide (CAS 68-12-2, InChIKey ZMXDDKWLCZADIW-UHFFFAOYSA-N), joined to the archive through identity.json -- never by name. CLASS (C284): both records are volatile molecular liquids (no role nonvolatile/solute, no dissociatesTo, no charge, no solid, a vapour-pressure model), and every measured temperature lies below both critical temperatures.

NO MEASURED SERIES IN THIS TREE. constant/archiveDatasets names each dataset (DOI, NIST URL, article sha256, extraction selector, extract sha256); `bin/choupo-curate fetch` writes the points into constant/archiveExtracts/, which is gitignored.

STUDIES: every bubble-temperature block of the pair at or below 110 kPa with a declared uncertainty (2 block(s)). FITTED, ALL OF THEM (2): 10.1016/j.fluid.2006.08.018 block 34 19 pts 95.5 kPa 355.35-411.55 K U 0.69-0.69 K 10.1021/je7001336 block 19 10 pts 94.4 kPa 355.55-406.45 K U 0.69-0.73 K SET ASIDE BY THE EXTRACTION RULES, BY NAME: (none)

THE RULE (DEV.md 4c C283 amended, C284). Every laboratory above is fitted; none is held out (the C16 partition and its band are retired as the acceptance test). NRTL, a\_ij = a\_ji = 0, alpha = 0.30 fixed: two parameters (b\_ij, b\_ji); each point weighted by 1/U, U its own stated expanded uncertainty. Objective twoStage (least squares, then minimax; T\_bubble absolute, in K). OUTLIERS: after a least-squares stage over every point, a point with |T\_model - T| > 3 K is discarded BY NAME (`discard` on its evidence entry, listed below) and the fit repeated without it. ACCEPTED when max |T\_model - T| over the KEPT points is below 1 K. C283's 5 \% relative clause judges composition; this fit has no composition residual, so only the temperature clause applies. A pair that splits ONE liquid inside its validity span is not promoted (arm f). Two starts: (0, 0) and (800, -400).

OUTLIERS DISCARDED AFTER THE LEAST-SQUARES STAGE: @@OUTLIERS-BEGIN@@ 5 point(s) of 29: doi 10.1021/je7001336 row 2: x 0.1192, P 94.400 kPa, T 395.05 K, dT +3.96 K; doi 10.1021/je7001336 row 3: x 0.2130, P 94.400 kPa, T 383.75 K, dT +4.35 K; doi 10.1021/je7001336 row 4: x 0.3512, P 94.400 kPa, T 373.25 K, dT +3.99 K; doi 10.1021/je7001336 row 5: x 0.4036, P 94.400 kPa, T 370.35 K, dT +3.74 K; doi 10.1021/je7001336 row 6: x 0.5199, P 94.400 kPa, T 365.25 K, dT +3.07 K @@OUTLIERS-END@@

@@RESULT-BEGIN@@ WHAT THE RUN FOUND (pinned in `expected`; measured, not argued). fit\_nrtl: b\_ij = 418.3 K, b\_ji = -300.5 K over 24 kept points of 2 laboratories (5 discarded): max |dT| 2.570 K, mean 1.380 K against the acceptance 1 K $\to$ notAccepted. fit\_nrtl\_start2 lands on the same pair (418.3, -300.5). NOT PROMOTABLE: max |dT| 2.570 K over the 24 kept points is not below 1 K @@RESULT-END@@

WHAT IS NOT CLAIMED. The pair is correct at these pressures and in this frame (ideal-gas vapour, the catalogue's own vapour-pressure records); the measurements were transcribed by the archive and have NOT been read back against the articles (the CHECKLIST that `fetch` writes is the aid).

\end{tutnarrative}
\begin{tutgolden}
\goldrow{diag}{fit\_nrtl}{bubble\_aad\_fit}{1.38043649811}
\goldrow{diag}{fit\_nrtl}{bubble\_max\_fit}{2.57017190694}
\goldrow{diag}{fit\_nrtl}{chi2}{128.535162952}
\goldrow{diag}{fit\_nrtl}{chi2\_in\_sample}{128.535162952}
\goldrow{diag}{fit\_nrtl}{chi2\_reduced}{5.84250740692}
\goldrow{diag}{fit\_nrtl}{cond\_JtJ}{310.072704903}
\goldrow{diag}{fit\_nrtl}{converged}{1}
\goldrow{diag}{fit\_nrtl}{corr\_computed}{1}
\goldrow{diag}{fit\_nrtl}{dof}{22}
\goldrow{diag}{fit\_nrtl}{identifiable}{1}
\goldrow{diag}{fit\_nrtl}{invertible}{1}
\goldrow{diag}{fit\_nrtl}{iter}{19}
\goldrow{diag}{fit\_nrtl}{max\_abs\_corr}{0.998391364865}
\goldrow{diag}{fit\_nrtl}{max\_abs\_resid}{2.57017190694}
\goldmore{54}
\end{tutgolden}

\clearpage
\tutentry{curate147\_nrtl\_dimethylCarbonate\_oXylene}{tutorials/props/curation/curate147\_nrtl\_dimethylCarbonate\_oXylene}
\begin{tutscheme}
\begin{tikzpicture}[>={Stealth[length=2mm]}, every node/.style={font=\scriptsize}]
  \node[draw, rounded corners=2pt, align=center, minimum width=3.1cm, minimum height=0.8cm, fill=black!4] at (0.0,-0.0) {\textbf{fit\_nrtl}\\[-1pt]{\tiny fitParameters}};
  \node[draw, rounded corners=2pt, align=center, minimum width=3.1cm, minimum height=0.8cm, fill=black!4] at (3.6,-0.0) {\textbf{fit\_nrtl\_start2}\\[-1pt]{\tiny fitParameters}};
  \node[draw, rounded corners=2pt, align=center, minimum width=3.1cm, minimum height=0.8cm, fill=black!4] at (7.2,-0.0) {\textbf{txy\_ideal}\\[-1pt]{\tiny propertyScan1D}};
  \node[draw, rounded corners=2pt, align=center, minimum width=3.1cm, minimum height=0.8cm, fill=black!4] at (0.0,-1.3) {\textbf{txy\_nrtl}\\[-1pt]{\tiny propertyScan1D}};
\end{tikzpicture}
\end{tutscheme}
\tutwhat{NRTL dimethylCarbonate-oXylene regressed by the curation pipeline to every laboratory's measured bubble temperatures (ThermoML locates, the articles source), accepted at max |dT| < 1 K}
\begin{tutnarrative}
AN NRTL PAIR FROM EVERY LABORATORY'S MEASUREMENTS.

WRITTEN BY bin/choupo-curate pair dimethylCarbonate oXylene (DEV.md 4c C284, the bulk run; record docs/design/binary-pairs-from-open-measurements.md section 11).

THE DECLARATION -- WRITTEN AND COMMITTED BEFORE ANY FIT RAN (several declarations of the bulk run share one commit).

IDENTITY. dimethylCarbonate (CAS 616-38-6, InChIKey IEJIGPNLZYLLBP-UHFFFAOYSA-N) + oXylene (CAS 95-47-6, InChIKey CTQNGGLPUBDAKN-UHFFFAOYSA-N), joined to the archive through identity.json -- never by name. CLASS (C284): both records are volatile molecular liquids (no role nonvolatile/solute, no dissociatesTo, no charge, no solid, a vapour-pressure model), and every measured temperature lies below both critical temperatures.

NO MEASURED SERIES IN THIS TREE. constant/archiveDatasets names each dataset (DOI, NIST URL, article sha256, extraction selector, extract sha256); `bin/choupo-curate fetch` writes the points into constant/archiveExtracts/, which is gitignored.

STUDIES: every bubble-temperature block of the pair at or below 110 kPa with a declared uncertainty (2 block(s)). FITTED, ALL OF THEM (2): 10.1021/acs.jced.9b00414 block 7 18 pts 101.3 kPa 364.68-414.43 K U 0.27-2.01 K 10.1021/acs.jced.7b00372 block 21 11 pts 93.13 kPa 361.48-404.50 K U 0.24-1.07 K (reports x of the other member: written as x[dimethylCarbonate] = 1 - x by --x-of) SET ASIDE BY THE EXTRACTION RULES, BY NAME: (none)

THE RULE (DEV.md 4c C283 amended, C284). Every laboratory above is fitted; none is held out (the C16 partition and its band are retired as the acceptance test). NRTL, a\_ij = a\_ji = 0, alpha = 0.30 fixed: two parameters (b\_ij, b\_ji); each point weighted by 1/U, U its own stated expanded uncertainty. Objective twoStage (least squares, then minimax; T\_bubble absolute, in K). OUTLIERS: after a least-squares stage over every point, a point with |T\_model - T| > 3 K is discarded BY NAME (`discard` on its evidence entry, listed below) and the fit repeated without it. ACCEPTED when max |T\_model - T| over the KEPT points is below 1 K. C283's 5 \% relative clause judges composition; this fit has no composition residual, so only the temperature clause applies. A pair that splits ONE liquid inside its validity span is not promoted (arm f). Two starts: (0, 0) and (800, -400).

OUTLIERS DISCARDED AFTER THE LEAST-SQUARES STAGE: @@OUTLIERS-BEGIN@@ 5 point(s) of 29: doi 10.1021/acs.jced.7b00372 row 7: x 0.4130, P 93.130 kPa, T 374.88 K, dT -3.03 K; doi 10.1021/acs.jced.7b00372 row 8: x 0.2655, P 93.130 kPa, T 382.02 K, dT -4.32 K; doi 10.1021/acs.jced.7b00372 row 9: x 0.1574, P 93.130 kPa, T 390.84 K, dT -5.37 K; doi 10.1021/acs.jced.7b00372 row 10: x 0.0975, P 93.130 kPa, T 397.46 K, dT -4.89 K; doi 10.1021/acs.jced.7b00372 row 11: x 0.0502, P 93.130 kPa, T 404.50 K, dT -3.61 K @@OUTLIERS-END@@

@@RESULT-BEGIN@@ WHAT THE RUN FOUND (pinned in `expected`; measured, not argued). fit\_nrtl: b\_ij = 36.1 K, b\_ji = 481.2 K over 24 kept points of 2 laboratories (5 discarded): max |dT| 4.241 K, mean 1.741 K against the acceptance 1 K $\to$ notAccepted. fit\_nrtl\_start2 lands on the same pair (36.1, 481.2). NOT PROMOTABLE: max |dT| 4.241 K over the 24 kept points is not below 1 K @@RESULT-END@@

WHAT IS NOT CLAIMED. The pair is correct at these pressures and in this frame (ideal-gas vapour, the catalogue's own vapour-pressure records); the measurements were transcribed by the archive and have NOT been read back against the articles (the CHECKLIST that `fetch` writes is the aid).

\end{tutnarrative}
\begin{tutgolden}
\goldrow{diag}{fit\_nrtl}{bubble\_aad\_fit}{1.74086147733}
\goldrow{diag}{fit\_nrtl}{bubble\_max\_fit}{4.24138405089}
\goldrow{diag}{fit\_nrtl}{chi2}{160.64542525}
\goldrow{diag}{fit\_nrtl}{chi2\_in\_sample}{160.64542525}
\goldrow{diag}{fit\_nrtl}{chi2\_reduced}{7.30206478411}
\goldrow{diag}{fit\_nrtl}{cond\_JtJ}{39.4396305986}
\goldrow{diag}{fit\_nrtl}{converged}{1}
\goldrow{diag}{fit\_nrtl}{corr\_computed}{1}
\goldrow{diag}{fit\_nrtl}{dof}{22}
\goldrow{diag}{fit\_nrtl}{identifiable}{1}
\goldrow{diag}{fit\_nrtl}{invertible}{1}
\goldrow{diag}{fit\_nrtl}{iter}{5}
\goldrow{diag}{fit\_nrtl}{max\_abs\_corr}{0.967834099998}
\goldrow{diag}{fit\_nrtl}{max\_abs\_resid}{4.24138405089}
\goldmore{54}
\end{tutgolden}

\clearpage
\tutentry{curate148\_nrtl\_nButanol\_nButylAcetate}{tutorials/props/curation/curate148\_nrtl\_nButanol\_nButylAcetate}
\begin{tutscheme}
\begin{tikzpicture}[>={Stealth[length=2mm]}, every node/.style={font=\scriptsize}]
  \node[draw, rounded corners=2pt, align=center, minimum width=3.1cm, minimum height=0.8cm, fill=black!4] at (0.0,-0.0) {\textbf{fit\_nrtl}\\[-1pt]{\tiny fitParameters}};
  \node[draw, rounded corners=2pt, align=center, minimum width=3.1cm, minimum height=0.8cm, fill=black!4] at (3.6,-0.0) {\textbf{fit\_nrtl\_start2}\\[-1pt]{\tiny fitParameters}};
  \node[draw, rounded corners=2pt, align=center, minimum width=3.1cm, minimum height=0.8cm, fill=black!4] at (7.2,-0.0) {\textbf{txy\_ideal}\\[-1pt]{\tiny propertyScan1D}};
  \node[draw, rounded corners=2pt, align=center, minimum width=3.1cm, minimum height=0.8cm, fill=black!4] at (0.0,-1.3) {\textbf{txy\_unifac}\\[-1pt]{\tiny propertyScan1D}};
  \node[draw, rounded corners=2pt, align=center, minimum width=3.1cm, minimum height=0.8cm, fill=black!4] at (3.6,-1.3) {\textbf{txy\_nrtl}\\[-1pt]{\tiny propertyScan1D}};
\end{tikzpicture}
\end{tutscheme}
\tutwhat{NRTL nButanol-nButylAcetate regressed by the curation pipeline to every laboratory's measured bubble temperatures (ThermoML locates, the articles source), accepted at max |dT| < 1 K}
\begin{tutnarrative}
AN NRTL PAIR FROM EVERY LABORATORY'S MEASUREMENTS.

WRITTEN BY bin/choupo-curate pair nButanol nButylAcetate (DEV.md 4c C284, the bulk run; record docs/design/binary-pairs-from-open-measurements.md section 11).

THE DECLARATION -- WRITTEN AND COMMITTED BEFORE ANY FIT RAN (several declarations of the bulk run share one commit).

IDENTITY. nButanol (CAS 71-36-3, InChIKey LRHPLDYGYMQRHN-UHFFFAOYSA-N) + nButylAcetate (CAS 123-86-4, InChIKey DKPFZGUDAPQIHT-UHFFFAOYSA-N), joined to the archive through identity.json -- never by name. CLASS (C284): both records are volatile molecular liquids (no role nonvolatile/solute, no dissociatesTo, no charge, no solid, a vapour-pressure model), and every measured temperature lies below both critical temperatures.

NO MEASURED SERIES IN THIS TREE. constant/archiveDatasets names each dataset (DOI, NIST URL, article sha256, extraction selector, extract sha256); `bin/choupo-curate fetch` writes the points into constant/archiveExtracts/, which is gitignored.

STUDIES: every bubble-temperature block of the pair at or below 110 kPa with a declared uncertainty (2 block(s)). FITTED, ALL OF THEM (2): 10.1021/je700411p block 9 19 pts 101.3 kPa 389.54-397.48 K U 0.17-0.17 K 10.1021/je5000296 block 2 9 pts 101.3 kPa 390.20-395.56 K U 0.18-0.26 K SET ASIDE BY THE EXTRACTION RULES, BY NAME: (none)

THE RULE (DEV.md 4c C283 amended, C284). Every laboratory above is fitted; none is held out (the C16 partition and its band are retired as the acceptance test). NRTL, a\_ij = a\_ji = 0, alpha = 0.30 fixed: two parameters (b\_ij, b\_ji); each point weighted by 1/U, U its own stated expanded uncertainty. Objective twoStage (least squares, then minimax; T\_bubble absolute, in K). OUTLIERS: after a least-squares stage over every point, a point with |T\_model - T| > 3 K is discarded BY NAME (`discard` on its evidence entry, listed below) and the fit repeated without it. ACCEPTED when max |T\_model - T| over the KEPT points is below 1 K. C283's 5 \% relative clause judges composition; this fit has no composition residual, so only the temperature clause applies. A pair that splits ONE liquid inside its validity span is not promoted (arm f). Two starts: (0, 0) and (800, -400).

OUTLIERS DISCARDED AFTER THE LEAST-SQUARES STAGE: @@OUTLIERS-BEGIN@@ none: every one of the 28 points lies within 3 K after the least-squares stage. @@OUTLIERS-END@@

@@RESULT-BEGIN@@ WHAT THE RUN FOUND (pinned in `expected`; measured, not argued). fit\_nrtl: b\_ij = 535.9 K, b\_ji = -226.7 K over 28 kept points of 2 laboratories (0 discarded): max |dT| 0.581 K, mean 0.235 K against the acceptance 1 K $\to$ accepted. fit\_nrtl\_start2 lands on the same pair (535.9, -226.7). PROMOTED on 2026-11-10 into data/standards/parameters/NRTL/nButanol-nButylAcetate.dat as reviewStatus autoPromoted, humanReviewed false, because: max |dT| 0.581 K < 1 K over the 28 kept points, one liquid across the validity span @@RESULT-END@@

WHAT IS NOT CLAIMED. The pair is correct at these pressures and in this frame (ideal-gas vapour, the catalogue's own vapour-pressure records); the measurements were transcribed by the archive and have NOT been read back against the articles (the CHECKLIST that `fetch` writes is the aid).

\end{tutnarrative}
\begin{tutgolden}
\goldrow{diag}{fit\_nrtl}{bubble\_aad\_fit}{0.2347592207}
\goldrow{diag}{fit\_nrtl}{bubble\_max\_fit}{0.581109565516}
\goldrow{diag}{fit\_nrtl}{chi2}{66.0960534393}
\goldrow{diag}{fit\_nrtl}{chi2\_in\_sample}{66.0960534393}
\goldrow{diag}{fit\_nrtl}{chi2\_reduced}{2.54215590151}
\goldrow{diag}{fit\_nrtl}{cond\_JtJ}{84.5764936285}
\goldrow{diag}{fit\_nrtl}{converged}{1}
\goldrow{diag}{fit\_nrtl}{corr\_computed}{1}
\goldrow{diag}{fit\_nrtl}{dof}{26}
\goldrow{diag}{fit\_nrtl}{identifiable}{1}
\goldrow{diag}{fit\_nrtl}{invertible}{1}
\goldrow{diag}{fit\_nrtl}{iter}{11}
\goldrow{diag}{fit\_nrtl}{max\_abs\_corr}{0.994140104865}
\goldrow{diag}{fit\_nrtl}{max\_abs\_resid}{0.581109565516}
\goldmore{71}
\end{tutgolden}

\clearpage
\tutentry{curate14\_nrtl\_butanol\_water}{tutorials/props/curation/curate14\_nrtl\_butanol\_water}
\begin{tutscheme}
\begin{tikzpicture}[>={Stealth[length=2mm]}, every node/.style={font=\scriptsize}]
  \node[draw, rounded corners=2pt, align=center, minimum width=3.1cm, minimum height=0.8cm, fill=black!4] at (0.0,-0.0) {\textbf{fit\_binodal}\\[-1pt]{\tiny fitParameters}};
  \node[draw, rounded corners=2pt, align=center, minimum width=3.1cm, minimum height=0.8cm, fill=black!4] at (3.6,-0.0) {\textbf{fit\_binodal\_start2}\\[-1pt]{\tiny fitParameters}};
  \node[draw, rounded corners=2pt, align=center, minimum width=3.1cm, minimum height=0.8cm, fill=black!4] at (7.2,-0.0) {\textbf{report\_homogeneous\_bubble}\\[-1pt]{\tiny fitParameters}};
  \node[draw, rounded corners=2pt, align=center, minimum width=3.1cm, minimum height=0.8cm, fill=black!4] at (0.0,-1.3) {\textbf{fit\_twoStage}\\[-1pt]{\tiny fitParameters}};
\end{tikzpicture}
\end{tutscheme}
\tutwhat{Water + 1-butanol: NRTL fitted to measured ThermoML bubble temperatures AND tie lines in ONE fit (kind T\_bubble+LLE), held out against other studies, with leave-one-study-out triage}
\begin{tutnarrative}
THE REAL WATER + 1-BUTANOL PAIR: BUBBLE TEMPERATURES AND MEASURED TIE LINES IN ONE FIT (DEV.md 4c C16, C225, C226).

The synthetic witness curate13 proves that the machinery recovers a pair it was given. This case runs the same machinery (`kind T\_bubble+LLE`) on MEASURED data from the NIST/TRC ThermoML Archive (doi 10.18434/mds2-2422, the copy on the curator's machine, sha256-verified), and asks the one question the synthetic case cannot: does an NRTL pair hold the gap the measurements show, its change with temperature AND the boiling liquid, and does it predict studies it was never handed?

WHAT THE ARCHIVE HOLDS FOR THIS PAIR (found by reading the XML, not by the indexer). No file of the archive is a pure binary water + 1-butanol study. Every datum below is a BINARY BLOCK inside a ternary or quaternary study (the authors report the binary edge of their system), so: (a) bin/choupo-thermoml vle-index finds ONE usable isobaric bubble series (doi 10.1021/je700411p, 15 interior points) and the heteroazeotrope temperatures of two more; (b) the liquid-liquid blocks are not reached by `extract-vle`, so the tie-line files were written by a small curator script (kept out of the tree) that reads the same XML, pairs the two liquids of each temperature, and converts mass fractions to mole fractions (M 74.123 and 18.015 g/mol, said in the file). Whether a given binary row was measured by its authors or taken from the literature is not recorded in the archive; the files say `reviewStatus transcribedNotCheckedAgainstArticle` and the run announces it.

THE PROMOTED OPERATION (C280, 2026-11-10; the sections below it are the negative results that led here, kept short). `fit\_binodal` fits the BINODAL and the HETEROAZEOTROPE temperature, not the homogeneous butanol-rich bubble series: least squares (declared: the water-rich tie lines carry a 15 \% scatter between laboratories, measured about a quadratic in T as mean 5.5 \%, max 18.6 \%, and a minimax objective equalises the largest error of both liquids at that scatter), binodal rows relative. FIT bubble: 10.1016/j.jct.2014.07.021 (2 heteroazeotrope points); tie lines: that study and 10.1021/acs.jced.7b00302, je5007218, j.fluid.2007.12.003, acs.jced.6b00941 -- 14 tie lines, 288-366 K. HELD OUT bubble: 10.1016/j.jct.2016.04.013 (2 heteroazeotrope points); 29 tie lines of 21 other studies including 10.1021/je700411p. The homogeneous bubble series of je700411p (15 points, x 0.56-0.99, 366-390 K) is scored by `report\_homogeneous\_bubble`, the SAME fit against that series only, so its miss is its own line and does not enter the heteroazeotrope verdict. (The mixed kind needed bubble evidence in the fit: the two heteroazeotrope points. Whether a fit with none is allowed was not tried.) RESULT (the engine's ERROR REPORT; five parameters b\_ij 1894.0429, b\_ji -163.8364, a\_ij -4.7108, a\_ji 4.0996, alpha 0.4459): VERDICT validated -- held-out tie-line AAD 0.0056 against 0.01, held-out heteroazeotrope bubble 0.376 K = 0.103 \% against 0.23 \%. Relative error mean / max, fit | held out: water-rich 6.9 \% / 17.3 \% | 5.2 \% / 16.6 \%; butanol-rich 1.9 \% / 5.8 \% | 2.2 \% / 11.5 \%. Solubility of butanol in water, model: 0.0194 (288 K), 0.0175 (320 K), 0.0173 (340 K), 0.0183 (366 K) -- the shallow minimum is drawn. fit\_binodal\_start2 lands on the same pair (b\_ij 1894.08). `fit\_twoStage` (the fitter's default objective on this partition): b\_ij 2132.88; rich-liquid relative 5.6 \% mean, 15 \% max, held-out tie-line AAD 0.0143 -- notValidated: minimax follows the water-rich scatter. WHAT THE PAIR IS NOT FOR: `report\_homogeneous\_bubble` scores the homogeneous butanol-rich series held out: mean 2.18 K = 0.59 \% (max 3.2 K) against the 0.23 \% band, NOT loosened -- notValidated, as found. The pair is for the liquid-liquid split and the heteroazeotrope; at 1 atm on the butanol-rich homogeneous side it reads about 2-3 K high. The catalogue record says so. The explicit data of this section replace the "WHAT THE RUN FOUND" blocks below for the promoted record; those describe the earlier partition (je700411p in the fit, loso\_$\ast$ operations since dropped).

THE FIRST PARTITION (superseded), DECLARED BEFORE THE FIRST FIT. FIT bubble temperatures: doi 10.1021/je700411p (15 points, 101.3 kPa); tie lines: the five studies that report THREE OR MORE temperatures -- 10.1021/je700411p (5), 10.1021/acs.jced.7b00302 (3), 10.1021/je5007218 (3), 10.1016/j.fluid.2007.12.003 (3), 10.1021/acs.jced.6b00941 (3): 17 tie lines, 288-362 K. HELD OUT bubble temperatures: 10.1016/j.jct.2014.07.021 and 10.1016/j.jct.2016.04.013 (two heteroazeotrope temperatures each, about 365.8 K); tie lines: the eighteen other studies, one or two temperatures each (21 tie lines, 288-366 K). No author name of a held-out study appears in a fit study (the archive carries no affiliations, so a shared author is the dependence the rule can check). BANDS written in `acceptance` below with their origin (maxAAD 0.23 \% of T, maxLLE 0.01 in mole fraction). TRIAGE `loso\_$\ast$` fit the pair with ONE tie-line study moved from fit to held-out in turn, to see whether the verdict is the property of the pair or of one study.

OBJECTIVE (C280). The operations below run the DEFAULT objective of `fitParameters`, twoStage, with tie-line rows relative; read the run header for the one that produced a number. The weights stated under THE FIT are the pre-C280 ones: the tie-line weight is now `sigmaRelative`, below. THE FIT. NRTL with tau\_ij = a\_ij + b\_ij/T and alpha free: FIVE parameters (b\_ij, b\_ji, a\_ij, a\_ji, alpha), i = nButanol, j = water. Bubble rows are weighted by 1/U (each point's own stated expanded uncertainty), tie-line mole fractions by 1/0.0005. The start (b\_ij 0, b\_ji 1000 K, a 0, alpha 0.30) splits the liquid at every measured tie line; `fit\_mixed\_start2` starts elsewhere (400 / 1400 / alpha 0.40) and must land on the same pair.

THE FIRST FORM, AND WHY IT WAS LEFT (kept as the record of a negative result). With a\_ij = a\_ji = 0, alpha = 0.30 fixed and the tie lines weighted by 1/0.005, the fit gave b\_ij = -39.14 K, b\_ji = 1178.21 K: held-out bubble AAD 1.72 K, held-out tie-line AAD 0.0165, and a butanol solubility in water that ROSE from 0.019 (288 K) to 0.067 (366 K) where the measurements keep it flat (0.016-0.021). Two b's and one alpha cannot draw a flat solubility line, and a weight of 0.005 on a mole fraction of 0.02 cannot see it. Both were changed, each with its reason in `residual` and `acceptance` below; the tie-line band (0.01) was not.

WHAT THE ABSOLUTE LEAST-SQUARES RUN FOUND -- the form this case had before C280 (measured, not argued; reproduced 2026-11-10 by `objective \{ method leastSquares; scale \{ binodal absolute; \} \}` with `liquidLiquid \{ sigma 0.0005; \}`, the numbers below to the last digit). VALIDATED under the bands below. It is NO LONGER what `expected` pins: the case now runs the default objective, and the section "WHAT THE MINIMAX STAGE FOUND" below says what that gave. fit\_mixed: b\_ij = 1486.66 K, b\_ji = -27.30 K, a\_ij = -3.6282, a\_ji = 3.6295, alpha = 0.4428; fit\_mixed\_start2 lands on the same five numbers. The four tau parameters are strongly correlated (max|corr| 0.998, cond JtJ 3.5e4): the pair is ONE set, to be used whole, and no single number of it carries a meaning. Solubility of butanol in water (mole fraction, model): 0.0188 at 288 K, 0.0186 at 300 K, 0.0189 at 320 K, 0.0197 at 340 K, 0.0215 at 366 K -- flat, as measured. HELD OUT (21 tie lines, 18 studies, all predicted as two liquids): AAD 0.0061 against 0.01; water-rich liquid AAD 0.0010 (max 0.0027), butanol-rich liquid AAD 0.0113 (max 0.060, recomputed outside the engine). The largest miss is one measurement, x = 0.550 at 298.15 K (10.1016/j.fluid.2014.07.011), where two other studies report 0.487 and 0.488 at the same temperature and the model gives 0.490. In the fit set the largest tie-line miss is 0.021. Bubble temperatures HELD OUT (4 heteroazeotrope points, two studies): AAD 0.53 K = 0.15 \% against 0.23 \%. Leave-one-study-out (4 of the 5 tie-line fit studies; the fifth, 10.1021/je700411p, is the only bubble series and the engine refuses its archive file in both roles): alpha 0.4385-0.4456, held-out tie-line AAD 0.0029-0.0075, held-out bubble AAD 0.54-0.59 K; every verdict is the same. The result is the pair's against these data, not one study's. IN RELATIVE TERMS (recomputed outside the engine), which is how a composition is judged: butanol-rich liquid, mean 1.1 \% (max 4.0 \%) in the fit set and 2.3 \% (max 10.9 \%, the one 298.15 K point above) held out; water-rich liquid, mean 12 \% (max 24 \%) in the fit set and 5.5 \% (max 17 \%) held out. The studies themselves disagree by that much on the water-rich liquid (0.0159 to 0.0207 at 298.15 K, 0.0170 to 0.0220 at 303 K), and below 320 K the model runs through the middle of them. ABOVE 320 K IT IS SYSTEMATICALLY HIGH, 0.019-0.021 against 0.016-0.019: the shallow minimum the measurements show near 333 K is not drawn. WHAT THE MINIMAX STAGE FOUND (C280; the default objective, twoStage, binodal rows RELATIVE with `sigmaRelative 0.025` -- the absolute 0.0005 expressed at the water-rich liquid's mole fraction of 0.02 -- and bubble rows absolute in kelvin at 1/U; 49 rows: 15 bubble, 34 tie-line; 27 s for the six operations). The practice Vitor asked for (2026-11-10): least squares for the first estimate, then minimise the LARGEST scaled error. VERDICT: notValidated, reported as found and with no band, weight or bound changed to reach another one: held-out tie-line AAD 0.0186 against the 0.01 band (bubble AAD 0.21 \% = 0.77 K, inside 0.23 \%). Relative mean / max error, fit set | held-out set: fit set held out mean max mean max water-rich liquid 6.7 \% 15.7 \% 6.7 \% 19.2 \% butanol-rich liquid 5.7 \% 15.7 \% 7.8 \% 17.2 \% bubble T (absolute) 0.60 K 1.07 K 0.77 K 0.96 K against the pre-C280 pair: water-rich 12 \% / 24 \% fit and 5.4 \% / 17 \% held out; butanol-rich 1.1 \% / 4.0 \% and 2.3 \% / 10.9 \%; bubble T fit mean 2.1 K, max 3.1 K. Solubility of butanol in water, model (mole fraction), measured 0.016-0.021 below 320 K and 0.016-0.019 above: 288 K 320 K 340 K 366 K pre-C280 0.0188 0.0189 0.0197 0.0215 (rises above 320 K) stage 1 (LS) 0.0191 0.0174 0.0171 0.0174 minimax 0.0197 0.0179 0.0175 0.0176 (flat, in the band) (outside the engine: NRTL, equal activities, from the printed parameters.) What it means: the systematic 15 \% excess of the water-rich liquid above 320 K is gone, and the largest relative miss of the fit set fell from 24 \% to 15.7 \%. The price is on the OTHER liquid: the relative scale lets a mole fraction of 0.5 miss by 15 \% -- 0.07 -- as the 0.02 one does, so the butanol-rich side went from AAD 0.011 to 0.036 held out and the verdict flipped on the binodal band, which is absolute. The stage-1 (relative least-squares) pair alone gives held-out tie-line AAD 0.0136, also above 0.01. Three things the run also showed: (a) the minimax objective equalises the relative errors of both liquids at the scatter level of the water-rich measurements (15 \%, the between-study scatter), which the butanol-rich liquid does not have; a per-property weight, or a bound below the data's own scatter, is the lever -- not changed here; (b) the maximum is not smooth, and `fit\_mixed\_start2` ends at a different point (b\_ij 1790 against 1680) where stage 1 alone lands on the same pair from both starts (1684.63 / 1684.64); (c) the correlation is 1.000 at the reported point, so the four tau numbers are still one set. The pair is NOT promoted. Independence: je700411p supplies the bubble series AND five tie lines of the fit (the run announces "one publication, counted twice"); the five held-out heteroazeotrope temperatures are three values (365.7-366.2 K).

WHAT IS NOT CLAIMED. The pair lives in this frame: ideal-gas vapour, the catalogue's Antoine records (water 0.70 K low at 1 atm, C226 item 7), one pressure for the bubble data (101.3 kPa), LLE data at about atmospheric pressure. Nothing here tests it at another pressure or with another vapour model.

\end{tutnarrative}
\begin{tutgolden}
\goldrow{diag}{fit\_binodal}{aad\_heldout\_K}{0.376320388878}
\goldrow{diag}{fit\_binodal}{aad\_heldout\_pct}{0.102898498545}
\goldrow{diag}{fit\_binodal}{bubble\_aad\_fit}{0.245855038585}
\goldrow{diag}{fit\_binodal}{bubble\_aad\_heldout}{0.376320388878}
\goldrow{diag}{fit\_binodal}{bubble\_max\_fit}{0.375302576021}
\goldrow{diag}{fit\_binodal}{bubble\_max\_heldout}{0.506233276606}
\goldrow{diag}{fit\_binodal}{chi2}{186.974832041}
\goldrow{diag}{fit\_binodal}{chi2\_in\_sample}{186.974832041}
\goldrow{diag}{fit\_binodal}{chi2\_reduced}{7.47899328162}
\goldrow{diag}{fit\_binodal}{cond\_JtJ}{291.796362448}
\goldrow{diag}{fit\_binodal}{converged}{1}
\goldrow{diag}{fit\_binodal}{corr\_computed}{1}
\goldrow{diag}{fit\_binodal}{dof}{25}
\goldrow{diag}{fit\_binodal}{identifiable}{1}
\goldmore{163}
\end{tutgolden}

\clearpage
\tutentry{curate150\_nrtl\_acetonitrile\_isopropanol}{tutorials/props/curation/curate150\_nrtl\_acetonitrile\_isopropanol}
\begin{tutscheme}
\begin{tikzpicture}[>={Stealth[length=2mm]}, every node/.style={font=\scriptsize}]
  \node[draw, rounded corners=2pt, align=center, minimum width=3.1cm, minimum height=0.8cm, fill=black!4] at (0.0,-0.0) {\textbf{fit\_nrtl}\\[-1pt]{\tiny fitParameters}};
  \node[draw, rounded corners=2pt, align=center, minimum width=3.1cm, minimum height=0.8cm, fill=black!4] at (3.6,-0.0) {\textbf{fit\_nrtl\_start2}\\[-1pt]{\tiny fitParameters}};
  \node[draw, rounded corners=2pt, align=center, minimum width=3.1cm, minimum height=0.8cm, fill=black!4] at (7.2,-0.0) {\textbf{txy\_ideal}\\[-1pt]{\tiny propertyScan1D}};
  \node[draw, rounded corners=2pt, align=center, minimum width=3.1cm, minimum height=0.8cm, fill=black!4] at (0.0,-1.3) {\textbf{txy\_nrtl}\\[-1pt]{\tiny propertyScan1D}};
\end{tikzpicture}
\end{tutscheme}
\tutwhat{NRTL acetonitrile-isopropanol regressed by the curation pipeline to every laboratory's measured bubble temperatures (ThermoML locates, the articles source), accepted at max |dT| < 1 K}
\begin{tutnarrative}
AN NRTL PAIR FROM EVERY LABORATORY'S MEASUREMENTS.

WRITTEN BY bin/choupo-curate pair acetonitrile isopropanol (DEV.md 4c C284, the bulk run; record docs/design/binary-pairs-from-open-measurements.md section 11).

THE DECLARATION -- WRITTEN AND COMMITTED BEFORE ANY FIT RAN (several declarations of the bulk run share one commit).

IDENTITY. acetonitrile (CAS 75-05-8, InChIKey WEVYAHXRMPXWCK-UHFFFAOYSA-N) + isopropanol (CAS 67-63-0, InChIKey KFZMGEQAYNKOFK-UHFFFAOYSA-N), joined to the archive through identity.json -- never by name. CLASS (C284): both records are volatile molecular liquids (no role nonvolatile/solute, no dissociatesTo, no charge, no solid, a vapour-pressure model), and every measured temperature lies below both critical temperatures.

NO MEASURED SERIES IN THIS TREE. constant/archiveDatasets names each dataset (DOI, NIST URL, article sha256, extraction selector, extract sha256); `bin/choupo-curate fetch` writes the points into constant/archiveExtracts/, which is gitignored.

STUDIES: every bubble-temperature block of the pair at or below 110 kPa with a declared uncertainty (2 block(s)). FITTED, ALL OF THEM (2): 10.1016/j.jct.2007.03.007 block 24 14 pts 95.5 kPa 350.35-353.15 K U 0.69-0.69 K 10.1016/j.fluid.2015.10.035 block 6 13 pts 101.3 kPa 348.40-353.03 K U 0.13-0.18 K (reports x of the other member: written as x[acetonitrile] = 1 - x by --x-of) SET ASIDE BY THE EXTRACTION RULES, BY NAME: (none)

THE RULE (DEV.md 4c C283 amended, C284). Every laboratory above is fitted; none is held out (the C16 partition and its band are retired as the acceptance test). NRTL, a\_ij = a\_ji = 0, alpha = 0.30 fixed: two parameters (b\_ij, b\_ji); each point weighted by 1/U, U its own stated expanded uncertainty. Objective twoStage (least squares, then minimax; T\_bubble absolute, in K). OUTLIERS: after a least-squares stage over every point, a point with |T\_model - T| > 3 K is discarded BY NAME (`discard` on its evidence entry, listed below) and the fit repeated without it. ACCEPTED when max |T\_model - T| over the KEPT points is below 1 K. C283's 5 \% relative clause judges composition; this fit has no composition residual, so only the temperature clause applies. A pair that splits ONE liquid inside its validity span is not promoted (arm f). Two starts: (0, 0) and (800, -400).

OUTLIERS DISCARDED AFTER THE LEAST-SQUARES STAGE: @@OUTLIERS-BEGIN@@ 7 point(s) of 27: doi 10.1016/j.jct.2007.03.007 row 3: x 0.2444, P 95.500 kPa, T 351.25 K, dT -3.02 K; doi 10.1016/j.jct.2007.03.007 row 4: x 0.3264, P 95.500 kPa, T 350.85 K, dT -3.29 K; doi 10.1016/j.jct.2007.03.007 row 5: x 0.3928, P 95.500 kPa, T 350.55 K, dT -3.33 K; doi 10.1016/j.jct.2007.03.007 row 6: x 0.4471, P 95.500 kPa, T 350.45 K, dT -3.40 K; doi 10.1016/j.jct.2007.03.007 row 7: x 0.5310, P 95.500 kPa, T 350.35 K, dT -3.40 K; doi 10.1016/j.jct.2007.03.007 row 8: x 0.6859, P 95.500 kPa, T 350.45 K, dT -3.16 K; doi 10.1016/j.jct.2007.03.007 row 9: x 0.7238, P 95.500 kPa, T 350.55 K, dT -3.04 K @@OUTLIERS-END@@

@@RESULT-BEGIN@@ WHAT THE RUN FOUND (pinned in `expected`; measured, not argued). fit\_nrtl: b\_ij = 1.6 K, b\_ji = 378.9 K over 20 kept points of 2 laboratories (7 discarded): max |dT| 2.597 K, mean 0.812 K against the acceptance 1 K $\to$ notAccepted. fit\_nrtl\_start2 lands on the same pair (1.6, 378.9). NOT PROMOTABLE: max |dT| 2.597 K over the 20 kept points is not below 1 K @@RESULT-END@@

WHAT IS NOT CLAIMED. The pair is correct at these pressures and in this frame (ideal-gas vapour, the catalogue's own vapour-pressure records); the measurements were transcribed by the archive and have NOT been read back against the articles (the CHECKLIST that `fetch` writes is the aid).

\end{tutnarrative}
\begin{tutgolden}
\goldrow{diag}{fit\_nrtl}{bubble\_aad\_fit}{0.811533490562}
\goldrow{diag}{fit\_nrtl}{bubble\_max\_fit}{2.59694463213}
\goldrow{diag}{fit\_nrtl}{chi2}{133.00846312}
\goldrow{diag}{fit\_nrtl}{chi2\_in\_sample}{133.00846312}
\goldrow{diag}{fit\_nrtl}{chi2\_reduced}{7.3893590622}
\goldrow{diag}{fit\_nrtl}{cond\_JtJ}{157.230962717}
\goldrow{diag}{fit\_nrtl}{converged}{1}
\goldrow{diag}{fit\_nrtl}{corr\_computed}{1}
\goldrow{diag}{fit\_nrtl}{dof}{18}
\goldrow{diag}{fit\_nrtl}{identifiable}{1}
\goldrow{diag}{fit\_nrtl}{invertible}{1}
\goldrow{diag}{fit\_nrtl}{iter}{5}
\goldrow{diag}{fit\_nrtl}{max\_abs\_corr}{0.993292481119}
\goldrow{diag}{fit\_nrtl}{max\_abs\_resid}{2.59694463213}
\goldmore{54}
\end{tutgolden}

\clearpage
\tutentry{curate151\_nrtl\_nNDimethylformamide\_aniline}{tutorials/props/curation/curate151\_nrtl\_nNDimethylformamide\_aniline}
\begin{tutscheme}
\begin{tikzpicture}[>={Stealth[length=2mm]}, every node/.style={font=\scriptsize}]
  \node[draw, rounded corners=2pt, align=center, minimum width=3.1cm, minimum height=0.8cm, fill=black!4] at (0.0,-0.0) {\textbf{fit\_nrtl}\\[-1pt]{\tiny fitParameters}};
  \node[draw, rounded corners=2pt, align=center, minimum width=3.1cm, minimum height=0.8cm, fill=black!4] at (3.6,-0.0) {\textbf{fit\_nrtl\_start2}\\[-1pt]{\tiny fitParameters}};
  \node[draw, rounded corners=2pt, align=center, minimum width=3.1cm, minimum height=0.8cm, fill=black!4] at (7.2,-0.0) {\textbf{txy\_ideal}\\[-1pt]{\tiny propertyScan1D}};
  \node[draw, rounded corners=2pt, align=center, minimum width=3.1cm, minimum height=0.8cm, fill=black!4] at (0.0,-1.3) {\textbf{txy\_nrtl}\\[-1pt]{\tiny propertyScan1D}};
\end{tikzpicture}
\end{tutscheme}
\tutwhat{NRTL nNDimethylformamide-aniline regressed by the curation pipeline to every laboratory's measured bubble temperatures (ThermoML locates, the articles source), accepted at max |dT| < 1 K}
\begin{tutnarrative}
AN NRTL PAIR FROM EVERY LABORATORY'S MEASUREMENTS.

WRITTEN BY bin/choupo-curate pair nNDimethylformamide aniline (DEV.md 4c C284, the bulk run; record docs/design/binary-pairs-from-open-measurements.md section 11).

THE DECLARATION -- WRITTEN AND COMMITTED BEFORE ANY FIT RAN (several declarations of the bulk run share one commit).

IDENTITY. nNDimethylformamide (CAS 68-12-2, InChIKey ZMXDDKWLCZADIW-UHFFFAOYSA-N) + aniline (CAS 62-53-3, InChIKey PAYRUJLWNCNPSJ-UHFFFAOYSA-N), joined to the archive through identity.json -- never by name. CLASS (C284): both records are volatile molecular liquids (no role nonvolatile/solute, no dissociatesTo, no charge, no solid, a vapour-pressure model), and every measured temperature lies below both critical temperatures.

NO MEASURED SERIES IN THIS TREE. constant/archiveDatasets names each dataset (DOI, NIST URL, article sha256, extraction selector, extract sha256); `bin/choupo-curate fetch` writes the points into constant/archiveExtracts/, which is gitignored.

STUDIES: every bubble-temperature block of the pair at or below 110 kPa with a declared uncertainty (2 block(s)). FITTED, ALL OF THEM (2): 10.1016/j.fluid.2008.07.018 block 20 16 pts 95.5 kPa 428.85-454.35 K U 0.69-0.69 K 10.1021/acs.jced.6b00700 block 12 11 pts 96.15 kPa 425.50-454.60 K U 0.2-0.2 K SET ASIDE BY THE EXTRACTION RULES, BY NAME: (none)

THE RULE (DEV.md 4c C283 amended, C284). Every laboratory above is fitted; none is held out (the C16 partition and its band are retired as the acceptance test). NRTL, a\_ij = a\_ji = 0, alpha = 0.30 fixed: two parameters (b\_ij, b\_ji); each point weighted by 1/U, U its own stated expanded uncertainty. Objective twoStage (least squares, then minimax; T\_bubble absolute, in K). OUTLIERS: after a least-squares stage over every point, a point with |T\_model - T| > 3 K is discarded BY NAME (`discard` on its evidence entry, listed below) and the fit repeated without it. ACCEPTED when max |T\_model - T| over the KEPT points is below 1 K. C283's 5 \% relative clause judges composition; this fit has no composition residual, so only the temperature clause applies. A pair that splits ONE liquid inside its validity span is not promoted (arm f). Two starts: (0, 0) and (800, -400).

OUTLIERS DISCARDED AFTER THE LEAST-SQUARES STAGE: @@OUTLIERS-BEGIN@@ none: every one of the 27 points lies within 3 K after the least-squares stage. @@OUTLIERS-END@@

@@RESULT-BEGIN@@ WHAT THE RUN FOUND (pinned in `expected`; measured, not argued). fit\_nrtl: b\_ij = 580.2 K, b\_ji = -629.6 K over 27 kept points of 2 laboratories (0 discarded): max |dT| 1.993 K, mean 0.960 K against the acceptance 1 K $\to$ notAccepted. fit\_nrtl\_start2 lands on the same pair (580.3, -629.7). NOT PROMOTABLE: max |dT| 1.993 K over the 27 kept points is not below 1 K @@RESULT-END@@

WHAT IS NOT CLAIMED. The pair is correct at these pressures and in this frame (ideal-gas vapour, the catalogue's own vapour-pressure records); the measurements were transcribed by the archive and have NOT been read back against the articles (the CHECKLIST that `fetch` writes is the aid).

\end{tutnarrative}
\begin{tutgolden}
\goldrow{diag}{fit\_nrtl}{bubble\_aad\_fit}{0.960253664069}
\goldrow{diag}{fit\_nrtl}{bubble\_max\_fit}{1.99312403092}
\goldrow{diag}{fit\_nrtl}{chi2}{251.141835067}
\goldrow{diag}{fit\_nrtl}{chi2\_in\_sample}{251.141835067}
\goldrow{diag}{fit\_nrtl}{chi2\_reduced}{10.0456734027}
\goldrow{diag}{fit\_nrtl}{cond\_JtJ}{6371.22053905}
\goldrow{diag}{fit\_nrtl}{converged}{1}
\goldrow{diag}{fit\_nrtl}{corr\_computed}{1}
\goldrow{diag}{fit\_nrtl}{dof}{25}
\goldrow{diag}{fit\_nrtl}{identifiable}{0}
\goldrow{diag}{fit\_nrtl}{invertible}{1}
\goldrow{diag}{fit\_nrtl}{iter}{13}
\goldrow{diag}{fit\_nrtl}{max\_abs\_corr}{0.999921531338}
\goldrow{diag}{fit\_nrtl}{max\_abs\_resid}{1.99312403092}
\goldmore{60}
\end{tutgolden}

\clearpage
\tutentry{curate152\_nrtl\_tetrahydrofuran\_methanol}{tutorials/props/curation/curate152\_nrtl\_tetrahydrofuran\_methanol}
\begin{tutscheme}
\begin{tikzpicture}[>={Stealth[length=2mm]}, every node/.style={font=\scriptsize}]
  \node[draw, rounded corners=2pt, align=center, minimum width=3.1cm, minimum height=0.8cm, fill=black!4] at (0.0,-0.0) {\textbf{fit\_nrtl}\\[-1pt]{\tiny fitParameters}};
  \node[draw, rounded corners=2pt, align=center, minimum width=3.1cm, minimum height=0.8cm, fill=black!4] at (3.6,-0.0) {\textbf{fit\_nrtl\_start2}\\[-1pt]{\tiny fitParameters}};
  \node[draw, rounded corners=2pt, align=center, minimum width=3.1cm, minimum height=0.8cm, fill=black!4] at (7.2,-0.0) {\textbf{txy\_ideal}\\[-1pt]{\tiny propertyScan1D}};
  \node[draw, rounded corners=2pt, align=center, minimum width=3.1cm, minimum height=0.8cm, fill=black!4] at (0.0,-1.3) {\textbf{txy\_nrtl}\\[-1pt]{\tiny propertyScan1D}};
\end{tikzpicture}
\end{tutscheme}
\tutwhat{NRTL tetrahydrofuran-methanol regressed by the curation pipeline to every laboratory's measured bubble temperatures (ThermoML locates, the articles source), accepted at max |dT| < 1 K}
\begin{tutnarrative}
AN NRTL PAIR FROM EVERY LABORATORY'S MEASUREMENTS.

WRITTEN BY bin/choupo-curate pair tetrahydrofuran methanol (DEV.md 4c C284, the bulk run; record docs/design/binary-pairs-from-open-measurements.md section 11).

THE DECLARATION -- WRITTEN AND COMMITTED BEFORE ANY FIT RAN (several declarations of the bulk run share one commit).

IDENTITY. tetrahydrofuran (CAS 109-99-9, InChIKey WYURNTSHIVDZCO-UHFFFAOYSA-N) + methanol (CAS 67-56-1, InChIKey OKKJLVBELUTLKV-UHFFFAOYSA-N), joined to the archive through identity.json -- never by name. CLASS (C284): both records are volatile molecular liquids (no role nonvolatile/solute, no dissociatesTo, no charge, no solid, a vapour-pressure model), and every measured temperature lies below both critical temperatures.

NO MEASURED SERIES IN THIS TREE. constant/archiveDatasets names each dataset (DOI, NIST URL, article sha256, extraction selector, extract sha256); `bin/choupo-curate fetch` writes the points into constant/archiveExtracts/, which is gitignored.

STUDIES: every bubble-temperature block of the pair at or below 110 kPa with a declared uncertainty (2 block(s)). FITTED, ALL OF THEM (2): 10.1016/j.fluid.2013.09.060 block 0 15 pts 101.3 kPa 332.40-336.40 K U 0.2-0.2 K 10.1016/j.fluid.2016.04.012 block 2 14 pts 101.3 kPa 332.55-335.85 K U 0.21-0.24 K (reports x of the other member: written as x[tetrahydrofuran] = 1 - x by --x-of) SET ASIDE BY THE EXTRACTION RULES, BY NAME: (none)

THE RULE (DEV.md 4c C283 amended, C284). Every laboratory above is fitted; none is held out (the C16 partition and its band are retired as the acceptance test). NRTL, a\_ij = a\_ji = 0, alpha = 0.30 fixed: two parameters (b\_ij, b\_ji); each point weighted by 1/U, U its own stated expanded uncertainty. Objective twoStage (least squares, then minimax; T\_bubble absolute, in K). OUTLIERS: after a least-squares stage over every point, a point with |T\_model - T| > 3 K is discarded BY NAME (`discard` on its evidence entry, listed below) and the fit repeated without it. ACCEPTED when max |T\_model - T| over the KEPT points is below 1 K. C283's 5 \% relative clause judges composition; this fit has no composition residual, so only the temperature clause applies. A pair that splits ONE liquid inside its validity span is not promoted (arm f). Two starts: (0, 0) and (800, -400).

OUTLIERS DISCARDED AFTER THE LEAST-SQUARES STAGE: @@OUTLIERS-BEGIN@@ none: every one of the 29 points lies within 3 K after the least-squares stage. @@OUTLIERS-END@@

@@RESULT-BEGIN@@ WHAT THE RUN FOUND (pinned in `expected`; measured, not argued). fit\_nrtl: b\_ij = 241.4 K, b\_ji = 61.4 K over 29 kept points of 2 laboratories (0 discarded): max |dT| 0.496 K, mean 0.207 K against the acceptance 1 K $\to$ accepted. fit\_nrtl\_start2 lands on the same pair (241.4, 61.4). PROMOTED on 2026-11-10 into data/standards/parameters/NRTL/methanol-tetrahydrofuran.dat as reviewStatus autoPromoted, humanReviewed false, because: max |dT| 0.496 K < 1 K over the 29 kept points, one liquid across the validity span @@RESULT-END@@

WHAT IS NOT CLAIMED. The pair is correct at these pressures and in this frame (ideal-gas vapour, the catalogue's own vapour-pressure records); the measurements were transcribed by the archive and have NOT been read back against the articles (the CHECKLIST that `fetch` writes is the aid).

\end{tutnarrative}
\begin{tutgolden}
\goldrow{diag}{fit\_nrtl}{bubble\_aad\_fit}{0.207420678277}
\goldrow{diag}{fit\_nrtl}{bubble\_max\_fit}{0.496156123199}
\goldrow{diag}{fit\_nrtl}{chi2}{41.2743983275}
\goldrow{diag}{fit\_nrtl}{chi2\_in\_sample}{41.2743983275}
\goldrow{diag}{fit\_nrtl}{chi2\_reduced}{1.52868141954}
\goldrow{diag}{fit\_nrtl}{cond\_JtJ}{97.6608692047}
\goldrow{diag}{fit\_nrtl}{converged}{1}
\goldrow{diag}{fit\_nrtl}{corr\_computed}{1}
\goldrow{diag}{fit\_nrtl}{dof}{27}
\goldrow{diag}{fit\_nrtl}{identifiable}{1}
\goldrow{diag}{fit\_nrtl}{invertible}{1}
\goldrow{diag}{fit\_nrtl}{iter}{5}
\goldrow{diag}{fit\_nrtl}{max\_abs\_corr}{0.99491922753}
\goldrow{diag}{fit\_nrtl}{max\_abs\_resid}{0.496156123199}
\goldmore{62}
\end{tutgolden}

\clearpage
\tutentry{curate153\_nrtl\_diisopropylEther\_224Trimethylpentane}{tutorials/props/curation/curate153\_nrtl\_diisopropylEther\_224Trimethylpentane}
\begin{tutscheme}
\begin{tikzpicture}[>={Stealth[length=2mm]}, every node/.style={font=\scriptsize}]
  \node[draw, rounded corners=2pt, align=center, minimum width=3.1cm, minimum height=0.8cm, fill=black!4] at (0.0,-0.0) {\textbf{fit\_nrtl}\\[-1pt]{\tiny fitParameters}};
  \node[draw, rounded corners=2pt, align=center, minimum width=3.1cm, minimum height=0.8cm, fill=black!4] at (3.6,-0.0) {\textbf{fit\_nrtl\_start2}\\[-1pt]{\tiny fitParameters}};
  \node[draw, rounded corners=2pt, align=center, minimum width=3.1cm, minimum height=0.8cm, fill=black!4] at (7.2,-0.0) {\textbf{txy\_ideal}\\[-1pt]{\tiny propertyScan1D}};
  \node[draw, rounded corners=2pt, align=center, minimum width=3.1cm, minimum height=0.8cm, fill=black!4] at (0.0,-1.3) {\textbf{txy\_nrtl}\\[-1pt]{\tiny propertyScan1D}};
\end{tikzpicture}
\end{tutscheme}
\tutwhat{NRTL diisopropylEther-224Trimethylpentane regressed by the curation pipeline to every laboratory's measured bubble temperatures (ThermoML locates, the articles source), accepted at max |dT| < 1 K}
\begin{tutnarrative}
AN NRTL PAIR FROM EVERY LABORATORY'S MEASUREMENTS.

WRITTEN BY bin/choupo-curate pair diisopropylEther 224Trimethylpentane (DEV.md 4c C284, the bulk run; record docs/design/binary-pairs-from-open-measurements.md section 11).

THE DECLARATION -- WRITTEN AND COMMITTED BEFORE ANY FIT RAN (several declarations of the bulk run share one commit).

IDENTITY. diisopropylEther (CAS 108-20-3, InChIKey ZAFNJMIOTHYJRJ-UHFFFAOYSA-N) + 224Trimethylpentane (CAS 540-84-1, InChIKey NHTMVDHEPJAVLT-UHFFFAOYSA-N), joined to the archive through identity.json -- never by name. CLASS (C284): both records are volatile molecular liquids (no role nonvolatile/solute, no dissociatesTo, no charge, no solid, a vapour-pressure model), and every measured temperature lies below both critical temperatures.

NO MEASURED SERIES IN THIS TREE. constant/archiveDatasets names each dataset (DOI, NIST URL, article sha256, extraction selector, extract sha256); `bin/choupo-curate fetch` writes the points into constant/archiveExtracts/, which is gitignored.

STUDIES: every bubble-temperature block of the pair at or below 110 kPa with a declared uncertainty (2 block(s)). FITTED, ALL OF THEM (2): 10.1016/j.fluid.2006.08.004 block 11 18 pts 101.3 kPa 342.57-368.15 K U 0.18-0.3 K 10.1016/j.fluid.2004.12.007 block 9 10 pts 95.96 kPa 340.85-364.55 K U 0.31-0.31 K SET ASIDE BY THE EXTRACTION RULES, BY NAME: (none)

THE RULE (DEV.md 4c C283 amended, C284). Every laboratory above is fitted; none is held out (the C16 partition and its band are retired as the acceptance test). NRTL, a\_ij = a\_ji = 0, alpha = 0.30 fixed: two parameters (b\_ij, b\_ji); each point weighted by 1/U, U its own stated expanded uncertainty. Objective twoStage (least squares, then minimax; T\_bubble absolute, in K). OUTLIERS: after a least-squares stage over every point, a point with |T\_model - T| > 3 K is discarded BY NAME (`discard` on its evidence entry, listed below) and the fit repeated without it. ACCEPTED when max |T\_model - T| over the KEPT points is below 1 K. C283's 5 \% relative clause judges composition; this fit has no composition residual, so only the temperature clause applies. A pair that splits ONE liquid inside its validity span is not promoted (arm f). Two starts: (0, 0) and (800, -400).

OUTLIERS DISCARDED AFTER THE LEAST-SQUARES STAGE: @@OUTLIERS-BEGIN@@ none: every one of the 28 points lies within 3 K after the least-squares stage. @@OUTLIERS-END@@

@@RESULT-BEGIN@@ WHAT THE RUN FOUND (pinned in `expected`; measured, not argued). fit\_nrtl: b\_ij = 289.9 K, b\_ji = -186.3 K over 28 kept points of 2 laboratories (0 discarded): max |dT| 1.714 K, mean 0.950 K against the acceptance 1 K $\to$ notAccepted. fit\_nrtl\_start2 lands on the same pair (289.9, -186.3). NOT PROMOTABLE: max |dT| 1.714 K over the 28 kept points is not below 1 K @@RESULT-END@@

WHAT IS NOT CLAIMED. The pair is correct at these pressures and in this frame (ideal-gas vapour, the catalogue's own vapour-pressure records); the measurements were transcribed by the archive and have NOT been read back against the articles (the CHECKLIST that `fetch` writes is the aid).

\end{tutnarrative}
\begin{tutgolden}
\goldrow{diag}{fit\_nrtl}{bubble\_aad\_fit}{0.949760771998}
\goldrow{diag}{fit\_nrtl}{bubble\_max\_fit}{1.71426740421}
\goldrow{diag}{fit\_nrtl}{chi2}{443.422429592}
\goldrow{diag}{fit\_nrtl}{chi2\_in\_sample}{443.422429592}
\goldrow{diag}{fit\_nrtl}{chi2\_reduced}{17.0547088304}
\goldrow{diag}{fit\_nrtl}{cond\_JtJ}{557.792966545}
\goldrow{diag}{fit\_nrtl}{converged}{1}
\goldrow{diag}{fit\_nrtl}{corr\_computed}{1}
\goldrow{diag}{fit\_nrtl}{dof}{26}
\goldrow{diag}{fit\_nrtl}{identifiable}{0}
\goldrow{diag}{fit\_nrtl}{invertible}{1}
\goldrow{diag}{fit\_nrtl}{iter}{9}
\goldrow{diag}{fit\_nrtl}{max\_abs\_corr}{0.999104813541}
\goldrow{diag}{fit\_nrtl}{max\_abs\_resid}{1.71426740421}
\goldmore{56}
\end{tutgolden}

\clearpage
\tutentry{curate155\_nrtl\_dimethylCarbonate\_pXylene}{tutorials/props/curation/curate155\_nrtl\_dimethylCarbonate\_pXylene}
\begin{tutscheme}
\begin{tikzpicture}[>={Stealth[length=2mm]}, every node/.style={font=\scriptsize}]
  \node[draw, rounded corners=2pt, align=center, minimum width=3.1cm, minimum height=0.8cm, fill=black!4] at (0.0,-0.0) {\textbf{fit\_nrtl}\\[-1pt]{\tiny fitParameters}};
  \node[draw, rounded corners=2pt, align=center, minimum width=3.1cm, minimum height=0.8cm, fill=black!4] at (3.6,-0.0) {\textbf{fit\_nrtl\_start2}\\[-1pt]{\tiny fitParameters}};
  \node[draw, rounded corners=2pt, align=center, minimum width=3.1cm, minimum height=0.8cm, fill=black!4] at (7.2,-0.0) {\textbf{txy\_ideal}\\[-1pt]{\tiny propertyScan1D}};
  \node[draw, rounded corners=2pt, align=center, minimum width=3.1cm, minimum height=0.8cm, fill=black!4] at (0.0,-1.3) {\textbf{txy\_nrtl}\\[-1pt]{\tiny propertyScan1D}};
\end{tikzpicture}
\end{tutscheme}
\tutwhat{NRTL dimethylCarbonate-pXylene regressed by the curation pipeline to every laboratory's measured bubble temperatures (ThermoML locates, the articles source), accepted at max |dT| < 1 K}
\begin{tutnarrative}
AN NRTL PAIR FROM EVERY LABORATORY'S MEASUREMENTS.

WRITTEN BY bin/choupo-curate pair dimethylCarbonate pXylene (DEV.md 4c C284, the bulk run; record docs/design/binary-pairs-from-open-measurements.md section 11).

THE DECLARATION -- WRITTEN AND COMMITTED BEFORE ANY FIT RAN (several declarations of the bulk run share one commit).

IDENTITY. dimethylCarbonate (CAS 616-38-6, InChIKey IEJIGPNLZYLLBP-UHFFFAOYSA-N) + pXylene (CAS 106-42-3, InChIKey URLKBWYHVLBVBO-UHFFFAOYSA-N), joined to the archive through identity.json -- never by name. CLASS (C284): both records are volatile molecular liquids (no role nonvolatile/solute, no dissociatesTo, no charge, no solid, a vapour-pressure model), and every measured temperature lies below both critical temperatures.

NO MEASURED SERIES IN THIS TREE. constant/archiveDatasets names each dataset (DOI, NIST URL, article sha256, extraction selector, extract sha256); `bin/choupo-curate fetch` writes the points into constant/archiveExtracts/, which is gitignored.

STUDIES: every bubble-temperature block of the pair at or below 110 kPa with a declared uncertainty (2 block(s)). FITTED, ALL OF THEM (2): 10.1016/j.jct.2019.01.027 block 5 14 pts 101.3 kPa 366.65-407.95 K U 0.28-0.75 K 10.1021/acs.jced.7b00372 block 15 11 pts 93.13 kPa 361.49-403.24 K U 0.23-1.12 K (reports x of the other member: written as x[dimethylCarbonate] = 1 - x by --x-of) SET ASIDE BY THE EXTRACTION RULES, BY NAME: (none)

THE RULE (DEV.md 4c C283 amended, C284). Every laboratory above is fitted; none is held out (the C16 partition and its band are retired as the acceptance test). NRTL, a\_ij = a\_ji = 0, alpha = 0.30 fixed: two parameters (b\_ij, b\_ji); each point weighted by 1/U, U its own stated expanded uncertainty. Objective twoStage (least squares, then minimax; T\_bubble absolute, in K). OUTLIERS: after a least-squares stage over every point, a point with |T\_model - T| > 3 K is discarded BY NAME (`discard` on its evidence entry, listed below) and the fit repeated without it. ACCEPTED when max |T\_model - T| over the KEPT points is below 1 K. C283's 5 \% relative clause judges composition; this fit has no composition residual, so only the temperature clause applies. A pair that splits ONE liquid inside its validity span is not promoted (arm f). Two starts: (0, 0) and (800, -400).

OUTLIERS DISCARDED AFTER THE LEAST-SQUARES STAGE: @@OUTLIERS-BEGIN@@ 6 point(s) of 25: doi 10.1016/j.jct.2019.01.027 row 10: x 0.4007, P 101.300 kPa, T 380.65 K, dT -3.06 K; doi 10.1016/j.jct.2019.01.027 row 11: x 0.4878, P 101.300 kPa, T 377.25 K, dT -3.26 K; doi 10.1021/acs.jced.7b00372 row 7: x 0.4312, P 93.130 kPa, T 369.75 K, dT +3.73 K; doi 10.1021/acs.jced.7b00372 row 8: x 0.2628, P 93.130 kPa, T 376.35 K, dT +6.06 K; doi 10.1021/acs.jced.7b00372 row 9: x 0.1584, P 93.130 kPa, T 384.79 K, dT +5.54 K; doi 10.1021/acs.jced.7b00372 row 10: x 0.0631, P 93.130 kPa, T 396.71 K, dT +3.34 K @@OUTLIERS-END@@

@@RESULT-BEGIN@@ WHAT THE RUN FOUND (pinned in `expected`; measured, not argued). fit\_nrtl: b\_ij = 666.1 K, b\_ji = -266.8 K over 19 kept points of 2 laboratories (6 discarded): max |dT| 3.093 K, mean 1.497 K against the acceptance 1 K $\to$ notAccepted. fit\_nrtl\_start2 lands on the same pair (666.1, -266.8). NOT PROMOTABLE: max |dT| 3.093 K over the 19 kept points is not below 1 K @@RESULT-END@@

WHAT IS NOT CLAIMED. The pair is correct at these pressures and in this frame (ideal-gas vapour, the catalogue's own vapour-pressure records); the measurements were transcribed by the archive and have NOT been read back against the articles (the CHECKLIST that `fetch` writes is the aid).

\end{tutnarrative}
\begin{tutgolden}
\goldrow{diag}{fit\_nrtl}{bubble\_aad\_fit}{1.49722969315}
\goldrow{diag}{fit\_nrtl}{bubble\_max\_fit}{3.09305180564}
\goldrow{diag}{fit\_nrtl}{chi2}{453.077209658}
\goldrow{diag}{fit\_nrtl}{chi2\_in\_sample}{453.077209658}
\goldrow{diag}{fit\_nrtl}{chi2\_reduced}{26.6516005681}
\goldrow{diag}{fit\_nrtl}{cond\_JtJ}{20.5722084907}
\goldrow{diag}{fit\_nrtl}{converged}{1}
\goldrow{diag}{fit\_nrtl}{corr\_computed}{1}
\goldrow{diag}{fit\_nrtl}{dof}{17}
\goldrow{diag}{fit\_nrtl}{identifiable}{1}
\goldrow{diag}{fit\_nrtl}{invertible}{1}
\goldrow{diag}{fit\_nrtl}{iter}{12}
\goldrow{diag}{fit\_nrtl}{max\_abs\_corr}{0.976547008656}
\goldrow{diag}{fit\_nrtl}{max\_abs\_resid}{3.09305180564}
\goldmore{56}
\end{tutgolden}

\clearpage
\tutentry{curate156\_nrtl\_cyclohexane\_tetrahydrofuran}{tutorials/props/curation/curate156\_nrtl\_cyclohexane\_tetrahydrofuran}
\begin{tutscheme}
\begin{tikzpicture}[>={Stealth[length=2mm]}, every node/.style={font=\scriptsize}]
  \node[draw, rounded corners=2pt, align=center, minimum width=3.1cm, minimum height=0.8cm, fill=black!4] at (0.0,-0.0) {\textbf{fit\_nrtl}\\[-1pt]{\tiny fitParameters}};
  \node[draw, rounded corners=2pt, align=center, minimum width=3.1cm, minimum height=0.8cm, fill=black!4] at (3.6,-0.0) {\textbf{fit\_nrtl\_start2}\\[-1pt]{\tiny fitParameters}};
  \node[draw, rounded corners=2pt, align=center, minimum width=3.1cm, minimum height=0.8cm, fill=black!4] at (7.2,-0.0) {\textbf{txy\_ideal}\\[-1pt]{\tiny propertyScan1D}};
  \node[draw, rounded corners=2pt, align=center, minimum width=3.1cm, minimum height=0.8cm, fill=black!4] at (0.0,-1.3) {\textbf{txy\_nrtl}\\[-1pt]{\tiny propertyScan1D}};
\end{tikzpicture}
\end{tutscheme}
\tutwhat{NRTL cyclohexane-tetrahydrofuran regressed by the curation pipeline to every laboratory's measured bubble temperatures (ThermoML locates, the articles source), accepted at max |dT| < 1 K}
\begin{tutnarrative}
AN NRTL PAIR FROM EVERY LABORATORY'S MEASUREMENTS.

WRITTEN BY bin/choupo-curate pair cyclohexane tetrahydrofuran (DEV.md 4c C284, the bulk run; record docs/design/binary-pairs-from-open-measurements.md section 11).

THE DECLARATION -- WRITTEN AND COMMITTED BEFORE ANY FIT RAN (several declarations of the bulk run share one commit).

IDENTITY. cyclohexane (CAS 110-82-7, InChIKey XDTMQSROBMDMFD-UHFFFAOYSA-N) + tetrahydrofuran (CAS 109-99-9, InChIKey WYURNTSHIVDZCO-UHFFFAOYSA-N), joined to the archive through identity.json -- never by name. CLASS (C284): both records are volatile molecular liquids (no role nonvolatile/solute, no dissociatesTo, no charge, no solid, a vapour-pressure model), and every measured temperature lies below both critical temperatures.

NO MEASURED SERIES IN THIS TREE. constant/archiveDatasets names each dataset (DOI, NIST URL, article sha256, extraction selector, extract sha256); `bin/choupo-curate fetch` writes the points into constant/archiveExtracts/, which is gitignored.

STUDIES: every bubble-temperature block of the pair at or below 110 kPa with a declared uncertainty (2 block(s)). FITTED, ALL OF THEM (1): 10.1021/je4004993 block 9 20 pts 101.3 kPa 339.04-351.31 K U 0.19-0.4 K SET ASIDE BY THE EXTRACTION RULES, BY NAME: - doi 10.1021/je030141r block 42: only 4 interior point(s) (< 5)

THE RULE (DEV.md 4c C283 amended, C284). Every laboratory above is fitted; none is held out (the C16 partition and its band are retired as the acceptance test). NRTL, a\_ij = a\_ji = 0, alpha = 0.30 fixed: two parameters (b\_ij, b\_ji); each point weighted by 1/U, U its own stated expanded uncertainty. Objective twoStage (least squares, then minimax; T\_bubble absolute, in K). OUTLIERS: after a least-squares stage over every point, a point with |T\_model - T| > 3 K is discarded BY NAME (`discard` on its evidence entry, listed below) and the fit repeated without it. ACCEPTED when max |T\_model - T| over the KEPT points is below 1 K. C283's 5 \% relative clause judges composition; this fit has no composition residual, so only the temperature clause applies. A pair that splits ONE liquid inside its validity span is not promoted (arm f). Two starts: (0, 0) and (800, -400).

OUTLIERS DISCARDED AFTER THE LEAST-SQUARES STAGE: @@OUTLIERS-BEGIN@@ none: every one of the 20 points lies within 3 K after the least-squares stage. @@OUTLIERS-END@@

@@RESULT-BEGIN@@ WHAT THE RUN FOUND (pinned in `expected`; measured, not argued). fit\_nrtl: b\_ij = 72.0 K, b\_ji = 106.2 K over 20 kept points of 1 laboratories (0 discarded): max |dT| 0.334 K, mean 0.136 K against the acceptance 1 K $\to$ accepted. fit\_nrtl\_start2 lands on the same pair (72.0, 106.2). PROMOTED on 2026-11-10 into data/standards/parameters/NRTL/cyclohexane-tetrahydrofuran.dat as reviewStatus autoPromoted, humanReviewed false, because: max |dT| 0.334 K < 1 K over the 20 kept points, one liquid across the validity span @@RESULT-END@@

WHAT IS NOT CLAIMED. The pair is correct at these pressures and in this frame (ideal-gas vapour, the catalogue's own vapour-pressure records); the measurements were transcribed by the archive and have NOT been read back against the articles (the CHECKLIST that `fetch` writes is the aid).

\end{tutnarrative}
\begin{tutgolden}
\goldrow{diag}{fit\_nrtl}{bubble\_aad\_fit}{0.13562540656}
\goldrow{diag}{fit\_nrtl}{bubble\_max\_fit}{0.334322874358}
\goldrow{diag}{fit\_nrtl}{chi2}{9.14674367179}
\goldrow{diag}{fit\_nrtl}{chi2\_in\_sample}{9.14674367179}
\goldrow{diag}{fit\_nrtl}{chi2\_reduced}{0.508152426211}
\goldrow{diag}{fit\_nrtl}{cond\_JtJ}{276.680711385}
\goldrow{diag}{fit\_nrtl}{converged}{1}
\goldrow{diag}{fit\_nrtl}{corr\_computed}{1}
\goldrow{diag}{fit\_nrtl}{dof}{18}
\goldrow{diag}{fit\_nrtl}{identifiable}{1}
\goldrow{diag}{fit\_nrtl}{invertible}{1}
\goldrow{diag}{fit\_nrtl}{iter}{6}
\goldrow{diag}{fit\_nrtl}{max\_abs\_corr}{0.998197746731}
\goldrow{diag}{fit\_nrtl}{max\_abs\_resid}{0.334322874358}
\goldmore{56}
\end{tutgolden}

\clearpage
\tutentry{curate157\_nrtl\_nOctane\_mXylene}{tutorials/props/curation/curate157\_nrtl\_nOctane\_mXylene}
\begin{tutscheme}
\begin{tikzpicture}[>={Stealth[length=2mm]}, every node/.style={font=\scriptsize}]
  \node[draw, rounded corners=2pt, align=center, minimum width=3.1cm, minimum height=0.8cm, fill=black!4] at (0.0,-0.0) {\textbf{fit\_nrtl}\\[-1pt]{\tiny fitParameters}};
  \node[draw, rounded corners=2pt, align=center, minimum width=3.1cm, minimum height=0.8cm, fill=black!4] at (3.6,-0.0) {\textbf{fit\_nrtl\_start2}\\[-1pt]{\tiny fitParameters}};
  \node[draw, rounded corners=2pt, align=center, minimum width=3.1cm, minimum height=0.8cm, fill=black!4] at (7.2,-0.0) {\textbf{txy\_ideal}\\[-1pt]{\tiny propertyScan1D}};
  \node[draw, rounded corners=2pt, align=center, minimum width=3.1cm, minimum height=0.8cm, fill=black!4] at (0.0,-1.3) {\textbf{txy\_unifac}\\[-1pt]{\tiny propertyScan1D}};
  \node[draw, rounded corners=2pt, align=center, minimum width=3.1cm, minimum height=0.8cm, fill=black!4] at (3.6,-1.3) {\textbf{txy\_nrtl}\\[-1pt]{\tiny propertyScan1D}};
\end{tikzpicture}
\end{tutscheme}
\tutwhat{NRTL nOctane-mXylene regressed by the curation pipeline to every laboratory's measured bubble temperatures (ThermoML locates, the articles source), accepted at max |dT| < 1 K}
\begin{tutnarrative}
AN NRTL PAIR FROM EVERY LABORATORY'S MEASUREMENTS.

WRITTEN BY bin/choupo-curate pair nOctane mXylene (DEV.md 4c C284, the bulk run; record docs/design/binary-pairs-from-open-measurements.md section 11).

THE DECLARATION -- WRITTEN AND COMMITTED BEFORE ANY FIT RAN (several declarations of the bulk run share one commit).

IDENTITY. nOctane (CAS 111-65-9, InChIKey TVMXDCGIABBOFY-UHFFFAOYSA-N) + mXylene (CAS 108-38-3, InChIKey IVSZLXZYQVIEFR-UHFFFAOYSA-N), joined to the archive through identity.json -- never by name. CLASS (C284): both records are volatile molecular liquids (no role nonvolatile/solute, no dissociatesTo, no charge, no solid, a vapour-pressure model), and every measured temperature lies below both critical temperatures.

NO MEASURED SERIES IN THIS TREE. constant/archiveDatasets names each dataset (DOI, NIST URL, article sha256, extraction selector, extract sha256); `bin/choupo-curate fetch` writes the points into constant/archiveExtracts/, which is gitignored.

STUDIES: every bubble-temperature block of the pair at or below 110 kPa with a declared uncertainty (2 block(s)). FITTED, ALL OF THEM (2): 10.1016/j.fluid.2013.01.015 block 13 13 pts 20 kPa 349.55-359.83 K U 0.13-0.13 K 10.1021/je3002283 block 14 11 pts 101.3 kPa 399.20-411.40 K U 0.2-0.3 K SET ASIDE BY THE EXTRACTION RULES, BY NAME: (none)

THE RULE (DEV.md 4c C283 amended, C284). Every laboratory above is fitted; none is held out (the C16 partition and its band are retired as the acceptance test). NRTL, a\_ij = a\_ji = 0, alpha = 0.30 fixed: two parameters (b\_ij, b\_ji); each point weighted by 1/U, U its own stated expanded uncertainty. Objective twoStage (least squares, then minimax; T\_bubble absolute, in K). OUTLIERS: after a least-squares stage over every point, a point with |T\_model - T| > 3 K is discarded BY NAME (`discard` on its evidence entry, listed below) and the fit repeated without it. ACCEPTED when max |T\_model - T| over the KEPT points is below 1 K. C283's 5 \% relative clause judges composition; this fit has no composition residual, so only the temperature clause applies. A pair that splits ONE liquid inside its validity span is not promoted (arm f). Two starts: (0, 0) and (800, -400).

OUTLIERS DISCARDED AFTER THE LEAST-SQUARES STAGE: @@OUTLIERS-BEGIN@@ none: every one of the 24 points lies within 3 K after the least-squares stage. @@OUTLIERS-END@@

@@RESULT-BEGIN@@ WHAT THE RUN FOUND (pinned in `expected`; measured, not argued). fit\_nrtl: b\_ij = -311.2 K, b\_ji = 477.8 K over 24 kept points of 2 laboratories (0 discarded): max |dT| 0.635 K, mean 0.200 K against the acceptance 1 K $\to$ accepted. fit\_nrtl\_start2 lands on the same pair (-311.2, 477.8). PROMOTED on 2026-11-10 into data/standards/parameters/NRTL/mXylene-nOctane.dat as reviewStatus autoPromoted, humanReviewed false, because: max |dT| 0.635 K < 1 K over the 24 kept points, one liquid across the validity span @@RESULT-END@@

WHAT IS NOT CLAIMED. The pair is correct at these pressures and in this frame (ideal-gas vapour, the catalogue's own vapour-pressure records); the measurements were transcribed by the archive and have NOT been read back against the articles (the CHECKLIST that `fetch` writes is the aid).

\end{tutnarrative}
\begin{tutgolden}
\goldrow{diag}{fit\_nrtl}{bubble\_aad\_fit}{0.200034995215}
\goldrow{diag}{fit\_nrtl}{bubble\_max\_fit}{0.634891498676}
\goldrow{diag}{fit\_nrtl}{chi2}{33.3713026435}
\goldrow{diag}{fit\_nrtl}{chi2\_in\_sample}{33.3713026435}
\goldrow{diag}{fit\_nrtl}{chi2\_reduced}{1.51687739289}
\goldrow{diag}{fit\_nrtl}{cond\_JtJ}{468.475075911}
\goldrow{diag}{fit\_nrtl}{converged}{1}
\goldrow{diag}{fit\_nrtl}{corr\_computed}{1}
\goldrow{diag}{fit\_nrtl}{dof}{22}
\goldrow{diag}{fit\_nrtl}{identifiable}{1}
\goldrow{diag}{fit\_nrtl}{invertible}{1}
\goldrow{diag}{fit\_nrtl}{iter}{11}
\goldrow{diag}{fit\_nrtl}{max\_abs\_corr}{0.995681145546}
\goldrow{diag}{fit\_nrtl}{max\_abs\_resid}{0.634891498676}
\goldmore{62}
\end{tutgolden}

\clearpage
\tutentry{curate158\_nrtl\_nOctane\_oXylene}{tutorials/props/curation/curate158\_nrtl\_nOctane\_oXylene}
\begin{tutscheme}
\begin{tikzpicture}[>={Stealth[length=2mm]}, every node/.style={font=\scriptsize}]
  \node[draw, rounded corners=2pt, align=center, minimum width=3.1cm, minimum height=0.8cm, fill=black!4] at (0.0,-0.0) {\textbf{fit\_nrtl}\\[-1pt]{\tiny fitParameters}};
  \node[draw, rounded corners=2pt, align=center, minimum width=3.1cm, minimum height=0.8cm, fill=black!4] at (3.6,-0.0) {\textbf{fit\_nrtl\_start2}\\[-1pt]{\tiny fitParameters}};
  \node[draw, rounded corners=2pt, align=center, minimum width=3.1cm, minimum height=0.8cm, fill=black!4] at (7.2,-0.0) {\textbf{txy\_ideal}\\[-1pt]{\tiny propertyScan1D}};
  \node[draw, rounded corners=2pt, align=center, minimum width=3.1cm, minimum height=0.8cm, fill=black!4] at (0.0,-1.3) {\textbf{txy\_unifac}\\[-1pt]{\tiny propertyScan1D}};
  \node[draw, rounded corners=2pt, align=center, minimum width=3.1cm, minimum height=0.8cm, fill=black!4] at (3.6,-1.3) {\textbf{txy\_nrtl}\\[-1pt]{\tiny propertyScan1D}};
\end{tikzpicture}
\end{tutscheme}
\tutwhat{NRTL nOctane-oXylene regressed by the curation pipeline to every laboratory's measured bubble temperatures (ThermoML locates, the articles source), accepted at max |dT| < 1 K}
\begin{tutnarrative}
AN NRTL PAIR FROM EVERY LABORATORY'S MEASUREMENTS.

WRITTEN BY bin/choupo-curate pair nOctane oXylene (DEV.md 4c C284, the bulk run; record docs/design/binary-pairs-from-open-measurements.md section 11).

THE DECLARATION -- WRITTEN AND COMMITTED BEFORE ANY FIT RAN (several declarations of the bulk run share one commit).

IDENTITY. nOctane (CAS 111-65-9, InChIKey TVMXDCGIABBOFY-UHFFFAOYSA-N) + oXylene (CAS 95-47-6, InChIKey CTQNGGLPUBDAKN-UHFFFAOYSA-N), joined to the archive through identity.json -- never by name. CLASS (C284): both records are volatile molecular liquids (no role nonvolatile/solute, no dissociatesTo, no charge, no solid, a vapour-pressure model), and every measured temperature lies below both critical temperatures.

NO MEASURED SERIES IN THIS TREE. constant/archiveDatasets names each dataset (DOI, NIST URL, article sha256, extraction selector, extract sha256); `bin/choupo-curate fetch` writes the points into constant/archiveExtracts/, which is gitignored.

STUDIES: every bubble-temperature block of the pair at or below 110 kPa with a declared uncertainty (2 block(s)). FITTED, ALL OF THEM (2): 10.1016/j.fluid.2013.01.015 block 11 12 pts 20 kPa 349.54-363.73 K U 0.13-0.14 K 10.1021/je3002283 block 12 11 pts 101.3 kPa 399.40-415.70 K U 0.3-0.5 K SET ASIDE BY THE EXTRACTION RULES, BY NAME: (none)

THE RULE (DEV.md 4c C283 amended, C284). Every laboratory above is fitted; none is held out (the C16 partition and its band are retired as the acceptance test). NRTL, a\_ij = a\_ji = 0, alpha = 0.30 fixed: two parameters (b\_ij, b\_ji); each point weighted by 1/U, U its own stated expanded uncertainty. Objective twoStage (least squares, then minimax; T\_bubble absolute, in K). OUTLIERS: after a least-squares stage over every point, a point with |T\_model - T| > 3 K is discarded BY NAME (`discard` on its evidence entry, listed below) and the fit repeated without it. ACCEPTED when max |T\_model - T| over the KEPT points is below 1 K. C283's 5 \% relative clause judges composition; this fit has no composition residual, so only the temperature clause applies. A pair that splits ONE liquid inside its validity span is not promoted (arm f). Two starts: (0, 0) and (800, -400).

OUTLIERS DISCARDED AFTER THE LEAST-SQUARES STAGE: @@OUTLIERS-BEGIN@@ none: every one of the 23 points lies within 3 K after the least-squares stage. @@OUTLIERS-END@@

@@RESULT-BEGIN@@ WHAT THE RUN FOUND (pinned in `expected`; measured, not argued). fit\_nrtl: b\_ij = -334.6 K, b\_ji = 516.4 K over 23 kept points of 2 laboratories (0 discarded): max |dT| 1.130 K, mean 0.410 K against the acceptance 1 K $\to$ notAccepted. fit\_nrtl\_start2 lands on A DIFFERENT pair (403.0, -282.1). NOT PROMOTABLE: max |dT| 1.130 K over the 23 kept points is not below 1 K @@RESULT-END@@

WHAT IS NOT CLAIMED. The pair is correct at these pressures and in this frame (ideal-gas vapour, the catalogue's own vapour-pressure records); the measurements were transcribed by the archive and have NOT been read back against the articles (the CHECKLIST that `fetch` writes is the aid).

\end{tutnarrative}
\begin{tutgolden}
\goldrow{diag}{fit\_nrtl}{bubble\_aad\_fit}{0.409979265099}
\goldrow{diag}{fit\_nrtl}{bubble\_max\_fit}{1.12965578876}
\goldrow{diag}{fit\_nrtl}{chi2}{110.648374871}
\goldrow{diag}{fit\_nrtl}{chi2\_in\_sample}{110.648374871}
\goldrow{diag}{fit\_nrtl}{chi2\_reduced}{5.26897023196}
\goldrow{diag}{fit\_nrtl}{cond\_JtJ}{362.926594531}
\goldrow{diag}{fit\_nrtl}{converged}{1}
\goldrow{diag}{fit\_nrtl}{corr\_computed}{1}
\goldrow{diag}{fit\_nrtl}{dof}{21}
\goldrow{diag}{fit\_nrtl}{identifiable}{1}
\goldrow{diag}{fit\_nrtl}{invertible}{1}
\goldrow{diag}{fit\_nrtl}{iter}{13}
\goldrow{diag}{fit\_nrtl}{max\_abs\_corr}{0.993735366391}
\goldrow{diag}{fit\_nrtl}{max\_abs\_resid}{1.12965578876}
\goldmore{62}
\end{tutgolden}

\clearpage
\tutentry{curate159\_nrtl\_benzene\_methanol}{tutorials/props/curation/curate159\_nrtl\_benzene\_methanol}
\begin{tutscheme}
\begin{tikzpicture}[>={Stealth[length=2mm]}, every node/.style={font=\scriptsize}]
  \node[draw, rounded corners=2pt, align=center, minimum width=3.1cm, minimum height=0.8cm, fill=black!4] at (0.0,-0.0) {\textbf{fit\_nrtl}\\[-1pt]{\tiny fitParameters}};
  \node[draw, rounded corners=2pt, align=center, minimum width=3.1cm, minimum height=0.8cm, fill=black!4] at (3.6,-0.0) {\textbf{fit\_nrtl\_start2}\\[-1pt]{\tiny fitParameters}};
  \node[draw, rounded corners=2pt, align=center, minimum width=3.1cm, minimum height=0.8cm, fill=black!4] at (7.2,-0.0) {\textbf{txy\_ideal}\\[-1pt]{\tiny propertyScan1D}};
  \node[draw, rounded corners=2pt, align=center, minimum width=3.1cm, minimum height=0.8cm, fill=black!4] at (0.0,-1.3) {\textbf{txy\_unifac}\\[-1pt]{\tiny propertyScan1D}};
  \node[draw, rounded corners=2pt, align=center, minimum width=3.1cm, minimum height=0.8cm, fill=black!4] at (3.6,-1.3) {\textbf{txy\_nrtl}\\[-1pt]{\tiny propertyScan1D}};
\end{tikzpicture}
\end{tutscheme}
\tutwhat{NRTL benzene-methanol regressed by the curation pipeline to every laboratory's measured bubble temperatures (ThermoML locates, the articles source), accepted at max |dT| < 1 K}
\begin{tutnarrative}
AN NRTL PAIR FROM EVERY LABORATORY'S MEASUREMENTS.

WRITTEN BY bin/choupo-curate pair benzene methanol (DEV.md 4c C284, the bulk run; record docs/design/binary-pairs-from-open-measurements.md section 11).

THE DECLARATION -- WRITTEN AND COMMITTED BEFORE ANY FIT RAN (several declarations of the bulk run share one commit).

IDENTITY. benzene (CAS 71-43-2, InChIKey UHOVQNZJYSORNB-UHFFFAOYSA-N) + methanol (CAS 67-56-1, InChIKey OKKJLVBELUTLKV-UHFFFAOYSA-N), joined to the archive through identity.json -- never by name. CLASS (C284): both records are volatile molecular liquids (no role nonvolatile/solute, no dissociatesTo, no charge, no solid, a vapour-pressure model), and every measured temperature lies below both critical temperatures.

NO MEASURED SERIES IN THIS TREE. constant/archiveDatasets names each dataset (DOI, NIST URL, article sha256, extraction selector, extract sha256); `bin/choupo-curate fetch` writes the points into constant/archiveExtracts/, which is gitignored.

STUDIES: every bubble-temperature block of the pair at or below 110 kPa with a declared uncertainty (2 block(s)). FITTED, ALL OF THEM (2): 10.1016/j.fluid.2019.03.022 block 0 17 pts 101.3 kPa 331.56-343.83 K U 0.09-1.08 K 10.1016/j.jct.2006.06.002 block 2 6 pts 5 isobars 99.5-100.5 kPa 331.83-336.09 K U 0.25-0.5 K (reports x of the other member: written as x[benzene] = 1 - x by --x-of) SET ASIDE BY THE EXTRACTION RULES, BY NAME: (none)

THE RULE (DEV.md 4c C283 amended, C284). Every laboratory above is fitted; none is held out (the C16 partition and its band are retired as the acceptance test). NRTL, a\_ij = a\_ji = 0, alpha = 0.30 fixed: two parameters (b\_ij, b\_ji); each point weighted by 1/U, U its own stated expanded uncertainty. Objective twoStage (least squares, then minimax; T\_bubble absolute, in K). OUTLIERS: after a least-squares stage over every point, a point with |T\_model - T| > 3 K is discarded BY NAME (`discard` on its evidence entry, listed below) and the fit repeated without it. ACCEPTED when max |T\_model - T| over the KEPT points is below 1 K. C283's 5 \% relative clause judges composition; this fit has no composition residual, so only the temperature clause applies. A pair that splits ONE liquid inside its validity span is not promoted (arm f). Two starts: (0, 0) and (800, -400).

OUTLIERS DISCARDED AFTER THE LEAST-SQUARES STAGE: @@OUTLIERS-BEGIN@@ none: every one of the 23 points lies within 3 K after the least-squares stage. @@OUTLIERS-END@@

@@RESULT-BEGIN@@ WHAT THE RUN FOUND (pinned in `expected`; measured, not argued). fit\_nrtl: b\_ij = 520.2 K, b\_ji = 239.4 K over 23 kept points of 2 laboratories (0 discarded): max |dT| 1.305 K, mean 0.429 K against the acceptance 1 K $\to$ notAccepted. fit\_nrtl\_start2 lands on the same pair (520.2, 239.4). NOT PROMOTABLE: max |dT| 1.305 K over the 23 kept points is not below 1 K @@RESULT-END@@

WHAT IS NOT CLAIMED. The pair is correct at these pressures and in this frame (ideal-gas vapour, the catalogue's own vapour-pressure records); the measurements were transcribed by the archive and have NOT been read back against the articles (the CHECKLIST that `fetch` writes is the aid).

\end{tutnarrative}
\begin{tutgolden}
\goldrow{diag}{fit\_nrtl}{bubble\_aad\_fit}{0.429208699045}
\goldrow{diag}{fit\_nrtl}{bubble\_max\_fit}{1.30533300165}
\goldrow{diag}{fit\_nrtl}{chi2}{81.1067242283}
\goldrow{diag}{fit\_nrtl}{chi2\_in\_sample}{81.1067242283}
\goldrow{diag}{fit\_nrtl}{chi2\_reduced}{3.86222496325}
\goldrow{diag}{fit\_nrtl}{cond\_JtJ}{10.7384038568}
\goldrow{diag}{fit\_nrtl}{converged}{1}
\goldrow{diag}{fit\_nrtl}{corr\_computed}{1}
\goldrow{diag}{fit\_nrtl}{dof}{21}
\goldrow{diag}{fit\_nrtl}{identifiable}{1}
\goldrow{diag}{fit\_nrtl}{invertible}{1}
\goldrow{diag}{fit\_nrtl}{iter}{5}
\goldrow{diag}{fit\_nrtl}{max\_abs\_corr}{0.956456764346}
\goldrow{diag}{fit\_nrtl}{max\_abs\_resid}{1.30533300165}
\goldmore{62}
\end{tutgolden}

\clearpage
\tutentry{curate15\_nrtl\_methanol\_dimethylCarbonate}{tutorials/props/curation/curate15\_nrtl\_methanol\_dimethylCarbonate}
\begin{tutscheme}
\begin{tikzpicture}[>={Stealth[length=2mm]}, every node/.style={font=\scriptsize}]
  \node[draw, rounded corners=2pt, align=center, minimum width=3.1cm, minimum height=0.8cm, fill=black!4] at (0.0,-0.0) {\textbf{fit\_nrtl}\\[-1pt]{\tiny fitParameters}};
  \node[draw, rounded corners=2pt, align=center, minimum width=3.1cm, minimum height=0.8cm, fill=black!4] at (3.6,-0.0) {\textbf{fit\_nrtl\_start2}\\[-1pt]{\tiny fitParameters}};
  \node[draw, rounded corners=2pt, align=center, minimum width=3.1cm, minimum height=0.8cm, fill=black!4] at (7.2,-0.0) {\textbf{txy\_ideal}\\[-1pt]{\tiny propertyScan1D}};
  \node[draw, rounded corners=2pt, align=center, minimum width=3.1cm, minimum height=0.8cm, fill=black!4] at (0.0,-1.3) {\textbf{txy\_nrtl}\\[-1pt]{\tiny propertyScan1D}};
\end{tikzpicture}
\end{tutscheme}
\tutwhat{NRTL methanol-dimethylCarbonate regressed by the curation pipeline to measured bubble temperatures (ThermoML locates, the article sources), held out against an independent study}
\begin{tutnarrative}
AN NRTL PAIR FROM PUBLISHED MEASUREMENTS, BY THE PIPELINE.

WRITTEN BY bin/choupo-curate pair methanol dimethylCarbonate (DEV.md 4c C282; record docs/design/binary-pairs-from-open-measurements.md section 8). Every rule below was decided before this tool existed (C16, C225, C226); the tool applies them and decides nothing.

@@PILOT-BEGIN@@ WHY THIS PAIR -- THE PILOT RANKING (C282; added after the run, but the ranking reads ONLY the archive index and was computed before any fit). Candidates: both components in the catalogue, NO record for the pair in any model directory of data/standards/parameters/, no acetic acid (vapour dimerisation, C16), no component whose vapour pressure is AmbroseWalton (C281), and (checked for the five chosen) no binary block in the archive with a second liquid. Ranked by the number of mutually author-disjoint isobaric bubble-T studies at or below 110 kPa with a declared uncertainty, then by their total points: 1 methanol + dimethylCarbonate 5 independent (9 studies, 274 pts) - acetone + methanol 4 -- passed over: curate05 holds its C16 regression and the two studies' disagreement, a curator's tie-break 2 ethanol + dimethylCarbonate 4 (4 studies, 156 pts) 3 diisopropylEther + isopropanol 3 (4 studies, 152 pts) - triethyleneGlycol + water 3 -- REFUSED by the tool: 2 of 3 studies disagree beyond their U 4 isopropanol + dimethylCarbonate 3 (3 studies, 128 pts) - 2Butanol + dimethylCarbonate 3 -- REFUSED by the tool: no author-disjoint held-out study reports x of the same component 5 ethanol + 224Trimethylpentane 3 (5 studies, 97 pts) THIS CASE: rank 1. @@PILOT-END@@

THE DECLARATION -- WRITTEN AND COMMITTED BEFORE ANY FIT RAN. (C16 slice 2: so the declaration can be read against the result.)

IDENTITY. methanol (CAS 67-56-1, InChIKey OKKJLVBELUTLKV-UHFFFAOYSA-N) + dimethylCarbonate (CAS 616-38-6, InChIKey IEJIGPNLZYLLBP-UHFFFAOYSA-N), joined to the archive through identity.json -- never by name.

NO MEASURED SERIES IN THIS TREE. constant/archiveDatasets names each dataset (DOI, NIST URL, article sha256, extraction selector, extract sha256); `bin/choupo-curate fetch` writes the points into constant/archiveExtracts/, which is gitignored.

STUDIES CONSIDERED: every bubble-temperature block of the pair at or below 110 kPa with a declared uncertainty (9 block(s)). USABLE after the screens: 8. SET ASIDE, BY NAME: - doi 10.1021/je200697m block 2: reports x of dimethylCarbonate, the others of methanol; extract-vle has no reorientation (C16 section 6.1) and a hand edit of a measured file is what the chain exists to prevent

TRIAGE (C226 item 3, leave one study out, fitted by the engine; a study whose held-out AAD exceeds its own largest stated U plus the others' largest stated U disagrees with them beyond their stated uncertainties): 10.1021/je0601309: AAD 0.131 K against the limit 2.400 K $\to$ agrees; 10.1021/acs.jced.7b00185: AAD 0.222 K against the limit 3.210 K $\to$ agrees; 10.1016/j.jct.2012.05.002: AAD 0.411 K against the limit 2.660 K $\to$ agrees; 10.1021/je301290c: AAD 0.260 K against the limit 2.310 K $\to$ agrees; 10.1021/acs.jced.9b00068: AAD 0.613 K against the limit 3.110 K $\to$ agrees; 10.1021/je2005402: AAD 0.453 K against the limit 2.420 K $\to$ agrees; 10.1021/je500443v: AAD 0.193 K against the limit 3.210 K $\to$ agrees; 10.1016/j.fluid.2018.09.005: AAD 0.315 K against the limit 2.570 K $\to$ agrees

PARTITION (the declared rule of C16 slice 2): the FIT set is the usable study with the most interior points; the HELD-OUT set is the next largest that shares NO author with it and was measured at one pressure. FIT doi 10.1021/je0601309 (65 interior points, 5 isobars 40-93.32 kPa) authors: Fukano, M.[Makoto]; Matsuda, H.[Hiroyuki]; Kurihara, K.[Kiyofumi]; Ochi, K.[Kenji] HELD OUT doi 10.1021/je301290c (19 interior points, 101.3 kPa) authors: Chen, X.[Xiaochun]; Cai, F.[Fufeng]; Wu, X.[Xinying]; Asumana, C.[Charles]; Yu, G.[Guangren]

BAND: maxAAD 0.088 \% of the held-out mean T. declared before the fit by bin/choupo-curate, the ONE rule of C16: the held-out set's largest stated expanded uncertainty (U95 = 0.3 K) plus the catalogue vapour-pressure model's own error at that set's pure endpoints (NONE: the held-out study reports no pure-component bubble point, so the term is ZERO, the strictest reading of the rule -- the band is the measurement's own uncertainty alone; measured with no pair involved), as a percentage of the held-out mean T (341.73 K) = 0.0878 \%, rounded UP to two significant figures

THE FIT. NRTL with a\_ij = a\_ji = 0 and alpha = 0.30 fixed: TWO parameters (b\_ij, b\_ji), the C16 frame (bubble temperatures over these ranges do not identify more). Each point weighted by 1/U, U its own stated expanded uncertainty. Two starts: (0, 0) and (800, -400).

@@RESULT-BEGIN@@ WHAT THE RUN FOUND (pinned in `expected`; measured, not argued). fit\_nrtl: b\_ij = 401.3 K, b\_ji = 89.4 K, in-sample rms 0.176 K over 65 points; fit\_nrtl\_start2 lands on the same pair (401.3, 89.4). HELD OUT: AAD 0.219 K (0.0641 \%) against the band 0.088 \% $\to$ validated. PROMOTED on 2026-11-10 into data/standards/parameters/NRTL/dimethylCarbonate-methanol.dat as reviewStatus autoPromoted, humanReviewed false, because: verdict validated, one liquid across the validity span, triage run @@RESULT-END@@

WHAT IS NOT CLAIMED. The pair is correct at these pressures and in this frame (ideal-gas vapour, the catalogue's own vapour-pressure records); the measurements were transcribed by the archive and have NOT been read back against the articles (the CHECKLIST that `fetch` writes beside the extracts is the aid for that reading).

\end{tutnarrative}
\begin{tutgolden}
\goldrow{diag}{fit\_nrtl}{aad\_heldout\_K}{0.218818170286}
\goldrow{diag}{fit\_nrtl}{aad\_heldout\_pct}{0.0640801029909}
\goldrow{diag}{fit\_nrtl}{bubble\_aad\_fit}{0.125600062097}
\goldrow{diag}{fit\_nrtl}{bubble\_aad\_heldout}{0.218818170286}
\goldrow{diag}{fit\_nrtl}{bubble\_max\_fit}{0.578243594749}
\goldrow{diag}{fit\_nrtl}{bubble\_max\_heldout}{0.50335641624}
\goldrow{diag}{fit\_nrtl}{chi2}{34.2428562333}
\goldrow{diag}{fit\_nrtl}{chi2\_in\_sample}{34.2428562333}
\goldrow{diag}{fit\_nrtl}{chi2\_reduced}{0.543537400529}
\goldrow{diag}{fit\_nrtl}{cond\_JtJ}{21.4954110232}
\goldrow{diag}{fit\_nrtl}{converged}{1}
\goldrow{diag}{fit\_nrtl}{corr\_computed}{1}
\goldrow{diag}{fit\_nrtl}{dof}{63}
\goldrow{diag}{fit\_nrtl}{identifiable}{1}
\goldmore{58}
\end{tutgolden}

\clearpage
\tutentry{curate160\_nrtl\_cyclohexane\_nHeptane}{tutorials/props/curation/curate160\_nrtl\_cyclohexane\_nHeptane}
\begin{tutscheme}
\begin{tikzpicture}[>={Stealth[length=2mm]}, every node/.style={font=\scriptsize}]
  \node[draw, rounded corners=2pt, align=center, minimum width=3.1cm, minimum height=0.8cm, fill=black!4] at (0.0,-0.0) {\textbf{fit\_nrtl}\\[-1pt]{\tiny fitParameters}};
  \node[draw, rounded corners=2pt, align=center, minimum width=3.1cm, minimum height=0.8cm, fill=black!4] at (3.6,-0.0) {\textbf{fit\_nrtl\_start2}\\[-1pt]{\tiny fitParameters}};
  \node[draw, rounded corners=2pt, align=center, minimum width=3.1cm, minimum height=0.8cm, fill=black!4] at (7.2,-0.0) {\textbf{txy\_ideal}\\[-1pt]{\tiny propertyScan1D}};
  \node[draw, rounded corners=2pt, align=center, minimum width=3.1cm, minimum height=0.8cm, fill=black!4] at (0.0,-1.3) {\textbf{txy\_unifac}\\[-1pt]{\tiny propertyScan1D}};
  \node[draw, rounded corners=2pt, align=center, minimum width=3.1cm, minimum height=0.8cm, fill=black!4] at (3.6,-1.3) {\textbf{txy\_nrtl}\\[-1pt]{\tiny propertyScan1D}};
\end{tikzpicture}
\end{tutscheme}
\tutwhat{NRTL cyclohexane-nHeptane regressed by the curation pipeline to every laboratory's measured bubble temperatures (ThermoML locates, the articles source), accepted at max |dT| < 1 K}
\begin{tutnarrative}
AN NRTL PAIR FROM EVERY LABORATORY'S MEASUREMENTS.

WRITTEN BY bin/choupo-curate pair cyclohexane nHeptane (DEV.md 4c C284, the bulk run; record docs/design/binary-pairs-from-open-measurements.md section 11).

THE DECLARATION -- WRITTEN AND COMMITTED BEFORE ANY FIT RAN (several declarations of the bulk run share one commit).

IDENTITY. cyclohexane (CAS 110-82-7, InChIKey XDTMQSROBMDMFD-UHFFFAOYSA-N) + nHeptane (CAS 142-82-5, InChIKey IMNFDUFMRHMDMM-UHFFFAOYSA-N), joined to the archive through identity.json -- never by name. CLASS (C284): both records are volatile molecular liquids (no role nonvolatile/solute, no dissociatesTo, no charge, no solid, a vapour-pressure model), and every measured temperature lies below both critical temperatures.

NO MEASURED SERIES IN THIS TREE. constant/archiveDatasets names each dataset (DOI, NIST URL, article sha256, extraction selector, extract sha256); `bin/choupo-curate fetch` writes the points into constant/archiveExtracts/, which is gitignored.

STUDIES: every bubble-temperature block of the pair at or below 110 kPa with a declared uncertainty (2 block(s)). FITTED, ALL OF THEM (2): 10.1016/j.jct.2019.02.006 block 11 10 pts 101.3 kPa 355.41-369.87 K U 1.04-1.1 K 10.1021/acs.jced.8b00929 block 6 9 pts 101.3 kPa 355.19-368.67 K U 0.41-0.42 K SET ASIDE BY THE EXTRACTION RULES, BY NAME: (none)

THE RULE (DEV.md 4c C283 amended, C284). Every laboratory above is fitted; none is held out (the C16 partition and its band are retired as the acceptance test). NRTL, a\_ij = a\_ji = 0, alpha = 0.30 fixed: two parameters (b\_ij, b\_ji); each point weighted by 1/U, U its own stated expanded uncertainty. Objective twoStage (least squares, then minimax; T\_bubble absolute, in K). OUTLIERS: after a least-squares stage over every point, a point with |T\_model - T| > 3 K is discarded BY NAME (`discard` on its evidence entry, listed below) and the fit repeated without it. ACCEPTED when max |T\_model - T| over the KEPT points is below 1 K. C283's 5 \% relative clause judges composition; this fit has no composition residual, so only the temperature clause applies. A pair that splits ONE liquid inside its validity span is not promoted (arm f). Two starts: (0, 0) and (800, -400).

OUTLIERS DISCARDED AFTER THE LEAST-SQUARES STAGE: @@OUTLIERS-BEGIN@@ none: every one of the 19 points lies within 3 K after the least-squares stage. @@OUTLIERS-END@@

@@RESULT-BEGIN@@ WHAT THE RUN FOUND (pinned in `expected`; measured, not argued). fit\_nrtl: b\_ij = 278.5 K, b\_ji = -197.6 K over 19 kept points of 2 laboratories (0 discarded): max |dT| 0.651 K, mean 0.308 K against the acceptance 1 K $\to$ accepted. fit\_nrtl\_start2 lands on the same pair (278.5, -197.6). PROMOTED on 2026-11-10 into data/standards/parameters/NRTL/cyclohexane-nHeptane.dat as reviewStatus autoPromoted, humanReviewed false, because: max |dT| 0.651 K < 1 K over the 19 kept points, one liquid across the validity span @@RESULT-END@@

WHAT IS NOT CLAIMED. The pair is correct at these pressures and in this frame (ideal-gas vapour, the catalogue's own vapour-pressure records); the measurements were transcribed by the archive and have NOT been read back against the articles (the CHECKLIST that `fetch` writes is the aid).

\end{tutnarrative}
\begin{tutgolden}
\goldrow{diag}{fit\_nrtl}{bubble\_aad\_fit}{0.308175974818}
\goldrow{diag}{fit\_nrtl}{bubble\_max\_fit}{0.651167679322}
\goldrow{diag}{fit\_nrtl}{chi2}{4.07878507649}
\goldrow{diag}{fit\_nrtl}{chi2\_in\_sample}{4.07878507649}
\goldrow{diag}{fit\_nrtl}{chi2\_reduced}{0.239928533911}
\goldrow{diag}{fit\_nrtl}{cond\_JtJ}{1085.31281544}
\goldrow{diag}{fit\_nrtl}{converged}{1}
\goldrow{diag}{fit\_nrtl}{corr\_computed}{1}
\goldrow{diag}{fit\_nrtl}{dof}{17}
\goldrow{diag}{fit\_nrtl}{identifiable}{0}
\goldrow{diag}{fit\_nrtl}{invertible}{1}
\goldrow{diag}{fit\_nrtl}{iter}{12}
\goldrow{diag}{fit\_nrtl}{max\_abs\_corr}{0.999539621448}
\goldrow{diag}{fit\_nrtl}{max\_abs\_resid}{0.651167679322}
\goldmore{71}
\end{tutgolden}

\clearpage
\tutentry{curate161\_nrtl\_nDecane\_mXylene}{tutorials/props/curation/curate161\_nrtl\_nDecane\_mXylene}
\begin{tutscheme}
\begin{tikzpicture}[>={Stealth[length=2mm]}, every node/.style={font=\scriptsize}]
  \node[draw, rounded corners=2pt, align=center, minimum width=3.1cm, minimum height=0.8cm, fill=black!4] at (0.0,-0.0) {\textbf{fit\_nrtl}\\[-1pt]{\tiny fitParameters}};
  \node[draw, rounded corners=2pt, align=center, minimum width=3.1cm, minimum height=0.8cm, fill=black!4] at (3.6,-0.0) {\textbf{fit\_nrtl\_start2}\\[-1pt]{\tiny fitParameters}};
  \node[draw, rounded corners=2pt, align=center, minimum width=3.1cm, minimum height=0.8cm, fill=black!4] at (7.2,-0.0) {\textbf{txy\_ideal}\\[-1pt]{\tiny propertyScan1D}};
  \node[draw, rounded corners=2pt, align=center, minimum width=3.1cm, minimum height=0.8cm, fill=black!4] at (0.0,-1.3) {\textbf{txy\_unifac}\\[-1pt]{\tiny propertyScan1D}};
  \node[draw, rounded corners=2pt, align=center, minimum width=3.1cm, minimum height=0.8cm, fill=black!4] at (3.6,-1.3) {\textbf{txy\_nrtl}\\[-1pt]{\tiny propertyScan1D}};
\end{tikzpicture}
\end{tutscheme}
\tutwhat{NRTL nDecane-mXylene regressed by the curation pipeline to every laboratory's measured bubble temperatures (ThermoML locates, the articles source), accepted at max |dT| < 1 K}
\begin{tutnarrative}
AN NRTL PAIR FROM EVERY LABORATORY'S MEASUREMENTS.

WRITTEN BY bin/choupo-curate pair nDecane mXylene (DEV.md 4c C284, the bulk run; record docs/design/binary-pairs-from-open-measurements.md section 11).

THE DECLARATION -- WRITTEN AND COMMITTED BEFORE ANY FIT RAN (several declarations of the bulk run share one commit).

IDENTITY. nDecane (CAS 124-18-5, InChIKey DIOQZVSQGTUSAI-UHFFFAOYSA-N) + mXylene (CAS 108-38-3, InChIKey IVSZLXZYQVIEFR-UHFFFAOYSA-N), joined to the archive through identity.json -- never by name. CLASS (C284): both records are volatile molecular liquids (no role nonvolatile/solute, no dissociatesTo, no charge, no solid, a vapour-pressure model), and every measured temperature lies below both critical temperatures.

NO MEASURED SERIES IN THIS TREE. constant/archiveDatasets names each dataset (DOI, NIST URL, article sha256, extraction selector, extract sha256); `bin/choupo-curate fetch` writes the points into constant/archiveExtracts/, which is gitignored.

STUDIES: every bubble-temperature block of the pair at or below 110 kPa with a declared uncertainty (2 block(s)). FITTED, ALL OF THEM (2): 10.1016/j.fluid.2012.09.019 block 14 12 pts 20 kPa 361.71-391.29 K U 0.1-0.17 K 10.1016/j.fluid.2004.10.027 block 25 5 pts 95.8 kPa 420.05-441.55 K U 1.01-1.01 K (reports x of the other member: written as x[nDecane] = 1 - x by --x-of) SET ASIDE BY THE EXTRACTION RULES, BY NAME: (none)

THE RULE (DEV.md 4c C283 amended, C284). Every laboratory above is fitted; none is held out (the C16 partition and its band are retired as the acceptance test). NRTL, a\_ij = a\_ji = 0, alpha = 0.30 fixed: two parameters (b\_ij, b\_ji); each point weighted by 1/U, U its own stated expanded uncertainty. Objective twoStage (least squares, then minimax; T\_bubble absolute, in K). OUTLIERS: after a least-squares stage over every point, a point with |T\_model - T| > 3 K is discarded BY NAME (`discard` on its evidence entry, listed below) and the fit repeated without it. ACCEPTED when max |T\_model - T| over the KEPT points is below 1 K. C283's 5 \% relative clause judges composition; this fit has no composition residual, so only the temperature clause applies. A pair that splits ONE liquid inside its validity span is not promoted (arm f). Two starts: (0, 0) and (800, -400).

OUTLIERS DISCARDED AFTER THE LEAST-SQUARES STAGE: @@OUTLIERS-BEGIN@@ 5 point(s) of 17: doi 10.1016/j.fluid.2004.10.027 row 1: x 0.7909, P 95.800 kPa, T 441.55 K, dT -8.02 K; doi 10.1016/j.fluid.2004.10.027 row 2: x 0.6021, P 95.800 kPa, T 436.75 K, dT -10.87 K; doi 10.1016/j.fluid.2004.10.027 row 3: x 0.5195, P 95.800 kPa, T 434.05 K, dT -10.94 K; doi 10.1016/j.fluid.2004.10.027 row 4: x 0.3062, P 95.800 kPa, T 425.25 K, dT -8.16 K; doi 10.1016/j.fluid.2004.10.027 row 5: x 0.2014, P 95.800 kPa, T 420.05 K, dT -5.48 K @@OUTLIERS-END@@

@@RESULT-BEGIN@@ WHAT THE RUN FOUND (pinned in `expected`; measured, not argued). fit\_nrtl: b\_ij = -263.0 K, b\_ji = 429.5 K over 12 kept points of 2 laboratories (5 discarded): max |dT| 0.107 K, mean 0.064 K against the acceptance 1 K $\to$ accepted. fit\_nrtl\_start2 lands on the same pair (-263.0, 429.5). PROMOTED on 2026-11-10 into data/standards/parameters/NRTL/mXylene-nDecane.dat as reviewStatus autoPromoted, humanReviewed false, because: max |dT| 0.107 K < 1 K over the 12 kept points, one liquid across the validity span @@RESULT-END@@

WHAT IS NOT CLAIMED. The pair is correct at these pressures and in this frame (ideal-gas vapour, the catalogue's own vapour-pressure records); the measurements were transcribed by the archive and have NOT been read back against the articles (the CHECKLIST that `fetch` writes is the aid).

\end{tutnarrative}
\begin{tutgolden}
\goldrow{diag}{fit\_nrtl}{bubble\_aad\_fit}{0.0641896452359}
\goldrow{diag}{fit\_nrtl}{bubble\_max\_fit}{0.10698292439}
\goldrow{diag}{fit\_nrtl}{chi2}{3.32846711286}
\goldrow{diag}{fit\_nrtl}{chi2\_in\_sample}{3.32846711286}
\goldrow{diag}{fit\_nrtl}{chi2\_reduced}{0.332846711286}
\goldrow{diag}{fit\_nrtl}{cond\_JtJ}{533.230446271}
\goldrow{diag}{fit\_nrtl}{converged}{1}
\goldrow{diag}{fit\_nrtl}{corr\_computed}{1}
\goldrow{diag}{fit\_nrtl}{dof}{10}
\goldrow{diag}{fit\_nrtl}{identifiable}{1}
\goldrow{diag}{fit\_nrtl}{invertible}{1}
\goldrow{diag}{fit\_nrtl}{iter}{11}
\goldrow{diag}{fit\_nrtl}{max\_abs\_corr}{0.997133766604}
\goldrow{diag}{fit\_nrtl}{max\_abs\_resid}{0.10698292439}
\goldmore{62}
\end{tutgolden}

\clearpage
\tutentry{curate162\_nrtl\_nDecane\_oXylene}{tutorials/props/curation/curate162\_nrtl\_nDecane\_oXylene}
\begin{tutscheme}
\begin{tikzpicture}[>={Stealth[length=2mm]}, every node/.style={font=\scriptsize}]
  \node[draw, rounded corners=2pt, align=center, minimum width=3.1cm, minimum height=0.8cm, fill=black!4] at (0.0,-0.0) {\textbf{fit\_nrtl}\\[-1pt]{\tiny fitParameters}};
  \node[draw, rounded corners=2pt, align=center, minimum width=3.1cm, minimum height=0.8cm, fill=black!4] at (3.6,-0.0) {\textbf{fit\_nrtl\_start2}\\[-1pt]{\tiny fitParameters}};
  \node[draw, rounded corners=2pt, align=center, minimum width=3.1cm, minimum height=0.8cm, fill=black!4] at (7.2,-0.0) {\textbf{txy\_ideal}\\[-1pt]{\tiny propertyScan1D}};
  \node[draw, rounded corners=2pt, align=center, minimum width=3.1cm, minimum height=0.8cm, fill=black!4] at (0.0,-1.3) {\textbf{txy\_unifac}\\[-1pt]{\tiny propertyScan1D}};
  \node[draw, rounded corners=2pt, align=center, minimum width=3.1cm, minimum height=0.8cm, fill=black!4] at (3.6,-1.3) {\textbf{txy\_nrtl}\\[-1pt]{\tiny propertyScan1D}};
\end{tikzpicture}
\end{tutscheme}
\tutwhat{NRTL nDecane-oXylene regressed by the curation pipeline to every laboratory's measured bubble temperatures (ThermoML locates, the articles source), accepted at max |dT| < 1 K}
\begin{tutnarrative}
AN NRTL PAIR FROM EVERY LABORATORY'S MEASUREMENTS.

WRITTEN BY bin/choupo-curate pair nDecane oXylene (DEV.md 4c C284, the bulk run; record docs/design/binary-pairs-from-open-measurements.md section 11).

THE DECLARATION -- WRITTEN AND COMMITTED BEFORE ANY FIT RAN (several declarations of the bulk run share one commit).

IDENTITY. nDecane (CAS 124-18-5, InChIKey DIOQZVSQGTUSAI-UHFFFAOYSA-N) + oXylene (CAS 95-47-6, InChIKey CTQNGGLPUBDAKN-UHFFFAOYSA-N), joined to the archive through identity.json -- never by name. CLASS (C284): both records are volatile molecular liquids (no role nonvolatile/solute, no dissociatesTo, no charge, no solid, a vapour-pressure model), and every measured temperature lies below both critical temperatures.

NO MEASURED SERIES IN THIS TREE. constant/archiveDatasets names each dataset (DOI, NIST URL, article sha256, extraction selector, extract sha256); `bin/choupo-curate fetch` writes the points into constant/archiveExtracts/, which is gitignored.

STUDIES: every bubble-temperature block of the pair at or below 110 kPa with a declared uncertainty (2 block(s)). FITTED, ALL OF THEM (2): 10.1016/j.fluid.2012.09.019 block 12 11 pts 20 kPa 366.11-391.33 K U 0.12-0.14 K 10.1016/j.fluid.2004.10.027 block 24 5 pts 95.8 kPa 417.45-434.95 K U 1.01-1.01 K (reports x of the other member: written as x[nDecane] = 1 - x by --x-of) SET ASIDE BY THE EXTRACTION RULES, BY NAME: (none)

THE RULE (DEV.md 4c C283 amended, C284). Every laboratory above is fitted; none is held out (the C16 partition and its band are retired as the acceptance test). NRTL, a\_ij = a\_ji = 0, alpha = 0.30 fixed: two parameters (b\_ij, b\_ji); each point weighted by 1/U, U its own stated expanded uncertainty. Objective twoStage (least squares, then minimax; T\_bubble absolute, in K). OUTLIERS: after a least-squares stage over every point, a point with |T\_model - T| > 3 K is discarded BY NAME (`discard` on its evidence entry, listed below) and the fit repeated without it. ACCEPTED when max |T\_model - T| over the KEPT points is below 1 K. C283's 5 \% relative clause judges composition; this fit has no composition residual, so only the temperature clause applies. A pair that splits ONE liquid inside its validity span is not promoted (arm f). Two starts: (0, 0) and (800, -400).

OUTLIERS DISCARDED AFTER THE LEAST-SQUARES STAGE: @@OUTLIERS-BEGIN@@ 1 point(s) of 16: doi 10.1016/j.fluid.2004.10.027 row 2: x 0.7124, P 95.800 kPa, T 430.95 K, dT +3.11 K @@OUTLIERS-END@@

@@RESULT-BEGIN@@ WHAT THE RUN FOUND (pinned in `expected`; measured, not argued). fit\_nrtl: b\_ij = -293.0 K, b\_ji = 408.0 K over 15 kept points of 2 laboratories (1 discarded): max |dT| 2.699 K, mean 0.416 K against the acceptance 1 K $\to$ notAccepted. fit\_nrtl\_start2 lands on the same pair (-293.0, 408.0). NOT PROMOTABLE: max |dT| 2.699 K over the 15 kept points is not below 1 K @@RESULT-END@@

WHAT IS NOT CLAIMED. The pair is correct at these pressures and in this frame (ideal-gas vapour, the catalogue's own vapour-pressure records); the measurements were transcribed by the archive and have NOT been read back against the articles (the CHECKLIST that `fetch` writes is the aid).

\end{tutnarrative}
\begin{tutgolden}
\goldrow{diag}{fit\_nrtl}{bubble\_aad\_fit}{0.416442946268}
\goldrow{diag}{fit\_nrtl}{bubble\_max\_fit}{2.69877249758}
\goldrow{diag}{fit\_nrtl}{chi2}{41.6359660005}
\goldrow{diag}{fit\_nrtl}{chi2\_in\_sample}{41.6359660005}
\goldrow{diag}{fit\_nrtl}{chi2\_reduced}{3.20276661543}
\goldrow{diag}{fit\_nrtl}{cond\_JtJ}{1091.18521384}
\goldrow{diag}{fit\_nrtl}{converged}{1}
\goldrow{diag}{fit\_nrtl}{corr\_computed}{1}
\goldrow{diag}{fit\_nrtl}{dof}{13}
\goldrow{diag}{fit\_nrtl}{identifiable}{1}
\goldrow{diag}{fit\_nrtl}{invertible}{1}
\goldrow{diag}{fit\_nrtl}{iter}{21}
\goldrow{diag}{fit\_nrtl}{max\_abs\_corr}{0.998581479335}
\goldrow{diag}{fit\_nrtl}{max\_abs\_resid}{2.69877249758}
\goldmore{62}
\end{tutgolden}

\clearpage
\tutentry{curate164\_vp\_2Ethoxyethanol}{tutorials/props/curation/curate164\_vp\_2Ethoxyethanol}
\begin{tutscheme}
\begin{tikzpicture}[>={Stealth[length=2mm]}, every node/.style={font=\scriptsize}]
  \node[draw, rounded corners=2pt, align=center, minimum width=3.1cm, minimum height=0.8cm, fill=black!4] at (0.0,-0.0) {\textbf{fit\_antoine}\\[-1pt]{\tiny vaporPressureFit}};
\end{tikzpicture}
\end{tutscheme}
\tutwhat{Antoine vapour pressure of 2Ethoxyethanol fitted by the curation pipeline to the measured vapour pressures of every usable laboratory (ThermoML locates, the articles source), outliers discarded by name}
\begin{tutnarrative}
the vapour pressure of 2Ethoxyethanol, REFITTED to measurements (DEV.md 4c C281 and C283). Written by `bin/choupo-curate vp 2Ethoxyethanol`; everything above the RESULT paragraph is the DECLARATION, committed before any fit ran, so it can be read against the result.

WHY. The catalogue record carried `vaporPressure \{ model AmbroseWalton; \}`, a corresponding-states ESTIMATE from Tc, Pc and omega that never reads the record's own Tb (C279 measured a median 19.6 \% AAD over the 25 such records with archive data). Here a three-parameter Antoine is fitted to EVERY usable laboratory's measurements of it.

STUDIES (one per article, extracted by `choupo-thermoml extract-vp`; ranked by points, ties by DOI): 10.1016/j.fluid.2007.03.030 25 pts 345.63-407.95 K Lladosa; Monton; Burguet; Munoz 10.1016/j.fluid.2012.08.030 11 pts 373.48-406.33 K Reddy; Kumar; Rao; Sri et al. 10.1021/je300202z 8 pts 388.6-407.8 K Rana; Bhate; Mahajani; Dabke SET ASIDE BY NAME: : the article file is not in the local cache 10.1021/je0255628: 4 pressure point(s), fewer than 5 10.1021/je025566c: 3 pressure point(s), fewer than 5 10.1021/je034153o: 3 pressure point(s), fewer than 5 10.1021/je200565u: 3 pressure point(s), fewer than 5

FIT AND GRADED ACCEPTANCE (DEV.md 4c C285, Vitor 2026-11-10; the evidence, objective and outliers of C283; declared before the fit). EVIDENCE: every study above is fitted; none is held out (C283 amended: laboratories disagree with each other by more than they claim, so the test is agreement across all of them, not one laboratory's claimed precision). OBJECTIVE: `objective \{ method twoStage; \}` -- least squares in log10 P for a first estimate, then the LARGEST relative error in P, (P\_model - P)/P, over the kept points minimised from it (Vitor's practice, C280); both answers kept. OUTLIERS: after the least-squares stage a point with a relative P error above 15 \% (three times the acceptance) is DISCARDED BY NAME (its study DOI, T and P), listed in the result and the record, and the fit repeated without it, once (`outliers \{ maxRelativeError 15 percent; \}`). ACCEPTANCE, GRADED, the DEFAULT of C285: the maximum relative P error over all kept points is graded clean up to 5 \%, warning up to 10 \%, severe up to 20 \% and refused above; where the normal-boiling-point condition below applies, the absolute gap is graded clean up to 1 K, warning up to 2 K, severe up to 4 K and refused above. The fit's grade is the WORSE of the two. A clean, warning or severe fit is ACCEPTED and promoted (autoPromoted, humanReviewed false) with `fitGrade <grade>;` and the deciding numbers in its provenance entry, which the component loader voices on every run that uses it (a warning or a severe grade); a refused one is not promoted. This case declares no `system/curationDict`, so the default applies; a case that needs another acceptance declares it there, with its reason. The AmbroseWalton curve's errors on the same points are reported beside the fit's, as information.

NORMAL BOILING POINT: the record's Tb 408.405 K is judged against the range of the KEPT points (all studies span 345.63-407.95 K before any outlier is discarded): if Tb lies OUTSIDE that range, the fitted Antoine's own normal boiling point (P = 101.325 kPa) is GRADED by its absolute gap to it (the temperature limits of the acceptance below); inside it, the gap is reported and is not a condition.

NO MEASURED POINT IS IN THIS TREE: constant/archiveDatasets declares every series by DOI, NIST URL and sha256; `bin/choupo-curate fetch` regenerates them into constant/archiveExtracts/ (gitignored).

RESULT @@RESULT-BEGIN@@ WHAT THE RUN FOUND (pinned in `expected`; measured, not argued). Objective twoStage. Antoine log10(Psat/bar) = 54.031976 - 182574.8370/(T + 2972.8847), in-sample R2 0.996203 over 44 kept points, 345.63-407.95 K (10.1016/j.fluid.2007.03.030 25/25 kept; 10.1016/j.fluid.2012.08.030 11/11 kept; 10.1021/je300202z 8/8 kept). Relative error in P over the kept points: mean 3.5407 \%, max 5.7089 \%; AmbroseWalton on the same points: mean 23.3634 \%, max 39.9725 \%. Outliers discarded: none. Its own normal boiling point (101.325 kPa): 406.49 K. PROMOTED on 2026-11-10 into data/standards/components/2Ethoxyethanol.dat as reviewStatus autoPromoted, humanReviewed false, because: fit grade warning under the DEFAULT graded acceptance of DEV.md 4c C285 (no system/curationDict), 3 studies, 44 points kept, 0 discarded as outliers; max relative error 5.709 \% (mean 3.541 \%), grade warning (limits 5 / 10 / 20 \%); Tb outside the kept points' range, normal boiling point 406.49 K, -1.92 K from the record's Tb 408.405 K, grade warning (limits 1 / 2 / 4 K); AmbroseWalton on the same points mean 23.36 \% / max 39.97 \% @@RESULT-END@@

\end{tutnarrative}
\begin{tutgolden}
\goldrow{diag}{fit\_antoine}{A}{54.0319756856}
\goldrow{diag}{fit\_antoine}{B}{182574.837039}
\goldrow{diag}{fit\_antoine}{C}{2972.88465544}
\goldrow{diag}{fit\_antoine}{R2}{0.996203203825}
\goldrow{diag}{fit\_antoine}{dHvap\_kJ\_per\_mol}{50.8928541401}
\goldrow{diag}{fit\_antoine}{minimax\_evaluations}{1501}
\goldrow{diag}{fit\_antoine}{minimax\_max\_scaled\_final}{0.0570892137688}
\goldrow{diag}{fit\_antoine}{minimax\_max\_scaled\_start}{0.0925668895744}
\goldrow{diag}{fit\_antoine}{n\_heldout}{0}
\goldrow{diag}{fit\_antoine}{n\_outliers\_dropped}{0}
\goldrow{diag}{fit\_antoine}{n\_points}{44}
\goldrow{diag}{fit\_antoine}{objective\_method}{2}
\goldrow{diag}{fit\_antoine}{relErr\_max\_fit\_pct}{5.70892137688}
\goldrow{diag}{fit\_antoine}{relErr\_max\_fit\_pct\_leastSquares}{9.25668895744}
\goldmore{3}
\end{tutgolden}

\clearpage
\tutentry{curate165\_vp\_ethylOctanoate}{tutorials/props/curation/curate165\_vp\_ethylOctanoate}
\begin{tutscheme}
\begin{tikzpicture}[>={Stealth[length=2mm]}, every node/.style={font=\scriptsize}]
  \node[draw, rounded corners=2pt, align=center, minimum width=3.1cm, minimum height=0.8cm, fill=black!4] at (0.0,-0.0) {\textbf{fit\_antoine}\\[-1pt]{\tiny vaporPressureFit}};
\end{tikzpicture}
\end{tutscheme}
\tutwhat{Antoine vapour pressure of ethylOctanoate fitted by the curation pipeline to the measured vapour pressures of every usable laboratory (ThermoML locates, the articles source), outliers discarded by name}
\begin{tutnarrative}
the vapour pressure of ethylOctanoate, REFITTED to measurements (DEV.md 4c C281 and C283). Written by `bin/choupo-curate vp ethylOctanoate`; everything above the RESULT paragraph is the DECLARATION, committed before any fit ran, so it can be read against the result.

WHY. The catalogue record carried `vaporPressure \{ model AmbroseWalton; \}`, a corresponding-states ESTIMATE from Tc, Pc and omega that never reads the record's own Tb (C279 measured a median 19.6 \% AAD over the 25 such records with archive data). Here a three-parameter Antoine is fitted to EVERY usable laboratory's measurements of it.

STUDIES (one per article, extracted by `choupo-thermoml extract-vp`; ranked by points, ties by DOI): 10.1021/je200730m 17 pts 273.26-452.3 K Benziane; Khimeche; Mokbel; Sawaya et al. SET ASIDE BY NAME: (none)

FIT AND GRADED ACCEPTANCE (DEV.md 4c C285, Vitor 2026-11-10; the evidence, objective and outliers of C283; declared before the fit). EVIDENCE: every study above is fitted; none is held out (C283 amended: laboratories disagree with each other by more than they claim, so the test is agreement across all of them, not one laboratory's claimed precision). OBJECTIVE: `objective \{ method twoStage; \}` -- least squares in log10 P for a first estimate, then the LARGEST relative error in P, (P\_model - P)/P, over the kept points minimised from it (Vitor's practice, C280); both answers kept. OUTLIERS: after the least-squares stage a point with a relative P error above 15 \% (three times the acceptance) is DISCARDED BY NAME (its study DOI, T and P), listed in the result and the record, and the fit repeated without it, once (`outliers \{ maxRelativeError 15 percent; \}`). ACCEPTANCE, GRADED, the DEFAULT of C285: the maximum relative P error over all kept points is graded clean up to 5 \%, warning up to 10 \%, severe up to 20 \% and refused above; where the normal-boiling-point condition below applies, the absolute gap is graded clean up to 1 K, warning up to 2 K, severe up to 4 K and refused above. The fit's grade is the WORSE of the two. A clean, warning or severe fit is ACCEPTED and promoted (autoPromoted, humanReviewed false) with `fitGrade <grade>;` and the deciding numbers in its provenance entry, which the component loader voices on every run that uses it (a warning or a severe grade); a refused one is not promoted. This case declares no `system/curationDict`, so the default applies; a case that needs another acceptance declares it there, with its reason. The AmbroseWalton curve's errors on the same points are reported beside the fit's, as information.

NORMAL BOILING POINT: no Tb declared in the record: no normal-boiling-point condition.

NO MEASURED POINT IS IN THIS TREE: constant/archiveDatasets declares every series by DOI, NIST URL and sha256; `bin/choupo-curate fetch` regenerates them into constant/archiveExtracts/ (gitignored).

RESULT @@RESULT-BEGIN@@ WHAT THE RUN FOUND (pinned in `expected`; measured, not argued). Objective twoStage. Antoine log10(Psat/bar) = 4.289283 - 1723.6415/(T + -80.3962), in-sample R2 0.999950 over 17 kept points, 273.26-452.3 K (10.1021/je200730m 17/17 kept). Relative error in P over the kept points: mean 2.0440 \%, max 2.7420 \%; AmbroseWalton on the same points: mean 9.6435 \%, max 18.0883 \%. Outliers discarded: none. Its own normal boiling point (101.325 kPa): 482.78 K. PROMOTED on 2026-11-10 into data/standards/components/ethylOctanoate.dat as reviewStatus autoPromoted, humanReviewed false, because: fit grade clean under the DEFAULT graded acceptance of DEV.md 4c C285 (no system/curationDict), 1 studies, 17 points kept, 0 discarded as outliers; max relative error 2.742 \% (mean 2.044 \%), grade clean (limits 5 / 10 / 20 \%); no Tb declared; AmbroseWalton on the same points mean 9.643 \% / max 18.09 \% @@RESULT-END@@

\end{tutnarrative}
\begin{tutgolden}
\goldrow{diag}{fit\_antoine}{A}{4.28928298603}
\goldrow{diag}{fit\_antoine}{B}{1723.64154729}
\goldrow{diag}{fit\_antoine}{C}{-80.396171833}
\goldrow{diag}{fit\_antoine}{R2}{0.999950262952}
\goldrow{diag}{fit\_antoine}{dHvap\_kJ\_per\_mol}{48.807740805}
\goldrow{diag}{fit\_antoine}{minimax\_evaluations}{449}
\goldrow{diag}{fit\_antoine}{minimax\_max\_scaled\_final}{0.027420096796}
\goldrow{diag}{fit\_antoine}{minimax\_max\_scaled\_start}{0.030746589966}
\goldrow{diag}{fit\_antoine}{n\_heldout}{0}
\goldrow{diag}{fit\_antoine}{n\_outliers\_dropped}{0}
\goldrow{diag}{fit\_antoine}{n\_points}{17}
\goldrow{diag}{fit\_antoine}{objective\_method}{2}
\goldrow{diag}{fit\_antoine}{relErr\_max\_fit\_pct}{2.7420096796}
\goldrow{diag}{fit\_antoine}{relErr\_max\_fit\_pct\_leastSquares}{3.0746589966}
\goldmore{3}
\end{tutgolden}

\clearpage
\tutentry{curate168\_vp\_13Propanediol}{tutorials/props/curation/curate168\_vp\_13Propanediol}
\begin{tutscheme}
\begin{tikzpicture}[>={Stealth[length=2mm]}, every node/.style={font=\scriptsize}]
  \node[draw, rounded corners=2pt, align=center, minimum width=3.1cm, minimum height=0.8cm, fill=black!4] at (0.0,-0.0) {\textbf{fit\_antoine}\\[-1pt]{\tiny vaporPressureFit}};
\end{tikzpicture}
\end{tutscheme}
\tutwhat{Antoine vapour pressure of 13Propanediol fitted by the curation pipeline to the measured vapour pressures of every usable laboratory (ThermoML locates, the articles source), outliers discarded by name}
\begin{tutnarrative}
the vapour pressure of 13Propanediol, REFITTED to measurements (DEV.md 4c C281 and C283). Written by `bin/choupo-curate vp 13Propanediol`; everything above the RESULT paragraph is the DECLARATION, committed before any fit ran, so it can be read against the result.

WHY. The catalogue record carried `vaporPressure \{ model AmbroseWalton; \}`, a corresponding-states ESTIMATE from Tc, Pc and omega that never reads the record's own Tb (C279 measured a median 19.6 \% AAD over the 25 such records with archive data). Here a three-parameter Antoine is fitted to EVERY usable laboratory's measurements of it.

STUDIES (one per article, extracted by `choupo-thermoml extract-vp`; ranked by points, ties by DOI): 10.1016/j.jct.2016.01.019 20 pts 414.2-465.45 K Rane; Kumari; Soujanya; Satyavathi 10.1021/je200766t 18 pts 313.2-461.85 K Mokbel; Sawaya; Zanota; Naccoul et al. 10.1021/je060419q 14 pts 293.5-342.3 K Verevkin 10.1016/j.jct.2015.01.014 7 pts 292.5-346.4 K Emel'yanenko; Verevkin SET ASIDE BY NAME: 10.1016/j.jct.2018.12.030: 1 pressure point(s), fewer than 5

FIT AND GRADED ACCEPTANCE (DEV.md 4c C285, Vitor 2026-11-10; the evidence, objective and outliers of C283; declared before the fit). EVIDENCE: every study above is fitted; none is held out (C283 amended: laboratories disagree with each other by more than they claim, so the test is agreement across all of them, not one laboratory's claimed precision). OBJECTIVE: `objective \{ method twoStage; \}` -- least squares in log10 P for a first estimate, then the LARGEST relative error in P, (P\_model - P)/P, over the kept points minimised from it (Vitor's practice, C280); both answers kept. OUTLIERS: after the least-squares stage a point with a relative P error above 15 \% (three times the acceptance) is DISCARDED BY NAME (its study DOI, T and P), listed in the result and the record, and the fit repeated without it, once (`outliers \{ maxRelativeError 15 percent; \}`). ACCEPTANCE, GRADED, the DEFAULT of C285: the maximum relative P error over all kept points is graded clean up to 5 \%, warning up to 10 \%, severe up to 20 \% and refused above; where the normal-boiling-point condition below applies, the absolute gap is graded clean up to 1 K, warning up to 2 K, severe up to 4 K and refused above. The fit's grade is the WORSE of the two. A clean, warning or severe fit is ACCEPTED and promoted (autoPromoted, humanReviewed false) with `fitGrade <grade>;` and the deciding numbers in its provenance entry, which the component loader voices on every run that uses it (a warning or a severe grade); a refused one is not promoted. This case declares no `system/curationDict`, so the default applies; a case that needs another acceptance declares it there, with its reason. The AmbroseWalton curve's errors on the same points are reported beside the fit's, as information.

NORMAL BOILING POINT: no Tb declared in the record: no normal-boiling-point condition.

NO MEASURED POINT IS IN THIS TREE: constant/archiveDatasets declares every series by DOI, NIST URL and sha256; `bin/choupo-curate fetch` regenerates them into constant/archiveExtracts/ (gitignored).

RESULT @@RESULT-BEGIN@@ WHAT THE RUN FOUND (pinned in `expected`; measured, not argued). Objective twoStage. Antoine log10(Psat/bar) = 6.171449 - 2687.6528/(T + -47.7069), in-sample R2 0.999851 over 59 kept points, 292.5-465.45 K (10.1016/j.jct.2016.01.019 20/20 kept; 10.1021/je200766t 18/18 kept; 10.1021/je060419q 14/14 kept; 10.1016/j.jct.2015.01.014 7/7 kept). Relative error in P over the kept points: mean 3.6118 \%, max 7.2487 \%; AmbroseWalton on the same points: mean 141.4260 \%, max 313.4007 \%. Outliers discarded: none. Its own normal boiling point (101.325 kPa): 483.61 K. PROMOTED on 2026-11-10 into data/standards/components/13Propanediol.dat as reviewStatus autoPromoted, humanReviewed false, because: fit grade warning under the DEFAULT graded acceptance of DEV.md 4c C285 (no system/curationDict), 4 studies, 59 points kept, 0 discarded as outliers; max relative error 7.249 \% (mean 3.612 \%), grade warning (limits 5 / 10 / 20 \%); no Tb declared; AmbroseWalton on the same points mean 141.4 \% / max 313.4 \% @@RESULT-END@@

\end{tutnarrative}
\begin{tutgolden}
\goldrow{diag}{fit\_antoine}{A}{6.1714490781}
\goldrow{diag}{fit\_antoine}{B}{2687.65283506}
\goldrow{diag}{fit\_antoine}{C}{-47.7069329261}
\goldrow{diag}{fit\_antoine}{R2}{0.999850829335}
\goldrow{diag}{fit\_antoine}{dHvap\_kJ\_per\_mol}{63.87789719}
\goldrow{diag}{fit\_antoine}{minimax\_evaluations}{468}
\goldrow{diag}{fit\_antoine}{minimax\_max\_scaled\_final}{0.0724869968582}
\goldrow{diag}{fit\_antoine}{minimax\_max\_scaled\_start}{0.0767995654552}
\goldrow{diag}{fit\_antoine}{n\_heldout}{0}
\goldrow{diag}{fit\_antoine}{n\_outliers\_dropped}{0}
\goldrow{diag}{fit\_antoine}{n\_points}{59}
\goldrow{diag}{fit\_antoine}{objective\_method}{2}
\goldrow{diag}{fit\_antoine}{relErr\_max\_fit\_pct}{7.24869968582}
\goldrow{diag}{fit\_antoine}{relErr\_max\_fit\_pct\_leastSquares}{7.67995654552}
\goldmore{3}
\end{tutgolden}

\clearpage
\tutentry{curate169\_vp\_13Butanediol}{tutorials/props/curation/curate169\_vp\_13Butanediol}
\begin{tutscheme}
\begin{tikzpicture}[>={Stealth[length=2mm]}, every node/.style={font=\scriptsize}]
  \node[draw, rounded corners=2pt, align=center, minimum width=3.1cm, minimum height=0.8cm, fill=black!4] at (0.0,-0.0) {\textbf{fit\_antoine}\\[-1pt]{\tiny vaporPressureFit}};
\end{tikzpicture}
\end{tutscheme}
\tutwhat{Antoine vapour pressure of 13Butanediol fitted by the curation pipeline to the measured vapour pressures of every usable laboratory (ThermoML locates, the articles source), outliers discarded by name}
\begin{tutnarrative}
the vapour pressure of 13Butanediol, REFITTED to measurements (DEV.md 4c C281 and C283). Written by `bin/choupo-curate vp 13Butanediol`; everything above the RESULT paragraph is the DECLARATION, committed before any fit ran, so it can be read against the result.

WHY. The catalogue record carried `vaporPressure \{ model AmbroseWalton; \}`, a corresponding-states ESTIMATE from Tc, Pc and omega that never reads the record's own Tb (C279 measured a median 19.6 \% AAD over the 25 such records with archive data). Here a three-parameter Antoine is fitted to EVERY usable laboratory's measurements of it.

STUDIES (one per article, extracted by `choupo-thermoml extract-vp`; ranked by points, ties by DOI): 10.1021/je400082c 22 pts 387.75-453.1 K Zhang; Wu; Sun; Li et al. 10.1021/je060419q 13 pts 288.3-332.3 K Verevkin 10.1021/je060269j 9 pts 481.4-688.2 K VonNiederhausern; Wilson; Giles SET ASIDE BY NAME: (none)

FIT AND GRADED ACCEPTANCE (DEV.md 4c C285, Vitor 2026-11-10; the evidence, objective and outliers of C283; declared before the fit). EVIDENCE: every study above is fitted; none is held out (C283 amended: laboratories disagree with each other by more than they claim, so the test is agreement across all of them, not one laboratory's claimed precision). OBJECTIVE: `objective \{ method twoStage; \}` -- least squares in log10 P for a first estimate, then the LARGEST relative error in P, (P\_model - P)/P, over the kept points minimised from it (Vitor's practice, C280); both answers kept. OUTLIERS: after the least-squares stage a point with a relative P error above 15 \% (three times the acceptance) is DISCARDED BY NAME (its study DOI, T and P), listed in the result and the record, and the fit repeated without it, once (`outliers \{ maxRelativeError 15 percent; \}`). ACCEPTANCE, GRADED, the DEFAULT of C285: the maximum relative P error over all kept points is graded clean up to 5 \%, warning up to 10 \%, severe up to 20 \% and refused above; where the normal-boiling-point condition below applies, the absolute gap is graded clean up to 1 K, warning up to 2 K, severe up to 4 K and refused above. The fit's grade is the WORSE of the two. A clean, warning or severe fit is ACCEPTED and promoted (autoPromoted, humanReviewed false) with `fitGrade <grade>;` and the deciding numbers in its provenance entry, which the component loader voices on every run that uses it (a warning or a severe grade); a refused one is not promoted. This case declares no `system/curationDict`, so the default applies; a case that needs another acceptance declares it there, with its reason. The AmbroseWalton curve's errors on the same points are reported beside the fit's, as information.

NORMAL BOILING POINT: no Tb declared in the record: no normal-boiling-point condition.

NO MEASURED POINT IS IN THIS TREE: constant/archiveDatasets declares every series by DOI, NIST URL and sha256; `bin/choupo-curate fetch` regenerates them into constant/archiveExtracts/ (gitignored).

RESULT @@RESULT-BEGIN@@ WHAT THE RUN FOUND (pinned in `expected`; measured, not argued). Objective twoStage. Antoine log10(Psat/bar) = 5.022835 - 1987.1123/(T + -87.7092), in-sample R2 0.999850 over 44 kept points, 288.3-688.2 K (10.1021/je400082c 22/22 kept; 10.1021/je060419q 13/13 kept; 10.1021/je060269j 9/9 kept). Relative error in P over the kept points: mean 4.7097 \%, max 7.1995 \%; AmbroseWalton on the same points: mean 4.9765 \%, max 10.8320 \%. Outliers discarded: none. Its own normal boiling point (101.325 kPa): 483.78 K. PROMOTED on 2026-11-10 into data/standards/components/13Butanediol.dat as reviewStatus autoPromoted, humanReviewed false, because: fit grade warning under the DEFAULT graded acceptance of DEV.md 4c C285 (no system/curationDict), 3 studies, 44 points kept, 0 discarded as outliers; max relative error 7.199 \% (mean 4.71 \%), grade warning (limits 5 / 10 / 20 \%); no Tb declared; AmbroseWalton on the same points mean 4.977 \% / max 10.83 \% @@RESULT-END@@

\end{tutnarrative}
\begin{tutgolden}
\goldrow{diag}{fit\_antoine}{A}{5.02283492679}
\goldrow{diag}{fit\_antoine}{B}{1987.11232708}
\goldrow{diag}{fit\_antoine}{C}{-87.709183589}
\goldrow{diag}{fit\_antoine}{R2}{0.999850457607}
\goldrow{diag}{fit\_antoine}{dHvap\_kJ\_per\_mol}{49.9676457582}
\goldrow{diag}{fit\_antoine}{minimax\_evaluations}{439}
\goldrow{diag}{fit\_antoine}{minimax\_max\_scaled\_final}{0.0719948561483}
\goldrow{diag}{fit\_antoine}{minimax\_max\_scaled\_start}{0.103947204459}
\goldrow{diag}{fit\_antoine}{n\_heldout}{0}
\goldrow{diag}{fit\_antoine}{n\_outliers\_dropped}{0}
\goldrow{diag}{fit\_antoine}{n\_points}{44}
\goldrow{diag}{fit\_antoine}{objective\_method}{2}
\goldrow{diag}{fit\_antoine}{relErr\_max\_fit\_pct}{7.19948561483}
\goldrow{diag}{fit\_antoine}{relErr\_max\_fit\_pct\_leastSquares}{10.3947204459}
\goldmore{3}
\end{tutgolden}

\clearpage
\tutentry{curate16\_nrtl\_ethanol\_dimethylCarbonate}{tutorials/props/curation/curate16\_nrtl\_ethanol\_dimethylCarbonate}
\begin{tutscheme}
\begin{tikzpicture}[>={Stealth[length=2mm]}, every node/.style={font=\scriptsize}]
  \node[draw, rounded corners=2pt, align=center, minimum width=3.1cm, minimum height=0.8cm, fill=black!4] at (0.0,-0.0) {\textbf{fit\_nrtl}\\[-1pt]{\tiny fitParameters}};
  \node[draw, rounded corners=2pt, align=center, minimum width=3.1cm, minimum height=0.8cm, fill=black!4] at (3.6,-0.0) {\textbf{fit\_nrtl\_start2}\\[-1pt]{\tiny fitParameters}};
  \node[draw, rounded corners=2pt, align=center, minimum width=3.1cm, minimum height=0.8cm, fill=black!4] at (7.2,-0.0) {\textbf{txy\_ideal}\\[-1pt]{\tiny propertyScan1D}};
  \node[draw, rounded corners=2pt, align=center, minimum width=3.1cm, minimum height=0.8cm, fill=black!4] at (0.0,-1.3) {\textbf{txy\_nrtl}\\[-1pt]{\tiny propertyScan1D}};
  \node[draw, rounded corners=2pt, align=center, minimum width=3.1cm, minimum height=0.8cm, fill=black!4] at (3.6,-1.3) {\textbf{txy\_ideal\_fitP}\\[-1pt]{\tiny propertyScan1D}};
  \node[draw, rounded corners=2pt, align=center, minimum width=3.1cm, minimum height=0.8cm, fill=black!4] at (7.2,-1.3) {\textbf{txy\_nrtl\_fitP}\\[-1pt]{\tiny propertyScan1D}};
\end{tikzpicture}
\end{tutscheme}
\tutwhat{NRTL ethanol-dimethylCarbonate regressed by the curation pipeline to measured bubble temperatures (ThermoML locates, the article sources), held out against an independent study}
\begin{tutnarrative}
AN NRTL PAIR FROM PUBLISHED MEASUREMENTS, BY THE PIPELINE.

WRITTEN BY bin/choupo-curate pair ethanol dimethylCarbonate (DEV.md 4c C282; record docs/design/binary-pairs-from-open-measurements.md section 8). Every rule below was decided before this tool existed (C16, C225, C226); the tool applies them and decides nothing.

@@PILOT-BEGIN@@ WHY THIS PAIR -- THE PILOT RANKING (C282; added after the run, but the ranking reads ONLY the archive index and was computed before any fit). Candidates: both components in the catalogue, NO record for the pair in any model directory of data/standards/parameters/, no acetic acid (vapour dimerisation, C16), no component whose vapour pressure is AmbroseWalton (C281), and (checked for the five chosen) no binary block in the archive with a second liquid. Ranked by the number of mutually author-disjoint isobaric bubble-T studies at or below 110 kPa with a declared uncertainty, then by their total points: 1 methanol + dimethylCarbonate 5 independent (9 studies, 274 pts) - acetone + methanol 4 -- passed over: curate05 holds its C16 regression and the two studies' disagreement, a curator's tie-break 2 ethanol + dimethylCarbonate 4 (4 studies, 156 pts) 3 diisopropylEther + isopropanol 3 (4 studies, 152 pts) - triethyleneGlycol + water 3 -- REFUSED by the tool: 2 of 3 studies disagree beyond their U 4 isopropanol + dimethylCarbonate 3 (3 studies, 128 pts) - 2Butanol + dimethylCarbonate 3 -- REFUSED by the tool: no author-disjoint held-out study reports x of the same component 5 ethanol + 224Trimethylpentane 3 (5 studies, 97 pts) THIS CASE: rank 2. @@PILOT-END@@

THE DECLARATION -- WRITTEN AND COMMITTED BEFORE ANY FIT RAN. (C16 slice 2: so the declaration can be read against the result.)

IDENTITY. ethanol (CAS 64-17-5, InChIKey LFQSCWFLJHTTHZ-UHFFFAOYSA-N) + dimethylCarbonate (CAS 616-38-6, InChIKey IEJIGPNLZYLLBP-UHFFFAOYSA-N), joined to the archive through identity.json -- never by name.

NO MEASURED SERIES IN THIS TREE. constant/archiveDatasets names each dataset (DOI, NIST URL, article sha256, extraction selector, extract sha256); `bin/choupo-curate fetch` writes the points into constant/archiveExtracts/, which is gitignored.

STUDIES CONSIDERED: every bubble-temperature block of the pair at or below 110 kPa with a declared uncertainty (4 block(s)). USABLE after the screens: 2. SET ASIDE, BY NAME: - doi 10.1021/je0601309 block 9: TRIAGE: disagrees with the other studies beyond their stated uncertainties (leave-one-study-out, below) - doi 10.1016/j.fluid.2018.05.024 block 0: TRIAGE: disagrees with the other studies beyond their stated uncertainties (leave-one-study-out, below)

TRIAGE (C226 item 3, leave one study out, fitted by the engine; a study whose held-out AAD exceeds its own largest stated U plus the others' largest stated U disagrees with them beyond their stated uncertainties): 10.1021/je0601309: AAD 0.896 K against the limit 0.800 K $\to$ DISAGREES -- set aside; 10.1021/acs.jced.9b00414: AAD 0.704 K against the limit 0.830 K $\to$ agrees; 10.1016/j.fluid.2018.05.024: AAD 0.945 K against the limit 0.750 K $\to$ DISAGREES -- set aside; 10.1021/je030131q: AAD 0.673 K against the limit 0.830 K $\to$ agrees

PARTITION (the declared rule of C16 slice 2): the FIT set is the usable study with the most interior points; the HELD-OUT set is the next largest that shares NO author with it and was measured at one pressure. FIT doi 10.1021/acs.jced.9b00414 (20 interior points, 101.3 kPa) authors: Liu, K.[Kai]; Xu, X.[Xin]; Zhang, Y.[Yi]; Xu, D.[Dongmei]; Gao, J.[Jun]; Zhang, L.[Lianzheng] HELD OUT doi 10.1021/je030131q (9 interior points, 95.8 kPa) authors: Vittal Prasad, T. E.; Prashanth, P. R.; Kumar, D. S.; Prasad, D.H.L.

BAND: maxAAD 0.24 \% of the held-out mean T. declared before the fit by bin/choupo-curate, the ONE rule of C16: the held-out set's largest stated expanded uncertainty (U95 = 0.32 K) plus the catalogue vapour-pressure model's own error at that set's pure endpoints (dimethylCarbonate +0.504 K; measured with no pair involved), as a percentage of the held-out mean T (347.42 K) = 0.2371 \%, rounded UP to two significant figures

THE FIT. NRTL with a\_ij = a\_ji = 0 and alpha = 0.30 fixed: TWO parameters (b\_ij, b\_ji), the C16 frame (bubble temperatures over these ranges do not identify more). Each point weighted by 1/U, U its own stated expanded uncertainty. Two starts: (0, 0) and (800, -400).

@@RESULT-BEGIN@@ WHAT THE RUN FOUND (pinned in `expected`; measured, not argued). fit\_nrtl: b\_ij = 427.6 K, b\_ji = -4.2 K, in-sample rms 0.029 K over 20 points; fit\_nrtl\_start2 lands on the same pair (427.6, -4.2). HELD OUT: AAD 1.184 K (0.3406 \%) against the band 0.24 \% $\to$ notValidated. NOT PROMOTABLE: held-out verdict notValidated (AAD 0.340557 \% against the band 0.24 \% declared before the fit) @@RESULT-END@@

WHAT IS NOT CLAIMED. The pair is correct at these pressures and in this frame (ideal-gas vapour, the catalogue's own vapour-pressure records); the measurements were transcribed by the archive and have NOT been read back against the articles (the CHECKLIST that `fetch` writes beside the extracts is the aid for that reading).

\end{tutnarrative}
\begin{tutgolden}
\goldrow{diag}{fit\_nrtl}{aad\_heldout\_K}{1.18370705596}
\goldrow{diag}{fit\_nrtl}{aad\_heldout\_pct}{0.340556701675}
\goldrow{diag}{fit\_nrtl}{bubble\_aad\_fit}{0.0171105070697}
\goldrow{diag}{fit\_nrtl}{bubble\_aad\_heldout}{1.18370705596}
\goldrow{diag}{fit\_nrtl}{bubble\_max\_fit}{0.0930111024659}
\goldrow{diag}{fit\_nrtl}{bubble\_max\_heldout}{3.91117758308}
\goldrow{diag}{fit\_nrtl}{chi2}{0.0995348764533}
\goldrow{diag}{fit\_nrtl}{chi2\_in\_sample}{0.0995348764533}
\goldrow{diag}{fit\_nrtl}{chi2\_reduced}{0.00552971535852}
\goldrow{diag}{fit\_nrtl}{cond\_JtJ}{49.2713871546}
\goldrow{diag}{fit\_nrtl}{converged}{1}
\goldrow{diag}{fit\_nrtl}{corr\_computed}{1}
\goldrow{diag}{fit\_nrtl}{dof}{18}
\goldrow{diag}{fit\_nrtl}{identifiable}{1}
\goldmore{66}
\end{tutgolden}

\clearpage
\tutentry{curate171\_vp\_16Hexanediol}{tutorials/props/curation/curate171\_vp\_16Hexanediol}
\begin{tutscheme}
\begin{tikzpicture}[>={Stealth[length=2mm]}, every node/.style={font=\scriptsize}]
  \node[draw, rounded corners=2pt, align=center, minimum width=3.1cm, minimum height=0.8cm, fill=black!4] at (0.0,-0.0) {\textbf{fit\_antoine}\\[-1pt]{\tiny vaporPressureFit}};
\end{tikzpicture}
\end{tutscheme}
\tutwhat{Antoine vapour pressure of 16Hexanediol fitted by the curation pipeline to the measured vapour pressures of every usable laboratory (ThermoML locates, the articles source), outliers discarded by name}
\begin{tutnarrative}
the vapour pressure of 16Hexanediol, REFITTED to measurements (DEV.md 4c C281 and C283). Written by `bin/choupo-curate vp 16Hexanediol`; everything above the RESULT paragraph is the DECLARATION, committed before any fit ran, so it can be read against the result.

WHY. The catalogue record carried `vaporPressure \{ model AmbroseWalton; \}`, a corresponding-states ESTIMATE from Tc, Pc and omega that never reads the record's own Tb (C279 measured a median 19.6 \% AAD over the 25 such records with archive data). Here a three-parameter Antoine is fitted to EVERY usable laboratory's measurements of it.

STUDIES (one per article, extracted by `choupo-thermoml extract-vp`; ranked by points, ties by DOI): 10.1021/je0602465 8 pts 523.2-740.4 K VonNiederhausern; Wilson; Giles SET ASIDE BY NAME: 10.1021/acs.jced.9b00306: 3 pressure point(s), fewer than 5

FIT AND GRADED ACCEPTANCE (DEV.md 4c C285, Vitor 2026-11-10; the evidence, objective and outliers of C283; declared before the fit). EVIDENCE: every study above is fitted; none is held out (C283 amended: laboratories disagree with each other by more than they claim, so the test is agreement across all of them, not one laboratory's claimed precision). OBJECTIVE: `objective \{ method twoStage; \}` -- least squares in log10 P for a first estimate, then the LARGEST relative error in P, (P\_model - P)/P, over the kept points minimised from it (Vitor's practice, C280); both answers kept. OUTLIERS: after the least-squares stage a point with a relative P error above 15 \% (three times the acceptance) is DISCARDED BY NAME (its study DOI, T and P), listed in the result and the record, and the fit repeated without it, once (`outliers \{ maxRelativeError 15 percent; \}`). ACCEPTANCE, GRADED, the DEFAULT of C285: the maximum relative P error over all kept points is graded clean up to 5 \%, warning up to 10 \%, severe up to 20 \% and refused above; where the normal-boiling-point condition below applies, the absolute gap is graded clean up to 1 K, warning up to 2 K, severe up to 4 K and refused above. The fit's grade is the WORSE of the two. A clean, warning or severe fit is ACCEPTED and promoted (autoPromoted, humanReviewed false) with `fitGrade <grade>;` and the deciding numbers in its provenance entry, which the component loader voices on every run that uses it (a warning or a severe grade); a refused one is not promoted. This case declares no `system/curationDict`, so the default applies; a case that needs another acceptance declares it there, with its reason. The AmbroseWalton curve's errors on the same points are reported beside the fit's, as information.

NORMAL BOILING POINT: no Tb declared in the record: no normal-boiling-point condition.

NO MEASURED POINT IS IN THIS TREE: constant/archiveDatasets declares every series by DOI, NIST URL and sha256; `bin/choupo-curate fetch` regenerates them into constant/archiveExtracts/ (gitignored).

RESULT @@RESULT-BEGIN@@ WHAT THE RUN FOUND (pinned in `expected`; measured, not argued). Objective twoStage. Antoine log10(Psat/bar) = 4.626258 - 1878.9662/(T + -119.1422), in-sample R2 0.999928 over 8 kept points, 523.2-740.4 K (10.1021/je0602465 8/8 kept). Relative error in P over the kept points: mean 1.0839 \%, max 1.3903 \%; AmbroseWalton on the same points: mean 3.3824 \%, max 5.4812 \%. Outliers discarded: none. Its own normal boiling point (101.325 kPa): 525.80 K. PROMOTED on 2026-11-10 into data/standards/components/16Hexanediol.dat as reviewStatus autoPromoted, humanReviewed false, because: fit grade clean under the DEFAULT graded acceptance of DEV.md 4c C285 (no system/curationDict), 1 studies, 8 points kept, 0 discarded as outliers; max relative error 1.39 \% (mean 1.084 \%), grade clean (limits 5 / 10 / 20 \%); no Tb declared; AmbroseWalton on the same points mean 3.382 \% / max 5.481 \% @@RESULT-END@@

\end{tutnarrative}
\begin{tutgolden}
\goldrow{diag}{fit\_antoine}{A}{4.62625817626}
\goldrow{diag}{fit\_antoine}{B}{1878.96615467}
\goldrow{diag}{fit\_antoine}{C}{-119.142166127}
\goldrow{diag}{fit\_antoine}{R2}{0.99992757219}
\goldrow{diag}{fit\_antoine}{dHvap\_kJ\_per\_mol}{51.092587228}
\goldrow{diag}{fit\_antoine}{minimax\_evaluations}{457}
\goldrow{diag}{fit\_antoine}{minimax\_max\_scaled\_final}{0.0139027114429}
\goldrow{diag}{fit\_antoine}{minimax\_max\_scaled\_start}{0.0159712042486}
\goldrow{diag}{fit\_antoine}{n\_heldout}{0}
\goldrow{diag}{fit\_antoine}{n\_outliers\_dropped}{0}
\goldrow{diag}{fit\_antoine}{n\_points}{8}
\goldrow{diag}{fit\_antoine}{objective\_method}{2}
\goldrow{diag}{fit\_antoine}{relErr\_max\_fit\_pct}{1.39027114429}
\goldrow{diag}{fit\_antoine}{relErr\_max\_fit\_pct\_leastSquares}{1.59712042486}
\goldmore{3}
\end{tutgolden}

\clearpage
\tutentry{curate172\_vp\_allylAlcohol}{tutorials/props/curation/curate172\_vp\_allylAlcohol}
\begin{tutscheme}
\begin{tikzpicture}[>={Stealth[length=2mm]}, every node/.style={font=\scriptsize}]
  \node[draw, rounded corners=2pt, align=center, minimum width=3.1cm, minimum height=0.8cm, fill=black!4] at (0.0,-0.0) {\textbf{fit\_antoine}\\[-1pt]{\tiny vaporPressureFit}};
\end{tikzpicture}
\end{tutscheme}
\tutwhat{Antoine vapour pressure of allylAlcohol fitted by the curation pipeline to the measured vapour pressures of every usable laboratory (ThermoML locates, the articles source), outliers discarded by name}
\begin{tutnarrative}
the vapour pressure of allylAlcohol, REFITTED to measurements (DEV.md 4c C281 and C283). Written by `bin/choupo-curate vp allylAlcohol`; everything above the RESULT paragraph is the DECLARATION, committed before any fit ran, so it can be read against the result.

WHY. The catalogue record carried `vaporPressure \{ model AmbroseWalton; \}`, a corresponding-states ESTIMATE from Tc, Pc and omega that never reads the record's own Tb (C279 measured a median 19.6 \% AAD over the 25 such records with archive data). Here a three-parameter Antoine is fitted to EVERY usable laboratory's measurements of it.

STUDIES (one per article, extracted by `choupo-thermoml extract-vp`; ranked by points, ties by DOI): 10.1021/je0499519 17 pts 311.42-355.7 K Lubomska; Malanowski 10.1021/je060088h 6 pts 438.1-546.5 K VonNiederhausern; Wilson; Giles 10.1021/je400885z 6 pts 373.15-423.15 K McDougal; Jasperson; Wilson SET ASIDE BY NAME: 10.1021/acs.jced.6b00893: 1 pressure point(s), fewer than 5 10.1021/acs.jced.8b00908: 1 pressure point(s), fewer than 5

FIT AND GRADED ACCEPTANCE (DEV.md 4c C285, Vitor 2026-11-10; the evidence, objective and outliers of C283; declared before the fit). EVIDENCE: every study above is fitted; none is held out (C283 amended: laboratories disagree with each other by more than they claim, so the test is agreement across all of them, not one laboratory's claimed precision). OBJECTIVE: `objective \{ method twoStage; \}` -- least squares in log10 P for a first estimate, then the LARGEST relative error in P, (P\_model - P)/P, over the kept points minimised from it (Vitor's practice, C280); both answers kept. OUTLIERS: after the least-squares stage a point with a relative P error above 15 \% (three times the acceptance) is DISCARDED BY NAME (its study DOI, T and P), listed in the result and the record, and the fit repeated without it, once (`outliers \{ maxRelativeError 15 percent; \}`). ACCEPTANCE, GRADED, the DEFAULT of C285: the maximum relative P error over all kept points is graded clean up to 5 \%, warning up to 10 \%, severe up to 20 \% and refused above; where the normal-boiling-point condition below applies, the absolute gap is graded clean up to 1 K, warning up to 2 K, severe up to 4 K and refused above. The fit's grade is the WORSE of the two. A clean, warning or severe fit is ACCEPTED and promoted (autoPromoted, humanReviewed false) with `fitGrade <grade>;` and the deciding numbers in its provenance entry, which the component loader voices on every run that uses it (a warning or a severe grade); a refused one is not promoted. This case declares no `system/curationDict`, so the default applies; a case that needs another acceptance declares it there, with its reason. The AmbroseWalton curve's errors on the same points are reported beside the fit's, as information.

NORMAL BOILING POINT: the record's Tb 369.925 K is judged against the range of the KEPT points (all studies span 311.42-546.5 K before any outlier is discarded): if Tb lies OUTSIDE that range, the fitted Antoine's own normal boiling point (P = 101.325 kPa) is GRADED by its absolute gap to it (the temperature limits of the acceptance below); inside it, the gap is reported and is not a condition.

NO MEASURED POINT IS IN THIS TREE: constant/archiveDatasets declares every series by DOI, NIST URL and sha256; `bin/choupo-curate fetch` regenerates them into constant/archiveExtracts/ (gitignored).

RESULT @@RESULT-BEGIN@@ WHAT THE RUN FOUND (pinned in `expected`; measured, not argued). Objective twoStage. Antoine log10(Psat/bar) = 4.704344 - 1398.8210/(T + -72.7855), in-sample R2 0.999985 over 29 kept points, 311.42-546.5 K (10.1021/je0499519 17/17 kept; 10.1021/je060088h 6/6 kept; 10.1021/je400885z 6/6 kept). Relative error in P over the kept points: mean 0.6417 \%, max 1.4168 \%; AmbroseWalton on the same points: mean 2.9021 \%, max 6.3822 \%. Outliers discarded: none. Its own normal boiling point (101.325 kPa): 370.49 K. PROMOTED on 2026-11-10 into data/standards/components/allylAlcohol.dat as reviewStatus autoPromoted, humanReviewed false, because: fit grade clean under the DEFAULT graded acceptance of DEV.md 4c C285 (no system/curationDict), 3 studies, 29 points kept, 0 discarded as outliers; max relative error 1.417 \% (mean 0.6417 \%), grade clean (limits 5 / 10 / 20 \%); Tb inside the kept points' range, normal boiling point 370.49 K, +0.57 K from the record's Tb 369.925 K (reported, not a condition); AmbroseWalton on the same points mean 2.902 \% / max 6.382 \% @@RESULT-END@@

\end{tutnarrative}
\begin{tutgolden}
\goldrow{diag}{fit\_antoine}{A}{4.70434351854}
\goldrow{diag}{fit\_antoine}{B}{1398.82103022}
\goldrow{diag}{fit\_antoine}{C}{-72.7854693401}
\goldrow{diag}{fit\_antoine}{R2}{0.999985397148}
\goldrow{diag}{fit\_antoine}{dHvap\_kJ\_per\_mol}{35.6417439945}
\goldrow{diag}{fit\_antoine}{minimax\_evaluations}{336}
\goldrow{diag}{fit\_antoine}{minimax\_max\_scaled\_final}{0.0141678622674}
\goldrow{diag}{fit\_antoine}{minimax\_max\_scaled\_start}{0.0202035867299}
\goldrow{diag}{fit\_antoine}{n\_heldout}{0}
\goldrow{diag}{fit\_antoine}{n\_outliers\_dropped}{0}
\goldrow{diag}{fit\_antoine}{n\_points}{29}
\goldrow{diag}{fit\_antoine}{objective\_method}{2}
\goldrow{diag}{fit\_antoine}{relErr\_max\_fit\_pct}{1.41678622674}
\goldrow{diag}{fit\_antoine}{relErr\_max\_fit\_pct\_leastSquares}{2.02035867299}
\goldmore{3}
\end{tutgolden}

\clearpage
\tutentry{curate174\_vp\_33Oxydipropanol}{tutorials/props/curation/curate174\_vp\_33Oxydipropanol}
\begin{tutscheme}
\begin{tikzpicture}[>={Stealth[length=2mm]}, every node/.style={font=\scriptsize}]
  \node[draw, rounded corners=2pt, align=center, minimum width=3.1cm, minimum height=0.8cm, fill=black!4] at (0.0,-0.0) {\textbf{fit\_antoine}\\[-1pt]{\tiny vaporPressureFit}};
\end{tikzpicture}
\end{tutscheme}
\tutwhat{Antoine vapour pressure of 33Oxydipropanol fitted by the curation pipeline to the measured vapour pressures of every usable laboratory (ThermoML locates, the articles source), outliers discarded by name}
\begin{tutnarrative}
the vapour pressure of 33Oxydipropanol, REFITTED to measurements (DEV.md 4c C281 and C283). Written by `bin/choupo-curate vp 33Oxydipropanol`; everything above the RESULT paragraph is the DECLARATION, committed before any fit ran, so it can be read against the result.

WHY. The catalogue record carried `vaporPressure \{ model AmbroseWalton; \}`, a corresponding-states ESTIMATE from Tc, Pc and omega that never reads the record's own Tb (C279 measured a median 19.6 \% AAD over the 25 such records with archive data). Here a three-parameter Antoine is fitted to EVERY usable laboratory's measurements of it.

STUDIES (one per article, extracted by `choupo-thermoml extract-vp`; ranked by points, ties by DOI): 10.1021/je4009014 13 pts 365.15-397.15 K Fendu; Oprea 10.1021/je060269j 8 pts 505-703.4 K VonNiederhausern; Wilson; Giles SET ASIDE BY NAME: (none)

FIT AND GRADED ACCEPTANCE (DEV.md 4c C285, Vitor 2026-11-10; the evidence, objective and outliers of C283; declared before the fit). EVIDENCE: every study above is fitted; none is held out (C283 amended: laboratories disagree with each other by more than they claim, so the test is agreement across all of them, not one laboratory's claimed precision). OBJECTIVE: `objective \{ method twoStage; \}` -- least squares in log10 P for a first estimate, then the LARGEST relative error in P, (P\_model - P)/P, over the kept points minimised from it (Vitor's practice, C280); both answers kept. OUTLIERS: after the least-squares stage a point with a relative P error above 15 \% (three times the acceptance) is DISCARDED BY NAME (its study DOI, T and P), listed in the result and the record, and the fit repeated without it, once (`outliers \{ maxRelativeError 15 percent; \}`). ACCEPTANCE, GRADED, the DEFAULT of C285: the maximum relative P error over all kept points is graded clean up to 5 \%, warning up to 10 \%, severe up to 20 \% and refused above; where the normal-boiling-point condition below applies, the absolute gap is graded clean up to 1 K, warning up to 2 K, severe up to 4 K and refused above. The fit's grade is the WORSE of the two. A clean, warning or severe fit is ACCEPTED and promoted (autoPromoted, humanReviewed false) with `fitGrade <grade>;` and the deciding numbers in its provenance entry, which the component loader voices on every run that uses it (a warning or a severe grade); a refused one is not promoted. This case declares no `system/curationDict`, so the default applies; a case that needs another acceptance declares it there, with its reason. The AmbroseWalton curve's errors on the same points are reported beside the fit's, as information.

NORMAL BOILING POINT: no Tb declared in the record: no normal-boiling-point condition.

NO MEASURED POINT IS IN THIS TREE: constant/archiveDatasets declares every series by DOI, NIST URL and sha256; `bin/choupo-curate fetch` regenerates them into constant/archiveExtracts/ (gitignored).

RESULT @@RESULT-BEGIN@@ WHAT THE RUN FOUND (pinned in `expected`; measured, not argued). Objective twoStage. Antoine log10(Psat/bar) = 4.843255 - 2057.0652/(T + -81.5527), in-sample R2 0.999922 over 21 kept points, 365.15-703.4 K (10.1021/je4009014 13/13 kept; 10.1021/je060269j 8/8 kept). Relative error in P over the kept points: mean 2.7003 \%, max 4.2239 \%; AmbroseWalton on the same points: mean 4.5240 \%, max 11.4663 \%. Outliers discarded: none. Its own normal boiling point (101.325 kPa): 506.78 K. PROMOTED on 2026-11-10 into data/standards/components/33Oxydipropanol.dat as reviewStatus autoPromoted, humanReviewed false, because: fit grade clean under the DEFAULT graded acceptance of DEV.md 4c C285 (no system/curationDict), 2 studies, 21 points kept, 0 discarded as outliers; max relative error 4.224 \% (mean 2.7 \%), grade clean (limits 5 / 10 / 20 \%); no Tb declared; AmbroseWalton on the same points mean 4.524 \% / max 11.47 \% @@RESULT-END@@

\end{tutnarrative}
\begin{tutgolden}
\goldrow{diag}{fit\_antoine}{A}{4.84325496907}
\goldrow{diag}{fit\_antoine}{B}{2057.06515624}
\goldrow{diag}{fit\_antoine}{C}{-81.5526761907}
\goldrow{diag}{fit\_antoine}{R2}{0.999922439834}
\goldrow{diag}{fit\_antoine}{dHvap\_kJ\_per\_mol}{50.3889269926}
\goldrow{diag}{fit\_antoine}{minimax\_evaluations}{420}
\goldrow{diag}{fit\_antoine}{minimax\_max\_scaled\_final}{0.0422387601092}
\goldrow{diag}{fit\_antoine}{minimax\_max\_scaled\_start}{0.0500632719709}
\goldrow{diag}{fit\_antoine}{n\_heldout}{0}
\goldrow{diag}{fit\_antoine}{n\_outliers\_dropped}{0}
\goldrow{diag}{fit\_antoine}{n\_points}{21}
\goldrow{diag}{fit\_antoine}{objective\_method}{2}
\goldrow{diag}{fit\_antoine}{relErr\_max\_fit\_pct}{4.22387601092}
\goldrow{diag}{fit\_antoine}{relErr\_max\_fit\_pct\_leastSquares}{5.00632719709}
\goldmore{3}
\end{tutgolden}

\clearpage
\tutentry{curate176\_vp\_methylLevulinate}{tutorials/props/curation/curate176\_vp\_methylLevulinate}
\begin{tutscheme}
\begin{tikzpicture}[>={Stealth[length=2mm]}, every node/.style={font=\scriptsize}]
  \node[draw, rounded corners=2pt, align=center, minimum width=3.1cm, minimum height=0.8cm, fill=black!4] at (0.0,-0.0) {\textbf{fit\_antoine}\\[-1pt]{\tiny vaporPressureFit}};
\end{tikzpicture}
\end{tutscheme}
\tutwhat{Antoine vapour pressure of methylLevulinate fitted by the curation pipeline to the measured vapour pressures of every usable laboratory (ThermoML locates, the articles source), outliers discarded by name}
\begin{tutnarrative}
the vapour pressure of methylLevulinate, REFITTED to measurements (DEV.md 4c C281 and C283). Written by `bin/choupo-curate vp methylLevulinate`; everything above the RESULT paragraph is the DECLARATION, committed before any fit ran, so it can be read against the result.

WHY. The catalogue record carried `vaporPressure \{ model AmbroseWalton; \}`, a corresponding-states ESTIMATE from Tc, Pc and omega that never reads the record's own Tb (C279 measured a median 19.6 \% AAD over the 25 such records with archive data). Here a three-parameter Antoine is fitted to EVERY usable laboratory's measurements of it.

STUDIES (one per article, extracted by `choupo-thermoml extract-vp`; ranked by points, ties by DOI): 10.1016/j.jct.2013.05.025 25 pts 333.15-453.15 K Lomba; Lafuente; Garcia-Mardones; Gascon et al. 10.1016/j.tca.2017.12.006 13 pts 293.5-353.2 K Emel'yanenko; Altuntepe; Held; Pimerzin et al. SET ASIDE BY NAME: (none)

FIT AND GRADED ACCEPTANCE (DEV.md 4c C285, Vitor 2026-11-10; the evidence, objective and outliers of C283; declared before the fit). EVIDENCE: every study above is fitted; none is held out (C283 amended: laboratories disagree with each other by more than they claim, so the test is agreement across all of them, not one laboratory's claimed precision). OBJECTIVE: `objective \{ method twoStage; \}` -- least squares in log10 P for a first estimate, then the LARGEST relative error in P, (P\_model - P)/P, over the kept points minimised from it (Vitor's practice, C280); both answers kept. OUTLIERS: after the least-squares stage a point with a relative P error above 15 \% (three times the acceptance) is DISCARDED BY NAME (its study DOI, T and P), listed in the result and the record, and the fit repeated without it, once (`outliers \{ maxRelativeError 15 percent; \}`). ACCEPTANCE, GRADED, the DEFAULT of C285: the maximum relative P error over all kept points is graded clean up to 5 \%, warning up to 10 \%, severe up to 20 \% and refused above; where the normal-boiling-point condition below applies, the absolute gap is graded clean up to 1 K, warning up to 2 K, severe up to 4 K and refused above. The fit's grade is the WORSE of the two. A clean, warning or severe fit is ACCEPTED and promoted (autoPromoted, humanReviewed false) with `fitGrade <grade>;` and the deciding numbers in its provenance entry, which the component loader voices on every run that uses it (a warning or a severe grade); a refused one is not promoted. This case declares no `system/curationDict`, so the default applies; a case that needs another acceptance declares it there, with its reason. The AmbroseWalton curve's errors on the same points are reported beside the fit's, as information.

NORMAL BOILING POINT: no Tb declared in the record: no normal-boiling-point condition.

NO MEASURED POINT IS IN THIS TREE: constant/archiveDatasets declares every series by DOI, NIST URL and sha256; `bin/choupo-curate fetch` regenerates them into constant/archiveExtracts/ (gitignored).

RESULT @@RESULT-BEGIN@@ WHAT THE RUN FOUND (pinned in `expected`; measured, not argued). Objective twoStage. Antoine log10(Psat/bar) = 4.618303 - 1844.8759/(T + -67.5116), in-sample R2 0.999817 over 38 kept points, 293.5-453.15 K (10.1016/j.jct.2013.05.025 25/25 kept; 10.1016/j.tca.2017.12.006 13/13 kept). Relative error in P over the kept points: mean 2.2705 \%, max 5.6578 \%; AmbroseWalton on the same points: mean 3.3549 \%, max 8.9404 \%. Outliers discarded: none. Its own normal boiling point (101.325 kPa): 467.48 K. PROMOTED on 2026-11-10 into data/standards/components/methylLevulinate.dat as reviewStatus autoPromoted, humanReviewed false, because: fit grade warning under the DEFAULT graded acceptance of DEV.md 4c C285 (no system/curationDict), 2 studies, 38 points kept, 0 discarded as outliers; max relative error 5.658 \% (mean 2.27 \%), grade warning (limits 5 / 10 / 20 \%); no Tb declared; AmbroseWalton on the same points mean 3.355 \% / max 8.94 \% @@RESULT-END@@

\end{tutnarrative}
\begin{tutgolden}
\goldrow{diag}{fit\_antoine}{A}{4.61830270357}
\goldrow{diag}{fit\_antoine}{B}{1844.87591677}
\goldrow{diag}{fit\_antoine}{C}{-67.5116488784}
\goldrow{diag}{fit\_antoine}{R2}{0.999817243581}
\goldrow{diag}{fit\_antoine}{dHvap\_kJ\_per\_mol}{48.7686315486}
\goldrow{diag}{fit\_antoine}{minimax\_evaluations}{903}
\goldrow{diag}{fit\_antoine}{minimax\_max\_scaled\_final}{0.0565783537168}
\goldrow{diag}{fit\_antoine}{minimax\_max\_scaled\_start}{0.0745783923445}
\goldrow{diag}{fit\_antoine}{n\_heldout}{0}
\goldrow{diag}{fit\_antoine}{n\_outliers\_dropped}{0}
\goldrow{diag}{fit\_antoine}{n\_points}{38}
\goldrow{diag}{fit\_antoine}{objective\_method}{2}
\goldrow{diag}{fit\_antoine}{relErr\_max\_fit\_pct}{5.65783537168}
\goldrow{diag}{fit\_antoine}{relErr\_max\_fit\_pct\_leastSquares}{7.45783923445}
\goldmore{3}
\end{tutgolden}

\clearpage
\tutentry{curate17\_nrtl\_diisopropylEther\_isopropanol}{tutorials/props/curation/curate17\_nrtl\_diisopropylEther\_isopropanol}
\begin{tutscheme}
\begin{tikzpicture}[>={Stealth[length=2mm]}, every node/.style={font=\scriptsize}]
  \node[draw, rounded corners=2pt, align=center, minimum width=3.1cm, minimum height=0.8cm, fill=black!4] at (0.0,-0.0) {\textbf{fit\_nrtl}\\[-1pt]{\tiny fitParameters}};
  \node[draw, rounded corners=2pt, align=center, minimum width=3.1cm, minimum height=0.8cm, fill=black!4] at (3.6,-0.0) {\textbf{fit\_nrtl\_start2}\\[-1pt]{\tiny fitParameters}};
  \node[draw, rounded corners=2pt, align=center, minimum width=3.1cm, minimum height=0.8cm, fill=black!4] at (7.2,-0.0) {\textbf{txy\_ideal}\\[-1pt]{\tiny propertyScan1D}};
  \node[draw, rounded corners=2pt, align=center, minimum width=3.1cm, minimum height=0.8cm, fill=black!4] at (0.0,-1.3) {\textbf{txy\_nrtl}\\[-1pt]{\tiny propertyScan1D}};
\end{tikzpicture}
\end{tutscheme}
\tutwhat{NRTL diisopropylEther-isopropanol regressed by the curation pipeline to measured bubble temperatures (ThermoML locates, the article sources), held out against an independent study}
\begin{tutnarrative}
AN NRTL PAIR FROM PUBLISHED MEASUREMENTS, BY THE PIPELINE.

WRITTEN BY bin/choupo-curate pair diisopropylEther isopropanol (DEV.md 4c C282; record docs/design/binary-pairs-from-open-measurements.md section 8). Every rule below was decided before this tool existed (C16, C225, C226); the tool applies them and decides nothing.

@@PILOT-BEGIN@@ WHY THIS PAIR -- THE PILOT RANKING (C282; added after the run, but the ranking reads ONLY the archive index and was computed before any fit). Candidates: both components in the catalogue, NO record for the pair in any model directory of data/standards/parameters/, no acetic acid (vapour dimerisation, C16), no component whose vapour pressure is AmbroseWalton (C281), and (checked for the five chosen) no binary block in the archive with a second liquid. Ranked by the number of mutually author-disjoint isobaric bubble-T studies at or below 110 kPa with a declared uncertainty, then by their total points: 1 methanol + dimethylCarbonate 5 independent (9 studies, 274 pts) - acetone + methanol 4 -- passed over: curate05 holds its C16 regression and the two studies' disagreement, a curator's tie-break 2 ethanol + dimethylCarbonate 4 (4 studies, 156 pts) 3 diisopropylEther + isopropanol 3 (4 studies, 152 pts) - triethyleneGlycol + water 3 -- REFUSED by the tool: 2 of 3 studies disagree beyond their U 4 isopropanol + dimethylCarbonate 3 (3 studies, 128 pts) - 2Butanol + dimethylCarbonate 3 -- REFUSED by the tool: no author-disjoint held-out study reports x of the same component 5 ethanol + 224Trimethylpentane 3 (5 studies, 97 pts) THIS CASE: rank 3. @@PILOT-END@@

THE DECLARATION -- WRITTEN AND COMMITTED BEFORE ANY FIT RAN. (C16 slice 2: so the declaration can be read against the result.)

IDENTITY. diisopropylEther (CAS 108-20-3, InChIKey ZAFNJMIOTHYJRJ-UHFFFAOYSA-N) + isopropanol (CAS 67-63-0, InChIKey KFZMGEQAYNKOFK-UHFFFAOYSA-N), joined to the archive through identity.json -- never by name.

NO MEASURED SERIES IN THIS TREE. constant/archiveDatasets names each dataset (DOI, NIST URL, article sha256, extraction selector, extract sha256); `bin/choupo-curate fetch` writes the points into constant/archiveExtracts/, which is gitignored.

STUDIES CONSIDERED: every bubble-temperature block of the pair at or below 110 kPa with a declared uncertainty (4 block(s)). USABLE after the screens: 4. SET ASIDE, BY NAME: (none)

TRIAGE (C226 item 3, leave one study out, fitted by the engine; a study whose held-out AAD exceeds its own largest stated U plus the others' largest stated U disagrees with them beyond their stated uncertainties): 10.1016/j.fluid.2014.09.014: AAD 1.829 K against the limit 2.770 K $\to$ agrees; 10.1016/j.fluid.2007.09.014: AAD 0.579 K against the limit 1.880 K $\to$ agrees; 10.1021/je300994y: AAD 0.298 K against the limit 2.240 K $\to$ agrees; 10.1016/j.jct.2018.11.013: AAD 0.395 K against the limit 2.770 K $\to$ agrees

PARTITION (the declared rule of C16 slice 2): the FIT set is the usable study with the most interior points; the HELD-OUT set is the next largest that shares NO author with it and was measured at one pressure. FIT doi 10.1016/j.fluid.2014.09.014 (67 interior points, 4 isobars 53.33-94.79 kPa) authors: Soujanya, J.; satyavathi, B.; Sankarshana, T. HELD OUT doi 10.1021/je300994y (19 interior points, 101.3 kPa) authors: Zhang, Z.[Zhigang]; Yang, L.[Lei]; Xing, Y.[Yuan]; Li, W.[Wenxiu]

BAND: maxAAD 0.24 \% of the held-out mean T. declared before the fit by bin/choupo-curate, the ONE rule of C16: the held-out set's largest stated expanded uncertainty (U95 = 0.61 K) plus the catalogue vapour-pressure model's own error at that set's pure endpoints (isopropanol -0.206 K, diisopropylEther +0.116 K; measured with no pair involved), as a percentage of the held-out mean T (342.77 K) = 0.2380 \%, rounded UP to two significant figures

THE FIT. NRTL with a\_ij = a\_ji = 0 and alpha = 0.30 fixed: TWO parameters (b\_ij, b\_ji), the C16 frame (bubble temperatures over these ranges do not identify more). Each point weighted by 1/U, U its own stated expanded uncertainty. Two starts: (0, 0) and (800, -400).

@@RESULT-BEGIN@@ WHAT THE RUN FOUND (pinned in `expected`; measured, not argued). fit\_nrtl: b\_ij = 169.0 K, b\_ji = 410.8 K, in-sample rms 0.815 K over 67 points; fit\_nrtl\_start2 lands on the same pair (168.9, 410.8). HELD OUT: AAD 1.894 K (0.5507 \%) against the band 0.24 \% $\to$ notValidated. NOT PROMOTABLE: held-out verdict notValidated (AAD 0.550724 \% against the band 0.24 \% declared before the fit) @@RESULT-END@@

WHAT IS NOT CLAIMED. The pair is correct at these pressures and in this frame (ideal-gas vapour, the catalogue's own vapour-pressure records); the measurements were transcribed by the archive and have NOT been read back against the articles (the CHECKLIST that `fetch` writes beside the extracts is the aid for that reading).

\end{tutnarrative}
\begin{tutgolden}
\goldrow{diag}{fit\_nrtl}{aad\_heldout\_K}{1.89391027641}
\goldrow{diag}{fit\_nrtl}{aad\_heldout\_pct}{0.550724059447}
\goldrow{diag}{fit\_nrtl}{bubble\_aad\_fit}{0.645273730052}
\goldrow{diag}{fit\_nrtl}{bubble\_aad\_heldout}{1.89391027641}
\goldrow{diag}{fit\_nrtl}{bubble\_max\_fit}{2.17775358224}
\goldrow{diag}{fit\_nrtl}{bubble\_max\_heldout}{3.19444093153}
\goldrow{diag}{fit\_nrtl}{chi2}{170.37382008}
\goldrow{diag}{fit\_nrtl}{chi2\_in\_sample}{170.37382008}
\goldrow{diag}{fit\_nrtl}{chi2\_reduced}{2.62113569354}
\goldrow{diag}{fit\_nrtl}{cond\_JtJ}{38.4725080237}
\goldrow{diag}{fit\_nrtl}{converged}{1}
\goldrow{diag}{fit\_nrtl}{corr\_computed}{1}
\goldrow{diag}{fit\_nrtl}{dof}{65}
\goldrow{diag}{fit\_nrtl}{identifiable}{1}
\goldmore{58}
\end{tutgolden}

\clearpage
\tutentry{curate18\_nrtl\_isopropanol\_dimethylCarbonate}{tutorials/props/curation/curate18\_nrtl\_isopropanol\_dimethylCarbonate}
\begin{tutscheme}
\begin{tikzpicture}[>={Stealth[length=2mm]}, every node/.style={font=\scriptsize}]
  \node[draw, rounded corners=2pt, align=center, minimum width=3.1cm, minimum height=0.8cm, fill=black!4] at (0.0,-0.0) {\textbf{fit\_nrtl}\\[-1pt]{\tiny fitParameters}};
  \node[draw, rounded corners=2pt, align=center, minimum width=3.1cm, minimum height=0.8cm, fill=black!4] at (3.6,-0.0) {\textbf{fit\_nrtl\_start2}\\[-1pt]{\tiny fitParameters}};
  \node[draw, rounded corners=2pt, align=center, minimum width=3.1cm, minimum height=0.8cm, fill=black!4] at (7.2,-0.0) {\textbf{txy\_ideal}\\[-1pt]{\tiny propertyScan1D}};
  \node[draw, rounded corners=2pt, align=center, minimum width=3.1cm, minimum height=0.8cm, fill=black!4] at (0.0,-1.3) {\textbf{txy\_nrtl}\\[-1pt]{\tiny propertyScan1D}};
  \node[draw, rounded corners=2pt, align=center, minimum width=3.1cm, minimum height=0.8cm, fill=black!4] at (3.6,-1.3) {\textbf{txy\_ideal\_fitP}\\[-1pt]{\tiny propertyScan1D}};
  \node[draw, rounded corners=2pt, align=center, minimum width=3.1cm, minimum height=0.8cm, fill=black!4] at (7.2,-1.3) {\textbf{txy\_nrtl\_fitP}\\[-1pt]{\tiny propertyScan1D}};
\end{tikzpicture}
\end{tutscheme}
\tutwhat{NRTL isopropanol-dimethylCarbonate regressed by the curation pipeline to measured bubble temperatures (ThermoML locates, the article sources), held out against an independent study}
\begin{tutnarrative}
AN NRTL PAIR FROM PUBLISHED MEASUREMENTS, BY THE PIPELINE.

WRITTEN BY bin/choupo-curate pair isopropanol dimethylCarbonate (DEV.md 4c C282; record docs/design/binary-pairs-from-open-measurements.md section 8). Every rule below was decided before this tool existed (C16, C225, C226); the tool applies them and decides nothing.

@@PILOT-BEGIN@@ WHY THIS PAIR -- THE PILOT RANKING (C282; added after the run, but the ranking reads ONLY the archive index and was computed before any fit). Candidates: both components in the catalogue, NO record for the pair in any model directory of data/standards/parameters/, no acetic acid (vapour dimerisation, C16), no component whose vapour pressure is AmbroseWalton (C281), and (checked for the five chosen) no binary block in the archive with a second liquid. Ranked by the number of mutually author-disjoint isobaric bubble-T studies at or below 110 kPa with a declared uncertainty, then by their total points: 1 methanol + dimethylCarbonate 5 independent (9 studies, 274 pts) - acetone + methanol 4 -- passed over: curate05 holds its C16 regression and the two studies' disagreement, a curator's tie-break 2 ethanol + dimethylCarbonate 4 (4 studies, 156 pts) 3 diisopropylEther + isopropanol 3 (4 studies, 152 pts) - triethyleneGlycol + water 3 -- REFUSED by the tool: 2 of 3 studies disagree beyond their U 4 isopropanol + dimethylCarbonate 3 (3 studies, 128 pts) - 2Butanol + dimethylCarbonate 3 -- REFUSED by the tool: no author-disjoint held-out study reports x of the same component 5 ethanol + 224Trimethylpentane 3 (5 studies, 97 pts) THIS CASE: rank 4. @@PILOT-END@@

THE DECLARATION -- WRITTEN AND COMMITTED BEFORE ANY FIT RAN. (C16 slice 2: so the declaration can be read against the result.)

IDENTITY. isopropanol (CAS 67-63-0, InChIKey KFZMGEQAYNKOFK-UHFFFAOYSA-N) + dimethylCarbonate (CAS 616-38-6, InChIKey IEJIGPNLZYLLBP-UHFFFAOYSA-N), joined to the archive through identity.json -- never by name.

NO MEASURED SERIES IN THIS TREE. constant/archiveDatasets names each dataset (DOI, NIST URL, article sha256, extraction selector, extract sha256); `bin/choupo-curate fetch` writes the points into constant/archiveExtracts/, which is gitignored.

STUDIES CONSIDERED: every bubble-temperature block of the pair at or below 110 kPa with a declared uncertainty (3 block(s)). USABLE after the screens: 2. SET ASIDE, BY NAME: - doi 10.1016/j.jct.2011.08.006 block 6: TRIAGE: disagrees with the other studies beyond their stated uncertainties (leave-one-study-out, below)

TRIAGE (C226 item 3, leave one study out, fitted by the engine; a study whose held-out AAD exceeds its own largest stated U plus the others' largest stated U disagrees with them beyond their stated uncertainties): 10.1016/j.jct.2011.08.006: AAD 0.636 K against the limit 0.520 K $\to$ DISAGREES -- set aside; 10.1016/j.jct.2004.09.009: AAD 0.629 K against the limit 0.670 K $\to$ agrees; 10.1021/je030131q: AAD 0.578 K against the limit 0.670 K $\to$ agrees

PARTITION (the declared rule of C16 slice 2): the FIT set is the usable study with the most interior points; the HELD-OUT set is the next largest that shares NO author with it and was measured at one pressure. FIT doi 10.1016/j.jct.2004.09.009 (32 interior points, 101.3 kPa) authors: Pereiro, A. B.; Rodriguez, A.; Canosa, J.; Tojo, J. HELD OUT doi 10.1021/je030131q (9 interior points, 95.8 kPa) authors: Vittal Prasad, T. E.; Prashanth, P. R.; Kumar, D. S.; Prasad, D.H.L.

BAND: maxAAD 0.24 \% of the held-out mean T. declared before the fit by bin/choupo-curate, the ONE rule of C16: the held-out set's largest stated expanded uncertainty (U95 = 0.32 K) plus the catalogue vapour-pressure model's own error at that set's pure endpoints (dimethylCarbonate +0.504 K, isopropanol -0.231 K; measured with no pair involved), as a percentage of the held-out mean T (351.83 K) = 0.2341 \%, rounded UP to two significant figures

THE FIT. NRTL with a\_ij = a\_ji = 0 and alpha = 0.30 fixed: TWO parameters (b\_ij, b\_ji), the C16 frame (bubble temperatures over these ranges do not identify more). Each point weighted by 1/U, U its own stated expanded uncertainty. Two starts: (0, 0) and (800, -400).

@@RESULT-BEGIN@@ WHAT THE RUN FOUND (pinned in `expected`; measured, not argued). fit\_nrtl: b\_ij = 284.4 K, b\_ji = 62.9 K, in-sample rms 0.132 K over 32 points; fit\_nrtl\_start2 lands on the same pair (284.4, 62.9). HELD OUT: AAD 0.184 K (0.0521 \%) against the band 0.24 \% $\to$ validated. PROMOTED on 2026-11-10 into data/standards/parameters/NRTL/dimethylCarbonate-isopropanol.dat as reviewStatus autoPromoted, humanReviewed false, because: verdict validated, one liquid across the validity span, triage run @@RESULT-END@@

WHAT IS NOT CLAIMED. The pair is correct at these pressures and in this frame (ideal-gas vapour, the catalogue's own vapour-pressure records); the measurements were transcribed by the archive and have NOT been read back against the articles (the CHECKLIST that `fetch` writes beside the extracts is the aid for that reading).

\end{tutnarrative}
\begin{tutgolden}
\goldrow{diag}{fit\_nrtl}{aad\_heldout\_K}{0.183696599893}
\goldrow{diag}{fit\_nrtl}{aad\_heldout\_pct}{0.0520996119818}
\goldrow{diag}{fit\_nrtl}{bubble\_aad\_fit}{0.0867280388939}
\goldrow{diag}{fit\_nrtl}{bubble\_aad\_heldout}{0.183696599893}
\goldrow{diag}{fit\_nrtl}{bubble\_max\_fit}{0.358804574762}
\goldrow{diag}{fit\_nrtl}{bubble\_max\_heldout}{0.542966592647}
\goldrow{diag}{fit\_nrtl}{chi2}{4.57923866199}
\goldrow{diag}{fit\_nrtl}{chi2\_in\_sample}{4.57923866199}
\goldrow{diag}{fit\_nrtl}{chi2\_reduced}{0.152641288733}
\goldrow{diag}{fit\_nrtl}{cond\_JtJ}{76.7456403943}
\goldrow{diag}{fit\_nrtl}{converged}{1}
\goldrow{diag}{fit\_nrtl}{corr\_computed}{1}
\goldrow{diag}{fit\_nrtl}{dof}{30}
\goldrow{diag}{fit\_nrtl}{identifiable}{1}
\goldmore{68}
\end{tutgolden}

\clearpage
\tutentry{curate19\_nrtl\_ethanol\_224Trimethylpentane}{tutorials/props/curation/curate19\_nrtl\_ethanol\_224Trimethylpentane}
\begin{tutscheme}
\begin{tikzpicture}[>={Stealth[length=2mm]}, every node/.style={font=\scriptsize}]
  \node[draw, rounded corners=2pt, align=center, minimum width=3.1cm, minimum height=0.8cm, fill=black!4] at (0.0,-0.0) {\textbf{fit\_nrtl}\\[-1pt]{\tiny fitParameters}};
  \node[draw, rounded corners=2pt, align=center, minimum width=3.1cm, minimum height=0.8cm, fill=black!4] at (3.6,-0.0) {\textbf{fit\_nrtl\_start2}\\[-1pt]{\tiny fitParameters}};
  \node[draw, rounded corners=2pt, align=center, minimum width=3.1cm, minimum height=0.8cm, fill=black!4] at (7.2,-0.0) {\textbf{txy\_ideal}\\[-1pt]{\tiny propertyScan1D}};
  \node[draw, rounded corners=2pt, align=center, minimum width=3.1cm, minimum height=0.8cm, fill=black!4] at (0.0,-1.3) {\textbf{txy\_unifac}\\[-1pt]{\tiny propertyScan1D}};
  \node[draw, rounded corners=2pt, align=center, minimum width=3.1cm, minimum height=0.8cm, fill=black!4] at (3.6,-1.3) {\textbf{txy\_nrtl}\\[-1pt]{\tiny propertyScan1D}};
  \node[draw, rounded corners=2pt, align=center, minimum width=3.1cm, minimum height=0.8cm, fill=black!4] at (7.2,-1.3) {\textbf{txy\_ideal\_fitP}\\[-1pt]{\tiny propertyScan1D}};
  \node[draw, rounded corners=2pt, align=center, minimum width=3.1cm, minimum height=0.8cm, fill=black!4] at (0.0,-2.6) {\textbf{txy\_unifac\_fitP}\\[-1pt]{\tiny propertyScan1D}};
  \node[draw, rounded corners=2pt, align=center, minimum width=3.1cm, minimum height=0.8cm, fill=black!4] at (3.6,-2.6) {\textbf{txy\_nrtl\_fitP}\\[-1pt]{\tiny propertyScan1D}};
\end{tikzpicture}
\end{tutscheme}
\tutwhat{NRTL ethanol-224Trimethylpentane regressed by the curation pipeline to measured bubble temperatures (ThermoML locates, the article sources), held out against an independent study}
\begin{tutnarrative}
AN NRTL PAIR FROM PUBLISHED MEASUREMENTS, BY THE PIPELINE.

WRITTEN BY bin/choupo-curate pair ethanol 224Trimethylpentane (DEV.md 4c C282; record docs/design/binary-pairs-from-open-measurements.md section 8). Every rule below was decided before this tool existed (C16, C225, C226); the tool applies them and decides nothing.

@@PILOT-BEGIN@@ WHY THIS PAIR -- THE PILOT RANKING (C282; added after the run, but the ranking reads ONLY the archive index and was computed before any fit). Candidates: both components in the catalogue, NO record for the pair in any model directory of data/standards/parameters/, no acetic acid (vapour dimerisation, C16), no component whose vapour pressure is AmbroseWalton (C281), and (checked for the five chosen) no binary block in the archive with a second liquid. Ranked by the number of mutually author-disjoint isobaric bubble-T studies at or below 110 kPa with a declared uncertainty, then by their total points: 1 methanol + dimethylCarbonate 5 independent (9 studies, 274 pts) - acetone + methanol 4 -- passed over: curate05 holds its C16 regression and the two studies' disagreement, a curator's tie-break 2 ethanol + dimethylCarbonate 4 (4 studies, 156 pts) 3 diisopropylEther + isopropanol 3 (4 studies, 152 pts) - triethyleneGlycol + water 3 -- REFUSED by the tool: 2 of 3 studies disagree beyond their U 4 isopropanol + dimethylCarbonate 3 (3 studies, 128 pts) - 2Butanol + dimethylCarbonate 3 -- REFUSED by the tool: no author-disjoint held-out study reports x of the same component 5 ethanol + 224Trimethylpentane 3 (5 studies, 97 pts) THIS CASE: rank 5. @@PILOT-END@@

THE DECLARATION -- WRITTEN AND COMMITTED BEFORE ANY FIT RAN. (C16 slice 2: so the declaration can be read against the result.)

IDENTITY. ethanol (CAS 64-17-5, InChIKey LFQSCWFLJHTTHZ-UHFFFAOYSA-N) + 224Trimethylpentane (CAS 540-84-1, InChIKey NHTMVDHEPJAVLT-UHFFFAOYSA-N), joined to the archive through identity.json -- never by name.

NO MEASURED SERIES IN THIS TREE. constant/archiveDatasets names each dataset (DOI, NIST URL, article sha256, extraction selector, extract sha256); `bin/choupo-curate fetch` writes the points into constant/archiveExtracts/, which is gitignored.

STUDIES CONSIDERED: every bubble-temperature block of the pair at or below 110 kPa with a declared uncertainty (5 block(s)). USABLE after the screens: 5. SET ASIDE, BY NAME: (none)

TRIAGE (C226 item 3, leave one study out, fitted by the engine; a study whose held-out AAD exceeds its own largest stated U plus the others' largest stated U disagrees with them beyond their stated uncertainties): 10.1021/je300847v: AAD 0.420 K against the limit 4.710 K $\to$ agrees; 10.1021/acs.jced.8b00680: AAD 0.372 K against the limit 4.710 K $\to$ agrees; 10.1016/j.fluid.2005.01.007: AAD 0.436 K against the limit 4.340 K $\to$ agrees; 10.1016/j.fluid.2007.06.005: AAD 0.411 K against the limit 4.310 K $\to$ agrees; 10.1016/j.jct.2005.04.006: AAD 2.145 K against the limit 3.420 K $\to$ agrees

PARTITION (the declared rule of C16 slice 2): the FIT set is the usable study with the most interior points; the HELD-OUT set is the next largest that shares NO author with it and was measured at one pressure. FIT doi 10.1021/je300847v (22 interior points, 101.3 kPa) authors: Pienaar, C.[Cornelia]; Schwarz, C.[Cara]; Knoetze, J. H.[Johannes H.]; Burger, A.[Andries] HELD OUT doi 10.1016/j.fluid.2005.01.007 (19 interior points, 101.32 kPa) authors: Ku, H. -C.[Hsu-Chen]; Tu, C. -H.[Chein-Hsiun]

BAND: maxAAD 0.57 \% of the held-out mean T. declared before the fit by bin/choupo-curate, the ONE rule of C16: the held-out set's largest stated expanded uncertainty (U95 = 1.93 K) plus the catalogue vapour-pressure model's own error at that set's pure endpoints (224Trimethylpentane -0.012 K, ethanol +0.012 K; measured with no pair involved), as a percentage of the held-out mean T (346.81 K) = 0.5600 \%, rounded UP to two significant figures

THE FIT. NRTL with a\_ij = a\_ji = 0 and alpha = 0.30 fixed: TWO parameters (b\_ij, b\_ji), the C16 frame (bubble temperatures over these ranges do not identify more). Each point weighted by 1/U, U its own stated expanded uncertainty. Two starts: (0, 0) and (800, -400).

@@RESULT-BEGIN@@ WHAT THE RUN FOUND (pinned in `expected`; measured, not argued). fit\_nrtl: b\_ij = 429.4 K, b\_ji = 493.2 K, in-sample rms 0.484 K over 22 points; fit\_nrtl\_start2 lands on the same pair (429.4, 493.2). HELD OUT: AAD 0.394 K (0.1132 \%) against the band 0.57 \% $\to$ validated. NOT PROMOTABLE: the pair splits ONE liquid inside its own validity span (check\_regressed\_pairs arm f) -- predicts a LIQUID-LIQUID SPLIT at 344.57 K (x\_ethanol 0.313 .. 0.648), inside its own validity span 344.57-347.65 K -- its evidence is bubble temperatures of ONE liquid, which a bubble scan cannot contradict; a pair that splits that liquid stays in its case, not the catalogue (or it DECLARES `liquidLiquid true;` with tie-line evidence and a gap band) @@RESULT-END@@

WHAT IS NOT CLAIMED. The pair is correct at these pressures and in this frame (ideal-gas vapour, the catalogue's own vapour-pressure records); the measurements were transcribed by the archive and have NOT been read back against the articles (the CHECKLIST that `fetch` writes beside the extracts is the aid for that reading).

\end{tutnarrative}
\begin{tutgolden}
\goldrow{diag}{fit\_nrtl}{aad\_heldout\_K}{0.394419156171}
\goldrow{diag}{fit\_nrtl}{aad\_heldout\_pct}{0.113232264022}
\goldrow{diag}{fit\_nrtl}{bubble\_aad\_fit}{0.408550722801}
\goldrow{diag}{fit\_nrtl}{bubble\_aad\_heldout}{0.394419156171}
\goldrow{diag}{fit\_nrtl}{bubble\_max\_fit}{0.911330416911}
\goldrow{diag}{fit\_nrtl}{bubble\_max\_heldout}{1.18628480853}
\goldrow{diag}{fit\_nrtl}{chi2}{17.6327244815}
\goldrow{diag}{fit\_nrtl}{chi2\_in\_sample}{17.6327244815}
\goldrow{diag}{fit\_nrtl}{chi2\_reduced}{0.881636224074}
\goldrow{diag}{fit\_nrtl}{cond\_JtJ}{15.0198048572}
\goldrow{diag}{fit\_nrtl}{converged}{1}
\goldrow{diag}{fit\_nrtl}{corr\_computed}{1}
\goldrow{diag}{fit\_nrtl}{dof}{20}
\goldrow{diag}{fit\_nrtl}{identifiable}{1}
\goldmore{82}
\end{tutgolden}

\clearpage
\tutentry{curate20\_vp\_diButylCarbonate}{tutorials/props/curation/curate20\_vp\_diButylCarbonate}
\begin{tutscheme}
\begin{tikzpicture}[>={Stealth[length=2mm]}, every node/.style={font=\scriptsize}]
  \node[draw, rounded corners=2pt, align=center, minimum width=3.1cm, minimum height=0.8cm, fill=black!4] at (0.0,-0.0) {\textbf{fit\_antoine}\\[-1pt]{\tiny vaporPressureFit}};
\end{tikzpicture}
\end{tutscheme}
\tutwhat{Antoine vapour pressure of diButylCarbonate fitted by the curation pipeline to measured vapour pressures (ThermoML locates, the article sources), held out against an independent study}
\begin{tutnarrative}
the vapour pressure of diButylCarbonate, REFITTED to measurements (DEV.md 4c C281 and C283). Written by `bin/choupo-curate vp diButylCarbonate`; everything above the RESULT paragraph is the DECLARATION, committed before any fit ran, so it can be read against the result.

WHY. The catalogue record carried `vaporPressure \{ model AmbroseWalton; \}`, a corresponding-states ESTIMATE from Tc, Pc and omega that never reads the record's own Tb (C279 measured a median 19.6 \% AAD over the 25 such records with archive data). Here a three-parameter Antoine is fitted to EVERY usable laboratory's measurements of it.

STUDIES (one per article, extracted by `choupo-thermoml extract-vp`; ranked by points, ties by DOI): 10.1021/acs.jced.6b00517 23 pts 383.16-478.18 K Yang; Xu; Yang 10.1016/j.jct.2008.02.012 21 pts 287.1-329.1 K Kozlova; Emel'yanenko; Georgieva; Verevkin et al. SET ASIDE BY NAME: (none)

FIT AND GRADED ACCEPTANCE (DEV.md 4c C285, Vitor 2026-11-10; the evidence, objective and outliers of C283; declared before the fit). EVIDENCE: every study above is fitted; none is held out (C283 amended: laboratories disagree with each other by more than they claim, so the test is agreement across all of them, not one laboratory's claimed precision). OBJECTIVE: `objective \{ method twoStage; \}` -- least squares in log10 P for a first estimate, then the LARGEST relative error in P, (P\_model - P)/P, over the kept points minimised from it (Vitor's practice, C280); both answers kept. OUTLIERS: after the least-squares stage a point with a relative P error above 15 \% (three times the acceptance) is DISCARDED BY NAME (its study DOI, T and P), listed in the result and the record, and the fit repeated without it, once (`outliers \{ maxRelativeError 15 percent; \}`). ACCEPTANCE, GRADED, the DEFAULT of C285: the maximum relative P error over all kept points is graded clean up to 5 \%, warning up to 10 \%, severe up to 20 \% and refused above; where the normal-boiling-point condition below applies, the absolute gap is graded clean up to 1 K, warning up to 2 K, severe up to 4 K and refused above. The fit's grade is the WORSE of the two. A clean, warning or severe fit is ACCEPTED and promoted (autoPromoted, humanReviewed false) with `fitGrade <grade>;` and the deciding numbers in its provenance entry, which the component loader voices on every run that uses it (a warning or a severe grade); a refused one is not promoted. This case declares no `system/curationDict`, so the default applies; a case that needs another acceptance declares it there, with its reason. The AmbroseWalton curve's errors on the same points are reported beside the fit's, as information.

NORMAL BOILING POINT: the record's Tb 480.6 K is judged against the range of the KEPT points (all studies span 287.1-478.18 K before any outlier is discarded): if Tb lies OUTSIDE that range, the fitted Antoine's own normal boiling point (P = 101.325 kPa) is GRADED by its absolute gap to it (the temperature limits of the acceptance below); inside it, the gap is reported and is not a condition.

EARLIER FITS, SUPERSEDED BY THIS DECLARATION @@EARLIER-BEGIN@@ C281 step 2, objective leastSquares (log10 P), the same partition and band: WHAT THE RUN FOUND (pinned in `expected`; measured, not argued). Antoine log10(Psat/bar) = 4.783236 - 2119.1665/(T + -34.8668), in-sample R2 0.998031 over 23 points, 383.16-478.18 K. HELD OUT: 21 points, AAD 105.1167 \% against the band 5.3 \% $\to$ notValidated. Its own normal boiling point (101.325 kPa): 478.44 K. NOT PROMOTED: held-out verdict notValidated (AAD 105.1 \% against the band 5.3 \% declared before the fit); the record keeps AmbroseWalton THEN C281 step 3, objective twoStage, one study fitted and one held out, judged by the band (superseded by C283): WHAT THE RUN FOUND (pinned in `expected`; measured, not argued). Objective twoStage. Antoine log10(Psat/bar) = 3.531990 - 1207.1399/(T + -135.3400), in-sample R2 0.995897 over 23 points, 383.16-478.18 K. HELD OUT: 21 points, AAD 34.1314 \% (max 55.8176 \%) against the band 5.3 \% $\to$ notValidated; the stage-1 least-squares answer on the same points: AAD 105.1167 \%, max 181.2842 \%. Its own normal boiling point (101.325 kPa): 477.67 K. NOT PROMOTED: held-out verdict notValidated (AAD 34.13 \% against the band 5.3 \% declared before the fit); the record keeps AmbroseWalton. THEN C283, every usable study fitted and none held out, accepted only at max relative error <= 5 \% and Tnb within 1 K (superseded by the graded acceptance of C285): WHAT THE RUN FOUND (pinned in `expected`; measured, not argued). Objective twoStage. Antoine log10(Psat/bar) = 4.018870 - 1504.3419/(T + -101.9840), in-sample R2 0.999789 over 44 kept points, 287.1-478.18 K (10.1021/acs.jced.6b00517 23/23 kept; 10.1016/j.jct.2008.02.012 21/21 kept). Relative error in P over the kept points: mean 4.2226 \%, max 8.8166 \%; AmbroseWalton on the same points: mean 10.8661 \%, max 27.2481 \%. Outliers discarded: none. Its own normal boiling point (101.325 kPa): 476.84 K. NOT PROMOTED by the DEFAULT acceptance of DEV.md 4c C283 (no system/curationDict): 2 studies, 44 points kept, 0 discarded as outliers; max relative error 8.817 \% (mean 4.223 \%) against 5 \%; AmbroseWalton on the same points mean 10.87 \% / max 27.25 \%; the record keeps AmbroseWalton @@EARLIER-END@@

NO MEASURED POINT IS IN THIS TREE: constant/archiveDatasets declares every series by DOI, NIST URL and sha256; `bin/choupo-curate fetch` regenerates them into constant/archiveExtracts/ (gitignored).

RESULT @@RESULT-BEGIN@@ WHAT THE RUN FOUND (pinned in `expected`; measured, not argued). Objective twoStage. Antoine log10(Psat/bar) = 4.018870 - 1504.3419/(T + -101.9840), in-sample R2 0.999789 over 44 kept points, 287.1-478.18 K (10.1021/acs.jced.6b00517 23/23 kept; 10.1016/j.jct.2008.02.012 21/21 kept). Relative error in P over the kept points: mean 4.2226 \%, max 8.8166 \%; AmbroseWalton on the same points: mean 10.8661 \%, max 27.2481 \%. Outliers discarded: none. Its own normal boiling point (101.325 kPa): 476.84 K. PROMOTED on 2026-11-10 into data/standards/components/diButylCarbonate.dat as reviewStatus autoPromoted, humanReviewed false, because: fit grade severe under the DEFAULT graded acceptance of DEV.md 4c C285 (no system/curationDict), 2 studies, 44 points kept, 0 discarded as outliers; max relative error 8.817 \% (mean 4.223 \%), grade warning (limits 5 / 10 / 20 \%); Tb outside the kept points' range, normal boiling point 476.84 K, -3.76 K from the record's Tb 480.6 K, grade severe (limits 1 / 2 / 4 K); AmbroseWalton on the same points mean 10.87 \% / max 27.25 \% @@RESULT-END@@

\end{tutnarrative}
\begin{tutgolden}
\goldrow{diag}{fit\_antoine}{A}{4.01886992528}
\goldrow{diag}{fit\_antoine}{B}{1504.34189732}
\goldrow{diag}{fit\_antoine}{C}{-101.983954832}
\goldrow{diag}{fit\_antoine}{R2}{0.99978907588}
\goldrow{diag}{fit\_antoine}{dHvap\_kJ\_per\_mol}{46.5319003141}
\goldrow{diag}{fit\_antoine}{n\_heldout}{0}
\goldrow{diag}{fit\_antoine}{n\_points}{44}
\goldrow{verdict}{fit\_antoine}{verdict}{validationRefused}
\goldrow{diag}{fit\_antoine}{minimax\_evaluations}{1172}
\goldrow{diag}{fit\_antoine}{minimax\_max\_scaled\_final}{0.088165879121}
\goldrow{diag}{fit\_antoine}{minimax\_max\_scaled\_start}{0.124705391729}
\goldrow{diag}{fit\_antoine}{objective\_method}{2}
\goldrow{diag}{fit\_antoine}{relErr\_max\_fit\_pct}{8.8165879121}
\goldrow{diag}{fit\_antoine}{relErr\_max\_fit\_pct\_leastSquares}{12.4705391729}
\goldmore{3}
\end{tutgolden}

\clearpage
\tutentry{curate21\_vp\_diEthylCarbonate}{tutorials/props/curation/curate21\_vp\_diEthylCarbonate}
\begin{tutscheme}
\begin{tikzpicture}[>={Stealth[length=2mm]}, every node/.style={font=\scriptsize}]
  \node[draw, rounded corners=2pt, align=center, minimum width=3.1cm, minimum height=0.8cm, fill=black!4] at (0.0,-0.0) {\textbf{fit\_antoine}\\[-1pt]{\tiny vaporPressureFit}};
\end{tikzpicture}
\end{tutscheme}
\tutwhat{Antoine vapour pressure of diEthylCarbonate fitted by the curation pipeline to measured vapour pressures (ThermoML locates, the article sources), held out against an independent study}
\begin{tutnarrative}
the vapour pressure of diEthylCarbonate, REFITTED to measurements (DEV.md 4c C281 and C283). Written by `bin/choupo-curate vp diEthylCarbonate`; everything above the RESULT paragraph is the DECLARATION, committed before any fit ran, so it can be read against the result.

WHY. The catalogue record carried `vaporPressure \{ model AmbroseWalton; \}`, a corresponding-states ESTIMATE from Tc, Pc and omega that never reads the record's own Tb (C279 measured a median 19.6 \% AAD over the 25 such records with archive data). Here a three-parameter Antoine is fitted to EVERY usable laboratory's measurements of it.

STUDIES (one per article, extracted by `choupo-thermoml extract-vp`; ranked by points, ties by DOI): 10.1016/j.jct.2008.02.012 44 pts 273.1-315.2 K Kozlova; Emel'yanenko; Georgieva; Verevkin et al. 10.1021/acs.jced.7b00295 42 pts 238.15-296.15 K Pokorny; Stejfa; Fulem; Cervinka et al. 10.1021/je200784y 33 pts 354.68-399.27 K Marrufo; Loras; Lladosa 10.1021/je800984q 21 pts 343.77-398.28 K Xing; Fang; Li; Guo et al. 10.1016/j.jct.2014.07.004 17 pts 277.96-353.15 K Dougassa; Jacquemin; Ouatani; Tessier et al. 10.1021/acs.jced.7b00236 10 pts 303.15-323.15 K Anugraha; Altway; Wibawa SET ASIDE BY NAME: 10.1016/j.fluid.2012.11.026: 2 pressure point(s), fewer than 5 10.1016/j.jct.2015.07.024: 3 pressure point(s), fewer than 5 10.1021/acs.jced.5b00064: 1 pressure point(s), fewer than 5 10.1021/acs.jced.6b00085: 1 pressure point(s), fewer than 5 10.1021/acs.jced.9b00068: 1 pressure point(s), fewer than 5 10.1021/acs.jced.9b00430: 1 pressure point(s), fewer than 5 10.1021/je049544x: 1 pressure point(s), fewer than 5

FIT AND GRADED ACCEPTANCE (DEV.md 4c C285, Vitor 2026-11-10; the evidence, objective and outliers of C283; declared before the fit). EVIDENCE: every study above is fitted; none is held out (C283 amended: laboratories disagree with each other by more than they claim, so the test is agreement across all of them, not one laboratory's claimed precision). OBJECTIVE: `objective \{ method twoStage; \}` -- least squares in log10 P for a first estimate, then the LARGEST relative error in P, (P\_model - P)/P, over the kept points minimised from it (Vitor's practice, C280); both answers kept. OUTLIERS: after the least-squares stage a point with a relative P error above 15 \% (three times the acceptance) is DISCARDED BY NAME (its study DOI, T and P), listed in the result and the record, and the fit repeated without it, once (`outliers \{ maxRelativeError 15 percent; \}`). ACCEPTANCE, GRADED, the DEFAULT of C285: the maximum relative P error over all kept points is graded clean up to 5 \%, warning up to 10 \%, severe up to 20 \% and refused above; where the normal-boiling-point condition below applies, the absolute gap is graded clean up to 1 K, warning up to 2 K, severe up to 4 K and refused above. The fit's grade is the WORSE of the two. A clean, warning or severe fit is ACCEPTED and promoted (autoPromoted, humanReviewed false) with `fitGrade <grade>;` and the deciding numbers in its provenance entry, which the component loader voices on every run that uses it (a warning or a severe grade); a refused one is not promoted. This case declares no `system/curationDict`, so the default applies; a case that needs another acceptance declares it there, with its reason. The AmbroseWalton curve's errors on the same points are reported beside the fit's, as information.

NORMAL BOILING POINT: the record's Tb 399.95 K is judged against the range of the KEPT points (all studies span 238.15-399.27 K before any outlier is discarded): if Tb lies OUTSIDE that range, the fitted Antoine's own normal boiling point (P = 101.325 kPa) is GRADED by its absolute gap to it (the temperature limits of the acceptance below); inside it, the gap is reported and is not a condition.

EARLIER FITS, SUPERSEDED BY THIS DECLARATION @@EARLIER-BEGIN@@ C281 step 2, objective leastSquares (log10 P), the same partition and band: WHAT THE RUN FOUND (pinned in `expected`; measured, not argued). Antoine log10(Psat/bar) = 7.478412 - 3270.7603/(T + 53.2804), in-sample R2 0.999007 over 44 points, 273.1-315.2 K. HELD OUT: 42 points, AAD 13.1956 \% against the band 2 \% $\to$ notValidated. Its own normal boiling point (101.325 kPa): 384.41 K. NOT PROMOTED: held-out verdict notValidated (AAD 13.2 \% against the band 2 \% declared before the fit); the record keeps AmbroseWalton THEN C281 step 3, objective twoStage, one study fitted and one held out, judged by the band (superseded by C283): WHAT THE RUN FOUND (pinned in `expected`; measured, not argued). Objective twoStage. Antoine log10(Psat/bar) = 11.408700 - 6651.8995/(T + 204.3733), in-sample R2 0.998387 over 44 points, 273.1-315.2 K. HELD OUT: 42 points, AAD 25.8563 \% (max 90.1072 \%) against the band 2 \% $\to$ notValidated; the stage-1 least-squares answer on the same points: AAD 13.1956 \%, max 43.6593 \%. Its own normal boiling point (101.325 kPa): 378.97 K. NOT PROMOTED: held-out verdict notValidated (AAD 25.86 \% against the band 2 \% declared before the fit); the record keeps AmbroseWalton. THEN C283, every usable study fitted and none held out, accepted only at max relative error <= 5 \% and Tnb within 1 K (superseded by the graded acceptance of C285): WHAT THE RUN FOUND (pinned in `expected`; measured, not argued). Objective twoStage. Antoine log10(Psat/bar) = 4.174165 - 1387.2771/(T + -67.0909), in-sample R2 0.999868 over 167 kept points, 238.15-399.27 K (10.1016/j.jct.2008.02.012 44/44 kept; 10.1021/acs.jced.7b00295 42/42 kept; 10.1021/je200784y 33/33 kept; 10.1021/je800984q 21/21 kept; 10.1016/j.jct.2014.07.004 17/17 kept; 10.1021/acs.jced.7b00236 10/10 kept). Relative error in P over the kept points: mean 2.2287 \%, max 7.5432 \%; AmbroseWalton on the same points: mean 76.2139 \%, max 177.2845 \%. Outliers discarded: none. Its own normal boiling point (101.325 kPa): 399.90 K. NOT PROMOTED by the DEFAULT acceptance of DEV.md 4c C283 (no system/curationDict): 6 studies, 167 points kept, 0 discarded as outliers; max relative error 7.543 \% (mean 2.229 \%) against 5 \%; AmbroseWalton on the same points mean 76.21 \% / max 177.3 \%; the record keeps AmbroseWalton @@EARLIER-END@@

NO MEASURED POINT IS IN THIS TREE: constant/archiveDatasets declares every series by DOI, NIST URL and sha256; `bin/choupo-curate fetch` regenerates them into constant/archiveExtracts/ (gitignored).

RESULT @@RESULT-BEGIN@@ WHAT THE RUN FOUND (pinned in `expected`; measured, not argued). Objective twoStage. Antoine log10(Psat/bar) = 4.174165 - 1387.2771/(T + -67.0909), in-sample R2 0.999868 over 167 kept points, 238.15-399.27 K (10.1016/j.jct.2008.02.012 44/44 kept; 10.1021/acs.jced.7b00295 42/42 kept; 10.1021/je200784y 33/33 kept; 10.1021/je800984q 21/21 kept; 10.1016/j.jct.2014.07.004 17/17 kept; 10.1021/acs.jced.7b00236 10/10 kept). Relative error in P over the kept points: mean 2.2287 \%, max 7.5432 \%; AmbroseWalton on the same points: mean 76.2139 \%, max 177.2845 \%. Outliers discarded: none. Its own normal boiling point (101.325 kPa): 399.90 K. PROMOTED on 2026-11-10 into data/standards/components/diEthylCarbonate.dat as reviewStatus autoPromoted, humanReviewed false, because: fit grade warning under the DEFAULT graded acceptance of DEV.md 4c C285 (no system/curationDict), 6 studies, 167 points kept, 0 discarded as outliers; max relative error 7.543 \% (mean 2.229 \%), grade warning (limits 5 / 10 / 20 \%); Tb outside the kept points' range, normal boiling point 399.90 K, -0.05 K from the record's Tb 399.95 K, grade clean (limits 1 / 2 / 4 K); AmbroseWalton on the same points mean 76.21 \% / max 177.3 \% @@RESULT-END@@

\end{tutnarrative}
\begin{tutgolden}
\goldrow{diag}{fit\_antoine}{A}{4.17416509512}
\goldrow{diag}{fit\_antoine}{B}{1387.27706685}
\goldrow{diag}{fit\_antoine}{C}{-67.090906648}
\goldrow{diag}{fit\_antoine}{R2}{0.999868166659}
\goldrow{diag}{fit\_antoine}{dHvap\_kJ\_per\_mol}{38.3708879568}
\goldrow{diag}{fit\_antoine}{n\_heldout}{0}
\goldrow{diag}{fit\_antoine}{n\_points}{167}
\goldrow{verdict}{fit\_antoine}{verdict}{validationRefused}
\goldrow{diag}{fit\_antoine}{minimax\_evaluations}{1052}
\goldrow{diag}{fit\_antoine}{minimax\_max\_scaled\_final}{0.0754322414701}
\goldrow{diag}{fit\_antoine}{minimax\_max\_scaled\_start}{0.0791737995068}
\goldrow{diag}{fit\_antoine}{objective\_method}{2}
\goldrow{diag}{fit\_antoine}{relErr\_max\_fit\_pct}{7.54322414701}
\goldrow{diag}{fit\_antoine}{relErr\_max\_fit\_pct\_leastSquares}{7.91737995068}
\goldmore{3}
\end{tutgolden}

\clearpage
\tutentry{curate22\_vp\_diNButylEther}{tutorials/props/curation/curate22\_vp\_diNButylEther}
\begin{tutscheme}
\begin{tikzpicture}[>={Stealth[length=2mm]}, every node/.style={font=\scriptsize}]
  \node[draw, rounded corners=2pt, align=center, minimum width=3.1cm, minimum height=0.8cm, fill=black!4] at (0.0,-0.0) {\textbf{fit\_antoine}\\[-1pt]{\tiny vaporPressureFit}};
\end{tikzpicture}
\end{tutscheme}
\tutwhat{Antoine vapour pressure of diNButylEther fitted by the curation pipeline to measured vapour pressures (ThermoML locates, the article sources), held out against an independent study}
\begin{tutnarrative}
the vapour pressure of diNButylEther, REFITTED to measurements (DEV.md 4c C281 and C283). Written by `bin/choupo-curate vp diNButylEther`; everything above the RESULT paragraph is the DECLARATION, committed before any fit ran, so it can be read against the result.

WHY. The catalogue record carried `vaporPressure \{ model AmbroseWalton; \}`, a corresponding-states ESTIMATE from Tc, Pc and omega that never reads the record's own Tb (C279 measured a median 19.6 \% AAD over the 25 such records with archive data). Here a three-parameter Antoine is fitted to EVERY usable laboratory's measurements of it.

STUDIES (one per article, extracted by `choupo-thermoml extract-vp`; ranked by points, ties by DOI): 10.1016/j.fluid.2006.06.017 26 pts 346.03-413.67 K Lladosa; Monton; Burguet; Munoz 10.1016/j.fluid.2015.05.002 14 pts 430.15-508.15 K Liebergesell; Kaminski; Pauls; de Loos et al. 10.1016/j.fluid.2007.09.002 10 pts 263.15-343.15 K Boukais-Belaribi; Belaribi; Lohmann; Jose 10.1021/je7007523 7 pts 313.15-313.15 K Alonso-Tristan; Villamanan; Chamorro; Segovia SET ASIDE BY NAME: 10.1016/j.fluid.2005.06.017: 1 pressure point(s), fewer than 5 10.1016/j.fluid.2005.11.026: 3 pressure point(s), fewer than 5 10.1016/j.fluid.2012.01.008: 1 pressure point(s), fewer than 5 10.1021/acs.jced.7b00119: 1 pressure point(s), fewer than 5 10.1021/je060237w: 4 pressure point(s), fewer than 5 10.1021/je100019r: 3 pressure point(s), fewer than 5

FIT AND GRADED ACCEPTANCE (DEV.md 4c C285, Vitor 2026-11-10; the evidence, objective and outliers of C283; declared before the fit). EVIDENCE: every study above is fitted; none is held out (C283 amended: laboratories disagree with each other by more than they claim, so the test is agreement across all of them, not one laboratory's claimed precision). OBJECTIVE: `objective \{ method twoStage; \}` -- least squares in log10 P for a first estimate, then the LARGEST relative error in P, (P\_model - P)/P, over the kept points minimised from it (Vitor's practice, C280); both answers kept. OUTLIERS: after the least-squares stage a point with a relative P error above 15 \% (three times the acceptance) is DISCARDED BY NAME (its study DOI, T and P), listed in the result and the record, and the fit repeated without it, once (`outliers \{ maxRelativeError 15 percent; \}`). ACCEPTANCE, GRADED, the DEFAULT of C285: the maximum relative P error over all kept points is graded clean up to 5 \%, warning up to 10 \%, severe up to 20 \% and refused above; where the normal-boiling-point condition below applies, the absolute gap is graded clean up to 1 K, warning up to 2 K, severe up to 4 K and refused above. The fit's grade is the WORSE of the two. A clean, warning or severe fit is ACCEPTED and promoted (autoPromoted, humanReviewed false) with `fitGrade <grade>;` and the deciding numbers in its provenance entry, which the component loader voices on every run that uses it (a warning or a severe grade); a refused one is not promoted. This case declares no `system/curationDict`, so the default applies; a case that needs another acceptance declares it there, with its reason. The AmbroseWalton curve's errors on the same points are reported beside the fit's, as information.

NORMAL BOILING POINT: the record's Tb 415.35 K is judged against the range of the KEPT points (all studies span 263.15-508.15 K before any outlier is discarded): if Tb lies OUTSIDE that range, the fitted Antoine's own normal boiling point (P = 101.325 kPa) is GRADED by its absolute gap to it (the temperature limits of the acceptance below); inside it, the gap is reported and is not a condition.

EARLIER FITS, SUPERSEDED BY THIS DECLARATION @@EARLIER-BEGIN@@ C281 step 2, objective leastSquares (log10 P), the same partition and band: WHAT THE RUN FOUND (pinned in `expected`; measured, not argued). Antoine log10(Psat/bar) = 4.546814 - 1697.1613/(T + -39.9108), in-sample R2 0.999982 over 26 points, 346.03-413.67 K. HELD OUT: 14 points, AAD 4.7526 \% against the band 0.36 \% $\to$ notValidated. Its own normal boiling point (101.325 kPa): 413.64 K. NOT PROMOTED: held-out verdict notValidated (AAD 4.753 \% against the band 0.36 \% declared before the fit); the record keeps AmbroseWalton THEN C281 step 3, objective twoStage, one study fitted and one held out, judged by the band (superseded by C283): WHAT THE RUN FOUND (pinned in `expected`; measured, not argued). Objective twoStage. Antoine log10(Psat/bar) = 4.644668 - 1764.4983/(T + -33.2244), in-sample R2 0.999980 over 26 points, 346.03-413.67 K. HELD OUT: 14 points, AAD 5.6372 \% (max 12.1202 \%) against the band 0.36 \% $\to$ notValidated; the stage-1 least-squares answer on the same points: AAD 4.7526 \%, max 10.3021 \%. Its own normal boiling point (101.325 kPa): 413.59 K. NOT PROMOTED: held-out verdict notValidated (AAD 5.637 \% against the band 0.36 \% declared before the fit); the record keeps AmbroseWalton. THEN C283, every usable study fitted and none held out, accepted only at max relative error <= 5 \% and Tnb within 1 K (superseded by the graded acceptance of C285): WHAT THE RUN FOUND (pinned in `expected`; measured, not argued). Objective twoStage. Antoine log10(Psat/bar) = 4.226260 - 1490.0735/(T + -60.6247), in-sample R2 0.999931 over 57 kept points, 263.15-508.15 K (10.1016/j.fluid.2006.06.017 26/26 kept; 10.1016/j.fluid.2015.05.002 14/14 kept; 10.1016/j.fluid.2007.09.002 10/10 kept; 10.1021/je7007523 7/7 kept). Relative error in P over the kept points: mean 1.4580 \%, max 3.9934 \%; AmbroseWalton on the same points: mean 4.5289 \%, max 8.2943 \%. Outliers discarded: none. Its own normal boiling point (101.325 kPa): 413.68 K. PROMOTED on 2026-11-10 into data/standards/components/diNButylEther.dat as reviewStatus autoPromoted, humanReviewed false, because: under the DEFAULT acceptance of DEV.md 4c C283 (no system/curationDict), 4 studies, 57 points kept, 0 discarded as outliers; max relative error 3.993 \% (mean 1.458 \%) against 5 \%; Tb inside the kept points' range, normal boiling point 413.68 K, -1.67 K from the record's Tb 415.35 K (reported, not a condition); AmbroseWalton on the same points mean 4.529 \% / max 8.294 \% @@EARLIER-END@@

NO MEASURED POINT IS IN THIS TREE: constant/archiveDatasets declares every series by DOI, NIST URL and sha256; `bin/choupo-curate fetch` regenerates them into constant/archiveExtracts/ (gitignored).

RESULT @@RESULT-BEGIN@@ WHAT THE RUN FOUND (pinned in `expected`; measured, not argued). Objective twoStage. Antoine log10(Psat/bar) = 4.226260 - 1490.0735/(T + -60.6247), in-sample R2 0.999931 over 57 kept points, 263.15-508.15 K (10.1016/j.fluid.2006.06.017 26/26 kept; 10.1016/j.fluid.2015.05.002 14/14 kept; 10.1016/j.fluid.2007.09.002 10/10 kept; 10.1021/je7007523 7/7 kept). Relative error in P over the kept points: mean 1.4580 \%, max 3.9934 \%; AmbroseWalton on the same points: mean 4.5289 \%, max 8.2943 \%. Outliers discarded: none. Its own normal boiling point (101.325 kPa): 413.68 K. PROMOTED on 2026-11-10 into data/standards/components/diNButylEther.dat as reviewStatus autoPromoted, humanReviewed false, because: fit grade clean under the DEFAULT graded acceptance of DEV.md 4c C285 (no system/curationDict), 4 studies, 57 points kept, 0 discarded as outliers; max relative error 3.993 \% (mean 1.458 \%), grade clean (limits 5 / 10 / 20 \%); Tb inside the kept points' range, normal boiling point 413.68 K, -1.67 K from the record's Tb 415.35 K (reported, not a condition); AmbroseWalton on the same points mean 4.529 \% / max 8.294 \% @@RESULT-END@@

\end{tutnarrative}
\begin{tutgolden}
\goldrow{diag}{fit\_antoine}{A}{4.22625991854}
\goldrow{diag}{fit\_antoine}{B}{1490.07349279}
\goldrow{diag}{fit\_antoine}{C}{-60.6246869427}
\goldrow{diag}{fit\_antoine}{R2}{0.999931261634}
\goldrow{diag}{fit\_antoine}{dHvap\_kJ\_per\_mol}{36.7795331664}
\goldrow{diag}{fit\_antoine}{n\_heldout}{0}
\goldrow{diag}{fit\_antoine}{n\_points}{57}
\goldrow{verdict}{fit\_antoine}{verdict}{validationRefused}
\goldrow{diag}{fit\_antoine}{minimax\_evaluations}{903}
\goldrow{diag}{fit\_antoine}{minimax\_max\_scaled\_final}{0.0399342646522}
\goldrow{diag}{fit\_antoine}{minimax\_max\_scaled\_start}{0.0467385305979}
\goldrow{diag}{fit\_antoine}{objective\_method}{2}
\goldrow{diag}{fit\_antoine}{relErr\_max\_fit\_pct}{3.99342646522}
\goldrow{diag}{fit\_antoine}{relErr\_max\_fit\_pct\_leastSquares}{4.67385305979}
\goldmore{3}
\end{tutgolden}

\clearpage
\tutentry{curate23\_vp\_diethylethanolamine}{tutorials/props/curation/curate23\_vp\_diethylethanolamine}
\begin{tutscheme}
\begin{tikzpicture}[>={Stealth[length=2mm]}, every node/.style={font=\scriptsize}]
  \node[draw, rounded corners=2pt, align=center, minimum width=3.1cm, minimum height=0.8cm, fill=black!4] at (0.0,-0.0) {\textbf{fit\_antoine}\\[-1pt]{\tiny vaporPressureFit}};
\end{tikzpicture}
\end{tutscheme}
\tutwhat{Antoine vapour pressure of diethylethanolamine fitted by the curation pipeline to measured vapour pressures (ThermoML locates, the article sources), held out against an independent study}
\begin{tutnarrative}
the vapour pressure of diethylethanolamine, REFITTED to measurements (DEV.md 4c C281 and C283). Written by `bin/choupo-curate vp diethylethanolamine`; everything above the RESULT paragraph is the DECLARATION, committed before any fit ran, so it can be read against the result.

WHY. The catalogue record carried `vaporPressure \{ model AmbroseWalton; \}`, a corresponding-states ESTIMATE from Tc, Pc and omega that never reads the record's own Tb (C279 measured a median 19.6 \% AAD over the 25 such records with archive data). Here a three-parameter Antoine is fitted to EVERY usable laboratory's measurements of it.

STUDIES (one per article, extracted by `choupo-thermoml extract-vp`; ranked by points, ties by DOI): 10.1021/je049761y 31 pts 278.2-318.3 K Kapteina; Slovik; Verevkin; Heintz 10.1016/j.fluid.2018.05.032 24 pts 273.64-308.15 K Soares; Stejfa; Ferreira; Pinho et al. 10.1021/je101259r 13 pts 333.2-449.5 K Klepacova; Huttenhuis; Derks; Versteeg SET ASIDE BY NAME: 10.1021/je301219c: 2 pressure point(s), fewer than 5 10.1021/je4008712: 2 pressure point(s), fewer than 5

FIT AND GRADED ACCEPTANCE (DEV.md 4c C285, Vitor 2026-11-10; the evidence, objective and outliers of C283; declared before the fit). EVIDENCE: every study above is fitted; none is held out (C283 amended: laboratories disagree with each other by more than they claim, so the test is agreement across all of them, not one laboratory's claimed precision). OBJECTIVE: `objective \{ method twoStage; \}` -- least squares in log10 P for a first estimate, then the LARGEST relative error in P, (P\_model - P)/P, over the kept points minimised from it (Vitor's practice, C280); both answers kept. OUTLIERS: after the least-squares stage a point with a relative P error above 15 \% (three times the acceptance) is DISCARDED BY NAME (its study DOI, T and P), listed in the result and the record, and the fit repeated without it, once (`outliers \{ maxRelativeError 15 percent; \}`). ACCEPTANCE, GRADED, the DEFAULT of C285: the maximum relative P error over all kept points is graded clean up to 5 \%, warning up to 10 \%, severe up to 20 \% and refused above; where the normal-boiling-point condition below applies, the absolute gap is graded clean up to 1 K, warning up to 2 K, severe up to 4 K and refused above. The fit's grade is the WORSE of the two. A clean, warning or severe fit is ACCEPTED and promoted (autoPromoted, humanReviewed false) with `fitGrade <grade>;` and the deciding numbers in its provenance entry, which the component loader voices on every run that uses it (a warning or a severe grade); a refused one is not promoted. This case declares no `system/curationDict`, so the default applies; a case that needs another acceptance declares it there, with its reason. The AmbroseWalton curve's errors on the same points are reported beside the fit's, as information.

NORMAL BOILING POINT: the record's Tb 434.15 K is judged against the range of the KEPT points (all studies span 273.64-449.5 K before any outlier is discarded): if Tb lies OUTSIDE that range, the fitted Antoine's own normal boiling point (P = 101.325 kPa) is GRADED by its absolute gap to it (the temperature limits of the acceptance below); inside it, the gap is reported and is not a condition.

EARLIER FITS, SUPERSEDED BY THIS DECLARATION @@EARLIER-BEGIN@@ C281 step 2, objective leastSquares (log10 P), the same partition and band: WHAT THE RUN FOUND (pinned in `expected`; measured, not argued). Antoine log10(Psat/bar) = 5.211269 - 1996.8790/(T + -43.7040), in-sample R2 0.999590 over 31 points, 278.2-318.3 K. HELD OUT: 24 points, AAD 3.5664 \% against the band 1.7 \% $\to$ notValidated. Its own normal boiling point (101.325 kPa): 427.31 K. NOT PROMOTED: held-out verdict notValidated (AAD 3.566 \% against the band 1.7 \% declared before the fit); the record keeps AmbroseWalton THEN C281 step 3, objective twoStage, one study fitted and one held out, judged by the band (superseded by C283): WHAT THE RUN FOUND (pinned in `expected`; measured, not argued). Objective twoStage. Antoine log10(Psat/bar) = 5.890606 - 2328.2686/(T + -24.9514), in-sample R2 0.999276 over 31 points, 278.2-318.3 K. HELD OUT: 24 points, AAD 2.6854 \% (max 3.5181 \%) against the band 1.7 \% $\to$ notValidated; the stage-1 least-squares answer on the same points: AAD 3.5664 \%, max 3.8501 \%. Its own normal boiling point (101.325 kPa): 420.59 K. NOT PROMOTED: held-out verdict notValidated (AAD 2.685 \% against the band 1.7 \% declared before the fit); the record keeps AmbroseWalton. THEN C283, every usable study fitted and none held out, accepted only at max relative error <= 5 \% and Tnb within 1 K (superseded by the graded acceptance of C285): WHAT THE RUN FOUND (pinned in `expected`; measured, not argued). Objective twoStage. Antoine log10(Psat/bar) = 4.618992 - 1719.9713/(T + -61.0754), in-sample R2 0.999872 over 68 kept points, 273.64-449.5 K (10.1021/je049761y 31/31 kept; 10.1016/j.fluid.2018.05.032 24/24 kept; 10.1021/je101259r 13/13 kept). Relative error in P over the kept points: mean 2.1349 \%, max 3.4740 \%; AmbroseWalton on the same points: mean 84.4745 \%, max 93.2983 \%. Outliers discarded: none. Its own normal boiling point (101.325 kPa): 433.91 K. PROMOTED on 2026-11-10 into data/standards/components/diethylethanolamine.dat as reviewStatus autoPromoted, humanReviewed false, because: under the DEFAULT acceptance of DEV.md 4c C283 (no system/curationDict), 3 studies, 68 points kept, 0 discarded as outliers; max relative error 3.474 \% (mean 2.135 \%) against 5 \%; Tb inside the kept points' range, normal boiling point 433.91 K, -0.24 K from the record's Tb 434.15 K (reported, not a condition); AmbroseWalton on the same points mean 84.47 \% / max 93.3 \% @@EARLIER-END@@

NO MEASURED POINT IS IN THIS TREE: constant/archiveDatasets declares every series by DOI, NIST URL and sha256; `bin/choupo-curate fetch` regenerates them into constant/archiveExtracts/ (gitignored).

RESULT @@RESULT-BEGIN@@ WHAT THE RUN FOUND (pinned in `expected`; measured, not argued). Objective twoStage. Antoine log10(Psat/bar) = 4.618992 - 1719.9713/(T + -61.0754), in-sample R2 0.999872 over 68 kept points, 273.64-449.5 K (10.1021/je049761y 31/31 kept; 10.1016/j.fluid.2018.05.032 24/24 kept; 10.1021/je101259r 13/13 kept). Relative error in P over the kept points: mean 2.1349 \%, max 3.4740 \%; AmbroseWalton on the same points: mean 84.4745 \%, max 93.2983 \%. Outliers discarded: none. Its own normal boiling point (101.325 kPa): 433.91 K. PROMOTED on 2026-11-10 into data/standards/components/diethylethanolamine.dat as reviewStatus autoPromoted, humanReviewed false, because: fit grade clean under the DEFAULT graded acceptance of DEV.md 4c C285 (no system/curationDict), 3 studies, 68 points kept, 0 discarded as outliers; max relative error 3.474 \% (mean 2.135 \%), grade clean (limits 5 / 10 / 20 \%); Tb inside the kept points' range, normal boiling point 433.91 K, -0.24 K from the record's Tb 434.15 K (reported, not a condition); AmbroseWalton on the same points mean 84.47 \% / max 93.3 \% @@RESULT-END@@

\end{tutnarrative}
\begin{tutgolden}
\goldrow{diag}{fit\_antoine}{A}{4.61899228234}
\goldrow{diag}{fit\_antoine}{B}{1719.97134688}
\goldrow{diag}{fit\_antoine}{C}{-61.0753858192}
\goldrow{diag}{fit\_antoine}{R2}{0.999872075818}
\goldrow{diag}{fit\_antoine}{dHvap\_kJ\_per\_mol}{44.0978069473}
\goldrow{diag}{fit\_antoine}{n\_heldout}{0}
\goldrow{diag}{fit\_antoine}{n\_points}{68}
\goldrow{verdict}{fit\_antoine}{verdict}{validationRefused}
\goldrow{diag}{fit\_antoine}{minimax\_evaluations}{810}
\goldrow{diag}{fit\_antoine}{minimax\_max\_scaled\_final}{0.0347397782803}
\goldrow{diag}{fit\_antoine}{minimax\_max\_scaled\_start}{0.0500051831682}
\goldrow{diag}{fit\_antoine}{objective\_method}{2}
\goldrow{diag}{fit\_antoine}{relErr\_max\_fit\_pct}{3.47397782803}
\goldrow{diag}{fit\_antoine}{relErr\_max\_fit\_pct\_leastSquares}{5.00051831682}
\goldmore{3}
\end{tutgolden}

\clearpage
\tutentry{curate24\_vp\_dimethylSulfoxide}{tutorials/props/curation/curate24\_vp\_dimethylSulfoxide}
\begin{tutscheme}
\begin{tikzpicture}[>={Stealth[length=2mm]}, every node/.style={font=\scriptsize}]
  \node[draw, rounded corners=2pt, align=center, minimum width=3.1cm, minimum height=0.8cm, fill=black!4] at (0.0,-0.0) {\textbf{fit\_antoine}\\[-1pt]{\tiny vaporPressureFit}};
\end{tikzpicture}
\end{tutscheme}
\tutwhat{Antoine vapour pressure of dimethylSulfoxide fitted by the curation pipeline to measured vapour pressures (ThermoML locates, the article sources), held out against an independent study}
\begin{tutnarrative}
the vapour pressure of dimethylSulfoxide, REFITTED to measurements (DEV.md 4c C281 and C283). Written by `bin/choupo-curate vp dimethylSulfoxide`; everything above the RESULT paragraph is the DECLARATION, committed before any fit ran, so it can be read against the result.

WHY. The catalogue record carried `vaporPressure \{ model AmbroseWalton; \}`, a corresponding-states ESTIMATE from Tc, Pc and omega that never reads the record's own Tb (C279 measured a median 19.6 \% AAD over the 25 such records with archive data). Here a three-parameter Antoine is fitted to EVERY usable laboratory's measurements of it.

STUDIES (one per article, extracted by `choupo-thermoml extract-vp`; ranked by points, ties by DOI): 10.1016/j.fluid.2011.02.002 21 pts 293.15-308.16 K Fulem; Ruzicka; Ruzicka 10.1021/je400531a 19 pts 401.95-463.45 K Zhang; Lv; Huang; Jia et al. 10.1016/j.jct.2018.04.014 9 pts 353.1-433.1 K de Pablo; Segovia; Martin; Martin et al. SET ASIDE BY NAME: 10.1016/j.fluid.2007.08.001: 1 pressure point(s), fewer than 5 10.1016/j.jct.2007.04.004: 1 pressure point(s), fewer than 5 10.1016/j.jct.2018.12.016: 1 pressure point(s), fewer than 5 10.1016/j.jct.2019.03.036: 1 pressure point(s), fewer than 5 10.1021/acs.jced.6b00040: 1 pressure point(s), fewer than 5 10.1021/acs.jced.6b00733: 1 pressure point(s), fewer than 5 10.1021/je700408x: 1 pressure point(s), fewer than 5

FIT AND GRADED ACCEPTANCE (DEV.md 4c C285, Vitor 2026-11-10; the evidence, objective and outliers of C283; declared before the fit). EVIDENCE: every study above is fitted; none is held out (C283 amended: laboratories disagree with each other by more than they claim, so the test is agreement across all of them, not one laboratory's claimed precision). OBJECTIVE: `objective \{ method twoStage; \}` -- least squares in log10 P for a first estimate, then the LARGEST relative error in P, (P\_model - P)/P, over the kept points minimised from it (Vitor's practice, C280); both answers kept. OUTLIERS: after the least-squares stage a point with a relative P error above 15 \% (three times the acceptance) is DISCARDED BY NAME (its study DOI, T and P), listed in the result and the record, and the fit repeated without it, once (`outliers \{ maxRelativeError 15 percent; \}`). ACCEPTANCE, GRADED, the DEFAULT of C285: the maximum relative P error over all kept points is graded clean up to 5 \%, warning up to 10 \%, severe up to 20 \% and refused above; where the normal-boiling-point condition below applies, the absolute gap is graded clean up to 1 K, warning up to 2 K, severe up to 4 K and refused above. The fit's grade is the WORSE of the two. A clean, warning or severe fit is ACCEPTED and promoted (autoPromoted, humanReviewed false) with `fitGrade <grade>;` and the deciding numbers in its provenance entry, which the component loader voices on every run that uses it (a warning or a severe grade); a refused one is not promoted. This case declares no `system/curationDict`, so the default applies; a case that needs another acceptance declares it there, with its reason. The AmbroseWalton curve's errors on the same points are reported beside the fit's, as information.

NORMAL BOILING POINT: the record's Tb 462.15 K is judged against the range of the KEPT points (all studies span 293.15-463.45 K before any outlier is discarded): if Tb lies OUTSIDE that range, the fitted Antoine's own normal boiling point (P = 101.325 kPa) is GRADED by its absolute gap to it (the temperature limits of the acceptance below); inside it, the gap is reported and is not a condition.

EARLIER FITS, SUPERSEDED BY THIS DECLARATION @@EARLIER-BEGIN@@ C281 step 2, objective leastSquares (log10 P), the same partition and band: WHAT THE RUN FOUND (pinned in `expected`; measured, not argued). Antoine log10(Psat/bar) = 4.482417 - 1835.6725/(T + -56.0230), in-sample R2 0.999999 over 21 points, 293.15-308.16 K. HELD OUT: 19 points, AAD 3.8716 \% against the band 3.2 \% $\to$ notValidated. Its own normal boiling point (101.325 kPa): 466.07 K. NOT PROMOTED: held-out verdict notValidated (AAD 3.872 \% against the band 3.2 \% declared before the fit); the record keeps AmbroseWalton THEN C281 step 3, objective twoStage, one study fitted and one held out, judged by the band (superseded by C283): WHAT THE RUN FOUND (pinned in `expected`; measured, not argued). Objective twoStage. Antoine log10(Psat/bar) = 4.515639 - 1850.8250/(T + -55.0893), in-sample R2 0.999999 over 21 points, 293.15-308.16 K. HELD OUT: 19 points, AAD 2.7312 \% (max 5.0537 \%) against the band 3.2 \% $\to$ validated; the stage-1 least-squares answer on the same points: AAD 3.8716 \%, max 6.4461 \%. Its own normal boiling point (101.325 kPa): 465.48 K. PROMOTED (and since REVERTED by `choupo-curate revert`) on 2026-11-10 into data/standards/components/dimethylSulfoxide.dat as reviewStatus autoPromoted, humanReviewed false, because: verdict validated (the record's Tb is inside the measured span or not declared; the fit's own normal boiling point is 465.48 K). THEN C283, every usable study fitted and none held out, accepted only at max relative error <= 5 \% and Tnb within 1 K (superseded by the graded acceptance of C285): WHAT THE RUN FOUND (pinned in `expected`; measured, not argued). Objective twoStage. Antoine log10(Psat/bar) = 4.693650 - 1946.0170/(T + -48.5201), in-sample R2 0.999995 over 49 kept points, 293.15-463.45 K (10.1016/j.fluid.2011.02.002 21/21 kept; 10.1021/je400531a 19/19 kept; 10.1016/j.jct.2018.04.014 9/9 kept). Relative error in P over the kept points: mean 0.5975 \%, max 0.7943 \%; AmbroseWalton on the same points: mean 21.3045 \%, max 35.9060 \%. Outliers discarded: none. Its own normal boiling point (101.325 kPa): 463.63 K. PROMOTED on 2026-11-10 into data/standards/components/dimethylSulfoxide.dat as reviewStatus autoPromoted, humanReviewed false, because: under the DEFAULT acceptance of DEV.md 4c C283 (no system/curationDict), 3 studies, 49 points kept, 0 discarded as outliers; max relative error 0.7943 \% (mean 0.5975 \%) against 5 \%; Tb inside the kept points' range, normal boiling point 463.63 K, +1.48 K from the record's Tb 462.15 K (reported, not a condition); AmbroseWalton on the same points mean 21.3 \% / max 35.91 \% @@EARLIER-END@@

NO MEASURED POINT IS IN THIS TREE: constant/archiveDatasets declares every series by DOI, NIST URL and sha256; `bin/choupo-curate fetch` regenerates them into constant/archiveExtracts/ (gitignored).

RESULT @@RESULT-BEGIN@@ WHAT THE RUN FOUND (pinned in `expected`; measured, not argued). Objective twoStage. Antoine log10(Psat/bar) = 4.693650 - 1946.0170/(T + -48.5201), in-sample R2 0.999995 over 49 kept points, 293.15-463.45 K (10.1016/j.fluid.2011.02.002 21/21 kept; 10.1021/je400531a 19/19 kept; 10.1016/j.jct.2018.04.014 9/9 kept). Relative error in P over the kept points: mean 0.5975 \%, max 0.7943 \%; AmbroseWalton on the same points: mean 21.3045 \%, max 35.9060 \%. Outliers discarded: none. Its own normal boiling point (101.325 kPa): 463.63 K. PROMOTED on 2026-11-10 into data/standards/components/dimethylSulfoxide.dat as reviewStatus autoPromoted, humanReviewed false, because: fit grade clean under the DEFAULT graded acceptance of DEV.md 4c C285 (no system/curationDict), 3 studies, 49 points kept, 0 discarded as outliers; max relative error 0.7943 \% (mean 0.5975 \%), grade clean (limits 5 / 10 / 20 \%); Tb inside the kept points' range, normal boiling point 463.63 K, +1.48 K from the record's Tb 462.15 K (reported, not a condition); AmbroseWalton on the same points mean 21.3 \% / max 35.91 \% @@RESULT-END@@

\end{tutnarrative}
\begin{tutgolden}
\goldrow{diag}{fit\_antoine}{A}{4.69364955777}
\goldrow{diag}{fit\_antoine}{B}{1946.01699011}
\goldrow{diag}{fit\_antoine}{C}{-48.5200814042}
\goldrow{diag}{fit\_antoine}{R2}{0.999995343342}
\goldrow{diag}{fit\_antoine}{dHvap\_kJ\_per\_mol}{46.478571606}
\goldrow{diag}{fit\_antoine}{n\_heldout}{0}
\goldrow{diag}{fit\_antoine}{n\_points}{49}
\goldrow{verdict}{fit\_antoine}{verdict}{validationRefused}
\goldrow{diag}{fit\_antoine}{minimax\_evaluations}{710}
\goldrow{diag}{fit\_antoine}{minimax\_max\_scaled\_final}{0.00794320890809}
\goldrow{diag}{fit\_antoine}{minimax\_max\_scaled\_start}{0.0102846733939}
\goldrow{diag}{fit\_antoine}{objective\_method}{2}
\goldrow{diag}{fit\_antoine}{relErr\_max\_fit\_pct}{0.794320890809}
\goldrow{diag}{fit\_antoine}{relErr\_max\_fit\_pct\_leastSquares}{1.02846733939}
\goldmore{3}
\end{tutgolden}

\clearpage
\tutentry{curate25\_vp\_dimethylethanolamine}{tutorials/props/curation/curate25\_vp\_dimethylethanolamine}
\begin{tutscheme}
\begin{tikzpicture}[>={Stealth[length=2mm]}, every node/.style={font=\scriptsize}]
  \node[draw, rounded corners=2pt, align=center, minimum width=3.1cm, minimum height=0.8cm, fill=black!4] at (0.0,-0.0) {\textbf{fit\_antoine}\\[-1pt]{\tiny vaporPressureFit}};
\end{tikzpicture}
\end{tutscheme}
\tutwhat{Antoine vapour pressure of dimethylethanolamine fitted by the curation pipeline to measured vapour pressures (ThermoML locates, the article sources), held out against an independent study}
\begin{tutnarrative}
the vapour pressure of dimethylethanolamine, REFITTED to measurements (DEV.md 4c C281 and C283). Written by `bin/choupo-curate vp dimethylethanolamine`; everything above the RESULT paragraph is the DECLARATION, committed before any fit ran, so it can be read against the result.

WHY. The catalogue record carried `vaporPressure \{ model AmbroseWalton; \}`, a corresponding-states ESTIMATE from Tc, Pc and omega that never reads the record's own Tb (C279 measured a median 19.6 \% AAD over the 25 such records with archive data). Here a three-parameter Antoine is fitted to EVERY usable laboratory's measurements of it.

STUDIES (one per article, extracted by `choupo-thermoml extract-vp`; ranked by points, ties by DOI): 10.1021/je049761y 29 pts 277.9-316.3 K Kapteina; Slovik; Verevkin; Heintz 10.1016/j.jct.2013.03.020 25 pts 273.15-363.15 K Chiali-Baba-Ahmed; Dergal; Negadi; Mokbel 10.1021/je101259r 15 pts 309.5-447.8 K Klepacova; Huttenhuis; Derks; Versteeg SET ASIDE BY NAME: (none)

FIT AND GRADED ACCEPTANCE (DEV.md 4c C285, Vitor 2026-11-10; the evidence, objective and outliers of C283; declared before the fit). EVIDENCE: every study above is fitted; none is held out (C283 amended: laboratories disagree with each other by more than they claim, so the test is agreement across all of them, not one laboratory's claimed precision). OBJECTIVE: `objective \{ method twoStage; \}` -- least squares in log10 P for a first estimate, then the LARGEST relative error in P, (P\_model - P)/P, over the kept points minimised from it (Vitor's practice, C280); both answers kept. OUTLIERS: after the least-squares stage a point with a relative P error above 15 \% (three times the acceptance) is DISCARDED BY NAME (its study DOI, T and P), listed in the result and the record, and the fit repeated without it, once (`outliers \{ maxRelativeError 15 percent; \}`). ACCEPTANCE, GRADED, the DEFAULT of C285: the maximum relative P error over all kept points is graded clean up to 5 \%, warning up to 10 \%, severe up to 20 \% and refused above; where the normal-boiling-point condition below applies, the absolute gap is graded clean up to 1 K, warning up to 2 K, severe up to 4 K and refused above. The fit's grade is the WORSE of the two. A clean, warning or severe fit is ACCEPTED and promoted (autoPromoted, humanReviewed false) with `fitGrade <grade>;` and the deciding numbers in its provenance entry, which the component loader voices on every run that uses it (a warning or a severe grade); a refused one is not promoted. This case declares no `system/curationDict`, so the default applies; a case that needs another acceptance declares it there, with its reason. The AmbroseWalton curve's errors on the same points are reported beside the fit's, as information.

NORMAL BOILING POINT: the record's Tb 404.15 K is judged against the range of the KEPT points (all studies span 273.15-447.8 K before any outlier is discarded): if Tb lies OUTSIDE that range, the fitted Antoine's own normal boiling point (P = 101.325 kPa) is GRADED by its absolute gap to it (the temperature limits of the acceptance below); inside it, the gap is reported and is not a condition.

EARLIER FITS, SUPERSEDED BY THIS DECLARATION @@EARLIER-BEGIN@@ C281 step 2, objective leastSquares (log10 P), the same partition and band: WHAT THE RUN FOUND (pinned in `expected`; measured, not argued). Antoine log10(Psat/bar) = 14.191855 - 9599.5633/(T + 289.4230), in-sample R2 0.998529 over 29 points, 277.9-316.3 K. HELD OUT: 25 points, AAD 12.9406 \% against the band 2.2 \% $\to$ notValidated. Its own normal boiling point (101.325 kPa): 387.26 K. NOT PROMOTED: held-out verdict notValidated (AAD 12.94 \% against the band 2.2 \% declared before the fit); the record keeps AmbroseWalton THEN C281 step 3, objective twoStage, one study fitted and one held out, judged by the band (superseded by C283): WHAT THE RUN FOUND (pinned in `expected`; measured, not argued). Objective twoStage. Antoine log10(Psat/bar) = 7.287831 - 3246.8966/(T + 46.1092), in-sample R2 0.998315 over 29 points, 277.9-316.3 K. HELD OUT: 25 points, AAD 8.3129 \% (max 11.9823 \%) against the band 2.2 \% $\to$ notValidated; the stage-1 least-squares answer on the same points: AAD 12.9406 \%, max 40.0319 \%. Its own normal boiling point (101.325 kPa): 399.76 K. NOT PROMOTED: held-out verdict notValidated (AAD 8.313 \% against the band 2.2 \% declared before the fit); the record keeps AmbroseWalton. THEN C283, every usable study fitted and none held out, accepted only at max relative error <= 5 \% and Tnb within 1 K (superseded by the graded acceptance of C285): WHAT THE RUN FOUND (pinned in `expected`; measured, not argued). Objective twoStage. Antoine log10(Psat/bar) = 5.281864 - 2003.2539/(T + -27.7427), in-sample R2 0.999092 over 69 kept points, 273.15-447.8 K (10.1021/je049761y 29/29 kept; 10.1016/j.jct.2013.03.020 25/25 kept; 10.1021/je101259r 15/15 kept). Relative error in P over the kept points: mean 5.3080 \%, max 8.5056 \%; AmbroseWalton on the same points: mean 37.1625 \%, max 63.9790 \%. Outliers discarded: none. Its own normal boiling point (101.325 kPa): 407.42 K. NOT PROMOTED by the DEFAULT acceptance of DEV.md 4c C283 (no system/curationDict): 3 studies, 69 points kept, 0 discarded as outliers; max relative error 8.506 \% (mean 5.308 \%) against 5 \%; AmbroseWalton on the same points mean 37.16 \% / max 63.98 \%; the record keeps AmbroseWalton @@EARLIER-END@@

NO MEASURED POINT IS IN THIS TREE: constant/archiveDatasets declares every series by DOI, NIST URL and sha256; `bin/choupo-curate fetch` regenerates them into constant/archiveExtracts/ (gitignored).

RESULT @@RESULT-BEGIN@@ WHAT THE RUN FOUND (pinned in `expected`; measured, not argued). Objective twoStage. Antoine log10(Psat/bar) = 5.281864 - 2003.2539/(T + -27.7427), in-sample R2 0.999092 over 69 kept points, 273.15-447.8 K (10.1021/je049761y 29/29 kept; 10.1016/j.jct.2013.03.020 25/25 kept; 10.1021/je101259r 15/15 kept). Relative error in P over the kept points: mean 5.3080 \%, max 8.5056 \%; AmbroseWalton on the same points: mean 37.1625 \%, max 63.9790 \%. Outliers discarded: none. Its own normal boiling point (101.325 kPa): 407.42 K. PROMOTED on 2026-11-10 into data/standards/components/dimethylethanolamine.dat as reviewStatus autoPromoted, humanReviewed false, because: fit grade warning under the DEFAULT graded acceptance of DEV.md 4c C285 (no system/curationDict), 3 studies, 69 points kept, 0 discarded as outliers; max relative error 8.506 \% (mean 5.308 \%), grade warning (limits 5 / 10 / 20 \%); Tb inside the kept points' range, normal boiling point 407.42 K, +3.27 K from the record's Tb 404.15 K (reported, not a condition); AmbroseWalton on the same points mean 37.16 \% / max 63.98 \% @@RESULT-END@@

\end{tutnarrative}
\begin{tutgolden}
\goldrow{diag}{fit\_antoine}{A}{5.2818640296}
\goldrow{diag}{fit\_antoine}{B}{2003.25390015}
\goldrow{diag}{fit\_antoine}{C}{-27.7426689283}
\goldrow{diag}{fit\_antoine}{R2}{0.999092262955}
\goldrow{diag}{fit\_antoine}{dHvap\_kJ\_per\_mol}{43.5849861054}
\goldrow{diag}{fit\_antoine}{n\_heldout}{0}
\goldrow{diag}{fit\_antoine}{n\_points}{69}
\goldrow{verdict}{fit\_antoine}{verdict}{validationRefused}
\goldrow{diag}{fit\_antoine}{minimax\_evaluations}{870}
\goldrow{diag}{fit\_antoine}{minimax\_max\_scaled\_final}{0.0850561402704}
\goldrow{diag}{fit\_antoine}{minimax\_max\_scaled\_start}{0.11674633769}
\goldrow{diag}{fit\_antoine}{objective\_method}{2}
\goldrow{diag}{fit\_antoine}{relErr\_max\_fit\_pct}{8.50561402704}
\goldrow{diag}{fit\_antoine}{relErr\_max\_fit\_pct\_leastSquares}{11.674633769}
\goldmore{3}
\end{tutgolden}

\clearpage
\tutentry{curate26\_vp\_furfural}{tutorials/props/curation/curate26\_vp\_furfural}
\begin{tutscheme}
\begin{tikzpicture}[>={Stealth[length=2mm]}, every node/.style={font=\scriptsize}]
  \node[draw, rounded corners=2pt, align=center, minimum width=3.1cm, minimum height=0.8cm, fill=black!4] at (0.0,-0.0) {\textbf{fit\_antoine}\\[-1pt]{\tiny vaporPressureFit}};
\end{tikzpicture}
\end{tutscheme}
\tutwhat{Antoine vapour pressure of furfural fitted by the curation pipeline to measured vapour pressures (ThermoML locates, the article sources), held out against an independent study}
\begin{tutnarrative}
the vapour pressure of furfural, REFITTED to measurements (DEV.md 4c C281 and C283). Written by `bin/choupo-curate vp furfural`; everything above the RESULT paragraph is the DECLARATION, committed before any fit ran, so it can be read against the result.

WHY. The catalogue record carried `vaporPressure \{ model AmbroseWalton; \}`, a corresponding-states ESTIMATE from Tc, Pc and omega that never reads the record's own Tb (C279 measured a median 19.6 \% AAD over the 25 such records with archive data). Here a three-parameter Antoine is fitted to EVERY usable laboratory's measurements of it.

STUDIES (one per article, extracted by `choupo-thermoml extract-vp`; ranked by points, ties by DOI): 10.1021/acs.jced.7b00574 19 pts 326.05-415.15 K Zheng; Luo; Yin; Chen et al. 10.1021/je060406c 16 pts 276.7-323.4 K Emel'yanenko; Dabrowska; Verevkin; Hertel et al. 10.1021/je020117y 9 pts 339.16-429.5 K Fele; Grilc SET ASIDE BY NAME: 10.1016/j.fluid.2014.10.037: 3 pressure point(s), fewer than 5 10.1016/j.jct.2018.03.019: 1 pressure point(s), fewer than 5 10.1021/acs.jced.8b00451: 1 pressure point(s), fewer than 5 10.1021/je400885z: 4 pressure point(s), fewer than 5

FIT AND GRADED ACCEPTANCE (DEV.md 4c C285, Vitor 2026-11-10; the evidence, objective and outliers of C283; declared before the fit). EVIDENCE: every study above is fitted; none is held out (C283 amended: laboratories disagree with each other by more than they claim, so the test is agreement across all of them, not one laboratory's claimed precision). OBJECTIVE: `objective \{ method twoStage; \}` -- least squares in log10 P for a first estimate, then the LARGEST relative error in P, (P\_model - P)/P, over the kept points minimised from it (Vitor's practice, C280); both answers kept. OUTLIERS: after the least-squares stage a point with a relative P error above 15 \% (three times the acceptance) is DISCARDED BY NAME (its study DOI, T and P), listed in the result and the record, and the fit repeated without it, once (`outliers \{ maxRelativeError 15 percent; \}`). ACCEPTANCE, GRADED, the DEFAULT of C285: the maximum relative P error over all kept points is graded clean up to 5 \%, warning up to 10 \%, severe up to 20 \% and refused above; where the normal-boiling-point condition below applies, the absolute gap is graded clean up to 1 K, warning up to 2 K, severe up to 4 K and refused above. The fit's grade is the WORSE of the two. A clean, warning or severe fit is ACCEPTED and promoted (autoPromoted, humanReviewed false) with `fitGrade <grade>;` and the deciding numbers in its provenance entry, which the component loader voices on every run that uses it (a warning or a severe grade); a refused one is not promoted. This case declares no `system/curationDict`, so the default applies; a case that needs another acceptance declares it there, with its reason. The AmbroseWalton curve's errors on the same points are reported beside the fit's, as information.

NORMAL BOILING POINT: the record's Tb 434.85 K is judged against the range of the KEPT points (all studies span 276.7-429.5 K before any outlier is discarded): if Tb lies OUTSIDE that range, the fitted Antoine's own normal boiling point (P = 101.325 kPa) is GRADED by its absolute gap to it (the temperature limits of the acceptance below); inside it, the gap is reported and is not a condition.

EARLIER FITS, SUPERSEDED BY THIS DECLARATION @@EARLIER-BEGIN@@ C281 step 2, objective leastSquares (log10 P), the same partition and band: WHAT THE RUN FOUND (pinned in `expected`; measured, not argued). Antoine log10(Psat/bar) = 3.638107 - 1181.7414/(T + -109.8817), in-sample R2 0.999802 over 19 points, 326.05-415.15 K. HELD OUT: 16 points, AAD 10.8740 \% against the band 5.6 \% $\to$ notValidated. Its own normal boiling point (101.325 kPa): 435.22 K. NOT PROMOTED: held-out verdict notValidated (AAD 10.87 \% against the band 5.6 \% declared before the fit); the record keeps AmbroseWalton THEN C281 step 3, objective twoStage, one study fitted and one held out, judged by the band (superseded by C283): WHAT THE RUN FOUND (pinned in `expected`; measured, not argued). Objective twoStage. Antoine log10(Psat/bar) = 3.762609 - 1251.4515/(T + -102.1178), in-sample R2 0.999771 over 19 points, 326.05-415.15 K. HELD OUT: 16 points, AAD 8.1701 \% (max 25.1514 \%) against the band 5.6 \% $\to$ notValidated; the stage-1 least-squares answer on the same points: AAD 10.8740 \%, max 31.7748 \%. Its own normal boiling point (101.325 kPa): 435.23 K. NOT PROMOTED: held-out verdict notValidated (AAD 8.17 \% against the band 5.6 \% declared before the fit); the record keeps AmbroseWalton. THEN C283, every usable study fitted and none held out, accepted only at max relative error <= 5 \% and Tnb within 1 K (superseded by the graded acceptance of C285): WHAT THE RUN FOUND (pinned in `expected`; measured, not argued). Objective twoStage. Antoine log10(Psat/bar) = 4.578554 - 1695.1936/(T + -61.6013), in-sample R2 0.999783 over 44 kept points, 276.7-429.5 K (10.1021/acs.jced.7b00574 19/19 kept; 10.1021/je060406c 16/16 kept; 10.1021/je020117y 9/9 kept). Relative error in P over the kept points: mean 2.2435 \%, max 5.0725 \%; AmbroseWalton on the same points: mean 4.2916 \%, max 9.3434 \%. Outliers discarded: none. Its own normal boiling point (101.325 kPa): 432.31 K. NOT PROMOTED by the DEFAULT acceptance of DEV.md 4c C283 (no system/curationDict): 3 studies, 44 points kept, 0 discarded as outliers; max relative error 5.073 \% (mean 2.243 \%) against 5 \%; AmbroseWalton on the same points mean 4.292 \% / max 9.343 \%; the record keeps AmbroseWalton @@EARLIER-END@@

NO MEASURED POINT IS IN THIS TREE: constant/archiveDatasets declares every series by DOI, NIST URL and sha256; `bin/choupo-curate fetch` regenerates them into constant/archiveExtracts/ (gitignored).

RESULT @@RESULT-BEGIN@@ WHAT THE RUN FOUND (pinned in `expected`; measured, not argued). Objective twoStage. Antoine log10(Psat/bar) = 4.578554 - 1695.1936/(T + -61.6013), in-sample R2 0.999783 over 44 kept points, 276.7-429.5 K (10.1021/acs.jced.7b00574 19/19 kept; 10.1021/je060406c 16/16 kept; 10.1021/je020117y 9/9 kept). Relative error in P over the kept points: mean 2.2435 \%, max 5.0725 \%; AmbroseWalton on the same points: mean 4.2916 \%, max 9.3434 \%. Outliers discarded: none. Its own normal boiling point (101.325 kPa): 432.31 K. PROMOTED on 2026-11-10 into data/standards/components/furfural.dat as reviewStatus autoPromoted, humanReviewed false, because: fit grade severe under the DEFAULT graded acceptance of DEV.md 4c C285 (no system/curationDict), 3 studies, 44 points kept, 0 discarded as outliers; max relative error 5.073 \% (mean 2.243 \%), grade warning (limits 5 / 10 / 20 \%); Tb outside the kept points' range, normal boiling point 432.31 K, -2.54 K from the record's Tb 434.85 K, grade severe (limits 1 / 2 / 4 K); AmbroseWalton on the same points mean 4.292 \% / max 9.343 \% @@RESULT-END@@

\end{tutnarrative}
\begin{tutgolden}
\goldrow{diag}{fit\_antoine}{A}{4.57855408926}
\goldrow{diag}{fit\_antoine}{B}{1695.19356892}
\goldrow{diag}{fit\_antoine}{C}{-61.6012877477}
\goldrow{diag}{fit\_antoine}{R2}{0.999782574573}
\goldrow{diag}{fit\_antoine}{dHvap\_kJ\_per\_mol}{44.2322438155}
\goldrow{diag}{fit\_antoine}{n\_heldout}{0}
\goldrow{diag}{fit\_antoine}{n\_points}{44}
\goldrow{verdict}{fit\_antoine}{verdict}{validationRefused}
\goldrow{diag}{fit\_antoine}{minimax\_evaluations}{436}
\goldrow{diag}{fit\_antoine}{minimax\_max\_scaled\_final}{0.0507252969903}
\goldrow{diag}{fit\_antoine}{minimax\_max\_scaled\_start}{0.0623602682698}
\goldrow{diag}{fit\_antoine}{objective\_method}{2}
\goldrow{diag}{fit\_antoine}{relErr\_max\_fit\_pct}{5.07252969903}
\goldrow{diag}{fit\_antoine}{relErr\_max\_fit\_pct\_leastSquares}{6.23602682698}
\goldmore{3}
\end{tutgolden}

\clearpage
\tutentry{curate27\_vp\_glycerol}{tutorials/props/curation/curate27\_vp\_glycerol}
\begin{tutscheme}
\begin{tikzpicture}[>={Stealth[length=2mm]}, every node/.style={font=\scriptsize}]
  \node[draw, rounded corners=2pt, align=center, minimum width=3.1cm, minimum height=0.8cm, fill=black!4] at (0.0,-0.0) {\textbf{fit\_antoine}\\[-1pt]{\tiny vaporPressureFit}};
\end{tikzpicture}
\end{tutscheme}
\tutwhat{Antoine vapour pressure of glycerol fitted by the curation pipeline to measured vapour pressures (ThermoML locates, the article sources), held out against an independent study}
\begin{tutnarrative}
the vapour pressure of glycerol, REFITTED to measurements (DEV.md 4c C281 and C283). Written by `bin/choupo-curate vp glycerol`; everything above the RESULT paragraph is the DECLARATION, committed before any fit ran, so it can be read against the result.

WHY. The catalogue record carried `vaporPressure \{ model AmbroseWalton; \}`, a corresponding-states ESTIMATE from Tc, Pc and omega that never reads the record's own Tb (C279 measured a median 19.6 \% AAD over the 25 such records with archive data). Here a three-parameter Antoine is fitted to EVERY usable laboratory's measurements of it.

STUDIES (one per article, extracted by `choupo-thermoml extract-vp`; ranked by points, ties by DOI): 10.1016/j.fluid.2015.03.038 113 pts 303.59-436.13 K Verevkin; Zaitsau; Emel'yanenko; Zhabina 10.1021/je200766t 11 pts 351.8-462.86 K Mokbel; Sawaya; Zanota; Naccoul et al. 10.1016/j.jct.2012.09.032 5 pts 477.25-520.65 K Veneral; Benazzi; Mazutti; Voll et al. 10.1021/je800070y 5 pts 463.15-543.15 K Yan; Suppes SET ASIDE BY NAME: 10.1016/j.jct.2013.09.046: 4 pressure point(s), fewer than 5 10.1021/acs.jced.5b01080: 3 pressure point(s), fewer than 5

FIT AND GRADED ACCEPTANCE (DEV.md 4c C285, Vitor 2026-11-10; the evidence, objective and outliers of C283; declared before the fit). EVIDENCE: every study above is fitted; none is held out (C283 amended: laboratories disagree with each other by more than they claim, so the test is agreement across all of them, not one laboratory's claimed precision). OBJECTIVE: `objective \{ method twoStage; \}` -- least squares in log10 P for a first estimate, then the LARGEST relative error in P, (P\_model - P)/P, over the kept points minimised from it (Vitor's practice, C280); both answers kept. OUTLIERS: after the least-squares stage a point with a relative P error above 15 \% (three times the acceptance) is DISCARDED BY NAME (its study DOI, T and P), listed in the result and the record, and the fit repeated without it, once (`outliers \{ maxRelativeError 15 percent; \}`). ACCEPTANCE, GRADED, the DEFAULT of C285: the maximum relative P error over all kept points is graded clean up to 5 \%, warning up to 10 \%, severe up to 20 \% and refused above; where the normal-boiling-point condition below applies, the absolute gap is graded clean up to 1 K, warning up to 2 K, severe up to 4 K and refused above. The fit's grade is the WORSE of the two. A clean, warning or severe fit is ACCEPTED and promoted (autoPromoted, humanReviewed false) with `fitGrade <grade>;` and the deciding numbers in its provenance entry, which the component loader voices on every run that uses it (a warning or a severe grade); a refused one is not promoted. This case declares no `system/curationDict`, so the default applies; a case that needs another acceptance declares it there, with its reason. The AmbroseWalton curve's errors on the same points are reported beside the fit's, as information.

NORMAL BOILING POINT: the record's Tb 563.15 K is judged against the range of the KEPT points (all studies span 303.59-543.15 K before any outlier is discarded): if Tb lies OUTSIDE that range, the fitted Antoine's own normal boiling point (P = 101.325 kPa) is GRADED by its absolute gap to it (the temperature limits of the acceptance below); inside it, the gap is reported and is not a condition.

EARLIER FITS, SUPERSEDED BY THIS DECLARATION @@EARLIER-BEGIN@@ C281 step 2, objective leastSquares (log10 P), the same partition and band: WHAT THE RUN FOUND (pinned in `expected`; measured, not argued). Antoine log10(Psat/bar) = 6.068473 - 2994.6061/(T + -63.1209), in-sample R2 0.999986 over 113 points, 303.59-436.13 K. HELD OUT: 11 points, AAD 8.3035 \% against the band 5 \% $\to$ notValidated. Its own normal boiling point (101.325 kPa): 557.06 K. NOT PROMOTED: held-out verdict notValidated (AAD 8.303 \% against the band 5 \% declared before the fit); the record keeps AmbroseWalton THEN C281 step 3, objective twoStage, one study fitted and one held out, judged by the band (superseded by C283): WHAT THE RUN FOUND (pinned in `expected`; measured, not argued). Objective twoStage. Antoine log10(Psat/bar) = 5.889944 - 2876.5912/(T + -69.3990), in-sample R2 0.999928 over 113 points, 303.59-436.13 K. HELD OUT: 11 points, AAD 6.9909 \% (max 13.2871 \%) against the band 5 \% $\to$ notValidated; the stage-1 least-squares answer on the same points: AAD 8.3035 \%, max 15.0182 \%. Its own normal boiling point (101.325 kPa): 558.26 K. NOT PROMOTED: held-out verdict notValidated (AAD 6.991 \% against the band 5 \% declared before the fit); the record keeps AmbroseWalton. THEN C283, every usable study fitted and none held out, accepted only at max relative error <= 5 \% and Tnb within 1 K (superseded by the graded acceptance of C285): WHAT THE RUN FOUND (pinned in `expected`; measured, not argued). Objective twoStage. Antoine log10(Psat/bar) = 5.808696 - 2818.0588/(T + -72.3003), in-sample R2 0.999561 over 134 kept points, 303.59-543.15 K (10.1016/j.fluid.2015.03.038 113/113 kept; 10.1021/je200766t 11/11 kept; 10.1016/j.jct.2012.09.032 5/5 kept; 10.1021/je800070y 5/5 kept). Relative error in P over the kept points: mean 6.2927 \%, max 8.6903 \%; AmbroseWalton on the same points: mean 99.8560 \%, max 100.0000 \%. Outliers discarded: none. Its own normal boiling point (101.325 kPa): 557.92 K. NOT PROMOTED by the DEFAULT acceptance of DEV.md 4c C283 (no system/curationDict): 4 studies, 134 points kept, 0 discarded as outliers; max relative error 8.69 \% (mean 6.293 \%) against 5 \%; AmbroseWalton on the same points mean 99.86 \% / max 100 \%; the record keeps AmbroseWalton @@EARLIER-END@@

NO MEASURED POINT IS IN THIS TREE: constant/archiveDatasets declares every series by DOI, NIST URL and sha256; `bin/choupo-curate fetch` regenerates them into constant/archiveExtracts/ (gitignored).

RESULT @@RESULT-BEGIN@@ WHAT THE RUN FOUND (pinned in `expected`; measured, not argued). Objective twoStage. Antoine log10(Psat/bar) = 5.808696 - 2818.0588/(T + -72.3003), in-sample R2 0.999561 over 134 kept points, 303.59-543.15 K (10.1016/j.fluid.2015.03.038 113/113 kept; 10.1021/je200766t 11/11 kept; 10.1016/j.jct.2012.09.032 5/5 kept; 10.1021/je800070y 5/5 kept). Relative error in P over the kept points: mean 6.2927 \%, max 8.6903 \%; AmbroseWalton on the same points: mean 99.8560 \%, max 100.0000 \%. Outliers discarded: none. Its own normal boiling point (101.325 kPa): 557.92 K. NOT PROMOTED by the DEFAULT graded acceptance of DEV.md 4c C285 (no system/curationDict): fit grade refused -- 4 studies, 134 points kept, 0 discarded as outliers; max relative error 8.69 \% (mean 6.293 \%), grade warning (limits 5 / 10 / 20 \%); Tb outside the kept points' range, normal boiling point 557.92 K, -5.23 K from the record's Tb 563.15 K, grade refused (limits 1 / 2 / 4 K); AmbroseWalton on the same points mean 99.86 \% / max 100 \%; the record keeps AmbroseWalton @@RESULT-END@@

\end{tutnarrative}
\begin{tutgolden}
\goldrow{diag}{fit\_antoine}{A}{5.80869633206}
\goldrow{diag}{fit\_antoine}{B}{2818.05883921}
\goldrow{diag}{fit\_antoine}{C}{-72.3002910192}
\goldrow{diag}{fit\_antoine}{R2}{0.999561297891}
\goldrow{diag}{fit\_antoine}{dHvap\_kJ\_per\_mol}{71.7918071366}
\goldrow{diag}{fit\_antoine}{n\_heldout}{0}
\goldrow{diag}{fit\_antoine}{n\_points}{134}
\goldrow{verdict}{fit\_antoine}{verdict}{validationRefused}
\goldrow{diag}{fit\_antoine}{minimax\_evaluations}{882}
\goldrow{diag}{fit\_antoine}{minimax\_max\_scaled\_final}{0.0869029077882}
\goldrow{diag}{fit\_antoine}{minimax\_max\_scaled\_start}{0.142936948672}
\goldrow{diag}{fit\_antoine}{objective\_method}{2}
\goldrow{diag}{fit\_antoine}{relErr\_max\_fit\_pct}{8.69029077882}
\goldrow{diag}{fit\_antoine}{relErr\_max\_fit\_pct\_leastSquares}{14.2936948672}
\goldmore{3}
\end{tutgolden}

\clearpage
\tutentry{curate287\_new\_components}{tutorials/props/curation/curate287\_new\_components}
\begin{tutscheme}
\begin{tikzpicture}[>={Stealth[length=2mm]}, every node/.style={font=\scriptsize}]
  \node[draw, rounded corners=2pt, align=center, minimum width=3.1cm, minimum height=0.8cm, fill=black!4] at (0.0,-0.0) {\textbf{joback\_2Ethoxyethanol}\\[-1pt]{\tiny estimateComponent}};
  \node[draw, rounded corners=2pt, align=center, minimum width=3.1cm, minimum height=0.8cm, fill=black!4] at (3.6,-0.0) {\textbf{joback\_ethylOctanoate}\\[-1pt]{\tiny estimateComponent}};
  \node[draw, rounded corners=2pt, align=center, minimum width=3.1cm, minimum height=0.8cm, fill=black!4] at (7.2,-0.0) {\textbf{joback\_13Propanediol}\\[-1pt]{\tiny estimateComponent}};
  \node[draw, rounded corners=2pt, align=center, minimum width=3.1cm, minimum height=0.8cm, fill=black!4] at (0.0,-1.3) {\textbf{joback\_13Butanediol}\\[-1pt]{\tiny estimateComponent}};
  \node[draw, rounded corners=2pt, align=center, minimum width=3.1cm, minimum height=0.8cm, fill=black!4] at (3.6,-1.3) {\textbf{joback\_16Hexanediol}\\[-1pt]{\tiny estimateComponent}};
  \node[draw, rounded corners=2pt, align=center, minimum width=3.1cm, minimum height=0.8cm, fill=black!4] at (7.2,-1.3) {\textbf{joback\_allylAlcohol}\\[-1pt]{\tiny estimateComponent}};
  \node[draw, rounded corners=2pt, align=center, minimum width=3.1cm, minimum height=0.8cm, fill=black!4] at (0.0,-2.6) {\textbf{joback\_33Oxydipropanol}\\[-1pt]{\tiny estimateComponent}};
  \node[draw, rounded corners=2pt, align=center, minimum width=3.1cm, minimum height=0.8cm, fill=black!4] at (3.6,-2.6) {\textbf{joback\_methylLevulinate}\\[-1pt]{\tiny estimateComponent}};
  \node[draw, rounded corners=2pt, align=center, minimum width=3.1cm, minimum height=0.8cm, fill=black!4] at (7.2,-2.6) {\textbf{smoke\_2Ethoxyethanol}\\[-1pt]{\tiny propertyScan1D}};
  \node[draw, rounded corners=2pt, align=center, minimum width=3.1cm, minimum height=0.8cm, fill=black!4] at (0.0,-3.9) {\textbf{smoke\_ethylOctanoate}\\[-1pt]{\tiny propertyScan1D}};
  \node[draw, rounded corners=2pt, align=center, minimum width=3.1cm, minimum height=0.8cm, fill=black!4] at (3.6,-3.9) {\textbf{smoke\_13Propanediol}\\[-1pt]{\tiny propertyScan1D}};
  \node[draw, rounded corners=2pt, align=center, minimum width=3.1cm, minimum height=0.8cm, fill=black!4] at (7.2,-3.9) {\textbf{smoke\_13Butanediol}\\[-1pt]{\tiny propertyScan1D}};
\end{tikzpicture}
\end{tutscheme}
\tutwhat{The eight C287 components (new from ThermoML measurements): their Joback estimates reproduced, and Psat, liquid Cp and liquid density evaluated}
\begin{tutnarrative}
the eight components DEV.md 4c C287 added to the catalogue from measurements the ThermoML archive locates, loaded and exercised (written by hand beside bin/curate/c287\_new\_components.py).

1. REPRODUCTION. One `estimateComponent` per component with the Joback groups its record declares (groups.joback): the run prints the ideal-gas Cp polynomial, dHf\_298 and dGf\_298 the record carries as ESTIMATES (s\_298 = (dHf - dGf)/298.15 + the elements' CODATA standard entropies, provenance.s\_298). The run also prints Joback's Tc, Pc, Tb and omega; the records do NOT use them -- their Tc and Pc are measured, their omega comes from the fitted vapour pressure.

2. SMOKE TEST. One `propertyScan1D` per component, 300-340 K at 1 bar: the vapour pressure (the Antoine its vp case promoted), the liquid heat capacity (fitted polynomial) and the liquid density (the record's fitted Vliq through ThermoPackage::density, liquid branch). Every value must evaluate; outside a fitted range the engine announces the extrapolation, and every autoPromoted value and every warning-graded fit is voiced on the console and in the caveat block.

\end{tutnarrative}
\begin{tutgolden}
\goldrow{diag}{joback\_2Ethoxyethanol}{Cp298}{125.482286837}
\goldrow{diag}{joback\_2Ethoxyethanol}{Hvap\_kJmol}{43.583}
\goldrow{diag}{joback\_2Ethoxyethanol}{MW}{90.122}
\goldrow{diag}{joback\_2Ethoxyethanol}{Pc\_bar}{42.8866304789}
\goldrow{diag}{joback\_2Ethoxyethanol}{Psat\_298\_bar}{0.00346439746144}
\goldrow{diag}{joback\_2Ethoxyethanol}{Psat\_Tb\_bar}{0.96676543001}
\goldrow{diag}{joback\_2Ethoxyethanol}{Tb\_K}{405.72}
\goldrow{diag}{joback\_2Ethoxyethanol}{Tc\_K}{568.319976978}
\goldrow{diag}{joback\_2Ethoxyethanol}{Vc\_cm3mol}{296.5}
\goldrow{diag}{joback\_2Ethoxyethanol}{Vliq298\_cm3mol}{73.0555736041}
\goldrow{diag}{joback\_2Ethoxyethanol}{dGf\_kJmol}{-259.02}
\goldrow{diag}{joback\_2Ethoxyethanol}{dHf\_kJmol}{-410.34}
\goldrow{diag}{joback\_2Ethoxyethanol}{omega}{0.769281250814}
\goldrow{diag}{joback\_ethylOctanoate}{Cp298}{250.766466582}
\goldmore{114}
\end{tutgolden}

\clearpage
\tutentry{curate28\_vp\_methylDiEthanolAmine}{tutorials/props/curation/curate28\_vp\_methylDiEthanolAmine}
\begin{tutscheme}
\begin{tikzpicture}[>={Stealth[length=2mm]}, every node/.style={font=\scriptsize}]
  \node[draw, rounded corners=2pt, align=center, minimum width=3.1cm, minimum height=0.8cm, fill=black!4] at (0.0,-0.0) {\textbf{fit\_antoine}\\[-1pt]{\tiny vaporPressureFit}};
\end{tikzpicture}
\end{tutscheme}
\tutwhat{Antoine vapour pressure of methylDiEthanolAmine fitted by the curation pipeline to measured vapour pressures (ThermoML locates, the article sources), held out against an independent study}
\begin{tutnarrative}
the vapour pressure of methylDiEthanolAmine, REFITTED to measurements (DEV.md 4c C281 and C283). Written by `bin/choupo-curate vp methylDiEthanolAmine`; everything above the RESULT paragraph is the DECLARATION, committed before any fit ran, so it can be read against the result.

WHY. The catalogue record carried `vaporPressure \{ model AmbroseWalton; \}`, a corresponding-states ESTIMATE from Tc, Pc and omega that never reads the record's own Tb (C279 measured a median 19.6 \% AAD over the 25 such records with archive data). Here a three-parameter Antoine is fitted to EVERY usable laboratory's measurements of it.

STUDIES (one per article, extracted by `choupo-thermoml extract-vp`; ranked by points, ties by DOI): 10.1021/je060269j 9 pts 519.7-738.4 K VonNiederhausern; Wilson; Giles 10.1021/je400373d 7 pts 467.34-520.28 K Yang; Feng; Cheng; Zhang et al. 10.1021/je800290k 7 pts 409.69-435.5 K Kim; Svendsen; Borresen 10.1021/acs.jced.8b00054 6 pts 430.52-488.42 K Soames; Helal; Iglauer; Barifcani et al. SET ASIDE BY NAME: 10.1016/j.fluid.2010.08.026: j.fluid.2010.08.026.xml: 1 of 1 pressure rows of methylDiEthanolAmine declare no stated uncertainty -- a partial uncertainty is no uncertainty for a held-out band. Nothing was written.

FIT AND GRADED ACCEPTANCE (DEV.md 4c C285, Vitor 2026-11-10; the evidence, objective and outliers of C283; declared before the fit). EVIDENCE: every study above is fitted; none is held out (C283 amended: laboratories disagree with each other by more than they claim, so the test is agreement across all of them, not one laboratory's claimed precision). OBJECTIVE: `objective \{ method twoStage; \}` -- least squares in log10 P for a first estimate, then the LARGEST relative error in P, (P\_model - P)/P, over the kept points minimised from it (Vitor's practice, C280); both answers kept. OUTLIERS: after the least-squares stage a point with a relative P error above 15 \% (three times the acceptance) is DISCARDED BY NAME (its study DOI, T and P), listed in the result and the record, and the fit repeated without it, once (`outliers \{ maxRelativeError 15 percent; \}`). ACCEPTANCE, GRADED, the DEFAULT of C285: the maximum relative P error over all kept points is graded clean up to 5 \%, warning up to 10 \%, severe up to 20 \% and refused above; where the normal-boiling-point condition below applies, the absolute gap is graded clean up to 1 K, warning up to 2 K, severe up to 4 K and refused above. The fit's grade is the WORSE of the two. A clean, warning or severe fit is ACCEPTED and promoted (autoPromoted, humanReviewed false) with `fitGrade <grade>;` and the deciding numbers in its provenance entry, which the component loader voices on every run that uses it (a warning or a severe grade); a refused one is not promoted. This case declares no `system/curationDict`, so the default applies; a case that needs another acceptance declares it there, with its reason. The AmbroseWalton curve's errors on the same points are reported beside the fit's, as information.

NORMAL BOILING POINT: the record's Tb 520.15 K is judged against the range of the KEPT points (all studies span 409.69-738.4 K before any outlier is discarded): if Tb lies OUTSIDE that range, the fitted Antoine's own normal boiling point (P = 101.325 kPa) is GRADED by its absolute gap to it (the temperature limits of the acceptance below); inside it, the gap is reported and is not a condition.

EARLIER FITS, SUPERSEDED BY THIS DECLARATION @@EARLIER-BEGIN@@ C281 step 2, objective leastSquares (log10 P), the same partition and band: WHAT THE RUN FOUND (pinned in `expected`; measured, not argued). Antoine log10(Psat/bar) = 4.717858 - 2013.4423/(T + -93.3716), in-sample R2 0.999965 over 9 points, 519.7-738.4 K. HELD OUT: 7 points, AAD 5.1341 \% against the band 2.3 \% $\to$ notValidated. Its own normal boiling point (101.325 kPa): 520.66 K. NOT PROMOTED: held-out verdict notValidated (AAD 5.134 \% against the band 2.3 \% declared before the fit); the record keeps AmbroseWalton THEN C281 step 3, objective twoStage, one study fitted and one held out, judged by the band (superseded by C283): WHAT THE RUN FOUND (pinned in `expected`; measured, not argued). Objective twoStage. Antoine log10(Psat/bar) = 4.639497 - 1925.8766/(T + -105.5834), in-sample R2 0.999948 over 9 points, 519.7-738.4 K. HELD OUT: 7 points, AAD 2.2571 \% (max 3.6738 \%) against the band 2.3 \% $\to$ validated; the stage-1 least-squares answer on the same points: AAD 5.1341 \%, max 8.0690 \%. Its own normal boiling point (101.325 kPa): 521.20 K. PROMOTED (and since REVERTED by `choupo-curate revert`) on 2026-11-10 into data/standards/components/methylDiEthanolAmine.dat as reviewStatus autoPromoted, humanReviewed false, because: verdict validated (the record's Tb is inside the measured span or not declared; the fit's own normal boiling point is 521.20 K). THEN C283, every usable study fitted and none held out, accepted only at max relative error <= 5 \% and Tnb within 1 K (superseded by the graded acceptance of C285): WHAT THE RUN FOUND (pinned in `expected`; measured, not argued). Objective twoStage. Antoine log10(Psat/bar) = 5.728795 - 3050.8503/(T + 8.0790), in-sample R2 0.998997 over 28 kept points, 409.69-738.4 K (10.1021/je060269j 9/9 kept; 10.1021/je400373d 7/7 kept; 10.1021/je800290k 7/7 kept; 10.1021/acs.jced.8b00054 5/6 kept). Relative error in P over the kept points: mean 6.2530 \%, max 11.0828 \%; AmbroseWalton on the same points: mean 68.2614 \%, max 94.8247 \%. Outliers discarded: 10.1021/acs.jced.8b00054 T=430.52 K P=5000 Pa err=16.59 \%. Its own normal boiling point (101.325 kPa): 525.00 K. NOT PROMOTED by the DEFAULT acceptance of DEV.md 4c C283 (no system/curationDict): 4 studies, 28 points kept, 1 discarded as outliers; max relative error 11.08 \% (mean 6.253 \%) against 5 \%; AmbroseWalton on the same points mean 68.26 \% / max 94.82 \%; the record keeps AmbroseWalton @@EARLIER-END@@

NO MEASURED POINT IS IN THIS TREE: constant/archiveDatasets declares every series by DOI, NIST URL and sha256; `bin/choupo-curate fetch` regenerates them into constant/archiveExtracts/ (gitignored).

RESULT @@RESULT-BEGIN@@ WHAT THE RUN FOUND (pinned in `expected`; measured, not argued). Objective twoStage. Antoine log10(Psat/bar) = 5.728795 - 3050.8503/(T + 8.0790), in-sample R2 0.998997 over 28 kept points, 409.69-738.4 K (10.1021/je060269j 9/9 kept; 10.1021/je400373d 7/7 kept; 10.1021/je800290k 7/7 kept; 10.1021/acs.jced.8b00054 5/6 kept). Relative error in P over the kept points: mean 6.2530 \%, max 11.0828 \%; AmbroseWalton on the same points: mean 68.2614 \%, max 94.8247 \%. Outliers discarded: 10.1021/acs.jced.8b00054 T=430.52 K P=5000 Pa err=16.59 \%. Its own normal boiling point (101.325 kPa): 525.00 K. PROMOTED on 2026-11-10 into data/standards/components/methylDiEthanolAmine.dat as reviewStatus autoPromoted, humanReviewed false, because: fit grade severe under the DEFAULT graded acceptance of DEV.md 4c C285 (no system/curationDict), 4 studies, 28 points kept, 1 discarded as outliers; max relative error 11.08 \% (mean 6.253 \%), grade severe (limits 5 / 10 / 20 \%); Tb inside the kept points' range, normal boiling point 525.00 K, +4.85 K from the record's Tb 520.15 K (reported, not a condition); AmbroseWalton on the same points mean 68.26 \% / max 94.82 \% @@RESULT-END@@

\end{tutnarrative}
\begin{tutgolden}
\goldrow{diag}{fit\_antoine}{A}{5.72879504246}
\goldrow{diag}{fit\_antoine}{B}{3050.85029386}
\goldrow{diag}{fit\_antoine}{C}{8.07897284389}
\goldrow{diag}{fit\_antoine}{R2}{0.998996747803}
\goldrow{diag}{fit\_antoine}{dHvap\_kJ\_per\_mol}{57.1503624869}
\goldrow{diag}{fit\_antoine}{n\_heldout}{0}
\goldrow{diag}{fit\_antoine}{n\_points}{28}
\goldrow{verdict}{fit\_antoine}{verdict}{validationRefused}
\goldrow{diag}{fit\_antoine}{minimax\_evaluations}{737}
\goldrow{diag}{fit\_antoine}{minimax\_max\_scaled\_final}{0.110828035074}
\goldrow{diag}{fit\_antoine}{minimax\_max\_scaled\_start}{0.132455103124}
\goldrow{diag}{fit\_antoine}{objective\_method}{2}
\goldrow{diag}{fit\_antoine}{relErr\_max\_fit\_pct}{11.0828035074}
\goldrow{diag}{fit\_antoine}{relErr\_max\_fit\_pct\_leastSquares}{13.2455103124}
\goldmore{3}
\end{tutgolden}

\clearpage
\tutentry{curate29\_vp\_methylTertPentylEther}{tutorials/props/curation/curate29\_vp\_methylTertPentylEther}
\begin{tutscheme}
\begin{tikzpicture}[>={Stealth[length=2mm]}, every node/.style={font=\scriptsize}]
  \node[draw, rounded corners=2pt, align=center, minimum width=3.1cm, minimum height=0.8cm, fill=black!4] at (0.0,-0.0) {\textbf{fit\_antoine}\\[-1pt]{\tiny vaporPressureFit}};
\end{tikzpicture}
\end{tutscheme}
\tutwhat{Antoine vapour pressure of methylTertPentylEther fitted by the curation pipeline to measured vapour pressures (ThermoML locates, the article sources), held out against an independent study}
\begin{tutnarrative}
the vapour pressure of methylTertPentylEther, REFITTED to measurements (DEV.md 4c C281 and C283). Written by `bin/choupo-curate vp methylTertPentylEther`; everything above the RESULT paragraph is the DECLARATION, committed before any fit ran, so it can be read against the result.

WHY. The catalogue record carried `vaporPressure \{ model AmbroseWalton; \}`, a corresponding-states ESTIMATE from Tc, Pc and omega that never reads the record's own Tb (C279 measured a median 19.6 \% AAD over the 25 such records with archive data). Here a three-parameter Antoine is fitted to EVERY usable laboratory's measurements of it.

STUDIES (one per article, extracted by `choupo-thermoml extract-vp`; ranked by points, ties by DOI): 10.1021/je101250d 21 pts 313.44-359.95 K Mejia; Segura; Cartes 10.1021/je0255980 11 pts 274.2-307.8 K Verevkin; Krasnykh; Vasiltsova; Heintz 10.1021/je049867t 11 pts 298.02-355.64 K Vorenberg; Raal; Ramjugernath 10.1021/je901026b 5 pts 318.13-338.28 K Bernatova; Pavlicek; Wichterle SET ASIDE BY NAME: 10.1016/j.fluid.2005.06.017: 1 pressure point(s), fewer than 5 10.1016/j.fluid.2005.10.024: 2 pressure point(s), fewer than 5 10.1021/je200215d: 1 pressure point(s), fewer than 5

FIT AND GRADED ACCEPTANCE (DEV.md 4c C285, Vitor 2026-11-10; the evidence, objective and outliers of C283; declared before the fit). EVIDENCE: every study above is fitted; none is held out (C283 amended: laboratories disagree with each other by more than they claim, so the test is agreement across all of them, not one laboratory's claimed precision). OBJECTIVE: `objective \{ method twoStage; \}` -- least squares in log10 P for a first estimate, then the LARGEST relative error in P, (P\_model - P)/P, over the kept points minimised from it (Vitor's practice, C280); both answers kept. OUTLIERS: after the least-squares stage a point with a relative P error above 15 \% (three times the acceptance) is DISCARDED BY NAME (its study DOI, T and P), listed in the result and the record, and the fit repeated without it, once (`outliers \{ maxRelativeError 15 percent; \}`). ACCEPTANCE, GRADED, the DEFAULT of C285: the maximum relative P error over all kept points is graded clean up to 5 \%, warning up to 10 \%, severe up to 20 \% and refused above; where the normal-boiling-point condition below applies, the absolute gap is graded clean up to 1 K, warning up to 2 K, severe up to 4 K and refused above. The fit's grade is the WORSE of the two. A clean, warning or severe fit is ACCEPTED and promoted (autoPromoted, humanReviewed false) with `fitGrade <grade>;` and the deciding numbers in its provenance entry, which the component loader voices on every run that uses it (a warning or a severe grade); a refused one is not promoted. This case declares no `system/curationDict`, so the default applies; a case that needs another acceptance declares it there, with its reason. The AmbroseWalton curve's errors on the same points are reported beside the fit's, as information.

NORMAL BOILING POINT: the record's Tb 359.51 K is judged against the range of the KEPT points (all studies span 274.2-359.95 K before any outlier is discarded): if Tb lies OUTSIDE that range, the fitted Antoine's own normal boiling point (P = 101.325 kPa) is GRADED by its absolute gap to it (the temperature limits of the acceptance below); inside it, the gap is reported and is not a condition.

EARLIER FITS, SUPERSEDED BY THIS DECLARATION @@EARLIER-BEGIN@@ C281 step 2, objective leastSquares (log10 P), the same partition and band: WHAT THE RUN FOUND (pinned in `expected`; measured, not argued). Antoine log10(Psat/bar) = 4.240937 - 1366.2891/(T + -36.8588), in-sample R2 0.999996 over 21 points, 313.44-359.95 K. HELD OUT: 11 points, AAD 9.9414 \% against the band 7.8 \% $\to$ notValidated. Its own normal boiling point (101.325 kPa): 359.46 K. NOT PROMOTED: held-out verdict notValidated (AAD 9.941 \% against the band 7.8 \% declared before the fit); the record keeps AmbroseWalton THEN C281 step 3, objective twoStage, one study fitted and one held out, judged by the band (superseded by C283): WHAT THE RUN FOUND (pinned in `expected`; measured, not argued). Objective twoStage. Antoine log10(Psat/bar) = 4.153458 - 1314.0100/(T + -42.6484), in-sample R2 0.999991 over 21 points, 313.44-359.95 K. HELD OUT: 11 points, AAD 9.1421 \% (max 13.3936 \%) against the band 7.8 \% $\to$ notValidated; the stage-1 least-squares answer on the same points: AAD 9.9414 \%, max 14.6229 \%. Its own normal boiling point (101.325 kPa): 359.45 K. NOT PROMOTED: held-out verdict notValidated (AAD 9.142 \% against the band 7.8 \% declared before the fit); the record keeps AmbroseWalton. THEN C283, every usable study fitted and none held out, accepted only at max relative error <= 5 \% and Tnb within 1 K (superseded by the graded acceptance of C285): WHAT THE RUN FOUND (pinned in `expected`; measured, not argued). Objective twoStage. Antoine log10(Psat/bar) = 3.697092 - 1055.8861/(T + -74.0592), in-sample R2 0.999639 over 48 kept points, 274.2-359.95 K (10.1021/je101250d 21/21 kept; 10.1021/je0255980 11/11 kept; 10.1021/je049867t 11/11 kept; 10.1021/je901026b 5/5 kept). Relative error in P over the kept points: mean 1.8755 \%, max 3.9517 \%; AmbroseWalton on the same points: mean 3.0983 \%, max 14.0397 \%. Outliers discarded: none. Its own normal boiling point (101.325 kPa): 360.10 K. PROMOTED on 2026-11-10 into data/standards/components/methylTertPentylEther.dat as reviewStatus autoPromoted, humanReviewed false, because: under the DEFAULT acceptance of DEV.md 4c C283 (no system/curationDict), 4 studies, 48 points kept, 0 discarded as outliers; max relative error 3.952 \% (mean 1.875 \%) against 5 \%; Tb inside the kept points' range, normal boiling point 360.10 K, +0.59 K from the record's Tb 359.51 K (reported, not a condition); AmbroseWalton on the same points mean 3.098 \% / max 14.04 \% @@EARLIER-END@@

NO MEASURED POINT IS IN THIS TREE: constant/archiveDatasets declares every series by DOI, NIST URL and sha256; `bin/choupo-curate fetch` regenerates them into constant/archiveExtracts/ (gitignored).

RESULT @@RESULT-BEGIN@@ WHAT THE RUN FOUND (pinned in `expected`; measured, not argued). Objective twoStage. Antoine log10(Psat/bar) = 3.697092 - 1055.8861/(T + -74.0592), in-sample R2 0.999639 over 48 kept points, 274.2-359.95 K (10.1021/je101250d 21/21 kept; 10.1021/je0255980 11/11 kept; 10.1021/je049867t 11/11 kept; 10.1021/je901026b 5/5 kept). Relative error in P over the kept points: mean 1.8755 \%, max 3.9517 \%; AmbroseWalton on the same points: mean 3.0983 \%, max 14.0397 \%. Outliers discarded: none. Its own normal boiling point (101.325 kPa): 360.10 K. PROMOTED on 2026-11-10 into data/standards/components/methylTertPentylEther.dat as reviewStatus autoPromoted, humanReviewed false, because: fit grade clean under the DEFAULT graded acceptance of DEV.md 4c C285 (no system/curationDict), 4 studies, 48 points kept, 0 discarded as outliers; max relative error 3.952 \% (mean 1.875 \%), grade clean (limits 5 / 10 / 20 \%); Tb inside the kept points' range, normal boiling point 360.10 K, +0.59 K from the record's Tb 359.51 K (reported, not a condition); AmbroseWalton on the same points mean 3.098 \% / max 14.04 \% @@RESULT-END@@

\end{tutnarrative}
\begin{tutgolden}
\goldrow{diag}{fit\_antoine}{A}{3.69709166155}
\goldrow{diag}{fit\_antoine}{B}{1055.88609712}
\goldrow{diag}{fit\_antoine}{C}{-74.0591883095}
\goldrow{diag}{fit\_antoine}{R2}{0.999639172927}
\goldrow{diag}{fit\_antoine}{dHvap\_kJ\_per\_mol}{32.0443086146}
\goldrow{diag}{fit\_antoine}{n\_heldout}{0}
\goldrow{diag}{fit\_antoine}{n\_points}{48}
\goldrow{verdict}{fit\_antoine}{verdict}{validationRefused}
\goldrow{diag}{fit\_antoine}{minimax\_evaluations}{1076}
\goldrow{diag}{fit\_antoine}{minimax\_max\_scaled\_final}{0.0395171095955}
\goldrow{diag}{fit\_antoine}{minimax\_max\_scaled\_start}{0.0541698569491}
\goldrow{diag}{fit\_antoine}{objective\_method}{2}
\goldrow{diag}{fit\_antoine}{relErr\_max\_fit\_pct}{3.95171095955}
\goldrow{diag}{fit\_antoine}{relErr\_max\_fit\_pct\_leastSquares}{5.41698569491}
\goldmore{3}
\end{tutgolden}

\clearpage
\tutentry{curate30\_vp\_methylethanolamine}{tutorials/props/curation/curate30\_vp\_methylethanolamine}
\begin{tutscheme}
\begin{tikzpicture}[>={Stealth[length=2mm]}, every node/.style={font=\scriptsize}]
  \node[draw, rounded corners=2pt, align=center, minimum width=3.1cm, minimum height=0.8cm, fill=black!4] at (0.0,-0.0) {\textbf{fit\_antoine}\\[-1pt]{\tiny vaporPressureFit}};
\end{tikzpicture}
\end{tutscheme}
\tutwhat{Antoine vapour pressure of methylethanolamine fitted by the curation pipeline to measured vapour pressures (ThermoML locates, the article sources), held out against an independent study}
\begin{tutnarrative}
the vapour pressure of methylethanolamine, REFITTED to measurements (DEV.md 4c C281 and C283). Written by `bin/choupo-curate vp methylethanolamine`; everything above the RESULT paragraph is the DECLARATION, committed before any fit ran, so it can be read against the result.

WHY. The catalogue record carried `vaporPressure \{ model AmbroseWalton; \}`, a corresponding-states ESTIMATE from Tc, Pc and omega that never reads the record's own Tb (C279 measured a median 19.6 \% AAD over the 25 such records with archive data). Here a three-parameter Antoine is fitted to EVERY usable laboratory's measurements of it.

STUDIES (one per article, extracted by `choupo-thermoml extract-vp`; ranked by points, ties by DOI): 10.1021/je049761y 16 pts 274.9-320.2 K Kapteina; Slovik; Verevkin; Heintz 10.1021/je101259r 16 pts 341.8-451.1 K Klepacova; Huttenhuis; Derks; Versteeg 10.1021/je400373d 7 pts 388-431.2 K Yang; Feng; Cheng; Zhang et al. SET ASIDE BY NAME: (none)

FIT AND GRADED ACCEPTANCE (DEV.md 4c C285, Vitor 2026-11-10; the evidence, objective and outliers of C283; declared before the fit). EVIDENCE: every study above is fitted; none is held out (C283 amended: laboratories disagree with each other by more than they claim, so the test is agreement across all of them, not one laboratory's claimed precision). OBJECTIVE: `objective \{ method twoStage; \}` -- least squares in log10 P for a first estimate, then the LARGEST relative error in P, (P\_model - P)/P, over the kept points minimised from it (Vitor's practice, C280); both answers kept. OUTLIERS: after the least-squares stage a point with a relative P error above 15 \% (three times the acceptance) is DISCARDED BY NAME (its study DOI, T and P), listed in the result and the record, and the fit repeated without it, once (`outliers \{ maxRelativeError 15 percent; \}`). ACCEPTANCE, GRADED, the DEFAULT of C285: the maximum relative P error over all kept points is graded clean up to 5 \%, warning up to 10 \%, severe up to 20 \% and refused above; where the normal-boiling-point condition below applies, the absolute gap is graded clean up to 1 K, warning up to 2 K, severe up to 4 K and refused above. The fit's grade is the WORSE of the two. A clean, warning or severe fit is ACCEPTED and promoted (autoPromoted, humanReviewed false) with `fitGrade <grade>;` and the deciding numbers in its provenance entry, which the component loader voices on every run that uses it (a warning or a severe grade); a refused one is not promoted. This case declares no `system/curationDict`, so the default applies; a case that needs another acceptance declares it there, with its reason. The AmbroseWalton curve's errors on the same points are reported beside the fit's, as information.

NORMAL BOILING POINT: the record's Tb 432.388 K is judged against the range of the KEPT points (all studies span 274.9-451.1 K before any outlier is discarded): if Tb lies OUTSIDE that range, the fitted Antoine's own normal boiling point (P = 101.325 kPa) is GRADED by its absolute gap to it (the temperature limits of the acceptance below); inside it, the gap is reported and is not a condition.

EARLIER FITS, SUPERSEDED BY THIS DECLARATION @@EARLIER-BEGIN@@ C281 step 2, objective leastSquares (log10 P), the same partition and band: WHAT THE RUN FOUND (pinned in `expected`; measured, not argued). Antoine log10(Psat/bar) = 5.008624 - 1878.4613/(T + -62.8829), in-sample R2 0.999948 over 16 points, 274.9-320.2 K. HELD OUT: 16 points, AAD 15.8476 \% against the band 2.6 \% $\to$ notValidated. Its own normal boiling point (101.325 kPa): 438.36 K. NOT PROMOTED: held-out verdict notValidated (AAD 15.85 \% against the band 2.6 \% declared before the fit); the record keeps AmbroseWalton THEN C281 step 3, objective twoStage, one study fitted and one held out, judged by the band (superseded by C283): WHAT THE RUN FOUND (pinned in `expected`; measured, not argued). Objective twoStage. Antoine log10(Psat/bar) = 4.716842 - 1748.1461/(T + -70.8495), in-sample R2 0.999935 over 16 points, 274.9-320.2 K. HELD OUT: 16 points, AAD 21.7538 \% (max 26.8497 \%) against the band 2.6 \% $\to$ notValidated; the stage-1 least-squares answer on the same points: AAD 15.8476 \%, max 18.4784 \%. Its own normal boiling point (101.325 kPa): 441.92 K. NOT PROMOTED: held-out verdict notValidated (AAD 21.75 \% against the band 2.6 \% declared before the fit); the record keeps AmbroseWalton. THEN C283, every usable study fitted and none held out, accepted only at max relative error <= 5 \% and Tnb within 1 K (superseded by the graded acceptance of C285): WHAT THE RUN FOUND (pinned in `expected`; measured, not argued). Objective twoStage. Antoine log10(Psat/bar) = 5.412086 - 2028.3357/(T + -56.3502), in-sample R2 0.999931 over 39 kept points, 274.9-451.1 K (10.1021/je049761y 16/16 kept; 10.1021/je101259r 16/16 kept; 10.1021/je400373d 7/7 kept). Relative error in P over the kept points: mean 2.2087 \%, max 4.3596 \%; AmbroseWalton on the same points: mean 8.2726 \%, max 18.6719 \%. Outliers discarded: none. Its own normal boiling point (101.325 kPa): 431.53 K. PROMOTED on 2026-11-10 into data/standards/components/methylethanolamine.dat as reviewStatus autoPromoted, humanReviewed false, because: under the DEFAULT acceptance of DEV.md 4c C283 (no system/curationDict), 3 studies, 39 points kept, 0 discarded as outliers; max relative error 4.36 \% (mean 2.209 \%) against 5 \%; Tb inside the kept points' range, normal boiling point 431.53 K, -0.86 K from the record's Tb 432.388 K (reported, not a condition); AmbroseWalton on the same points mean 8.273 \% / max 18.67 \% @@EARLIER-END@@

NO MEASURED POINT IS IN THIS TREE: constant/archiveDatasets declares every series by DOI, NIST URL and sha256; `bin/choupo-curate fetch` regenerates them into constant/archiveExtracts/ (gitignored).

RESULT @@RESULT-BEGIN@@ WHAT THE RUN FOUND (pinned in `expected`; measured, not argued). Objective twoStage. Antoine log10(Psat/bar) = 5.412086 - 2028.3357/(T + -56.3502), in-sample R2 0.999931 over 39 kept points, 274.9-451.1 K (10.1021/je049761y 16/16 kept; 10.1021/je101259r 16/16 kept; 10.1021/je400373d 7/7 kept). Relative error in P over the kept points: mean 2.2087 \%, max 4.3596 \%; AmbroseWalton on the same points: mean 8.2726 \%, max 18.6719 \%. Outliers discarded: none. Its own normal boiling point (101.325 kPa): 431.53 K. PROMOTED on 2026-11-10 into data/standards/components/methylethanolamine.dat as reviewStatus autoPromoted, humanReviewed false, because: fit grade clean under the DEFAULT graded acceptance of DEV.md 4c C285 (no system/curationDict), 3 studies, 39 points kept, 0 discarded as outliers; max relative error 4.36 \% (mean 2.209 \%), grade clean (limits 5 / 10 / 20 \%); Tb inside the kept points' range, normal boiling point 431.53 K, -0.86 K from the record's Tb 432.388 K (reported, not a condition); AmbroseWalton on the same points mean 8.273 \% / max 18.67 \% @@RESULT-END@@

\end{tutnarrative}
\begin{tutgolden}
\goldrow{diag}{fit\_antoine}{A}{5.41208619177}
\goldrow{diag}{fit\_antoine}{B}{2028.33565299}
\goldrow{diag}{fit\_antoine}{C}{-56.3501865138}
\goldrow{diag}{fit\_antoine}{R2}{0.999930851341}
\goldrow{diag}{fit\_antoine}{dHvap\_kJ\_per\_mol}{50.7097518717}
\goldrow{diag}{fit\_antoine}{n\_heldout}{0}
\goldrow{diag}{fit\_antoine}{n\_points}{39}
\goldrow{verdict}{fit\_antoine}{verdict}{validationRefused}
\goldrow{diag}{fit\_antoine}{minimax\_evaluations}{496}
\goldrow{diag}{fit\_antoine}{minimax\_max\_scaled\_final}{0.0435955460638}
\goldrow{diag}{fit\_antoine}{minimax\_max\_scaled\_start}{0.0493325468889}
\goldrow{diag}{fit\_antoine}{objective\_method}{2}
\goldrow{diag}{fit\_antoine}{relErr\_max\_fit\_pct}{4.35955460638}
\goldrow{diag}{fit\_antoine}{relErr\_max\_fit\_pct\_leastSquares}{4.93325468889}
\goldmore{3}
\end{tutgolden}

\clearpage
\tutentry{curate31\_vp\_monoethanolamine}{tutorials/props/curation/curate31\_vp\_monoethanolamine}
\begin{tutscheme}
\begin{tikzpicture}[>={Stealth[length=2mm]}, every node/.style={font=\scriptsize}]
  \node[draw, rounded corners=2pt, align=center, minimum width=3.1cm, minimum height=0.8cm, fill=black!4] at (0.0,-0.0) {\textbf{fit\_antoine}\\[-1pt]{\tiny vaporPressureFit}};
\end{tikzpicture}
\end{tutscheme}
\tutwhat{Antoine vapour pressure of monoethanolamine fitted by the curation pipeline to measured vapour pressures (ThermoML locates, the article sources), held out against an independent study}
\begin{tutnarrative}
the vapour pressure of monoethanolamine, REFITTED to measurements (DEV.md 4c C281 and C283). Written by `bin/choupo-curate vp monoethanolamine`; everything above the RESULT paragraph is the DECLARATION, committed before any fit ran, so it can be read against the result.

WHY. The catalogue record carried `vaporPressure \{ model AmbroseWalton; \}`, a corresponding-states ESTIMATE from Tc, Pc and omega that never reads the record's own Tb (C279 measured a median 19.6 \% AAD over the 25 such records with archive data). Here a three-parameter Antoine is fitted to EVERY usable laboratory's measurements of it.

STUDIES (one per article, extracted by `choupo-thermoml extract-vp`; ranked by points, ties by DOI): 10.1016/j.fluid.2015.06.045 38 pts 273.65-303.16 K Ruzicka; Fulem; Mahnel; Cervinka 10.1021/je800530u 26 pts 283.15-363.15 K Belabbaci; Razzouk; Mokbel; Jose et al. 10.1021/je800290k 16 pts 357.46-435.61 K Kim; Svendsen; Borresen 10.1021/je049761y 13 pts 279-324.5 K Kapteina; Slovik; Verevkin; Heintz 10.1021/je101259r 7 pts 358.6-441.9 K Klepacova; Huttenhuis; Derks; Versteeg 10.1021/je400885z 6 pts 373.15-423.15 K McDougal; Jasperson; Wilson SET ASIDE BY NAME: 10.1021/je3002078: 1 pressure point(s), fewer than 5 10.1021/je7001336: 1 pressure point(s), fewer than 5

FIT AND GRADED ACCEPTANCE (DEV.md 4c C285, Vitor 2026-11-10; the evidence, objective and outliers of C283; declared before the fit). EVIDENCE: every study above is fitted; none is held out (C283 amended: laboratories disagree with each other by more than they claim, so the test is agreement across all of them, not one laboratory's claimed precision). OBJECTIVE: `objective \{ method twoStage; \}` -- least squares in log10 P for a first estimate, then the LARGEST relative error in P, (P\_model - P)/P, over the kept points minimised from it (Vitor's practice, C280); both answers kept. OUTLIERS: after the least-squares stage a point with a relative P error above 15 \% (three times the acceptance) is DISCARDED BY NAME (its study DOI, T and P), listed in the result and the record, and the fit repeated without it, once (`outliers \{ maxRelativeError 15 percent; \}`). ACCEPTANCE, GRADED, the DEFAULT of C285: the maximum relative P error over all kept points is graded clean up to 5 \%, warning up to 10 \%, severe up to 20 \% and refused above; where the normal-boiling-point condition below applies, the absolute gap is graded clean up to 1 K, warning up to 2 K, severe up to 4 K and refused above. The fit's grade is the WORSE of the two. A clean, warning or severe fit is ACCEPTED and promoted (autoPromoted, humanReviewed false) with `fitGrade <grade>;` and the deciding numbers in its provenance entry, which the component loader voices on every run that uses it (a warning or a severe grade); a refused one is not promoted. This case declares no `system/curationDict`, so the default applies; a case that needs another acceptance declares it there, with its reason. The AmbroseWalton curve's errors on the same points are reported beside the fit's, as information.

NORMAL BOILING POINT: the record's Tb 444.15 K is judged against the range of the KEPT points (all studies span 273.65-441.9 K before any outlier is discarded): if Tb lies OUTSIDE that range, the fitted Antoine's own normal boiling point (P = 101.325 kPa) is GRADED by its absolute gap to it (the temperature limits of the acceptance below); inside it, the gap is reported and is not a condition.

EARLIER FITS, SUPERSEDED BY THIS DECLARATION @@EARLIER-BEGIN@@ C281 step 2, objective leastSquares (log10 P), the same partition and band: WHAT THE RUN FOUND (pinned in `expected`; measured, not argued). Antoine log10(Psat/bar) = 6.766025 - 2807.2440/(T + -17.4236), in-sample R2 0.999993 over 38 points, 273.65-303.16 K. HELD OUT: 26 points, AAD 2.8176 \% against the band 5.4 \% $\to$ validated. Its own normal boiling point (101.325 kPa): 432.68 K. NOT PROMOTED: validated, but the fitted Antoine's normal boiling point 432.68 K misses the record's Tb 444.15 K by -11.47 K, outside the declared 2 K -- an extrapolation that contradicts the record; the record keeps AmbroseWalton THEN C281 step 3, objective twoStage, one study fitted and one held out, judged by the band (superseded by C283): WHAT THE RUN FOUND (pinned in `expected`; measured, not argued). Objective twoStage. Antoine log10(Psat/bar) = 6.301246 - 2556.1410/(T + -30.0521), in-sample R2 0.999986 over 38 points, 273.65-303.16 K. HELD OUT: 26 points, AAD 1.4219 \% (max 2.7692 \%) against the band 5.4 \% $\to$ validated; the stage-1 least-squares answer on the same points: AAD 2.8176 \%, max 7.3052 \%. Its own normal boiling point (101.325 kPa): 436.08 K. NOT PROMOTED: validated, but the fitted Antoine's normal boiling point 436.08 K misses the record's Tb 444.15 K by -8.07 K, outside the declared 2 K -- an extrapolation that contradicts the record; the record keeps AmbroseWalton. THEN C283, every usable study fitted and none held out, accepted only at max relative error <= 5 \% and Tnb within 1 K (superseded by the graded acceptance of C285): WHAT THE RUN FOUND (pinned in `expected`; measured, not argued). Objective twoStage. Antoine log10(Psat/bar) = 5.379683 - 2089.0887/(T + -55.2179), in-sample R2 0.999903 over 93 kept points, 273.65-441.9 K (10.1016/j.fluid.2015.06.045 38/38 kept; 10.1021/je800530u 26/26 kept; 10.1021/je800290k 16/16 kept; 10.1021/je049761y 0/13 kept; 10.1021/je101259r 7/7 kept; 10.1021/je400885z 6/6 kept). Relative error in P over the kept points: mean 2.8553 \%, max 4.5292 \%; AmbroseWalton on the same points: mean 6.3644 \%, max 15.8433 \%. Outliers discarded: 10.1021/je049761y T=279 K P=8.48 Pa err=23.33 \%; 10.1021/je049761y T=281.1 K P=10.55 Pa err=20.67 \%; 10.1021/je049761y T=294 K P=32.21 Pa err=23.25 \%; 10.1021/je049761y T=297.2 K P=40.9 Pa err=26.47 \%; 10.1021/je049761y T=300.2 K P=53.03 Pa err=24.29 \%; 10.1021/je049761y T=303.2 K P=63.3 Pa err=31.95 \%; 10.1021/je049761y T=306.3 K P=83.66 Pa err=26.82 \%; 10.1021/je049761y T=309.4 K P=105.6 Pa err=26.93 \%; 10.1021/je049761y T=312.4 K P=133.7 Pa err=25.07 \%; 10.1021/je049761y T=315.4 K P=168.9 Pa err=22.93 \%; 10.1021/je049761y T=318.4 K P=202.7 Pa err=26.61 \%; 10.1021/je049761y T=321.4 K P=255.9 Pa err=23.4 \%; 10.1021/je049761y T=324.5 K P=313.4 Pa err=24.27 \%. Its own normal boiling point (101.325 kPa): 443.96 K. PROMOTED on 2026-11-10 into data/standards/components/monoethanolamine.dat as reviewStatus autoPromoted, humanReviewed false, because: under the DEFAULT acceptance of DEV.md 4c C283 (no system/curationDict), 6 studies, 93 points kept, 13 discarded as outliers; max relative error 4.529 \% (mean 2.855 \%) against 5 \%; Tb outside the kept points' range, normal boiling point 443.96 K, -0.19 K from the record's Tb 444.15 K (within 1 K); AmbroseWalton on the same points mean 6.364 \% / max 15.84 \% @@EARLIER-END@@

NO MEASURED POINT IS IN THIS TREE: constant/archiveDatasets declares every series by DOI, NIST URL and sha256; `bin/choupo-curate fetch` regenerates them into constant/archiveExtracts/ (gitignored).

RESULT @@RESULT-BEGIN@@ WHAT THE RUN FOUND (pinned in `expected`; measured, not argued). Objective twoStage. Antoine log10(Psat/bar) = 5.379683 - 2089.0887/(T + -55.2179), in-sample R2 0.999903 over 93 kept points, 273.65-441.9 K (10.1016/j.fluid.2015.06.045 38/38 kept; 10.1021/je800530u 26/26 kept; 10.1021/je800290k 16/16 kept; 10.1021/je049761y 0/13 kept; 10.1021/je101259r 7/7 kept; 10.1021/je400885z 6/6 kept). Relative error in P over the kept points: mean 2.8553 \%, max 4.5292 \%; AmbroseWalton on the same points: mean 6.3644 \%, max 15.8433 \%. Outliers discarded: 10.1021/je049761y T=279 K P=8.48 Pa err=23.33 \%; 10.1021/je049761y T=281.1 K P=10.55 Pa err=20.67 \%; 10.1021/je049761y T=294 K P=32.21 Pa err=23.25 \%; 10.1021/je049761y T=297.2 K P=40.9 Pa err=26.47 \%; 10.1021/je049761y T=300.2 K P=53.03 Pa err=24.29 \%; 10.1021/je049761y T=303.2 K P=63.3 Pa err=31.95 \%; 10.1021/je049761y T=306.3 K P=83.66 Pa err=26.82 \%; 10.1021/je049761y T=309.4 K P=105.6 Pa err=26.93 \%; 10.1021/je049761y T=312.4 K P=133.7 Pa err=25.07 \%; 10.1021/je049761y T=315.4 K P=168.9 Pa err=22.93 \%; 10.1021/je049761y T=318.4 K P=202.7 Pa err=26.61 \%; 10.1021/je049761y T=321.4 K P=255.9 Pa err=23.4 \%; 10.1021/je049761y T=324.5 K P=313.4 Pa err=24.27 \%. Its own normal boiling point (101.325 kPa): 443.96 K. PROMOTED on 2026-11-10 into data/standards/components/monoethanolamine.dat as reviewStatus autoPromoted, humanReviewed false, because: fit grade clean under the DEFAULT graded acceptance of DEV.md 4c C285 (no system/curationDict), 6 studies, 93 points kept, 13 discarded as outliers; max relative error 4.529 \% (mean 2.855 \%), grade clean (limits 5 / 10 / 20 \%); Tb outside the kept points' range, normal boiling point 443.96 K, -0.19 K from the record's Tb 444.15 K, grade clean (limits 1 / 2 / 4 K); AmbroseWalton on the same points mean 6.364 \% / max 15.84 \% @@RESULT-END@@

\end{tutnarrative}
\begin{tutgolden}
\goldrow{diag}{fit\_antoine}{A}{5.37968312661}
\goldrow{diag}{fit\_antoine}{B}{2089.08865441}
\goldrow{diag}{fit\_antoine}{C}{-55.2179317189}
\goldrow{diag}{fit\_antoine}{R2}{0.999903075597}
\goldrow{diag}{fit\_antoine}{dHvap\_kJ\_per\_mol}{52.2332058277}
\goldrow{diag}{fit\_antoine}{n\_heldout}{0}
\goldrow{diag}{fit\_antoine}{n\_points}{93}
\goldrow{verdict}{fit\_antoine}{verdict}{validationRefused}
\goldrow{diag}{fit\_antoine}{minimax\_evaluations}{502}
\goldrow{diag}{fit\_antoine}{minimax\_max\_scaled\_final}{0.0452918335739}
\goldrow{diag}{fit\_antoine}{minimax\_max\_scaled\_start}{0.0581378330687}
\goldrow{diag}{fit\_antoine}{objective\_method}{2}
\goldrow{diag}{fit\_antoine}{relErr\_max\_fit\_pct}{4.52918335739}
\goldrow{diag}{fit\_antoine}{relErr\_max\_fit\_pct\_leastSquares}{5.81378330687}
\goldmore{3}
\end{tutgolden}

\clearpage
\tutentry{curate32\_vp\_nEicosane}{tutorials/props/curation/curate32\_vp\_nEicosane}
\begin{tutscheme}
\begin{tikzpicture}[>={Stealth[length=2mm]}, every node/.style={font=\scriptsize}]
  \node[draw, rounded corners=2pt, align=center, minimum width=3.1cm, minimum height=0.8cm, fill=black!4] at (0.0,-0.0) {\textbf{fit\_antoine}\\[-1pt]{\tiny vaporPressureFit}};
\end{tikzpicture}
\end{tutscheme}
\tutwhat{Antoine vapour pressure of nEicosane fitted by the curation pipeline to measured vapour pressures (ThermoML locates, the article sources), held out against an independent study}
\begin{tutnarrative}
the vapour pressure of nEicosane, REFITTED to measurements (DEV.md 4c C281 and C283). Written by `bin/choupo-curate vp nEicosane`; everything above the RESULT paragraph is the DECLARATION, committed before any fit ran, so it can be read against the result.

WHY. The catalogue record carried `vaporPressure \{ model AmbroseWalton; \}`, a corresponding-states ESTIMATE from Tc, Pc and omega that never reads the record's own Tb (C279 measured a median 19.6 \% AAD over the 25 such records with archive data). Here a three-parameter Antoine is fitted to EVERY usable laboratory's measurements of it.

STUDIES (one per article, extracted by `choupo-thermoml extract-vp`; ranked by points, ties by DOI): 10.1021/acs.jced.6b00821 8 pts 313.15-343.15 K Widegren; Beall; Tolbert; Lovestead et al. 10.1021/je800534x 6 pts 312.76-373.17 K Razzouk; Naccoul; Mokbel; Saab et al. SET ASIDE BY NAME: (none)

FIT AND GRADED ACCEPTANCE (DEV.md 4c C285, Vitor 2026-11-10; the evidence, objective and outliers of C283; declared before the fit). EVIDENCE: every study above is fitted; none is held out (C283 amended: laboratories disagree with each other by more than they claim, so the test is agreement across all of them, not one laboratory's claimed precision). OBJECTIVE: `objective \{ method twoStage; \}` -- least squares in log10 P for a first estimate, then the LARGEST relative error in P, (P\_model - P)/P, over the kept points minimised from it (Vitor's practice, C280); both answers kept. OUTLIERS: after the least-squares stage a point with a relative P error above 15 \% (three times the acceptance) is DISCARDED BY NAME (its study DOI, T and P), listed in the result and the record, and the fit repeated without it, once (`outliers \{ maxRelativeError 15 percent; \}`). ACCEPTANCE, GRADED, the DEFAULT of C285: the maximum relative P error over all kept points is graded clean up to 5 \%, warning up to 10 \%, severe up to 20 \% and refused above; where the normal-boiling-point condition below applies, the absolute gap is graded clean up to 1 K, warning up to 2 K, severe up to 4 K and refused above. The fit's grade is the WORSE of the two. A clean, warning or severe fit is ACCEPTED and promoted (autoPromoted, humanReviewed false) with `fitGrade <grade>;` and the deciding numbers in its provenance entry, which the component loader voices on every run that uses it (a warning or a severe grade); a refused one is not promoted. This case declares no `system/curationDict`, so the default applies; a case that needs another acceptance declares it there, with its reason. The AmbroseWalton curve's errors on the same points are reported beside the fit's, as information.

NORMAL BOILING POINT: the record's Tb 616.95 K is judged against the range of the KEPT points (all studies span 312.76-373.17 K before any outlier is discarded): if Tb lies OUTSIDE that range, the fitted Antoine's own normal boiling point (P = 101.325 kPa) is GRADED by its absolute gap to it (the temperature limits of the acceptance below); inside it, the gap is reported and is not a condition.

EARLIER FITS, SUPERSEDED BY THIS DECLARATION @@EARLIER-BEGIN@@ C281 step 2, objective leastSquares (log10 P), the same partition and band: WHAT THE RUN FOUND (pinned in `expected`; measured, not argued). Antoine log10(Psat/bar) = 3.607758 - 1978.6123/(T + -123.7264), in-sample R2 0.999592 over 8 points, 313.15-343.15 K. HELD OUT: 6 points, AAD 28.3788 \% against the band 4.2 \% $\to$ notValidated. Its own normal boiling point (101.325 kPa): 673.03 K. NOT PROMOTED: held-out verdict notValidated (AAD 28.38 \% against the band 4.2 \% declared before the fit); the record keeps AmbroseWalton THEN C281 step 3, objective twoStage, one study fitted and one held out, judged by the band (superseded by C283): WHAT THE RUN FOUND (pinned in `expected`; measured, not argued). Objective twoStage. Antoine log10(Psat/bar) = 4.280221 - 2250.4296/(T + -110.8773), in-sample R2 0.999299 over 8 points, 313.15-343.15 K. HELD OUT: 6 points, AAD 26.3951 \% (max 34.2660 \%) against the band 4.2 \% $\to$ notValidated; the stage-1 least-squares answer on the same points: AAD 28.3788 \%, max 36.2230 \%. Its own normal boiling point (101.325 kPa): 637.35 K. NOT PROMOTED: held-out verdict notValidated (AAD 26.4 \% against the band 4.2 \% declared before the fit); the record keeps AmbroseWalton. THEN C283, every usable study fitted and none held out, accepted only at max relative error <= 5 \% and Tnb within 1 K (superseded by the graded acceptance of C285): WHAT THE RUN FOUND (pinned in `expected`; measured, not argued). Objective twoStage. Antoine log10(Psat/bar) = 5.335640 - 2719.8070/(T + -89.8307), in-sample R2 0.999806 over 9 kept points, 313.15-373.17 K (10.1021/acs.jced.6b00821 8/8 kept; 10.1021/je800534x 1/6 kept). Relative error in P over the kept points: mean 2.1274 \%, max 2.8272 \%; AmbroseWalton on the same points: mean 2.9478 \%, max 6.8461 \%. Outliers discarded: 10.1021/je800534x T=312.76 K P=0.0113 Pa err=17.95 \%; 10.1021/je800534x T=322.82 K P=0.0345 Pa err=21.53 \%; 10.1021/je800534x T=327.82 K P=0.0601 Pa err=20.25 \%; 10.1021/je800534x T=337.82 K P=0.174 Pa err=17.81 \%; 10.1021/je800534x T=342.89 K P=0.293 Pa err=16.07 \%. Its own normal boiling point (101.325 kPa): 600.12 K. NOT PROMOTED by the DEFAULT acceptance of DEV.md 4c C283 (no system/curationDict): 2 studies, 9 points kept, 5 discarded as outliers; max relative error 2.827 \% (mean 2.127 \%) against 5 \% passes, but Tb lies outside the kept points' range and the normal boiling point 600.12 K, -16.83 K from the record's Tb 616.95 K is not within 1 K -- the gap is announced; AmbroseWalton on the same points mean 2.948 \% / max 6.846 \%; the record keeps AmbroseWalton @@EARLIER-END@@

NO MEASURED POINT IS IN THIS TREE: constant/archiveDatasets declares every series by DOI, NIST URL and sha256; `bin/choupo-curate fetch` regenerates them into constant/archiveExtracts/ (gitignored).

RESULT @@RESULT-BEGIN@@ WHAT THE RUN FOUND (pinned in `expected`; measured, not argued). Objective twoStage. Antoine log10(Psat/bar) = 5.335640 - 2719.8070/(T + -89.8307), in-sample R2 0.999806 over 9 kept points, 313.15-373.17 K (10.1021/acs.jced.6b00821 8/8 kept; 10.1021/je800534x 1/6 kept). Relative error in P over the kept points: mean 2.1274 \%, max 2.8272 \%; AmbroseWalton on the same points: mean 2.9478 \%, max 6.8461 \%. Outliers discarded: 10.1021/je800534x T=312.76 K P=0.0113 Pa err=17.95 \%; 10.1021/je800534x T=322.82 K P=0.0345 Pa err=21.53 \%; 10.1021/je800534x T=327.82 K P=0.0601 Pa err=20.25 \%; 10.1021/je800534x T=337.82 K P=0.174 Pa err=17.81 \%; 10.1021/je800534x T=342.89 K P=0.293 Pa err=16.07 \%. Its own normal boiling point (101.325 kPa): 600.12 K. NOT PROMOTED by the DEFAULT graded acceptance of DEV.md 4c C285 (no system/curationDict): fit grade refused -- 2 studies, 9 points kept, 5 discarded as outliers; max relative error 2.827 \% (mean 2.127 \%), grade clean (limits 5 / 10 / 20 \%); Tb outside the kept points' range, normal boiling point 600.12 K, -16.83 K from the record's Tb 616.95 K, grade refused (limits 1 / 2 / 4 K); AmbroseWalton on the same points mean 2.948 \% / max 6.846 \%; the record keeps AmbroseWalton @@RESULT-END@@

\end{tutnarrative}
\begin{tutgolden}
\goldrow{diag}{fit\_antoine}{A}{5.33564034105}
\goldrow{diag}{fit\_antoine}{B}{2719.80695797}
\goldrow{diag}{fit\_antoine}{C}{-89.8306524404}
\goldrow{diag}{fit\_antoine}{R2}{0.999805961711}
\goldrow{diag}{fit\_antoine}{dHvap\_kJ\_per\_mol}{90.3207676901}
\goldrow{diag}{fit\_antoine}{n\_heldout}{0}
\goldrow{diag}{fit\_antoine}{n\_points}{9}
\goldrow{verdict}{fit\_antoine}{verdict}{validationRefused}
\goldrow{diag}{fit\_antoine}{minimax\_evaluations}{352}
\goldrow{diag}{fit\_antoine}{minimax\_max\_scaled\_final}{0.0282722513095}
\goldrow{diag}{fit\_antoine}{minimax\_max\_scaled\_start}{0.0410720246362}
\goldrow{diag}{fit\_antoine}{objective\_method}{2}
\goldrow{diag}{fit\_antoine}{relErr\_max\_fit\_pct}{2.82722513095}
\goldrow{diag}{fit\_antoine}{relErr\_max\_fit\_pct\_leastSquares}{4.10720246362}
\goldmore{3}
\end{tutgolden}

\clearpage
\tutentry{curate33\_vp\_nHexylAcetate}{tutorials/props/curation/curate33\_vp\_nHexylAcetate}
\begin{tutscheme}
\begin{tikzpicture}[>={Stealth[length=2mm]}, every node/.style={font=\scriptsize}]
  \node[draw, rounded corners=2pt, align=center, minimum width=3.1cm, minimum height=0.8cm, fill=black!4] at (0.0,-0.0) {\textbf{fit\_antoine}\\[-1pt]{\tiny vaporPressureFit}};
\end{tikzpicture}
\end{tutscheme}
\tutwhat{Antoine vapour pressure of nHexylAcetate fitted by the curation pipeline to measured vapour pressures (ThermoML locates, the article sources), held out against an independent study}
\begin{tutnarrative}
the vapour pressure of nHexylAcetate, REFITTED to measurements (DEV.md 4c C281 and C283). Written by `bin/choupo-curate vp nHexylAcetate`; everything above the RESULT paragraph is the DECLARATION, committed before any fit ran, so it can be read against the result.

WHY. The catalogue record carried `vaporPressure \{ model AmbroseWalton; \}`, a corresponding-states ESTIMATE from Tc, Pc and omega that never reads the record's own Tb (C279 measured a median 19.6 \% AAD over the 25 such records with archive data). Here a three-parameter Antoine is fitted to EVERY usable laboratory's measurements of it.

STUDIES (one per article, extracted by `choupo-thermoml extract-vp`; ranked by points, ties by DOI): 10.1021/je050141m 14 pts 397.97-439.98 K Schmitt; Hasse 10.1016/j.jct.2005.08.003 13 pts 274.5-309.1 K Krasnykh; Verevkin; Koutek; Doubsky SET ASIDE BY NAME: (none)

FIT AND GRADED ACCEPTANCE (DEV.md 4c C285, Vitor 2026-11-10; the evidence, objective and outliers of C283; declared before the fit). EVIDENCE: every study above is fitted; none is held out (C283 amended: laboratories disagree with each other by more than they claim, so the test is agreement across all of them, not one laboratory's claimed precision). OBJECTIVE: `objective \{ method twoStage; \}` -- least squares in log10 P for a first estimate, then the LARGEST relative error in P, (P\_model - P)/P, over the kept points minimised from it (Vitor's practice, C280); both answers kept. OUTLIERS: after the least-squares stage a point with a relative P error above 15 \% (three times the acceptance) is DISCARDED BY NAME (its study DOI, T and P), listed in the result and the record, and the fit repeated without it, once (`outliers \{ maxRelativeError 15 percent; \}`). ACCEPTANCE, GRADED, the DEFAULT of C285: the maximum relative P error over all kept points is graded clean up to 5 \%, warning up to 10 \%, severe up to 20 \% and refused above; where the normal-boiling-point condition below applies, the absolute gap is graded clean up to 1 K, warning up to 2 K, severe up to 4 K and refused above. The fit's grade is the WORSE of the two. A clean, warning or severe fit is ACCEPTED and promoted (autoPromoted, humanReviewed false) with `fitGrade <grade>;` and the deciding numbers in its provenance entry, which the component loader voices on every run that uses it (a warning or a severe grade); a refused one is not promoted. This case declares no `system/curationDict`, so the default applies; a case that needs another acceptance declares it there, with its reason. The AmbroseWalton curve's errors on the same points are reported beside the fit's, as information.

NORMAL BOILING POINT: the record's Tb 444.65 K is judged against the range of the KEPT points (all studies span 274.5-439.98 K before any outlier is discarded): if Tb lies OUTSIDE that range, the fitted Antoine's own normal boiling point (P = 101.325 kPa) is GRADED by its absolute gap to it (the temperature limits of the acceptance below); inside it, the gap is reported and is not a condition.

EARLIER FITS, SUPERSEDED BY THIS DECLARATION @@EARLIER-BEGIN@@ C281 step 2, objective leastSquares (log10 P), the same partition and band: WHAT THE RUN FOUND (pinned in `expected`; measured, not argued). Antoine log10(Psat/bar) = 4.198146 - 1539.5515/(T + -77.2388), in-sample R2 0.999994 over 14 points, 397.97-439.98 K. HELD OUT: 13 points, AAD 5.8948 \% against the band 7.2 \% $\to$ validated. Its own normal boiling point (101.325 kPa): 444.46 K. PROMOTED (and since REVERTED by `choupo-curate revert`) on 2026-11-10 into data/standards/components/nHexylAcetate.dat as reviewStatus autoPromoted, humanReviewed false, because: verdict validated, normal boiling point 444.46 K within 2 K of the record's Tb 444.65 K THEN C281 step 3, objective twoStage, one study fitted and one held out, judged by the band (superseded by C283): WHAT THE RUN FOUND (pinned in `expected`; measured, not argued). Objective twoStage. Antoine log10(Psat/bar) = 4.228762 - 1562.4512/(T + -74.4887), in-sample R2 0.999992 over 14 points, 397.97-439.98 K. HELD OUT: 13 points, AAD 9.3708 \% (max 14.0038 \%) against the band 7.2 \% $\to$ notValidated; the stage-1 least-squares answer on the same points: AAD 5.8948 \%, max 10.2435 \%. Its own normal boiling point (101.325 kPa): 444.47 K. NOT PROMOTED: held-out verdict notValidated (AAD 9.371 \% against the band 7.2 \% declared before the fit); the record keeps AmbroseWalton. THEN C283, every usable study fitted and none held out, accepted only at max relative error <= 5 \% and Tnb within 1 K (superseded by the graded acceptance of C285): WHAT THE RUN FOUND (pinned in `expected`; measured, not argued). Objective twoStage. Antoine log10(Psat/bar) = 4.797190 - 1877.4971/(T + -51.1724), in-sample R2 0.999960 over 27 kept points, 274.5-439.98 K (10.1021/je050141m 14/14 kept; 10.1016/j.jct.2005.08.003 13/13 kept). Relative error in P over the kept points: mean 1.7082 \%, max 3.2973 \%; AmbroseWalton on the same points: mean 19.5636 \%, max 43.5862 \%. Outliers discarded: none. Its own normal boiling point (101.325 kPa): 443.01 K. NOT PROMOTED by the DEFAULT acceptance of DEV.md 4c C283 (no system/curationDict): 2 studies, 27 points kept, 0 discarded as outliers; max relative error 3.297 \% (mean 1.708 \%) against 5 \% passes, but Tb lies outside the kept points' range and the normal boiling point 443.01 K, -1.64 K from the record's Tb 444.65 K is not within 1 K -- the gap is announced; AmbroseWalton on the same points mean 19.56 \% / max 43.59 \%; the record keeps AmbroseWalton @@EARLIER-END@@

NO MEASURED POINT IS IN THIS TREE: constant/archiveDatasets declares every series by DOI, NIST URL and sha256; `bin/choupo-curate fetch` regenerates them into constant/archiveExtracts/ (gitignored).

RESULT @@RESULT-BEGIN@@ WHAT THE RUN FOUND (pinned in `expected`; measured, not argued). Objective twoStage. Antoine log10(Psat/bar) = 4.797190 - 1877.4971/(T + -51.1724), in-sample R2 0.999960 over 27 kept points, 274.5-439.98 K (10.1021/je050141m 14/14 kept; 10.1016/j.jct.2005.08.003 13/13 kept). Relative error in P over the kept points: mean 1.7082 \%, max 3.2973 \%; AmbroseWalton on the same points: mean 19.5636 \%, max 43.5862 \%. Outliers discarded: none. Its own normal boiling point (101.325 kPa): 443.01 K. PROMOTED on 2026-11-10 into data/standards/components/nHexylAcetate.dat as reviewStatus autoPromoted, humanReviewed false, because: fit grade warning under the DEFAULT graded acceptance of DEV.md 4c C285 (no system/curationDict), 2 studies, 27 points kept, 0 discarded as outliers; max relative error 3.297 \% (mean 1.708 \%), grade clean (limits 5 / 10 / 20 \%); Tb outside the kept points' range, normal boiling point 443.01 K, -1.64 K from the record's Tb 444.65 K, grade warning (limits 1 / 2 / 4 K); AmbroseWalton on the same points mean 19.56 \% / max 43.59 \% @@RESULT-END@@

\end{tutnarrative}
\begin{tutgolden}
\goldrow{diag}{fit\_antoine}{A}{4.79719030133}
\goldrow{diag}{fit\_antoine}{B}{1877.49712626}
\goldrow{diag}{fit\_antoine}{C}{-51.1723664762}
\goldrow{diag}{fit\_antoine}{R2}{0.999959904224}
\goldrow{diag}{fit\_antoine}{dHvap\_kJ\_per\_mol}{46.0283555971}
\goldrow{diag}{fit\_antoine}{n\_heldout}{0}
\goldrow{diag}{fit\_antoine}{n\_points}{27}
\goldrow{verdict}{fit\_antoine}{verdict}{validationRefused}
\goldrow{diag}{fit\_antoine}{minimax\_evaluations}{436}
\goldrow{diag}{fit\_antoine}{minimax\_max\_scaled\_final}{0.0329732755178}
\goldrow{diag}{fit\_antoine}{minimax\_max\_scaled\_start}{0.0439792127434}
\goldrow{diag}{fit\_antoine}{objective\_method}{2}
\goldrow{diag}{fit\_antoine}{relErr\_max\_fit\_pct}{3.29732755178}
\goldrow{diag}{fit\_antoine}{relErr\_max\_fit\_pct\_leastSquares}{4.39792127434}
\goldmore{3}
\end{tutgolden}

\clearpage
\tutentry{curate34\_vp\_diethanolamine}{tutorials/props/curation/curate34\_vp\_diethanolamine}
\begin{tutscheme}
\begin{tikzpicture}[>={Stealth[length=2mm]}, every node/.style={font=\scriptsize}]
  \node[draw, rounded corners=2pt, align=center, minimum width=3.1cm, minimum height=0.8cm, fill=black!4] at (0.0,-0.0) {\textbf{fit\_antoine}\\[-1pt]{\tiny vaporPressureFit}};
\end{tikzpicture}
\end{tutscheme}
\tutwhat{Antoine vapour pressure of diethanolamine fitted by the curation pipeline to the measured vapour pressures of every usable laboratory (ThermoML locates, the articles source), outliers discarded by name}
\begin{tutnarrative}
the vapour pressure of diethanolamine, REFITTED to measurements (DEV.md 4c C281 and C283). Written by `bin/choupo-curate vp diethanolamine`; everything above the RESULT paragraph is the DECLARATION, committed before any fit ran, so it can be read against the result.

WHY. The catalogue record carried `vaporPressure \{ model AmbroseWalton; \}`, a corresponding-states ESTIMATE from Tc, Pc and omega that never reads the record's own Tb (C279 measured a median 19.6 \% AAD over the 25 such records with archive data). Here a three-parameter Antoine is fitted to EVERY usable laboratory's measurements of it.

STUDIES (one per article, extracted by `choupo-thermoml extract-vp`; ranked by points, ties by DOI): 10.1021/je101259r 5 pts 427.5-452.4 K Klepacova; Huttenhuis; Derks; Versteeg SET ASIDE BY NAME: 10.1016/j.jct.2018.09.002: 1 pressure point(s), fewer than 5

FIT AND GRADED ACCEPTANCE (DEV.md 4c C285, Vitor 2026-11-10; the evidence, objective and outliers of C283; declared before the fit). EVIDENCE: every study above is fitted; none is held out (C283 amended: laboratories disagree with each other by more than they claim, so the test is agreement across all of them, not one laboratory's claimed precision). OBJECTIVE: `objective \{ method twoStage; \}` -- least squares in log10 P for a first estimate, then the LARGEST relative error in P, (P\_model - P)/P, over the kept points minimised from it (Vitor's practice, C280); both answers kept. OUTLIERS: after the least-squares stage a point with a relative P error above 15 \% (three times the acceptance) is DISCARDED BY NAME (its study DOI, T and P), listed in the result and the record, and the fit repeated without it, once (`outliers \{ maxRelativeError 15 percent; \}`). ACCEPTANCE, GRADED, the DEFAULT of C285: the maximum relative P error over all kept points is graded clean up to 5 \%, warning up to 10 \%, severe up to 20 \% and refused above; where the normal-boiling-point condition below applies, the absolute gap is graded clean up to 1 K, warning up to 2 K, severe up to 4 K and refused above. The fit's grade is the WORSE of the two. A clean, warning or severe fit is ACCEPTED and promoted (autoPromoted, humanReviewed false) with `fitGrade <grade>;` and the deciding numbers in its provenance entry, which the component loader voices on every run that uses it (a warning or a severe grade); a refused one is not promoted. This case declares no `system/curationDict`, so the default applies; a case that needs another acceptance declares it there, with its reason. The AmbroseWalton curve's errors on the same points are reported beside the fit's, as information.

NORMAL BOILING POINT: the record's Tb 542 K is judged against the range of the KEPT points (all studies span 427.5-452.4 K before any outlier is discarded): if Tb lies OUTSIDE that range, the fitted Antoine's own normal boiling point (P = 101.325 kPa) is GRADED by its absolute gap to it (the temperature limits of the acceptance below); inside it, the gap is reported and is not a condition.

EARLIER FITS, SUPERSEDED BY THIS DECLARATION @@EARLIER-BEGIN@@ C283, every usable study fitted and none held out, accepted only at max relative error <= 5 \% and Tnb within 1 K (superseded by the graded acceptance of C285): WHAT THE RUN FOUND (pinned in `expected`; measured, not argued). Objective twoStage. Antoine log10(Psat/bar) = 4.246563 - 1576.6883/(T + -168.4642), in-sample R2 0.999956 over 5 kept points, 427.5-452.4 K (10.1021/je101259r 5/5 kept). Relative error in P over the kept points: mean 0.2724 \%, max 0.3326 \%; AmbroseWalton on the same points: mean 15.9429 \%, max 16.5744 \%. Outliers discarded: none. Its own normal boiling point (101.325 kPa): 540.25 K. NOT PROMOTED by the DEFAULT acceptance of DEV.md 4c C283 (no system/curationDict): 1 studies, 5 points kept, 0 discarded as outliers; max relative error 0.3326 \% (mean 0.2724 \%) against 5 \% passes, but Tb lies outside the kept points' range and the normal boiling point 540.25 K, -1.75 K from the record's Tb 542 K is not within 1 K -- the gap is announced; AmbroseWalton on the same points mean 15.94 \% / max 16.57 \%; the record keeps AmbroseWalton @@EARLIER-END@@

NO MEASURED POINT IS IN THIS TREE: constant/archiveDatasets declares every series by DOI, NIST URL and sha256; `bin/choupo-curate fetch` regenerates them into constant/archiveExtracts/ (gitignored).

RESULT @@RESULT-BEGIN@@ WHAT THE RUN FOUND (pinned in `expected`; measured, not argued). Objective twoStage. Antoine log10(Psat/bar) = 4.246563 - 1576.6883/(T + -168.4642), in-sample R2 0.999956 over 5 kept points, 427.5-452.4 K (10.1021/je101259r 5/5 kept). Relative error in P over the kept points: mean 0.2724 \%, max 0.3326 \%; AmbroseWalton on the same points: mean 15.9429 \%, max 16.5744 \%. Outliers discarded: none. Its own normal boiling point (101.325 kPa): 540.25 K. PROMOTED on 2026-11-10 into data/standards/components/diethanolamine.dat as reviewStatus autoPromoted, humanReviewed false, because: fit grade warning under the DEFAULT graded acceptance of DEV.md 4c C285 (no system/curationDict), 1 studies, 5 points kept, 0 discarded as outliers; max relative error 0.3326 \% (mean 0.2724 \%), grade clean (limits 5 / 10 / 20 \%); Tb outside the kept points' range, normal boiling point 540.25 K, -1.75 K from the record's Tb 542 K, grade warning (limits 1 / 2 / 4 K); AmbroseWalton on the same points mean 15.94 \% / max 16.57 \% @@RESULT-END@@

\end{tutnarrative}
\begin{tutgolden}
\goldrow{diag}{fit\_antoine}{A}{4.24656334076}
\goldrow{diag}{fit\_antoine}{B}{1576.68827443}
\goldrow{diag}{fit\_antoine}{C}{-168.464233131}
\goldrow{diag}{fit\_antoine}{R2}{0.99995640559}
\goldrow{diag}{fit\_antoine}{dHvap\_kJ\_per\_mol}{76.6303303196}
\goldrow{diag}{fit\_antoine}{minimax\_evaluations}{852}
\goldrow{diag}{fit\_antoine}{minimax\_max\_scaled\_final}{0.00332559355479}
\goldrow{diag}{fit\_antoine}{minimax\_max\_scaled\_start}{0.00419305125669}
\goldrow{diag}{fit\_antoine}{n\_heldout}{0}
\goldrow{diag}{fit\_antoine}{n\_outliers\_dropped}{0}
\goldrow{diag}{fit\_antoine}{n\_points}{5}
\goldrow{diag}{fit\_antoine}{objective\_method}{2}
\goldrow{diag}{fit\_antoine}{relErr\_max\_fit\_pct}{0.332559355479}
\goldrow{diag}{fit\_antoine}{relErr\_max\_fit\_pct\_leastSquares}{0.419305125669}
\goldmore{3}
\end{tutgolden}

\clearpage
\tutentry{curate35\_vp\_mNitrotoluene}{tutorials/props/curation/curate35\_vp\_mNitrotoluene}
\begin{tutscheme}
\begin{tikzpicture}[>={Stealth[length=2mm]}, every node/.style={font=\scriptsize}]
  \node[draw, rounded corners=2pt, align=center, minimum width=3.1cm, minimum height=0.8cm, fill=black!4] at (0.0,-0.0) {\textbf{fit\_antoine}\\[-1pt]{\tiny vaporPressureFit}};
\end{tikzpicture}
\end{tutscheme}
\tutwhat{Antoine vapour pressure of mNitrotoluene fitted by the curation pipeline to the measured vapour pressures of every usable laboratory (ThermoML locates, the articles source), outliers discarded by name}
\begin{tutnarrative}
the vapour pressure of mNitrotoluene, REFITTED to measurements (DEV.md 4c C281 and C283). Written by `bin/choupo-curate vp mNitrotoluene`; everything above the RESULT paragraph is the DECLARATION, committed before any fit ran, so it can be read against the result.

WHY. The catalogue record carried `vaporPressure \{ model AmbroseWalton; \}`, a corresponding-states ESTIMATE from Tc, Pc and omega that never reads the record's own Tb (C279 measured a median 19.6 \% AAD over the 25 such records with archive data). Here a three-parameter Antoine is fitted to EVERY usable laboratory's measurements of it.

STUDIES (one per article, extracted by `choupo-thermoml extract-vp`; ranked by points, ties by DOI): 10.1016/j.jct.2017.03.029 34 pts 290.6-338 K Bikelyte; Hartel; Stierstorfer; Klapotke et al. SET ASIDE BY NAME: 10.1021/je900293j: 3 pressure point(s), fewer than 5

FIT AND GRADED ACCEPTANCE (DEV.md 4c C285, Vitor 2026-11-10; the evidence, objective and outliers of C283; declared before the fit). EVIDENCE: every study above is fitted; none is held out (C283 amended: laboratories disagree with each other by more than they claim, so the test is agreement across all of them, not one laboratory's claimed precision). OBJECTIVE: `objective \{ method twoStage; \}` -- least squares in log10 P for a first estimate, then the LARGEST relative error in P, (P\_model - P)/P, over the kept points minimised from it (Vitor's practice, C280); both answers kept. OUTLIERS: after the least-squares stage a point with a relative P error above 15 \% (three times the acceptance) is DISCARDED BY NAME (its study DOI, T and P), listed in the result and the record, and the fit repeated without it, once (`outliers \{ maxRelativeError 15 percent; \}`). ACCEPTANCE, GRADED, the DEFAULT of C285: the maximum relative P error over all kept points is graded clean up to 5 \%, warning up to 10 \%, severe up to 20 \% and refused above; where the normal-boiling-point condition below applies, the absolute gap is graded clean up to 1 K, warning up to 2 K, severe up to 4 K and refused above. The fit's grade is the WORSE of the two. A clean, warning or severe fit is ACCEPTED and promoted (autoPromoted, humanReviewed false) with `fitGrade <grade>;` and the deciding numbers in its provenance entry, which the component loader voices on every run that uses it (a warning or a severe grade); a refused one is not promoted. This case declares no `system/curationDict`, so the default applies; a case that needs another acceptance declares it there, with its reason. The AmbroseWalton curve's errors on the same points are reported beside the fit's, as information.

NORMAL BOILING POINT: the record's Tb 504 K is judged against the range of the KEPT points (all studies span 290.6-338 K before any outlier is discarded): if Tb lies OUTSIDE that range, the fitted Antoine's own normal boiling point (P = 101.325 kPa) is GRADED by its absolute gap to it (the temperature limits of the acceptance below); inside it, the gap is reported and is not a condition.

EARLIER FITS, SUPERSEDED BY THIS DECLARATION @@EARLIER-BEGIN@@ C283, every usable study fitted and none held out, accepted only at max relative error <= 5 \% and Tnb within 1 K (superseded by the graded acceptance of C285): WHAT THE RUN FOUND (pinned in `expected`; measured, not argued). Objective twoStage. Antoine log10(Psat/bar) = 6.592207 - 3238.6044/(T + 9.9450), in-sample R2 0.999449 over 34 kept points, 290.6-338 K (10.1016/j.jct.2017.03.029 34/34 kept). Relative error in P over the kept points: mean 1.6175 \%, max 2.9342 \%; AmbroseWalton on the same points: mean 48.1334 \%, max 57.8067 \%. Outliers discarded: none. Its own normal boiling point (101.325 kPa): 481.76 K. NOT PROMOTED by the DEFAULT acceptance of DEV.md 4c C283 (no system/curationDict): 1 studies, 34 points kept, 0 discarded as outliers; max relative error 2.934 \% (mean 1.617 \%) against 5 \% passes, but Tb lies outside the kept points' range and the normal boiling point 481.76 K, -22.24 K from the record's Tb 504 K is not within 1 K -- the gap is announced; AmbroseWalton on the same points mean 48.13 \% / max 57.81 \%; the record keeps AmbroseWalton @@EARLIER-END@@

NO MEASURED POINT IS IN THIS TREE: constant/archiveDatasets declares every series by DOI, NIST URL and sha256; `bin/choupo-curate fetch` regenerates them into constant/archiveExtracts/ (gitignored).

RESULT @@RESULT-BEGIN@@ WHAT THE RUN FOUND (pinned in `expected`; measured, not argued). Objective twoStage. Antoine log10(Psat/bar) = 6.592207 - 3238.6044/(T + 9.9450), in-sample R2 0.999449 over 34 kept points, 290.6-338 K (10.1016/j.jct.2017.03.029 34/34 kept). Relative error in P over the kept points: mean 1.6175 \%, max 2.9342 \%; AmbroseWalton on the same points: mean 48.1334 \%, max 57.8067 \%. Outliers discarded: none. Its own normal boiling point (101.325 kPa): 481.76 K. NOT PROMOTED by the DEFAULT graded acceptance of DEV.md 4c C285 (no system/curationDict): fit grade refused -- 1 studies, 34 points kept, 0 discarded as outliers; max relative error 2.934 \% (mean 1.617 \%), grade clean (limits 5 / 10 / 20 \%); Tb outside the kept points' range, normal boiling point 481.76 K, -22.24 K from the record's Tb 504 K, grade refused (limits 1 / 2 / 4 K); AmbroseWalton on the same points mean 48.13 \% / max 57.81 \%; the record keeps AmbroseWalton @@RESULT-END@@

\end{tutnarrative}
\begin{tutgolden}
\goldrow{diag}{fit\_antoine}{A}{6.59220680973}
\goldrow{diag}{fit\_antoine}{B}{3238.60435642}
\goldrow{diag}{fit\_antoine}{C}{9.94498925265}
\goldrow{diag}{fit\_antoine}{R2}{0.999449002049}
\goldrow{diag}{fit\_antoine}{dHvap\_kJ\_per\_mol}{58.5086391272}
\goldrow{diag}{fit\_antoine}{minimax\_evaluations}{494}
\goldrow{diag}{fit\_antoine}{minimax\_max\_scaled\_final}{0.0293415504217}
\goldrow{diag}{fit\_antoine}{minimax\_max\_scaled\_start}{0.0346746185625}
\goldrow{diag}{fit\_antoine}{n\_heldout}{0}
\goldrow{diag}{fit\_antoine}{n\_outliers\_dropped}{0}
\goldrow{diag}{fit\_antoine}{n\_points}{34}
\goldrow{diag}{fit\_antoine}{objective\_method}{2}
\goldrow{diag}{fit\_antoine}{relErr\_max\_fit\_pct}{2.93415504217}
\goldrow{diag}{fit\_antoine}{relErr\_max\_fit\_pct\_leastSquares}{3.46746185625}
\goldmore{3}
\end{tutgolden}

\clearpage
\tutentry{curate36\_vp\_nAminoethylEthanolamine}{tutorials/props/curation/curate36\_vp\_nAminoethylEthanolamine}
\begin{tutscheme}
\begin{tikzpicture}[>={Stealth[length=2mm]}, every node/.style={font=\scriptsize}]
  \node[draw, rounded corners=2pt, align=center, minimum width=3.1cm, minimum height=0.8cm, fill=black!4] at (0.0,-0.0) {\textbf{fit\_antoine}\\[-1pt]{\tiny vaporPressureFit}};
\end{tikzpicture}
\end{tutscheme}
\tutwhat{Antoine vapour pressure of nAminoethylEthanolamine fitted by the curation pipeline to the measured vapour pressures of every usable laboratory (ThermoML locates, the articles source), outliers discarded by name}
\begin{tutnarrative}
the vapour pressure of nAminoethylEthanolamine, REFITTED to measurements (DEV.md 4c C281 and C283). Written by `bin/choupo-curate vp nAminoethylEthanolamine`; everything above the RESULT paragraph is the DECLARATION, committed before any fit ran, so it can be read against the result.

WHY. The catalogue record carried `vaporPressure \{ model AmbroseWalton; \}`, a corresponding-states ESTIMATE from Tc, Pc and omega that never reads the record's own Tb (C279 measured a median 19.6 \% AAD over the 25 such records with archive data). Here a three-parameter Antoine is fitted to EVERY usable laboratory's measurements of it.

STUDIES (one per article, extracted by `choupo-thermoml extract-vp`; ranked by points, ties by DOI): 10.1021/je2002489 33 pts 323.2-374.2 K Verevkin; Tong; Welz-Biermann; Chernyak SET ASIDE BY NAME: (none)

FIT AND GRADED ACCEPTANCE (DEV.md 4c C285, Vitor 2026-11-10; the evidence, objective and outliers of C283; declared before the fit). EVIDENCE: every study above is fitted; none is held out (C283 amended: laboratories disagree with each other by more than they claim, so the test is agreement across all of them, not one laboratory's claimed precision). OBJECTIVE: `objective \{ method twoStage; \}` -- least squares in log10 P for a first estimate, then the LARGEST relative error in P, (P\_model - P)/P, over the kept points minimised from it (Vitor's practice, C280); both answers kept. OUTLIERS: after the least-squares stage a point with a relative P error above 15 \% (three times the acceptance) is DISCARDED BY NAME (its study DOI, T and P), listed in the result and the record, and the fit repeated without it, once (`outliers \{ maxRelativeError 15 percent; \}`). ACCEPTANCE, GRADED, the DEFAULT of C285: the maximum relative P error over all kept points is graded clean up to 5 \%, warning up to 10 \%, severe up to 20 \% and refused above; where the normal-boiling-point condition below applies, the absolute gap is graded clean up to 1 K, warning up to 2 K, severe up to 4 K and refused above. The fit's grade is the WORSE of the two. A clean, warning or severe fit is ACCEPTED and promoted (autoPromoted, humanReviewed false) with `fitGrade <grade>;` and the deciding numbers in its provenance entry, which the component loader voices on every run that uses it (a warning or a severe grade); a refused one is not promoted. This case declares no `system/curationDict`, so the default applies; a case that needs another acceptance declares it there, with its reason. The AmbroseWalton curve's errors on the same points are reported beside the fit's, as information.

NORMAL BOILING POINT: the record's Tb 517 K is judged against the range of the KEPT points (all studies span 323.2-374.2 K before any outlier is discarded): if Tb lies OUTSIDE that range, the fitted Antoine's own normal boiling point (P = 101.325 kPa) is GRADED by its absolute gap to it (the temperature limits of the acceptance below); inside it, the gap is reported and is not a condition.

EARLIER FITS, SUPERSEDED BY THIS DECLARATION @@EARLIER-BEGIN@@ C283, every usable study fitted and none held out, accepted only at max relative error <= 5 \% and Tnb within 1 K (superseded by the graded acceptance of C285): WHAT THE RUN FOUND (pinned in `expected`; measured, not argued). Objective twoStage. Antoine log10(Psat/bar) = 6.274747 - 3062.5355/(T + -30.3642), in-sample R2 0.999411 over 33 kept points, 323.2-374.2 K (10.1021/je2002489 33/33 kept). Relative error in P over the kept points: mean 2.2635 \%, max 4.4704 \%; AmbroseWalton on the same points: mean 68.6798 \%, max 79.7274 \%. Outliers discarded: none. Its own normal boiling point (101.325 kPa): 518.88 K. NOT PROMOTED by the DEFAULT acceptance of DEV.md 4c C283 (no system/curationDict): 1 studies, 33 points kept, 0 discarded as outliers; max relative error 4.47 \% (mean 2.264 \%) against 5 \% passes, but Tb lies outside the kept points' range and the normal boiling point 518.88 K, +1.88 K from the record's Tb 517 K is not within 1 K -- the gap is announced; AmbroseWalton on the same points mean 68.68 \% / max 79.73 \%; the record keeps AmbroseWalton @@EARLIER-END@@

NO MEASURED POINT IS IN THIS TREE: constant/archiveDatasets declares every series by DOI, NIST URL and sha256; `bin/choupo-curate fetch` regenerates them into constant/archiveExtracts/ (gitignored).

RESULT @@RESULT-BEGIN@@ WHAT THE RUN FOUND (pinned in `expected`; measured, not argued). Objective twoStage. Antoine log10(Psat/bar) = 6.274747 - 3062.5355/(T + -30.3642), in-sample R2 0.999411 over 33 kept points, 323.2-374.2 K (10.1021/je2002489 33/33 kept). Relative error in P over the kept points: mean 2.2635 \%, max 4.4704 \%; AmbroseWalton on the same points: mean 68.6798 \%, max 79.7274 \%. Outliers discarded: none. Its own normal boiling point (101.325 kPa): 518.88 K. PROMOTED on 2026-11-10 into data/standards/components/nAminoethylEthanolamine.dat as reviewStatus autoPromoted, humanReviewed false, because: fit grade warning under the DEFAULT graded acceptance of DEV.md 4c C285 (no system/curationDict), 1 studies, 33 points kept, 0 discarded as outliers; max relative error 4.47 \% (mean 2.264 \%), grade clean (limits 5 / 10 / 20 \%); Tb outside the kept points' range, normal boiling point 518.88 K, +1.88 K from the record's Tb 517 K, grade warning (limits 1 / 2 / 4 K); AmbroseWalton on the same points mean 68.68 \% / max 79.73 \% @@RESULT-END@@

\end{tutnarrative}
\begin{tutgolden}
\goldrow{diag}{fit\_antoine}{A}{6.27474745869}
\goldrow{diag}{fit\_antoine}{B}{3062.53548288}
\goldrow{diag}{fit\_antoine}{C}{-30.3641631241}
\goldrow{diag}{fit\_antoine}{R2}{0.999410712628}
\goldrow{diag}{fit\_antoine}{dHvap\_kJ\_per\_mol}{69.4442496492}
\goldrow{diag}{fit\_antoine}{minimax\_evaluations}{459}
\goldrow{diag}{fit\_antoine}{minimax\_max\_scaled\_final}{0.0447039504695}
\goldrow{diag}{fit\_antoine}{minimax\_max\_scaled\_start}{0.0490640428544}
\goldrow{diag}{fit\_antoine}{n\_heldout}{0}
\goldrow{diag}{fit\_antoine}{n\_outliers\_dropped}{0}
\goldrow{diag}{fit\_antoine}{n\_points}{33}
\goldrow{diag}{fit\_antoine}{objective\_method}{2}
\goldrow{diag}{fit\_antoine}{relErr\_max\_fit\_pct}{4.47039504695}
\goldrow{diag}{fit\_antoine}{relErr\_max\_fit\_pct\_leastSquares}{4.90640428544}
\goldmore{3}
\end{tutgolden}

\clearpage
\tutentry{curate37\_vp\_nAminoethylPiperazine}{tutorials/props/curation/curate37\_vp\_nAminoethylPiperazine}
\begin{tutscheme}
\begin{tikzpicture}[>={Stealth[length=2mm]}, every node/.style={font=\scriptsize}]
  \node[draw, rounded corners=2pt, align=center, minimum width=3.1cm, minimum height=0.8cm, fill=black!4] at (0.0,-0.0) {\textbf{fit\_antoine}\\[-1pt]{\tiny vaporPressureFit}};
\end{tikzpicture}
\end{tutscheme}
\tutwhat{Antoine vapour pressure of nAminoethylPiperazine fitted by the curation pipeline to the measured vapour pressures of every usable laboratory (ThermoML locates, the articles source), outliers discarded by name}
\begin{tutnarrative}
the vapour pressure of nAminoethylPiperazine, REFITTED to measurements (DEV.md 4c C281 and C283). Written by `bin/choupo-curate vp nAminoethylPiperazine`; everything above the RESULT paragraph is the DECLARATION, committed before any fit ran, so it can be read against the result.

WHY. The catalogue record carried `vaporPressure \{ model AmbroseWalton; \}`, a corresponding-states ESTIMATE from Tc, Pc and omega that never reads the record's own Tb (C279 measured a median 19.6 \% AAD over the 25 such records with archive data). Here a three-parameter Antoine is fitted to EVERY usable laboratory's measurements of it.

STUDIES (one per article, extracted by `choupo-thermoml extract-vp`; ranked by points, ties by DOI): 10.1016/j.jct.2009.09.003 22 pts 295.7-338.1 K Efimova; Emel'yanenko; Verevkin; Chernyak SET ASIDE BY NAME: (none)

FIT AND GRADED ACCEPTANCE (DEV.md 4c C285, Vitor 2026-11-10; the evidence, objective and outliers of C283; declared before the fit). EVIDENCE: every study above is fitted; none is held out (C283 amended: laboratories disagree with each other by more than they claim, so the test is agreement across all of them, not one laboratory's claimed precision). OBJECTIVE: `objective \{ method twoStage; \}` -- least squares in log10 P for a first estimate, then the LARGEST relative error in P, (P\_model - P)/P, over the kept points minimised from it (Vitor's practice, C280); both answers kept. OUTLIERS: after the least-squares stage a point with a relative P error above 15 \% (three times the acceptance) is DISCARDED BY NAME (its study DOI, T and P), listed in the result and the record, and the fit repeated without it, once (`outliers \{ maxRelativeError 15 percent; \}`). ACCEPTANCE, GRADED, the DEFAULT of C285: the maximum relative P error over all kept points is graded clean up to 5 \%, warning up to 10 \%, severe up to 20 \% and refused above; where the normal-boiling-point condition below applies, the absolute gap is graded clean up to 1 K, warning up to 2 K, severe up to 4 K and refused above. The fit's grade is the WORSE of the two. A clean, warning or severe fit is ACCEPTED and promoted (autoPromoted, humanReviewed false) with `fitGrade <grade>;` and the deciding numbers in its provenance entry, which the component loader voices on every run that uses it (a warning or a severe grade); a refused one is not promoted. This case declares no `system/curationDict`, so the default applies; a case that needs another acceptance declares it there, with its reason. The AmbroseWalton curve's errors on the same points are reported beside the fit's, as information.

NORMAL BOILING POINT: the record's Tb 495.15 K is judged against the range of the KEPT points (all studies span 295.7-338.1 K before any outlier is discarded): if Tb lies OUTSIDE that range, the fitted Antoine's own normal boiling point (P = 101.325 kPa) is GRADED by its absolute gap to it (the temperature limits of the acceptance below); inside it, the gap is reported and is not a condition.

EARLIER FITS, SUPERSEDED BY THIS DECLARATION @@EARLIER-BEGIN@@ C283, every usable study fitted and none held out, accepted only at max relative error <= 5 \% and Tnb within 1 K (superseded by the graded acceptance of C285): WHAT THE RUN FOUND (pinned in `expected`; measured, not argued). Objective twoStage. Antoine log10(Psat/bar) = 6.171252 - 2614.8819/(T + -43.5764), in-sample R2 0.999615 over 22 kept points, 295.7-338.1 K (10.1016/j.jct.2009.09.003 22/22 kept). Relative error in P over the kept points: mean 1.7273 \%, max 3.6564 \%; AmbroseWalton on the same points: mean 14.4097 \%, max 20.8514 \%. Outliers discarded: none. Its own normal boiling point (101.325 kPa): 467.69 K. NOT PROMOTED by the DEFAULT acceptance of DEV.md 4c C283 (no system/curationDict): 1 studies, 22 points kept, 0 discarded as outliers; max relative error 3.656 \% (mean 1.727 \%) against 5 \% passes, but Tb lies outside the kept points' range and the normal boiling point 467.69 K, -27.46 K from the record's Tb 495.15 K is not within 1 K -- the gap is announced; AmbroseWalton on the same points mean 14.41 \% / max 20.85 \%; the record keeps AmbroseWalton @@EARLIER-END@@

NO MEASURED POINT IS IN THIS TREE: constant/archiveDatasets declares every series by DOI, NIST URL and sha256; `bin/choupo-curate fetch` regenerates them into constant/archiveExtracts/ (gitignored).

RESULT @@RESULT-BEGIN@@ WHAT THE RUN FOUND (pinned in `expected`; measured, not argued). Objective twoStage. Antoine log10(Psat/bar) = 6.171252 - 2614.8819/(T + -43.5764), in-sample R2 0.999615 over 22 kept points, 295.7-338.1 K (10.1016/j.jct.2009.09.003 22/22 kept). Relative error in P over the kept points: mean 1.7273 \%, max 3.6564 \%; AmbroseWalton on the same points: mean 14.4097 \%, max 20.8514 \%. Outliers discarded: none. Its own normal boiling point (101.325 kPa): 467.69 K. NOT PROMOTED by the DEFAULT graded acceptance of DEV.md 4c C285 (no system/curationDict): fit grade refused -- 1 studies, 22 points kept, 0 discarded as outliers; max relative error 3.656 \% (mean 1.727 \%), grade clean (limits 5 / 10 / 20 \%); Tb outside the kept points' range, normal boiling point 467.69 K, -27.46 K from the record's Tb 495.15 K, grade refused (limits 1 / 2 / 4 K); AmbroseWalton on the same points mean 14.41 \% / max 20.85 \%; the record keeps AmbroseWalton @@RESULT-END@@

\end{tutnarrative}
\begin{tutgolden}
\goldrow{diag}{fit\_antoine}{A}{6.17125225758}
\goldrow{diag}{fit\_antoine}{B}{2614.88187204}
\goldrow{diag}{fit\_antoine}{C}{-43.5763690216}
\goldrow{diag}{fit\_antoine}{R2}{0.999615215845}
\goldrow{diag}{fit\_antoine}{dHvap\_kJ\_per\_mol}{65.9708375265}
\goldrow{diag}{fit\_antoine}{minimax\_evaluations}{670}
\goldrow{diag}{fit\_antoine}{minimax\_max\_scaled\_final}{0.0365635781836}
\goldrow{diag}{fit\_antoine}{minimax\_max\_scaled\_start}{0.0401734432781}
\goldrow{diag}{fit\_antoine}{n\_heldout}{0}
\goldrow{diag}{fit\_antoine}{n\_outliers\_dropped}{0}
\goldrow{diag}{fit\_antoine}{n\_points}{22}
\goldrow{diag}{fit\_antoine}{objective\_method}{2}
\goldrow{diag}{fit\_antoine}{relErr\_max\_fit\_pct}{3.65635781836}
\goldrow{diag}{fit\_antoine}{relErr\_max\_fit\_pct\_leastSquares}{4.01734432781}
\goldmore{3}
\end{tutgolden}

\clearpage
\tutentry{curate38\_vp\_nPentylAcetate}{tutorials/props/curation/curate38\_vp\_nPentylAcetate}
\begin{tutscheme}
\begin{tikzpicture}[>={Stealth[length=2mm]}, every node/.style={font=\scriptsize}]
  \node[draw, rounded corners=2pt, align=center, minimum width=3.1cm, minimum height=0.8cm, fill=black!4] at (0.0,-0.0) {\textbf{fit\_antoine}\\[-1pt]{\tiny vaporPressureFit}};
\end{tikzpicture}
\end{tutscheme}
\tutwhat{Antoine vapour pressure of nPentylAcetate fitted by the curation pipeline to the measured vapour pressures of every usable laboratory (ThermoML locates, the articles source), outliers discarded by name}
\begin{tutnarrative}
the vapour pressure of nPentylAcetate, REFITTED to measurements (DEV.md 4c C281 and C283). Written by `bin/choupo-curate vp nPentylAcetate`; everything above the RESULT paragraph is the DECLARATION, committed before any fit ran, so it can be read against the result.

WHY. The catalogue record carried `vaporPressure \{ model AmbroseWalton; \}`, a corresponding-states ESTIMATE from Tc, Pc and omega that never reads the record's own Tb (C279 measured a median 19.6 \% AAD over the 25 such records with archive data). Here a three-parameter Antoine is fitted to EVERY usable laboratory's measurements of it.

STUDIES (one per article, extracted by `choupo-thermoml extract-vp`; ranked by points, ties by DOI): 10.1016/j.fluid.2012.08.003 15 pts 274.2-313.4 K Samarov; Nazmutdinov; Verevkin SET ASIDE BY NAME: 10.1021/acs.jced.9b00068: 1 pressure point(s), fewer than 5 10.1021/je0301850: 3 pressure point(s), fewer than 5 10.1021/je101301s: 4 pressure point(s), fewer than 5 10.1021/je900812v: je900812v.xml: 1 of 1 pressure rows of nPentylAcetate declare no stated uncertainty -- a partial uncertainty is no uncertainty for a held-out band. Nothing was written.

FIT AND GRADED ACCEPTANCE (DEV.md 4c C285, Vitor 2026-11-10; the evidence, objective and outliers of C283; declared before the fit). EVIDENCE: every study above is fitted; none is held out (C283 amended: laboratories disagree with each other by more than they claim, so the test is agreement across all of them, not one laboratory's claimed precision). OBJECTIVE: `objective \{ method twoStage; \}` -- least squares in log10 P for a first estimate, then the LARGEST relative error in P, (P\_model - P)/P, over the kept points minimised from it (Vitor's practice, C280); both answers kept. OUTLIERS: after the least-squares stage a point with a relative P error above 15 \% (three times the acceptance) is DISCARDED BY NAME (its study DOI, T and P), listed in the result and the record, and the fit repeated without it, once (`outliers \{ maxRelativeError 15 percent; \}`). ACCEPTANCE, GRADED, the DEFAULT of C285: the maximum relative P error over all kept points is graded clean up to 5 \%, warning up to 10 \%, severe up to 20 \% and refused above; where the normal-boiling-point condition below applies, the absolute gap is graded clean up to 1 K, warning up to 2 K, severe up to 4 K and refused above. The fit's grade is the WORSE of the two. A clean, warning or severe fit is ACCEPTED and promoted (autoPromoted, humanReviewed false) with `fitGrade <grade>;` and the deciding numbers in its provenance entry, which the component loader voices on every run that uses it (a warning or a severe grade); a refused one is not promoted. This case declares no `system/curationDict`, so the default applies; a case that needs another acceptance declares it there, with its reason. The AmbroseWalton curve's errors on the same points are reported beside the fit's, as information.

NORMAL BOILING POINT: the record's Tb 422.15 K is judged against the range of the KEPT points (all studies span 274.2-313.4 K before any outlier is discarded): if Tb lies OUTSIDE that range, the fitted Antoine's own normal boiling point (P = 101.325 kPa) is GRADED by its absolute gap to it (the temperature limits of the acceptance below); inside it, the gap is reported and is not a condition.

EARLIER FITS, SUPERSEDED BY THIS DECLARATION @@EARLIER-BEGIN@@ C283, every usable study fitted and none held out, accepted only at max relative error <= 5 \% and Tnb within 1 K (superseded by the graded acceptance of C285): WHAT THE RUN FOUND (pinned in `expected`; measured, not argued). Objective twoStage. Antoine log10(Psat/bar) = 7.284960 - 3220.3538/(T + 36.8556), in-sample R2 0.999934 over 15 kept points, 274.2-313.4 K (10.1016/j.fluid.2012.08.003 15/15 kept). Relative error in P over the kept points: mean 0.6998 \%, max 1.0682 \%; AmbroseWalton on the same points: mean 35.9095 \%, max 42.0745 \%. Outliers discarded: none. Its own normal boiling point (101.325 kPa): 405.55 K. NOT PROMOTED by the DEFAULT acceptance of DEV.md 4c C283 (no system/curationDict): 1 studies, 15 points kept, 0 discarded as outliers; max relative error 1.068 \% (mean 0.6998 \%) against 5 \% passes, but Tb lies outside the kept points' range and the normal boiling point 405.55 K, -16.60 K from the record's Tb 422.15 K is not within 1 K -- the gap is announced; AmbroseWalton on the same points mean 35.91 \% / max 42.07 \%; the record keeps AmbroseWalton @@EARLIER-END@@

NO MEASURED POINT IS IN THIS TREE: constant/archiveDatasets declares every series by DOI, NIST URL and sha256; `bin/choupo-curate fetch` regenerates them into constant/archiveExtracts/ (gitignored).

RESULT @@RESULT-BEGIN@@ WHAT THE RUN FOUND (pinned in `expected`; measured, not argued). Objective twoStage. Antoine log10(Psat/bar) = 7.284960 - 3220.3538/(T + 36.8556), in-sample R2 0.999934 over 15 kept points, 274.2-313.4 K (10.1016/j.fluid.2012.08.003 15/15 kept). Relative error in P over the kept points: mean 0.6998 \%, max 1.0682 \%; AmbroseWalton on the same points: mean 35.9095 \%, max 42.0745 \%. Outliers discarded: none. Its own normal boiling point (101.325 kPa): 405.55 K. NOT PROMOTED by the DEFAULT graded acceptance of DEV.md 4c C285 (no system/curationDict): fit grade refused -- 1 studies, 15 points kept, 0 discarded as outliers; max relative error 1.068 \% (mean 0.6998 \%), grade clean (limits 5 / 10 / 20 \%); Tb outside the kept points' range, normal boiling point 405.55 K, -16.60 K from the record's Tb 422.15 K, grade refused (limits 1 / 2 / 4 K); AmbroseWalton on the same points mean 35.91 \% / max 42.07 \%; the record keeps AmbroseWalton @@RESULT-END@@

\end{tutnarrative}
\begin{tutgolden}
\goldrow{diag}{fit\_antoine}{A}{7.28495972567}
\goldrow{diag}{fit\_antoine}{B}{3220.35376765}
\goldrow{diag}{fit\_antoine}{C}{36.855645453}
\goldrow{diag}{fit\_antoine}{R2}{0.999934369942}
\goldrow{diag}{fit\_antoine}{dHvap\_kJ\_per\_mol}{49.3606821797}
\goldrow{diag}{fit\_antoine}{minimax\_evaluations}{425}
\goldrow{diag}{fit\_antoine}{minimax\_max\_scaled\_final}{0.0106821667998}
\goldrow{diag}{fit\_antoine}{minimax\_max\_scaled\_start}{0.0133442743264}
\goldrow{diag}{fit\_antoine}{n\_heldout}{0}
\goldrow{diag}{fit\_antoine}{n\_outliers\_dropped}{0}
\goldrow{diag}{fit\_antoine}{n\_points}{15}
\goldrow{diag}{fit\_antoine}{objective\_method}{2}
\goldrow{diag}{fit\_antoine}{relErr\_max\_fit\_pct}{1.06821667998}
\goldrow{diag}{fit\_antoine}{relErr\_max\_fit\_pct\_leastSquares}{1.33442743264}
\goldmore{3}
\end{tutgolden}

\clearpage
\tutentry{curate39\_vp\_phenanthrene}{tutorials/props/curation/curate39\_vp\_phenanthrene}
\begin{tutscheme}
\begin{tikzpicture}[>={Stealth[length=2mm]}, every node/.style={font=\scriptsize}]
  \node[draw, rounded corners=2pt, align=center, minimum width=3.1cm, minimum height=0.8cm, fill=black!4] at (0.0,-0.0) {\textbf{fit\_antoine}\\[-1pt]{\tiny vaporPressureFit}};
\end{tikzpicture}
\end{tutscheme}
\tutwhat{Antoine vapour pressure of phenanthrene fitted by the curation pipeline to the measured vapour pressures of every usable laboratory (ThermoML locates, the articles source), outliers discarded by name}
\begin{tutnarrative}
the vapour pressure of phenanthrene, REFITTED to measurements (DEV.md 4c C281 and C283). Written by `bin/choupo-curate vp phenanthrene`; everything above the RESULT paragraph is the DECLARATION, committed before any fit ran, so it can be read against the result.

WHY. The catalogue record carried `vaporPressure \{ model AmbroseWalton; \}`, a corresponding-states ESTIMATE from Tc, Pc and omega that never reads the record's own Tb (C279 measured a median 19.6 \% AAD over the 25 such records with archive data). Here a three-parameter Antoine is fitted to EVERY usable laboratory's measurements of it.

STUDIES (one per article, extracted by `choupo-thermoml extract-vp`; ranked by points, ties by DOI): 10.1016/j.jct.2014.07.006 68 pts 400-480 K Ximello; Flores; Rojas; Camarillo et al. 10.1016/j.tca.2016.03.013 10 pts 384-420 K Sanchez-Bulas; Cruz-Vasquez; Hernandez-Obregon; Rojas SET ASIDE BY NAME: (none)

FIT AND GRADED ACCEPTANCE (DEV.md 4c C285, Vitor 2026-11-10; the evidence, objective and outliers of C283; declared before the fit). EVIDENCE: every study above is fitted; none is held out (C283 amended: laboratories disagree with each other by more than they claim, so the test is agreement across all of them, not one laboratory's claimed precision). OBJECTIVE: `objective \{ method twoStage; \}` -- least squares in log10 P for a first estimate, then the LARGEST relative error in P, (P\_model - P)/P, over the kept points minimised from it (Vitor's practice, C280); both answers kept. OUTLIERS: after the least-squares stage a point with a relative P error above 15 \% (three times the acceptance) is DISCARDED BY NAME (its study DOI, T and P), listed in the result and the record, and the fit repeated without it, once (`outliers \{ maxRelativeError 15 percent; \}`). ACCEPTANCE, GRADED, the DEFAULT of C285: the maximum relative P error over all kept points is graded clean up to 5 \%, warning up to 10 \%, severe up to 20 \% and refused above; where the normal-boiling-point condition below applies, the absolute gap is graded clean up to 1 K, warning up to 2 K, severe up to 4 K and refused above. The fit's grade is the WORSE of the two. A clean, warning or severe fit is ACCEPTED and promoted (autoPromoted, humanReviewed false) with `fitGrade <grade>;` and the deciding numbers in its provenance entry, which the component loader voices on every run that uses it (a warning or a severe grade); a refused one is not promoted. This case declares no `system/curationDict`, so the default applies; a case that needs another acceptance declares it there, with its reason. The AmbroseWalton curve's errors on the same points are reported beside the fit's, as information.

NORMAL BOILING POINT: the record's Tb 610.03 K is judged against the range of the KEPT points (all studies span 384-480 K before any outlier is discarded): if Tb lies OUTSIDE that range, the fitted Antoine's own normal boiling point (P = 101.325 kPa) is GRADED by its absolute gap to it (the temperature limits of the acceptance below); inside it, the gap is reported and is not a condition.

EARLIER FITS, SUPERSEDED BY THIS DECLARATION @@EARLIER-BEGIN@@ C283, every usable study fitted and none held out, accepted only at max relative error <= 5 \% and Tnb within 1 K (superseded by the graded acceptance of C285): WHAT THE RUN FOUND (pinned in `expected`; measured, not argued). Objective twoStage. Antoine log10(Psat/bar) = 7.753722 - 5573.9627/(T + 126.3843), in-sample R2 0.998979 over 77 kept points, 388-480 K (10.1016/j.jct.2014.07.006 68/68 kept; 10.1016/j.tca.2016.03.013 9/10 kept). Relative error in P over the kept points: mean 2.9881 \%, max 5.9412 \%; AmbroseWalton on the same points: mean 13.2720 \%, max 31.9634 \%. Outliers discarded: 10.1016/j.tca.2016.03.013 T=384 K P=56.7 Pa err=15.65 \%. Its own normal boiling point (101.325 kPa): 593.02 K. NOT PROMOTED by the DEFAULT acceptance of DEV.md 4c C283 (no system/curationDict): 2 studies, 77 points kept, 1 discarded as outliers; max relative error 5.941 \% (mean 2.988 \%) against 5 \%; AmbroseWalton on the same points mean 13.27 \% / max 31.96 \%; the record keeps AmbroseWalton @@EARLIER-END@@

NO MEASURED POINT IS IN THIS TREE: constant/archiveDatasets declares every series by DOI, NIST URL and sha256; `bin/choupo-curate fetch` regenerates them into constant/archiveExtracts/ (gitignored).

RESULT @@RESULT-BEGIN@@ WHAT THE RUN FOUND (pinned in `expected`; measured, not argued). Objective twoStage. Antoine log10(Psat/bar) = 7.753722 - 5573.9627/(T + 126.3843), in-sample R2 0.998979 over 77 kept points, 388-480 K (10.1016/j.jct.2014.07.006 68/68 kept; 10.1016/j.tca.2016.03.013 9/10 kept). Relative error in P over the kept points: mean 2.9881 \%, max 5.9412 \%; AmbroseWalton on the same points: mean 13.2720 \%, max 31.9634 \%. Outliers discarded: 10.1016/j.tca.2016.03.013 T=384 K P=56.7 Pa err=15.65 \%. Its own normal boiling point (101.325 kPa): 593.02 K. NOT PROMOTED by the DEFAULT graded acceptance of DEV.md 4c C285 (no system/curationDict): fit grade refused -- 2 studies, 77 points kept, 1 discarded as outliers; max relative error 5.941 \% (mean 2.988 \%), grade warning (limits 5 / 10 / 20 \%); Tb outside the kept points' range, normal boiling point 593.02 K, -17.01 K from the record's Tb 610.03 K, grade refused (limits 1 / 2 / 4 K); AmbroseWalton on the same points mean 13.27 \% / max 31.96 \%; the record keeps AmbroseWalton @@RESULT-END@@

\end{tutnarrative}
\begin{tutgolden}
\goldrow{diag}{fit\_antoine}{A}{7.75372175025}
\goldrow{diag}{fit\_antoine}{B}{5573.96266648}
\goldrow{diag}{fit\_antoine}{C}{126.384289078}
\goldrow{diag}{fit\_antoine}{R2}{0.998979199706}
\goldrow{diag}{fit\_antoine}{dHvap\_kJ\_per\_mol}{66.8652577344}
\goldrow{diag}{fit\_antoine}{minimax\_evaluations}{793}
\goldrow{diag}{fit\_antoine}{minimax\_max\_scaled\_final}{0.0594120768538}
\goldrow{diag}{fit\_antoine}{minimax\_max\_scaled\_start}{0.0804068053883}
\goldrow{diag}{fit\_antoine}{n\_heldout}{0}
\goldrow{diag}{fit\_antoine}{n\_outliers\_dropped}{1}
\goldrow{diag}{fit\_antoine}{n\_points}{77}
\goldrow{diag}{fit\_antoine}{objective\_method}{2}
\goldrow{diag}{fit\_antoine}{relErr\_max\_fit\_pct}{5.94120768538}
\goldrow{diag}{fit\_antoine}{relErr\_max\_fit\_pct\_leastSquares}{8.04068053883}
\goldmore{3}
\end{tutgolden}

\clearpage
\tutentry{curate40\_vp\_pyrene}{tutorials/props/curation/curate40\_vp\_pyrene}
\begin{tutscheme}
\begin{tikzpicture}[>={Stealth[length=2mm]}, every node/.style={font=\scriptsize}]
  \node[draw, rounded corners=2pt, align=center, minimum width=3.1cm, minimum height=0.8cm, fill=black!4] at (0.0,-0.0) {\textbf{fit\_antoine}\\[-1pt]{\tiny vaporPressureFit}};
\end{tikzpicture}
\end{tutscheme}
\tutwhat{Antoine vapour pressure of pyrene fitted by the curation pipeline to the measured vapour pressures of every usable laboratory (ThermoML locates, the articles source), outliers discarded by name}
\begin{tutnarrative}
the vapour pressure of pyrene, REFITTED to measurements (DEV.md 4c C281 and C283). Written by `bin/choupo-curate vp pyrene`; everything above the RESULT paragraph is the DECLARATION, committed before any fit ran, so it can be read against the result.

WHY. The catalogue record carried `vaporPressure \{ model AmbroseWalton; \}`, a corresponding-states ESTIMATE from Tc, Pc and omega that never reads the record's own Tb (C279 measured a median 19.6 \% AAD over the 25 such records with archive data). Here a three-parameter Antoine is fitted to EVERY usable laboratory's measurements of it.

STUDIES (one per article, extracted by `choupo-thermoml extract-vp`; ranked by points, ties by DOI): 10.1021/je800300x 23 pts 298.15-510 K Hanshaw; Nutt; Chickos SET ASIDE BY NAME: (none)

FIT AND GRADED ACCEPTANCE (DEV.md 4c C285, Vitor 2026-11-10; the evidence, objective and outliers of C283; declared before the fit). EVIDENCE: every study above is fitted; none is held out (C283 amended: laboratories disagree with each other by more than they claim, so the test is agreement across all of them, not one laboratory's claimed precision). OBJECTIVE: `objective \{ method twoStage; \}` -- least squares in log10 P for a first estimate, then the LARGEST relative error in P, (P\_model - P)/P, over the kept points minimised from it (Vitor's practice, C280); both answers kept. OUTLIERS: after the least-squares stage a point with a relative P error above 15 \% (three times the acceptance) is DISCARDED BY NAME (its study DOI, T and P), listed in the result and the record, and the fit repeated without it, once (`outliers \{ maxRelativeError 15 percent; \}`). ACCEPTANCE, GRADED, the DEFAULT of C285: the maximum relative P error over all kept points is graded clean up to 5 \%, warning up to 10 \%, severe up to 20 \% and refused above; where the normal-boiling-point condition below applies, the absolute gap is graded clean up to 1 K, warning up to 2 K, severe up to 4 K and refused above. The fit's grade is the WORSE of the two. A clean, warning or severe fit is ACCEPTED and promoted (autoPromoted, humanReviewed false) with `fitGrade <grade>;` and the deciding numbers in its provenance entry, which the component loader voices on every run that uses it (a warning or a severe grade); a refused one is not promoted. This case declares no `system/curationDict`, so the default applies; a case that needs another acceptance declares it there, with its reason. The AmbroseWalton curve's errors on the same points are reported beside the fit's, as information.

NORMAL BOILING POINT: the record's Tb 667.95 K is judged against the range of the KEPT points (all studies span 298.15-510 K before any outlier is discarded): if Tb lies OUTSIDE that range, the fitted Antoine's own normal boiling point (P = 101.325 kPa) is GRADED by its absolute gap to it (the temperature limits of the acceptance below); inside it, the gap is reported and is not a condition.

EARLIER FITS, SUPERSEDED BY THIS DECLARATION @@EARLIER-BEGIN@@ C283, every usable study fitted and none held out, accepted only at max relative error <= 5 \% and Tnb within 1 K (superseded by the graded acceptance of C285): WHAT THE RUN FOUND (pinned in `expected`; measured, not argued). Objective twoStage. Antoine log10(Psat/bar) = 4.304067 - 2443.8287/(T + -91.2424), in-sample R2 0.999964 over 23 kept points, 298.15-510 K (10.1021/je800300x 23/23 kept). Relative error in P over the kept points: mean 2.3269 \%, max 3.5140 \%; AmbroseWalton on the same points: mean 49.2860 \%, max 59.8483 \%. Outliers discarded: none. Its own normal boiling point (101.325 kPa): 659.79 K. NOT PROMOTED by the DEFAULT acceptance of DEV.md 4c C283 (no system/curationDict): 1 studies, 23 points kept, 0 discarded as outliers; max relative error 3.514 \% (mean 2.327 \%) against 5 \% passes, but Tb lies outside the kept points' range and the normal boiling point 659.79 K, -8.16 K from the record's Tb 667.95 K is not within 1 K -- the gap is announced; AmbroseWalton on the same points mean 49.29 \% / max 59.85 \%; the record keeps AmbroseWalton @@EARLIER-END@@

NO MEASURED POINT IS IN THIS TREE: constant/archiveDatasets declares every series by DOI, NIST URL and sha256; `bin/choupo-curate fetch` regenerates them into constant/archiveExtracts/ (gitignored).

RESULT @@RESULT-BEGIN@@ WHAT THE RUN FOUND (pinned in `expected`; measured, not argued). Objective twoStage. Antoine log10(Psat/bar) = 4.304067 - 2443.8287/(T + -91.2424), in-sample R2 0.999964 over 23 kept points, 298.15-510 K (10.1021/je800300x 23/23 kept). Relative error in P over the kept points: mean 2.3269 \%, max 3.5140 \%; AmbroseWalton on the same points: mean 49.2860 \%, max 59.8483 \%. Outliers discarded: none. Its own normal boiling point (101.325 kPa): 659.79 K. NOT PROMOTED by the DEFAULT graded acceptance of DEV.md 4c C285 (no system/curationDict): fit grade refused -- 1 studies, 23 points kept, 0 discarded as outliers; max relative error 3.514 \% (mean 2.327 \%), grade clean (limits 5 / 10 / 20 \%); Tb outside the kept points' range, normal boiling point 659.79 K, -8.16 K from the record's Tb 667.95 K, grade refused (limits 1 / 2 / 4 K); AmbroseWalton on the same points mean 49.29 \% / max 59.85 \%; the record keeps AmbroseWalton @@RESULT-END@@

\end{tutnarrative}
\begin{tutgolden}
\goldrow{diag}{fit\_antoine}{A}{4.30406708414}
\goldrow{diag}{fit\_antoine}{B}{2443.82866644}
\goldrow{diag}{fit\_antoine}{C}{-91.2423952632}
\goldrow{diag}{fit\_antoine}{R2}{0.999963763315}
\goldrow{diag}{fit\_antoine}{dHvap\_kJ\_per\_mol}{69.3961803103}
\goldrow{diag}{fit\_antoine}{minimax\_evaluations}{487}
\goldrow{diag}{fit\_antoine}{minimax\_max\_scaled\_final}{0.0351395931044}
\goldrow{diag}{fit\_antoine}{minimax\_max\_scaled\_start}{0.0484218164156}
\goldrow{diag}{fit\_antoine}{n\_heldout}{0}
\goldrow{diag}{fit\_antoine}{n\_outliers\_dropped}{0}
\goldrow{diag}{fit\_antoine}{n\_points}{23}
\goldrow{diag}{fit\_antoine}{objective\_method}{2}
\goldrow{diag}{fit\_antoine}{relErr\_max\_fit\_pct}{3.51395931044}
\goldrow{diag}{fit\_antoine}{relErr\_max\_fit\_pct\_leastSquares}{4.84218164156}
\goldmore{3}
\end{tutgolden}

\clearpage
\tutentry{curate41\_vp\_squalane}{tutorials/props/curation/curate41\_vp\_squalane}
\begin{tutscheme}
\begin{tikzpicture}[>={Stealth[length=2mm]}, every node/.style={font=\scriptsize}]
  \node[draw, rounded corners=2pt, align=center, minimum width=3.1cm, minimum height=0.8cm, fill=black!4] at (0.0,-0.0) {\textbf{fit\_antoine}\\[-1pt]{\tiny vaporPressureFit}};
\end{tikzpicture}
\end{tutscheme}
\tutwhat{Antoine vapour pressure of squalane fitted by the curation pipeline to the measured vapour pressures of every usable laboratory (ThermoML locates, the articles source), outliers discarded by name}
\begin{tutnarrative}
the vapour pressure of squalane, REFITTED to measurements (DEV.md 4c C281 and C283). Written by `bin/choupo-curate vp squalane`; everything above the RESULT paragraph is the DECLARATION, committed before any fit ran, so it can be read against the result.

WHY. The catalogue record carried `vaporPressure \{ model AmbroseWalton; \}`, a corresponding-states ESTIMATE from Tc, Pc and omega that never reads the record's own Tb (C279 measured a median 19.6 \% AAD over the 25 such records with archive data). Here a three-parameter Antoine is fitted to EVERY usable laboratory's measurements of it.

STUDIES (one per article, extracted by `choupo-thermoml extract-vp`; ranked by points, ties by DOI): 10.1016/j.jct.2019.03.032 32 pts 298.15-686.4 K Zafar; Chickos SET ASIDE BY NAME: (none)

FIT AND GRADED ACCEPTANCE (DEV.md 4c C285, Vitor 2026-11-10; the evidence, objective and outliers of C283; declared before the fit). EVIDENCE: every study above is fitted; none is held out (C283 amended: laboratories disagree with each other by more than they claim, so the test is agreement across all of them, not one laboratory's claimed precision). OBJECTIVE: `objective \{ method twoStage; \}` -- least squares in log10 P for a first estimate, then the LARGEST relative error in P, (P\_model - P)/P, over the kept points minimised from it (Vitor's practice, C280); both answers kept. OUTLIERS: after the least-squares stage a point with a relative P error above 15 \% (three times the acceptance) is DISCARDED BY NAME (its study DOI, T and P), listed in the result and the record, and the fit repeated without it, once (`outliers \{ maxRelativeError 15 percent; \}`). ACCEPTANCE, GRADED, the DEFAULT of C285: the maximum relative P error over all kept points is graded clean up to 5 \%, warning up to 10 \%, severe up to 20 \% and refused above; where the normal-boiling-point condition below applies, the absolute gap is graded clean up to 1 K, warning up to 2 K, severe up to 4 K and refused above. The fit's grade is the WORSE of the two. A clean, warning or severe fit is ACCEPTED and promoted (autoPromoted, humanReviewed false) with `fitGrade <grade>;` and the deciding numbers in its provenance entry, which the component loader voices on every run that uses it (a warning or a severe grade); a refused one is not promoted. This case declares no `system/curationDict`, so the default applies; a case that needs another acceptance declares it there, with its reason. The AmbroseWalton curve's errors on the same points are reported beside the fit's, as information.

NORMAL BOILING POINT: the record's Tb 623.15 K is judged against the range of the KEPT points (all studies span 298.15-686.4 K before any outlier is discarded): if Tb lies OUTSIDE that range, the fitted Antoine's own normal boiling point (P = 101.325 kPa) is GRADED by its absolute gap to it (the temperature limits of the acceptance below); inside it, the gap is reported and is not a condition.

EARLIER FITS, SUPERSEDED BY THIS DECLARATION @@EARLIER-BEGIN@@ C283, every usable study fitted and none held out, accepted only at max relative error <= 5 \% and Tnb within 1 K (superseded by the graded acceptance of C285): WHAT THE RUN FOUND (pinned in `expected`; measured, not argued). Objective twoStage. Antoine log10(Psat/bar) = 5.524774 - 3196.0462/(T + -103.1706), in-sample R2 0.999906 over 32 kept points, 298.15-686.4 K (10.1016/j.jct.2019.03.032 32/32 kept). Relative error in P over the kept points: mean 5.4320 \%, max 9.4322 \%; AmbroseWalton on the same points: mean 69.1733 \%, max 446.1594 \%. Outliers discarded: none. Its own normal boiling point (101.325 kPa): 682.26 K. NOT PROMOTED by the DEFAULT acceptance of DEV.md 4c C283 (no system/curationDict): 1 studies, 32 points kept, 0 discarded as outliers; max relative error 9.432 \% (mean 5.432 \%) against 5 \%; AmbroseWalton on the same points mean 69.17 \% / max 446.2 \%; the record keeps AmbroseWalton @@EARLIER-END@@

NO MEASURED POINT IS IN THIS TREE: constant/archiveDatasets declares every series by DOI, NIST URL and sha256; `bin/choupo-curate fetch` regenerates them into constant/archiveExtracts/ (gitignored).

RESULT @@RESULT-BEGIN@@ WHAT THE RUN FOUND (pinned in `expected`; measured, not argued). Objective twoStage. Antoine log10(Psat/bar) = 5.524774 - 3196.0462/(T + -103.1706), in-sample R2 0.999906 over 32 kept points, 298.15-686.4 K (10.1016/j.jct.2019.03.032 32/32 kept). Relative error in P over the kept points: mean 5.4320 \%, max 9.4322 \%; AmbroseWalton on the same points: mean 69.1733 \%, max 446.1594 \%. Outliers discarded: none. Its own normal boiling point (101.325 kPa): 682.26 K. PROMOTED on 2026-11-10 into data/standards/components/squalane.dat as reviewStatus autoPromoted, humanReviewed false, because: fit grade warning under the DEFAULT graded acceptance of DEV.md 4c C285 (no system/curationDict), 1 studies, 32 points kept, 0 discarded as outliers; max relative error 9.432 \% (mean 5.432 \%), grade warning (limits 5 / 10 / 20 \%); Tb inside the kept points' range, normal boiling point 682.26 K, +59.11 K from the record's Tb 623.15 K (reported, not a condition); AmbroseWalton on the same points mean 69.17 \% / max 446.2 \% @@RESULT-END@@

\end{tutnarrative}
\begin{tutgolden}
\goldrow{diag}{fit\_antoine}{A}{5.52477413003}
\goldrow{diag}{fit\_antoine}{B}{3196.04623316}
\goldrow{diag}{fit\_antoine}{C}{-103.170561023}
\goldrow{diag}{fit\_antoine}{R2}{0.999906346688}
\goldrow{diag}{fit\_antoine}{dHvap\_kJ\_per\_mol}{84.7497844015}
\goldrow{diag}{fit\_antoine}{minimax\_evaluations}{477}
\goldrow{diag}{fit\_antoine}{minimax\_max\_scaled\_final}{0.094321980772}
\goldrow{diag}{fit\_antoine}{minimax\_max\_scaled\_start}{0.130819654301}
\goldrow{diag}{fit\_antoine}{n\_heldout}{0}
\goldrow{diag}{fit\_antoine}{n\_outliers\_dropped}{0}
\goldrow{diag}{fit\_antoine}{n\_points}{32}
\goldrow{diag}{fit\_antoine}{objective\_method}{2}
\goldrow{diag}{fit\_antoine}{relErr\_max\_fit\_pct}{9.4321980772}
\goldrow{diag}{fit\_antoine}{relErr\_max\_fit\_pct\_leastSquares}{13.0819654301}
\goldmore{3}
\end{tutgolden}

\clearpage
\tutentry{curate42\_vp\_triethanolamine}{tutorials/props/curation/curate42\_vp\_triethanolamine}
\begin{tutscheme}
\begin{tikzpicture}[>={Stealth[length=2mm]}, every node/.style={font=\scriptsize}]
  \node[draw, rounded corners=2pt, align=center, minimum width=3.1cm, minimum height=0.8cm, fill=black!4] at (0.0,-0.0) {\textbf{fit\_antoine}\\[-1pt]{\tiny vaporPressureFit}};
\end{tikzpicture}
\end{tutscheme}
\tutwhat{Antoine vapour pressure of triethanolamine fitted by the curation pipeline to the measured vapour pressures of every usable laboratory (ThermoML locates, the articles source), outliers discarded by name}
\begin{tutnarrative}
the vapour pressure of triethanolamine, REFITTED to measurements (DEV.md 4c C281 and C283). Written by `bin/choupo-curate vp triethanolamine`; everything above the RESULT paragraph is the DECLARATION, committed before any fit ran, so it can be read against the result.

WHY. The catalogue record carried `vaporPressure \{ model AmbroseWalton; \}`, a corresponding-states ESTIMATE from Tc, Pc and omega that never reads the record's own Tb (C279 measured a median 19.6 \% AAD over the 25 such records with archive data). Here a three-parameter Antoine is fitted to EVERY usable laboratory's measurements of it.

STUDIES (one per article, extracted by `choupo-thermoml extract-vp`; ranked by points, ties by DOI): 10.1016/j.tca.2018.05.007 55 pts 365-515 K Zghal; Farjas; Camps; Dammak et al. SET ASIDE BY NAME: (none)

FIT AND GRADED ACCEPTANCE (DEV.md 4c C285, Vitor 2026-11-10; the evidence, objective and outliers of C283; declared before the fit). EVIDENCE: every study above is fitted; none is held out (C283 amended: laboratories disagree with each other by more than they claim, so the test is agreement across all of them, not one laboratory's claimed precision). OBJECTIVE: `objective \{ method twoStage; \}` -- least squares in log10 P for a first estimate, then the LARGEST relative error in P, (P\_model - P)/P, over the kept points minimised from it (Vitor's practice, C280); both answers kept. OUTLIERS: after the least-squares stage a point with a relative P error above 15 \% (three times the acceptance) is DISCARDED BY NAME (its study DOI, T and P), listed in the result and the record, and the fit repeated without it, once (`outliers \{ maxRelativeError 15 percent; \}`). ACCEPTANCE, GRADED, the DEFAULT of C285: the maximum relative P error over all kept points is graded clean up to 5 \%, warning up to 10 \%, severe up to 20 \% and refused above; where the normal-boiling-point condition below applies, the absolute gap is graded clean up to 1 K, warning up to 2 K, severe up to 4 K and refused above. The fit's grade is the WORSE of the two. A clean, warning or severe fit is ACCEPTED and promoted (autoPromoted, humanReviewed false) with `fitGrade <grade>;` and the deciding numbers in its provenance entry, which the component loader voices on every run that uses it (a warning or a severe grade); a refused one is not promoted. This case declares no `system/curationDict`, so the default applies; a case that needs another acceptance declares it there, with its reason. The AmbroseWalton curve's errors on the same points are reported beside the fit's, as information.

NORMAL BOILING POINT: the record's Tb 613 K is judged against the range of the KEPT points (all studies span 365-515 K before any outlier is discarded): if Tb lies OUTSIDE that range, the fitted Antoine's own normal boiling point (P = 101.325 kPa) is GRADED by its absolute gap to it (the temperature limits of the acceptance below); inside it, the gap is reported and is not a condition.

EARLIER FITS, SUPERSEDED BY THIS DECLARATION @@EARLIER-BEGIN@@ C283, every usable study fitted and none held out, accepted only at max relative error <= 5 \% and Tnb within 1 K (superseded by the graded acceptance of C285): WHAT THE RUN FOUND (pinned in `expected`; measured, not argued). Objective twoStage. Antoine log10(Psat/bar) = 6.379786 - 3486.4829/(T + -59.3907), in-sample R2 0.999017 over 54 kept points, 365-515 K (10.1016/j.tca.2018.05.007 54/55 kept). Relative error in P over the kept points: mean 4.8557 \%, max 12.1496 \%; AmbroseWalton on the same points: mean 49.4296 \%, max 83.3204 \%. Outliers discarded: 10.1016/j.tca.2018.05.007 T=465 K P=511.6 Pa err=18.34 \%. Its own normal boiling point (101.325 kPa): 606.37 K. NOT PROMOTED by the DEFAULT acceptance of DEV.md 4c C283 (no system/curationDict): 1 studies, 54 points kept, 1 discarded as outliers; max relative error 12.15 \% (mean 4.856 \%) against 5 \%; AmbroseWalton on the same points mean 49.43 \% / max 83.32 \%; the record keeps AmbroseWalton @@EARLIER-END@@

NO MEASURED POINT IS IN THIS TREE: constant/archiveDatasets declares every series by DOI, NIST URL and sha256; `bin/choupo-curate fetch` regenerates them into constant/archiveExtracts/ (gitignored).

RESULT @@RESULT-BEGIN@@ WHAT THE RUN FOUND (pinned in `expected`; measured, not argued). Objective twoStage. Antoine log10(Psat/bar) = 6.379786 - 3486.4829/(T + -59.3907), in-sample R2 0.999017 over 54 kept points, 365-515 K (10.1016/j.tca.2018.05.007 54/55 kept). Relative error in P over the kept points: mean 4.8557 \%, max 12.1496 \%; AmbroseWalton on the same points: mean 49.4296 \%, max 83.3204 \%. Outliers discarded: 10.1016/j.tca.2018.05.007 T=465 K P=511.6 Pa err=18.34 \%. Its own normal boiling point (101.325 kPa): 606.37 K. NOT PROMOTED by the DEFAULT graded acceptance of DEV.md 4c C285 (no system/curationDict): fit grade refused -- 1 studies, 54 points kept, 1 discarded as outliers; max relative error 12.15 \% (mean 4.856 \%), grade severe (limits 5 / 10 / 20 \%); Tb outside the kept points' range, normal boiling point 606.37 K, -6.63 K from the record's Tb 613 K, grade refused (limits 1 / 2 / 4 K); AmbroseWalton on the same points mean 49.43 \% / max 83.32 \%; the record keeps AmbroseWalton @@RESULT-END@@

\end{tutnarrative}
\begin{tutgolden}
\goldrow{diag}{fit\_antoine}{A}{6.37978569309}
\goldrow{diag}{fit\_antoine}{B}{3486.48293915}
\goldrow{diag}{fit\_antoine}{C}{-59.3907102323}
\goldrow{diag}{fit\_antoine}{R2}{0.999017201306}
\goldrow{diag}{fit\_antoine}{dHvap\_kJ\_per\_mol}{85.2838376743}
\goldrow{diag}{fit\_antoine}{minimax\_evaluations}{455}
\goldrow{diag}{fit\_antoine}{minimax\_max\_scaled\_final}{0.121495803417}
\goldrow{diag}{fit\_antoine}{minimax\_max\_scaled\_start}{0.134483391072}
\goldrow{diag}{fit\_antoine}{n\_heldout}{0}
\goldrow{diag}{fit\_antoine}{n\_outliers\_dropped}{1}
\goldrow{diag}{fit\_antoine}{n\_points}{54}
\goldrow{diag}{fit\_antoine}{objective\_method}{2}
\goldrow{diag}{fit\_antoine}{relErr\_max\_fit\_pct}{12.1495803417}
\goldrow{diag}{fit\_antoine}{relErr\_max\_fit\_pct\_leastSquares}{13.4483391072}
\goldmore{3}
\end{tutgolden}

\clearpage
\tutentry{curate43\_vp\_triethylamine}{tutorials/props/curation/curate43\_vp\_triethylamine}
\begin{tutscheme}
\begin{tikzpicture}[>={Stealth[length=2mm]}, every node/.style={font=\scriptsize}]
  \node[draw, rounded corners=2pt, align=center, minimum width=3.1cm, minimum height=0.8cm, fill=black!4] at (0.0,-0.0) {\textbf{fit\_antoine}\\[-1pt]{\tiny vaporPressureFit}};
\end{tikzpicture}
\end{tutscheme}
\tutwhat{Antoine vapour pressure of triethylamine fitted by the curation pipeline to the measured vapour pressures of every usable laboratory (ThermoML locates, the articles source), outliers discarded by name}
\begin{tutnarrative}
the vapour pressure of triethylamine, REFITTED to measurements (DEV.md 4c C281 and C283). Written by `bin/choupo-curate vp triethylamine`; everything above the RESULT paragraph is the DECLARATION, committed before any fit ran, so it can be read against the result.

WHY. The catalogue record carried `vaporPressure \{ model AmbroseWalton; \}`, a corresponding-states ESTIMATE from Tc, Pc and omega that never reads the record's own Tb (C279 measured a median 19.6 \% AAD over the 25 such records with archive data). Here a three-parameter Antoine is fitted to EVERY usable laboratory's measurements of it.

STUDIES (one per article, extracted by `choupo-thermoml extract-vp`; ranked by points, ties by DOI): 10.1021/je800603z 5 pts 273.04-352.77 K Mokbel; Razzouk; Sawaya; Jose SET ASIDE BY NAME: (none)

FIT AND GRADED ACCEPTANCE (DEV.md 4c C285, Vitor 2026-11-10; the evidence, objective and outliers of C283; declared before the fit). EVIDENCE: every study above is fitted; none is held out (C283 amended: laboratories disagree with each other by more than they claim, so the test is agreement across all of them, not one laboratory's claimed precision). OBJECTIVE: `objective \{ method twoStage; \}` -- least squares in log10 P for a first estimate, then the LARGEST relative error in P, (P\_model - P)/P, over the kept points minimised from it (Vitor's practice, C280); both answers kept. OUTLIERS: after the least-squares stage a point with a relative P error above 15 \% (three times the acceptance) is DISCARDED BY NAME (its study DOI, T and P), listed in the result and the record, and the fit repeated without it, once (`outliers \{ maxRelativeError 15 percent; \}`). ACCEPTANCE, GRADED, the DEFAULT of C285: the maximum relative P error over all kept points is graded clean up to 5 \%, warning up to 10 \%, severe up to 20 \% and refused above; where the normal-boiling-point condition below applies, the absolute gap is graded clean up to 1 K, warning up to 2 K, severe up to 4 K and refused above. The fit's grade is the WORSE of the two. A clean, warning or severe fit is ACCEPTED and promoted (autoPromoted, humanReviewed false) with `fitGrade <grade>;` and the deciding numbers in its provenance entry, which the component loader voices on every run that uses it (a warning or a severe grade); a refused one is not promoted. This case declares no `system/curationDict`, so the default applies; a case that needs another acceptance declares it there, with its reason. The AmbroseWalton curve's errors on the same points are reported beside the fit's, as information.

NORMAL BOILING POINT: the record's Tb 362.9 K is judged against the range of the KEPT points (all studies span 273.04-352.77 K before any outlier is discarded): if Tb lies OUTSIDE that range, the fitted Antoine's own normal boiling point (P = 101.325 kPa) is GRADED by its absolute gap to it (the temperature limits of the acceptance below); inside it, the gap is reported and is not a condition.

EARLIER FITS, SUPERSEDED BY THIS DECLARATION @@EARLIER-BEGIN@@ C283, every usable study fitted and none held out, accepted only at max relative error <= 5 \% and Tnb within 1 K (superseded by the graded acceptance of C285): WHAT THE RUN FOUND (pinned in `expected`; measured, not argued). Objective twoStage. Antoine log10(Psat/bar) = 4.187638 - 1340.7812/(T + -41.8128), in-sample R2 0.999996 over 5 kept points, 273.04-352.77 K (10.1021/je800603z 5/5 kept). Relative error in P over the kept points: mean 0.2398 \%, max 0.2609 \%; AmbroseWalton on the same points: mean 1.0386 \%, max 1.6667 \%. Outliers discarded: none. Its own normal boiling point (101.325 kPa): 362.43 K. PROMOTED on 2026-11-10 into data/standards/components/triethylamine.dat as reviewStatus autoPromoted, humanReviewed false, because: under the DEFAULT acceptance of DEV.md 4c C283 (no system/curationDict), 1 studies, 5 points kept, 0 discarded as outliers; max relative error 0.2609 \% (mean 0.2398 \%) against 5 \%; Tb outside the kept points' range, normal boiling point 362.43 K, -0.47 K from the record's Tb 362.9 K (within 1 K); AmbroseWalton on the same points mean 1.039 \% / max 1.667 \% @@EARLIER-END@@

NO MEASURED POINT IS IN THIS TREE: constant/archiveDatasets declares every series by DOI, NIST URL and sha256; `bin/choupo-curate fetch` regenerates them into constant/archiveExtracts/ (gitignored).

RESULT @@RESULT-BEGIN@@ WHAT THE RUN FOUND (pinned in `expected`; measured, not argued). Objective twoStage. Antoine log10(Psat/bar) = 4.187638 - 1340.7812/(T + -41.8128), in-sample R2 0.999996 over 5 kept points, 273.04-352.77 K (10.1021/je800603z 5/5 kept). Relative error in P over the kept points: mean 0.2398 \%, max 0.2609 \%; AmbroseWalton on the same points: mean 1.0386 \%, max 1.6667 \%. Outliers discarded: none. Its own normal boiling point (101.325 kPa): 362.43 K. PROMOTED on 2026-11-10 into data/standards/components/triethylamine.dat as reviewStatus autoPromoted, humanReviewed false, because: fit grade clean under the DEFAULT graded acceptance of DEV.md 4c C285 (no system/curationDict), 1 studies, 5 points kept, 0 discarded as outliers; max relative error 0.2609 \% (mean 0.2398 \%), grade clean (limits 5 / 10 / 20 \%); Tb outside the kept points' range, normal boiling point 362.43 K, -0.47 K from the record's Tb 362.9 K, grade clean (limits 1 / 2 / 4 K); AmbroseWalton on the same points mean 1.039 \% / max 1.667 \% @@RESULT-END@@

\end{tutnarrative}
\begin{tutgolden}
\goldrow{diag}{fit\_antoine}{A}{4.1876380157}
\goldrow{diag}{fit\_antoine}{B}{1340.78117874}
\goldrow{diag}{fit\_antoine}{C}{-41.8127969415}
\goldrow{diag}{fit\_antoine}{R2}{0.999995953663}
\goldrow{diag}{fit\_antoine}{dHvap\_kJ\_per\_mol}{33.0361801862}
\goldrow{diag}{fit\_antoine}{minimax\_evaluations}{611}
\goldrow{diag}{fit\_antoine}{minimax\_max\_scaled\_final}{0.00260919533853}
\goldrow{diag}{fit\_antoine}{minimax\_max\_scaled\_start}{0.00371493301982}
\goldrow{diag}{fit\_antoine}{n\_heldout}{0}
\goldrow{diag}{fit\_antoine}{n\_outliers\_dropped}{0}
\goldrow{diag}{fit\_antoine}{n\_points}{5}
\goldrow{diag}{fit\_antoine}{objective\_method}{2}
\goldrow{diag}{fit\_antoine}{relErr\_max\_fit\_pct}{0.260919533853}
\goldrow{diag}{fit\_antoine}{relErr\_max\_fit\_pct\_leastSquares}{0.371493301981}
\goldmore{3}
\end{tutgolden}

\subsection{electrolyte}
\grouptheory{electrolyte}
\clearpage
\tutentry{archer01\_nacl\_cold\_to\_hot}{tutorials/props/electrolyte/archer01\_nacl\_cold\_to\_hot}
\begin{tutscheme}
\begin{tikzpicture}[>={Stealth[length=2mm]}, every node/.style={font=\scriptsize}]
  \node[draw, rounded corners=2pt, align=center, minimum width=3.1cm, minimum height=0.8cm, fill=black!4] at (0.0,-0.0) {\textbf{nacl\_0C}\\[-1pt]{\tiny pitzerActivity}};
  \node[draw, rounded corners=2pt, align=center, minimum width=3.1cm, minimum height=0.8cm, fill=black!4] at (3.6,-0.0) {\textbf{nacl\_25C}\\[-1pt]{\tiny pitzerActivity}};
  \node[draw, rounded corners=2pt, align=center, minimum width=3.1cm, minimum height=0.8cm, fill=black!4] at (7.2,-0.0) {\textbf{nacl\_50C}\\[-1pt]{\tiny pitzerActivity}};
  \node[draw, rounded corners=2pt, align=center, minimum width=3.1cm, minimum height=0.8cm, fill=black!4] at (0.0,-1.3) {\textbf{nacl\_100C}\\[-1pt]{\tiny pitzerActivity}};
\end{tikzpicture}
\end{tutscheme}
\tutwhat{NaCl Pitzer 0-100 C vs Archer 1992 -- the cold edge the freezing curve stands on}
\begin{tutnarrative}
The catalogue's NaCl Pitzer surface (Appelo 25 C parameters + SP77 T-functions) against Archer, J. Phys. Chem. Ref. Data 21 (1992) 793, DOI 10.1063/1.555915, Tables 9 (phi) and 10 (gamma\_pm) -- HIS fitted equation's check values ("presented only as a means to test the validity of coding", the paper's own words), at 273, 298, 323 and 373 K.

WHY THE 273 K ROW IS THE POINT: the SP77 temperature extension was fitted ABOVE 25 C, and the freezing-point curve (fpd01\_nacl\_freezing) leans on that surface BELOW 0 C as an announced extrapolation. This case measures what the extrapolation costs at the last anchor before the ice field -- the op announces the T-window violation at 273 K, and the AAD against Archer says whether the announced risk is a real one at the cold edge.

WHAT THIS COMPARISON IS, said precisely: Archer's tables are values of an independent CRITICAL FIT (his global equation, 250-600 K), not raw measurements, and his database overlaps the corpora SP77 drew on. Agreement here is CROSS-EVALUATION CONSISTENCY between two independent fits of overlapping data -- stronger than fit-consistency (our parameters never saw Archer's equation), weaker than agreement with an independent measurement. The Hamer \& Wu case (pitzer\_gamma\_hamer\_wu) plays the same role at 25 C; this case extends it across T.

Temperatures are the table's printed values (273 K, not 273.15). If Archer meant T+0.15 K the shift in gamma\_pm is below 2e-4 -- inside the quoted AADs either way, and stated here so the reading is a recorded choice, not an accident.

\end{tutnarrative}
\begin{tutgolden}
\goldrow{diag}{nacl\_0C}{aw\_1m}{0.967382054689}
\goldrow{diag}{nacl\_0C}{gamma\_AAD\_pct}{1.05461170236}
\goldrow{diag}{nacl\_0C}{gamma\_nMeas}{5}
\goldrow{diag}{nacl\_0C}{gamma\_pm\_1m}{0.641839583122}
\goldrow{diag}{nacl\_0C}{gamma\_pm\_3m}{0.664861748786}
\goldrow{diag}{nacl\_0C}{phi\_1m}{0.920393250284}
\goldrow{diag}{nacl\_0C}{phi\_AAD\_pct}{0.370645531615}
\goldrow{diag}{nacl\_0C}{phi\_nMeas}{5}
\goldrow{diag}{nacl\_25C}{aw\_1m}{0.966828534868}
\goldrow{diag}{nacl\_25C}{gamma\_AAD\_pct}{0.186509289297}
\goldrow{diag}{nacl\_25C}{gamma\_nMeas}{5}
\goldrow{diag}{nacl\_25C}{gamma\_pm\_1m}{0.657130174321}
\goldrow{diag}{nacl\_25C}{gamma\_pm\_3m}{0.71384939713}
\goldrow{diag}{nacl\_25C}{phi\_1m}{0.936278540959}
\goldmore{18}
\end{tutgolden}

\clearpage
\tutentry{bjerrum01\_carbonate\_pH}{tutorials/props/electrolyte/bjerrum01\_carbonate\_pH}
\begin{tutscheme}
\begin{tikzpicture}[>={Stealth[length=2mm]}, every node/.style={font=\scriptsize}]
  \node[draw, rounded corners=2pt, align=center, minimum width=3.1cm, minimum height=0.8cm, fill=black!4] at (0.0,-0.0) {\textbf{acid\_5p0}\\[-1pt]{\tiny speciate}};
  \node[draw, rounded corners=2pt, align=center, minimum width=3.1cm, minimum height=0.8cm, fill=black!4] at (3.6,-0.0) {\textbf{acid\_3p0}\\[-1pt]{\tiny speciate}};
  \node[draw, rounded corners=2pt, align=center, minimum width=3.1cm, minimum height=0.8cm, fill=black!4] at (7.2,-0.0) {\textbf{acid\_2p0}\\[-1pt]{\tiny speciate}};
  \node[draw, rounded corners=2pt, align=center, minimum width=3.1cm, minimum height=0.8cm, fill=black!4] at (0.0,-1.3) {\textbf{acid\_1p5}\\[-1pt]{\tiny speciate}};
  \node[draw, rounded corners=2pt, align=center, minimum width=3.1cm, minimum height=0.8cm, fill=black!4] at (3.6,-1.3) {\textbf{acid\_1p0}\\[-1pt]{\tiny speciate}};
  \node[draw, rounded corners=2pt, align=center, minimum width=3.1cm, minimum height=0.8cm, fill=black!4] at (7.2,-1.3) {\textbf{acid\_0p6}\\[-1pt]{\tiny speciate}};
  \node[draw, rounded corners=2pt, align=center, minimum width=3.1cm, minimum height=0.8cm, fill=black!4] at (0.0,-2.6) {\textbf{acid\_0p35}\\[-1pt]{\tiny speciate}};
  \node[draw, rounded corners=2pt, align=center, minimum width=3.1cm, minimum height=0.8cm, fill=black!4] at (3.6,-2.6) {\textbf{acid\_0p2}\\[-1pt]{\tiny speciate}};
  \node[draw, rounded corners=2pt, align=center, minimum width=3.1cm, minimum height=0.8cm, fill=black!4] at (7.2,-2.6) {\textbf{acid\_0p1}\\[-1pt]{\tiny speciate}};
  \node[draw, rounded corners=2pt, align=center, minimum width=3.1cm, minimum height=0.8cm, fill=black!4] at (0.0,-3.9) {\textbf{acid\_0p05}\\[-1pt]{\tiny speciate}};
  \node[draw, rounded corners=2pt, align=center, minimum width=3.1cm, minimum height=0.8cm, fill=black!4] at (3.6,-3.9) {\textbf{base\_0p0}\\[-1pt]{\tiny speciate}};
  \node[draw, rounded corners=2pt, align=center, minimum width=3.1cm, minimum height=0.8cm, fill=black!4] at (7.2,-3.9) {\textbf{base\_0p1}\\[-1pt]{\tiny speciate}};
\end{tikzpicture}
\end{tutscheme}
\tutwhat{Carbonate speciation across the pH range -- the Bjerrum plot, titrated}
\begin{tutnarrative}
THE SPECIATION OF A CARBONATE SOLUTION ACROSS THE WHOLE pH RANGE -- the diagram every textbook draws, built the way Choupo has to build it.

THE ONE THING THAT MAKES THIS DIFFERENT FROM THE TEXTBOOK VERSION: here the pH is NOT an axis you set -- and that is a CHOICE this case makes, in one word, which you can read three screens down. A textbook draws the fractions against pH by substituting pH into the mass-action expressions, which is arithmetic and always works. Choupo can do that too: `pH 7.8;` is a legal declaration, it fixes [H+], and it DROPS the electroneutrality row -- the run then reports the net charge the composition carries (`pH GIVEN -- this is the net charge the analysis carries`) instead of forcing it to zero. That mode exists because a laboratory sheet reports a MEASURED pH, which is a datum. It is not what this case wants. Every op below declares `pH solve;` instead, so [H+] joins the unknowns and the charge balance is the equation that closes the system. The price of that choice is the whole lesson: each point had to be a beaker somebody could weigh out, the 44 of them differ only in what is IN the beaker, and the x-axis is a RESULT, exactly as it is in a laboratory.

WHAT IS ACTUALLY IN EACH BEAKER. Total carbonate is held at 1 mmol/kg throughout. The 44 points differ only in how much strong acid or strong base sits beside it, and each one is a real neutral mixture somebody could weigh out:

Cl / C\_T = 5 ... 0.05 hydrochloric acid added to dissolved CO2 Na / C\_T = 0 dissolved CO2 alone (H2CO3) Na / C\_T = 1 sodium bicarbonate (NaHCO3) Na / C\_T = 2 sodium carbonate (Na2CO3) Na / C\_T = 2 ... 8 sodium carbonate + caustic soda

That is what `totalsBasis stoichiometric;` declares below: these totals are neutral-molecule feeds, so the water-analysis charge-imbalance check -- which is right for a LAB SHEET, where an imbalance means somebody mismeasured -- does not apply. Here the imbalance IS the titrant, and the electroneutrality equation IS the weak-acid pH. Declaring it changes no number: it changes what the engine says about them.

WHAT TO LOOK FOR, in the order a student should look:

1. THE TWO CROSSOVERS LAND ON THE TWO pK's, AND NOBODY PUT THEM THERE. CO2(aq) and HCO3- are equal at pH 6.341; the record's log K1 is 6.352. HCO3- and CO3-- are equal at pH 10.262; log K2 is 10.329. Both crossover pH's are INTERPOLATED between the two beakers that bracket them -- the axis is 44 sampled points, not a continuum, and no beaker lands exactly on a crossing. Neither pK was used as an input to any point -- they are the mass- action constants inside the network, and the crossover is where the titration happens to arrive. The small shift DOWN in both is the activity coefficients, while a textbook plot takes every gamma as 1. The first is moved by gamma\_HCO3 alone (about 0.975 there; CO2(aq) is neutral) -- gamma\_CO3 is 0.90 at that crossover but carbonate is a trace there and moves nothing. The second is moved by the ratio gamma\_CO3 / gamma\_HCO3, 0.81 against 0.95.

2. THE THREE FRACTIONS DO NOT SUM TO ONE. They dip to 0.99966 through the middle of the range, and the missing carbon is NOT round-off: it is `NaHCO3aq`, the sodium-bicarbonate ION PAIR, which the curated chemistry carries and the three-curve drawing has no place for. The titrant is not a spectator. Look at the `m\_NaHCO3aq` diagnostic -- and note where it is ABSENT: the two beakers with no sodium in them (the acid side, and Na/C\_T = 0) have no such species to report, and the engine REFUSES to report a species that is not in the table rather than printing a zero. Absent and zero are different claims.

3. THE ACTIVITY COEFFICIENT OF THE DIVALENT ION MAKES A U. gamma\_CO3 runs 0.74 at the acid end, up to 0.98 near Na/C\_T = 0 where the ionic strength is lowest, and back down to 0.67 at the alkaline end -- because BOTH ends of a titration are loaded with strong electrolyte and only the middle is dilute. A textbook diagram takes every gamma as 1 at every pH, which is the one assumption that is worst exactly where the diagram is usually read.

4. THE SATURATION INDICES ARE THERE ALL ALONG. Every point reports SI\_nahcolite / SI\_natron / SI\_trona -- how far the solution is from precipitating a sodium carbonate mineral. At 1 mmol/kg they stay far below zero, which is why this stays a clean acid-base diagram; make C\_T a hundred times larger and the diagram acquires a ceiling.

THE CONSTANTS, and their lineage: log K1 = 6.352, log K2 = 10.329 and Kw = 1e-14 at 25 C, from the USGS PHREEQC database (public domain), which is where data/standards/chemistry/\{CO2aq,CO3\}-formation.dat and water-dissociation.dat come from. Each carries the van't Hoff dH and the full analytic log K(T), so this whole diagram can be redrawn at another temperature and both pK's will move -- change `temperature` below and watch which crossover moves further.

NOT VALIDATED AGAINST MEASUREMENT. This case is a STRUCTURAL witness: it shows that the engine's speciation network, its activity model and its charge balance compose into the diagram the chemistry predicts. No titration data is compared against; the two crossovers landing near the two pK's is an internal consistency check on the assembly, not agreement with an experiment.

\end{tutnarrative}
\begin{tutgolden}
\goldrow{diag}{acid\_5p0}{I}{0.00499268}
\goldrow{diag}{acid\_5p0}{aw}{0.999802}
\goldrow{diag}{acid\_5p0}{chargeImposed}{1}
\goldrow{diag}{acid\_5p0}{chargeResidual}{-3.02883e-16}
\goldrow{diag}{acid\_5p0}{gamma\_CO2aq}{1}
\goldrow{diag}{acid\_5p0}{gamma\_CO3}{0.738626}
\goldrow{diag}{acid\_5p0}{gamma\_HCO3}{0.927056}
\goldrow{diag}{acid\_5p0}{m\_CO2aq}{0.000999896}
\goldrow{diag}{acid\_5p0}{m\_CO3}{1.31694e-15}
\goldrow{diag}{acid\_5p0}{m\_HCO3}{1.03591e-07}
\goldrow{diag}{acid\_5p0}{pH}{2.33456}
\goldrow{diag}{acid\_3p0}{I}{0.0029974}
\goldrow{diag}{acid\_3p0}{aw}{0.999874}
\goldrow{diag}{acid\_3p0}{chargeImposed}{1}
\goldmore{635}
\end{tutgolden}

\clearpage
\tutentry{composition01\_nacl}{tutorials/props/electrolyte/composition01\_nacl}
\begin{tutscheme}
\begin{tikzpicture}[>={Stealth[length=2mm]}, every node/.style={font=\scriptsize}]
  \node[draw, rounded corners=2pt, align=center, minimum width=3.1cm, minimum height=0.8cm, fill=black!4] at (0.0,-0.0) {\textbf{naclWater}\\[-1pt]{\tiny speciate}};
\end{tikzpicture}
\end{tutscheme}
\tutwhat{Apparent-salt input: the analysis is given as composition \{ NaCl 2.0 mol/kg \}, expanded to Na + Cl ion totals through NaCl's dissociatesTo, with electroneutrality validated (Sum nu$\ast$charge = 0). Byte-identical to writing totals \{ Na 2.0; Cl 2.0 \} by hand.}
\tutreadme{The speciate op expands each salt to ion totals through its dissociatesTo (NaCl $\to$ Na 1, Cl 1), and validates electroneutrality (Sum nu$\ast$charge = 0) so a formulated-salts input can never silently unbalance charge. This is byte-identical to writing the ions by hand (totals \{ Na 2.0; Cl 2.0 \}): SI\_halite -0.953, feed charge imbalance 0.00\%.}
\tutreadme{composition\{\} is the apparent-salt input mode; analyticalTotals\{\} / totals\{\} stay as the mode for waters measured directly in ions/masters (roadmap Phase B).}
\begin{tutgolden}
\goldrow{diag}{naclWater}{I}{2}
\goldrow{diag}{naclWater}{SI\_halite}{-0.953442}
\goldrow{diag}{naclWater}{aw}{0.930475}
\goldrow{diag}{naclWater}{gamma\_Cl}{1.01683}
\goldrow{diag}{naclWater}{gamma\_Na}{1.01683}
\goldrow{diag}{naclWater}{m\_Cl}{2}
\goldrow{diag}{naclWater}{m\_Na}{2}
\goldrow{diag}{naclWater}{pH}{7.13303}
\goldrow{diag}{naclWater}{chargeImposed}{1}
\goldrow{diag}{naclWater}{chargeResidual}{-6.39009e-17}
\end{tutgolden}

\clearpage
\tutentry{davies01\_band\_edge}{tutorials/props/electrolyte/davies01\_band\_edge}
\begin{tutscheme}
\begin{tikzpicture}[>={Stealth[length=2mm]}, every node/.style={font=\scriptsize}]
  \node[draw, rounded corners=2pt, align=center, minimum width=3.1cm, minimum height=0.8cm, fill=black!4] at (0.0,-0.0) {\textbf{davies\_m0p05}\\[-1pt]{\tiny speciate}};
  \node[draw, rounded corners=2pt, align=center, minimum width=3.1cm, minimum height=0.8cm, fill=black!4] at (3.6,-0.0) {\textbf{hmw\_m0p05}\\[-1pt]{\tiny speciate}};
  \node[draw, rounded corners=2pt, align=center, minimum width=3.1cm, minimum height=0.8cm, fill=black!4] at (7.2,-0.0) {\textbf{davies\_m0p1}\\[-1pt]{\tiny speciate}};
  \node[draw, rounded corners=2pt, align=center, minimum width=3.1cm, minimum height=0.8cm, fill=black!4] at (0.0,-1.3) {\textbf{hmw\_m0p1}\\[-1pt]{\tiny speciate}};
  \node[draw, rounded corners=2pt, align=center, minimum width=3.1cm, minimum height=0.8cm, fill=black!4] at (3.6,-1.3) {\textbf{davies\_m0p25}\\[-1pt]{\tiny speciate}};
  \node[draw, rounded corners=2pt, align=center, minimum width=3.1cm, minimum height=0.8cm, fill=black!4] at (7.2,-1.3) {\textbf{hmw\_m0p25}\\[-1pt]{\tiny speciate}};
  \node[draw, rounded corners=2pt, align=center, minimum width=3.1cm, minimum height=0.8cm, fill=black!4] at (0.0,-2.6) {\textbf{davies\_m0p5}\\[-1pt]{\tiny speciate}};
  \node[draw, rounded corners=2pt, align=center, minimum width=3.1cm, minimum height=0.8cm, fill=black!4] at (3.6,-2.6) {\textbf{hmw\_m0p5}\\[-1pt]{\tiny speciate}};
  \node[draw, rounded corners=2pt, align=center, minimum width=3.1cm, minimum height=0.8cm, fill=black!4] at (7.2,-2.6) {\textbf{davies\_m1p0}\\[-1pt]{\tiny speciate}};
  \node[draw, rounded corners=2pt, align=center, minimum width=3.1cm, minimum height=0.8cm, fill=black!4] at (0.0,-3.9) {\textbf{hmw\_m1p0}\\[-1pt]{\tiny speciate}};
  \node[draw, rounded corners=2pt, align=center, minimum width=3.1cm, minimum height=0.8cm, fill=black!4] at (3.6,-3.9) {\textbf{davies\_m2p0}\\[-1pt]{\tiny speciate}};
  \node[draw, rounded corners=2pt, align=center, minimum width=3.1cm, minimum height=0.8cm, fill=black!4] at (7.2,-3.9) {\textbf{hmw\_m2p0}\\[-1pt]{\tiny speciate}};
\end{tikzpicture}
\end{tutscheme}
\tutwhat{The Davies band edge, MEASURED: gamma\_pm and a\_w for NaCl under Davies beside PitzerHMW on identical brines, m = 0.05..6 mol/kg -- the cross-model curve behind every 'indicative beyond I \textasciitilde{} 0.5' advisory}
\begin{tutnarrative}
THE DAVIES BAND EDGE, measured instead of remembered.

Same NaCl brine, same totals, same solved pH -- speciated twice per rung: once under Davies (charge-only, the teaching rung; phi = 1 water activity) and once under PitzerHMW (per-ion virials; rigorous phi). What this measures is CROSS-MODEL divergence between two of Choupo's own engines -- weaker than agreement with an independent measurement, and said so; the Pitzer surface itself is anchored elsewhere (pitzer01\_nacl, archer01\_nacl\_cold\_to\_hot, pitzer\_gamma\_hamer\_wu).

WHY IT EXISTS: every electrolyte run above I \textasciitilde{} 0.5 mol/kg carries the advisory "Davies ... INDICATIVE"; the marcilla01/02 anchor contacts put numbers on what that costs at saturation (halite at 3.65 vs the measured 6.14 mol/kg). This case turns the advisory into a CURVE -- per rung, gamma\_Na, gamma\_Cl and a\_w under both closures, side by side. No threshold is invented: the divergence grows smoothly and there is no principled cut, so the table reports and the reader judges (the same posture as the neglected-dCp advisory in SolidPhase).

\end{tutnarrative}
\tutreadme{Same NaCl brine speciated twice per rung: Davies (teaching rung, charge-only, $\varphi$ = 1) beside PitzerHMW (per-ion virials), 298.15 K. Cross-model divergence between two in-tree engines --- weaker than an independent measurement and said so; the Pitzer surface itself is anchored elsewhere (pitzer01\_nacl, archer01\_nacl\_cold\_to\_hot, pitzer\_gamma\_hamer\_wu).}
\tutreadme{$\ast$ The "I \textasciitilde{} 0.5 trust band" every advisory quotes is real: the $\gamma$$\pm$ error crosses \textasciitilde{}2--8 \% over m = 0.25--0.5 and explodes beyond it (Davies' +0.3$\cdot$I term has no physics above the band --- it is the fit's tail). $\ast$ At m = 6, Davies reports SI\_halite = +1.10 where PitzerHMW reports -0.022 --- the rigorous rung sits essentially AT saturation beside Archer's peritectic (6.096 mol/kg), and its a\_w = 0.7592 beside the literature $\approx$ 0.753. This is the whole quantitative story of the marcilla01/02 anchor gaps (Davies saturating at 3.65 vs 6.14 mol/kg; the too-wet organic from a\_w $\approx$ 0.88): the model, not the machinery.}
\begin{tutgolden}
\goldrow{diag}{davies\_m0p05}{I}{0.0500001}
\goldrow{diag}{davies\_m0p05}{SI\_halite}{-4.34316}
\goldrow{diag}{davies\_m0p05}{aw}{0.9982}
\goldrow{diag}{davies\_m0p05}{gamma\_Cl}{0.821203}
\goldrow{diag}{davies\_m0p05}{gamma\_Na}{0.821203}
\goldrow{diag}{davies\_m0p05}{m\_Cl}{0.05}
\goldrow{diag}{davies\_m0p05}{m\_Na}{0.05}
\goldrow{diag}{davies\_m0p05}{pH}{7.00292}
\goldrow{diag}{hmw\_m0p05}{I}{0.0500001}
\goldrow{diag}{hmw\_m0p05}{SI\_halite}{-4.34351}
\goldrow{diag}{hmw\_m0p05}{aw}{0.998302}
\goldrow{diag}{hmw\_m0p05}{gamma\_Cl}{0.820867}
\goldrow{diag}{hmw\_m0p05}{gamma\_Na}{0.820867}
\goldrow{diag}{hmw\_m0p05}{m\_Cl}{0.05}
\goldmore{146}
\end{tutgolden}

\clearpage
\tutentry{edwards01\_sour\_water\_activity}{tutorials/props/electrolyte/edwards01\_sour\_water\_activity}
\begin{tutscheme}
\begin{tikzpicture}[>={Stealth[length=2mm]}, every node/.style={font=\scriptsize}]
  \node[draw, rounded corners=2pt, align=center, minimum width=3.1cm, minimum height=0.8cm, fill=black!4] at (0.0,-0.0) {\textbf{sourWater\_edwards}\\[-1pt]{\tiny speciate}};
  \node[draw, rounded corners=2pt, align=center, minimum width=3.1cm, minimum height=0.8cm, fill=black!4] at (3.6,-0.0) {\textbf{sourWater\_davies}\\[-1pt]{\tiny speciate}};
\end{tikzpicture}
\end{tutscheme}
\tutwhat{Edwards et al. (1978) activity model on a sour-water NH3/CO2 brine -- the molecular solutes get a real gamma from the model their parameters were regressed in, where Davies can only return 1.}
\begin{tutnarrative}
the WITNESS for the `edwardsPitzer` activity model.

WHAT THIS CASE IS FOR. Choupo already carries two aqueous activity models, and each fails this fluid in its own way:

`davies` is a charge-only correlation. Every singly-charged ion gets the same gamma and every NEUTRAL solute gets gamma = 1 -- exactly, by construction, at every ionic strength. A sour-water stripper's whole difficulty is the volatility of dissolved NH3 and CO2, i.e. of two NEUTRAL molecules, and Davies cannot say anything about either. `pitzerHMW` is the Harvie-Moller-Weare parameterisation: a DIFFERENT truncation, with a different parameter set (beta2, C-phi, theta, psi), fitted to strong electrolytes. Its numbers are not Edwards' numbers, and Edwards' numbers are not valid inside its equation.

`edwardsPitzer` is the truncation Edwards, Maurer, Newman \& Prausnitz (AIChE J. 24(6):966-976, 1978) REGRESSED HIS NH3/CO2/H2S/SO2 PARAMETERS IN: Pitzer's expansion cut after the second virial coefficient, extended to molecular solutes. Using his parameters anywhere else would be quoting his numbers under someone else's equation.

THE COMPOSITION is the sour-water region his Table 7 anchors: ammonia in large excess over carbon dioxide, the carbonate and carbamate networks both live, at an ionic strength around 1 mol/kg where a dilute correlation has no business being trusted.

WHAT TO READ IN THE OUTPUT. Run the two operations side by side and compare `gamma\_NH3aq` and `gamma\_CO2aq`:

under `davies` they are 1.000000 and 1.000000. Not approximately -- exactly, because the model's gamma is a function of charge alone and the charge is zero. A student who has only ever seen Davies may not realise this is a STRUCTURAL silence rather than a small correction. under `edwardsPitzer` they are 1.0224 and 1.0101: both SALTED OUT, NH3 rather more than CO2. Both rise because both Table 4 self-terms are positive at 25 C -- Eq 14 gives beta0(NH3,NH3) = 0.01522 and beta0(CO2,CO2) = 0.00822 kg/mol -- and NH3 rises further because it is the more abundant of the two by three orders of magnitude, so its own self-term dominates its Eq 8 sum.

THE SPECIATION MOVES TOO, and by far more than the gammas: the answer is I = 2.78 against Davies' 1.78, with m\_NH3aq 0.41 against 1.01. That is not a molecular effect at all -- it is the IONIC one. Davies at I \textasciitilde{} 2 returns gamma\_NH4 = 0.956, because its bracket saturates and then turns back up; Edwards' Eq 8 returns 0.301 at the same charge. The NH4+/NH3aq equilibrium reads those gammas, so the split follows. The lesson is the one the advisory prints on the Davies run: a dilute correlation used far outside its band does not merely lose accuracy, it changes the answer's SHAPE.

Neither number here is validated against measurement -- see below.

WHAT THIS CASE IS $\ast$NOT$\ast$. It is not the Table 7 anchor. Reproducing Table 7 needs the vapour side -- Henry's constants on the paper's own convention, the two-phase equilibrium, the total pressure -- and that slice is BUILT: `edwards02\_table7\_vle` (2026-08-23; record in docs/design/sour-water-stripper-scope.md section 6c). This case validates the ACTIVITY SURFACE alone, which is the piece the model owns.

WHAT IS AND IS NOT VALIDATED HERE. This case is a STRUCTURAL witness: it proves the model is wired, that its five parameter rules fire and are labelled, and that a neutral solute gets a gamma. It does NOT validate any number against data -- there is no measurement in it. The kernel itself is separately checked against two closed forms the paper supplies (the Eq 8a low-dissociation limit and the Debye-Huckel limiting law), and the measured-vs-computed comparison is Table 7 -- edwards02\_table7\_vle. A green run here means "the machine is right", never "the answer is right".

HONESTY ABOUT THE PARAMETERS. Edwards measured exactly one family: beta0 for LIKE molecule pairs (Eq 14, Table 4). Every other entry in the matrix Eq 8 sums over is ESTIMATED from those by his Eqs 21-24, and the paper calls both of its parameter tables "preliminary results subject to change as more and better data become available". The engine prints, at verbosity >= 3, WHICH RULE produced each number -- one measured family and four estimation routes -- so nothing here can be mistaken for a measured constant.

DATA SCOPE: SEALED (bin/choupo-import). The case carries its own species, reactions and EdwardsPitzer parameter records under constant/, each named and sha256-verified by constant/propertyManifest, and the runtime is FORBIDDEN the installation catalogue.

Sealing this case is where the importer's dependency closure was found not to know the EdwardsPitzer parameter home. The first seal kept 9 of the 28 pair parameters -- the Bronsted like-sign zeros, which need no record at all -- and the case still RAN, still converged, and answered with 19 terms missing: gamma\_HCO3 and gamma\_NH4 identical (the charge-only Debye-Huckel fallback) and both neutral gammas exactly 1, which is what a Davies run prints. Nothing refused. The importer now compares the staged sealed run against this case's own `expected` golden, because running is not agreeing.

\end{tutnarrative}
\begin{tutgolden}
\goldrow{diag}{sourWater\_edwards}{I}{2.77652}
\goldrow{diag}{sourWater\_edwards}{aw}{0.954591}
\goldrow{diag}{sourWater\_edwards}{gamma\_CO2aq}{1.01011}
\goldrow{diag}{sourWater\_edwards}{gamma\_CO3}{0.0181147}
\goldrow{diag}{sourWater\_edwards}{gamma\_HCO3}{0.227994}
\goldrow{diag}{sourWater\_edwards}{gamma\_NH3aq}{1.02239}
\goldrow{diag}{sourWater\_edwards}{gamma\_NH4}{0.30091}
\goldrow{diag}{sourWater\_edwards}{m\_CO2aq}{0.000484785}
\goldrow{diag}{sourWater\_edwards}{m\_CO3}{0.537903}
\goldrow{diag}{sourWater\_edwards}{m\_HCO3}{0.911612}
\goldrow{diag}{sourWater\_edwards}{m\_NH3aq}{0.41019}
\goldrow{diag}{sourWater\_edwards}{m\_NH4}{2.48981}
\goldrow{diag}{sourWater\_edwards}{pH}{9}
\goldrow{diag}{sourWater\_davies}{I}{1.77822}
\goldmore{16}
\end{tutgolden}

\clearpage
\tutentry{edwards02\_table7\_vle}{tutorials/props/electrolyte/edwards02\_table7\_vle}
\begin{tutscheme}
\begin{tikzpicture}[>={Stealth[length=2mm]}, every node/.style={font=\scriptsize}]
  \node[draw, rounded corners=2pt, align=center, minimum width=3.1cm, minimum height=0.8cm, fill=black!4] at (0.0,-0.0) {\textbf{table7\_row1}\\[-1pt]{\tiny speciate}};
  \node[draw, rounded corners=2pt, align=center, minimum width=3.1cm, minimum height=0.8cm, fill=black!4] at (3.6,-0.0) {\textbf{table7\_row2}\\[-1pt]{\tiny speciate}};
  \node[draw, rounded corners=2pt, align=center, minimum width=3.1cm, minimum height=0.8cm, fill=black!4] at (7.2,-0.0) {\textbf{table7\_row3}\\[-1pt]{\tiny speciate}};
  \node[draw, rounded corners=2pt, align=center, minimum width=3.1cm, minimum height=0.8cm, fill=black!4] at (0.0,-1.3) {\textbf{row1\_nocarbamate}\\[-1pt]{\tiny speciate}};
\end{tikzpicture}
\end{tutscheme}
\tutwhat{Edwards (1978) Table 7 anchor: NH3-CO2-H2O VLE at 100 degC -- the paper's own predictions reproduced from its own equations, the experiment quoted as context.}
\begin{tutnarrative}
the LITERATURE ANCHOR for the edwardsPitzer model: Table 7 of Edwards, Maurer, Newman \& Prausnitz, AIChE J. 24(6):966-976 (1978) -- NH3-CO2-H2O at 100 degC, three points.

edwards01 is the STRUCTURAL witness: it proves the activity surface is wired, and its own header says none of its numbers is validated against anything. This case is the validation it points at. For each Table 7 point the paper states the total molalities of ammonia and carbon dioxide in the LIQUID and predicts the speciation (molecular NH3 and CO2, ionic strength) and the VAPOUR above it (y\_NH3, y\_CO2, total pressure). The whole model stack here is the paper's own:

liquid activities Eq 8/10 (edwardsPitzer, the case surface) equilibria K(T) Eq 7 / Table 1 + Eq 20 (case-local chemistry/) Henry's constants Eq 13 / Table 3 (parameters/EdwardsPitzer/henry-$\ast$) pressure correction Eq 11 / Table 2 (vInfinity in the same records) vapour side Eq 5/6 (the speciate op's vapour\{\} block)

WHAT IS COMPARED, AND AGAINST WHAT. Table 7 carries TWO different columns and conflating them would be dishonest:

the paper's own PREDICTIONS (column I) -- a code-to-code comparison. Those are the golden values here, with ONE declared gap: the paper computes vapour-phase fugacity coefficients from a model its transcription does not carry, so this case declares `fugacity ideal;` (phi = 1) and the residual disagreement is expected to be of the size of that correction (a few percent at 2-3 atm), plus the numerical difference between two implementations of Eqs 7-24.

the EXPERIMENT (Otsuka et al. 1960) -- quoted below as CONTEXT, never as a target. The paper's OWN agreement with it is rough (first row: y\_NH3 0.066 measured vs 0.123 predicted -- a factor of two) and the paper says so plainly. A student reading this case sees a correlation reproduced closely and its real-world agreement stated honestly, which is the opposite of a tool that prints one number.

Table 7, the paper's column I (the golden targets here): row mNH3tot mCO2tot m\_NH3 m\_CO2 I y\_NH3 y\_CO2 P[atm] 1 2.90 1.45 1.245 0.014 1.50 0.123 0.494 2.51 2 3.71 1.14 2.279 0.0056 1.26 0.293 0.230 1.98 3 4.30 0.907 3.087 0.0030 1.01 0.407 0.116 1.97 The experiment (Otsuka 1960, context only): 1 0.066 0.506 3.15 2 0.274 0.202 2.08 3 0.355 0.095 2.00

THE CARBAMATE ENTERS THE CATALOGUE HERE (case-local): NH2COO- with Eq 20's equilibrium. Without it, an ammonia-rich carbonated water reports molecular molalities the paper's own predictions contradict -- it holds a substantial share of the dissolved CO2 at these conditions. Eq 20 is fitted 20-60 degC and used at 100 degC exactly as the paper uses it, with the extrapolation ANNOUNCED and the paper's own sensitivity argument in the record header (halving K20 moves P by under 2 \%).

DECLARED APPROXIMATIONS, each announced by the run: `fugacity ideal;` -- phi = 1 on all vapour species (see above); the solvent Poynting factor of Eq 6 is omitted (< 0.2 \% below 4 atm); Eq 20 extrapolated 60 $\to$ 100 degC (the paper's own move); water's Psat comes from the case catalogue's own curated Antoine, not from a value the paper does not print.

DATA SCOPE: sealed (bin/choupo-import); every record named and verified by constant/propertyManifest.

\end{tutnarrative}
\begin{tutgolden}
\goldrow{diag}{table7\_row1}{I}{1.547}
\goldrow{diag}{table7\_row1}{P\_atm}{2.51337}
\goldrow{diag}{table7\_row1}{aw}{0.946304}
\goldrow{diag}{table7\_row1}{chargeImposed}{1}
\goldrow{diag}{table7\_row1}{chargeResidual}{2.5541e-13}
\goldrow{diag}{table7\_row1}{gamma\_CO2aq}{0.919147}
\goldrow{diag}{table7\_row1}{gamma\_CO3}{0.022685}
\goldrow{diag}{table7\_row1}{gamma\_HCO3}{0.268386}
\goldrow{diag}{table7\_row1}{gamma\_NH2COO}{0.528413}
\goldrow{diag}{table7\_row1}{gamma\_NH3aq}{0.989831}
\goldrow{diag}{table7\_row1}{gamma\_NH4}{0.318905}
\goldrow{diag}{table7\_row1}{m\_CO2aq}{0.0141073}
\goldrow{diag}{table7\_row1}{m\_CO3}{0.0555284}
\goldrow{diag}{table7\_row1}{m\_HCO3}{1.09498}
\goldmore{99}
\end{tutgolden}

\clearpage
\tutentry{edwards03\_table8\_vle}{tutorials/props/electrolyte/edwards03\_table8\_vle}
\begin{tutscheme}
\begin{tikzpicture}[>={Stealth[length=2mm]}, every node/.style={font=\scriptsize}]
  \node[draw, rounded corners=2pt, align=center, minimum width=3.1cm, minimum height=0.8cm, fill=black!4] at (0.0,-0.0) {\textbf{table8\_row1}\\[-1pt]{\tiny speciate}};
  \node[draw, rounded corners=2pt, align=center, minimum width=3.1cm, minimum height=0.8cm, fill=black!4] at (3.6,-0.0) {\textbf{table8\_row2}\\[-1pt]{\tiny speciate}};
  \node[draw, rounded corners=2pt, align=center, minimum width=3.1cm, minimum height=0.8cm, fill=black!4] at (7.2,-0.0) {\textbf{table8\_row3}\\[-1pt]{\tiny speciate}};
\end{tikzpicture}
\end{tutscheme}
\tutwhat{Edwards (1978) Table 8 anchor: NH3-H2S-H2O VLE at 80 degC -- the paper's own predictions reproduced from its own equations, the experiment quoted as context.}
\begin{tutnarrative}
the SECOND literature anchor for the Edwards (1978) model stack: Table 8 of the paper, ammonia + hydrogen sulfide in water at 80 degC, against the paper's OWN prediction column I.

WHY A SECOND ANCHOR. edwards02 (Table 7) pins the NH3-CO2 system, where the difficulty is the carbonate/carbamate network. Table 8 exercises what Table 7 cannot: the SULFIDE family (the mi-H2Saq-$\ast$ molecule-ion records curated 2026-08-24 from the paper's Table 6, including its two TABULATED ZEROS), a different temperature (80 degC, so every T-polynomial is read off its 25 degC calibration point), and ionic strengths up to 7.8 mol/kg -- beyond the 6-molal Pitzer validity the paper itself quotes. Rows 2 and 3 are code-to-code targets and NOT physical claims; the paper runs them and says so ("if ionization is extensive, the upper limit is necessarily lower").

WHAT IS COMPARED, AND AGAINST WHAT. Exactly the edwards02 posture:

the paper's own PREDICTIONS (column I: all beta0 from Tables 4 and 6, the configuration these curated records carry) are the anchor targets; the EXPERIMENT (Miles \& Wilson 1975) is CONTEXT, never a target. The paper's own worst multisolute disagreement lives here: rows 2 and 3 measure 16.50 and 13.20 atm against predictions of 12.91 and 12.64 atm (22-28 \%), and its Table 9 sensitivity analysis shows the liquid-composition error needed to close such gaps.

THE CHEMISTRY IS THE PAPER'S OWN: case-local constant/chemistry/ carries Table 1 (H2O, NH3 and H2S rows) in the formation direction, each header recording the derivation. The H2S formation logK25 is 7.0033 against the independently curated PHREEQC 6.994 -- two primaries 0.009 apart, a finding, not an error.

DECLARED GAPS, the same list as edwards02 plus none: `fugacity ideal;` -- the paper's own vapour-phi model is not transcribed; the residual is expected to carry that correction (largest on rows 2-3 at 13-16 atm); the solvent Poynting factor of Eq 6 is omitted; water's Psat comes from the case catalogue's own curated Antoine.

DATA SCOPE: sealed (bin/choupo-import); every record named and verified by constant/propertyManifest.

\end{tutnarrative}
\begin{tutgolden}
\goldrow{diag}{table8\_row1}{I}{1.13632}
\goldrow{diag}{table8\_row1}{P\_atm}{1.10768}
\goldrow{diag}{table8\_row1}{aw}{0.896601}
\goldrow{diag}{table8\_row1}{chargeImposed}{1}
\goldrow{diag}{table8\_row1}{chargeResidual}{-4.35684e-17}
\goldrow{diag}{table8\_row1}{gamma\_H2Saq}{1.08209}
\goldrow{diag}{table8\_row1}{gamma\_HS}{0.516824}
\goldrow{diag}{table8\_row1}{gamma\_NH3aq}{1.06377}
\goldrow{diag}{table8\_row1}{gamma\_NH4}{0.531883}
\goldrow{diag}{table8\_row1}{m\_H2Saq}{0.0038303}
\goldrow{diag}{table8\_row1}{m\_HS}{1.13617}
\goldrow{diag}{table8\_row1}{m\_NH3aq}{3.97368}
\goldrow{diag}{table8\_row1}{m\_NH4}{1.13632}
\goldrow{diag}{table8\_row1}{pH}{8.67302}
\goldmore{64}
\end{tutgolden}

\clearpage
\tutentry{enrtl\_mixed\_nacl\_ethanol\_esteso}{tutorials/props/electrolyte/enrtl\_mixed\_nacl\_ethanol\_esteso}
\begin{tutscheme}
\begin{tikzpicture}[>={Stealth[length=2mm]}, every node/.style={font=\scriptsize}]
  \node[draw, rounded corners=2pt, align=center, minimum width=3.1cm, minimum height=0.8cm, fill=black!4] at (0.0,-0.0) {\textbf{esteso}\\[-1pt]{\tiny enrtlMixedSolvent}};
\end{tikzpicture}
\end{tutscheme}
\tutwhat{VALIDATION of the generalized segment-based eNRTL (Chen \& Song 2004, AIChE J 50:1928, DOI 10.1002/aic.10151) against the PRIMARY measured gamma\_pm of NaCl in ethanol-water at 25 C (Esteso et al., J. Solution Chem. 18 (1989) 277, emf cell, 0-90 wt\% ethanol). ZERO refit: every tau/X:Z is the published Table 2/4/7 value; A\_phi per composition is Esteso's measured Table II. Headline AAD \textasciitilde{}4.3\% over 47 points (paper quotes RMS 0.026 on gamma); \textasciitilde{}1\% in pure water, \textasciitilde{}1.2\% at 70 wt\%. The +8\% at 80/90 wt\% is DECLARED physics the complete-dissociation model cannot carry: Esteso measures ion association K\_A = 30-46 there (his Table V).}
\begin{tutnarrative}
can eNRTL predict a salt's activity in the antisolvent mixture?

The SAME physics the antisolvent crystalliser leans on -- gamma\_pm of NaCl as ethanol pushes the water out -- validated against the PRIMARY dataset:

MODEL Chen \& Song, "Generalized electrolyte-NRTL model for mixed-solvent electrolyte systems", AIChE J 50(8):1928-1941 (2004), DOI 10.1002/aic.10151. Ethanol = 1.811 (-C2H4-) + 0.609 (-OH) segments (Table 4); ion<$\to$OH taus = the aqueous NaCl pair (Table 2); ion<$\to$C2H4 taus from Table 7. ZERO refit here. DATA Esteso, Gonzalez-Diaz, Hernandez-Luis \& Fernandez-Merida, "Activity coefficients for NaCl in ethanol-water mixtures at 25 C", J. Solution Chem. 18(3):277-288 (1989). Table I (emf, molal scale, per-solvent reference). A\_phi per composition = Esteso Table II (measured dielectric + density). 0 wt\% gate: Robinson \& Stokes, Electrolyte Solutions, 2nd ed. (1959), NaCl appendix table; A\_phi(water, 25 C) = 0.391.

WHAT THE NUMBERS SAY (locked by the golden): total AAD \textasciitilde{}4.3\% over 47 points; pure water \textasciitilde{}1.0\%; 70 wt\% ethanol \textasciitilde{}1.2\%. 80/90 wt\% sit +8\% systematically -- Esteso's own chemical-model fit finds ion-pair association K\_A = 30-46 there (Table V), physics a complete- dissociation model cannot represent. Declared, not tuned away.

\end{tutnarrative}
\begin{tutgolden}
\goldrow{diag}{esteso}{AAD\_total\_pct}{4.30645321132}
\goldrow{diag}{esteso}{AAD\_wt0\_pct}{1.0101269464}
\goldrow{diag}{esteso}{AAD\_wt20\_pct}{5.27484474055}
\goldrow{diag}{esteso}{AAD\_wt40\_pct}{5.27091506549}
\goldrow{diag}{esteso}{AAD\_wt60\_pct}{2.69902360213}
\goldrow{diag}{esteso}{AAD\_wt70\_pct}{1.18874152393}
\goldrow{diag}{esteso}{AAD\_wt80\_pct}{7.80889140849}
\goldrow{diag}{esteso}{AAD\_wt90\_pct}{7.78650421238}
\goldrow{diag}{esteso}{n\_points}{47}
\end{tutgolden}

\clearpage
\tutentry{enrtl\_multisalt\_nacl\_kcl}{tutorials/props/electrolyte/enrtl\_multisalt\_nacl\_kcl}
\begin{tutscheme}
\begin{tikzpicture}[>={Stealth[length=2mm]}, every node/.style={font=\scriptsize}]
  \node[draw, rounded corners=2pt, align=center, minimum width=3.1cm, minimum height=0.8cm, fill=black!4] at (0.0,-0.0) {\textbf{harned\_nacl\_kcl}\\[-1pt]{\tiny enrtlMultiSalt}};
  \node[draw, rounded corners=2pt, align=center, minimum width=3.1cm, minimum height=0.8cm, fill=black!4] at (3.6,-0.0) {\textbf{harned\_nacl\_kcl\_refined}\\[-1pt]{\tiny enrtlMultiSalt}};
\end{tikzpicture}
\end{tutscheme}
\tutwhat{Multi-salt symmetric eNRTL (Song \& Chen 2009) -- NaCl+KCl Harned sweep at 4 mol/kg}
\begin{tutnarrative}
THE MULTI-SALT eNRTL ARRIVES: the multicomponent symmetric electrolyte NRTL of Song \& Chen 2009 (Ind. Eng. Chem. Res. 48, 7788, DOI 10.1021/ie9004578) on the classic Harned experiment -- NaCl + KCl sharing water at a CONSTANT 4 mol/kg total molality, gamma\_pm of BOTH salts as KCl replaces NaCl. This is the paper's own demonstration system (its Figures 1-2).

What the run shows (all golden-locked): gamma\_pm(NaCl) slides from its pure-solution value to its TRACE value in 4 m KCl, and vice versa -- both ln gamma curves are NEARLY LINEAR in the salt fraction: Harned's rule, EMERGING from the model rather than being assumed (R\textasciicircum{}2 > 0.99 with zero electrolyte-electrolyte parameters). The PURE endpoints reproduce the validated single-salt eNRTL kernel exactly (the reduction pin \textasciitilde{}1e-15): the multicomponent formulation collapses correctly. The pair parameters are the EXACT Chen \& Evans 1986 Table 1 values the paper uses (case-local constant/electrolyte/enrtl.dat, cited). The standards K-Cl.dat carries an INDEPENDENT refit to Hamer \& Wu 1972 that lands within 0.2 \% of the 1986 values -- two roads, one answer. The 1986 paper itself (Table 2) validates tauEE = 0 for NaCl-KCl: sigma = 0.019 \% on the Covington 1968 mixture osmotic data with the salt-salt energies set to zero -- this case's default is the canonical choice. HONESTY, printed by the op: the published multi-electrolyte eNRTL activity-coefficient expressions are NOT exactly Gibbs-Duhem consistent (the composition-dependent mixing rules are differentiated as constants). The op's FD oracle quantifies that inconsistency at the mid-sweep state (gibbsDuhem\_dev\_mid) next to the pure-salt state where the same oracle sits at \textasciitilde{}1e-9. THE REFINED FORMULATION SHIPS TOO (second op below): gamma as the EXACT derivative of the same Gex -- the chain rule kept through the composition-dependent mixing rules (the Bollas, Chen \& Barton 2008 idea on the symmetric form). Its oracle closes (\textasciitilde{}1e-9) at every composition; both reduce EXACTLY to the same single-salt kernel at the endpoints. Mid-sweep they differ by |dln gamma\_pm| \textasciitilde{} 0.041 (\textasciitilde{}4 \% in gamma), and the trace limits differ visibly (NaCl trace 0.617 published vs 0.678 consistent). WHICH is closer to nature is a question for MIXTURE data (Covington 1968 / Robinson 1961) -- this case states the facts and locks both; it does not pick a winner without measurements.

tauEE (the same-anion electrolyte-electrolyte parameter tau\_NaCl,KCl of the paper's Figures 1-2) defaults to 0; set it to +-1 to reproduce the paper's sensitivity fan.

THE EXPERIMENTAL VERDICT (Robinson 1961, J. Phys. Chem. 65, 662, DOI 10.1021/j100822a016 -- isopiestic NaCl-KCl mixtures, Harned coefficients at M \textasciitilde{} 4.4 mol/kg: alpha\_NaCl = +0.0243, alpha\_KCl = -0.0091; the op now emits harned\_alpha\_1/2 on the same convention, golden-locked): alpha\_NaCl: published +0.0285 (+17 \%) consistent +0.0180 (-26 \%) alpha\_KCl: published -0.0256 (2.8x) consistent -0.0187 (2.1x) With tauEE = 0, PUBLISHED wins on the dominant NaCl coefficient; BOTH formulations overstate the (small) effect on KCl.

CROSS-CHECK vs THE REFINED-eNRTL PAPER (Bollas, Chen \& Barton 2008, AIChE J 54, 1608, DOI 10.1002/aic.11485 -- their Figure 1 is THIS system, 4 m, same 1986 parameters, Robinson's points overlaid): their Fig 1a (original eNRTL) matches our `published` to the pixel (NaCl solids at x=1 for tau = 0/0.25/0.50: ours -0.216/-0.186/-0.157 vs read -0.215/-0.185/-0.15; KCl dashed at x=0: -0.144/-0.183 vs -0.145/-0.185); their Fig 1b (refined) shows the SIGNATURE our `consistent` shows: the tau fan nearly COLLAPSES (exact derivatives are almost insensitive to tau\_EE: ours moves only -0.176 $\to$ -0.184 for tau 0 $\to$ 0.5, where published moves -0.216 $\to$ -0.157), and the KCl trace lands at \textasciitilde{}-0.16 (ours -0.163). Their refinement differentiates the 1986 UNSYMMETRIC Gex; ours differentiates the 2009 SYMMETRIC one -- same idea, sibling surfaces, same qualitative physics, near-identical aqueous numbers. A single fitted tauEE cannot close both (best compromise tau \textasciitilde{} 0.29 trades the NaCl error up to -23 \% to halve the KCl one) -- the residual is structural, not a missing parameter. Note the consistency of the record: Chen \& Evans 1986 validated tauEE = 0 against mixture OSMOTIC data (sigma 0.019 \%) -- the solvent sees the mixture mean, where the individual-gamma asymmetries largely cancel; Robinson's per-salt gammas are where they show.

FIGURE-1 REPRODUCTION -- TO THE PIXEL (checked 2026-07-02, figure read at 200 dpi): the paper's y-axis says "ln gamma" but the numbers are LOG10 (log10 of the pure-NaCl-4m gamma is -0.106 = the figure's -0.105; ln would be -0.24, off the plotted range). With the 1986 parameters this kernel reproduces the full tau fan to <= 0.001 in log10: x\_KCl=0 (all curves): -0.106 (figure \textasciitilde{} -0.105) x\_KCl=1: tau=0 $\to$ -0.216 (figure \textasciitilde{} -0.215) tau=+1 $\to$ flat, dip to -0.101, back to -0.106 (the figure's circle-marked curve shows the same shallow dip) tau=-1 $\to$ -0.364 (figure \textasciitilde{} -0.365)

THE PARAGRAPH ABOVE IS A DATED READING, NOT THE LIVE STATE (2026-09-03). The golden this case ships today pins gamma1\_pure = 0.801 for pure NaCl at 4 m, i.e. log10 = -0.096, where the reading of 2026-07-02 recorded -0.106 against the figure's -0.105: the kernel's answer has moved by \textasciitilde{}0.01 in log10 since the figure was read (the 2026-08-29 fleet audit found the drift). Whether the CURRENT curves still sit on the figure can only be settled by reading the figure again, and that has not been done -- so "to the pixel" is what was true on 2026-07-02, kept as the record of that check, and is not a claim about this golden. Nothing here was adjusted toward the figure.

\end{tutnarrative}
\begin{tutgolden}
\goldrow{diag}{harned\_nacl\_kcl}{consistent}{0}
\goldrow{diag}{harned\_nacl\_kcl}{formulation\_gap\_mid}{0.0413132048052}
\goldrow{diag}{harned\_nacl\_kcl}{gamma1\_pure}{0.801115128273}
\goldrow{diag}{harned\_nacl\_kcl}{gamma1\_trace}{0.616514114921}
\goldrow{diag}{harned\_nacl\_kcl}{gamma2\_pure}{0.575990027563}
\goldrow{diag}{harned\_nacl\_kcl}{gamma2\_trace}{0.728881856731}
\goldrow{diag}{harned\_nacl\_kcl}{gibbsDuhem\_dev\_mid}{0.471752659198}
\goldrow{diag}{harned\_nacl\_kcl}{harned\_R2\_min}{0.998954550075}
\goldrow{diag}{harned\_nacl\_kcl}{harned\_alpha\_1}{0.0284936089821}
\goldrow{diag}{harned\_nacl\_kcl}{harned\_alpha\_2}{-0.0256162918946}
\goldrow{diag}{harned\_nacl\_kcl}{harned\_slope\_1}{-0.262435837152}
\goldrow{diag}{harned\_nacl\_kcl}{harned\_slope\_2}{-0.235934767418}
\goldrow{diag}{harned\_nacl\_kcl}{oracle\_pure\_dev}{4.03345401523e-09}
\goldrow{diag}{harned\_nacl\_kcl}{reduction\_pin\_dev}{1.11022302463e-15}
\goldmore{18}
\end{tutgolden}

\clearpage
\tutentry{enthalpy\_naoh\_water}{tutorials/props/electrolyte/enthalpy\_naoh\_water}
\begin{tutscheme}
\begin{tikzpicture}[>={Stealth[length=2mm]}, every node/.style={font=\scriptsize}]
  \node[draw, rounded corners=2pt, align=center, minimum width=3.1cm, minimum height=0.8cm, fill=black!4] at (0.0,-0.0) {\textbf{gamma\_naoh}\\[-1pt]{\tiny pitzerActivity}};
\end{tikzpicture}
\end{tutscheme}
\tutwhat{NaOH-water heat of dilution: Pitzer with calorimetric T-slots vs Parker (1965)}
\begin{tutnarrative}
the calorimetric home of the electrolyte-enthalpy build (docs/electrolyte-enthalpy-spec.md).

constant/experiments/phiL\_25C\_parker1965.dat carries the MEASURED Phi\_L (Parker NSRDS-NBS 2, 1965) used by the curation refit (bin/curate/fit\_lphi\_pitzer.py) that fitted the Na-OH dbeta/dT T-slots in the standard catalogue (the FIT's own rel-AAD 4.97\%, gate <5\%; the in-engine L\_phi scan pinned by this case lands at rel-AAD 4.77\% -- two measurements, not one). The op below sweeps the ACTIVITY surface whose T-derivative is L\_phi; the in-engine L\_phi scan + AAD overlay land with slices 4-5 (the 3-curve chart).

\end{tutnarrative}
\begin{tutgolden}
\goldrow{diag}{gamma\_naoh}{Lphi\_AAD\_abs\_Jmol}{22.9776}
\goldrow{diag}{gamma\_naoh}{Lphi\_AAD\_rel\_pct}{4.7712}
\goldrow{diag}{gamma\_naoh}{Lphi\_extrap\_maxdev\_Jmol}{15774.0028}
\goldrow{diag}{gamma\_naoh}{Lphi\_nMeas}{29.0000}
\goldrow{diag}{gamma\_naoh}{aw\_1m}{0.9665}
\goldrow{diag}{gamma\_naoh}{gamma\_pm\_1m}{0.6679}
\goldrow{diag}{gamma\_naoh}{gamma\_pm\_3m}{0.7838}
\goldrow{diag}{gamma\_naoh}{phi\_1m}{0.9471}
\end{tutgolden}

\clearpage
\tutentry{farelo\_licl\_range}{tutorials/props/electrolyte/farelo\_licl\_range}
\begin{tutscheme}
\begin{tikzpicture}[>={Stealth[length=2mm]}, every node/.style={font=\scriptsize}]
  \node[draw, rounded corners=2pt, align=center, minimum width=3.1cm, minimum height=0.8cm, fill=black!4] at (0.0,-0.0) {\textbf{licl\_1.0}\\[-1pt]{\tiny speciate}};
  \node[draw, rounded corners=2pt, align=center, minimum width=3.1cm, minimum height=0.8cm, fill=black!4] at (3.6,-0.0) {\textbf{licl\_6.0}\\[-1pt]{\tiny speciate}};
  \node[draw, rounded corners=2pt, align=center, minimum width=3.1cm, minimum height=0.8cm, fill=black!4] at (7.2,-0.0) {\textbf{licl\_13.0}\\[-1pt]{\tiny speciate}};
  \node[draw, rounded corners=2pt, align=center, minimum width=3.1cm, minimum height=0.8cm, fill=black!4] at (0.0,-1.3) {\textbf{licl\_19.7}\\[-1pt]{\tiny speciate}};
  \node[draw, rounded corners=2pt, align=center, minimum width=3.1cm, minimum height=0.8cm, fill=black!4] at (3.6,-1.3) {\textbf{farelo\_LiCl\_sat}\\[-1pt]{\tiny speciate}};
\end{tikzpicture}
\end{tutscheme}
\tutwhat{LiCl Pitzer REFITTED to Hamer-Wu + Farelo saturation: gamma\_pm to \textasciitilde{}11\% over 0.5-19 mol/kg (was 774 at 19.7m), and the LiCl.H2O hydrate Ksp reproduces Farelo's 19.7 mol/kg via the HMW osmotic water activity (a\_w=0.11).}
\tutreadme{Prof. F?tima Farelo measured LiCl solubility at 19.7 $\to$ 23.5 mol/kg (solid: LiCl$\cdot$H$_2$O) --- $\ast$J. Chem. Eng. Data$\ast$ 50 (2005) 1470, DOI 10.1021/je050111j. That is one of the most concentrated aqueous electrolytes there is.}
\tutreadme{The catalogue's earlier Li--Cl Pitzer parameters (Appelo, Parkhurst \& Post, $\ast$Geochim. Cosmochim. Acta$\ast$ 125 (2014) 49), fitted below \textasciitilde{}6 mol/kg, blow up at high concentration --- $\gamma$$\pm$ = 774 at 19.7 mol/kg vs \textasciitilde{}45 measured. The two-parameter B-term is linear in $\ast$m$\ast$, so $\gamma$ grows exponentially.}
\tutreadme{Rather than declare defeat, the parameters were refitted to Hamer \& Wu (1972) $\gamma$$\pm$ and Farelo's saturation. A negative C$\varphi$ (-0.0026) tames the high-$\ast$m$\ast$ blow-up, and the three-parameter form then reaches saturation:}
\tutreadme{carries a water-activity leg. At 19.7 mol/kg the model gives a\_w = 0.11 (not 1!) --- which only the rigorous Pitzer-HMW osmotic coefficient delivers, not the $\varphi$ = 1 dilute shortcut. The lithiumChlorideH2O mineral Ksp (logK$_2$$_5$ = 4.984) is calibrated to Farelo's saturation $\ast$through$\ast$ that a\_w. The two pieces meet: the refitted $\gamma$ and the osmotic a\_w together reproduce the measured solubility.}
\tutreadme{- electrolyte/pairs.dat: the Li--Cl Pitzer pair, refitted to Hamer--Wu + Farelo, valid to \textasciitilde{}19 mol/kg (supersedes the Appelo low-$\ast$m$\ast$ fit for concentrated LiCl). - electrolyte/minerals.dat: the LiCl$\cdot$H$_2$O hydrate Ksp, calibrated to Farelo. - constant/species/Li.dat: Li? aqueous formation datum (hfAq = -278.49 kJ/mol, CODATA).}
\tutreadme{- The full measured dataset (this and the other five ternaries) is tabulated, cited, in constant/experimental/farelo\_solubility.csv. - The equilibrate round-trip (precipitate LiCl$\cdot$H$_2$O from a supersaturated feed) diverges at I $\approx$ 20 --- a known active-set / $\gamma$ fixed-point numerical limit at extreme ionic strength, $\ast$independent of the activity model$\ast$. The SI = 0 calibration at the saturation composition is exact; reaching it by precipitation is the convergence problem, not the thermodynamics. - The AlCl$_3$ systems from the same paper remain deferred: Al$^3$? hydrolyses (the paper notes pH < 2.75) and needs an Al speciation model Choupo does not yet carry.}
\begin{tutgolden}
\goldrow{diag}{farelo\_LiCl\_sat}{I}{19.7}
\goldrow{diag}{farelo\_LiCl\_sat}{SI\_lithiumChlorideH2O}{1.18638e-05}
\goldrow{diag}{farelo\_LiCl\_sat}{aw}{0.113466}
\goldrow{diag}{farelo\_LiCl\_sat}{chargeImposed}{1}
\goldrow{diag}{farelo\_LiCl\_sat}{chargeResidual}{5.04184e-14}
\goldrow{diag}{farelo\_LiCl\_sat}{gamma\_Cl}{46.7905}
\goldrow{diag}{farelo\_LiCl\_sat}{gamma\_Li}{46.7905}
\goldrow{diag}{farelo\_LiCl\_sat}{m\_Cl}{19.7}
\goldrow{diag}{farelo\_LiCl\_sat}{m\_Li}{19.7}
\goldrow{diag}{farelo\_LiCl\_sat}{pH}{7.99525}
\end{tutgolden}

\clearpage
\tutentry{farelo\_nacl\_nh4cl}{tutorials/props/electrolyte/farelo\_nacl\_nh4cl}
\begin{tutscheme}
\begin{tikzpicture}[>={Stealth[length=2mm]}, every node/.style={font=\scriptsize}]
  \node[draw, rounded corners=2pt, align=center, minimum width=3.1cm, minimum height=0.8cm, fill=black!4] at (0.0,-0.0) {\textbf{farelo\_NaCl\_pure}\\[-1pt]{\tiny speciate}};
  \node[draw, rounded corners=2pt, align=center, minimum width=3.1cm, minimum height=0.8cm, fill=black!4] at (3.6,-0.0) {\textbf{farelo\_NaCl\_nh4\_093}\\[-1pt]{\tiny speciate}};
  \node[draw, rounded corners=2pt, align=center, minimum width=3.1cm, minimum height=0.8cm, fill=black!4] at (7.2,-0.0) {\textbf{farelo\_NaCl\_nh4\_187}\\[-1pt]{\tiny speciate}};
  \node[draw, rounded corners=2pt, align=center, minimum width=3.1cm, minimum height=0.8cm, fill=black!4] at (0.0,-1.3) {\textbf{farelo\_NaCl\_nh4\_280}\\[-1pt]{\tiny speciate}};
  \node[draw, rounded corners=2pt, align=center, minimum width=3.1cm, minimum height=0.8cm, fill=black!4] at (3.6,-1.3) {\textbf{farelo\_NH4Cl\_pure}\\[-1pt]{\tiny speciate}};
\end{tikzpicture}
\end{tutscheme}
\tutwhat{Pitzer-HMW vs Farelo (JCED 2005): NaCl+NH4Cl common-ion salting-out -- model SI \textasciitilde{} 0 at each measured saturation point (NaCl 6.15$\to$5.5$\to$5.0; NH4Cl 7.35 mol/kg, 298 K).}
\tutreadme{A glass-box validation of Choupo's multi-ion Pitzer electrolyte model against the solubility measurements of Prof. F?tima Farelo (Chemical Engineering Dept., IST --- T?cnico Lisboa).}
\tutreadme{M. Farelo, C. Fernandes \& A. Avelino, $\ast$"Solubilities for Six Ternary Systems: NaCl + NH$_4$Cl + H$_2$O, KCl + NH$_4$Cl + H$_2$O, NaCl + LiCl + H$_2$O, KCl + LiCl + H$_2$O, NaCl + AlCl$_3$ + H$_2$O, and KCl + AlCl$_3$ + H$_2$O at T = (298 to 333) K"$\ast$, J. Chem. Eng. Data 50 (2005) 1470--1477. DOI [10.1021/je050111j](https://doi.org/10.1021/je050111j).}
\tutreadme{Saturation molalities measured to $\pm$0.0007 mol/kg (visual polythermal / isothermal methods). The pure-salt anchors at 298.15 K: NaCl 6.15, NH$_4$Cl 7.35 mol/kg; the salting-out series (Table 1): NaCl solubility falls 6.15 $\to$ 5.5 $\to$ 5.0 $\to$ 4.95 as NH$_4$Cl rises 0 $\to$ 0.93 $\to$ 1.87 $\to$ 2.80 mol/kg.}
\tutreadme{The model is calibrated to only the two pure-salt solubilities --- halite's Ksp from PHREEQC, sal-ammoniac's Ksp pinned to Farelo's NH$_4$Cl point. Everything else is a prediction. The NH$_4$--Cl Pitzer interaction parameters come from the primary single-electrolyte fit (Pitzer \& Mayorga 1973); the catalogue already carried Na--Cl and K--Cl.}
\tutreadme{At each of Farelo's measured saturation compositions the model's saturation index}
\tutreadme{should be $\approx$ 0 if the model reproduces the measurement. It is, to $\pm$0.03 log units:}
\begin{tutgolden}
\goldrow{diag}{farelo\_NaCl\_pure}{I}{6.15}
\goldrow{diag}{farelo\_NaCl\_pure}{SI\_halite}{0.0168246}
\goldrow{diag}{farelo\_NaCl\_pure}{aw}{0.751868}
\goldrow{diag}{farelo\_NaCl\_pure}{gamma\_Cl}{1.0105}
\goldrow{diag}{farelo\_NaCl\_pure}{gamma\_Na}{1.0105}
\goldrow{diag}{farelo\_NaCl\_pure}{m\_Cl}{6.15}
\goldrow{diag}{farelo\_NaCl\_pure}{m\_Na}{6.15}
\goldrow{diag}{farelo\_NaCl\_pure}{pH}{7.11704}
\goldrow{diag}{farelo\_NaCl\_nh4\_093}{I}{6.42998}
\goldrow{diag}{farelo\_NaCl\_nh4\_093}{SI\_halite}{-0.022239}
\goldrow{diag}{farelo\_NaCl\_nh4\_093}{SI\_salammoniac}{-0.70594}
\goldrow{diag}{farelo\_NaCl\_nh4\_093}{aw}{0.746427}
\goldrow{diag}{farelo\_NaCl\_nh4\_093}{gamma\_NH4}{0.591469}
\goldrow{diag}{farelo\_NaCl\_nh4\_093}{gamma\_Na}{1.04074}
\goldmore{39}
\end{tutgolden}

\clearpage
\tutentry{fpd01\_nacl\_freezing}{tutorials/props/electrolyte/fpd01\_nacl\_freezing}
\begin{tutscheme}
\begin{tikzpicture}[>={Stealth[length=2mm]}, every node/.style={font=\scriptsize}]
  \node[draw, rounded corners=2pt, align=center, minimum width=3.1cm, minimum height=0.8cm, fill=black!4] at (0.0,-0.0) {\textbf{fpd\_nacl}\\[-1pt]{\tiny freezingPoint}};
\end{tikzpicture}
\end{tutscheme}
\tutwhat{NaCl freezing-point depression -- ice as a SolidPhase, the crystal's first case}
\begin{tutnarrative}
THE ICE WITNESS. The freezing-point curve of NaCl(aq), solved as chemical-potential equality between an ACTUAL SolidPhase (the crystallising mode, its first case-reachable consumer) and the Pitzer surface's water activity:

f\_solid(T) = a\_w(m, T) $\ast$ Psat(T)

Until this case the crystal was exercised only by check\_ice\_freezing's probe; here its refusals and announced approximations (the neglected liquid/solid heat-capacity term, the sub-freezing Psat extrapolation) all fire on a real run. The dilute-limit KPI closes the K\_f loop: the derived cryoscopic constant against the slope this same phase produces. ANCHOR HONESTY: the validation dataset carries ONE measured point with an INTERIM citation (see its header) -- a thin anchor, said plainly, and the enrichment path (a primary FPD table read) is on the curation ledger.

\end{tutnarrative}
\begin{tutgolden}
\goldrow{diag}{fpd\_nacl}{Tf\_lastM\_K}{266.325}
\goldrow{diag}{fpd\_nacl}{depression\_lastM\_K}{6.83528}
\goldrow{diag}{fpd\_nacl}{slope\_dilute\_K\_kg\_mol}{3.46546}
\goldrow{diag}{fpd\_nacl}{slope\_over\_nuKf}{0.931447}
\goldrow{diag}{fpd\_nacl}{theta\_AAD\_pct}{0.0460433}
\goldrow{diag}{fpd\_nacl}{theta\_nMeas}{1}
\end{tutgolden}

\clearpage
\tutentry{ksp\_temperature}{tutorials/props/electrolyte/ksp\_temperature}
\begin{tutscheme}
\begin{tikzpicture}[>={Stealth[length=2mm]}, every node/.style={font=\scriptsize}]
  \node[draw, rounded corners=2pt, align=center, minimum width=3.1cm, minimum height=0.8cm, fill=black!4] at (0.0,-0.0) {\textbf{T10}\\[-1pt]{\tiny speciate}};
  \node[draw, rounded corners=2pt, align=center, minimum width=3.1cm, minimum height=0.8cm, fill=black!4] at (3.6,-0.0) {\textbf{T25}\\[-1pt]{\tiny speciate}};
  \node[draw, rounded corners=2pt, align=center, minimum width=3.1cm, minimum height=0.8cm, fill=black!4] at (7.2,-0.0) {\textbf{T40}\\[-1pt]{\tiny speciate}};
  \node[draw, rounded corners=2pt, align=center, minimum width=3.1cm, minimum height=0.8cm, fill=black!4] at (0.0,-1.3) {\textbf{T60}\\[-1pt]{\tiny speciate}};
  \node[draw, rounded corners=2pt, align=center, minimum width=3.1cm, minimum height=0.8cm, fill=black!4] at (3.6,-1.3) {\textbf{T80}\\[-1pt]{\tiny speciate}};
\end{tikzpicture}
\end{tutscheme}
\tutwhat{K(T) MATTERS: a temperature scan of mineral saturation for ONE fixed water (a hard, carbonate-bearing groundwater at pH 7.6), speciated at 10, 25, 40, 60 and 80 C. The point is the SIGN of dK/dT. CALCITE dissolution is EXOTHERMIC, so its equilibrium constant FALLS as temperature rises and its saturation index SI\_calcite RISES -- RETROGRADE SOLUBILITY: calcite is LESS soluble hot. This is why hot heat-exchanger surfaces and warm RO concentrate scale with calcium carbonate even when the cold feed looks safe. Contrast HALITE (NaCl): no analytic K(T) in the database, so it falls back to constant-dH van't Hoff with a small POSITIVE dH (endothermic) -- normal solubility, SI\_halite drifts gently DOWN with T. At 25 C nothing in the K(T) machinery moves (every form returns logK25 exactly, byte-stable); off 25 C the run announces which K(T) form each entry uses (analytic / van't Hoff / flat) and warns if any analytic is used beyond its fitted validity range. K(T): analytic = PHREEQC -analytical\_expression anchored on logK25 (logK(T)=logK25+[ana(T)-ana(298.15)]); USGS PHREEQC, public domain.}
\begin{tutnarrative}
why hot surfaces scale with calcite (K(T) MATTERS).

ONE fixed water, speciated at five temperatures. A hard, carbonate-bearing groundwater (gypsum- and calcite-relevant) at a GIVEN pH 7.6 -- holding pH and the totals fixed isolates the lesson: the only thing that changes from point to point is the equilibrium constant K(T).

ion mol/kg w ion mol/kg w Ca 0.010 Cl 0.050 Na 0.030 SO4 0.012 HCO3 0.005 pH 7.6 (given)

THE LESSON -- the SIGN of dK/dT:

calcite dissolution is EXOTHERMIC (dH = -28.1 kJ/mol), so log K FALLS as T rises. SI\_calcite = log(IAP) - log K(T) therefore RISES with temperature -- RETROGRADE SOLUBILITY. Calcite is LESS soluble hot, which is exactly why heat-exchanger tubes and warm RO concentrate scale with CaCO3 although the cold feed looks safe. K(T) here is the PHREEQC ANALYTIC expression (-analytical\_ expression A1..A5), anchored on logK25 so logK(25 C) = logK25 EXACTLY.

gypsum analytic too, but only mildly retrograde (its solubility peaks near 40 C) -- a gentler curve for contrast.

halite NaCl carries NO analytic K(T) in phreeqc.dat (the db's -analytic line is commented out), so it falls back to constant-dH van't Hoff. dH = +1.37 kJ/mol > 0 (endothermic) $\to$ NORMAL solubility: log K(T) RISES and SI\_halite drifts DOWN. The honest contrast with calcite: the scaling verdict flips with temperature for the retrograde salt, not for the normal one.

At 25 C the K(T) machinery returns logK25 for EVERY entry (byte-stable); off 25 C the run announces the form in use per entry (analytic / van't Hoff / flat) and warns loudly if an analytic is used beyond its fitted validity range. (Calcite's analytic is valid 0 - 300 C, so the whole 10 - 80 C scan is in range.)

Speciation/mineral data: case-local constant/electrolyte/ -- a self- contained SNAPSHOT of the standards catalogue (USGS PHREEQC, public domain; see the .dat headers). The case files ECLIPSE data/standards/electrolyte/.

\end{tutnarrative}
\begin{tutgolden}
\goldrow{diag}{T10}{I}{0.0795197}
\goldrow{diag}{T10}{SI\_calcite}{0.56588}
\goldrow{diag}{T10}{SI\_gypsum}{-0.335309}
\goldrow{diag}{T10}{SI\_halite}{-4.56674}
\goldrow{diag}{T10}{aw}{0.998115}
\goldrow{diag}{T10}{pH}{7.6}
\goldrow{diag}{T25}{I}{0.0772864}
\goldrow{diag}{T25}{SI\_calcite}{0.789002}
\goldrow{diag}{T25}{SI\_gypsum}{-0.365046}
\goldrow{diag}{T25}{SI\_halite}{-4.60657}
\goldrow{diag}{T25}{aw}{0.998126}
\goldrow{diag}{T25}{pH}{7.6}
\goldrow{diag}{T40}{I}{0.074682}
\goldrow{diag}{T40}{SI\_calcite}{0.981885}
\goldmore{91}
\end{tutgolden}

\clearpage
\tutentry{overlay01\_nacl\_ksp}{tutorials/props/electrolyte/overlay01\_nacl\_ksp}
\begin{tutscheme}
\begin{tikzpicture}[>={Stealth[length=2mm]}, every node/.style={font=\scriptsize}]
  \node[draw, rounded corners=2pt, align=center, minimum width=3.1cm, minimum height=0.8cm, fill=black!4] at (0.0,-0.0) {\textbf{naclWater}\\[-1pt]{\tiny speciate}};
\end{tikzpicture}
\end{tutscheme}
\tutwhat{PARTIAL deep-merge overlay: a case recalibrates ONLY the halite equilibrium logK25 via a case-local overlayOf NaCl; every other datum (dH, dissolutionReaction, calorimetric, crystal) still comes from the standard components/NaCl.dat record. The [overlay] line names exactly which leaf came from the case tier.}
\tutreadme{A case recalibrates ONLY the halite equilibrium logK25 (1.60 here) to its own value, via a case-local constant/components/NaCl.dat declaring overlayOf NaCl;.}
\tutreadme{The reader deep-merges this delta over the standard components/NaCl.dat: only solidPhases.halite.equilibrium.logK25 is replaced; dH, the dissolutionReaction, the calorimetric anchor and the crystal block all still resolve from the standard record. The run prints the provenance:}
\tutreadme{This is the roadmap Phase A acceptance test: a case ships ONLY its deltas, never a second competing complete record, and per-value provenance (tier + file) is preserved.}
\begin{tutgolden}
\goldrow{diag}{naclWater}{I}{2}
\goldrow{diag}{naclWater}{SI\_halite}{-0.983442}
\goldrow{diag}{naclWater}{aw}{0.930475}
\goldrow{diag}{naclWater}{gamma\_Cl}{1.01683}
\goldrow{diag}{naclWater}{gamma\_Na}{1.01683}
\goldrow{diag}{naclWater}{m\_Cl}{2}
\goldrow{diag}{naclWater}{m\_Na}{2}
\goldrow{diag}{naclWater}{pH}{7.13303}
\goldrow{diag}{naclWater}{chargeImposed}{1}
\goldrow{diag}{naclWater}{chargeResidual}{-6.39009e-17}
\end{tutgolden}

\clearpage
\tutentry{pb82\_calcite\_open\_co2}{tutorials/props/electrolyte/pb82\_calcite\_open\_co2}
\begin{tutscheme}
\begin{tikzpicture}[>={Stealth[length=2mm]}, every node/.style={font=\scriptsize}]
  \node[draw, rounded corners=2pt, align=center, minimum width=3.1cm, minimum height=0.8cm, fill=black!4] at (0.0,-0.0) {\textbf{calciteEq\_25C}\\[-1pt]{\tiny speciate}};
\end{tikzpicture}
\end{tutscheme}
\tutwhat{Calcite in CO2-saturated water (Plummer \& Busenberg 1982): open pCO2 0.956 atm + equilibrate calcite at 25 C -- the engine's pH at calcite equilibrium beside the paper's junction-corrected 6.004 +/- 0.005 and its model value 6.011}
\begin{tutnarrative}
THE PLUMMER \& BUSENBERG (1982) EXPERIMENT, as the engine models it: calcite at equilibrium with CO2-saturated water -- an OPEN system at pCO2 = 0.956 atm, 25 C (their electrode runs' mean CO2 pressure), the gas pin, the full carbonate speciation, the CaHCO3+/CaCO3(aq) ion pairs and the mineral ceiling all engaged at once.

PRIMARY: L. N. Plummer \& E. Busenberg, Geochim. Cosmochim. Acta 46 (1982) 1011-1040, DOI 10.1016/0016-7037(82)90056-4. Owner-provided PDF.

THE ANCHOR, and what it independently tests: their MEASURED pH at calcite equilibrium under exactly these conditions, corrected for liquid-junction potential, is 6.004 +/- 0.005 (4 M KCl filling, p. 1022; uncorrected Orion runs read near 5.984); their own aqueous model puts the thermodynamic pH at 6.011. The engine's number below comes out of Henry + K1/K2 + the ion pairs + Davies + electroneutrality + the SI = 0 ceiling END TO END -- none of which was fitted to this pH.

LINEAGE, stated honestly: the calcite log K(T), K\_H, K1, K2 and ion-pair constants in our curated records descend from the USGS PHREEQC database, whose carbonate block IS this paper's Table 3 and Eqs -- so agreement in the CONSTANTS is lineage consistency, not validation. What the measured pH tests independently is the ASSEMBLY: that the engine composes those constants (gas pin, chained carbonate master, ion pairs, activity model, charge balance, mineral ceiling) into the same physical answer the electrode saw.

WHAT IT ACTUALLY LANDS AT, and where that is written. The engine gives pH 6.02946 -- 0.025 ABOVE the measurement, 0.018 above their own model's 6.011 -- and the case attributes that offset to an activity-model difference (Davies here against their Truesdell-Jones, at I = 0.027). SINCE 2026-08-12 THAT COMPARISON IS DATA, NOT PROSE: `expected` carries an `anchor` row on the measured 6.004 with a band set to the explained offset, beside the self-recorded row on 6.02946 -- so a drift AWAY from the electrode fails as an anchor disagreement, and an unexplained drift TOWARD it fails as a moved golden. Both ends, on purpose. This header used to end with "a wrong composition would not land within 0.01 pH", which reads as though the engine does; it does not, and the number is above.

APPROACH FROM ABOVE: the engine's equilibrate is a precipitation CEILING (SI $\to$ 0 from supersaturation; it does not dissolve a solid it was never given). The seed is therefore a SUPERSATURATED Ca(HCO3)2 solution under the pin; the ceiling relaxes it down to the same calcite saturation P\&B's dissolution runs approached from below -- one fixed point at (T, pCO2), whichever side you arrive from. The HCO3 seed is an initial guess only: the CO2 pin REPLACES the carbonate mole balance, so dissolved carbonate is a solved outcome (announced by the run).

\end{tutnarrative}
\begin{tutgolden}
\goldrow{diag}{calciteEq\_25C}{I}{0.0265778}
\goldrow{diag}{calciteEq\_25C}{SI\_aragonite}{-0.114}
\goldrow{diag}{calciteEq\_25C}{SI\_calcite}{3.22147e-11}
\goldrow{diag}{calciteEq\_25C}{SI\_portlandite}{-13.0919}
\goldrow{diag}{calciteEq\_25C}{aw}{0.998916}
\goldrow{diag}{calciteEq\_25C}{n\_calcite}{0.00543277}
\goldrow{diag}{calciteEq\_25C}{pH}{6.02946}
\goldrow{diag}{calciteEq\_25C}{pH}{6.004}
\goldrow{diag}{calciteEq\_25C}{chargeImposed}{1}
\goldrow{diag}{calciteEq\_25C}{chargeResidual}{2.22179e-18}
\end{tutgolden}

\clearpage
\tutentry{pitzer01\_nacl}{tutorials/props/electrolyte/pitzer01\_nacl}
\begin{tutscheme}
\begin{tikzpicture}[>={Stealth[length=2mm]}, every node/.style={font=\scriptsize}]
  \node[draw, rounded corners=2pt, align=center, minimum width=3.1cm, minimum height=0.8cm, fill=black!4] at (0.0,-0.0) {\textbf{gamma\_nacl}\\[-1pt]{\tiny pitzerActivity}};
\end{tikzpicture}
\end{tutscheme}
\tutwhat{Pitzer mean ionic activity coefficient gamma\_pm(m) for NaCl in water at 25 C --- validated vs Robinson \& Stokes (AAD 0.17\%). The glass-box electrolyte foundation: see the famous gamma\_pm dip (\textasciitilde{}0.66 near 1 m) and the rise above 1 at high concentration.}
\begin{tutnarrative}
the first electrolyte case.

gamma\_pm, the osmotic coefficient phi and the water activity a\_w for NaCl, by the Pitzer ion-interaction model. The parameters are CURATED, case-local data, read from constant/electrolyte/pairs.dat (PRIMARY: Appelo, Appl. Geochem. 55 (2015) 62 --- a 3-parameter set, beta0/beta1/Cphi). Open pitzer\_nacl.csv: gamma\_pm dips to \textasciitilde{}0.657 near 1 mol/kg, then rises above 1 by \textasciitilde{}6 mol/kg --- the textbook NaCl signature, from a closed form.

NOTE on the two L\_phi columns: pitzer\_nacl.csv reports L\_phi\_fit and L\_phi\_anchored, and their SPLIT is the lesson. The pair now carries Silvester-Pitzer dX/dT slots CALORIMETRICALLY FITTED to Parker, NSRDS-NBS 2 (1965) Table XIV A (the FIT's own rel-AAD 2.5 \%, m 0.1-6; bin/curate/fit\_lphi\_pitzer.py): the fit column tracks the measured heat-of-dilution curve -- the sign change near 1 mol/kg, the deep negative branch to -480 cal/mol -- while the anchored column (slots zeroed, Gibbs-Helmholtz-consistent but uncalibrated) misses it entirely. The validation block below pins the IN-ENGINE AAD vs the same curated Parker series, golden-locked (rel-AAD 10.24 \% at the 25 measured points with |Phi\_L| >= 50 cal/mol in 0.1-6 mol/kg).

WHY 10.24 \% HERE AND \textasciitilde{}2.5 \% IN THE FIT TOOL -- adjudicated 2026-09-03. Both compute the SAME statistic on the SAME 25 points, and the tool's Python replica reproduces this kernel's L\_phi to 7e-6 J/mol given the same parameters -- so it is not two implementations of one equation. It is two PARAMETERISATIONS in one record: Na-Cl.dat declares tCorrelationForm silvesterPitzer1977 and carries the SP77 T-functions (sp\_q1..q14, fitted by Silvester and Pitzer over 0-300 C) AND the linear calorimetric slots fitted to Parker. The engine honours the DECLARED form, so the 10.24 \% is SP77's disagreement with Parker at 25 C; the fit tool evaluates the linear slots it fitted to exactly these points, which is where 2.5 \% (2.59 \% with the stored slots) comes from. The engine now announces, when it reads this record, that the slots are present and not read. Nothing was adjusted toward Parker. The gamma\_pm / phi / a\_w values come from the untouched 25 C surface and are identical in both.

\end{tutnarrative}
\begin{tutgolden}
\goldrow{diag}{gamma\_nacl}{aw\_1m}{0.9668}
\goldrow{diag}{gamma\_nacl}{gamma\_pm\_1m}{0.6572}
\goldrow{diag}{gamma\_nacl}{gamma\_pm\_3m}{0.7140}
\goldrow{diag}{gamma\_nacl}{phi\_1m}{0.9363}
\goldrow{diag}{gamma\_nacl}{Lphi\_AAD\_abs\_Jmol}{97.6414}
\goldrow{diag}{gamma\_nacl}{Lphi\_AAD\_rel\_pct}{10.2355}
\goldrow{diag}{gamma\_nacl}{Lphi\_nMeas}{29.0000}
\end{tutgolden}

\clearpage
\tutentry{pitzer02\_cacl2}{tutorials/props/electrolyte/pitzer02\_cacl2}
\begin{tutscheme}
\begin{tikzpicture}[>={Stealth[length=2mm]}, every node/.style={font=\scriptsize}]
  \node[draw, rounded corners=2pt, align=center, minimum width=3.1cm, minimum height=0.8cm, fill=black!4] at (0.0,-0.0) {\textbf{gamma\_cacl2}\\[-1pt]{\tiny pitzerActivity}};
\end{tikzpicture}
\end{tutscheme}
\tutwhat{Pitzer gamma\_pm(m) for CaCl2 (a 2:1 salt) in water at 25 C --- the SAME general formula as NaCl, only the charges/stoichiometry differ. Validated vs Robinson \& Stokes (AAD 0.49\%, 0.1-2 m).}
\begin{tutnarrative}
the same Pitzer engine on a 2:1 electrolyte (Ca2+, 2 Cl-). Parameters curated case-local in constant/electrolyte/pairs.dat (PRIMARY: Appelo, Appl. Geochem. 55 (2015) 62). gamma\_pm dips below 0.45 then climbs steeply --- the divalent-cation signature. Valid to \textasciitilde{}2 mol/kg; CaCl2 carries a non-zero beta2 (-1.13, with alpha2 = 12) on top of beta0/beta1/Cphi --- the extra association term a 2:1 salt needs (NaCl, by contrast, has beta2 = 0).

\end{tutnarrative}
\begin{tutgolden}
\goldrow{diag}{gamma\_cacl2}{aw\_1m}{0.9449}
\goldrow{diag}{gamma\_cacl2}{gamma\_pm\_1m}{0.4985}
\goldrow{diag}{gamma\_cacl2}{gamma\_pm\_3m}{1.4771}
\goldrow{diag}{gamma\_cacl2}{phi\_1m}{1.0483}
\end{tutgolden}

\clearpage
\tutentry{pitzer02\_nacl\_package}{tutorials/props/electrolyte/pitzer02\_nacl\_package}
\begin{tutscheme}
\begin{tikzpicture}[>={Stealth[length=2mm]}, every node/.style={font=\scriptsize}]
  \node[draw, rounded corners=2pt, align=center, minimum width=3.1cm, minimum height=0.8cm, fill=black!4] at (0.0,-0.0) {\textbf{gamma\_nacl}\\[-1pt]{\tiny electrolyteActivity}};
\end{tikzpicture}
\end{tutscheme}
\tutwhat{NaCl Pitzer via the NEW property-package builder (reads zero old salt files)}
\begin{tutnarrative}
No `cation`/`anion` keys: the salt is resolved by the SELECTED propertyPackage (constant/thermoPhysPropDict $\to$ aqueousNaCl\_pitzer). The op consumes the assembled package via thermo.electrolyte() and reports gamma\_pm / phi / a\_w.

\end{tutnarrative}
\begin{tutgolden}
\goldrow{diag}{gamma\_nacl}{aw\_1m}{0.966826}
\goldrow{diag}{gamma\_nacl}{gamma\_pm\_1m}{0.657172}
\goldrow{diag}{gamma\_nacl}{gamma\_pm\_3m}{0.714026}
\goldrow{diag}{gamma\_nacl}{phi\_1m}{0.93634}
\end{tutgolden}

\clearpage
\tutentry{pitzer\_calcite\_brine}{tutorials/props/electrolyte/pitzer\_calcite\_brine}
\begin{tutscheme}
\begin{tikzpicture}[>={Stealth[length=2mm]}, every node/.style={font=\scriptsize}]
  \node[draw, rounded corners=2pt, align=center, minimum width=3.1cm, minimum height=0.8cm, fill=black!4] at (0.0,-0.0) {\textbf{seawater\_pitzer}\\[-1pt]{\tiny speciate}};
  \node[draw, rounded corners=2pt, align=center, minimum width=3.1cm, minimum height=0.8cm, fill=black!4] at (3.6,-0.0) {\textbf{seawater\_davies}\\[-1pt]{\tiny speciate}};
\end{tikzpicture}
\end{tutscheme}
\tutwhat{Pitzer-HMW carbonate system on surface seawater -- pins SI\_calcite / SI\_gypsum against the published oceanographic saturation state, captures the neutral CO2 salting-out (gamma\_CO2aq > 1 via lambda), and the Davies-vs-Pitzer SI divergence.}
\begin{tutnarrative}
the slice-S4 GATE for the Pitzer CARBONATE system.

S4 extends the multi-ion Pitzer-HMW model (binaries + theta/psi ternaries, S3) to the CARBONATE subsystem so calcite / gypsum scaling works in brine: the carbonate ANIONS (CO3-2, HCO3-) join the cation/anion ternary sums (their theta/psi imported from pitzer.dat into mixing.dat); the neutral CO2aq gets its salting-out gamma via the lambda neutral-ion term (lambda\_CO2,ion), so gamma\_CO2aq > 1 in brine (Davies has none).

COMPOSITION: standard surface seawater (S = 35) WITH the carbonate system -- the major-ion subset of the S3 verify case + DIC \textasciitilde{} 2.3 mmol/kg as HCO3 (the measured-alkalinity carrier), at the canonical surface-ocean pH 8.2. Na+ 0.4861 Mg+2 0.0547 Ca+2 0.01065 K+ 0.01058 Cl- 0.5658 SO4-2 0.02926 HCO3- 0.00230 (I \textasciitilde{} 0.68 mol/kg)

THE PUBLISHED PIN (cited; E\_theta before/after recorded below):

SI\_calcite (surface seawater, 25 C, pH 8.2) \textasciitilde{} +0.6 to +0.8. Surface ocean water is calcite-SUPERSATURATED: the calcite saturation ratio Omega\_calcite = IAP/Ksp \textasciitilde{} 4-6, i.e. SI = log10(Omega) \textasciitilde{} +0.6 to +0.8. Sources: Mucci, A., Am. J. Sci. 283 (1983) 780 (the seawater calcite/ aragonite solubility); Millero, F.J., Chem. Rev. 107 (2007) 308 (the seawater carbonate system); a well-documented oceanographic number (the ocean-acidification baseline). Choupo (Pitzer-HMW) computes SI\_calcite = +0.655 -- INSIDE the published band. SI\_gypsum (seawater) \textasciitilde{} -0.9 to -0.6 (strongly UNDERsaturated -- seawater is nowhere near gypsum saturation); Choupo computes -0.964 -- a touch below the -0.9 floor (this carbonate-bearing brine sits at slightly higher I than the carbonate-free seawater\_verify case, which gives -0.687 in-band). Still firmly undersaturated, the physically-correct verdict.

E\_theta(I) BEFORE $\to$ AFTER (constant-theta $\to$ higher-order mixing active): SI\_calcite +0.724 $\to$ +0.655 (IN band +0.6..+0.8) -- E\_theta lowers the 2+/2- ion gammas, nudging the carbonate IAP SI\_gypsum -0.875 $\to$ -0.964 gamma\_CO2aq 1.105 $\to$ 1.106 (salting-out kept) gamma\_Ca 0.227 $\to$ 0.194 (E\_theta on Ca pairs) Davies (untouched -- E\_theta is Pitzer-only): SI\_calcite +0.868, byte-ident.

The golden master locks Choupo's exact output; the band above validates it is PHYSICALLY right (a wrong pin is worse than no pin -- the band is cited).

DAVIES vs PITZER (the differentiator, also pinned): the SAME composition under Davies gives SI\_calcite +0.868 and gamma\_CO2aq = 1.000 (no neutral salting- out) -- Davies is out of its trustworthy band (\textasciitilde{}0.5) at this I, overshoots the calcite SI by \textasciitilde{}0.21 vs Pitzer-HMW-with-E\_theta and misses the CO2 salting-out entirely. Pitzer is the rigorous tool for brine. Both entries run here so the divergence is VISIBLE.

DATA SCOPE: self-DESCRIBING (inline manifest); the ions, pairs, mixing WITH carbonate theta/psi/lambda, speciation and minerals resolve from the curated data/standards/ tree at runtime -- seal with bin/choupo-import for a runtime-self-contained copy.

\end{tutnarrative}
\begin{tutgolden}
\goldrow{diag}{seawater\_pitzer}{I}{0.682714}
\goldrow{diag}{seawater\_pitzer}{SI\_anhydrite}{-1.24804}
\goldrow{diag}{seawater\_pitzer}{SI\_aragonite}{0.540696}
\goldrow{diag}{seawater\_pitzer}{SI\_calcite}{0.654696}
\goldrow{diag}{seawater\_pitzer}{SI\_gypsum}{-0.964342}
\goldrow{diag}{seawater\_pitzer}{SI\_halite}{-2.48855}
\goldrow{diag}{seawater\_pitzer}{SI\_sylvite}{-3.51959}
\goldrow{diag}{seawater\_pitzer}{aw}{0.981412}
\goldrow{diag}{seawater\_pitzer}{gamma\_CO2aq}{1.10585}
\goldrow{diag}{seawater\_pitzer}{gamma\_CO3}{0.105443}
\goldrow{diag}{seawater\_pitzer}{gamma\_Ca}{0.194358}
\goldrow{diag}{seawater\_pitzer}{gamma\_HCO3}{0.659478}
\goldrow{diag}{seawater\_pitzer}{m\_CO2aq}{1.4318e-05}
\goldrow{diag}{seawater\_pitzer}{m\_CO3}{7.71637e-05}
\goldmore{94}
\end{tutgolden}

\clearpage
\tutentry{pitzer\_gamma\_hamer\_wu}{tutorials/props/electrolyte/pitzer\_gamma\_hamer\_wu}
\begin{tutscheme}
\begin{tikzpicture}[>={Stealth[length=2mm]}, every node/.style={font=\scriptsize}]
  \node[draw, rounded corners=2pt, align=center, minimum width=3.1cm, minimum height=0.8cm, fill=black!4] at (0.0,-0.0) {\textbf{gamma\_nacl}\\[-1pt]{\tiny pitzerActivity}};
  \node[draw, rounded corners=2pt, align=center, minimum width=3.1cm, minimum height=0.8cm, fill=black!4] at (3.6,-0.0) {\textbf{gamma\_kcl}\\[-1pt]{\tiny pitzerActivity}};
  \node[draw, rounded corners=2pt, align=center, minimum width=3.1cm, minimum height=0.8cm, fill=black!4] at (7.2,-0.0) {\textbf{gamma\_kbr}\\[-1pt]{\tiny pitzerActivity}};
  \node[draw, rounded corners=2pt, align=center, minimum width=3.1cm, minimum height=0.8cm, fill=black!4] at (0.0,-1.3) {\textbf{gamma\_licl\_consistency}\\[-1pt]{\tiny pitzerActivity}};
  \node[draw, rounded corners=2pt, align=center, minimum width=3.1cm, minimum height=0.8cm, fill=black!4] at (3.6,-1.3) {\textbf{gamma\_hcl}\\[-1pt]{\tiny pitzerActivity}};
  \node[draw, rounded corners=2pt, align=center, minimum width=3.1cm, minimum height=0.8cm, fill=black!4] at (7.2,-1.3) {\textbf{gamma\_naoh}\\[-1pt]{\tiny pitzerActivity}};
\end{tikzpicture}
\end{tutscheme}
\tutwhat{Pitzer gamma\_pm/phi vs Hamer \& Wu 1972 -- 6 salts, evaluated NBS data}
\begin{tutnarrative}
THE BREADTH VALIDATION: the catalogue's Pitzer pairs against the NBS critically-evaluated gamma\_pm and osmotic-coefficient tables -- Hamer \& Wu, J. Phys. Chem. Ref. Data 1 (1972) 1047 -- for SIX 1-1 salts at 25 C. Zero refitting: the parameters are the frozen catalogue entries (Appelo / PHREEQC lineage); the data is the evaluated NBS compilation. Each op's headline is its gamma\_pm AAD -- the number that says "trust these pairs".

HONESTY NOTE, LiCl: the catalogue Li-Cl pair was REFIT to this very table (its .dat says so), so its AAD here is a FIT-CONSISTENCY check, not independent validation -- and at \textasciitilde{}3 \% it is the WORST of the six: the two-anchor compromise (this table + Farelo saturation at 19 mol/kg) stretches 3 parameters past what the most non-ideal 1-1 salt allows. The five pairs that never saw this data land at 0.1-0.7 \%.

Windows stay inside each pair's declared validity (KCl saturates at 4.803 mol/kg; KBr data used to 5).

\end{tutnarrative}
\begin{tutgolden}
\goldrow{diag}{gamma\_nacl}{aw\_1m}{0.966826401114}
\goldrow{diag}{gamma\_nacl}{gamma\_AAD\_pct}{0.111724891644}
\goldrow{diag}{gamma\_nacl}{gamma\_nMeas}{8}
\goldrow{diag}{gamma\_nacl}{gamma\_pm\_1m}{0.657171576559}
\goldrow{diag}{gamma\_nacl}{gamma\_pm\_3m}{0.71402584056}
\goldrow{diag}{gamma\_nacl}{phi\_1m}{0.936339794474}
\goldrow{diag}{gamma\_nacl}{phi\_AAD\_pct}{0.0869907167063}
\goldrow{diag}{gamma\_nacl}{phi\_nMeas}{8}
\goldrow{diag}{gamma\_kcl}{aw\_1m}{0.968139226009}
\goldrow{diag}{gamma\_kcl}{gamma\_AAD\_pct}{0.13395574493}
\goldrow{diag}{gamma\_kcl}{gamma\_nMeas}{7}
\goldrow{diag}{gamma\_kcl}{gamma\_pm\_1m}{0.604281386529}
\goldrow{diag}{gamma\_kcl}{gamma\_pm\_3m}{0.569777816543}
\goldrow{diag}{gamma\_kcl}{phi\_1m}{0.898678143951}
\goldmore{49}
\end{tutgolden}

\clearpage
\tutentry{pitzer\_hmw\_hot\_agreement}{tutorials/props/electrolyte/pitzer\_hmw\_hot\_agreement}
\begin{tutscheme}
\begin{tikzpicture}[>={Stealth[length=2mm]}, every node/.style={font=\scriptsize}]
  \node[draw, rounded corners=2pt, align=center, minimum width=3.1cm, minimum height=0.8cm, fill=black!4] at (0.0,-0.0) {\textbf{hmw\_nacl\_100C}\\[-1pt]{\tiny speciate}};
  \node[draw, rounded corners=2pt, align=center, minimum width=3.1cm, minimum height=0.8cm, fill=black!4] at (3.6,-0.0) {\textbf{kernel\_nacl\_100C}\\[-1pt]{\tiny pitzerActivity}};
\end{tikzpicture}
\end{tutscheme}
\tutwhat{HMW vs single-salt Pitzer at 100 C, same NaCl 4 m state: the two surfaces must AGREE now that both apply the SP77 T-functions}
\begin{tutnarrative}
one salt, one temperature, TWO Pitzer engines: they must agree (forum 100.2/\#101, the HMW-virials(T) fix).

NaCl at 4 mol/kg, 100 C. Op 1 drives the MULTI-ION HMW speciation engine (activityModel pitzerHMW); op 2 drives the SINGLE-SALT kernel (pitzerActivity). Both now read the SAME Na-Cl.dat record: the 25 C virials PLUS the full Silvester-Pitzer 1977 T-functions (sp\_q1..q14) and the same A\_phi(T). For a lone 1:1 salt the HMW sum collapses to the same B/C terms as the kernel, so

gamma\_pm(HMW) = sqrt(gamma\_Na $\ast$ gamma\_Cl) == gamma\_pm(kernel)

to numerical noise. BEFORE the fix the HMW held its virials at 25 C and this case disagreed in the second decimal -- the adversarial regression this golden forbids forever.

(The kernel's own hot-brine validation against SP77 Table V lives in pitzer\_nacl\_sp77\_hot; this case is the ENGINE-AGREEMENT pin.)

\end{tutnarrative}
\begin{tutgolden}
\goldrow{diag}{hmw\_nacl\_100C}{I}{4.0000}
\goldrow{diag}{hmw\_nacl\_100C}{SI\_halite}{-0.7734}
\goldrow{diag}{hmw\_nacl\_100C}{aw}{0.8533}
\goldrow{diag}{hmw\_nacl\_100C}{gamma\_Cl}{0.7309}
\goldrow{diag}{hmw\_nacl\_100C}{gamma\_Na}{0.7309}
\goldrow{diag}{hmw\_nacl\_100C}{m\_Cl}{4.0000}
\goldrow{diag}{hmw\_nacl\_100C}{m\_Na}{4.0000}
\goldrow{diag}{hmw\_nacl\_100C}{pH}{7.0000}
\goldrow{diag}{kernel\_nacl\_100C}{aw\_1m}{0.9670}
\goldrow{diag}{kernel\_nacl\_100C}{gamma\_pm\_1m}{0.6221}
\goldrow{diag}{kernel\_nacl\_100C}{gamma\_pm\_3m}{0.6760}
\goldrow{diag}{kernel\_nacl\_100C}{phi\_1m}{0.9327}
\goldrow{diag}{hmw\_nacl\_100C}{chargeImposed}{0.0000}
\goldrow{diag}{hmw\_nacl\_100C}{chargeResidual}{-0.0000}
\end{tutgolden}

\clearpage
\tutentry{pitzer\_nacl\_sp77\_hot}{tutorials/props/electrolyte/pitzer\_nacl\_sp77\_hot}
\begin{tutscheme}
\begin{tikzpicture}[>={Stealth[length=2mm]}, every node/.style={font=\scriptsize}]
  \node[draw, rounded corners=2pt, align=center, minimum width=3.1cm, minimum height=0.8cm, fill=black!4] at (0.0,-0.0) {\textbf{nacl\_100C}\\[-1pt]{\tiny pitzerActivity}};
  \node[draw, rounded corners=2pt, align=center, minimum width=3.1cm, minimum height=0.8cm, fill=black!4] at (3.6,-0.0) {\textbf{nacl\_150C}\\[-1pt]{\tiny pitzerActivity}};
  \node[draw, rounded corners=2pt, align=center, minimum width=3.1cm, minimum height=0.8cm, fill=black!4] at (7.2,-0.0) {\textbf{nacl\_200C}\\[-1pt]{\tiny pitzerActivity}};
\end{tikzpicture}
\end{tutscheme}
\tutwhat{NaCl Pitzer to 200 C -- Silvester \& Pitzer 1977 T-functions, validated vs their Table V}
\begin{tutnarrative}
HOT BRINE: the NaCl Pitzer surface from 25 to 200 C on the FULL Silvester \& Pitzer 1977 temperature functions (J. Phys. Chem. 81, 1822 -- eqs 30-32, Table IV coefficients VERBATIM in the standards Na-Cl.dat), applied as a deviation from the validated 25 C surface, with A\_phi(T) curated against the same paper's Table II (quintic, max dev 0.17 \%, 0-300 C; the old closed form drifted +67 \% at 300 C).

Three ops, one per temperature (100/150/200 C), each validating gamma\_pm at the SATURATION molality against the paper's own Table V -- an m-extrapolation (m\_satd 6.7-8.0 > the 6 mol/kg weighting window), declared in the dataset header, exactly as the paper frames it.

What to look at: gamma\_pm(1 m) falls 0.657 $\to$ \textasciitilde{}0.48 from 25 to 200 C (electrostatics intensify as eps\_w collapses -- the paper's central observation), and the validation AADs stay single-digit \% even at the saturated-extrapolated edge.

\end{tutnarrative}
\begin{tutgolden}
\goldrow{diag}{nacl\_100C}{aw\_1m}{0.96695360248}
\goldrow{diag}{nacl\_100C}{gamma\_AAD\_pct}{0.285062879686}
\goldrow{diag}{nacl\_100C}{gamma\_nMeas}{1}
\goldrow{diag}{nacl\_100C}{gamma\_pm\_1m}{0.622051060857}
\goldrow{diag}{nacl\_100C}{gamma\_pm\_3m}{0.675983294046}
\goldrow{diag}{nacl\_100C}{phi\_1m}{0.932688469875}
\goldrow{diag}{nacl\_150C}{aw\_1m}{0.967903934433}
\goldrow{diag}{nacl\_150C}{gamma\_AAD\_pct}{0.0424658694305}
\goldrow{diag}{nacl\_150C}{gamma\_nMeas}{1}
\goldrow{diag}{nacl\_150C}{gamma\_pm\_1m}{0.559281383488}
\goldrow{diag}{nacl\_150C}{gamma\_pm\_3m}{0.573568787797}
\goldrow{diag}{nacl\_150C}{phi\_1m}{0.905424310805}
\goldrow{diag}{nacl\_200C}{aw\_1m}{0.969344388695}
\goldrow{diag}{nacl\_200C}{gamma\_AAD\_pct}{0.431744113932}
\goldmore{4}
\end{tutgolden}

\clearpage
\tutentry{pitzer\_seawater\_verify}{tutorials/props/electrolyte/pitzer\_seawater\_verify}
\begin{tutscheme}
\begin{tikzpicture}[>={Stealth[length=2mm]}, every node/.style={font=\scriptsize}]
  \node[draw, rounded corners=2pt, align=center, minimum width=3.1cm, minimum height=0.8cm, fill=black!4] at (0.0,-0.0) {\textbf{reduction\_nacl}\\[-1pt]{\tiny speciate}};
  \node[draw, rounded corners=2pt, align=center, minimum width=3.1cm, minimum height=0.8cm, fill=black!4] at (3.6,-0.0) {\textbf{seawater}\\[-1pt]{\tiny speciate}};
\end{tikzpicture}
\end{tutscheme}
\tutwhat{Pitzer-HMW verification: (binaries + ternary theta/psi) on standard synthetic seawater (I \textasciitilde{} 0.7). Pins per-ion + mean activity coefficients against PUBLISHED HMW-1984 / seawater-Pitzer values, and the single-salt reduction (ternaries $\to$ 0).}
\begin{tutnarrative}
the TIER-2 verification weapon for the Pitzer-HMW ternary-mixing slice (S3).

The forum gate: "NO published-seawater golden masters => NO-GO." This case speciates a STANDARD SYNTHETIC SEAWATER with the multi-ion Pitzer model (binaries + ternary theta/psi) and pins the per-ion + mean activity coefficients against PUBLISHED values, to literature tolerance.

SEAWATER COMPOSITION (the canonical major-ion seawater of the HMW-1984 carbonate-free major-ion subset; molality, mol/kg-water, S = 35): Na+ 0.4861 Mg+2 0.0547 Ca+2 0.01065 K+ 0.01058 Cl- 0.5658 SO4-2 0.02926 giving ionic strength I \textasciitilde{} 0.70 mol/kg -- the classic seawater benchmark. (The carbonate system -- HCO3/CO3, \textasciitilde{}2.3 mmol/kg, < 0.4\% of I -- is OMITTED: it is slice S4, and its absence shifts I by < 0.3\% and the major-ion gammas by < 0.2\%, well inside the pins' tolerance.) pH is GIVEN at 8.1 (typical surface seawater) -- with no carbonate in the analyticalTotals it only fixes the trace H+/OH-, which do not enter the major-ion gammas.

WHY PITZER, NOT DAVIES: at I \textasciitilde{} 0.7 Davies is already out of its trustworthy band (\textasciitilde{}0.5). The Pitzer specific-interaction model -- WITH the ternary theta/psi mixing terms AND the E\_theta(I) higher-order electrostatic mixing term (Pitzer 1975, now ACTIVE for unlike-charge same-sign pairs) -- is the standard tool for seawater and brines to I \textasciitilde{} 6. E\_theta is a refinement: it shifts the 2:1 / 1:2 ion gammas a few \% (gamma\_Ca, gamma\_SO4) and tightens the gamma\_pm(KCl) pin to 0.638 (inside 0.62-0.64); the 1:1 NaCl gamma\_pm stays \textasciitilde{}0.66. (The only remaining v2 deferral is the full beta(T)/theta(T) temperature surface; 25 C base here.)

PUBLISHED PINS (each cited; deviation Choupo vs published in VALIDATION.md):

gamma\_pm(NaCl) in seawater \textasciitilde{} 0.66-0.68 Pitzer's seawater model / Millero \& Pierrot; classic textbook value gamma\_pm(NaCl, seawater I\textasciitilde{}0.7) \textasciitilde{} 0.67. gamma\_pm(KCl) \textasciitilde{} 0.62-0.64 ; gamma\_pm(MgCl2), gamma\_pm(CaCl2) etc. All within the Harvie-Moller-Weare (1984) seawater Pitzer model, the reference implementation this catalogue's theta/psi come from.

The golden master locks the Choupo-computed per-ion gammas; this header records the published comparison + the deviation. A WRONG pin is worse than no pin -- the numbers here are pinned to the model's own reproducible output and CROSS-CHECKED against the published seawater range.

FINDING, 2026-08-10: "ALL PINS STILL IN BAND" WAS CHECKED ON TWO OF SIX.

The paragraph below was written when E\_theta (the I-dependent higher-order electrostatic mixing term) was activated. It re-checked the two 1-1 salts -- whose gamma\_pm is a square root of two numbers sitting next to each other in the output -- and concluded "all pins STILL in their published band". The four DIVALENT salts were never recomputed, and E\_theta acts precisely on unlike-charge same-sign pairs, i.e. on exactly those four.

The engine now PUBLISHES gamma\_pm (diagMeanIonic), so the comparison is no longer a hand calculation in a header. Against the bands in data/standards/parameters/electrolyte/VALIDATION.md:

NaCl 0.6646 band 0.66-0.69 IN KCl 0.6381 band 0.62-0.65 IN MgCl2 0.4608 band 0.47-0.50 BELOW CaCl2 0.4459 band 0.46-0.50 BELOW Na2SO4 0.3489 band 0.36-0.40 BELOW MgSO4 0.1460 band 0.15-0.20 BELOW

Four of six are outside, by 2-4 \%. VALIDATION.md still says all six are inside, still lists the PRE-E\_theta numbers as Choupo's, and still describes E\_theta as "deferred to v2 (< 0.5 \% on the major-ion gamma at this I)" -- while this run announces "E\_theta now ACTIVE" and the divalent shift is \textasciitilde{}5 \%. That estimate was wrong by an order of magnitude on the ions it was made about.

NOT RESOLVED, and deliberately NOT TUNED. Two readings are open and this case cannot choose between them: the quoted bands are loose secondary citations ("\textasciitilde{} 0.47-0.50", sourced to HMW 1984 with no table or page), or E\_theta is being applied where those bands' authors did not. Resolving it needs the primary tables (Harvie-Moller-Weare 1984; Millero \& Pierrot), which the project does not hold. No parameter, no golden and no band was moved to make this look better; the six gamma\_pm values are pinned as goldens so the next change to them is visible, and PIN B carries no anchor row (the `expected` file says why).

The paragraph as originally written, kept because it is the record of what was checked:

E\_theta(I) BEFORE $\to$ AFTER (constant-theta $\to$ higher-order mixing active), all pins STILL in their published band: gamma\_pm(NaCl) 0.6714 $\to$ 0.6646 (band 0.66-0.68; 1:1, E\_theta=0 for Na-Cl, shift via coupling) -- IN BAND gamma\_pm(KCl) 0.6447 $\to$ 0.6381 (band 0.62-0.64) -- E\_theta moved it INTO band gamma\_Ca 0.2141 $\to$ 0.1856 (2+ ion; E\_theta acts via Ca-Na/Ca-K/Ca-Mg) gamma\_SO4 0.1227 $\to$ 0.1040 (2- ion; E\_theta acts via SO4-Cl/SO4-OH) gamma\_Mg 0.2363 $\to$ 0.2049 gamma\_Na 0.6495 $\to$ 0.6390 gamma\_Cl 0.6941 $\to$ 0.6911 Single-salt reduction (PIN A, NaCl alone): gamma\_pm UNCHANGED at 0.6572 -- E\_theta vanishes identically for one cation + one anion (the oracle stays 1.5658e-14 over all 54 binaries).

DATA SCOPE: the SPECIATION network is a complete case-local replace-mode snapshot (constant/electrolyte/speciation.dat) -- that claim covers the speciation data only. Ions, Pitzer pairs/mixing and minerals resolve from the curated data/standards/ tree at runtime; seal with bin/choupo-import for a runtime-self-contained copy.

\end{tutnarrative}
\begin{tutgolden}
\goldrow{diag}{reduction\_nacl}{I}{1}
\goldrow{diag}{reduction\_nacl}{SI\_halite}{-1.93464}
\goldrow{diag}{reduction\_nacl}{aw}{0.966826}
\goldrow{diag}{reduction\_nacl}{gamma\_Cl}{0.657172}
\goldrow{diag}{reduction\_nacl}{gamma\_Na}{0.657172}
\goldrow{diag}{reduction\_nacl}{gamma\_pm\_Na\_Cl}{0.657172}
\goldrow{diag}{reduction\_nacl}{m\_Cl}{1}
\goldrow{diag}{reduction\_nacl}{m\_Na}{1}
\goldrow{diag}{reduction\_nacl}{pH}{7}
\goldrow{diag}{seawater}{I}{0.720461}
\goldrow{diag}{seawater}{SI\_anhydrite}{-0.970682}
\goldrow{diag}{seawater}{SI\_gypsum}{-0.687047}
\goldrow{diag}{seawater}{SI\_halite}{-2.48552}
\goldrow{diag}{seawater}{SI\_sylvite}{-3.51313}
\goldmore{50}
\end{tutgolden}

\clearpage
\tutentry{pitzer\_vs\_davies\_ro\_concentrate}{tutorials/props/electrolyte/pitzer\_vs\_davies\_ro\_concentrate}
\begin{tutscheme}
\begin{tikzpicture}[>={Stealth[length=2mm]}, every node/.style={font=\scriptsize}]
  \node[draw, rounded corners=2pt, align=center, minimum width=3.1cm, minimum height=0.8cm, fill=black!4] at (0.0,-0.0) {\textbf{scan\_davies}\\[-1pt]{\tiny scalingScan}};
  \node[draw, rounded corners=2pt, align=center, minimum width=3.1cm, minimum height=0.8cm, fill=black!4] at (3.6,-0.0) {\textbf{scan\_pitzer}\\[-1pt]{\tiny scalingScan}};
\end{tikzpicture}
\end{tutscheme}
\tutwhat{Davies vs Pitzer-HMW on a seawater RO concentrate: gypsum/calcite SI vs recovery from the same feed under two activity models (full narrative in propsDict).}
\begin{tutnarrative}
THE FLAGSHIP of the Pitzer x speciation feature (slice S6): the payoff that makes the whole multi-ion Pitzer-HMW build worth it. ONE feed, TWO activity models, the engineering answer FORKS.

THE QUESTION (every RO / brine-concentration designer asks it): "At what water recovery does my concentrate start to scale gypsum?" The answer DEPENDS ON THE ACTIVITY MODEL -- and that is the whole point.

THE FEED: standard surface seawater (S = 35), the major-ion subset: Na+ 0.4861 Mg+2 0.0547 Ca+2 0.01065 K+ 0.01058 (cations) Cl- 0.5658 SO4-2 0.02926 HCO3- 0.0024 (anions, DIC carrier) I \textasciitilde{} 0.7 mol/kg at the feed. Recovery r scales every total as feed/(1 - r); scanning r = 0 $\to$ 0.80 drives I from \textasciitilde{}0.7 up to \textasciitilde{}3.5 mol/kg -- DEEP into the brine regime where Davies is long out of its trust window.

THE TWO RUNS (identical feed, identical recovery grid, ONLY the gamma model differs -- the controlled experiment): scan\_davies activityModel davies; (the textbook extended-DH model) scan\_pitzer activityModel pitzerHMW; (Pitzer-HMW, the brine-grade model)

THE PAYOFF ARTEFACTS (read straight off the two CSVs + the KPIs): 1. SI\_gypsum vs recovery / I, BOTH models -- the HEADLINE. They track together at low I and FORK visibly past I \textasciitilde{} 1: Davies keeps climbing, Pitzer's gypsum SI is held DOWN by the strong Ca-SO4 specific interaction (beta2 = -59.3, the 2:2 ion-pairing virial Davies cannot see). The gypsum-onset recovery is LATER under Pitzer. 2. SI\_calcite vs recovery / I, both models -- same fork, carbonate flavour. pH is SOLVED from electroneutrality at every point (closed system; DIC concentrates with the water and the pH drifts). 3. gamma\_Ca2+ and gamma\_SO4-2 vs I, both models -- the MECHANISM plot. Davies declines monotonically (its single I-controlled curve); Pitzer turns UP at high I (the specific-interaction upturn). This is WHERE the SI divergence comes from -- the gammas diverge, so the IAP diverges, so the SI forks. (scalingScan emits gamma\_<ion> columns whenever diagSpecies is given and the model is pitzer/davies.)

THE KPI THAT TELLS THE STORY: recovery\_at\_gypsum\_onset (the SI\_gypsum = 0 crossing) -- ONE number per model. Davies predicts gypsum scaling at one recovery; Pitzer at another. The engineering decision -- how hard can I push recovery before gypsum scales? -- flips on the activity model. Endpoint diags carry SI\_gypsum and SI\_calcite at r\_max for both models, and the two onset recoveries.

WHY PITZER IS THE ANSWER (cite the credentials, do not assert): Davies is ANNOUNCED-untrustworthy past I \textasciitilde{} 0.5 mol/kg (the run says so). Pitzer-HMW is built for exactly this regime -- valid to I \textasciitilde{} 6 mol/kg for the seawater system, with the E\_theta(I) higher-order electrostatic mixing term ACTIVE (Pitzer 1975) -- and is CREDENTIALED in this repo by two pinned cases: pitzer\_seawater\_verify -- single-salt oracle 1.57e-14 + HMW seawater mean activity coefficients inside the published bands; pitzer\_calcite\_brine -- SI\_calcite / SI\_gypsum on surface seawater inside the published oceanographic saturation band. So here we TRUST the Pitzer fork and read Davies as the cautionary curve.

E\_theta(I) BEFORE $\to$ AFTER (the scan\_pitzer fork only; scan\_davies is BYTE- IDENTICAL -- E\_theta is Pitzer-only). At r\_max (I \textasciitilde{} 3.3 mol/kg, deep brine where higher-order mixing bites hardest): SI\_gypsum\_rmax -0.144 $\to$ -0.317 SI\_calcite\_rmax 2.026 $\to$ 1.938 SI\_gypsum\_r0 -0.877 $\to$ -0.967 SI\_calcite\_r0 1.240 $\to$ 1.183 I\_rmax 3.271 $\to$ 3.322 gypsum stays undersaturated end-to-end The engineering verdict (gypsum the limiting scale) is UNCHANGED in direction; E\_theta refines the SI magnitude in the high-I tail.

DATA SCOPE: self-DESCRIBING (inline manifest); the ions, pairs, mixing with the carbonate theta/psi/lambda, speciation and minerals resolve from the curated data/standards/ tree at runtime -- seal with bin/choupo-import for a runtime-self-contained copy.

\end{tutnarrative}
\begin{tutgolden}
\goldrow{diag}{scan\_davies}{I\_r0}{0.657308}
\goldrow{diag}{scan\_davies}{I\_rmax}{3.03589}
\goldrow{diag}{scan\_davies}{SI\_anhydrite\_r0}{-0.889841}
\goldrow{diag}{scan\_davies}{SI\_anhydrite\_rmax}{0.269649}
\goldrow{diag}{scan\_davies}{SI\_aragonite\_r0}{1.26146}
\goldrow{diag}{scan\_davies}{SI\_aragonite\_rmax}{2.0927}
\goldrow{diag}{scan\_davies}{SI\_calcite\_r0}{1.37546}
\goldrow{diag}{scan\_davies}{SI\_calcite\_rmax}{2.2067}
\goldrow{diag}{scan\_davies}{SI\_gypsum\_r0}{-0.607638}
\goldrow{diag}{scan\_davies}{SI\_gypsum\_rmax}{0.481394}
\goldrow{diag}{scan\_davies}{SI\_halite\_r0}{-2.3956}
\goldrow{diag}{scan\_davies}{SI\_halite\_rmax}{-0.45497}
\goldrow{diag}{scan\_davies}{SI\_sylvite\_r0}{-3.3946}
\goldrow{diag}{scan\_davies}{SI\_sylvite\_rmax}{-1.44942}
\goldmore{192}
\end{tutgolden}

\clearpage
\tutentry{precipitation\_ro\_brackish}{tutorials/props/electrolyte/precipitation\_ro\_brackish}
\begin{tutscheme}
\begin{tikzpicture}[>={Stealth[length=2mm]}, every node/.style={font=\scriptsize}]
  \node[draw, rounded corners=2pt, align=center, minimum width=3.1cm, minimum height=0.8cm, fill=black!4] at (0.0,-0.0) {\textbf{pin\_gypsum\_amount}\\[-1pt]{\tiny speciate}};
  \node[draw, rounded corners=2pt, align=center, minimum width=3.1cm, minimum height=0.8cm, fill=black!4] at (3.6,-0.0) {\textbf{feed}\\[-1pt]{\tiny speciate}};
  \node[draw, rounded corners=2pt, align=center, minimum width=3.1cm, minimum height=0.8cm, fill=black!4] at (7.2,-0.0) {\textbf{feed\_eq}\\[-1pt]{\tiny speciate}};
  \node[draw, rounded corners=2pt, align=center, minimum width=3.1cm, minimum height=0.8cm, fill=black!4] at (0.0,-1.3) {\textbf{scan\_eq}\\[-1pt]{\tiny scalingScan}};
\end{tikzpicture}
\end{tutscheme}
\tutwhat{Equilibrium PRECIPITATION (the SI = 0 ceiling) of a brackish RO concentrate: aqueous speciation (Davies) + augmented-Newton chemical equilibrium with an announced active set. Sibling of scaling\_ro\_brackish (which reports SI only) -- here named minerals are ALLOWED to precipitate to saturation and the engine returns the thermodynamic MAXIMUM amount (SI $\to$ 0, infinite time, no nucleation barrier -- a CEILING, NOT a deposit forecast; real scale is kinetically limited). Story: (1) pin\_gypsum\_amount -- a CaSO4 water with Ca and SO4 DOUBLED past gypsum saturation; the machine hands back EXACTLY the over-added excess as solid (n\_gypsum). (2) feed / feed\_eq -- the brackish water raw and equilibrated side by side: calcite reaches its ceiling, releases H+ (pH drops, solved closure), and consuming shared Ca pushes gypsum further from saturation. (3) scan\_eq -- calcite + gypsum allowed across recovery to r = 0.85 with a feedFlow, giving a kg/day calcite curve; gypsum stays n = 0 throughout while its SIeq drifts BELOW the raw free-water SI once calcite consumes the shared Ca -- the coupling, visible. Cross-check tool: PHREEQC EQUILIBRIUM\_PHASES.}
\begin{tutnarrative}
the EQUILIBRIUM-PRECIPITATION ceiling.

Sibling of scaling\_ro\_brackish (which is SI-ONLY -- "how supersaturated?"). Here the `equilibrate \{ minerals ( ... ); \}` block names the ALLOWED solids, and the engine drives each to its SI = 0 saturation, returning the amount precipitated. THE FRAMING (printed loudly at run time, non-suppressible):

EQUILIBRIUM CEILING -- not a deposit prediction. The amounts are the thermodynamic MAXIMUM (SI $\to$ 0, infinite time, no nucleation barrier). Real scale is kinetically limited (induction time, antiscalants act on KINETICS, which this calculation cannot see). Actual deposit <= ceiling, often far less. The safe direction is rigorous: ceiling \textasciitilde{} 0 => no driving force.

ALGORITHM (visible in the trace): augmented Newton in (ln m\_j [, ln a\_H]) plus one raw-linear n\_p per ACTIVE mineral; each master mole balance gains a sink -sum nu\_pj n\_p; one SI = 0 row per active phase; an OUTER active-set loop ADMITS the most-supersaturated allowed mineral and EVICTS any with n\_p < 0, re-solving between, until complementarity holds (SI <= 0 with n = 0, OR SI = 0 with n > 0). A carbonate mineral's H+ leg lands the released H+ in the electroneutrality row (pH solve) -- so precipitating calcite ACIDIFIES.

CROSS-CHECK TOOL: PHREEQC EQUILIBRIUM\_PHASES (same ceiling framing). Data: curated USGS PHREEQC records (public domain) resolved from data/standards/ at runtime; self-DESCRIBING (inline manifest) -- seal with bin/choupo-import for a runtime-self-contained copy.

GOLDEN PINS (locked by the golden master): pin\_gypsum\_amount Ca = SO4 = 0.030542 mol/kg (DOUBLE the gypsum-saturation value 0.015271) + equilibrate gypsum $\to$ SI\_gypsum = 0.000 and n\_gypsum = 0.015271 (the machine returns EXACTLY the over-added excess). feed\_eq calcite + gypsum allowed; calcite reaches SI = 0, releases H+ (pH 7.72 $\to$ \textasciitilde{}7.50, solved), gypsum stays dissolved. scan\_eq r 0 $\to$ 0.85 with feedFlow: calcite kg/day ceiling; gypsum n = 0 throughout while SIeq\_gypsum drifts BELOW the raw SI\_gypsum (calcite eats the shared Ca).

\end{tutnarrative}
\begin{tutgolden}
\goldrow{diag}{pin\_gypsum\_amount}{I}{0.0455084}
\goldrow{diag}{pin\_gypsum\_amount}{SI\_anhydrite}{-0.299583}
\goldrow{diag}{pin\_gypsum\_amount}{SI\_gypsum}{1.33618e-10}
\goldrow{diag}{pin\_gypsum\_amount}{aw}{0.99952}
\goldrow{diag}{pin\_gypsum\_amount}{n\_gypsum}{0.0152707}
\goldrow{diag}{pin\_gypsum\_amount}{pH}{7}
\goldrow{diag}{feed}{I}{0.0266002}
\goldrow{diag}{feed}{SI\_anhydrite}{-1.68532}
\goldrow{diag}{feed}{SI\_aragonite}{0.111093}
\goldrow{diag}{feed}{SI\_calcite}{0.225093}
\goldrow{diag}{feed}{SI\_gypsum}{-1.38593}
\goldrow{diag}{feed}{SI\_halite}{-5.35323}
\goldrow{diag}{feed}{SI\_sylvite}{-6.39779}
\goldrow{diag}{feed}{aw}{0.999302}
\goldmore{193}
\end{tutgolden}

\clearpage
\tutentry{rainwater\_air}{tutorials/props/electrolyte/rainwater\_air}
\begin{tutscheme}
\begin{tikzpicture}[>={Stealth[length=2mm]}, every node/.style={font=\scriptsize}]
  \node[draw, rounded corners=2pt, align=center, minimum width=3.1cm, minimum height=0.8cm, fill=black!4] at (0.0,-0.0) {\textbf{air25}\\[-1pt]{\tiny speciate}};
  \node[draw, rounded corners=2pt, align=center, minimum width=3.1cm, minimum height=0.8cm, fill=black!4] at (3.6,-0.0) {\textbf{air10}\\[-1pt]{\tiny speciate}};
\end{tikzpicture}
\end{tutscheme}
\tutwhat{WATER OPEN TO AIR: the generic multi-gas Henry pin. Pure water equilibrated with the atmosphere (pCO2 4.2e-4, pO2 0.2095, pN2 0.7808 atm) at 25 C and at 10 C. Each listed gas PINS its dissolved species by Henry's law and that master's mole balance is REPLACED by the pin -- three pins at once: CO2 chains through the carbonate family (CO2aq is a computed species, the HCO3 row is swapped), while O2 and N2 are inert neutral masters pinned directly. TWO CLASSIC NUMBERS FALL OUT: (1) unpolluted RAINWATER pH \textasciitilde{} 5.6 -- carbonic acid from 420 ppm CO2 alone, no pollutant needed; (2) DISSOLVED-OXYGEN saturation \textasciitilde{} 8.5 mg/L at 25 C rising to \textasciitilde{} 11 mg/L at 10 C -- gas solubility is RETROGRADE with temperature (K\_H analytic from USGS PHREEQC PHASES, public domain), which is why cold streams hold more oxygen than warm ones. The `totals` entries for HCO3 / O2 / N2 are initial guesses only (announced): in an OPEN system each pinned total is a SOLVED OUTCOME of gas-liquid equilibrium.}
\begin{tutnarrative}
water open to AIR: the generic multi-gas Henry pin.

Pure water equilibrated with the ATMOSPHERE -- three gases pinned at once:

atmosphere \{ pCO2 4.2e-4 atm; O2 0.2095 atm; N2 0.7808 atm; \}

Each gas pins its dissolved species by Henry's law, a = K\_H(T) $\ast$ p, and the pinned master's mole balance is REPLACED by the pin (the OPEN system of SpeciationSolver, generalised beyond CO2). Two pin shapes appear side by side, which is the point of the case:

CO2 dissolves into a COMPUTED species (CO2aq, formed from the HCO3 master): the pin swaps the HCO3 row and the whole carbonate family (CO2aq / HCO3- / CO3--) equilibrates with the gas -- DIC is a solved OUTCOME; O2,N2 dissolve into INERT NEUTRAL MASTERS (species/aqueous O2, N2 -- identity records, charge 0): the pin replaces each master's own row directly -- the dissolved concentration is set by the gas phase alone.

THE TWO CLASSIC NUMBERS (textbook checks, no fitting anywhere):

pH \textasciitilde{} 5.6 unpolluted rainwater: carbonic acid from 420 ppm CO2 only (pH solve; -- H+ from electroneutrality against HCO3-); DO 8.5 $\to$ 11 mg/L dissolved-oxygen saturation, 25 C $\to$ 10 C: K\_H(T) analytic (USGS PHREEQC PHASES) makes gas solubility RETROGRADE -- cold water holds more oxygen. N2 saturates near 14 mg/L (25 C) for the same reason.

The `totals` entries below are INITIAL GUESSES only (the run announces it): in an OPEN system every pinned total is a solved outcome. The feed charge-imbalance warning is EXPECTED here -- an HCO3 guess with no cation is not a measured analysis; the solved H+ is its physical counter-ion.

Diagnostics locked by the golden master: pH, I, m\_CO2aq, and the dissolved O2 / N2 molalities at both temperatures.

\end{tutnarrative}
\begin{tutgolden}
\goldrow{diag}{air25}{I}{2.52807e-06}
\goldrow{diag}{air25}{aw}{0.999986}
\goldrow{diag}{air25}{gamma\_CO2aq}{1}
\goldrow{diag}{air25}{gamma\_OH}{0.998138}
\goldrow{diag}{air25}{m\_CO2aq}{1.42971e-05}
\goldrow{diag}{air25}{m\_OH}{3.97038e-09}
\goldrow{diag}{air25}{pH}{5.59803}
\goldrow{diag}{air10}{I}{2.79349e-06}
\goldrow{diag}{air10}{aw}{0.999981}
\goldrow{diag}{air10}{gamma\_CO2aq}{1}
\goldrow{diag}{air10}{gamma\_OH}{0.998089}
\goldrow{diag}{air10}{m\_CO2aq}{2.2591e-05}
\goldrow{diag}{air10}{m\_OH}{1.04436e-09}
\goldrow{diag}{air10}{pH}{5.55469}
\goldmore{12}
\end{tutgolden}

\clearpage
\tutentry{scaling\_ro\_brackish}{tutorials/props/electrolyte/scaling\_ro\_brackish}
\begin{tutscheme}
\begin{tikzpicture}[>={Stealth[length=2mm]}, every node/.style={font=\scriptsize}]
  \node[draw, rounded corners=2pt, align=center, minimum width=3.1cm, minimum height=0.8cm, fill=black!4] at (0.0,-0.0) {\textbf{pin\_kw}\\[-1pt]{\tiny speciate}};
  \node[draw, rounded corners=2pt, align=center, minimum width=3.1cm, minimum height=0.8cm, fill=black!4] at (3.6,-0.0) {\textbf{pin\_gypsum}\\[-1pt]{\tiny speciate}};
  \node[draw, rounded corners=2pt, align=center, minimum width=3.1cm, minimum height=0.8cm, fill=black!4] at (7.2,-0.0) {\textbf{pin\_neutral}\\[-1pt]{\tiny speciate}};
  \node[draw, rounded corners=2pt, align=center, minimum width=3.1cm, minimum height=0.8cm, fill=black!4] at (0.0,-1.3) {\textbf{feed}\\[-1pt]{\tiny speciate}};
  \node[draw, rounded corners=2pt, align=center, minimum width=3.1cm, minimum height=0.8cm, fill=black!4] at (3.6,-1.3) {\textbf{scan}\\[-1pt]{\tiny scalingScan}};
  \node[draw, rounded corners=2pt, align=center, minimum width=3.1cm, minimum height=0.8cm, fill=black!4] at (7.2,-1.3) {\textbf{scan\_open}\\[-1pt]{\tiny scalingScan}};
\end{tikzpicture}
\end{tutscheme}
\tutwhat{RO membrane-scaling audit of a brackish groundwater: aqueous speciation (Davies) with pH SOLVED from electroneutrality (the feed charge imbalance is announced first) + saturation indices SI = log10(IAP/K), scanned over water recovery to r = 0.85 (concentrate = feed/(1-r)) TWICE: closed system (DIC conserved and concentrating) vs open to the atmosphere (a(CO2aq) pinned by Henry at pCO2 = 4.2e-4 atm, DIC a solved outcome -- this CO2-rich water degasses, pH and calcite SI jump above the closed curve) -- the calcite SI curves differ; that contrast is the lesson. Three analytic pins lock the physics: pure water at given pH 7 reproduces Kw (m\_OH \textasciitilde{} 1e-7), a back-computed gypsum-saturated CaSO4 solution gives SI\_gypsum \textasciitilde{} 0.00, and pure water with pH solved lands at exactly 7.000.}
\begin{tutnarrative}
will my RO concentrate scale?

A REPRESENTATIVE brackish groundwater (generic textbook-style composition, \textasciitilde{}1500 mg/L TDS; no licensed table reproduced):

ion mg/L mol/kg w ion mg/L mol/kg w Ca+2 84 0.0021 Cl- 510 0.0144 Mg+2 27 0.0011 SO4-2 250 0.0026 Na+ 363 0.0158 HCO3- 183 0.0030 K+ 12 0.0003 lab pH 7.8 (nominal -- see below)

(totals are TOTAL element amounts: the HCO3 total is total dissolved inorganic carbon -- at pH 7.8 the reported HCO3 is \textasciitilde{}97\% of it. Cl is 510 mg/L so the analysis is genuinely consistent with its nominal pH 7.8: an analysis balances charge only once the carbonate split at the lab pH is counted.)

pH IS SOLVED, NOT GIVEN (`pH solve;`): H+ joins the Newton unknowns and ELECTRONEUTRALITY sum z\_i m\_i = 0 closes the system. Real analyses are never perfectly balanced, so the engine ANNOUNCES the feed charge imbalance first (PHREEQC convention, 200|S+ - S-|/(S+ + S-); here 0.35\%) -- the solved pH absorbs exactly that error, and above \textasciitilde{}5\% the solver warns loudly that it is not trustworthy. This feed solves to pH \textasciitilde{}7.72, near the lab's nominal 7.8: the analysis is consistent.

UNITS: every `totals` entry carries its basis (bare numbers are refused). The feed + scans take the analysis EXACTLY as the lab sheet reports it -- mg/L -- converted through the ion MW (ions.dat) at rho \textasciitilde{} 1 kg/L (dilute aqueous, announced once per run). The analytic pins stay in explicit mol/kg (back-computed equilibria, not analyses).

THE CLOSED-vs-OPEN LESSON (two scans, same feed): scan CLOSED system -- every total (DIC included) concentrates as feed/(1-r); the charge-balance ratio is scale-invariant, so the pH only drifts gently (7.72 $\to$ 7.58, the Davies gammas at work) while calcite SI climbs 0.23 $\to$ 1.38. scan\_open OPEN to the atmosphere (`atmosphere \{ pCO2 4.2e-4 atm; \}`) -- a(CO2aq) is PINNED by Henry, the HCO3 mole balance is REPLACED by the pin, and DIC becomes a solved OUTCOME. This groundwater is \textasciitilde{}8x CO2-supersaturated vs the atmosphere, so opening it DEGASES CO2: the pH jumps to 8.54 and climbs to 9.08, and calcite SI lands ABOVE the closed curve (2.59 vs 1.38 at r = 0.85) -- why concentrate pH control (acid dosing) matters. Compare the two CSVs (both carry a pH column right after I); the contrast IS the lesson.

VALIDATION PINS (locked by the golden master): pin\_kw pure water, pH 7 GIVEN $\to$ m\_OH = Kw a\_w/(a\_H gamma\_OH) \textasciitilde{} 1.0004e-7 (unchanged) pin\_gypsum totals Ca = SO4 = 0.015271 mol/kg, pH 7 GIVEN, BACK-COMPUTED from the catalogue's logK25(gypsum) = -4.55 + CaSO4aq logK25 = 2.14 + Davies at the resulting I = 0.0455 (26\% ion pairing) so that SI\_gypsum = 0.00 +- 0.001 (unchanged) pin\_neutral pure water, pH SOLVED $\to$ electroneutrality m\_H = m\_OH must give pH 7.000 (locks the new charge-balance closure)

Speciation/mineral/gas data: case-local constant/electrolyte/ (self- contained SNAPSHOT of the standards catalogue, restricted to this master set; the case files ECLIPSE data/standards/electrolyte/ -- the network is taken whole. Provenance: USGS PHREEQC, public domain -- see the .dat headers; the CO2(g) Henry constant is the db's PHASES CO2(g) entry).

\end{tutnarrative}
\begin{tutgolden}
\goldrow{diag}{pin\_kw}{I}{1.00037e-07}
\goldrow{diag}{pin\_kw}{aw}{1}
\goldrow{diag}{pin\_kw}{m\_OH}{1.00037e-07}
\goldrow{diag}{pin\_kw}{pH}{7}
\goldrow{diag}{pin\_gypsum}{I}{0.0455076}
\goldrow{diag}{pin\_gypsum}{SI\_anhydrite}{-0.299593}
\goldrow{diag}{pin\_gypsum}{SI\_gypsum}{-1.04349e-05}
\goldrow{diag}{pin\_gypsum}{aw}{0.99952}
\goldrow{diag}{pin\_gypsum}{m\_CaSO4aq}{0.0038941}
\goldrow{diag}{pin\_gypsum}{pH}{7}
\goldrow{diag}{pin\_neutral}{I}{1.00037e-07}
\goldrow{diag}{pin\_neutral}{aw}{1}
\goldrow{diag}{pin\_neutral}{pH}{7}
\goldrow{diag}{feed}{I}{0.0266002}
\goldmore{244}
\end{tutgolden}

\clearpage
\tutentry{softener\_brackish}{tutorials/props/electrolyte/softener\_brackish}
\begin{tutscheme}
\begin{tikzpicture}[>={Stealth[length=2mm]}, every node/.style={font=\scriptsize}]
  \node[draw, rounded corners=2pt, align=center, minimum width=3.1cm, minimum height=0.8cm, fill=black!4] at (0.0,-0.0) {\textbf{raw}\\[-1pt]{\tiny scalingScan}};
  \node[draw, rounded corners=2pt, align=center, minimum width=3.1cm, minimum height=0.8cm, fill=black!4] at (3.6,-0.0) {\textbf{softened}\\[-1pt]{\tiny exchange}};
  \node[draw, rounded corners=2pt, align=center, minimum width=3.1cm, minimum height=0.8cm, fill=black!4] at (7.2,-0.0) {\textbf{softened\_scaling}\\[-1pt]{\tiny scalingScan}};
\end{tikzpicture}
\end{tutscheme}
\tutwhat{Ion-exchange softening of a brackish RO feed: removes the calcite-scaling driver (raw vs softened vs softened-scaling; full narrative in propsDict).}
\begin{tutnarrative}
ion-exchange SOFTENING removes the scaling driver.

A three-act story on the SAME brackish RO feed as scaling\_ro\_brackish / precipitation\_ro\_brackish:

(1) raw scalingScan of the raw water: calcite SI > 0 already at the feed and climbing with recovery -- the SCALING PROBLEM (a CaCO3-supersaturated concentrate). (2) softened the `exchange` op equilibrates the SAME water with a strong-acid cation resin (SAC\_Na, Gaines-Thomas). Ca and Mg drop \textasciitilde{}99\% onto the resin; Na rises eq-for-eq (the salt penalty, printed in meq/L); hardness\_removal\_pct \textasciitilde{} 99\%. (3) softened\_scaling feed the SOFTENED composition (Ca/Mg gone, Na up) into a scalingScan: SI\_calcite is now NEGATIVE across the whole recovery range -- the scaling DRIVER is gone. Removing the hardness, not chasing the supersaturation, is the cure.

MODEL (visible in the verbosity>=3 trace): the SAME log-space Newton as speciate, PLUS one master (the free resin site X-), one balance (CEC capacity = sum of equivalents on the sites), and the exchange half-reactions Me(z+) + z X- = MeX (Gaines-Thomas equivalent-fraction activity for the bound species, Davies for the aqueous ions). Divalent Ca2+ (logK 0.8) is equivalent-fraction-favoured and outcompetes Na+; watch the per-iteration selectivity ratio beta(CaX2)/beta(NaX)\textasciicircum{}2 climb.

CROSS-CHECK TOOL: PHREEQC EXCHANGE (same Gaines-Thomas half-reactions).

THE FRESH-RESIN CAVEAT (printed loudly, non-suppressible): the `exchange` op returns the LIMITING (fully-equilibrated) effluent -- the best a fresh / fully-loaded contactor does per pass. A real fixed bed leaks rising hardness as the exchange front migrates to breakthrough, then must be regenerated; cycle length and breakthrough are TRANSIENT and not modelled here.

Brackish water (I \textasciitilde{} 0.03-0.06): Davies-valid, no Pitzer needed; the honesty banner + the I>0.5 Davies advisory cover the boundary.

Data: curated USGS PHREEQC records (public domain) resolved from data/standards/ at runtime, + the SAC\_Na resin nameplate; self-DESCRIBING (inline manifest) -- seal with bin/choupo-import for a runtime-self-contained copy.

\end{tutnarrative}
\begin{tutgolden}
\goldrow{diag}{raw}{I\_r0}{0.0266002}
\goldrow{diag}{raw}{I\_rmax}{0.165631}
\goldrow{diag}{raw}{SI\_anhydrite\_r0}{-1.68532}
\goldrow{diag}{raw}{SI\_anhydrite\_rmax}{-0.646016}
\goldrow{diag}{raw}{SI\_aragonite\_r0}{0.111093}
\goldrow{diag}{raw}{SI\_aragonite\_rmax}{1.26704}
\goldrow{diag}{raw}{SI\_calcite\_r0}{0.225093}
\goldrow{diag}{raw}{SI\_calcite\_rmax}{1.38104}
\goldrow{diag}{raw}{SI\_gypsum\_r0}{-1.38593}
\goldrow{diag}{raw}{SI\_gypsum\_rmax}{-0.349983}
\goldrow{diag}{raw}{SI\_halite\_r0}{-5.35323}
\goldrow{diag}{raw}{SI\_halite\_rmax}{-3.82475}
\goldrow{diag}{raw}{SI\_sylvite\_r0}{-6.39779}
\goldrow{diag}{raw}{SI\_sylvite\_rmax}{-4.87569}
\goldmore{193}
\end{tutgolden}

\clearpage
\tutentry{tartaricAcid\_acidulation}{tutorials/props/electrolyte/tartaricAcid\_acidulation}
\begin{tutscheme}
\begin{tikzpicture}[>={Stealth[length=2mm]}, every node/.style={font=\scriptsize}]
  \node[draw, rounded corners=2pt, align=center, minimum width=3.1cm, minimum height=0.8cm, fill=black!4] at (0.0,-0.0) {\textbf{acidulation\_pitzer}\\[-1pt]{\tiny speciate}};
  \node[draw, rounded corners=2pt, align=center, minimum width=3.1cm, minimum height=0.8cm, fill=black!4] at (3.6,-0.0) {\textbf{acidulation\_davies}\\[-1pt]{\tiny speciate}};
  \node[draw, rounded corners=2pt, align=center, minimum width=3.1cm, minimum height=0.8cm, fill=black!4] at (7.2,-0.0) {\textbf{gypsum\_ceiling}\\[-1pt]{\tiny speciate}};
  \node[draw, rounded corners=2pt, align=center, minimum width=3.1cm, minimum height=0.8cm, fill=black!4] at (0.0,-1.3) {\textbf{pKa2\_davies}\\[-1pt]{\tiny speciate}};
  \node[draw, rounded corners=2pt, align=center, minimum width=3.1cm, minimum height=0.8cm, fill=black!4] at (3.6,-1.3) {\textbf{pKa2\_pitzer}\\[-1pt]{\tiny speciate}};
\end{tikzpicture}
\end{tutscheme}
\tutwhat{Tartaric-acid acidulation liquor: rigorous Pitzer brine chemistry (Ca-SO4-K-Cl) + the gypsum precipitation ceiling, with tartrate honestly on the Debye-Hueckel limit (no published Pitzer params). Davies vs Pitzer shown side by side.}
\begin{tutgolden}
\goldrow{diag}{acidulation\_pitzer}{I}{1.6414}
\goldrow{diag}{acidulation\_pitzer}{SI\_anhydrite}{1.07386}
\goldrow{diag}{acidulation\_pitzer}{SI\_arcanite}{-2.80366}
\goldrow{diag}{acidulation\_pitzer}{SI\_gypsum}{1.36221}
\goldrow{diag}{acidulation\_pitzer}{SI\_sylvite}{-3.52403}
\goldrow{diag}{acidulation\_pitzer}{aw}{0.986675}
\goldrow{diag}{acidulation\_pitzer}{gamma\_Ca}{0.0716778}
\goldrow{diag}{acidulation\_pitzer}{gamma\_H2Tart}{1}
\goldrow{diag}{acidulation\_pitzer}{gamma\_HTart}{0.409404}
\goldrow{diag}{acidulation\_pitzer}{gamma\_SO4}{0.0662673}
\goldrow{diag}{acidulation\_pitzer}{gamma\_Tart}{0.036756}
\goldrow{diag}{acidulation\_pitzer}{m\_Ca}{0.397714}
\goldrow{diag}{acidulation\_pitzer}{m\_H2Tart}{0.0656066}
\goldrow{diag}{acidulation\_pitzer}{m\_HTart}{0.0314016}
\goldmore{74}
\end{tutgolden}

\subsection{estimate}
\grouptheory{estimate}
\clearpage
\tutentry{estimate\_acetone}{tutorials/props/estimate/estimate\_acetone}
\begin{tutscheme}
\begin{tikzpicture}[>={Stealth[length=2mm]}, every node/.style={font=\scriptsize}]
  \node[draw, rounded corners=2pt, align=center, minimum width=3.1cm, minimum height=0.8cm, fill=black!4] at (0.0,-0.0) {\textbf{estimate\_acetone}\\[-1pt]{\tiny estimateComponent}};
\end{tikzpicture}
\end{tutscheme}
\tutwhat{CREATE a component from its groups: Joback estimate of acetone (2 CH3 + 1 ketone), validated against NIST --- see the estimate built up from groups}
\begin{tutnarrative}
The FIRST stage of the props lifecycle (CREATE $\to$ curate $\to$ wire $\to$ simulate): build a pure-component property set from its molecular GROUPS by Joback group contribution, and SEE each group's contribution add up.

Acetone CH3-CO-CH3 = 2 x CH3 + 1 x ketone (>C=O).

The reference block below validates FOUR numbers -- Tb, Tc, Pc, omega -- and no others: Joback lands 2.2 \% low on Tb, 1.5 \% low on Tc, 2.2 \% high on Pc (read off the golden against the block, 2026-09-02), and you SEE the error, the whole point of "estimation, then trust". This header used to claim "\textasciitilde{}2 \% on Tc/Pc/Vc and \textasciitilde{}0.3 \% on dHf"; the block carries no reference Vc and no reference dHf, so those two were claims nothing in the case checks. omega comes from (Tb,Tc,Pc) by Lee-Kesler, not a group sum.

\end{tutnarrative}
\tutreadme{Acetone is 2 $\times$ CH3 + 1 $\times$ ketone. This case runs under choupoProps and asks Joback's group-contribution method to create a property set from those three groups, then lays the estimate beside the reference values the case itself declares. The golden: Tb 322.11 K, Tc 500.56 K, Pc 48.02 bar, ? 0.30, Vc 209.5 cm$^3$/mol, $\Delta$Hf -217.83 kJ/mol.}
\tutreadme{1. Estimation is a resolution problem, not a runtime one. Choupo does not estimate properties while a flowsheet runs; it estimates them HERE, at curation time, so a student can look at the number before trusting it. output \{ proposal auto; \} writes the estimate as a record you can read, edit, and only then wire into a case --- the run prints where: constant/components/acetone.estimated.dat, with the GAPS the method left listed in its header. 2. The error is the point. Against the reference block in system/propsDict (Tb 329.2 K, Tc 508.1 K, Pc 47.0 bar, ? 0.307): Tb is 2.2 \% low, Tc 1.5 \% low, Pc 2.2 \% high --- Joback's usual band for a small ketone, and you can see it rather than be told. 3. ? is not a group sum. It comes from (Tb, Tc, Pc) through Lee-Kesler; an error in Tb propagates into it. That chain is printed in the run's log. 4. A reference is a citation, not an endorsement. The block names where its four numbers were read; the run compares against those four and nothing else.}
\tutreadme{Open estimate\_ethanol\_benzene: the same method on a hydrogen-bonding molecule, where Joback's Tb goes noticeably wrong --- the honest limit of a group method, shown rather than warned about.}
\begin{tutgolden}
\goldrow{diag}{estimate\_acetone}{Cp298}{74.969755157}
\goldrow{diag}{estimate\_acetone}{Hvap\_kJmol}{27.534}
\goldrow{diag}{estimate\_acetone}{MW}{58.08}
\goldrow{diag}{estimate\_acetone}{Pc\_bar}{48.0249960499}
\goldrow{diag}{estimate\_acetone}{Psat\_298\_bar}{0.387196368809}
\goldrow{diag}{estimate\_acetone}{Psat\_Tb\_bar}{0.97953832278}
\goldrow{diag}{estimate\_acetone}{Tb\_K}{322.11}
\goldrow{diag}{estimate\_acetone}{Tc\_K}{500.559004953}
\goldrow{diag}{estimate\_acetone}{Vc\_cm3mol}{209.5}
\goldrow{diag}{estimate\_acetone}{Vliq298\_cm3mol}{82.1064043414}
\goldrow{diag}{estimate\_acetone}{dGf\_kJmol}{-154.54}
\goldrow{diag}{estimate\_acetone}{dHf\_kJmol}{-217.83}
\goldrow{diag}{estimate\_acetone}{omega}{0.29678401624}
\end{tutgolden}

\clearpage
\tutentry{estimate\_ethanol\_benzene}{tutorials/props/estimate/estimate\_ethanol\_benzene}
\begin{tutscheme}
\begin{tikzpicture}[>={Stealth[length=2mm]}, every node/.style={font=\scriptsize}]
  \node[draw, rounded corners=2pt, align=center, minimum width=3.1cm, minimum height=0.8cm, fill=black!4] at (0.0,-0.0) {\textbf{estimate\_ethanol}\\[-1pt]{\tiny estimateComponent}};
  \node[draw, rounded corners=2pt, align=center, minimum width=3.1cm, minimum height=0.8cm, fill=black!4] at (3.6,-0.0) {\textbf{estimate\_benzene}\\[-1pt]{\tiny estimateComponent}};
\end{tikzpicture}
\end{tutscheme}
\tutwhat{Joback estimates: ethanol (H-bonding $\to$ Tb weak, the honest limit) and benzene (aromatic groups $\to$ good) --- see where group contribution wins and loses}
\begin{tutnarrative}
Two Joback estimates, side by side, to SEE the method's spread:

ethanol (1 CH3 + 1 CH2 + 1 OH) -- a strong H-bonder; Joback's Tb comes out \textasciitilde{}14 K LOW (337 vs 351 K). The honest limit, in plain sight. benzene (6 arCH) -- aromatic ring; Joback Tc/Tb within \textasciitilde{}1.5\%. Group contribution at its best.

The lesson is the CONTRAST: an estimate is a starting point whose error you must SEE, not a number to trust blindly.

\end{tutnarrative}
\begin{tutgolden}
\goldrow{diag}{estimate\_ethanol}{Cp298}{64.6209120993}
\goldrow{diag}{estimate\_ethanol}{Hvap\_kJmol}{36.725}
\goldrow{diag}{estimate\_ethanol}{MW}{46.069}
\goldrow{diag}{estimate\_ethanol}{Pc\_bar}{57.5664143723}
\goldrow{diag}{estimate\_ethanol}{Psat\_298\_bar}{0.155662207837}
\goldrow{diag}{estimate\_ethanol}{Psat\_Tb\_bar}{0.972926042801}
\goldrow{diag}{estimate\_ethanol}{Tb\_K}{337.54}
\goldrow{diag}{estimate\_ethanol}{Tc\_K}{499.407373566}
\goldrow{diag}{estimate\_ethanol}{Vc\_cm3mol}{166.5}
\goldrow{diag}{estimate\_ethanol}{Vliq298\_cm3mol}{57.6954539184}
\goldrow{diag}{estimate\_ethanol}{dGf\_kJmol}{-170.86}
\goldrow{diag}{estimate\_ethanol}{dHf\_kJmol}{-236.84}
\goldrow{diag}{estimate\_ethanol}{omega}{0.573129252591}
\goldrow{diag}{estimate\_benzene}{Cp298}{81.8236158005}
\goldmore{12}
\end{tutgolden}

\clearpage
\tutentry{estimate\_petroleum\_cut}{tutorials/props/estimate/estimate\_petroleum\_cut}
\begin{tutscheme}
\begin{tikzpicture}[>={Stealth[length=2mm]}, every node/.style={font=\scriptsize}]
  \node[draw, rounded corners=2pt, align=center, minimum width=3.1cm, minimum height=0.8cm, fill=black!4] at (0.0,-0.0) {\textbf{estimate\_petroleum\_cut}\\[-1pt]{\tiny estimateComponent}};
\end{tikzpicture}
\end{tutscheme}
\tutwhat{ESTIMATE a petroleum PSEUDO-COMPONENT by name from (Tb, SG) -- Riazi-Daubert; see the SCREAMING ESTIMATE provenance and that the lump FLASHES}
\begin{tutnarrative}
CREATE a PETROLEUM PSEUDO-COMPONENT from the two bulk measurements every crude assay reports: the normal boiling point Tb and the specific gravity SG. The Riazi-Daubert (1987) correlations turn (Tb, SG) into the corresponding-states constants (MW, Tc, Pc, Vc); omega follows by Lee-Kesler, and the ideal-gas Cp by Kesler-Lee (1976).

This is the SCALAR-input sibling of the Joback group path: a petroleum cut is a LUMP of hundreds of species with NO single molecular structure, so it has NO group decomposition -- it is anchored on bulk (Tb, SG), not groups. The estimator is a registered ConstantEstimator (`model RiaziDaubert;`), chosen exactly like Joback -- the SAME explicit factory.

SEE in the console: the (Tb, SG) anchors and the Watson K factor; the constants built from them, validated against n-decane (a real species at this Tb/SG: Tc \textasciitilde{} exact, omega within \textasciitilde{}2 \%, Pc/MW within the correlation's stated \textasciitilde{}5-7 \%); the proposal .dat whose provenance SCREAMS status ESTIMATE / lumped true / "NOT a real species".

Promote with the printed `mv`, then it FLASHES (see the sibling tutorial steady/flash/flash\_pseudoComponent\_petCut, which ships this very lump and closes an energy balance with it).

\end{tutnarrative}
\begin{tutgolden}
\goldrow{diag}{estimate\_petroleum\_cut}{Cp298}{248.058663071}
\goldrow{diag}{estimate\_petroleum\_cut}{MW}{151.514068352}
\goldrow{diag}{estimate\_petroleum\_cut}{Pc\_bar}{19.7757194366}
\goldrow{diag}{estimate\_petroleum\_cut}{Psat\_Tb\_bar}{0.965291883924}
\goldrow{diag}{estimate\_petroleum\_cut}{SG}{0.73}
\goldrow{diag}{estimate\_petroleum\_cut}{Tb\_K}{447.3}
\goldrow{diag}{estimate\_petroleum\_cut}{Tc\_K}{617.804443469}
\goldrow{diag}{estimate\_petroleum\_cut}{Vc\_cm3mol}{627.45599905}
\goldrow{diag}{estimate\_petroleum\_cut}{Vliq298\_cm3mol}{203.429494662}
\goldrow{diag}{estimate\_petroleum\_cut}{omega}{0.481302966295}
\goldrow{diag}{estimate\_petroleum\_cut}{Psat\_298\_bar}{0.00181901199888}
\end{tutgolden}

\clearpage
\tutentry{petroleumCuts01\_riazi\_table36\_tbp}{tutorials/props/estimate/petroleumCuts01\_riazi\_table36\_tbp}
\begin{tutscheme}
\begin{tikzpicture}[>={Stealth[length=2mm]}, every node/.style={font=\scriptsize}]
  \node[draw, rounded corners=2pt, align=center, minimum width=3.1cm, minimum height=0.8cm, fill=black!4] at (0.0,-0.0) {\textbf{cuts5}\\[-1pt]{\tiny petroleumCuts}};
  \node[draw, rounded corners=2pt, align=center, minimum width=3.1cm, minimum height=0.8cm, fill=black!4] at (3.6,-0.0) {\textbf{cuts3}\\[-1pt]{\tiny petroleumCuts}};
\end{tikzpicture}
\end{tutscheme}
\tutwhat{CUT a TBP curve into petroleum PSEUDO-COMPONENTS (Riazi ASTM MNL50): Tb per cut, SG by constant Watson K, API and Watson K per cut and for the blend, one ESTIMATE record per cut}
\begin{tutnarrative}
CUT a TRUE-BOILING-POINT curve into PSEUDO-COMPONENTS. A crude fraction is a continuum of species; a process model needs a handful of lumps. The op `petroleumCuts` divides the curve into cuts, takes each cut's Tb at the mid-point of its distilled interval, gives it a specific gravity, and hands the pair (Tb, SG) to the SAME Riazi-Daubert estimator `estimateComponent` uses (MW, Tc, Pc, Vc; omega by Lee-Kesler; Cp by Kesler-Lee). One ESTIMATE record per cut is written into THIS case (constant/components/ petCut\_NN.estimated.dat, git-ignored: a proposal you review and promote with `mv`, never a standard record).

THE CURVE AND WHAT IS -- AND IS NOT -- VALIDATED. The six TBP points are the naphtha-kerosene blend printed in Riazi, "Characterization and Properties of Petroleum Fractions" (ASTM MNL50, 2005), Table 3.6, p. 103 (a TBP column, vol\% distilled vs C, quoted there from the API Technical Data Book). The book prints NO cut Tb, cut SG, MW or Tc for it: it uses the curve to illustrate the D86/TBP/EFV interconversions (Example 3.2), and there it estimates the blend's SG as 0.7862 from Eq. 3.17 -- the one printed number this case can reproduce, and does (feed SG 0.7862 from `specificGravityFrom TBP`). EVERYTHING ELSE IS ILLUSTRATIVE: the cut boundaries are this case's choice, and no number is validated against a distillation unit or a measured cut. The ASTM D86 and EFV columns of the same table are NOT used (the op refuses D86 input: the conversion is not implemented).

HOW EACH CUT IS CHARACTERISED (all equations are Riazi's, by number): Tb of a cut = the TBP temperature at the mid-point of its distilled interval (p. 83, footnote a: a mid-point temperature can serve as Tb for a narrow fraction), piecewise-linear between the curve points. SG of a cut = (1.8 Tb)\textasciicircum{}(1/3) / Kw with ONE Watson factor for the whole feed (Eq. 2.13 rearranged). THIS IS AN ASSUMPTION OF THE OP, NOT A RECOMMENDATION OF THE BOOK; `method declared;` lets you give the SG of each cut instead. Kw is Eq. 2.13 at the feed SG and the feed's MeABP (Eq. 3.3-3.5), solved by fixed point. API gravity (Eq. 2.4) and Watson K (Eq. 2.13) are published per cut and for the blend. CLOSURE: the cuts re-blended by volume (Eq. 3.45) must reproduce the declared feed SG within the DECLARED tolerance. A miss is announced, never adjusted. Both ops below re-blend about 1 \% above the feed SG (the curve is convex and the mid-point Tb is not the cut's mean); op `cuts5` declares 2 \% and closes, op `cuts3` declares 0.5 \% on purpose and prints the announcement. Note that the miss does NOT shrink monotonically with the number of cuts.

SEE: the cut table, the blend line, the closure line, the records.

\end{tutnarrative}
\begin{tutgolden}
\goldrow{diag}{cuts5}{closure\_SGdiff\_rel}{0.0100149194599}
\goldrow{diag}{cuts5}{closure\_SGtolerance}{0.02}
\goldrow{diag}{cuts5}{closure\_ok}{1}
\goldrow{diag}{cuts5}{cut01\_API}{70.6831235495}
\goldrow{diag}{cuts5}{cut01\_Cp298}{103.794999098}
\goldrow{diag}{cuts5}{cut01\_MW}{72.5545199844}
\goldrow{diag}{cuts5}{cut01\_Pc\_bar}{40.7350563483}
\goldrow{diag}{cuts5}{cut01\_SG}{0.69986058933}
\goldrow{diag}{cuts5}{cut01\_Tb\_K}{313.7}
\goldrow{diag}{cuts5}{cut01\_Tc\_K}{493.845044921}
\goldrow{diag}{cuts5}{cut01\_Vc\_cm3mol}{289.2234904}
\goldrow{diag}{cuts5}{cut01\_WatsonK}{11.8100246242}
\goldrow{diag}{cuts5}{cut01\_fraction}{0.1}
\goldrow{diag}{cuts5}{cut01\_massFraction}{0.0881217698017}
\goldmore{138}
\end{tutgolden}

\clearpage
\tutentry{userEstimate01\_linear\_tb}{tutorials/props/estimate/userEstimate01\_linear\_tb}
\begin{tutscheme}
\begin{tikzpicture}[>={Stealth[length=2mm]}, every node/.style={font=\scriptsize}]
  \node[draw, rounded corners=2pt, align=center, minimum width=3.1cm, minimum height=0.8cm, fill=black!4] at (0.0,-0.0) {\textbf{tb\_acetone}\\[-1pt]{\tiny linearTbEstimate}};
\end{tikzpicture}
\end{tutscheme}
\tutwhat{User-code property estimation: a student's OWN linear group-contribution Tb estimator (linearTbEstimate), compiled from code/ and linked into choupoProps.}
\begin{tutnarrative}
A student's OWN property-estimation op (linearTbEstimate), defined in code/ and linked into choupoProps by bin/buildCode. Tb = base + sum(count$\ast$incr). Here: 198.2 + 2$\ast$23.58 + 1$\ast$76.75 = 322.11 K (the Joback CH3,CH3,ketone sum). ---

\end{tutnarrative}
\begin{tutgolden}
\goldrow{diag}{tb\_acetone}{Tb\_K}{322.11}
\end{tutgolden}

\clearpage
\tutentry{vanKrevelen01\_polystyrene\_density}{tutorials/props/estimate/vanKrevelen01\_polystyrene\_density}
\begin{tutscheme}
\begin{tikzpicture}[>={Stealth[length=2mm]}, every node/.style={font=\scriptsize}]
  \node[draw, rounded corners=2pt, align=center, minimum width=3.1cm, minimum height=0.8cm, fill=black!4] at (0.0,-0.0) {\textbf{vanKrevelen01\_polystyrene\_density}\\[-1pt]{\tiny estimateComponent}};
\end{tikzpicture}
\end{tutscheme}
\tutwhat{ESTIMATE the density of polystyrene by Van Krevelen group contribution on the repeat unit -- glass-box additive sum, Vw from Bondi 1964 (licence-clean via UNIFAC R)}
\begin{tutnarrative}
ESTIMATE a POLYMER's repeat-unit DENSITY by group contribution -- the polymer sibling of the Joback small-molecule path. Polystyrene's repeat unit -[CH2-CH(C6H5)]- is decomposed into groups; each adds its molar mass and its Bondi van der Waals volume Vw; the sums give

M0 = sum n\_i MW\_i repeat-unit molar mass Vw = sum n\_i Vw\_i van der Waals volume (Bondi 1964) V = k $\ast$ Vw molar volume (k = packing factor) rho = M0 / V density

GLASS-BOX: the console prints every group's contribution -- you can redo the whole estimate by hand. M0 is an exact mass-balance fact (104.15 g/mol); rho carries the packing-factor uncertainty (the teaching point below).

LICENCE (the honest part): the Vw group values are Bondi, J. Phys. Chem. 68 (1964) 441 -- the PRIMARY per value -- derived from the UNIFAC R\_k volume parameters Choupo already ships (Vw = R\_k x 15.17), NOT copied from the Van Krevelen book table (copyright). See data/standards/parameters/vanKrevelen.dat. Slice 1 ships DENSITY only; Tg / solubility-parameter / Tm are DEFERRED (their open parameters are a separate, correctly-attributed thing).

THE TEACHING POINT -- the packing factor k: rho depends on k = V/Vw, the fraction of space the chains do NOT pack into. Default k = 1.60 (amorphous/glassy); crystalline packs denser (\textasciitilde{}1.43). k is YOURS: change `polymer \{ packing k; \}` and watch rho move. Atactic PS is amorphous (rho \textasciitilde{} 1.05) -- the estimate lands there; a crystalline polymer would need the smaller k. That k is visible (never hidden) is the lesson.

M0 = 104.15 also feeds the step-growth reactor KPIs (Mn = M0/(1-p)) -- the same repeat-unit mass, two places.

\end{tutnarrative}
\begin{tutgolden}
\goldrow{diag}{vanKrevelen01\_polystyrene\_density}{M0\_g\_per\_mol}{104.152}
\goldrow{diag}{vanKrevelen01\_polystyrene\_density}{V\_cm3\_per\_mol}{100.5616}
\goldrow{diag}{vanKrevelen01\_polystyrene\_density}{Vw\_cm3\_per\_mol}{62.851}
\goldrow{diag}{vanKrevelen01\_polystyrene\_density}{density\_g\_cm3}{1.0357034892}
\goldrow{diag}{vanKrevelen01\_polystyrene\_density}{packing\_k}{1.6}
\end{tutgolden}

\clearpage
\tutentry{yang2020\_Tg\_polystyrene}{tutorials/props/estimate/yang2020\_Tg\_polystyrene}
\begin{tutscheme}
\begin{tikzpicture}[>={Stealth[length=2mm]}, every node/.style={font=\scriptsize}]
  \node[draw, rounded corners=2pt, align=center, minimum width=3.1cm, minimum height=0.8cm, fill=black!4] at (0.0,-0.0) {\textbf{yang2020\_Tg\_polystyrene}\\[-1pt]{\tiny estimateComponent}};
\end{tikzpicture}
\end{tutscheme}
\tutwhat{ESTIMATE the glass-transition Tg of polystyrene by the Yang 2020 modified group-contribution scheme (CC-BY) -- the repeat unit -CH2-CH(phenyl)- gives Tg(inf)=444 K, an HONEST +19\% over-estimate of the measured \textasciitilde{}373 K (a teaching point about group-contribution limits)}
\begin{tutnarrative}
Yang-2020 Tg estimate + an HONEST limitation

ESTIMATE the GLASS-TRANSITION temperature Tg of polystyrene by the YANG 2020 modified group-contribution scheme. The repeat unit -CH2-CH(phenyl)- decomposes into two Yang main-chain groups:

CH2 : MW 14 Yg 4.026 CHphenyl: MW 90 Yg 42.153

M0 = 14 + 90 = 104 g/mol Yg(inf) = 4.026 + 42.153 = 46.179 (1e3 g.K/mol) Tg(inf) = 46.179$\ast$1e3 / 104 = 444 K

>>> THE HONEST TEACHING POINT <<< Measured polystyrene Tg \textasciitilde{} 373 K (atactic, high Mw). Yang 2020 gives 444 K -- a +19\% OVER-ESTIMATE. This is NOT a bug: it is the model telling the truth about itself. The CHphenyl contribution (Yg 42.153) was regressed across the whole training set, and the bulky monosubstituted phenyl pendant is one where an additive main-chain scheme over-predicts the limiting Tg. Box: "all models are wrong, some are useful". The student SEES the deviation (the console prints the per-cent error against the measured anchor) and learns where a group-contribution method earns trust (PVC: 351 vs 354 K, -0.8\%) and where it does not (polystyrene: 444 vs 373 K, +19\%).

Contrast the flagship `yang2020\_Tg\_pvc`, where Tg(inf)=351 K reproduces the paper's Fig 6 value EXACTLY -- that case goldens Tg as validated physics. Here Tg is SHOWN (the machinery is identical) but the case exists to expose the over-estimate, so only M0 (an exact mass-balance fact) is goldened, not Tg.

This is Yang 2020 -- NOT bare Van Krevelen. Nb (backbone-in-the-side-chain) = 0 for this vinyl polymer.

SOURCE + LICENCE: Yang, Zou, Ye, Zhu, Dong, Bi, ACS Omega 5 (2020) 19655, doi:10.1021/acsomega.0c04499 -- CC-BY 4.0; values cited per group as "Yang 2020", never "Van Krevelen".

\end{tutnarrative}
\begin{tutgolden}
\goldrow{diag}{yang2020\_Tg\_polystyrene}{M0\_g\_per\_mol}{104}
\goldrow{diag}{yang2020\_Tg\_polystyrene}{Tg\_K}{444}
\goldrow{diag}{yang2020\_Tg\_polystyrene}{YgSum\_1e3gKmol}{46.18}
\end{tutgolden}

\clearpage
\tutentry{yang2020\_Tg\_pvc}{tutorials/props/estimate/yang2020\_Tg\_pvc}
\begin{tutscheme}
\begin{tikzpicture}[>={Stealth[length=2mm]}, every node/.style={font=\scriptsize}]
  \node[draw, rounded corners=2pt, align=center, minimum width=3.1cm, minimum height=0.8cm, fill=black!4] at (0.0,-0.0) {\textbf{yang2020\_Tg\_pvc}\\[-1pt]{\tiny estimateComponent}};
\end{tikzpicture}
\end{tutscheme}
\tutwhat{ESTIMATE the glass-transition Tg of poly(vinyl chloride) (PVC) by the Yang 2020 modified group-contribution scheme (CC-BY) -- the repeat unit -CH2-CHCl- gives Tg(inf)=351 K, EXACT vs the paper (Fig 6)}
\begin{tutnarrative}
FLAGSHIP Yang-2020 glass-transition estimate

ESTIMATE the GLASS-TRANSITION temperature Tg of poly(vinyl chloride) (PVC) by the YANG 2020 modified group-contribution scheme -- the polymer sibling of the Joback small-molecule path, for Tg. The repeat unit -CH2-CHCl- is decomposed into two Yang main-chain groups; each adds its limiting Tg contribution Yg(inf):

M0 = sum n\_i MW\_i repeat-unit molar mass [g/mol] Yg(inf) = sum n\_i Yg\_i(inf) (+ Nb$\ast$Ygb) limiting Tg function [1e3 g.K/mol] Tg(inf) = Yg(inf) $\ast$ 1e3 / M0 Tg at infinite Mw [K]

THE ADDITIVE SUM (shown on the console): CH2 : MW 14 Yg 4.026 CHCl: MW 48.5 Yg 17.911

M0 = 14 + 48.5 = 62.5 g/mol Yg(inf) = 4.026 + 17.911 = 21.937 (1e3 g.K/mol) Tg(inf) = 21.937$\ast$1e3 / 62.5 = 351 K <-- EXACT vs Yang 2020, Fig 6

This is Yang 2020 -- NOT bare Van Krevelen. The scheme predicts the INFINITE- Mw limit Tg(inf); the molecular-weight dependence Tg(Mn) = Tg(inf) - K/Mn is a later slice. Nb (backbone-in-the-side-chain) = 0 for this vinyl polymer, so the bracket is just the sum of the two main-chain group contributions.

SOURCE + LICENCE (the honest part): Yang, Zou, Ye, Zhu, Dong, Bi, ACS Omega 5 (2020) 19655, doi:10.1021/acsomega.0c04499 -- CC-BY 4.0, so the parameter values ARE redistributable. Cited per value as "Yang 2020", never "Van Krevelen".

THE GOLDEN: M0 (an exact mass-balance fact) AND Tg(inf)=351 K are checked -- Tg here is VALIDATED PHYSICS (it reproduces the paper's Fig 6 value exactly), not just a mass balance. The measured anchor (PVC Tg \textasciitilde{} 354 K) confirms it.

\end{tutnarrative}
\begin{tutgolden}
\goldrow{diag}{yang2020\_Tg\_pvc}{M0\_g\_per\_mol}{62.5}
\goldrow{diag}{yang2020\_Tg\_pvc}{Tg\_K}{351}
\goldrow{diag}{yang2020\_Tg\_pvc}{YgSum\_1e3gKmol}{21.937}
\end{tutgolden}

\subsection{fit}
\grouptheory{fit}
\clearpage
\tutentry{cp01\_polynomial\_recovery}{tutorials/props/fit/cp01\_polynomial\_recovery}
\begin{tutscheme}
\begin{tikzpicture}[>={Stealth[length=2mm]}, every node/.style={font=\scriptsize}]
  \node[draw, rounded corners=2pt, align=center, minimum width=3.1cm, minimum height=0.8cm, fill=black!4] at (0.0,-0.0) {\textbf{cp\_dce\_liquid\_fit}\\[-1pt]{\tiny heatCapacityFit}};
\end{tikzpicture}
\end{tutscheme}
\tutwhat{Quadratic Cp fit on data generated from a quadratic -- exact recovery, with the Kirchhoff Cp\_ig - Cp\_liq cross-check that ties Cp to dHvap}
\begin{tutnarrative}
A LINEAR FIT, AND THE PHYSICS THAT CHECKS IT.

`heatCapacityFit` regresses Cp(T) = A0 + A1 T + A2 T\textasciicircum{}2 by linear least squares -- normal equations, no iteration. Fed data generated FROM a quadratic, it must return that quadratic, so the expected file pins the coefficients rather than a tolerance band.

The interesting number is not the fit. It is KIRCHHOFF:

d(dHvap)/dT = Cp\_ig - Cp\_liq

The op takes the FITTED liquid Cp and the component's CATALOGUE ideal-gas Cp and reports that difference at the mid temperature. It should be NEGATIVE -- the enthalpy of vaporisation falls as T rises, to zero at the critical point -- and its size is the slope of dHvap(T). That is the link tying Cp to dHvap to the Antoine fit next door (psat01\_antoine\_recovery, which reports dHvap from B): three ops, one consistency network, and a student can check one against the other by hand.

Why the case exists: `heatCapacityFit` had NO case and no gate. Writing the first one for its sibling vaporPressureFit segfaulted choupoProps twice and found a thermoPhysPropDict being silently ignored.

\end{tutnarrative}
\begin{tutgolden}
\goldrow{diag}{cp\_dce\_liquid\_fit}{A0}{194}
\goldrow{diag}{cp\_dce\_liquid\_fit}{A1}{-0.618063}
\goldrow{diag}{cp\_dce\_liquid\_fit}{A2}{0.00170318}
\goldrow{diag}{cp\_dce\_liquid\_fit}{Cp\_at\_Tmid}{166.076}
\goldrow{diag}{cp\_dce\_liquid\_fit}{Cp\_ig\_at\_Tmid}{84.2326}
\goldrow{diag}{cp\_dce\_liquid\_fit}{R2}{1}
\goldrow{diag}{cp\_dce\_liquid\_fit}{dHvap\_dT\_kirchhoff}{-81.8435}
\goldrow{diag}{cp\_dce\_liquid\_fit}{degree}{2}
\goldrow{diag}{cp\_dce\_liquid\_fit}{n\_points}{21}
\goldrow{verdict}{cp\_dce\_liquid\_fit}{verdict}{notClaimed}
\end{tutgolden}

\clearpage
\tutentry{psat01\_antoine\_recovery}{tutorials/props/fit/psat01\_antoine\_recovery}
\begin{tutscheme}
\begin{tikzpicture}[>={Stealth[length=2mm]}, every node/.style={font=\scriptsize}]
  \node[draw, rounded corners=2pt, align=center, minimum width=3.1cm, minimum height=0.8cm, fill=black!4] at (0.0,-0.0) {\textbf{psat\_water\_fit}\\[-1pt]{\tiny vaporPressureFit}};
\end{tikzpicture}
\end{tutscheme}
\tutwhat{Antoine fit on data generated from Choupo's own water correlation: the coefficients must come back exactly, so any deviation is the op's arithmetic}
\begin{tutnarrative}
FITTING ANTOINE, AND KNOWING THE ANSWER BEFOREHAND.

`vaporPressureFit` regresses log10(Psat[bar]) = A - B/(T + C) on measured Psat(T). The fit is linear in (A, B) at fixed C, so the op golden-sections C and solves (A, B) by least squares at each one -- derivative-free, and worth reading in the source.

This case feeds it data generated FROM THE SAME EQUATION, by evaluating Choupo's own curated water correlation at twenty temperatures. That is deliberate. Against laboratory data a fit can only be judged to the scatter of the measurements; against data the equation itself produced, the fit must RETURN THE COEFFICIENTS IT WAS GIVEN -- so any deviation is the op's arithmetic and nothing else, and the expected file can pin it to the digit.

Why the case exists: `vaporPressureFit` had NO case and no gate. It could have been broken for a year and the suite would have stayed green; the first to find out would be a student fitting their own data. The psychrometric chart was in exactly that position last week, and writing its first test found it reporting nothing.

What to read in the output: A, B, C -- the recovered coefficients, against 5.40221 / 1838.675 / -31.737; R2 -- 1 to floating point, because the data has no scatter; dHvap -- the Clausius-Clapeyron reading of B WITH the (T/(T+C))\textasciicircum{}2 correction, at the highest T. Compare it to water's tabulated enthalpy of vaporisation: this is the physical cross-check on a fit that is otherwise just a straight line through transformed numbers.

\end{tutnarrative}
\begin{tutgolden}
\goldrow{diag}{psat\_water\_fit}{A}{5.40221000142}
\goldrow{diag}{psat\_water\_fit}{B}{1838.67500085}
\goldrow{diag}{psat\_water\_fit}{C}{-31.7369999323}
\goldrow{diag}{psat\_water\_fit}{R2}{1}
\goldrow{diag}{psat\_water\_fit}{dHvap\_kJ\_per\_mol}{42.1162117513}
\goldrow{diag}{psat\_water\_fit}{n\_points}{20}
\goldrow{verdict}{psat\_water\_fit}{verdict}{notClaimed}
\end{tutgolden}

\subsection{gibbs}
\grouptheory{gibbs}
\clearpage
\tutentry{landscape01\_ammonia\_synthesis}{tutorials/props/gibbs/landscape01\_ammonia\_synthesis}
\begin{tutscheme}
\begin{tikzpicture}[>={Stealth[length=2mm]}, every node/.style={font=\scriptsize}]
  \node[draw, rounded corners=2pt, align=center, minimum width=3.1cm, minimum height=0.8cm, fill=black!4] at (0.0,-0.0) {\textbf{ammonia}\\[-1pt]{\tiny gibbsMap}};
\end{tikzpicture}
\end{tutscheme}
\tutwhat{Equilibrium landscape: ammonia synthesis, NH3 mole fraction against T at 1, 10, 100 bar -- and the 298 K / 700 K contrast}
\begin{tutnarrative}
THE AMMONIA LANDSCAPE -- N2 + 3 H2 at stoichiometric feed, equilibrium NH3 mole fraction against temperature at 1, 10 and 100 bar (the ideal-gas element-potential kernel, the same one nh3\_equilibrium\_map contours over the whole plane). This is the first witness of the EduTool `equilibrium-landscapes` (gui/src/ui/methods/EquilibriumLandscapesTool.tsx): a SURVEY over reaction systems the engine already solves, read live from this case's `map.csv` and its anchor KPIs.

WHAT THE THREE CURVES SHOW. Exothermic AND mole-reducing (4 $\to$ 2), so the equilibrium loves COLD and PRESSURE: every curve falls with T, and the 100 bar curve sits above the 1 bar one wherever the gas is dry. The grid starts at 300 K on purpose: from 300 to 380 K the 100 bar cells are a SATURATED gas -- the kernel finds NH3 supersaturated there and condenses it (ElementPotential.cpp, the one-condensable two-phase branch), and the mole fraction it reports is the GAS-phase Psat(T)/P. That is why the 100 bar curve RISES with T at its cold end: that leg is ammonia's vapour pressure, not its chemistry, and the curve joins the dry equilibrium at the dew point near 400 K (measured: 0.1047 at 300 K, 0.9311 at 400 K).

THE CONTRAST THE PAGE CLOSES ON, computed here and never typed. Two anchor cells: 298.15 K / 1 bar -- ambient, where the reaction is thermodynamically favourable and nothing in this tree can make it go (the only ammonia rate law here, Dyson \& Simon 1968, written in the cases that use it, was fitted over 150-300 atm and its effectiveness factor is a table at 150, 225 and 300 atm that the engine REFUSES to extrapolate: tutorials/props/kinetics/ ammoniaRate01\_dyson\_simon/constant/reactions) -- and 700 K / 200 bar, the gibbs10\_ammonia\_fugacity operating point, which this ideal-gas kernel answers as the SHAPE and that SRK case answers with the fugacity correction folded in. The third anchor, 700 K / 1 bar, is what the same chemistry reaches at converter temperature without the pressure.

Anchor KPI names carry the approach magnitude when one is declared (`metric\_T25C\_P1atm\_dT<n>`), so a changed approach breaks the golden loudly. Nothing here is validated against a measurement: the Larson \& Dodge (1923) pins live in nh3\_equilibrium\_map and are not repeated.

WHY THE PRESSURES ARE 1, 10 AND 100 BAR. `Pgrid` is `from / to / n / log` (GibbsMapOp.cpp reads no list form), and three log-spaced points from 1 to 100 bar land on the decades a textbook draws. Ideal-gas validity is the same caveat as nh3\_equilibrium\_map's: the high-pressure curve is the SHAPE of the pressure dependence, not a quantitative prediction.

temperatureApproach IS DECLARED, AT ZERO. The EduTool's slider replaces the NUMBER of a declared scalar and can add no key (gui/src/case/ methodRun.ts), so the key must be here for the page to move it. Zero is true equilibrium and prints nothing; a magnitude > 0 makes the run assign the sign at EVERY cell from the thermicity of the overall transformation at that cell's T (GibbsMapOp.cpp, `resolveDirection`) and print the tally.

\end{tutnarrative}

\clearpage
\tutentry{landscape02\_water\_gas\_shift}{tutorials/props/gibbs/landscape02\_water\_gas\_shift}
\begin{tutscheme}
\begin{tikzpicture}[>={Stealth[length=2mm]}, every node/.style={font=\scriptsize}]
  \node[draw, rounded corners=2pt, align=center, minimum width=3.1cm, minimum height=0.8cm, fill=black!4] at (0.0,-0.0) {\textbf{wgs}\\[-1pt]{\tiny gibbsMap}};
\end{tikzpicture}
\end{tutscheme}
\tutwhat{Equilibrium landscape: water-gas shift, CO conversion (carbon in CO2) against T -- pressure-indifferent}
\begin{tutnarrative}
THE SHIFT LANDSCAPE -- equimolar CO + H2O (gibbs01\_water\_gas\_shift's own feed and species set), the fraction of the feed carbon leaving as CO2 against temperature at 1, 10 and 100 bar. `elementYield` on carbon with CO2 as the product IS the CO conversion here, because every carbon atom fed arrives as CO -- and it is the forum-permitted metric (species "conversion" is banned where the species is also produced; docs/design/gibbs-map-forum-2026-07-02.md).

WHAT THE THREE CURVES SHOW, and it is one curve. Exothermic (the conversion falls with T) and mole-CONSERVING (2 $\to$ 2), so the three pressures give the SAME equilibrium to every printed digit: the ideal-gas equilibrium constant of a Delta\_n = 0 reaction carries no P. The page draws the three curves on top of each other and says why. MEASURED, not assumed: no cell on this grid condenses -- the three pressures agree to every printed digit at all 46 temperatures, which they could not if water had left the gas anywhere (at 400 K the equilibrium leaves 2.5 \% water in the gas, 2.5 bar at 100 bar against a 2.46 bar Psat -- and the run's own identity of the three columns is the proof, not this estimate).

WHY THE PRESSURES ARE 1, 10 AND 100 BAR. `Pgrid` is `from / to / n / log` (GibbsMapOp.cpp reads no list form), and three log-spaced points from 1 to 100 bar land on the decades a textbook draws. Ideal-gas validity is the same caveat as nh3\_equilibrium\_map's: the high-pressure curve is the SHAPE of the pressure dependence, not a quantitative prediction.

temperatureApproach IS DECLARED, AT ZERO. The EduTool's slider replaces the NUMBER of a declared scalar and can add no key (gui/src/case/ methodRun.ts), so the key must be here for the page to move it. Zero is true equilibrium and prints nothing; a magnitude > 0 makes the run assign the sign at EVERY cell from the thermicity of the overall transformation at that cell's T (GibbsMapOp.cpp, `resolveDirection`) and print the tally.

\end{tutnarrative}

\clearpage
\tutentry{landscape03\_steam\_reforming}{tutorials/props/gibbs/landscape03\_steam\_reforming}
\begin{tutscheme}
\begin{tikzpicture}[>={Stealth[length=2mm]}, every node/.style={font=\scriptsize}]
  \node[draw, rounded corners=2pt, align=center, minimum width=3.1cm, minimum height=0.8cm, fill=black!4] at (0.0,-0.0) {\textbf{reformer}\\[-1pt]{\tiny gibbsMap}};
\end{tikzpicture}
\end{tutscheme}
\tutwhat{Equilibrium landscape: methane steam reforming (S/C 3), carbon left in CH4 against T -- pressure hurts}
\begin{tutnarrative}
THE REFORMING LANDSCAPE -- CH4 + 3 H2O (gibbs02\_steam\_reforming's own feed and species set), the fraction of the feed carbon STILL in methane against temperature at 1, 10 and 100 bar. Read it as 1 minus the methane conversion: LESS is MORE converted, and the page labels it so rather than subtracting -- the metric on the axis is exactly the number the engine wrote. (`elementYield` with CH4 as the "product" is the honest way to ask for the unconverted feed; the metric slot accepts any listed species.)

WHAT THE THREE CURVES SHOW. Endothermic (the slip falls as T rises) and mole-INCREASING (CH4 + H2O $\to$ CO + 3 H2, 2 $\to$ 4), so pressure HURTS: the 100 bar curve sits far above the 1 bar one, which is why industrial reformers run hot and why the shift on the same catalyst (exothermic, Delta\_n = 0) is the one the approach caveat calls GLOBAL -- one dT for two reactions of opposite thermicity. The engine reads the OVERALL thermicity from this feed (endothermic at every cell of this grid, the run's own tally) and evaluates the chemistry at T - dT.

WHY THE PRESSURES ARE 1, 10 AND 100 BAR. `Pgrid` is `from / to / n / log` (GibbsMapOp.cpp reads no list form), and three log-spaced points from 1 to 100 bar land on the decades a textbook draws. Ideal-gas validity is the same caveat as nh3\_equilibrium\_map's: the high-pressure curve is the SHAPE of the pressure dependence, not a quantitative prediction.

temperatureApproach IS DECLARED, AT ZERO. The EduTool's slider replaces the NUMBER of a declared scalar and can add no key (gui/src/case/ methodRun.ts), so the key must be here for the page to move it. Zero is true equilibrium and prints nothing; a magnitude > 0 makes the run assign the sign at EVERY cell from the thermicity of the overall transformation at that cell's T (GibbsMapOp.cpp, `resolveDirection`) and print the tally.

\end{tutnarrative}

\clearpage
\tutentry{landscape04\_methane\_combustion}{tutorials/props/gibbs/landscape04\_methane\_combustion}
\begin{tutscheme}
\begin{tikzpicture}[>={Stealth[length=2mm]}, every node/.style={font=\scriptsize}]
  \node[draw, rounded corners=2pt, align=center, minimum width=3.1cm, minimum height=0.8cm, fill=black!4] at (0.0,-0.0) {\textbf{burner}\\[-1pt]{\tiny gibbsMap}};
\end{tikzpicture}
\end{tutscheme}
\tutwhat{Equilibrium landscape: methane combustion in stoichiometric air, carbon in CO2 against T -- complete until dissociation}
\begin{tutnarrative}
THE COMBUSTION LANDSCAPE -- CH4 in stoichiometric air (gibbs03's own feed and species set), the fraction of the feed carbon fully oxidised to CO2 against temperature at 1, 10 and 100 bar, 600 to 2500 K.

WHAT THE THREE CURVES SHOW, and it is where the approach is INDIFFERENT. Below about 1500 K the equilibrium is complete conversion to better than 0.2 \% (0.99996 at 600 K, 0.9987 at 1500 K and 1 bar) -- a flat line at 1 on any axis a plot can draw -- and a 50 K approach moves it by nothing visible: the whole conversion-versus-temperature curve is flat, so a HORIZONTAL shift along it changes nothing. Only above \textasciitilde{}1800 K, where CO2 starts to dissociate to CO, does the curve acquire a slope and the approach a visible effect, and pressure holds the CO2 together (the dissociation makes moles). This is the page's counterexample to the ammonia and reforming systems, where the same 50 K moves the answer by a great deal.

The grid starts at 600 K rather than 300 K by measurement: on a 300 K start the cells at 300 and 350 K (every pressure), 400 K / 10 bar and 500 K / 100 bar did NOT converge -- the product water is far below its dew point there and the gas-only Newton the kernel runs first does not reach a root -- and a witness that ships hatched cells at its cold end teaches nothing the 600 K plateau does not. From 600 K up every cell converges dry (y\_H2O P stays below Psat on the whole grid).

WHY THE PRESSURES ARE 1, 10 AND 100 BAR. `Pgrid` is `from / to / n / log` (GibbsMapOp.cpp reads no list form), and three log-spaced points from 1 to 100 bar land on the decades a textbook draws. Ideal-gas validity is the same caveat as nh3\_equilibrium\_map's: the high-pressure curve is the SHAPE of the pressure dependence, not a quantitative prediction.

temperatureApproach IS DECLARED, AT ZERO. The EduTool's slider replaces the NUMBER of a declared scalar and can add no key (gui/src/case/ methodRun.ts), so the key must be here for the page to move it. Zero is true equilibrium and prints nothing; a magnitude > 0 makes the run assign the sign at EVERY cell from the thermicity of the overall transformation at that cell's T (GibbsMapOp.cpp, `resolveDirection`) and print the tally.

\end{tutnarrative}

\clearpage
\tutentry{landscape05\_thermal\_no}{tutorials/props/gibbs/landscape05\_thermal\_no}
\begin{tutscheme}
\begin{tikzpicture}[>={Stealth[length=2mm]}, every node/.style={font=\scriptsize}]
  \node[draw, rounded corners=2pt, align=center, minimum width=3.1cm, minimum height=0.8cm, fill=black!4] at (0.0,-0.0) {\textbf{burner}\\[-1pt]{\tiny gibbsMap}};
\end{tikzpicture}
\end{tutscheme}
\tutwhat{Equilibrium landscape: thermal NO in methane-air combustion products against T -- the GLOBAL caveat, drawn}
\begin{tutnarrative}
THE THERMAL-NO LANDSCAPE -- CH4 in stoichiometric air with NO in the species list (gibbs06\_thermal\_nox's own feed and species set), the equilibrium NO mole fraction against temperature at 1, 10 and 100 bar, 1000 to 2500 K. The page draws it on a LOG axis: it spans many decades (read the two ends off the CSV rather than off this sentence).

WHAT IT IS HERE FOR, and it is the sharpest thing on the page. NO formation (N2 + O2 $\to$ 2 NO) is ENDOTHERMIC -- the curve RISES with T -- but it rides inside a feed whose OVERALL transformation, methane burning, is exothermic. The engine reads the direction from the overall thermicity of THIS feed at each cell (GibbsReactor::approachDirection, the one home; GibbsMapOp calls it per cell), finds every cell exothermic, and evaluates the chemistry at T + dT. For the combustion that is the conservative direction. For the NO riding inside it, T + dT sits HIGHER on a rising curve: the declared approach reports MORE NO than the equilibrium at the physical T -- an outlet beyond equilibrium for the minor product. That is the engine's own GLOBAL caveat ("it cannot resolve per-reaction approaches") drawn as a curve rather than printed as a sentence, and the page says it is a limitation of the single-number model, not a finding about flames.

Pressure: 2 $\to$ 2 for NO formation, so the 1 and 10 bar curves nearly coincide; the small spread that remains comes from the combustion products' own dissociation equilibria, which do carry P.

WHY THE PRESSURES ARE 1, 10 AND 100 BAR. `Pgrid` is `from / to / n / log` (GibbsMapOp.cpp reads no list form), and three log-spaced points from 1 to 100 bar land on the decades a textbook draws. Ideal-gas validity is the same caveat as nh3\_equilibrium\_map's: the high-pressure curve is the SHAPE of the pressure dependence, not a quantitative prediction.

temperatureApproach IS DECLARED, AT ZERO. The EduTool's slider replaces the NUMBER of a declared scalar and can add no key (gui/src/case/ methodRun.ts), so the key must be here for the page to move it. Zero is true equilibrium and prints nothing; a magnitude > 0 makes the run assign the sign at EVERY cell from the thermicity of the overall transformation at that cell's T (GibbsMapOp.cpp, `resolveDirection`) and print the tally.

\end{tutnarrative}

\clearpage
\tutentry{landscape06\_nitrogen\_fixation\_routes}{tutorials/props/gibbs/landscape06\_nitrogen\_fixation\_routes}
\begin{tutscheme}
\begin{tikzpicture}[>={Stealth[length=2mm]}, every node/.style={font=\scriptsize}]
  \node[draw, rounded corners=2pt, align=center, minimum width=3.1cm, minimum height=0.8cm, fill=black!4] at (0.0,-0.0) {\textbf{routes298}\\[-1pt]{\tiny reactionGibbs}};
  \node[draw, rounded corners=2pt, align=center, minimum width=3.1cm, minimum height=0.8cm, fill=black!4] at (3.6,-0.0) {\textbf{routes310}\\[-1pt]{\tiny reactionGibbs}};
  \node[draw, rounded corners=2pt, align=center, minimum width=3.1cm, minimum height=0.8cm, fill=black!4] at (7.2,-0.0) {\textbf{haber700}\\[-1pt]{\tiny reactionGibbs}};
\end{tikzpicture}
\end{tutscheme}
\tutwhat{Nitrogen fixation, three overall routes on one formation surface: heterotrophic (sugar as the electron donor), phototrophic (water), and Haber-Bosch -- dG as the engine prices it}
\begin{tutnarrative}
THREE OVERALL ROUTES TO FIXED NITROGEN, PRICED ON ONE SURFACE. The sixth witness of the EduTool `equilibrium-landscapes` (gui/src/ui/methods/EquilibriumLandscapesTool.tsx) and the first case of the `reactionGibbs` property operation (src/propertyOps/ReactionGibbs.cpp), which exists because of it.

WHY A NEW OP WAS NEEDED, measured before it was written. The engine's one route to a reaction's Gibbs energy, `Reaction::equilibrium` (src/thermo/reaction/Reaction.cpp), prices every species on the ideal-gas rung with `g\_pure\_ig(T)`, and that rung REFUSES a record tabulated on the solid standard state. A gibbsMap with glucose beside the gases was run and refused, verbatim: "Component 'glucose': h\_pure\_ig(T) evaluates the IDEAL-GAS rung, but this record declares standardThermochemistry. referenceState pureSolid -- dHf\_298 and s\_298 are tabulated on the solid standard state, so reading them here would be wrong by a heat of sublimation." The same refusal binds `equilibriumReactor` and `gibbsReactor`. The biological routes are exactly a solid sugar and liquid water beside gases, so no reactor and no existing op could price them; `reactionGibbs` sums a DECLARED reaction over the phase-aware `h\_formation` / `s\_formation` surface the Component already carried.

THE THREE ROUTES (constant/reactions), balanced by ELECTRON COUNT -- the biological ones carry NO H2 (Vitor's correction, 2026-09-27: nitrogenase reduces N2 with protons and electrons that come from the organism's substrate): heterotrophic C6H12O6(s) + 4 N2 + 6 H2O(l) $\to$ 6 CO2 + 8 NH3 phototrophic N2 + 3 H2O(l) $\to$ 2 NH3 + 3/2 O2 haberBosch N2 + 3 H2 $\to$ 2 NH3 Every species is priced on a DECLARED standard state (`phases \{\}`): the sugar on its own pureSolid rung, water as a pure liquid reached by the vaporisation crossing from its ideal-gas datum (announced), the gases on idealGas. The lesson's expectation, stated before the run: the sugar route is EXERGONIC (respiration pays for fixation), the water route is ENDERGONIC (light pays), the industrial route is exergonic at 298 K and loses that with temperature -- the landscape witness beside this case draws the loss.

TWO TEMPERATURES, AND WHY ONLY THE FIRST IS ASKED OF THE SUGAR ROUTE. 298.15 K is the datum temperature, where a declared rung answers with no heat capacity integrated. glucose.dat carries NO heat-capacity block on any rung ("FLAGGED, NOT FABRICATED", its own words), so asking the heterotrophic route at 310.15 K REFUSES by name -- measured, verbatim: "Component 'glucose': h\_formation liquid leg needs liquidHeatCapacity" (the solid leg falls back to the liquid Cp and finds none). The refusal is the finding and is quoted on the page; this case therefore prices all three routes at 298.15 K in one op and the two routes that can move in temperature at 310.15 K in a second op, so that the run passes and the gap is stated rather than papered over; a third op prices the industrial route alone at 700 K, the converter temperature of gibbs10, where its dG has changed sign -- which is the whole reason the converter needs its pressure, and the landscape witness beside this case draws it. Liquid water is not asked at 700 K: a pure-liquid standard state above the critical temperature is a number with no state behind it.

NOTHING HERE IS A RATE. An exergonic route is one the equilibrium permits; what makes each one go -- the enzyme's ATP, the light, the iron catalyst's temperature -- is outside every record in this tree, and the ATP stoichiometry of nitrogenase is deliberately NOT cited because no reference for it exists in the repository.

\end{tutnarrative}

\clearpage
\tutentry{landscape07\_sugar\_nitrogen\_all\_pathways}{tutorials/props/gibbs/landscape07\_sugar\_nitrogen\_all\_pathways}
\begin{tutscheme}
\begin{tikzpicture}[>={Stealth[length=2mm]}, every node/.style={font=\scriptsize}]
  \node[draw, rounded corners=2pt, align=center, minimum width=3.1cm, minimum height=0.8cm, fill=black!4] at (0.0,-0.0) {\textbf{sugarNitrogen}\\[-1pt]{\tiny gibbsMap}};
  \node[draw, rounded corners=2pt, align=center, minimum width=3.1cm, minimum height=0.8cm, fill=black!4] at (3.6,-0.0) {\textbf{boudouard}\\[-1pt]{\tiny gibbsMap}};
\end{tikzpicture}
\end{tutscheme}
\tutwhat{Gibbs maps with a pure solid: the sucrose nitrogen-fixation reactants with every pathway open, and Boudouard carbon}
\begin{tutnarrative}
TWO GIBBS MAPS THAT CARRY A PURE SOLID, the first cases to (2026-09-29). Until that day `gibbsMap` called the fluid-only solve, so a species whose record declares `referenceState pureSolid;` was priced on the ideal-gas rung and its record refused the whole map (found by Vitor, mapping sucrose beside the gases). The op now calls `GibbsMethod::solve`, the one home the gibbsReactor already used: a declared pure solid is its own phase at activity one, deposited only where the solid-free equilibrium would supersaturate it, and its amount is the CSV column `s\_<name>`.

sugarNitrogen -- the REACTANTS of the sucrose nitrogen-fixation route C12H22O11(s) + 8 N2 + 13 H2O $\to$ 12 CO2 + 16 NH3 (balanced by electron count; its dG at 298.15 K, 1 bar is about -355 kJ per sucrose on `reactionGibbs`, so the ROUTE is permitted), handed to a Gibbs map with every species open, methane included. Where the atoms go at minimum G is NOT that route: measured, the carbon ends as methane and CO2 (the anaerobic digester's answer), the nitrogen stays N2 and ammonia is near 1e-6 at ambient. A declared route that is permitted is not the route equilibrium picks; the organism picks it with an enzyme.

boudouard -- 2 CO <$\to$ C(graphite) + CO2, CO fed alone: carbon DEPOSITS at low temperature and not at high, pinned as the fraction of the fed carbon that leaves as graphite (elementYield on the solid).

NOT HERE, said: a solid that must REMAIN for the atoms to close. With the six species sucrose, CO2, N2, NH3, H2, water and no methane, carbon has one fluid carrier (CO2) and too little oxygen for it, so the solid-free equilibrium the solid test starts from does not exist; every cell reports unconverged, with the solver's own cause printed. That is a limit of GibbsMethod::solve, the same in the gibbsReactor, recorded in DEV.md section 5, not fixed here.

\end{tutnarrative}
\begin{tutgolden}
\goldrow{diag}{sugarNitrogen}{metric\_T25C\_P1atm}{1.27109422855e-06}
\goldrow{diag}{sugarNitrogen}{metric\_T37C\_P1atm}{2.67461125111e-06}
\goldrow{diag}{sugarNitrogen}{n\_cells}{39}
\goldrow{diag}{sugarNitrogen}{n\_unconverged}{0}
\goldrow{diag}{boudouard}{metric\_T427C\_P1atm}{0.496091670405}
\goldrow{diag}{boudouard}{metric\_T927C\_P1atm}{0.0181328889912}
\goldrow{diag}{boudouard}{n\_cells}{30}
\goldrow{diag}{boudouard}{n\_unconverged}{0}
\end{tutgolden}

\clearpage
\tutentry{nh3\_equilibrium\_map}{tutorials/props/gibbs/nh3\_equilibrium\_map}
\begin{tutscheme}
\begin{tikzpicture}[>={Stealth[length=2mm]}, every node/.style={font=\scriptsize}]
  \node[draw, rounded corners=2pt, align=center, minimum width=3.1cm, minimum height=0.8cm, fill=black!4] at (0.0,-0.0) {\textbf{haber\_map}\\[-1pt]{\tiny gibbsMap}};
\end{tikzpicture}
\end{tutscheme}
\tutwhat{Ammonia equilibrium map (T x P): why Haber runs at 400-500 C and 150-300 bar}
\begin{tutnarrative}
THE HABER MAP -- the forum-ratified Gibbs equilibrium map (docs/design/gibbs-map-forum-2026-07-02.md) on its canonical pilot: N2 + 3 H2 at stoichiometric feed, equilibrium NH3 mole fraction contoured over 300-700 C x 1-300 bar. The map shows the whole Haber story at a glance: equilibrium loves COLD and PRESSURE (exothermic, Delta\_n < 0); the industrial window (400-500 C, 150-300 bar) sits where the catalyst can actually work -- the tension IS the lesson.

VALIDATION ANCHORS (golden-locked; verified 2026-07-03): the classic measurements of Larson \& Dodge, J. Am. Chem. Soc. 45 (1923) 2918 -- 400 C / 10 atm: engine 4.09 \% vs L\&D 3.85 \% (+6 \%) 450 C / 10 atm: engine 2.22 \% vs L\&D 2.11 \% (+5 \%) The \textasciitilde{}5 \% offset is consistent with propagating the JANAF formation-data uncertainty of NH3 into Kp (exp(-dG/RT) amplifies \textasciitilde{}0.2 kJ/mol into \textasciitilde{}5 \%) -- honest agreement for an equilibrium constant spanning decades. IDEAL-GAS VALIDITY: the kernel is ideal-gas, so pins stay at <= 50 atm, 350-500 C (the forum's chair: at 300 bar a temperatureApproach would silently absorb the missing fugacity corrections -- announced, and the high-P corner of this map is drawn as EQUILIBRIUM SHAPE, not quantitative prediction).

temperatureApproach: 0 here (true equilibrium). Re-run with e.g. `temperatureApproach 25;` to see the contours slide -- the empirical closeness-to-equilibrium parameter industry quotes in C (announced loud, the anchor KPI names then carry \_dT25 so the golden breaks visibly).

\end{tutnarrative}
\begin{tutgolden}
\goldrow{diag}{haber\_map}{metric\_T400C\_P10atm}{0.0409234}
\goldrow{diag}{haber\_map}{metric\_T400C\_P50atm}{0.157785}
\goldrow{diag}{haber\_map}{metric\_T450C\_P10atm}{0.0221983}
\goldrow{diag}{haber\_map}{metric\_T500C\_P100atm}{0.104969}
\goldrow{diag}{haber\_map}{n\_cells}{625}
\goldrow{diag}{haber\_map}{n\_unconverged}{0}
\end{tutgolden}

\clearpage
\tutentry{reactionEnthalpy01\_haber\_kirchhoff}{tutorials/props/gibbs/reactionEnthalpy01\_haber\_kirchhoff}
\begin{tutscheme}
\begin{tikzpicture}[>={Stealth[length=2mm]}, every node/.style={font=\scriptsize}]
  \node[draw, rounded corners=2pt, align=center, minimum width=3.1cm, minimum height=0.8cm, fill=black!4] at (0.0,-0.0) {\textbf{haberT}\\[-1pt]{\tiny reactionGibbs}};
\end{tikzpicture}
\end{tutscheme}
\tutwhat{The reaction enthalpy is a function of temperature: N2 + 3 H2 $\to$ 2 NH3 priced from 25 C to 527 C on the formation surface -- the witness of the reaction-enthalpy EduTool}
\begin{tutnarrative}
THE REACTION ENTHALPY IS A FUNCTION OF TEMPERATURE. The ammonia synthesis N2 + 3 H2 $\to$ 2 NH3, every species an ideal gas at the standard pressure, priced by `reactionGibbs` at seven temperatures from 298.15 K (where formation enthalpies are tabulated) to 800 K (above an ammonia converter):

dH$^\circ$(T) = sum\_i nu\_i h\_i(T), h\_i(T) = dHf\_i(298.15 K) + int Cp\_i dT

The "$^\circ$" fixes the standard state (1 bar, ideal gas); it does not fix the temperature. At 298.15 K the result is -91.88 kJ per mole of reaction, i.e. -45.94 kJ per mole of NH3; at a converter's 700-800 K it is about -53 kJ per mole of NH3 -- the 25 C value under-states the heat a converter releases by about 14 \%. The witness of the EduTool `reaction-enthalpy` (DEV.md 4c, C31). The constant/ tree is landscape06's, copied whole: the same formation surface, so the 298.15 K and 700 K rows reproduce that case's own `routes298` and `haber700` Haber-Bosch rows.

\end{tutnarrative}
\tutreadme{The witness of the EduTool Reaction enthalpy and temperature (reaction-enthalpy). One reactionGibbs operation prices N$_2$ + 3 H$_2$ $\to$ 2 NH$_3$ (ideal gases at 1 bar) at seven temperatures, from 298.15 K to 800 K, on the formation surface of the catalogue records.}
\tutreadme{The degree sign fixes the standard state (1 bar, each species in its reference phase), not the temperature: tables are printed at 298.15 K by convention. A converter running at 700--800 K releases about 16 \% more heat per mole of ammonia than the 25 $^\circ$C value says.}
\tutreadme{constant/ is copied whole from [landscape06\_nitrogen\_fixation\_routes](../landscape06\_nitrogen\_fixation\_routes), so the formation surface is the same one; this case reads only N$_2$, H$_2$ and NH$_3$.}
\begin{tutgolden}
\goldrow{diag}{haberT}{haberBosch\_T298K\_dG\_kJ\_mol}{-32.8135035}
\goldrow{diag}{haberT}{haberBosch\_T298K\_dH\_kJ\_mol}{-91.88}
\goldrow{diag}{haberT}{haberBosch\_T298K\_dS\_J\_molK}{-198.11}
\goldrow{diag}{haberT}{haberBosch\_T298K\_lnK}{13.2368181896}
\goldrow{diag}{haberT}{haberBosch\_T400K\_dG\_kJ\_mol}{-11.9635351712}
\goldrow{diag}{haberT}{haberBosch\_T400K\_dH\_kJ\_mol}{-96.1395343725}
\goldrow{diag}{haberT}{haberBosch\_T400K\_dS\_J\_molK}{-210.439998003}
\goldrow{diag}{haberT}{haberBosch\_T400K\_lnK}{3.59720637425}
\goldrow{diag}{haberT}{haberBosch\_T500K\_dG\_kJ\_mol}{9.51125665348}
\goldrow{diag}{haberT}{haberBosch\_T500K\_dH\_kJ\_mol}{-99.7619775808}
\goldrow{diag}{haberT}{haberBosch\_T500K\_dS\_J\_molK}{-218.546468469}
\goldrow{diag}{haberT}{haberBosch\_T500K\_lnK}{-2.2878824743}
\goldrow{diag}{haberT}{haberBosch\_T600K\_dG\_kJ\_mol}{31.6629415717}
\goldrow{diag}{haberT}{haberBosch\_T600K\_dH\_kJ\_mol}{-102.826551039}
\goldmore{14}
\end{tutgolden}

\subsection{heat}
\grouptheory{heat}
\clearpage
\tutentry{bellDelaware01\_hewitt\_example62}{tutorials/props/heat/bellDelaware01\_hewitt\_example62}
\begin{tutscheme}
\begin{tikzpicture}[>={Stealth[length=2mm]}, every node/.style={font=\scriptsize}]
  \node[draw, rounded corners=2pt, align=center, minimum width=3.1cm, minimum height=0.8cm, fill=black!4] at (0.0,-0.0) {\textbf{book}\\[-1pt]{\tiny bellDelawareBench}};
  \node[draw, rounded corners=2pt, align=center, minimum width=3.1cm, minimum height=0.8cm, fill=black!4] at (3.6,-0.0) {\textbf{exact}\\[-1pt]{\tiny bellDelawareBench}};
\end{tikzpicture}
\end{tutscheme}
\tutwhat{Bell-Delaware shell side -- Hewitt, Shires \& Bott Example 6.2 reproduced step by step at the book's own declared fluid and geometry (J\_L, J\_B, R\_L, R\_B read from Bell's charts and declared, as the textbook does; `auto` computes them instead).}
\begin{tutnarrative}
the Bell-Delaware shell-side rating on the textbook's own worked example.

SOURCE. G. F. Hewitt, G. L. Shires, T. R. Bott, $\ast$Process Heat Transfer$\ast$ (CRC Press, 1994), Example 6.1 (the geometry and the light-hydrocarbon stream) and Example 6.2 (the Bell-Delaware solution, pp. 280-285), with the method of Section 6.3.2 (Eqs 6.8-6.19, 6.20, 6.22, 6.26-6.31) and the ideal bank of Eq 2.218 / Table 2.5.

WHAT IS DECLARED, AS THE BOOK DOES. Bell's J\_L, J\_B (Figs 6.8, 6.9) and R\_L, R\_B (Figs 6.10, 6.11) are CHART READINGS in this method; the book reads J\_L = 0.58, J\_B = 0.92, R\_L = 0.33, R\_B = 0.76 and so does this case. Each factor is a declared number (the chart reading) or the word `auto` (Taborek's closed form, HEDH 3.3; the gate runs this geometry with `auto` and holds the four values to the book's readings within the reading error); an absent one refuses by name. The fluid properties are the book's too (the bench is the textbook's inputs, not a package's), which is why this is a props bench.

WHAT THE BOOK PRINTS, AND WHAT THIS RUN GIVES (the gate holds every row): S\_m 0.01667 m2 (book) 0.016622 from Eq 6.9 as written: the book's S\_m is 0.3 \% high. Op `book` starts from the PRINTED S\_m and reproduces every digit; op `exact` starts from the geometry Re 2.097e4 Pr 7.481 ideal h 1418 W/m2K F\_c 0.6470 J\_c 1.016 (S\_sb+S\_tb)/S\_m 0.4187 F\_bp 0.2666 h 769 W/m2K K\_f 0.3657 N\_c 8.5 dp\_c 233 Pa dp\_w 165 Pa shell dP 4501 Pa

\end{tutnarrative}
\begin{tutgolden}
\goldrow{diag}{book}{Fbp}{0.266646670666}
\goldrow{diag}{book}{Fc}{0.646768246556}
\goldrow{diag}{book}{JB}{0.92}
\goldrow{diag}{book}{JL}{0.58}
\goldrow{diag}{book}{J\_combined}{0.541963186181}
\goldrow{diag}{book}{Jc}{1.01567313752}
\goldrow{diag}{book}{Kf}{0.365744906832}
\goldrow{diag}{book}{Nc}{8.50078740157}
\goldrow{diag}{book}{Ncw}{3.40031496063}
\goldrow{diag}{book}{Pr}{7.48089123867}
\goldrow{diag}{book}{Re}{20969.998519}
\goldrow{diag}{book}{Sm}{0.01667}
\goldrow{diag}{book}{Ssb}{0.00282638619068}
\goldrow{diag}{book}{Stb}{0.0041524379945}
\goldmore{34}
\end{tutgolden}

\clearpage
\tutentry{bellDelaware02\_serth\_example61}{tutorials/props/heat/bellDelaware02\_serth\_example61}
\begin{tutscheme}
\begin{tikzpicture}[>={Stealth[length=2mm]}, every node/.style={font=\scriptsize}]
  \node[draw, rounded corners=2pt, align=center, minimum width=3.1cm, minimum height=0.8cm, fill=black!4] at (0.0,-0.0) {\textbf{chart}\\[-1pt]{\tiny bellDelawareBench}};
  \node[draw, rounded corners=2pt, align=center, minimum width=3.1cm, minimum height=0.8cm, fill=black!4] at (3.6,-0.0) {\textbf{equations}\\[-1pt]{\tiny bellDelawareBench}};
\end{tikzpicture}
\end{tutscheme}
\tutwhat{Bell-Delaware shell side, form Serth -- Serth Example 6.1 with every correction factor computed (J\_L, J\_B, R\_L, R\_B auto) and the ideal bank read from Fig 6.1 (op chart) or from Eqs 6.1-6.4 (op equations).}
\begin{tutnarrative}
the Bell-Delaware shell side, FORM SERTH, on the textbook's own worked example.

SOURCE. R. W. Serth, $\ast$Process Heat Transfer: Principles, Applications and Rules of Thumb$\ast$ (Academic Press / Elsevier, 2007), Example 6.1 (the final kerosene/crude-oil exchanger of Example 5.1, 124 tubes of 1.0 in OD on a 1.25 in square pitch in a 19.25 in shell), with Sections 6.4-6.7 (Eqs 6.6-6.57) and Table 6.1. Inputs below are the book's, converted to SI (1 in = 0.0254 m exactly; 1 lb = 0.45359237 kg; 1 Btu/lb F = 4186.8 J/kg K).

UNLIKE HEWITT'S EXAMPLE, every correction factor here is COMPUTED (`auto`): J\_L, J\_B, R\_L, R\_B from Eqs 6.48-6.53, J\_S = R\_S = 1 (equal spacing, Eqs 6.54-6.55) and J\_R = 1 (Re > 100, Eq 6.57). Op `chart` declares the ideal bank as the book's example reads it from Fig 6.1 (j = 0.0060, f = 0.09); op `equations` computes j and f from Eqs 6.1-6.4 (the book prints 0.0063 and 0.089 for them).

CLEARANCES. Serth's Eqs 6.22 and 6.36 use d\_tb and d\_sb as the gap itself and his Example 6.1 inserts the TEMA values (0.4 mm; 0.8 + 0.002 D\_s mm plus a 0.75 mm allowance) in them, so `clearanceBasis radial;` is the reading that reproduces his printed S\_tb, S\_sb, J\_L and R\_L; `diametral` would halve both areas and the gate shows it does NOT reproduce them.

PRESSURE DROP. The printed summary Eq 6.16 omits R\_L on the central spaces while Eq 6.9 has it; the example's own arithmetic (dP\_f = 159.4 lbf/ft2) needs it, so Eq 6.9 is used and Eq 6.16 is a candidate erratum.

\end{tutnarrative}
\begin{tutgolden}
\goldrow{diag}{chart}{Fbp}{0.283778060144}
\goldrow{diag}{chart}{Fc}{0.796270244553}
\goldrow{diag}{chart}{JB}{0.863052673184}
\goldrow{diag}{chart}{JL}{0.479455411602}
\goldrow{diag}{chart}{JR}{1}
\goldrow{diag}{chart}{JS}{1}
\goldrow{diag}{chart}{J\_combined}{0.464822263533}
\goldrow{diag}{chart}{Jc}{1.12331457608}
\goldrow{diag}{chart}{Nc}{9.24}
\goldrow{diag}{chart}{Ncw}{2.464}
\goldrow{diag}{chart}{Pr}{7.23999985478}
\goldrow{diag}{chart}{RB}{0.64665168579}
\goldrow{diag}{chart}{RL}{0.265542709084}
\goldrow{diag}{chart}{RS}{1}
\goldmore{48}
\end{tutgolden}

\clearpage
\tutentry{htbench01\_correlations}{tutorials/props/heat/htbench01\_correlations}
\begin{tutscheme}
\begin{tikzpicture}[>={Stealth[length=2mm]}, every node/.style={font=\scriptsize}]
  \node[draw, rounded corners=2pt, align=center, minimum width=3.1cm, minimum height=0.8cm, fill=black!4] at (0.0,-0.0) {\textbf{bench}\\[-1pt]{\tiny heatTransferBench}};
\end{tikzpicture}
\end{tutscheme}
\tutwhat{Heat-transfer correlation bench -- each convective correlation (Dittus-Boelter, Gnielinski, Kern) verified against its published Nusselt anchor. Dittus-Boelter and Gnielinski are pinned to Incropera \& DeWitt 6th ed. Example 8.4; Kern's closed j\_H form is self-consistency-checked (its original pin is a chart). AAD tolerance 5\% honours textbook rounding + the correlations' own scatter.}
\begin{tutnarrative}
the heat-transfer correlation self-check bench.

Mirrors the membrane mass-transfer story (Sherwood) on the heat side (Nusselt): each registered convective correlation RUNS its verify() against a published Nu anchor and the deviation is printed + locked.

DittusBoelter Nu = 0.023 Re\textasciicircum{}0.8 Pr\textasciicircum{}0.4 (Incropera Ex 8.4, Nu=205) Gnielinski Nu = (f/8)(Re-1000)Pr/[..] (Incropera Ex 8.4 alt, Nu=224) Kern (shell) Nu = 0.36 Re\_s\textasciicircum{}0.55 Pr\textasciicircum{}1/3 (Kern PHT; closed-form check)

The per-correlation deviation diags are the golden KPIs.

\end{tutnarrative}
\begin{tutgolden}
\goldrow{diag}{bench}{Nu\_choupo\_DittusBoelter}{200.475}
\goldrow{diag}{bench}{Nu\_choupo\_Gnielinski}{223.734}
\goldrow{diag}{bench}{Nu\_choupo\_Kern}{97.5647}
\goldrow{diag}{bench}{Nu\_published\_DittusBoelter}{205}
\goldrow{diag}{bench}{Nu\_published\_Gnielinski}{224}
\goldrow{diag}{bench}{Nu\_published\_Kern}{97.5647}
\goldrow{diag}{bench}{allPass}{1}
\goldrow{diag}{bench}{dev\_DittusBoelter}{0.0220718}
\goldrow{diag}{bench}{dev\_Gnielinski}{0.00118956}
\goldrow{diag}{bench}{dev\_Kern}{0}
\goldrow{diag}{bench}{dev\_NusseltFilm}{0.0245797}
\goldrow{diag}{bench}{dev\_Rohsenow}{0.000270621}
\goldrow{diag}{bench}{dev\_ZuberCHF}{0.000684704}
\goldrow{diag}{bench}{h\_choupo\_NusseltFilm}{6680.26}
\goldmore{5}
\end{tutgolden}

\clearpage
\tutentry{psychro01\_n2\_water}{tutorials/props/heat/psychro01\_n2\_water}
\begin{tutscheme}
\begin{tikzpicture}[>={Stealth[length=2mm]}, every node/.style={font=\scriptsize}]
  \node[draw, rounded corners=2pt, align=center, minimum width=3.1cm, minimum height=0.8cm, fill=black!4] at (0.0,-0.0) {\textbf{airWater}\\[-1pt]{\tiny psychrometricChart}};
\end{tikzpicture}
\end{tutscheme}
\tutwhat{Psychrometric chart for N2/water at 1 atm, 0-60 C: saturation + four RH lines + three adiabatic-saturation lines, from the vapour-pressure model a flash would use}
\begin{tutnarrative}
THE PSYCHROMETRIC CHART, computed rather than looked up.

Nitrogen + water at 1 atm, 0 $\to$ 60 C: the saturation curve, four relative-humidity lines and three adiabatic-saturation (wet-bulb) lines, all from the SAME vapour-pressure model a flash would use. Nothing here is a chart digitised from a handbook -- the curve is Y\_sat(T) evaluated point by point, so the student can change the model and watch the chart move.

Why this case exists at all: `psychrometricChart` is reachable from the GUI's Property Explorer, and until now NOTHING exercised it. It could break and the whole suite would stay green; the first to find out would be a student opening the Explorer. An op a user can reach and a test cannot is a promise with no guard behind it.

What to read in the output (long-format CSV: T\_C, Y, curve): `saturation` -- Y\_sat = 0.622 $\ast$ Psat/(P - Psat), the chart's ceiling; `RH=...` -- the same expression at a fraction of Psat, so the lines fan out and steepen with T exactly as the printed chart does; `wetBulb=...` -- the adiabatic-saturation lines, which are NOT parallel to the RH lines: they carry the latent heat of the water evaporated, so their slope is set by Hvap/cp of the moist air, not by geometry.

\end{tutnarrative}
\begin{tutgolden}
\goldrow{diag}{airWater}{Y\_adiabatic\_last}{0.0150791}
\goldrow{diag}{airWater}{Y\_sat\_first}{0.00390022}
\goldrow{diag}{airWater}{Y\_sat\_last}{0.158606}
\goldrow{diag}{airWater}{nRhCurves}{4}
\goldrow{diag}{airWater}{nSaturationPoints}{61}
\goldrow{diag}{airWater}{Lewis\_last}{0.816578}
\goldrow{diag}{airWater}{nAdiabaticCurves}{3}
\goldrow{diag}{airWater}{nTrueWetBulbCurves}{3}
\end{tutgolden}

\clearpage
\tutentry{psychro02\_air\_water}{tutorials/props/heat/psychro02\_air\_water}
\begin{tutscheme}
\begin{tikzpicture}[>={Stealth[length=2mm]}, every node/.style={font=\scriptsize}]
  \node[draw, rounded corners=2pt, align=center, minimum width=3.1cm, minimum height=0.8cm, fill=black!4] at (0.0,-0.0) {\textbf{airWater}\\[-1pt]{\tiny psychrometricChart}};
\end{tikzpicture}
\end{tutscheme}
\tutwhat{Psychrometric chart with dry AIR as the carrier (the textbook object): saturation, RH, adiabatic-saturation and true wet-bulb families at 101.325 kPa, Lewis \textasciitilde{} 0.851 -- the case-local air component deliberately shadows the catalogue mixture alias}
\begin{tutnarrative}
THE TEXTBOOK CHART: dry AIR as the carrier, not a stand-in. Twin of psychro01\_n2\_water with the REAL carrier -- the lumped pseudo-component air (MW 28.96, ChemSep v8.3), which this case declares CASE-LOCALLY in constant/components/air.dat. That declaration is load-bearing, not a convenience: the catalogue also holds mixtures/air.dat (an expansion table, air $\to$ N2 + O2 + Ar), and the settled precedence -- the case always overrides -- is what makes the token load as the COMPONENT here; the run announces the shadowed expansion. Until 2026-08-31 that tie went to the mixture unconditionally and this exact case refused with "carrier 'air' is not a component".

What the twin buys over psychro01: Y\_sat(60 C) = 0.1534 against the \textasciitilde{}0.152 a printed air chart reads (N2 gives 0.1586 -- the 3.4 \% MW error the registry note predicted), and Lewis = 0.851, the handbook air/water value the Le\textasciicircum{}(2/3) wet-bulb family exists to teach.

\end{tutnarrative}
\begin{tutgolden}
\goldrow{diag}{airWater}{Lewis\_last}{0.850543}
\goldrow{diag}{airWater}{Y\_adiabatic\_last}{0.0146521}
\goldrow{diag}{airWater}{Y\_sat\_first}{0.00377268}
\goldrow{diag}{airWater}{Y\_sat\_last}{0.15342}
\goldrow{diag}{airWater}{nAdiabaticCurves}{3}
\goldrow{diag}{airWater}{nRhCurves}{4}
\goldrow{diag}{airWater}{nSaturationPoints}{61}
\goldrow{diag}{airWater}{nTrueWetBulbCurves}{3}
\end{tutgolden}

\subsection{hydraulics}
\grouptheory{hydraulics}
\clearpage
\tutentry{letdown02\_valve\_fluid\_properties}{tutorials/props/hydraulics/letdown02\_valve\_fluid\_properties}
\begin{tutscheme}
\begin{tikzpicture}[>={Stealth[length=2mm]}, every node/.style={font=\scriptsize}]
  \node[draw, rounded corners=2pt, align=center, minimum width=3.1cm, minimum height=0.8cm, fill=black!4] at (0.0,-0.0) {\textbf{purge}\\[-1pt]{\tiny propertyPoint}};
  \node[draw, rounded corners=2pt, align=center, minimum width=3.1cm, minimum height=0.8cm, fill=black!4] at (3.6,-0.0) {\textbf{liquidLP}\\[-1pt]{\tiny propertyPoint}};
\end{tikzpicture}
\end{tutscheme}
\tutwhat{The two valve inlets of letdown01: the loop purge's Cp/Cv and Z at 140 bar (is the gas letdown choked?) and liquid ammonia's vapour pressure at 253 K (does the liquid letdown cavitate or flash?)}
\begin{tutnarrative}
THE TWO FLUID PROPERTIES A LETDOWN VALVE IS JUDGED BY (DEV.md 4c C64, the "Large pressure letdown" EduTool).

The throttles themselves run in tutorials/steady/hydraulics/ letdown01\_ammonia\_flashing. A valve's hydraulic regime is decided by two more numbers, which a flowsheet unit does not publish, so this props case computes them at the two valve inlets of that case:

the GAS: the loop purge at 140 bar and 253.15 K, let down to a fuel header. Whether its flow is choked is decided by the ratio of specific heats k = Cp/Cv -- propertyPoint's `gamma`, Cp\_ig/(Cp\_ig - R), the ideal-gas ratio the critical-pressure-ratio formula is derived with -- and the SRK compressibility Z. the LIQUID: the 20 bar drum liquid of letdown01, at the temperature its first throttle left it (253.01420618 K, copied from that case's converged/LetdownLiquid; gui/tests/letdown.test.ts holds it to the ValveHP T\_out golden). Whether its second throttle cavitates or flashes is judged by the ammonia's VAPOUR PRESSURE at that temperature, Psat\_NH3 (the record's Antoine curve), against the pressure at the valve's vena contracta -- the PURE ammonia's, so for a liquid holding dissolved gas a didactic approximation.

Compositions are the component flows in kmol/h exactly as letdown01's converged streams print them; propertyPoint normalises them. The liquid's dissolved H2, N2 and Ar do not enter Psat\_NH3 -- that key is the pure ammonia's vapour pressure, the P\_v of the valve-sizing equations. The drum liquid is in equilibrium with its gas at 20 bar, so even a small drop below it can begin to release the gas it still holds, and that is letdown01's to show, not this case's.

\end{tutnarrative}
\begin{tutgolden}
\goldrow{diag}{purge}{Cp\_ig}{28.9429405633}
\goldrow{diag}{purge}{H\_R}{-143.566682255}
\goldrow{diag}{purge}{H\_ig}{-2457.83921394}
\goldrow{diag}{purge}{H\_real}{-2601.40589619}
\goldrow{diag}{purge}{P}{14000000}
\goldrow{diag}{purge}{S\_R}{-1.02704601342}
\goldrow{diag}{purge}{S\_ig}{103.874258602}
\goldrow{diag}{purge}{S\_real}{102.847212588}
\goldrow{diag}{purge}{T}{253.15}
\goldrow{diag}{purge}{Z}{1.067126493}
\goldrow{diag}{purge}{gamma}{1.40305749363}
\goldrow{diag}{purge}{v\_molar}{0.00016043531937}
\goldrow{diag}{liquidLP}{Cp\_ig}{34.2364932879}
\goldrow{diag}{liquidLP}{H\_R}{-2189.62838136}
\goldmore{13}
\end{tutgolden}

\clearpage
\tutentry{moody01\_friction\_correlations}{tutorials/props/hydraulics/moody01\_friction\_correlations}
\begin{tutscheme}
\begin{tikzpicture}[>={Stealth[length=2mm]}, every node/.style={font=\scriptsize}]
  \node[draw, rounded corners=2pt, align=center, minimum width=3.1cm, minimum height=0.8cm, fill=black!4] at (0.0,-0.0) {\textbf{moody}\\[-1pt]{\tiny frictionBench}};
\end{tikzpicture}
\end{tutscheme}
\tutwhat{The Moody chart's four correlations, each verified against its own published anchor and then all asked the same question -- so the disagreement is visible instead of hidden behind a default}
\begin{tutnarrative}
moody01 -- FOUR CORRELATIONS, ONE QUESTION, FOUR ANSWERS.

The Darcy friction factor is the most-used correlated quantity in chemical engineering, and every simulator computes it. What none of them shows is that the answer DEPENDS ON WHICH CORRELATION YOU PICKED.

Choupo has computed Blasius, Colebrook-White, Haaland and Churchill since the pipe unit was written -- correctly, with their citations in comments, as free functions nobody could name or reach. This case is what happens when they become objects: each states the window its own authors gave it, each carries its primary citation, each reproduces a published anchor on every run, and all four can be asked the same question at once.

WHAT THE RUN DOES, IN TWO PARTS.

(1) VERIFY. Each correlation reproduces its own anchor, and the anchors are deliberately of DIFFERENT KINDS, because what can be checked differs:

Blasius against its own closed form. This checks the ARITHMETIC and nothing else, and the bench says so rather than letting a 0.000 \% deviation look like a validation. Colebrook against the fully-rough von Karman limit, where the implicit equation has a closed form. This checks that the FIXED-POINT ITERATION reaches the answer it must. Haaland against COLEBROOK at the same point -- because Haaland's own published claim is "I reproduce that equation within about 2 \%". Anchoring it to anything else would be checking a claim nobody made. Churchill against the EXACT laminar law f = 64/Re, which its single expression contains as its own low-Re limit. This is the strongest of the four: a mis-typed exponent fails here while the turbulent branch still looks plausible.

(2) COMPARE. All four are evaluated at Re = 1e5, eps/D = 1e-3 -- an ordinary turbulent water line in a commercial steel pipe -- and the spread is published as a number.

WHAT TO READ IN THE OUTPUT. Blasius is the one to watch: it carries NO roughness term at all, so at eps/D = 1e-3 it answers as though the wall were smooth, and it is simply wrong. It says so on its own line. That is the whole lesson of a validity window -- not that the correlation is bad, but that it was asked a question it was never fitted for, and only the window can tell you.

WHAT THIS CASE DOES NOT DO. It does not say which correlation is right. That depends on the pipe -- its roughness, whether the flow is developed, whether the duct is even circular -- and a bench that ranked them would hide the choice the engineer has to make, which is exactly what a default does today in every tool that has one.

\end{tutnarrative}
\begin{tutgolden}
\goldrow{diag}{moody}{allPass}{1}
\goldrow{diag}{moody}{dev\_Blasius}{0}
\goldrow{diag}{moody}{dev\_Churchill}{1.73472e-16}
\goldrow{diag}{moody}{dev\_Colebrook}{1.61329e-05}
\goldrow{diag}{moody}{dev\_Haaland}{0.0134393}
\goldrow{diag}{moody}{f\_Blasius}{0.01777}
\goldrow{diag}{moody}{f\_Churchill}{0.0223432}
\goldrow{diag}{moody}{f\_Colebrook}{0.0221745}
\goldrow{diag}{moody}{f\_Haaland}{0.0219662}
\goldrow{diag}{moody}{f\_max}{0.0223432}
\goldrow{diag}{moody}{f\_min}{0.01777}
\goldrow{diag}{moody}{n\_outside\_window}{1}
\goldrow{diag}{moody}{spread\_inWindow\_pct}{1.71637}
\goldrow{diag}{moody}{spread\_pct}{25.7358}
\end{tutgolden}

\clearpage
\tutentry{moody02\_correlation\_sweeps}{tutorials/props/hydraulics/moody02\_correlation\_sweeps}
\begin{tutscheme}
\begin{tikzpicture}[>={Stealth[length=2mm]}, every node/.style={font=\scriptsize}]
  \node[draw, rounded corners=2pt, align=center, minimum width=3.1cm, minimum height=0.8cm, fill=black!4] at (0.0,-0.0) {\textbf{moody}\\[-1pt]{\tiny frictionBench}};
  \node[draw, rounded corners=2pt, align=center, minimum width=3.1cm, minimum height=0.8cm, fill=black!4] at (3.6,-0.0) {\textbf{nusselt}\\[-1pt]{\tiny heatTransferBench}};
\end{tikzpicture}
\end{tutscheme}
\tutwhat{The friction-factor and Nusselt correlations swept across five decades of Reynolds number -- the engine's own curves, which the EduTool on why correlations are simple draws, differentiates and fits}
\begin{tutnarrative}
moody02 -- THE CORRELATIONS AS CURVES, NOT AS POINTS (DEV.md 4c C119).

moody01 asks four friction correlations ONE question and publishes their disagreement. This case asks them the same question at every Reynolds number from 1e2 to 1e8, and asks the Nusselt correlations the same at every Reynolds number from 2e3 to 2e7, and writes the answers as two CSV files. It is the witness of the EduTool "Why engineering correlations are simple", which draws, differentiates and fits these curves in the browser -- so every curve on that page is the engine's own correlation, the one the pipe and the heat exchanger use, never a re-typed copy.

WHAT IS SYNTHETIC, said plainly. The CSVs hold no measurement. They are the correlations evaluated on a grid, and the EduTool treats them as the "truth" a student fits a power law to -- which is the exercise, because the point it teaches is how simple a curve looks over a narrow range, not how well any correlation matches a laboratory.

WHAT THE TWO SWEEPS WRITE.

moodySweep.csv one row per (eps/D, Re): every friction correlation's f and its own window flag (1 inside, 0 outside), the regime the engine names (laminar below 2300, transition to 4000, turbulent above), and Churchill's two asymptotes fL = 64/Re and fT -- the pieces its single expression blends as (fL\textasciicircum{}12 + fT\textasciicircum{}12)\textasciicircum{}(1/12).

nusseltSweep.csv one row per (Pr, Re): every single-phase Nusselt correlation's Nu and window flag. Each correlation is handed a unit fluid whose Re and Pr ARE the grid point; heating side. Kern is SHELL-side and is evaluated on the same numbers -- it is not a tube correlation.

The grid is logarithmic, 20 points per decade, so every decade holds 21 points, its two ends included -- the EduTool's one-decade power-law fit uses exactly those. The Nusselt grid starts at 2e3 because Gnielinski's expression carries (Re - 1000) and is zero at Re = 1000.

\end{tutnarrative}
\begin{tutgolden}
\goldrow{diag}{moody}{allPass}{1}
\goldrow{diag}{moody}{dev\_Blasius}{0}
\goldrow{diag}{moody}{dev\_Churchill}{1.73472347598e-16}
\goldrow{diag}{moody}{dev\_Colebrook}{1.61328517053e-05}
\goldrow{diag}{moody}{dev\_Haaland}{0.013439281758}
\goldrow{diag}{moody}{sweep\_rows}{726}
\goldrow{diag}{nusselt}{Nu\_choupo\_DittusBoelter}{200.475286267}
\goldrow{diag}{nusselt}{Nu\_choupo\_Gnielinski}{223.73353907}
\goldrow{diag}{nusselt}{Nu\_choupo\_Kern}{97.5646525378}
\goldrow{diag}{nusselt}{Nu\_published\_DittusBoelter}{205}
\goldrow{diag}{nusselt}{Nu\_published\_Gnielinski}{224}
\goldrow{diag}{nusselt}{Nu\_published\_Kern}{97.5646525378}
\goldrow{diag}{nusselt}{allPass}{1}
\goldrow{diag}{nusselt}{dev\_DittusBoelter}{0.0220717743085}
\goldmore{12}
\end{tutgolden}

\subsection{kinetics}
\grouptheory{kinetics}
\clearpage
\tutentry{ammoniaRate01\_dyson\_simon}{tutorials/props/kinetics/ammoniaRate01\_dyson\_simon}
\begin{tutscheme}
\begin{tikzpicture}[>={Stealth[length=2mm]}, every node/.style={font=\scriptsize}]
  \node[draw, rounded corners=2pt, align=center, minimum width=3.1cm, minimum height=0.8cm, fill=black!4] at (0.0,-0.0) {\textbf{dysonSimon}\\[-1pt]{\tiny rateLawBench}};
\end{tikzpicture}
\end{tutscheme}
\tutwhat{The ammonia synthesis rate of Dyson \& Simon (1968), WRITTEN in the case (constant/reactions, `type expression`, every constant with its unit and equation number) and put on the rate-law bench: Eq 19's pseudo-homogeneous rate, Eq 39's effectiveness factor as a declared table over Table I's three pressures, Eq 40, a declared 150-300 atm window and an announced pressure interpolation; the paper's anchors are held by a gate on every case's copy of the law. It locates the LOCUS OF MAXIMUM RATE at three conversions and shows it descending -- the curve ammonia03\_quench\_converter zigzags about -- and evaluates the rate at that converter's own bed states.}
\begin{tutnarrative}
propsDict -- the ammonia synthesis rate of Dyson \& Simon (1968), on the bench. The law itself is constant/reactions' `ammoniaSynthesis`, every number of it written there (`type expression;`); this file only says where to evaluate it and what to print. Two questions, and they are different questions.

THE LOCUS. Sweeping temperature at a FIXED conversion locates the temperature of greatest net rate. Three conversions locate three points of the locus, and the point is that it DESCENDS. The conversion is of N2, walked from the paper's own reference mixture (3:1 H2/N2, 12.7 \% inerts) along the reaction's stoichiometry. The sweep runs 600-900 K because that is roughly the band an iron catalyst lives in; at the lowest conversion the maximum is pushed ONTO the window's upper edge, and the bench says so rather than reporting an edge as a maximum.

THE DECLARED STATES. Three gases TRANSCRIBED from the converged answer of tutorials/steady/reactors/ammonia03\_quench\_converter, at the digits that case's own golden pins. A transcription is a SECOND HOME for somebody else's number and it will go stale if that case moves, so `check\_ammonia\_rate` RUNS ammonia03 and holds these three blocks to its converged result.

The catalyst particle is the law's own parameter (`dp 8 mm;` in constant/reactions), inside the 6-10 mm range Eq 39 was fitted for; the law REFUSES one outside it.

\end{tutnarrative}
\tutreadme{tutorials/steady/reactors/ammonia03\_quench\_converter builds a three-bed quench converter out of equilibrium beds and then says, in its own README, what it cannot do: it cannot size a bed, it cannot tell a good injection point from a bad one, and it rewards over-quenching into a region where a real catalyst is too slow to reach the equilibrium the model hands it. All three gaps are the same gap --- there is no rate law.}
\tutreadme{This case is the rate law --- and since 2026-10-12 (DEV.md 4c C67) the law is WRITTEN here, in constant/reactions, every constant with its unit and its equation number, as kinetics \{ type expression; \}. The engine compiles it once, checks its dimensions, prints it back as PARSED (the top of the run log), and the rateLawBench op evaluates it. Nothing about ammonia is compiled into the engine; the paper's anchors live in a gate.}
\tutreadme{The source. D. C. Dyson and J. M. Simon, $\ast$A kinetic expression with diffusion correction for ammonia synthesis on industrial catalyst$\ast$, Ind. Eng. Chem. Fundam. 7 (1968) 605--610. Equations 2, 6--8, 18, 19, 39 with Table I, and 40.}
\tutreadme{Nothing in this case is measured by Choupo: the data these equations were fitted to are Nielsen, Kj?r and Hansen's (J. Catalysis 3 (1964) 68).}
\tutreadme{The same two letters, meaning opposite things. This case keeps the paper's convention because every number in it is read off the paper's own equations, and the engine says so at every surface that prints either letter. Carry $\eta$ across from the reactor side and you will read a conversion as an effectiveness factor.}
\tutreadme{A correlation in this tree carries three things a reader can check, and this one carries them in the CASE:}
\begin{tutgolden}
\goldrow{diag}{dysonSimon}{bed1Inlet\_Eq40}{166.458716538}
\goldrow{diag}{dysonSimon}{bed1Inlet\_V3}{386.705388618}
\goldrow{diag}{dysonSimon}{bed1Inlet\_eta}{0.0365798726419}
\goldrow{diag}{dysonSimon}{bed1Inlet\_xi}{0.430453573799}
\goldrow{diag}{dysonSimon}{bed1Outlet\_Eq40}{2.87605463668}
\goldrow{diag}{dysonSimon}{bed1Outlet\_V3}{5.33952792223}
\goldrow{diag}{dysonSimon}{bed1Outlet\_eta}{0.224255825635}
\goldrow{diag}{dysonSimon}{bed1Outlet\_xi}{0.53863462811}
\goldrow{diag}{dysonSimon}{bed3Outlet\_Eq40}{0.374882283122}
\goldrow{diag}{dysonSimon}{bed3Outlet\_V3}{0.466226453532}
\goldrow{diag}{dysonSimon}{bed3Outlet\_eta}{0.336428411115}
\goldrow{diag}{dysonSimon}{bed3Outlet\_xi}{0.804077675735}
\goldrow{diag}{dysonSimon}{locus1\_Eq40}{-889.947776496}
\goldrow{diag}{dysonSimon}{locus1\_T\_maxRate\_K}{900}
\goldmore{14}
\end{tutgolden}

\clearpage
\tutentry{rate01\_arrhenius\_recovery}{tutorials/props/kinetics/rate01\_arrhenius\_recovery}
\begin{tutscheme}
\begin{tikzpicture}[>={Stealth[length=2mm]}, every node/.style={font=\scriptsize}]
  \node[draw, rounded corners=2pt, align=center, minimum width=3.1cm, minimum height=0.8cm, fill=black!4] at (0.0,-0.0) {\textbf{rate\_arrhenius}\\[-1pt]{\tiny kinetics1D}};
\end{tikzpicture}
\end{tutscheme}
\tutwhat{Second-order decay at three temperatures, generated from k0 = 4e6, Ea = 60 kJ/mol: the Arrhenius line must give both back}
\begin{tutnarrative}
THE ACTIVATION ENERGY COMES OUT OF A SLOPE.

`kinetics1D` is a PROPERTY evaluation, not a reactor: it uses the integrated rate law in closed form, so a student can overlay competing orders on the same lab data and SEE which one is straight. Given isotherms at more than one temperature it goes further -- fits k on each one, then the line

ln k = ln k0 - (Ea/R) (1/T)

and Ea falls out of the SLOPE. That is the whole of Arrhenius, and it is worth watching happen on data whose answer is known.

So the data here is generated from k0 = 4.0e6 L/(mol s), Ea = 60 kJ/mol, second order, at 40/60/80 C. The fit must return those two numbers; the expected file pins them, and R2 = 1 because the points carry no scatter.

What to read: the per-isotherm k values (3.93e-4, 1.57e-3, 5.34e-3), which span a factor of 13.6 over 40 K -- the reason a rate constant is never quoted without its temperature.

Why the case exists: `kinetics1D` had NO case and no gate. Its sibling vaporPressureFit was in the same position, and writing the first case for THAT one segfaulted choupoProps twice.

\end{tutnarrative}
\begin{tutgolden}
\goldrow{diag}{rate\_arrhenius}{Ea}{60000}
\goldrow{diag}{rate\_arrhenius}{Ea\_kJ\_per\_mol}{60}
\goldrow{diag}{rate\_arrhenius}{R2}{1}
\goldrow{diag}{rate\_arrhenius}{arrhenius}{1}
\goldrow{diag}{rate\_arrhenius}{fitted}{1}
\goldrow{diag}{rate\_arrhenius}{k0}{4e+06}
\goldrow{diag}{rate\_arrhenius}{n\_temps}{3}
\goldrow{diag}{rate\_arrhenius}{order}{2}
\end{tutgolden}

\subsection{membrane}
\grouptheory{membrane}
\clearpage
\tutentry{dspmde01\_water\_treatment\_salts}{tutorials/props/membrane/dspmde01\_water\_treatment\_salts}
\begin{tutscheme}
\begin{tikzpicture}[>={Stealth[length=2mm]}, every node/.style={font=\scriptsize}]
  \node[draw, rounded corners=2pt, align=center, minimum width=3.1cm, minimum height=0.8cm, fill=black!4] at (0.0,-0.0) {\textbf{monovalentWater}\\[-1pt]{\tiny poreTransport}};
  \node[draw, rounded corners=2pt, align=center, minimum width=3.1cm, minimum height=0.8cm, fill=black!4] at (3.6,-0.0) {\textbf{mixedWater}\\[-1pt]{\tiny poreTransport}};
  \node[draw, rounded corners=2pt, align=center, minimum width=3.1cm, minimum height=0.8cm, fill=black!4] at (7.2,-0.0) {\textbf{hardWater}\\[-1pt]{\tiny poreTransport}};
\end{tikzpicture}
\end{tutscheme}
\tutwhat{A water-treatment feed through the DSPM-DE pore: nitrate, fluoride, bromide, carbonate and ammonium salts from the standard catalogue, on NanoFiltran's worked-example pore and the catalogue's NF270 archetype -- every ion's radius and D0 cited on its species record}
\begin{tutnarrative}
THE RESEARCHER'S PATH, PROVED: a water that carries nitrate, fluoride, bromide, carbonate and ammonium, fed as SALTS from the standard catalogue, through the DSPM-DE pore (DEV.md 4c C157/C161/C170).

WHAT KIND OF CHECK THIS IS. A STRUCTURAL witness, not a validation. It shows that every one of these salts decomposes into its ions and that every ion finds a CITED hydrated radius (Nightingale 1959, Table I) and a CITED D0 (USGS PHREEQC phreeqc.dat, `-dw`) on its species record, so the pore can be solved. No number here is a measurement. The two pores are NanoFiltran's worked example (constant/assets/NanoFiltranExample.dat, which describes no real membrane) and the catalogue's NF270\_dspmde (an archetype whose pore is set, not fitted to this water): a researcher replaces them with the record his own fit writes (tutorials/steady/membranes/membrane24\_fit\_then\_simulate is the template).

WHAT IS NOT HERE, AND WHY. Bicarbonate (HCO3-), phosphate and the hydrogen phosphates have a PHREEQC D0 but NO hydrated radius: Nightingale's Table I does not list them, and a Stokes radius from D0 is a different convention that would shift every steric hindrance in silence -- so DSPM-DE REFUSES them by name (DEV.md 4c C161). NaHCO3 in this feed would stop the run with that refusal. The carbonate below is the CO3-- of Na2CO3 taken fully dissociated: nothing here speciates it, so it stands for a high-pH water; at a natural pH most of it would be bicarbonate.

THE THREE OPERATIONS, each a `poreTransport` bench at an IMPOSED J\_v, ideal activities, NanoFiltran's "no polarisation" film (k = 10 m/s), the solved pore at 801 nodes: monovalentWater on NANOFILTRAN'S WORKED-EXAMPLE pore (X -30 mol/m3, r\_p 0.40 nm, 2.5 um, eps\_p 60): NaCl 5, NaNO3 1, NaF 0.5, KBr 0.5, NH4Cl 0.5 mol/m3 (Na K NH4 Cl NO3 F Br) mixedWater the same water plus Na2CO3 0.5 mol/m3 (+ CO3--), on the catalogue's NF270\_dspmde pore (r\_p 0.60 nm, X -5 mol/m3) hardWater Ca(NO3)2 2, KNO3 2, MgCl2 1, NaBr 1, MgSO4 1 mol/m3 (Ca Mg K Na NO3 Cl Br SO4), on NF270\_dspmde WHY TWO PORES. NanoFiltran's example pore was set against NanoFiltran's own ion database, whose radii are STOKES radii (Na+ 1.84, Cl- 1.21, SO4-- 2.31 A; nanofiltran\_geraldesAlves2008/constant/species/). Against the catalogue's HYDRATED radii a 0.40 nm pore nearly closes to the divalent ions: lambda = r/r\_p is 0.985 for CO3--, 0.95 for SO4--, 1.03 for Ca++ and 1.07 for Mg++. Measured 2026-10-26: Na2CO3 (alone, or in this water) and MgSO4 + NaCl do NOT converge there (the solved pore refuses, naming the stalled continuation), and MgCl2 or Ca(NO3)2 alone are refused because no cation enters the pore. So the divalent waters run on NF270\_dspmde, whose 0.60 nm pore is the one its record states is consistent with the hydrated-radius tier. A pore and its ion radii are ONE convention: a pore fitted on one cannot be read with the other.

\end{tutnarrative}
\tutreadme{A water-treatment feed carries more than sodium chloride: nitrate, fluoride, bromide, carbonate, ammonium, calcium, magnesium and sulfate. A case feeds them as salts --- each a component record with an ion map (dissociatesTo) --- and the DSPM-DE pore works ion by ion: every ion needs its hydrated radius (steric partition and hindrance) and its free-solution diffusivity D0 (diffusion and convection in the pore).}
\tutreadme{This case proves that path end to end from the standard catalogue. Its salts are the standard records NaCl, NaNO3, NaF, KBr, Na2CO3, NH4Cl, calciumNitrate (Ca(NO3)2), KNO3, MgCl2, NaBr and MgSO4 (sealed into constant/components/ by bin/choupo-import, so the case is self-contained), and every ion resolves:}
\tutreadme{$\ast$ its hydrated radius from Nightingale, J. Phys. Chem. 63 (1959) 1381, Table I (column r\_H, 25 $^\circ$C), read back against the primary page; $\ast$ its D0 from the public-domain USGS PHREEQC database phreeqc.dat, the 25 $^\circ$C Dw of its -dw line.}
\tutreadme{Each species record (constant/species/<ion>.dat) cites its page, table, database line and file hash in its own transport \{ provenance \{\} \} block.}
\tutreadme{Three operations, each a poreTransport bench at an IMPOSED permeate flux (J\_v = 1e-6, 5e-6, 2e-5 m/s), ideal activities, no polarisation:}
\tutreadme{Each publishes, per J\_v and per ion, the bulk, wall and permeate concentrations and the observed and real rejections. At 2e-5 m/s, for example, mixedWater rejects CO3-- at 0.99985 and NO3- at 0.741, and hardWater rejects SO4-- at 0.99989 against Cl- at 0.778 and Br- at 0.767 -- the divalent co-ion held back by the negative pore, as NF does.}
\begin{tutgolden}
\goldrow{diag}{monovalentWater}{Jv\_p1}{1e-06}
\goldrow{diag}{monovalentWater}{Jv\_p2}{5e-06}
\goldrow{diag}{monovalentWater}{Jv\_p3}{2e-05}
\goldrow{diag}{monovalentWater}{Kc\_Br}{1.19061896103}
\goldrow{diag}{monovalentWater}{Kc\_Cl}{1.18629086518}
\goldrow{diag}{monovalentWater}{Kc\_F}{1.13952821128}
\goldrow{diag}{monovalentWater}{Kc\_K}{1.18846271114}
\goldrow{diag}{monovalentWater}{Kc\_NH4}{1.18846271114}
\goldrow{diag}{monovalentWater}{Kc\_NO3}{1.17968143659}
\goldrow{diag}{monovalentWater}{Kc\_Na}{1.12423947195}
\goldrow{diag}{monovalentWater}{Kd\_Br}{0.0143696929036}
\goldrow{diag}{monovalentWater}{Kd\_Cl}{0.0133665877414}
\goldrow{diag}{monovalentWater}{Kd\_F}{0.0054600222726}
\goldrow{diag}{monovalentWater}{Kd\_K}{0.0138630195272}
\goldmore{1202}
\end{tutgolden}

\clearpage
\tutentry{nanofiltran\_geraldesAlves2008}{tutorials/props/membrane/nanofiltran\_geraldesAlves2008}
\begin{tutscheme}
\begin{tikzpicture}[>={Stealth[length=2mm]}, every node/.style={font=\scriptsize}]
  \node[draw, rounded corners=2pt, align=center, minimum width=3.1cm, minimum height=0.8cm, fill=black!4] at (0.0,-0.0) {\textbf{example\_ideal}\\[-1pt]{\tiny poreTransport}};
  \node[draw, rounded corners=2pt, align=center, minimum width=3.1cm, minimum height=0.8cm, fill=black!4] at (3.6,-0.0) {\textbf{example\_noBorn}\\[-1pt]{\tiny poreTransport}};
  \node[draw, rounded corners=2pt, align=center, minimum width=3.1cm, minimum height=0.8cm, fill=black!4] at (7.2,-0.0) {\textbf{example\_davies}\\[-1pt]{\tiny poreTransport}};
  \node[draw, rounded corners=2pt, align=center, minimum width=3.1cm, minimum height=0.8cm, fill=black!4] at (0.0,-1.3) {\textbf{NaCl\_ideal}\\[-1pt]{\tiny poreTransport}};
  \node[draw, rounded corners=2pt, align=center, minimum width=3.1cm, minimum height=0.8cm, fill=black!4] at (3.6,-1.3) {\textbf{Na2SO4\_ideal}\\[-1pt]{\tiny poreTransport}};
  \node[draw, rounded corners=2pt, align=center, minimum width=3.1cm, minimum height=0.8cm, fill=black!4] at (7.2,-1.3) {\textbf{mixed\_polarised}\\[-1pt]{\tiny poreTransport}};
  \node[draw, rounded corners=2pt, align=center, minimum width=3.1cm, minimum height=0.8cm, fill=black!4] at (0.0,-2.6) {\textbf{dye4\_p1}\\[-1pt]{\tiny poreTransport}};
  \node[draw, rounded corners=2pt, align=center, minimum width=3.1cm, minimum height=0.8cm, fill=black!4] at (3.6,-2.6) {\textbf{dye4\_p2}\\[-1pt]{\tiny poreTransport}};
  \node[draw, rounded corners=2pt, align=center, minimum width=3.1cm, minimum height=0.8cm, fill=black!4] at (7.2,-2.6) {\textbf{dye4\_p3}\\[-1pt]{\tiny poreTransport}};
  \node[draw, rounded corners=2pt, align=center, minimum width=3.1cm, minimum height=0.8cm, fill=black!4] at (0.0,-3.9) {\textbf{dye4\_p4}\\[-1pt]{\tiny poreTransport}};
  \node[draw, rounded corners=2pt, align=center, minimum width=3.1cm, minimum height=0.8cm, fill=black!4] at (3.6,-3.9) {\textbf{dye4\_p5}\\[-1pt]{\tiny poreTransport}};
  \node[draw, rounded corners=2pt, align=center, minimum width=3.1cm, minimum height=0.8cm, fill=black!4] at (7.2,-3.9) {\textbf{dye37\_p1}\\[-1pt]{\tiny poreTransport}};
\end{tikzpicture}
\end{tutscheme}
\tutwhat{The NanoFiltran paper (Geraldes \& Alves 2008) in one case: Choupo against NanoFiltran code to code, against Bowen \& Mohammad's MEASURED CA30 rejections (Fig. 3) and on the stiff case of Table 2 and Fig. 4 -- its tables and figures generated by the views of its propsDict}
\begin{tutnarrative}
V. Geraldes, A.M. Brites Alves, "Computer program for simulation of mass transport in nanofiltration membranes", J. Membr. Sci. 321 (2008) 172-182 (doi:10.1016/j.memsci.2008.04.054) --- NanoFiltran, and the two comparisons its paper makes, run by Choupo's DSPM-DE pore in one case.


\textbf{Where the theory is}

The model and the method are the Theory Guide's DSPM-DE chapter (section 19): the whole boundary-value problem one pore solves --- the extended Nernst-Planck equations, electroneutrality at every node, zero current, the steric, Born and Donnan partitions at both faces, the multi-ionic wall --- in 19.3; the Peclet numbers that decide how hard it is in 19.4; the numerical method (Scharfetter-Gummel fluxes, the block-tridiagonal Newton, the local Donnan seed, the continuation in J\_v) in 19.5; and the node counts below, with what they mean, in 19.8. In the GUI, F1 on any operation of this case opens that chapter at the model, with the numerical method offered beside it.


\textbf{What this case imposes, and what it computes}

It IMPOSES the permeate flux J\_v; it does not compute the flux from the pressure. Every operation is a poreTransport bench: it is handed J\_v (a list; in study B, Bowen \& Mohammad's fitted function evaluated at each operation's feed and written in the operation) and solves the extended Nernst-Planck equations across the pore, with electroneutrality at every node, steric, Born and Donnan partitions at both faces and the multi-ionic film at the wall. So the membrane record's A\_w (the pure-water permeance, J\_w = A\_w (dP - dpi)) is carried because the record grammar requires it and is not read here.

What it publishes, per J\_v point and per ion: the bulk, wall and permeate concentrations, the observed rejection R\_obs = 1 - c\_p/c\_b and the real rejection R\_real = 1 - c\_p/c\_wall, the potential across the membrane and the pore Newton's iteration count.


\textbf{The three studies: what each teaches, and which one is stiff}

Whether a pore is stiff is read off one number per ion, its pore Peclet number Pe\_i = K\_c,i J\_v / (K\_d,i D\_i A\_k/dx) --- convection against hindered diffusion across the layer --- computed here from the case's own ion records and membranes.

\begin{itemize}
\item A --- Code to code: Choupo against NanoFiltran on one model. The same DSPM-DE model and the same inputs, solved by two independent programs (NanoFiltran: Fortran 77, central differences, under-relaxed Newton; Choupo: C++, exponentially fitted fluxes, full Newton). It teaches that the equations are solved right --- the two programs' grid limits agree within 8e-7 on 35 permeates --- and it is NOT a validation: no number here is a measurement, and the membrane is NanoFiltran's worked example, which describes no real membrane. The pore runs from nearly linear (J\_v 1e-6 m/s, Pe of Mg2+ 0.7) to stiff (5e-5 m/s, Pe of Mg2+ 33).
\item B --- Experiment: Bowen \& Mohammad's CA30 rejections (the paper's Fig. 3). It teaches what agreement with measurements means when the inputs were FITTED to those very data by Bowen \& Mohammad (AIChE J. 44 (1998) 1799): the pore radius and thickness, the membrane charge as a function of the feed (Eq. 43) and the permeate flux as a function of the feed (Eqs. 44-45). Agreement here says the pore model, carried by inputs fitted on these data, reproduces them --- not that it predicts a membrane it was not fitted on. This study is NOT stiff: Pe is 0.06-0.14 for Na+ and 0.03-0.07 for Cl-, and the pore profiles are nearly straight. Its physics is at the WALL: the trivalent dye is as large as the pore (lambda = 1), so the pore excludes it; it accumulates in the film and sets the wall's electroneutrality, and that is what drives Cl- to R\_obs = -2.
\item C --- The stiff case of Table 2 and Fig. 4, and its grid study. A numerical study, nothing measured: the paper's "very stiff" pore solved on six node counts, and Choupo's answer against Table 2 accounted for one change at a time. It IS stiff: Pe is 33 for Mg2+ against 0.56 for Na+ and 0.20 for Cl-, because the two counter-ions' hindered diffusivities K\_d D differ 70 times (Mg2+ at lambda = 0.875 barely diffuses). Mg2+ is carried, the trace Na+ diffuses, and the pore holds an interior front where Na+ hands the fixed charge over to Mg2+ (Na+ falls from 12 to 0.16 mol/m3 between a third of the way in and the permeate face). The run prints the stiffness --- the electromigration flux against the net flux --- as 1.65e3 for every node count.
\end{itemize}

\textbf{How many nodes: it depends on what you report}

Measured on study C (re-run 2026-10-25 on a copy of this case whose six N operations were joined by 801 and 1601 nodes, every operation writing its pore profile). Errors against the case's finest grid, 350 nodes, which is itself within 1e-4 relative of the Richardson limit of 801 and 1601 nodes on every row:


{\footnotesize\noindent\begin{tabular}{@{}lrrrrr@{}}
quantity (value at 350 nodes) & 10 nodes & 20 & 50 & 100 & 200 \\
\hline
R\_obs of Na+ (99.4457 \%), in points & -0.073 & -0.017 & -0.0025 & -0.0006 & -0.0001 \\
permeate Na+ (0.05543 mol/m3) & +13.2 \% & +3.03 \% & +0.45 \% & +0.10 \% & +0.019 \% \\
permeate Mg2+ (8.6317 mol/m3) & -0.64 \% & -0.15 \% & -0.023 \% & -0.005 \% & -0.001 \% \\
permeate Cl- (17.3188 mol/m3) & -0.59 \% & -0.14 \% & -0.021 \% & -0.005 \% & -0.001 \% \\
Mg2+ at the pore's permeate end (15.0235 mol/m3) & -0.077 \% & -0.018 \% & -0.0026 \% & -0.0006 \% & -0.0001 \% \\
psi at the pore's permeate end (0.07217 V) & -4.52 \% & -1.08 \% & -0.16 \% & -0.038 \% & -0.007 \% \\
potential across the membrane, dPsi\_m (0.14496 V) & -2.21 \% & -0.53 \% & -0.080 \% & -0.018 \% & -0.003 \% \\
Newton iterations (12 at 350 nodes) & 7 & 9 & 9 & 10 & 11 \\
\end{tabular}}

The lesson: the rejection forgives a coarse grid; the permeate of a trace ion --- and so any selectivity built on it --- does not. Ten nodes put the Na+ rejection within 0.1 point and its permeate 13 \% high. The error falls as the square of the cell size (double the nodes, a quarter of the error), so declare nodes at n and 2n and compare the quantity you report.


\textbf{Read the comparisons for what they are}

\begin{itemize}
\item Nothing was fitted by this run. In study B the parameters were fitted by Bowen \& Mohammad, before this case existed; Choupo fits nothing.
\item The error bars of Fig. 3 are READING uncertainty. The measured symbols were digitised from the paper's figure (calibration checked against the tick labels; reading error +-1.5 px = 0.32 mol/m3 in c\_Cl and 0.011 in R\_obs). They are not the original experimental uncertainty.
\item The feed chloride of each study-B operation is the mean of its two markers' readings --- the Na+ and Cl- symbols of one measurement sit at one feed, and were read at slightly different abscissae.
\item The two programs count N differently. Choupo's N nodes have N - 1 intervals; NanoFiltran's N counts the permeate node, so it has N - 2 intervals inside the membrane.
\item Table 2 carries a factor of ln 10 in its Davies activities, and Choupo deliberately reproduces no such option, so Choupo is NOT pinned to Table 2: it is pinned to NanoFiltran with that constant (and its Born constant and diffusivity input) corrected. The grid-study figure's "Davies corrected" NanoFiltran line has ONLY the Davies constant corrected, so the gap left between it and Choupo is not discretisation alone.
\item 12 against 88 iterations is not a speed measurement. Choupo runs a full Newton with an analytic Jacobian to its own tolerance; NanoFiltran an under-relaxed (0.2) Newton to TOL 1e-4. The criteria, the relaxation and the work per iteration differ.
\end{itemize}
The full account --- every difference between the two programs and how each was handled, the Davies finding, and the term-by-term accounting of Table 2 --- is the leading comment of system/propsDict.

\end{tutnarrative}
\begin{tutgolden}
\goldrow{diag}{example\_ideal}{Robs\_Cl\_p1}{0.202989025795}
\goldrow{diag}{example\_ideal}{Robs\_Cl\_p2}{0.42410107917}
\goldrow{diag}{example\_ideal}{Robs\_Cl\_p3}{0.51695014572}
\goldrow{diag}{example\_ideal}{Robs\_Cl\_p4}{0.603832218837}
\goldrow{diag}{example\_ideal}{Robs\_Cl\_p5}{0.656025668338}
\goldrow{diag}{example\_ideal}{Robs\_Mg\_p1}{0.277594917868}
\goldrow{diag}{example\_ideal}{Robs\_Mg\_p2}{0.518891301308}
\goldrow{diag}{example\_ideal}{Robs\_Mg\_p3}{0.561965456186}
\goldrow{diag}{example\_ideal}{Robs\_Mg\_p4}{0.57192034085}
\goldrow{diag}{example\_ideal}{Robs\_Mg\_p5}{0.571194624339}
\goldrow{diag}{example\_ideal}{Robs\_Na\_p1}{-0.095434542496}
\goldrow{diag}{example\_ideal}{Robs\_Na\_p2}{0.0449401906182}
\goldrow{diag}{example\_ideal}{Robs\_Na\_p3}{0.336888903859}
\goldrow{diag}{example\_ideal}{Robs\_Na\_p4}{0.731479730785}
\goldmore{1909}
\end{tutgolden}

\clearpage
\tutentry{polarisation01\_geraldes\_afonso\_table1}{tutorials/props/membrane/polarisation01\_geraldes\_afonso\_table1}
\begin{tutscheme}
\begin{tikzpicture}[>={Stealth[length=2mm]}, every node/.style={font=\scriptsize}]
  \node[draw, rounded corners=2pt, align=center, minimum width=3.1cm, minimum height=0.8cm, fill=black!4] at (0.0,-0.0) {\textbf{table1}\\[-1pt]{\tiny polarisationIndex}};
  \node[draw, rounded corners=2pt, align=center, minimum width=3.1cm, minimum height=0.8cm, fill=black!4] at (3.6,-0.0) {\textbf{gamma\_dye4}\\[-1pt]{\tiny series}};
  \node[draw, rounded corners=2pt, align=center, minimum width=3.1cm, minimum height=0.8cm, fill=black!4] at (7.2,-0.0) {\textbf{gamma\_dye37}\\[-1pt]{\tiny series}};
  \node[draw, rounded corners=2pt, align=center, minimum width=3.1cm, minimum height=0.8cm, fill=black!4] at (0.0,-1.3) {\textbf{rejection\_dye4}\\[-1pt]{\tiny series}};
  \node[draw, rounded corners=2pt, align=center, minimum width=3.1cm, minimum height=0.8cm, fill=black!4] at (3.6,-1.3) {\textbf{rejection\_dye37}\\[-1pt]{\tiny series}};
\end{tikzpicture}
\end{tutscheme}
\tutwhat{Concentration polarisation of a Na+/Cl-/dye3- solution (Geraldes \& Afonso 2007, Table 1): per-ion k, one interface field, Gamma per ion -- no film thickness}
\begin{tutnarrative}
polarisation01 -- THE CONCENTRATION POLARISATION OF A MULTI-IONIC SOLUTION, computed the way Geraldes \& Afonso did it (J. Membr. Sci. 300 (2007) 20-27, doi:10.1016/j.memsci.2007.04.025) for the case study of their Section 2.2: nanofiltration of Na+ / Cl- / dye3- solutions in a stirred cell (experimental data of Bowen \& Mohammad, AIChE J. 44 (1998) 1799).

THE QUESTION. How much does each ion pile up at the membrane wall when three ions of different diffusivities and charges are rejected to different degrees? The classical film model cannot answer it: it needs ONE effective diffusivity for the solution, and a mixture of ions has none. The extended Nernst-Planck approach answers it only after choosing a film thickness, which nobody measures.

THE MODEL. At the feed-solution/membrane interface, for every ion, the flux J\_v C\_p equals back-diffusion plus convection plus MIGRATION in the electrical potential gradient xi that the ions' different mobilities set up (Eqs. 2-7); electroneutrality at the interface (Eq. 8) fixes xi. The mass-transfer coefficient of each ion comes from the cell's conventional correlation on the ion's OWN diffusivity, Sh\_i = 0.23 Re\textasciicircum{}0.567 Sc\_i\textasciicircum{}0.33 (Bowen et al. 1997), corrected for the suction effect of the permeation by Xi\_i = phi\_i + (1 + 0.26 phi\_i\textasciicircum{}1.4)\textasciicircum{}-1.7 (Geraldes \& Afonso 2006). Four algebraic equations per operating point (Eqs. 25-28), NO film thickness, NO new parameter.

WHAT TO READ IN THE OUTPUT. The polarisation index Gamma\_i = (C\_m - C\_b)/C\_b per ion and per row -- the open symbols of the paper's Fig. 7. The dye (slowest, most rejected) polarises most; Na+ is dragged along by the field; Cl-, less rejected than the dye and pulled by the same field, shows a NEGATIVE Gamma at the low Cl- feeds -- its wall concentration is BELOW the bulk, which the film model applied ion by ion cannot produce. The residual sum z\_i C\_m,i / sum |z\_i| C\_m,i printed under each row is the electroneutrality the field enforces.

WHAT IS AND IS NOT MEASURED. C\_b, C\_p and J\_v are the table's; Na+ is derived from electroneutrality as the table's own footnote says. The Gamma are PREDICTIONS of the model from those inputs and the ions' D -- nothing is fitted. The paper compares them with an extended Nernst-Planck integration in the dye's film (its Fig. 7, filled symbols), which this case does not reproduce.

The gate check\_polarisation\_coupling recomputes the same four equations in Python from this case's own data file and holds every Gamma to 1e-8.

\end{tutnarrative}
\tutreadme{The question. Three ions -- Na+, Cl- and a trivalent dye -- are nanofiltered in a stirred cell. The dye is rejected almost completely, the chloride barely. How much does each pile up at the membrane wall?}
\tutreadme{The lesson. Geraldes \& Afonso, J. Membr. Sci. 300 (2007) 20-27 (doi:10.1016/j.memsci.2007.04.025), answer it with four algebraic equations per operating point and NO film thickness: the cell's conventional correlation Sh\_i = 0.23 Re\textasciicircum{}0.567 Sc\_i\textasciicircum{}0.33 on each ion's OWN diffusivity, the suction correction Xi\_i = phi\_i + (1 + 0.26 phi\_i\textasciicircum{}1.4)\textasciicircum{}-1.7, and one electrical potential gradient at the interface fixed by electroneutrality at the wall. The polarisationIndex bench prints, for each of the ten rows of the paper's Table 1 (data of Bowen \& Mohammad, AIChE J. 44 (1998) 1799), the per-ion k, phi, Xi, wall concentration, polarisation index Gamma\_i = (C\_m - C\_b)/C\_b and the field xi -- the open symbols of the paper's Fig. 7. The dye (slowest, most rejected) polarises most; Na+ is dragged after it by the field; Cl- sits BELOW its bulk concentration (negative Gamma), which the film model applied ion by ion cannot produce.}
\tutreadme{Headline numbers (Gamma): row 1 (dye 4, Cl 5.2 mol/m3) dye +0.504, Na +0.310, Cl -0.137; row 10 (dye 37, Cl 87.3) dye +0.273, Na +0.120, Cl -0.074. Re = 26906 against the paper's 26907.}
\tutreadme{The figures. The run draws them (the GUI's review screen opens on them). The op publishes its DEFAULT views with nothing declared -- Gamma of every ion row by row, the rejection the cell MEASURES (1 - C\_p/C\_b, from Table 1) beside the one the membrane makes (1 - C\_p/C\_m, Choupo's wall), and Table 1 row by row with Gamma\_i and the field -- and the propsDict's views split the same rows by dye level and draw them against the feed chloride at 4 and at 37 mol/m3 of dye (the paper's Fig. 7, open symbols). Every number is the bench's published rows table; the paper's own Gamma values are not in the tree, so none is drawn as if it were.}
\tutreadme{What is measured and what is not. C\_b, C\_p and J\_v are the table's; Na+ is derived by electroneutrality as the table's footnote says (closeBy Na;). Every Gamma is a PREDICTION from those inputs, the ions' D and the correlation; nothing is fitted. The paper's filled symbols (an extended Nernst-Planck integration in the dye's film) are not reproduced. The dye is a case-local species (constant/species/dye.dat, z = -3, D = 0.36e-9 from the paper's footnote); the sources do not identify the compound, so its MW is a flagged nominal placeholder that nothing here reads.}
\tutreadme{The gate. check\_polarisation\_coupling re-solves Eqs. 25-28 in Python from this case's own data file and holds every Gamma AND xi to 1e-8, holds a single salt to the paper's closed form Eq. 20, and holds the filmTheory + none policy to the classical film model to round-off.}
\begin{tutgolden}
\goldrow{diag}{table1}{D\_Cl}{2.03e-09}
\goldrow{diag}{table1}{D\_Na}{1.33e-09}
\goldrow{diag}{table1}{D\_dye}{3.6e-10}
\goldrow{diag}{table1}{Gamma\_Cl\_r1}{-0.136799297507}
\goldrow{diag}{table1}{Gamma\_Cl\_r10}{-0.0737363093742}
\goldrow{diag}{table1}{Gamma\_Cl\_r2}{-0.095324747833}
\goldrow{diag}{table1}{Gamma\_Cl\_r3}{-0.0498126163785}
\goldrow{diag}{table1}{Gamma\_Cl\_r4}{-0.0178510515183}
\goldrow{diag}{table1}{Gamma\_Cl\_r5}{-0.00198770276507}
\goldrow{diag}{table1}{Gamma\_Cl\_r6}{-0.24271484424}
\goldrow{diag}{table1}{Gamma\_Cl\_r7}{-0.176886232679}
\goldrow{diag}{table1}{Gamma\_Cl\_r8}{-0.13368256672}
\goldrow{diag}{table1}{Gamma\_Cl\_r9}{-0.10815205541}
\goldrow{diag}{table1}{Gamma\_Na\_r1}{0.310336981862}
\goldmore{55}
\end{tutgolden}

\clearpage
\tutentry{shortlist01\_class\_to\_modules}{tutorials/props/membrane/shortlist01\_class\_to\_modules}
\begin{tutscheme}
\begin{tikzpicture}[>={Stealth[length=2mm]}, every node/.style={font=\scriptsize}]
  \node[draw, rounded corners=2pt, align=center, minimum width=3.1cm, minimum height=0.8cm, fill=black!4] at (0.0,-0.0) {\textbf{shortlist}\\[-1pt]{\tiny membraneShortlist}};
\end{tikzpicture}
\end{tutscheme}
\tutwhat{From a membrane class to the commercial elements that do it: hard filters, each survivor re-rated at the student's conditions, every criterion side by side, no score}
\begin{tutnarrative}
shortlist01 -- THE SHORTLIST (DEV.md 4c C189 slice 5): from a membrane class to the commercial elements that do what the class does.

THE QUESTION. membrane26 chose a membrane by WHAT IT DOES (NF\_DIVALENT\_SELECTIVE: hold sulphate, let sodium chloride through). Which products in the catalogue do that, at MY pressure, MY temperature, MY feed and MY recovery -- not at the data sheet's?

WHAT THE OP DOES. (1) Hard filters on each module record: the membrane is of the class's type (NF here), the record carries ONE rated test of MgSO4, the sheet's rated rejection lies inside the class's band (>= 90 \%), the declared limits hold at this feed. (2) Each survivor is re-run with the transport law and film of the `operation` block below, its intrinsic rejection DERIVED from its OWN rated test (slice 1), at the conditions below, with the feed flow found so ONE element delivers the recovery asked. Then the limits that need the run (feed flow, pressure drop) filter again.

WHAT IT DOES NOT DO. It scores nothing and ranks nothing. Every criterion is published side by side, the candidates in alphabetical order of their names, and a record's commercial-relation fields (`partner`, `dataVerifiedBy`) are never read. What matters most -- permeate flow, permeate quality, specific pressure -- is the student's decision.

WHAT TO READ. The printed rows (and shortlist.csv): permeate flow and concentration, flux, P/flux (bar per L/m2/h), the hydraulic energy per m3 of permeate (feed pressure x feed flow / permeate flow, no recovery device), the feed flow, the element's pressure drop and each limit's margin ((limit - value)/limit; pH in pH units; "n/s" = the record states no such limit). EXCLUDED rows say why, in the record's own numbers.

\end{tutnarrative}
\tutreadme{The question. membrane26 chose a membrane by what it DOES (NF\_DIVALENT\_SELECTIVE: hold sulphate, let sodium chloride through). Which products in the catalogue do that -- at MY pressure, MY temperature, MY feed and MY recovery, not at the data sheet's?}
\tutreadme{The op. membraneShortlist (DEV.md 4c C189 slice 5). It does two things and nothing else.}
\tutreadme{1. $\ast$Hard filters$\ast$, on each module record: the membrane is of the class's type (NF here), the record carries ONE rated test of the asked solute, that test's rated rejection lies inside the class's band, the declared limits (P\_max, T\_max, pH) hold at the student's feed. A limit the record does not state is not a filter (n/s). 2. $\ast$Re-rating$\ast$: every survivor is re-run with the film and spacer of the operation block, its intrinsic rejection DERIVED from its OWN rated test (slice 1), at the student's conditions, with the feed flow found so ONE element delivers the recovery asked. The limits that need the run (feed flow, pressure drop) then filter again.}
\tutreadme{What it will not do. Score, weight or rank. Every criterion is printed side by side (permeate flow, observed rejection, permeate concentration, flux, P/flux, hydraulic energy per m3, feed flow, pressure drop, margins to each limit), the candidates in alphabetical order of their names, and what matters most is the student's decision. A record's commercial-relation fields (partner, dataVerifiedBy) are never read: the gate proves it.}
\tutreadme{What to expect here. tiers ( standards case ) keeps the golden independent of one machine's private data/local/assets/ records (the importer's home); drop the line and an imported catalogue joins the table, each of its records announced [local] UNVERIFIED. Of the shipped records NF270-4040 passes; SW30HR-380 is excluded (an RO membrane, rated on NaCl); the Generic elements are excluded (they carry no rated test).}
\tutreadme{Read it with care. The re-rating is the transport law of a lumped salt with its intrinsic rejection fixed at the sheet's ONE test -- a model's reading of a data sheet away from its test, not a guarantee of the product.}
\begin{tutgolden}
\goldrow{diag}{shortlist}{NF270\_4040\_\_A\_w\_LMH\_per\_bar}{14.8806573771}
\goldrow{diag}{shortlist}{NF270\_4040\_\_R\_int}{0.979346761192}
\goldrow{diag}{shortlist}{NF270\_4040\_\_R\_obs}{0.969702755801}
\goldrow{diag}{shortlist}{NF270\_4040\_\_dP\_bar}{0.430832573708}
\goldrow{diag}{shortlist}{NF270\_4040\_\_feedFlow\_m3h}{3.46318793226}
\goldrow{diag}{shortlist}{NF270\_4040\_\_flux\_LMH}{68.3523937423}
\goldrow{diag}{shortlist}{NF270\_4040\_\_margin\_P}{0.829268292683}
\goldrow{diag}{shortlist}{NF270\_4040\_\_margin\_T}{0.0628634292001}
\goldrow{diag}{shortlist}{NF270\_4040\_\_margin\_dP}{0.569167426292}
\goldrow{diag}{shortlist}{NF270\_4040\_\_margin\_feedFlow}{0.0380033521509}
\goldrow{diag}{shortlist}{NF270\_4040\_\_margin\_pH}{2.5}
\goldrow{diag}{shortlist}{NF270\_4040\_\_permeateConc\_ppm}{60.5944883985}
\goldrow{diag}{shortlist}{NF270\_4040\_\_permeateFlow\_m3d}{12.4674766186}
\goldrow{diag}{shortlist}{NF270\_4040\_\_recovery}{0.150000000752}
\goldmore{10}
\end{tutgolden}

\subsection{molecular}
\grouptheory{molecular}
\clearpage
\tutentry{componentLookup01\_catalogue\_ladder}{tutorials/props/molecular/componentLookup01\_catalogue\_ladder}
\begin{tutscheme}
\begin{tikzpicture}[>={Stealth[length=2mm]}, every node/.style={font=\scriptsize}]
  \node[draw, rounded corners=2pt, align=center, minimum width=3.1cm, minimum height=0.8cm, fill=black!4] at (0.0,-0.0) {\textbf{c\_name}\\[-1pt]{\tiny componentLookup}};
  \node[draw, rounded corners=2pt, align=center, minimum width=3.1cm, minimum height=0.8cm, fill=black!4] at (3.6,-0.0) {\textbf{c\_cas}\\[-1pt]{\tiny componentLookup}};
  \node[draw, rounded corners=2pt, align=center, minimum width=3.1cm, minimum height=0.8cm, fill=black!4] at (7.2,-0.0) {\textbf{c\_lake}\\[-1pt]{\tiny componentLookup}};
  \node[draw, rounded corners=2pt, align=center, minimum width=3.1cm, minimum height=0.8cm, fill=black!4] at (0.0,-1.3) {\textbf{c\_none}\\[-1pt]{\tiny componentLookup}};
\end{tikzpicture}
\end{tutscheme}
\tutwhat{componentLookup: a catalogue record by name and by CAS, a group-estimate-lake compound, and an unknown}
\begin{tutnarrative}
WHAT IT IS. "Does Choupo have this compound, and what does its record carry?" as ONE operation with a fixed byte budget (2048 bytes) and the same evidence-card format as pairLookup (DEV.md 4c C227, slice a). The query resolves by RECORD -- a catalogue file name, a declared alias or a CAS -- never by name similarity. Four operations:

c\_name ethanol the catalogue record, by name c\_cas 64-17-5 the SAME record, by CAS c\_lake 1Pentadecanol only in the group-estimate lake: identity and UNIFAC groups, NO values, "not reviewed" c\_none unobtainium "I do not have this compound"

The card lists which properties the record has and which it lacks; it never prints the values (they are in the record, and a case uses the name). Pinned: the tier word of each card and its byte count. The thermoPhysPropDict is a placeholder the bench requires; no op reads it.

\end{tutnarrative}
\begin{tutgolden}
\goldrow{diag}{c\_name}{budget\_bytes}{2048}
\goldrow{diag}{c\_name}{card\_bytes}{807}
\goldrow{diag}{c\_cas}{budget\_bytes}{2048}
\goldrow{diag}{c\_cas}{card\_bytes}{806}
\goldrow{diag}{c\_lake}{budget\_bytes}{2048}
\goldrow{diag}{c\_lake}{card\_bytes}{663}
\goldrow{diag}{c\_none}{budget\_bytes}{2048}
\goldrow{diag}{c\_none}{card\_bytes}{535}
\goldrow{verdict}{c\_name}{tier}{catalogue}
\goldrow{verdict}{c\_cas}{tier}{catalogue}
\goldrow{verdict}{c\_lake}{tier}{estimateLake}
\goldrow{verdict}{c\_none}{tier}{none}
\end{tutgolden}

\clearpage
\tutentry{cosmoSAC01\_water\_ethanol}{tutorials/props/molecular/cosmoSAC01\_water\_ethanol}
\begin{tutscheme}
\begin{tikzpicture}[>={Stealth[length=2mm]}, every node/.style={font=\scriptsize}]
  \node[draw, rounded corners=2pt, align=center, minimum width=3.1cm, minimum height=0.8cm, fill=black!4] at (0.0,-0.0) {\textbf{pure\_water}\\[-1pt]{\tiny propertyScan1D}};
  \node[draw, rounded corners=2pt, align=center, minimum width=3.1cm, minimum height=0.8cm, fill=black!4] at (3.6,-0.0) {\textbf{bin\_wet}\\[-1pt]{\tiny propertyScan1D}};
  \node[draw, rounded corners=2pt, align=center, minimum width=3.1cm, minimum height=0.8cm, fill=black!4] at (7.2,-0.0) {\textbf{ternary}\\[-1pt]{\tiny propertyScan1D}};
  \node[draw, rounded corners=2pt, align=center, minimum width=3.1cm, minimum height=0.8cm, fill=black!4] at (0.0,-1.3) {\textbf{water\_hexane}\\[-1pt]{\tiny propertyScan1D}};
\end{tikzpicture}
\end{tutscheme}
\tutwhat{COSMO-SAC 2002 (NIST benchmark, VT-2005 profiles): pure$\to$lngamma=0, binary/ternary multicomponent gammas, no NaN.}
\begin{tutgolden}
\goldrow{diag}{pure\_water}{n\_failures}{0}
\goldrow{diag}{pure\_water}{n\_points}{2}
\goldrow{diag}{bin\_wet}{n\_failures}{0}
\goldrow{diag}{bin\_wet}{n\_points}{2}
\goldrow{diag}{ternary}{n\_failures}{0}
\goldrow{diag}{ternary}{n\_points}{2}
\goldrow{diag}{water\_hexane}{n\_failures}{0}
\goldrow{diag}{water\_hexane}{n\_points}{2}
\goldrow{diag}{pure\_water}{n\_emptyCurves}{0}
\goldrow{diag}{bin\_wet}{n\_emptyCurves}{0}
\goldrow{diag}{ternary}{n\_emptyCurves}{0}
\goldrow{diag}{water\_hexane}{n\_emptyCurves}{0}
\end{tutgolden}

\clearpage
\tutentry{entropy01\_air\_ledger}{tutorials/props/molecular/entropy01\_air\_ledger}
\begin{tutscheme}
\begin{tikzpicture}[>={Stealth[length=2mm]}, every node/.style={font=\scriptsize}]
  \node[draw, rounded corners=2pt, align=center, minimum width=3.1cm, minimum height=0.8cm, fill=black!4] at (0.0,-0.0) {\textbf{ledger}\\[-1pt]{\tiny propertyPoint}};
  \node[draw, rounded corners=2pt, align=center, minimum width=3.1cm, minimum height=0.8cm, fill=black!4] at (3.6,-0.0) {\textbf{explainS}\\[-1pt]{\tiny explainProperty}};
  \node[draw, rounded corners=2pt, align=center, minimum width=3.1cm, minimum height=0.8cm, fill=black!4] at (7.2,-0.0) {\textbf{explainH}\\[-1pt]{\tiny explainProperty}};
  \node[draw, rounded corners=2pt, align=center, minimum width=3.1cm, minimum height=0.8cm, fill=black!4] at (0.0,-1.3) {\textbf{temperatureLine}\\[-1pt]{\tiny propertyScan1D}};
  \node[draw, rounded corners=2pt, align=center, minimum width=3.1cm, minimum height=0.8cm, fill=black!4] at (3.6,-1.3) {\textbf{pressureLine}\\[-1pt]{\tiny propertyScan1D}};
\end{tikzpicture}
\end{tutscheme}
\tutwhat{The entropy ledger of 79/21 N2/O2 at 400 K, 2 bar on SRK: every line of s(T,P,y) published separately, so a reader can rebuild the total by hand}
\begin{tutnarrative}
THE ENTROPY LEDGER, published line by line -- the witness behind the "What is entropy?" EduTool page.

The engine's mixture entropy is assembled in ThermoPackage::S\_ig (src/thermo/ThermoPackage.cpp) as

S\_ig(T,P,y) = sum\_i y\_i $\ast$ s\_ig\_i(T) (pure-state line) - R $\ast$ sum\_i y\_i $\ast$ ln y\_i (mixing line) - R $\ast$ ln(P / 1 bar) (pressure line)

with each s\_ig\_i(T) = s\_298\_i + integral Cp\_i/T dT off the component's own .dat (Component::s\_formation), and S\_real = S\_ig + S\_residual(SRK) (the model line, Sandler eq. 6.4-31). The `ledger` op below publishes the per-component s\_ig\_i AND the assembled S\_ig / S\_R / S\_real at one state, so a reader -- or the EduTool page in the browser -- can rebuild the total from the lines and check it against the engine's own number. The two scans draw the temperature line (integral Cp/T at fixed P) and the pressure line (-R ln P at fixed T, with the residual gap growing).

Nothing here is a validation against measurement: every number is the engine's own answer, decomposed. The value of the case is that the decomposition CANNOT drift from the simulator, because it is published by the same call that computes it.

\end{tutnarrative}
\begin{tutgolden}
\goldrow{diag}{ledger}{Cp\_ig}{29.46367008}
\goldrow{diag}{ledger}{H\_R}{-5.48583758685}
\goldrow{diag}{ledger}{H\_ig}{2986.16894826}
\goldrow{diag}{ledger}{H\_real}{2980.68311067}
\goldrow{diag}{ledger}{P}{200000}
\goldrow{diag}{ledger}{S\_R}{-0.0191625058917}
\goldrow{diag}{ledger}{S\_ig}{201.577774992}
\goldrow{diag}{ledger}{S\_real}{201.558612487}
\goldrow{diag}{ledger}{T}{400}
\goldrow{diag}{ledger}{Z}{1.00065691172}
\goldrow{diag}{ledger}{gamma}{1.3931335315}
\goldrow{diag}{ledger}{s\_ig\_N2}{200.188720505}
\goldrow{diag}{ledger}{s\_ig\_O2}{213.897898046}
\goldrow{diag}{ledger}{v\_molar}{0.0166398489719}
\goldmore{36}
\end{tutgolden}

\clearpage
\tutentry{envelope01\_natural\_gas}{tutorials/props/molecular/envelope01\_natural\_gas}
\begin{tutscheme}
\begin{tikzpicture}[>={Stealth[length=2mm]}, every node/.style={font=\scriptsize}]
  \node[draw, rounded corners=2pt, align=center, minimum width=3.1cm, minimum height=0.8cm, fill=black!4] at (0.0,-0.0) {\textbf{envelope}\\[-1pt]{\tiny phaseEnvelope}};
\end{tikzpicture}
\end{tutscheme}
\tutwhat{P-T phase envelope of a lean natural gas: the dew branch over its cricondentherm, and a bubble branch that never starts}
\begin{tutgolden}
\goldrow{diag}{envelope}{bubble\_points}{0}
\goldrow{diag}{envelope}{bubble\_stopped\_bar}{1}
\goldrow{diag}{envelope}{dew\_cricondentherm\_K}{275.946403666}
\goldrow{diag}{envelope}{dew\_points}{40}
\end{tutgolden}

\clearpage
\tutentry{envelope02\_c3c4c5\_branch\_stop}{tutorials/props/molecular/envelope02\_c3c4c5\_branch\_stop}
\begin{tutscheme}
\begin{tikzpicture}[>={Stealth[length=2mm]}, every node/.style={font=\scriptsize}]
  \node[draw, rounded corners=2pt, align=center, minimum width=3.1cm, minimum height=0.8cm, fill=black!4] at (0.0,-0.0) {\textbf{envelope}\\[-1pt]{\tiny phaseEnvelope}};
\end{tikzpicture}
\end{tutscheme}
\tutwhat{P-T envelope of a C3/C4/C5 feed marched through propane's Tc: the reference factor is held continuous above 0.99 Tc, so both branches cross Tc; with the cubic's vapour pseudo-root both reach 40 bar (C245), and no cricondentherm is claimed}
\begin{tutgolden}
\goldrow{diag}{envelope}{bubble\_points}{40}
\goldrow{diag}{envelope}{dew\_points}{40}
\end{tutgolden}

\clearpage
\tutentry{eval\_nrtl\_catalogue}{tutorials/props/molecular/eval\_nrtl\_catalogue}
\begin{tutscheme}
\begin{tikzpicture}[>={Stealth[length=2mm]}, every node/.style={font=\scriptsize}]
  \node[draw, rounded corners=2pt, align=center, minimum width=3.1cm, minimum height=0.8cm, fill=black!4] at (0.0,-0.0) {\textbf{catalogue\_as\_assembled}\\[-1pt]{\tiny fitParameters}};
\end{tikzpicture}
\end{tutscheme}
\tutwhat{EVALUATE mode: score the CATALOGUE NRTL pair against your own data -- the scale before any fit}
\begin{tutnarrative}
`mode evaluate`: the SCALE, not the fit.

ONE operation, one question: how good is the ethanol-water NRTL pair against this 1-atm bubble-T dataset? Nothing is adjusted, nothing is written. The thermoPhysPropDict declares NO inline pairs, so what gets scored is a REAL provenance-bearing pair resolved through the search path (from this case's own vendored constant/parameters/ copy, identical parameters and full primary citation) -- not numbers typed inline into the dict.

Read the verdict (rms/aad in kelvin) and THEN decide whether a refit is worth it: compare with fitNRTL01's fitted rms (0.31 K) on this same data. Failures are honest: a non-converged point becomes status=nonconverged in the parity CSV and never enters the statistics; no chi2\_reduced is printed -- without sigmas and degrees of freedom from THIS act, the name would lie.

The pinned-what-if grammar and the leakage sentinel live in the sibling case eval\_nrtl\_pinned (a pinned path must EXIST in the dict, so what-ifs run on an INLINE scaffold and say so).

\end{tutnarrative}
\begin{tutgolden}
\goldrow{diag}{catalogue\_as\_assembled}{aad}{2.0400}
\goldrow{diag}{catalogue\_as\_assembled}{chi2}{60.1945}
\goldrow{diag}{catalogue\_as\_assembled}{max\_abs\_resid}{3.4379}
\goldrow{diag}{catalogue\_as\_assembled}{max\_resid\_x1}{0.5000}
\goldrow{diag}{catalogue\_as\_assembled}{mode\_evaluate}{1.0000}
\goldrow{diag}{catalogue\_as\_assembled}{n\_data}{11.0000}
\goldrow{diag}{catalogue\_as\_assembled}{n\_feasible}{11.0000}
\goldrow{diag}{catalogue\_as\_assembled}{partial}{0.0000}
\goldrow{diag}{catalogue\_as\_assembled}{penalised\_points}{0.0000}
\goldrow{diag}{catalogue\_as\_assembled}{rms}{2.3393}
\goldrow{diag}{catalogue\_as\_assembled}{sse}{60.1945}
\goldrow{diag}{catalogue\_as\_assembled}{validity\_declared}{1.0000}
\goldrow{diag}{catalogue\_as\_assembled}{n\_outside\_validity}{0.0000}
\end{tutgolden}

\clearpage
\tutentry{eval\_nrtl\_pinned}{tutorials/props/molecular/eval\_nrtl\_pinned}
\begin{tutscheme}
\begin{tikzpicture}[>={Stealth[length=2mm]}, every node/.style={font=\scriptsize}]
  \node[draw, rounded corners=2pt, align=center, minimum width=3.1cm, minimum height=0.8cm, fill=black!4] at (0.0,-0.0) {\textbf{scaffold\_as\_is}\\[-1pt]{\tiny fitParameters}};
  \node[draw, rounded corners=2pt, align=center, minimum width=3.1cm, minimum height=0.8cm, fill=black!4] at (3.6,-0.0) {\textbf{rough\_paper\_values}\\[-1pt]{\tiny fitParameters}};
  \node[draw, rounded corners=2pt, align=center, minimum width=3.1cm, minimum height=0.8cm, fill=black!4] at (7.2,-0.0) {\textbf{scaffold\_again}\\[-1pt]{\tiny fitParameters}};
\end{tikzpicture}
\end{tutscheme}
\tutwhat{EVALUATE a pinned what-if: a pinned run scores the package but mutates NOTHING (the leakage sentinel)}
\begin{tutnarrative}
the pinned WHAT-IF, and the LEAKAGE SENTINEL.

TRANSITIONAL SCAFFOLD: pinning by dict path requires an inline pair; the permanent surface (ParameterOverrides on the RESOLVED parameter registry, design forum \#67/\#69) replaces this case when it lands.

Three evaluations of the same dataset: 1. scaffold\_as\_is -- the inline pair, untouched (the baseline). 2. rough\_paper\_values -- the SAME dict with four coefficients PINNED to deliberately rough values. The pin operates on a DEEP COPY (Dictionary::deepCopy, forum \#63-7); nothing on disk or in memory survives the op. 3. scaffold\_again -- repeats 1 AFTER the pinned run. Its verdict must be IDENTICAL to 1's: the golden pins both, so the working-tree isolation is guarded by the corpus forever. Before deepCopy existed, the old alias bug would make THIS operation report the rough values' score.

\end{tutnarrative}
\begin{tutgolden}
\goldrow{diag}{scaffold\_as\_is}{aad}{2.0400}
\goldrow{diag}{scaffold\_as\_is}{chi2}{60.1945}
\goldrow{diag}{scaffold\_as\_is}{max\_abs\_resid}{3.4379}
\goldrow{diag}{scaffold\_as\_is}{max\_resid\_x1}{0.5000}
\goldrow{diag}{scaffold\_as\_is}{mode\_evaluate}{1.0000}
\goldrow{diag}{scaffold\_as\_is}{n\_data}{11.0000}
\goldrow{diag}{scaffold\_as\_is}{n\_feasible}{11.0000}
\goldrow{diag}{scaffold\_as\_is}{partial}{0.0000}
\goldrow{diag}{scaffold\_as\_is}{penalised\_points}{0.0000}
\goldrow{diag}{scaffold\_as\_is}{rms}{2.3393}
\goldrow{diag}{scaffold\_as\_is}{sse}{60.1945}
\goldrow{diag}{rough\_paper\_values}{aad}{1.5942}
\goldrow{diag}{rough\_paper\_values}{chi2}{61.7621}
\goldrow{diag}{rough\_paper\_values}{max\_abs\_resid}{5.5866}
\goldmore{21}
\end{tutgolden}

\clearpage
\tutentry{exergy01\_air\_dead\_state}{tutorials/props/molecular/exergy01\_air\_dead\_state}
\begin{tutscheme}
\begin{tikzpicture}[>={Stealth[length=2mm]}, every node/.style={font=\scriptsize}]
  \node[draw, rounded corners=2pt, align=center, minimum width=3.1cm, minimum height=0.8cm, fill=black!4] at (0.0,-0.0) {\textbf{b\_state}\\[-1pt]{\tiny exergy}};
  \node[draw, rounded corners=2pt, align=center, minimum width=3.1cm, minimum height=0.8cm, fill=black!4] at (3.6,-0.0) {\textbf{b\_dead}\\[-1pt]{\tiny exergy}};
\end{tikzpicture}
\end{tutscheme}
\tutwhat{Physical exergy of 79/21 N2/O2 at 400 K, 2 bar against a DECLARED 298.15 K / 1 bar dead state -- b = (h-h0) - T0(s-s0) from the engine's own H\_real/S\_real, with the structural zero at the dead state itself as the second row}
\begin{tutnarrative}
PHYSICAL EXERGY, the number the entropy ledger buys -- the witness behind the exergy work that follows the "What is entropy?" page.

b = (h - h0) - T0$\ast$(s - s0), every term the engine's own H\_real/S\_real at the SAME composition, so the enthalpy datum, the s\_298 anchors, the mixing term and the reference pressure ALL cancel: the exergy is datum-independent. The dead state is DECLARED (an exergy is a statement about a state AND an environment; the engine refuses to assume one).

Two rows on purpose: the same state the entropy01 ledger decomposes, and the dead state itself -- whose exergy is 0 identically, by construction, which the golden pins as the structural zero.

Nothing here is validated against measurement: every number is the engine's own answer, decomposed. The published dh and T0ds legs re-add to b on the reader's desk.

\end{tutnarrative}
\begin{tutgolden}
\goldrow{diag}{b\_state}{P}{200000}
\goldrow{diag}{b\_state}{P0}{100000}
\goldrow{diag}{b\_state}{T}{400}
\goldrow{diag}{b\_state}{T0}{298.15}
\goldrow{diag}{b\_state}{T0ds}{850.601}
\goldrow{diag}{b\_state}{b\_physical}{2136.71}
\goldrow{diag}{b\_state}{dh}{2987.32}
\goldrow{diag}{b\_state}{h}{2980.68}
\goldrow{diag}{b\_state}{h0}{-6.63211}
\goldrow{diag}{b\_state}{s}{201.559}
\goldrow{diag}{b\_state}{s0}{198.706}
\goldrow{diag}{b\_dead}{P}{100000}
\goldrow{diag}{b\_dead}{P0}{100000}
\goldrow{diag}{b\_dead}{T}{298.15}
\goldmore{8}
\end{tutgolden}

\clearpage
\tutentry{flash01\_operating\_line}{tutorials/props/molecular/flash01\_operating\_line}
\begin{tutscheme}
\begin{tikzpicture}[>={Stealth[length=2mm]}, every node/.style={font=\scriptsize}]
  \node[draw, rounded corners=2pt, align=center, minimum width=3.1cm, minimum height=0.8cm, fill=black!4] at (0.0,-0.0) {\textbf{txy}\\[-1pt]{\tiny propertyScan1D}};
\end{tikzpicture}
\end{tutscheme}
\tutwhat{The equilibrium curve a flash operating line is drawn on}
\begin{tutnarrative}
THE CURVE, AND ONLY THE CURVE.

A binary flash is fixed by two of \{T, P, V/F\} -- Duhem -- so at a frozen pressure the equilibrium curve y$\ast$(x) is frozen too. Everything the flash tool draws on top of it (the operating line, its intersection, the lever rule) is GRAPHICAL construction on this one table, which is why turning a knob there redraws at once instead of re-solving.

So this case publishes exactly one thing: the curve. It does NOT flash anything. The flash the student watches is the intersection, performed on these rows -- and the corpus already carries the ENGINE's own answer for the same pair (tutorials/steady/flash/flash01\_benzene\_toluene, V/F = 0.30398 at 370 K), which is the number to check the construction against.

51 points is not a rounding of taste: the construction interpolates linearly between rows (no spline -- an overshoot near an azeotrope would invent equilibrium the engine never computed), so the sampling IS the resolution of every answer read off it.

\end{tutnarrative}
\begin{tutgolden}
\goldrow{diag}{txy}{n\_emptyCurves}{0}
\goldrow{diag}{txy}{n\_failures}{0}
\goldrow{diag}{txy}{n\_points}{51}
\end{tutgolden}

\clearpage
\tutentry{mixrules01\_natural\_gas}{tutorials/props/molecular/mixrules01\_natural\_gas}
\begin{tutscheme}
\begin{tikzpicture}[>={Stealth[length=2mm]}, every node/.style={font=\scriptsize}]
  \node[draw, rounded corners=2pt, align=center, minimum width=3.1cm, minimum height=0.8cm, fill=black!4] at (0.0,-0.0) {\textbf{mix\_30bar}\\[-1pt]{\tiny mixingRules}};
  \node[draw, rounded corners=2pt, align=center, minimum width=3.1cm, minimum height=0.8cm, fill=black!4] at (3.6,-0.0) {\textbf{mix\_1bar}\\[-1pt]{\tiny mixingRules}};
\end{tikzpicture}
\end{tutscheme}
\tutwhat{The van der Waals one-fluid mixing rule, decomposed: SRK a\_mix/b\_mix of a 90/5/5 CH4/N2/CO2 gas at 250 K, pair by pair -- the one curated kij (N2-CH4, DECHEMA) beside the pairs that run the kij = 0 default, and the same decomposition at 1 bar showing a\_mix is a (T, y) fact while only A, B, Z carry the pressure}
\begin{tutnarrative}
THE MIXING RULE, DECOMPOSED -- the multicomponent lesson. A cubic prices a MIXTURE through two numbers, and the one-fluid rule says how n pure components and n(n-1)/2 binary corrections build them:

a\_mix(T) = sum\_i sum\_j y\_i y\_j (1 - kij) sqrt(a\_i a\_j) b\_mix = sum\_i y\_i b\_i

Every term is on the table: the six pair terms of this ternary re-add to the a\_mix the flash would run (the op REFUSES if they do not), the curated N2-CH4 kij sits beside the two CO2 pairs running the announced kij = 0 default, and the pair table is where a student sees that a ternary needs THREE binary decisions -- there is no ternary parameter.

Two states on purpose, same composition and temperature: a\_mix and b\_mix are IDENTICAL at 30 bar and 1 bar (the rule sees only T and y), while A, B and Z carry the pressure -- at 1 bar the same gas is nearly ideal (Z \textasciitilde{} 0.997) and at 30 bar it is not (Z \textasciitilde{} 0.93), priced by the SAME two numbers through the cubic.

Nothing here is validated against measurement: every number is the engine's own arithmetic, decomposed. The one cited constant is the N2-CH4 kij record (Knapp et al., DECHEMA VI, 1982).

\end{tutnarrative}
\begin{tutgolden}
\goldrow{diag}{mix\_30bar}{A}{0.140523}
\goldrow{diag}{mix\_30bar}{B}{0.0428434}
\goldrow{diag}{mix\_30bar}{P\_Pa}{3e+06}
\goldrow{diag}{mix\_30bar}{T\_K}{250}
\goldrow{diag}{mix\_30bar}{Z\_vapour}{0.901046}
\goldrow{diag}{mix\_30bar}{a\_CH4}{0.200506}
\goldrow{diag}{mix\_30bar}{a\_CO2}{0.429536}
\goldrow{diag}{mix\_30bar}{a\_N2}{0.0837883}
\goldrow{diag}{mix\_30bar}{a\_mix}{0.202383}
\goldrow{diag}{mix\_30bar}{b\_CH4}{2.98484e-05}
\goldrow{diag}{mix\_30bar}{b\_CO2}{2.96742e-05}
\goldrow{diag}{mix\_30bar}{b\_N2}{2.6754e-05}
\goldrow{diag}{mix\_30bar}{b\_mix}{2.9685e-05}
\goldrow{diag}{mix\_30bar}{dadT\_mix}{-0.000515994}
\goldmore{46}
\end{tutgolden}

\clearpage
\tutentry{nrtl01\_ethanol\_water\_package}{tutorials/props/molecular/nrtl01\_ethanol\_water\_package}
\begin{tutscheme}
\begin{tikzpicture}[>={Stealth[length=2mm]}, every node/.style={font=\scriptsize}]
  \node[draw, rounded corners=2pt, align=center, minimum width=3.1cm, minimum height=0.8cm, fill=black!4] at (0.0,-0.0) {\textbf{g\_nrtl}\\[-1pt]{\tiny activityCoefficients}};
\end{tikzpicture}
\end{tutscheme}
\tutwhat{ethanol-water NRTL via the SAME property-package builder (molecular U4)}
\begin{tutgolden}
\goldrow{diag}{g\_nrtl}{gamma\_ethanol}{1.26245}
\goldrow{diag}{g\_nrtl}{gamma\_water}{1.4349}
\end{tutgolden}

\clearpage
\tutentry{pairLookup01\_catalogue\_ladder}{tutorials/props/molecular/pairLookup01\_catalogue\_ladder}
\begin{tutscheme}
\begin{tikzpicture}[>={Stealth[length=2mm]}, every node/.style={font=\scriptsize}]
  \node[draw, rounded corners=2pt, align=center, minimum width=3.1cm, minimum height=0.8cm, fill=black!4] at (0.0,-0.0) {\textbf{p\_curated}\\[-1pt]{\tiny pairLookup}};
  \node[draw, rounded corners=2pt, align=center, minimum width=3.1cm, minimum height=0.8cm, fill=black!4] at (3.6,-0.0) {\textbf{p\_regressed}\\[-1pt]{\tiny pairLookup}};
  \node[draw, rounded corners=2pt, align=center, minimum width=3.1cm, minimum height=0.8cm, fill=black!4] at (7.2,-0.0) {\textbf{p\_estimate}\\[-1pt]{\tiny pairLookup}};
  \node[draw, rounded corners=2pt, align=center, minimum width=3.1cm, minimum height=0.8cm, fill=black!4] at (0.0,-1.3) {\textbf{p\_none}\\[-1pt]{\tiny pairLookup}};
\end{tikzpicture}
\end{tutscheme}
\tutwhat{pairLookup: the honesty ladder for four pairs (curated, regressed, UNIFAC estimate, none)}
\begin{tutnarrative}
WHAT IT IS. The student's FIRST question about a binary -- "what does Choupo have for this pair?" -- asked as ONE operation that answers in a fixed byte budget (2048 bytes of card text, held by check\_lookup\_budget) with an EVIDENCE CARD, never a number invented to fill a gap. The four operations walk the honesty ladder (DEV.md 4c C227, slice a):

p\_curated ethanol + water a catalogue pair record: the paste-ready source entry, validity, provenance p\_regressed acetone + ethanol Choupo's own regression (fittedInCase, interim) with its held-out verdict p\_estimate acetone + cyclo- no pair record, both carry UNIFAC hexane groups: SAID to be an estimate, with the measured warning (DEV.md C219) p\_none 1Pentadecanol + a component only in the group- water estimate lake: "I do not have this pair"

HOW IT READS. The operations read the INSTALLATION catalogue (a question about the catalogue), never the case's own constant/. The card carries what the records carry today; what the C16 chain will add (laboratories, archive sha256, extraction filter, fit AAD) is listed under announcedAbsent and never filled. The rung "servable from the measurement archive but not yet fitted" is slice (b) and the card says it does not exist yet. Pinned: the TIER WORD of each card and its byte count; the numbers a card quotes live in the records, which have their own gates. The thermoPhysPropDict is a placeholder the bench requires; no op reads it.

\end{tutnarrative}
\begin{tutgolden}
\goldrow{diag}{p\_curated}{budget\_bytes}{2048}
\goldrow{diag}{p\_curated}{card\_bytes}{947}
\goldrow{diag}{p\_curated}{n\_records}{2}
\goldrow{diag}{p\_regressed}{budget\_bytes}{2048}
\goldrow{diag}{p\_regressed}{card\_bytes}{957}
\goldrow{diag}{p\_regressed}{n\_records}{1}
\goldrow{diag}{p\_estimate}{budget\_bytes}{2048}
\goldrow{diag}{p\_estimate}{card\_bytes}{800}
\goldrow{diag}{p\_estimate}{n\_records}{0}
\goldrow{diag}{p\_none}{budget\_bytes}{2048}
\goldrow{diag}{p\_none}{card\_bytes}{779}
\goldrow{diag}{p\_none}{n\_records}{0}
\goldrow{verdict}{p\_curated}{tier}{curated}
\goldrow{verdict}{p\_regressed}{tier}{regressed}
\goldmore{3}
\end{tutgolden}

\clearpage
\tutentry{pairLookup03\_provenance\_stale\_seal}{tutorials/props/molecular/pairLookup03\_provenance\_stale\_seal}
\begin{tutscheme}
\begin{tikzpicture}[>={Stealth[length=2mm]}, every node/.style={font=\scriptsize}]
  \node[draw, rounded corners=2pt, align=center, minimum width=3.1cm, minimum height=0.8cm, fill=black!4] at (0.0,-0.0) {\textbf{p\_curated}\\[-1pt]{\tiny pairLookup}};
\end{tikzpicture}
\end{tutscheme}
\tutwhat{pairLookup provenance line: case sealed from an OLDER release than the installation}
\begin{tutnarrative}
WHAT IT IS. Witness of the PROVENANCE LINE every lookup card carries (DEV.md 4c C227 (a)) (the card carries a provenance line and the JSON a `catalogue` object).

THIS CASE: SEALED, and its manifest names Choupo-2607 -- a legitimate old seal (the catalogue identity a seal records is the catalogue as it stood when the case was frozen; CATALOGUE.dat explains why 394 manifests keep that name) -- while the installation is Choupo-2608. The lookups still read the installation, so the card is from a release that is NOT the case's own, and the provenance line SAYS so: `case sealed Choupo-2607; differs "..."`. The result JSON carries the same fact as `releaseDiffers true`. Pinned: the tier word, the card bytes.

\end{tutnarrative}
\begin{tutgolden}
\goldrow{diag}{p\_curated}{budget\_bytes}{2048}
\goldrow{diag}{p\_curated}{card\_bytes}{1039}
\goldrow{diag}{p\_curated}{n\_records}{2}
\goldrow{verdict}{p\_curated}{caseSeal}{sealed}
\goldrow{verdict}{p\_curated}{releaseDiffers}{true}
\end{tutgolden}

\clearpage
\tutentry{pcsaft01\_pure\_nhexane}{tutorials/props/molecular/pcsaft01\_pure\_nhexane}
\begin{tutscheme}
\begin{tikzpicture}[>={Stealth[length=2mm]}, every node/.style={font=\scriptsize}]
  \node[draw, rounded corners=2pt, align=center, minimum width=3.1cm, minimum height=0.8cm, fill=black!4] at (0.0,-0.0) {\textbf{liquid\_298K\_1bar}\\[-1pt]{\tiny propertyPoint}};
\end{tikzpicture}
\end{tutscheme}
\tutwhat{PC-SAFT pure n-hexane: Z + liquid molar volume at 298 K, 1 bar}
\begin{tutnarrative}
the NON-ASSOCIATING PC-SAFT core, validated on a pure n-alkane. n-hexane's segment parameters (m, sigma, epsilon/k) are the Gross \& Sadowski 2001 Table 1 fit; the case evaluates the compressibility factor Z and the liquid molar volume at 298.15 K / 1 bar and locks them as a golden. The engine PROVES the model self-consistent on every run (PCSAFT::verify: a\_res $\to$ 0 in the ideal limit, and the residual chemical potentials reproduce the bulk residual Gibbs energy to \textasciitilde{}1e-9). Physical cross-check (see README): the PC-SAFT liquid density here is \textasciitilde{}650 kg/m3, within \textasciitilde{}0.8 \% of the \textasciitilde{}655 kg/m3 experimental value.

\end{tutnarrative}
\begin{tutgolden}
\goldrow{diag}{liquid\_298K\_1bar}{Cp\_ig}{1.43106e+02}
\goldrow{diag}{liquid\_298K\_1bar}{H\_R}{-3.91460e+02}
\goldrow{diag}{liquid\_298K\_1bar}{H\_ig}{-1.67200e+05}
\goldrow{diag}{liquid\_298K\_1bar}{H\_real}{-1.67591e+05}
\goldrow{diag}{liquid\_298K\_1bar}{P}{1.00000e+05}
\goldrow{diag}{liquid\_298K\_1bar}{S\_R}{-0.846425525134}
\goldrow{diag}{liquid\_298K\_1bar}{S\_ig}{3.88510e+02}
\goldrow{diag}{liquid\_298K\_1bar}{S\_real}{387.663574475}
\goldrow{diag}{liquid\_298K\_1bar}{T}{2.98150e+02}
\goldrow{diag}{liquid\_298K\_1bar}{Z}{9.42530e-01}
\goldrow{diag}{liquid\_298K\_1bar}{Z\_liquid}{5.35107e-03}
\goldrow{diag}{liquid\_298K\_1bar}{gamma}{1.06168e+00}
\goldrow{diag}{liquid\_298K\_1bar}{v\_molar}{2.33649e-02}
\goldrow{diag}{liquid\_298K\_1bar}{v\_molar\_liq}{1.32651e-04}
\end{tutgolden}

\clearpage
\tutentry{pcsaft02\_states\_crosscheck}{tutorials/props/molecular/pcsaft02\_states\_crosscheck}
\begin{tutscheme}
\begin{tikzpicture}[>={Stealth[length=2mm]}, every node/.style={font=\scriptsize}]
  \node[draw, rounded corners=2pt, align=center, minimum width=3.1cm, minimum height=0.8cm, fill=black!4] at (0.0,-0.0) {\textbf{CH4\_vapour\_pcsaft}\\[-1pt]{\tiny propertyPoint}};
  \node[draw, rounded corners=2pt, align=center, minimum width=3.1cm, minimum height=0.8cm, fill=black!4] at (3.6,-0.0) {\textbf{CH4\_vapour\_srk}\\[-1pt]{\tiny propertyPoint}};
  \node[draw, rounded corners=2pt, align=center, minimum width=3.1cm, minimum height=0.8cm, fill=black!4] at (7.2,-0.0) {\textbf{propane\_liquid}\\[-1pt]{\tiny propertyPoint}};
  \node[draw, rounded corners=2pt, align=center, minimum width=3.1cm, minimum height=0.8cm, fill=black!4] at (0.0,-1.3) {\textbf{nButane\_liquid}\\[-1pt]{\tiny propertyPoint}};
\end{tikzpicture}
\end{tutscheme}
\tutwhat{PC-SAFT across states: CH4 Z vs SRK cross-check + propane/n-butane liquid density vs handbook}
\begin{tutnarrative}
PC-SAFT validated ACROSS STATES and AGAINST AN INDEPENDENT MODEL, before the model is trusted in a flowsheet.

Three probes, all goldened: (1) CH4 vapour Z at 300 K / 50 bar -- PC-SAFT vs SRK at the SAME state (the per-op `thermo \{\}` override): two independently-validated models must agree closely for a simple near-spherical vapour, and the run prints both so the student SEES the agreement (not told to assume it). (2) propane compressed liquid at 298.15 K / 15 bar -- rho\_liq vs the saturated-liquid handbook value (0.493 g/mL at 25 degC). (3) n-butane compressed liquid at 298.15 K / 5 bar -- rho\_liq vs 0.573 g/mL at 25 degC. Deviations sit at the 1-2 \% liquid-density AAD Gross \& Sadowski 2001 (Table 2) report for the n-alkane series -- these parameters PREDICT the density; it was never refitted here.

\end{tutnarrative}
\begin{tutgolden}
\goldrow{diag}{CH4\_vapour\_pcsaft}{Cp\_ig}{3.56607e+01}
\goldrow{diag}{CH4\_vapour\_pcsaft}{H\_R}{-8.28601e+02}
\goldrow{diag}{CH4\_vapour\_pcsaft}{H\_ig}{-7.44541e+04}
\goldrow{diag}{CH4\_vapour\_pcsaft}{H\_real}{-7.52827e+04}
\goldrow{diag}{CH4\_vapour\_pcsaft}{P}{5.00000e+06}
\goldrow{diag}{CH4\_vapour\_pcsaft}{S\_R}{-2.01270526714}
\goldrow{diag}{CH4\_vapour\_pcsaft}{S\_ig}{1.53964e+02}
\goldrow{diag}{CH4\_vapour\_pcsaft}{S\_real}{151.951191741}
\goldrow{diag}{CH4\_vapour\_pcsaft}{T}{3.00000e+02}
\goldrow{diag}{CH4\_vapour\_pcsaft}{Z}{9.12141e-01}
\goldrow{diag}{CH4\_vapour\_pcsaft}{gamma}{1.30404e+00}
\goldrow{diag}{CH4\_vapour\_pcsaft}{v\_molar}{4.55038e-04}
\goldrow{diag}{CH4\_vapour\_srk}{Cp\_ig}{3.56607e+01}
\goldrow{diag}{CH4\_vapour\_srk}{H\_R}{-8.31158e+02}
\goldmore{42}
\end{tutgolden}

\clearpage
\tutentry{pcsaft03\_association\_pure}{tutorials/props/molecular/pcsaft03\_association\_pure}
\begin{tutscheme}
\begin{tikzpicture}[>={Stealth[length=2mm]}, every node/.style={font=\scriptsize}]
  \node[draw, rounded corners=2pt, align=center, minimum width=3.1cm, minimum height=0.8cm, fill=black!4] at (0.0,-0.0) {\textbf{water\_liquid\_298K}\\[-1pt]{\tiny propertyPoint}};
  \node[draw, rounded corners=2pt, align=center, minimum width=3.1cm, minimum height=0.8cm, fill=black!4] at (3.6,-0.0) {\textbf{ethanol\_liquid\_298K}\\[-1pt]{\tiny propertyPoint}};
\end{tikzpicture}
\end{tutscheme}
\tutwhat{PC-SAFT association (Wertheim TPT1): 2B water + 2B ethanol, pure liquids}
\begin{tutnarrative}
PC-SAFT WITH ASSOCIATION (Wertheim TPT1; Gross \& Sadowski, Ind. Eng. Chem. Res. 41 (2002) 5510) -- the hydrogen-bonded liquids the non-associating core cannot represent, as PURE liquids at 298.15 K, 1 bar.

The association parameters live on the component records (2B water, 2B ethanol -- both G\&S 2002 Table 1, fitted WITH the association term; the paper lists water with TWO sites, and the scheme is PART of the fit: this same set mis-curated as 4C passed the pure-density check by coincidence and collapsed the ethanol/water mixture flash); the case only declares `model PCSAFT;`, and the engine announces the roster it found on the records:

[pcsaft] associating: water (2B), ethanol (2B); cross: Wolbach-Sandler

Liquid densities PREDICTED here (goldened), vs the CRC Handbook 25 degC values -- each deviation stated against what the SET is fitted for:

water v\_molar\_liq $\to$ rho\_liq = 921.9 kg/m3 vs 997.05 (-7.5 \% -- the PUBLISHED trade-off of the 2-site water fit, which is Psat-accurate instead: -0.2 \% at 358 K, +4 \% at 298 K) ethanol v\_molar\_liq $\to$ rho\_liq = 779.8 kg/m3 vs 789.3 (-1.2 \%)

What the student sees at verbosity 3 (the widened validation, ratified 2026-08-02): the construction-time verify() now ALSO echoes, per associating component as a pure fluid, the converged site fractions X\textasciicircum{}D / X\textasciicircum{}A, the a\_assoc share, and the CLOSED-FORM-vs-ITERATIVE oracle deviation (for nD = nA the mass-action closure collapses to a quadratic; the damped fixed point must land on its root -- observed machine zero). X\textasciicircum{}A(water) \textasciitilde{} 0.039: 96 \% of water's sites are hydrogen-bonded at ambient conditions -- the physics the term adds.

The two ops publish the standard vapour-root family (Z, v\_molar) plus the labelled liquid root (v\_molar\_liq); at 1 bar both fluids carry a metastable vapour root beside the stable liquid.

\end{tutnarrative}
\begin{tutgolden}
\goldrow{diag}{water\_liquid\_298K}{Cp\_ig}{3.36596e+01}
\goldrow{diag}{water\_liquid\_298K}{H\_R}{-2.18780e+03}
\goldrow{diag}{water\_liquid\_298K}{H\_ig}{-2.41826e+05}
\goldrow{diag}{water\_liquid\_298K}{H\_real}{-2.44014e+05}
\goldrow{diag}{water\_liquid\_298K}{P}{1.00000e+05}
\goldrow{diag}{water\_liquid\_298K}{S\_R}{-6.48369067192}
\goldrow{diag}{water\_liquid\_298K}{S\_ig}{1.88834e+02}
\goldrow{diag}{water\_liquid\_298K}{S\_real}{182.350309328}
\goldrow{diag}{water\_liquid\_298K}{T}{2.98150e+02}
\goldrow{diag}{water\_liquid\_298K}{Z}{9.00945e-01}
\goldrow{diag}{water\_liquid\_298K}{gamma}{1.32805e+00}
\goldrow{diag}{water\_liquid\_298K}{v\_molar}{2.23340e-02}
\goldrow{diag}{ethanol\_liquid\_298K}{Cp\_ig}{6.54261e+01}
\goldrow{diag}{ethanol\_liquid\_298K}{H\_R}{-4.01997e+03}
\goldmore{16}
\end{tutgolden}

\clearpage
\tutentry{ppr78Kij01\_ch4\_co2\_h2}{tutorials/props/molecular/ppr78Kij01\_ch4\_co2\_h2}
\begin{tutscheme}
\begin{tikzpicture}[>={Stealth[length=2mm]}, every node/.style={font=\scriptsize}]
  \node[draw, rounded corners=2pt, align=center, minimum width=3.1cm, minimum height=0.8cm, fill=black!4] at (0.0,-0.0) {\textbf{kij\_280K}\\[-1pt]{\tiny ppr78Kij}};
  \node[draw, rounded corners=2pt, align=center, minimum width=3.1cm, minimum height=0.8cm, fill=black!4] at (3.6,-0.0) {\textbf{kij\_500K}\\[-1pt]{\tiny ppr78Kij}};
\end{tikzpicture}
\end{tutscheme}
\tutwhat{The E-PPR78 group-contribution kij of a CH4/CO2/H2 gas, evaluated once at a declared temperature and written as a kij SET record a phi-phi package reads: at 280 K every pair lies inside the T window of the data its group parameters were fitted on, at 500 K every pair is an extrapolation and the set says so pair by pair}
\begin{tutnarrative}
A kij FROM GROUPS, AT A DECLARED TEMPERATURE. A cubic needs one binary correction kij per pair, and a measured one exists for few pairs. PPR78 (Jaubert \& Mutelet) predicts it from the GROUPS each molecule is made of:

k\_ij(T) = [ E\_ij(T) - (delta\_i - delta\_j)\textasciicircum{}2 ] / ( 2 delta\_i delta\_j ) delta\_i = sqrt(a\_i(T)) / b\_i (Choupo's own PR78) E\_ij(T) = -1/2 SUM\_k SUM\_l (alpha\_ik - alpha\_jk)(alpha\_il - alpha\_jl) A\_kl (298.15 / T)\textasciicircum{}(B\_kl / A\_kl - 1)

so kij DEPENDS ON T. The `ppr78Kij` op evaluates it ONCE, at the temperature it is told, and writes the matrix to a FILE -- a kij SET record (`recordType eosBinaryInteractionSet;`) that a phi-phi package then declares -- with the three terms of the quotient beside every value, so the division can be redone by hand.

WHAT THIS CASE SHOWS. The same three pairs at two temperatures: 280 K: every pair (CH4-CO2, CH4-H2, CO2-H2) lies inside the T range of the binary data its group parameters were fitted on (Xu et al. 2017, Table 2) -- nothing is marked; 500 K: every pair lies above that range, and the set writes `extrapolated yes;` with the window on each of them. The components are each their own group, so alpha is 1 or 0 and E\_ij is the single group-pair term.

THE GROUP TABLE is constant/parameters/PPR78/xu2017.dat, the E-PPR78 subset (Xu, Lasala, Privat \& Jaubert, Int. J. Greenh. Gas Control 56 (2017) 126-154, Table 3); its own header explains its transcription and that it is `reviewStatus interim`.

WHAT IS NOT VALIDATED. Nothing here is compared with a measurement: the numbers are Eq. 4 evaluated on the cited table and the engine's own PR78 a\_i(T), b\_i. Whether those kij reproduce a phase diagram is not asked here.

THE TWO SET FILES are run outputs: they carry the date they were written, so they are named `$\ast$.proposal.dat`, which git ignores.

\end{tutnarrative}
\begin{tutgolden}
\goldrow{diag}{kij\_280K}{T\_K}{280}
\goldrow{diag}{kij\_280K}{kij\_CH4\_CO2}{0.108020158785}
\goldrow{diag}{kij\_280K}{kij\_CH4\_H2}{0.29190260859}
\goldrow{diag}{kij\_280K}{kij\_CO2\_H2}{0.0990306279107}
\goldrow{diag}{kij\_280K}{nPairs}{3}
\goldrow{diag}{kij\_500K}{T\_K}{500}
\goldrow{diag}{kij\_500K}{kij\_CH4\_CO2}{0.147492367254}
\goldrow{diag}{kij\_500K}{kij\_CH4\_H2}{0.644015046404}
\goldrow{diag}{kij\_500K}{kij\_CO2\_H2}{0.424699449845}
\goldrow{diag}{kij\_500K}{nPairs}{3}
\end{tutgolden}

\clearpage
\tutentry{solubility01\_hildebrand\_ladder}{tutorials/props/molecular/solubility01\_hildebrand\_ladder}
\tutwhat{Hildebrand: one number for like-dissolves-like, and where it stops}
\tutreadme{Seven liquids on a ladder of cohesive energy density, from hexane to water.}
\tutreadme{$\delta$ is derived, never stored: sqrt(($\Delta$Hvap(T) - RT)/V\_m), with the latent heat from the Watson correlation the enthalpy legs already use. Every input is already in the component record, so a stored $\delta$ would be a second home for a fact the record fixes --- free to drift from it.}
\tutreadme{The published column is an $\ast$anchor$\ast$, not an input. It is what ChemSep computed independently, and its only job is to give the derivation something to be wrong against. Worst deviation here is 0.85 \%.}
\tutreadme{That matters because the three ways this arithmetic fails are all silent:}
\tutreadme{$\ast$ the RT term dropped --- $\delta$ a few per cent high, and a solubility parameter has no intuitive magnitude to make that look wrong; $\ast$ V\_m in the wrong unit; $\ast$ Pa\textasciicircum{}0.5 reported where the literature uses MPa\textasciicircum{}0.5 --- a factor of 31.6, and the number still reads as a solubility parameter.}
\tutreadme{Same device as the Zc column in a property data bank: a redundant determination turns a plausible wrong number into a visible disagreement.}
\begin{tutgolden}
\goldrow{diag}{solubilityParameter\_0}{delta\_acetone}{19673.2}
\goldrow{diag}{solubilityParameter\_0}{delta\_benzene}{18648.2}
\goldrow{diag}{solubilityParameter\_0}{delta\_cyclohexane}{16733.6}
\goldrow{diag}{solubilityParameter\_0}{delta\_ethanol}{26301.9}
\goldrow{diag}{solubilityParameter\_0}{delta\_nHexane}{14859}
\goldrow{diag}{solubilityParameter\_0}{delta\_toluene}{18183.8}
\goldrow{diag}{solubilityParameter\_0}{delta\_water}{48266.4}
\goldrow{diag}{solubilityParameter\_0}{dev\_acetone}{0.00287938}
\goldrow{diag}{solubilityParameter\_0}{dev\_benzene}{0.00277184}
\goldrow{diag}{solubilityParameter\_0}{dev\_cyclohexane}{0.000380181}
\goldrow{diag}{solubilityParameter\_0}{dev\_ethanol}{0.00619352}
\goldrow{diag}{solubilityParameter\_0}{dev\_nHexane}{0.000736807}
\goldrow{diag}{solubilityParameter\_0}{dev\_toluene}{0.003629}
\goldrow{diag}{solubilityParameter\_0}{dev\_water}{0.00849075}
\goldmore{6}
\end{tutgolden}

\clearpage
\tutentry{srk01\_h2n2\_density}{tutorials/props/molecular/srk01\_h2n2\_density}
\begin{tutscheme}
\begin{tikzpicture}[>={Stealth[length=2mm]}, every node/.style={font=\scriptsize}]
  \node[draw, rounded corners=2pt, align=center, minimum width=3.1cm, minimum height=0.8cm, fill=black!4] at (0.0,-0.0) {\textbf{p\_240K\_191bar\_xN2\_0p50008}\\[-1pt]{\tiny propertyPoint}};
  \node[draw, rounded corners=2pt, align=center, minimum width=3.1cm, minimum height=0.8cm, fill=black!4] at (3.6,-0.0) {\textbf{p\_250K\_50bar\_xN2\_0p50008}\\[-1pt]{\tiny propertyPoint}};
  \node[draw, rounded corners=2pt, align=center, minimum width=3.1cm, minimum height=0.8cm, fill=black!4] at (7.2,-0.0) {\textbf{p\_250K\_100bar\_xN2\_0p50008}\\[-1pt]{\tiny propertyPoint}};
  \node[draw, rounded corners=2pt, align=center, minimum width=3.1cm, minimum height=0.8cm, fill=black!4] at (0.0,-1.3) {\textbf{p\_250K\_150bar\_xN2\_0p50008}\\[-1pt]{\tiny propertyPoint}};
  \node[draw, rounded corners=2pt, align=center, minimum width=3.1cm, minimum height=0.8cm, fill=black!4] at (3.6,-1.3) {\textbf{p\_250K\_199bar\_xN2\_0p50008}\\[-1pt]{\tiny propertyPoint}};
  \node[draw, rounded corners=2pt, align=center, minimum width=3.1cm, minimum height=0.8cm, fill=black!4] at (7.2,-1.3) {\textbf{p\_350K\_199bar\_xN2\_0p50008}\\[-1pt]{\tiny propertyPoint}};
  \node[draw, rounded corners=2pt, align=center, minimum width=3.1cm, minimum height=0.8cm, fill=black!4] at (0.0,-2.6) {\textbf{p\_250K\_50bar\_xN2\_0p90001}\\[-1pt]{\tiny propertyPoint}};
  \node[draw, rounded corners=2pt, align=center, minimum width=3.1cm, minimum height=0.8cm, fill=black!4] at (3.6,-2.6) {\textbf{p\_250K\_100bar\_xN2\_0p90001}\\[-1pt]{\tiny propertyPoint}};
  \node[draw, rounded corners=2pt, align=center, minimum width=3.1cm, minimum height=0.8cm, fill=black!4] at (7.2,-2.6) {\textbf{p\_250K\_150bar\_xN2\_0p90001}\\[-1pt]{\tiny propertyPoint}};
  \node[draw, rounded corners=2pt, align=center, minimum width=3.1cm, minimum height=0.8cm, fill=black!4] at (0.0,-3.9) {\textbf{p\_250K\_197bar\_xN2\_0p90001}\\[-1pt]{\tiny propertyPoint}};
  \node[draw, rounded corners=2pt, align=center, minimum width=3.1cm, minimum height=0.8cm, fill=black!4] at (3.6,-3.9) {\textbf{p\_240K\_200bar\_xN2\_0p95000}\\[-1pt]{\tiny propertyPoint}};
  \node[draw, rounded corners=2pt, align=center, minimum width=3.1cm, minimum height=0.8cm, fill=black!4] at (7.2,-3.9) {\textbf{p\_250K\_50bar\_xN2\_0p95000}\\[-1pt]{\tiny propertyPoint}};
\end{tikzpicture}
\end{tutscheme}
\tutwhat{Does SRK describe the ammonia synthesis loop's gas? Sixteen MEASURED (T, P, x, rho) points on the nitrogen + hydrogen binary -- Hernandez-Gomez, Tuma, Gomez-Hernandez and Chamorro, J. Chem. Eng. Data (2017), doi:10.1021/acs.jced.7b00694 -- evaluated with the same SRK the green-ammonia plant runs on, and with kij = 0 because the catalogue carries no N2-H2 pair. The engine computes the molar volume; the case's own evidence file carries the measurement; the comparison is the lesson.}
\begin{tutnarrative}
IS SRK GOOD ENOUGH FOR THE AMMONIA SYNTHESIS LOOP?

Vitor's question, 2026-09-09: "dadas as pressoes extremas a que isto funciona, achas que os modelos termodinamicos que estamos a usar sao razoaveis e estao backed up por dados experimentais?" This case is the answer in numbers rather than in an opinion.

WHAT IT DOES. Sixteen `propertyPoint` operations, one per MEASURED state of the nitrogen + hydrogen binary, at exactly the (T, P, x) the experimenters reported. The engine computes the molar volume with the SAME declaration the green ammonia plant carries -- SRK, van der Waals one-fluid, kij = 0 -- and prints it. The measurement sits beside it in `constant/experiments/N2-H2-density-jced-7b00694.dat`, cited to its primary publication.

HOW TO READ THE ANSWER. Each operation reports `v\_molar` in m3/mol. The mass density the experimenters measured is

rho = MW\_mix / v\_molar , MW\_mix = x$\ast$MW(N2) + (1-x)$\ast$MW(H2)

and the two molar masses are in the components' own records (28.013 and 2.016 kg/kmol). Do that division for the four points at x(N2) = 0.50008 and 250 K and you have SRK's error climbing with pressure, which is the whole lesson: a cubic equation of state fitted to critical constants is a correlation, and 200 bar is far outside where its two parameters were determined.

`bin/curate/check\_srk\_h2n2\_aad.py` does exactly that arithmetic, from THIS run's own output and THIS case's own evidence file, and refuses if either moves. The number is not written down anywhere -- it is recomputed, because a transcribed agreement statistic is a second home for a fact the run owns.

WHAT THIS CASE DOES NOT SAY. It does not say SRK is wrong, and it does not say a better model exists in this tree -- none does for this binary today. It does not test the loop's COMPOSITION: see the caveat in the evidence file's own header. And it says nothing about ammonia, which is the strongly associating component the loop actually has to condense.

\end{tutnarrative}
\tutreadme{The green ammonia plant in tutorials/plant/greenAmmoniaPEQ\_baseCase runs its synthesis loop at around 200 bar and condenses its product near 250 K, and it prices every stream with a Soave--Redlich--Kwong equation of state whose only inputs are each component's critical temperature, critical pressure and acentric factor. Two questions follow, and only the second one has an answer a student can check:}
\tutreadme{1. Is that reasonable? --- an engineering judgement. 2. How far from measurement is it? --- a number.}
\tutreadme{Sixteen propertyPoint operations, one per measured state of the nitrogen + hydrogen binary, at exactly the temperature, pressure and composition the experimenters reported. The engine computes the molar volume with the plant's own declaration --- SRK, van der Waals one-fluid mixing, and kij = 0, which the run announces rather than assuming in silence:}
\tutreadme{The measurement sits beside it in constant/experiments/, cited to its primary publication:}
\tutreadme{R. Hern?ndez-G?mez, D. Tuma, A. G?mez-Hern?ndez, C. R. Chamorro, $\ast$Accurate Experimental (p, ?, T) Data for the Introduction of Hydrogen into the Natural Gas Grid: Thermodynamic Characterization of the Nitrogen--Hydrogen Binary System from 240 to 350 K and Pressures up to 20 MPa$\ast$, J. Chem. Eng. Data (2017). [doi:10.1021/acs.jced.7b00694](https://doi.org/10.1021/acs.jced.7b00694)}
\tutreadme{Each operation reports v\_molar in m$^3$/mol. Turn it into a mass density with}
\begin{tutgolden}
\goldrow{diag}{p\_240K\_191bar\_xN2\_0p50008}{Cp\_ig}{28.9758015655}
\goldrow{diag}{p\_240K\_191bar\_xN2\_0p50008}{H\_R}{-456.193455635}
\goldrow{diag}{p\_240K\_191bar\_xN2\_0p50008}{H\_ig}{-1683.90872502}
\goldrow{diag}{p\_240K\_191bar\_xN2\_0p50008}{H\_real}{-2140.10218065}
\goldrow{diag}{p\_240K\_191bar\_xN2\_0p50008}{P}{19066000}
\goldrow{diag}{p\_240K\_191bar\_xN2\_0p50008}{S\_R}{-2.3117177058}
\goldrow{diag}{p\_240K\_191bar\_xN2\_0p50008}{S\_ig}{116.976587739}
\goldrow{diag}{p\_240K\_191bar\_xN2\_0p50008}{S\_real}{114.664870033}
\goldrow{diag}{p\_240K\_191bar\_xN2\_0p50008}{T}{240.083}
\goldrow{diag}{p\_240K\_191bar\_xN2\_0p50008}{Z}{1.08683166645}
\goldrow{diag}{p\_240K\_191bar\_xN2\_0p50008}{gamma}{1.40241644741}
\goldrow{diag}{p\_240K\_191bar\_xN2\_0p50008}{v\_molar}{0.000113788478235}
\goldrow{diag}{p\_250K\_50bar\_xN2\_0p50008}{Cp\_ig}{28.9824821078}
\goldrow{diag}{p\_250K\_50bar\_xN2\_0p50008}{H\_R}{-156.040876541}
\goldmore{178}
\end{tutgolden}

\subsection{powder}
\grouptheory{powder}
\clearpage
\tutentry{compaction01\_mcc}{tutorials/props/powder/compaction01\_mcc}
\begin{tutscheme}
\begin{tikzpicture}[>={Stealth[length=2mm]}, every node/.style={font=\scriptsize}]
  \node[draw, rounded corners=2pt, align=center, minimum width=3.1cm, minimum height=0.8cm, fill=black!4] at (0.0,-0.0) {\textbf{mcc\_inDie}\\[-1pt]{\tiny compactionFit}};
  \node[draw, rounded corners=2pt, align=center, minimum width=3.1cm, minimum height=0.8cm, fill=black!4] at (3.6,-0.0) {\textbf{mcc\_inDie\_highWindow}\\[-1pt]{\tiny compactionFit}};
  \node[draw, rounded corners=2pt, align=center, minimum width=3.1cm, minimum height=0.8cm, fill=black!4] at (7.2,-0.0) {\textbf{mcc\_lubricated}\\[-1pt]{\tiny compactionFit}};
\end{tikzpicture}
\end{tutscheme}
\tutwhat{Microcrystalline cellulose on the compaction bench: Heckel, Kawakita and Ryshkewitch-Duckworth fitted to two open datasets, in die and out of die, with the declared Heckel window shown to move Py}
\begin{tutnarrative}
compaction01 -- MICROCRYSTALLINE CELLULOSE ON THE COMPACTION BENCH.

MCC is the main excipient of the apalutamide tablet case (tutorials/batch/plant/apalutamide01\_tablets: MCC 295.5 + silicified MCC 116.9 mg of a 700 mg core). This case fits the three compaction models to OPEN measurements of MCC, so the numbers a tablet-press design needs -- how hard the press must push, how much the bed shrinks, how strong the tablet comes out -- have a source a reader can open.

WHY A NEW FOLDER, `tutorials/props/powder/`. The tutorial folders of `props/` are named by the subject a student studies, and powder compaction is none of the existing ones (it is not thermodynamics, transport or a membrane): it is solids processing, the first props bench of its kind.

THREE OPERATIONS, TWO SOURCES (both CC BY 4.0; datasets in constant/experimental/, each with its own provenance header):

mcc\_inDie Avicel PH200, van der Haven et al., Sci. Data 13 (2026) 1468 -- 45 tablets compressed once on a compaction simulator, IN-DIE density at maximum force; 22 tablets crushed diametrally. Heckel on the declared low-pressure window, Kawakita on all pressures, Ryshkewitch-Duckworth on the crushing data, compared with the sigma\_0 and k\_b the dataset itself reports. mcc\_inDie\_highWindow the SAME 45 tablets, Heckel on a high-pressure window: the mean yield pressure roughly doubles. The window IS part of the answer, which is why the case declares it and the op never chooses it. mcc\_lubricated CEOLUS UF-711 + 2 \% sodium stearyl fumarate + 0.3 \% silica (Kukuls \& Mohylyuk 2025 dataset; protocol of Mohylyuk et al., Pharmaceutics 16 (2024) 362), the mean of ten tablets at each of six pressures, OUT OF DIE, compared with the Py = 106 MPa that paper reports.

WHAT TO READ IN THE OUTPUT. Py of MCC is low (tens to \textasciitilde{}100 MPa in die on the low-pressure line): it deforms PLASTICALLY, which is why it binds so well and why it is the filler of choice for direct compression. Compare it with compaction02 (dicalcium phosphate, brittle, several hundred MPa). The out-of-die Py is HIGHER than the in-die one on similar pressures because the ejected tablet has recovered the elastic part of its strain.

THE PAPER'S Py IS NOT REPRODUCED, AND THE REASON IS THE LESSON. Mohylyuk et al. 2024 state their linear region as 100-150 MPa. The deposited data carries ONE of its six pressures in that window (137 MPa), and a line through fewer than 4 points refuses here, so this case declares 90-340 MPa (4 points) and gets a different, larger, Py. (The two-point slope between the 91 and 137 MPa rows would give about 100 MPa -- close to the paper's -- but a two-point line has no residual to judge it by.) Both numbers are printed; nothing is tuned.

A FINDING, PRINTED AND NOT TUNED: on both materials of the van der Haven dataset (here and in the other powder case) the fitted k\_b agrees with the dataset's own constant within 10 \%, but the fitted sigma\_0 comes out about 1.5 times the dataset's. A fit in sigma\_t instead of ln(sigma\_t) does not close it (measured: 30.9 and 28.4 MPa). The files read do not say how the dataset's "Eq. 4" was fitted or on which porosity; its report PDF was not read. sigma\_0 is an extrapolation to zero porosity, so a small difference in slope or in the porosity basis moves it a lot.

HOW TO USE THIS FOR YOUR FORMULATION. 1. Find an OPEN dataset (CC BY: MDPI Pharmaceutics, Int. J. Pharm. X, Sci. Data, Zenodo, institutional dataverses) for the excipient that dominates your blend -- the same grade if you can. Check its licence on the article page; NonCommercial or "available on request" is not usable. 2. Write its rows in the experiments grammar (copy one of the files in constant/experimental/): P in MPa and ONE of D, porosity or the compact's density (then declare `trueDensity`, with its source). Say whether the density is in-die or out-of-die. 3. Plot the Heckel line first (run with a wide window and look at the `heckel` table's residuals), then DECLARE the linear region you defend, with a `reason`. Report Py WITH its window and its in-die or out-of-die basis; a Py without them cannot be compared. 4. If the source reports its own values, put them in `reported \{\}`. A difference is a finding (window, density basis, true density), never something to tune away. LIMITS, said plainly: these are data for MCC, not for your blend. NO mixture rule (for Py, sigma\_0 or k\_b) is implemented; the source's own ternary data shows mixtures do not follow a simple average. Lubricant (magnesium stearate) weakens MCC tablets and the lubricated dataset here is the closer stand-in for a real formulation. For another MCC grade the literature gives in-die Py = 52.2 +/- 4.2 MPa for Vivapur 102 at 11.28 mm (Lura et al., Int. J. Pharm. X 2 (2020) 100061, Table 2, CC BY) -- quoted for comparison, not fitted here.

\end{tutnarrative}
\begin{tutgolden}
\goldrow{diag}{mcc\_inDie}{dev\_ryshkewitch\_kb\_pct}{8.68066281422}
\goldrow{diag}{mcc\_inDie}{dev\_ryshkewitch\_sigma0\_MPa\_pct}{57.4687063406}
\goldrow{diag}{mcc\_inDie}{heckel\_A}{0.657612232409}
\goldrow{diag}{mcc\_inDie}{heckel\_A\_ci95}{0.063671311004}
\goldrow{diag}{mcc\_inDie}{heckel\_DA}{0.481913070072}
\goldrow{diag}{mcc\_inDie}{heckel\_K\_ci95\_perMPa}{0.000609029851532}
\goldrow{diag}{mcc\_inDie}{heckel\_K\_perMPa}{0.00906082550054}
\goldrow{diag}{mcc\_inDie}{heckel\_Py\_MPa}{110.365220028}
\goldrow{diag}{mcc\_inDie}{heckel\_Py\_ci95\_MPa}{7.41827701722}
\goldrow{diag}{mcc\_inDie}{heckel\_R2}{0.990980625546}
\goldrow{diag}{mcc\_inDie}{kawakita\_R2}{0.999983473773}
\goldrow{diag}{mcc\_inDie}{kawakita\_a}{0.780867436427}
\goldrow{diag}{mcc\_inDie}{kawakita\_b\_perMPa}{0.111550591336}
\goldrow{diag}{mcc\_inDie}{kawakita\_invb\_MPa}{8.96454234825}
\goldmore{39}
\end{tutgolden}

\clearpage
\tutentry{compaction02\_dcpd\_brittle}{tutorials/props/powder/compaction02\_dcpd\_brittle}
\begin{tutscheme}
\begin{tikzpicture}[>={Stealth[length=2mm]}, every node/.style={font=\scriptsize}]
  \node[draw, rounded corners=2pt, align=center, minimum width=3.1cm, minimum height=0.8cm, fill=black!4] at (0.0,-0.0) {\textbf{dcpd\_inDie}\\[-1pt]{\tiny compactionFit}};
  \node[draw, rounded corners=2pt, align=center, minimum width=3.1cm, minimum height=0.8cm, fill=black!4] at (3.6,-0.0) {\textbf{dcpa\_lubricated}\\[-1pt]{\tiny compactionFit}};
\end{tikzpicture}
\end{tutscheme}
\tutwhat{A brittle excipient on the compaction bench: dicalcium phosphate fitted to two open datasets (no open lactose dataset was found, and none is invented), against MCC's plastic behaviour in compaction01}
\begin{tutnarrative}
compaction02 -- A BRITTLE EXCIPIENT ON THE COMPACTION BENCH: DICALCIUM PHOSPHATE, AND WHY IT IS NOT LACTOSE.

The donepezil tablet case (tutorials/batch/plant/donepezil01\_tablets) is 62 \% lactose monohydrate, a FRAGMENTING excipient. This case was meant to fit the compaction models to open lactose data, and it does not, because NO open dataset with tabulated lactose pressure-density or strength-porosity rows was found on 2026-11-02 (CC BY sources searched: MDPI Pharmaceutics, Int. J. Pharm. X, PMC open access, Zenodo, Data in Brief; the lactose datasets located were NonCommercial, figure-only, or "available on request"). No lactose numbers were invented in its place.

What it fits instead is the open brittle reference that WAS found, from the same two sources as compaction01, so the plastic-against-brittle contrast can be read off two cases built the same way:

dcpd\_inDie DCPD Emcompress premium (dihydrate), van der Haven et al., Sci. Data 13 (2026) 1468, CC BY 4.0: 42 tablets, in-die density at maximum force; 19 crushed diametrally. Compared with the sigma\_0 and k\_b the dataset itself reports. dcpa\_lubricated CaHPO4 anhydrous DI-CAFOS A60 + 2 \% sodium stearyl fumarate + 0.3 \% silica (Kukuls \& Mohylyuk 2025, CC BY 4.0; protocol of Mohylyuk et al., Pharmaceutics 16 (2024) 362), out of die, compared with that paper's Py = 434 MPa.

WHAT TO READ. Py here is several hundred MPa -- several times MCC's (compaction01) -- the signature of a material that densifies by FRACTURE, not by flow. Its tablets are weak at a given porosity (low tensile strength) and k\_b is large: the strength collapses fast as the pores open. Lactose sits between the two: for Tablettose 80 the literature gives an in-die Py of 120.8 +/- 1.5 MPa at 11.28 mm, against 52.2 +/- 4.2 MPa for MCC Vivapur 102 in the same study (Lura et al., Int. J. Pharm. X 2 (2020) 100061, Table 2, CC BY) -- quoted for the donepezil group to compare with, NOT fitted here, because the study gives the curves only as figures.

A FINDING, PRINTED AND NOT TUNED: on both materials of the van der Haven dataset (here and in the other powder case) the fitted k\_b agrees with the dataset's own constant within 10 \%, but the fitted sigma\_0 comes out about 1.5 times the dataset's. A fit in sigma\_t instead of ln(sigma\_t) does not close it (measured: 30.9 and 28.4 MPa). The files read do not say how the dataset's "Eq. 4" was fitted or on which porosity; its report PDF was not read. sigma\_0 is an extrapolation to zero porosity, so a small difference in slope or in the porosity basis moves it a lot.

FOR THE DONEPEZIL GROUP. To fit your own lactose: find an open dataset (or digitise a CC BY figure, stating the reading error in the dataset header), write it in the grammar of the files in constant/experimental/, declare its Heckel window with a reason, and put the paper's own Py in `reported \{\}`. Say which lactose (monohydrate milled, sieved, agglomerated, spray-dried, anhydrous): Py differs between grades by a factor of two. A wet-granulated blend (your case's route) compacts differently again from the dry excipient, and NO mixture rule is implemented -- the excipient stands in for the blend, and you say so.

THE PAPER'S Py AND THIS CASE'S. Mohylyuk et al. 2024 analyse the F 0-100 formulation between 250 and 450 MPa. The deposited averages have three rows there (243 MPa falls just below it); a 3-point line refuses here, so the case declares 140-440 MPa (four rows) and both values are printed.

\end{tutnarrative}
\begin{tutgolden}
\goldrow{diag}{dcpd\_inDie}{dev\_ryshkewitch\_kb\_pct}{2.33048956558}
\goldrow{diag}{dcpd\_inDie}{dev\_ryshkewitch\_sigma0\_MPa\_pct}{49.4763706254}
\goldrow{diag}{dcpd\_inDie}{heckel\_A}{1.33186263307}
\goldrow{diag}{dcpd\_inDie}{heckel\_A\_ci95}{0.0229519605641}
\goldrow{diag}{dcpd\_inDie}{heckel\_DA}{0.736014904288}
\goldrow{diag}{dcpd\_inDie}{heckel\_K\_ci95\_perMPa}{7.74836836182e-05}
\goldrow{diag}{dcpd\_inDie}{heckel\_K\_perMPa}{0.00196611284247}
\goldrow{diag}{dcpd\_inDie}{heckel\_Py\_MPa}{508.617805854}
\goldrow{diag}{dcpd\_inDie}{heckel\_Py\_ci95\_MPa}{20.0444146948}
\goldrow{diag}{dcpd\_inDie}{heckel\_R2}{0.99330881447}
\goldrow{diag}{dcpd\_inDie}{kawakita\_R2}{0.999736957603}
\goldrow{diag}{dcpd\_inDie}{kawakita\_a}{0.63714261759}
\goldrow{diag}{dcpd\_inDie}{kawakita\_b\_perMPa}{0.0655574652353}
\goldrow{diag}{dcpd\_inDie}{kawakita\_invb\_MPa}{15.2537929344}
\goldmore{30}
\end{tutgolden}

\subsection{reactor}
\grouptheory{reactor}
\clearpage
\tutentry{thiele01\_sphere\_first\_order}{tutorials/props/reactor/thiele01\_sphere\_first\_order}
\begin{tutscheme}
\begin{tikzpicture}[>={Stealth[length=2mm]}, every node/.style={font=\scriptsize}]
  \node[draw, rounded corners=2pt, align=center, minimum width=3.1cm, minimum height=0.8cm, fill=black!4] at (0.0,-0.0) {\textbf{thiele\_CO\_alumina}\\[-1pt]{\tiny thielePellet}};
\end{tikzpicture}
\end{tutscheme}
\tutwhat{Single catalyst pellet: intraparticle concentration field, effectiveness factor, and the eta(phi) family, each verified against the first-order isothermal closed form}
\begin{tutnarrative}
One thielePellet operation.

THE RATE CONSTANT IS DECLARED HERE, and not on the catalyst record, because kinetics are equipment- and reaction-specific: the asset carries the SOLID, never the chemistry run on it. Two reactors loaded with the same catalyst would otherwise carry two copies of one rate constant and drift apart.

k = 40 1/s is a TEACHING VALUE chosen so the generalised modulus lands near 1 -- the transition region, where the field has a visible gradient and eta is neither \textasciitilde{}1 nor \textasciitilde{}1/phi. It is not measured and is not claimed to be.

\end{tutnarrative}
\tutreadme{Solves the intraparticle boundary-value problem for one porous catalyst pellet, publishes the concentration field, the effectiveness factor and the Thiele modulus, and checks every one of them against the closed form.}
\tutreadme{With xi the normalised coordinate from the pellet centre (0) to its external surface (1) and u = c/c\_s:}
\tutreadme{Both moduli are published because both conventions are in circulation and they differ by a factor of three for a sphere. The engine's own announcement in src/unitOperations/reactor/CatalystPellet.H uses the generalised one.}
\tutreadme{The equation is discretised with second-order central differences and handed to the existing solver::newtonND (finite-difference Jacobian, LU with partial pivoting) --- no new dependency, no new solver, every step readable in src/unitOperations/reactor/pellet/PelletDiffusion.cpp.}
\tutreadme{xi is one column, not two aliases: it is r/R for a sphere or a cylinder and x/L for a slab.}
\tutreadme{First-order isothermal is the rare case that has a closed form in all three geometries, so the numerical answer is not merely produced --- it is checked:}
\begin{tutgolden}
\goldrow{diag}{thiele\_CO\_alumina}{phi}{1.0161e+00}
\goldrow{diag}{thiele\_CO\_alumina}{phi\_char}{3.0483e+00}
\goldrow{diag}{thiele\_CO\_alumina}{Lambda}{5.0000e-04}
\goldrow{diag}{thiele\_CO\_alumina}{dimension}{1.5000e-03}
\goldrow{diag}{thiele\_CO\_alumina}{epsilon\_p}{4.5000e-01}
\goldrow{diag}{thiele\_CO\_alumina}{tau}{3.0000e+00}
\goldrow{diag}{thiele\_CO\_alumina}{rateConstant}{4.0000e+01}
\goldrow{diag}{thiele\_CO\_alumina}{T}{5.7315e+02}
\goldrow{diag}{thiele\_CO\_alumina}{P}{1.01325e+05}
\goldrow{diag}{thiele\_CO\_alumina}{D\_molecular}{6.45709e-05}
\goldrow{diag}{thiele\_CO\_alumina}{D\_eff}{9.68564e-06}
\goldrow{diag}{thiele\_CO\_alumina}{eta}{6.65742e-01}
\goldrow{diag}{thiele\_CO\_alumina}{eta\_closedForm}{6.65742e-01}
\goldrow{diag}{thiele\_CO\_alumina}{eta\_flux}{6.65731e-01}
\goldmore{17}
\end{tutgolden}

\clearpage
\tutentry{thiele02\_prater\_temperature}{tutorials/props/reactor/thiele02\_prater\_temperature}
\begin{tutscheme}
\begin{tikzpicture}[>={Stealth[length=2mm]}, every node/.style={font=\scriptsize}]
  \node[draw, rounded corners=2pt, align=center, minimum width=3.1cm, minimum height=0.8cm, fill=black!4] at (0.0,-0.0) {\textbf{thiele\_CO\_alumina\_prater}\\[-1pt]{\tiny thielePellet}};
\end{tikzpicture}
\end{tutscheme}
\tutwhat{Single catalyst pellet: intraparticle concentration field, effectiveness factor, and the eta(phi) family, each verified against the first-order isothermal closed form}
\begin{tutnarrative}
One thielePellet operation.

THE RATE CONSTANT IS DECLARED HERE, and not on the catalyst record, because kinetics are equipment- and reaction-specific: the asset carries the SOLID, never the chemistry run on it. Two reactors loaded with the same catalyst would otherwise carry two copies of one rate constant and drift apart.

k = 40 1/s is a TEACHING VALUE chosen so the generalised modulus lands near 1 -- the transition region, where the field has a visible gradient and eta is neither \textasciitilde{}1 nor \textasciitilde{}1/phi. It is not measured and is not claimed to be.

\end{tutnarrative}
\tutreadme{The SAME pellet as thiele01\_sphere\_first\_order, with one block added, so the difference between the two runs is the temperature.}
\tutreadme{Everything else --- catalyst, species, T, P, rate constant, diffusion, grid, verification tolerances --- is thiele01 unchanged. Run both and diff the output: nothing else moves.}
\tutreadme{The reaction is exothermic, so the pellet is hottest where the reactant is most depleted, which is the centre:}
\tutreadme{The field CSV gains two columns, T\_over\_Ts and T\_K, one value per node --- so the temperature is a FIELD, not a headline number, and the EduTool paints the same discs by it.}
\tutreadme{The species and energy balances inside a pellet carry the same reaction term. Eliminate it between them and, for constant D\_eff and k\_eff and any kinetics at all, what is left is a first integral --- Prater's relation:}
\tutreadme{So the temperature field costs no second boundary-value problem. It is algebra over the concentration field the operation already solved.}
\begin{tutgolden}
\goldrow{diag}{thiele\_CO\_alumina\_prater}{D\_eff}{9.68564e-06}
\goldrow{diag}{thiele\_CO\_alumina\_prater}{D\_molecular}{6.45709e-05}
\goldrow{diag}{thiele\_CO\_alumina\_prater}{Lambda}{0.0005}
\goldrow{diag}{thiele\_CO\_alumina\_prater}{P}{101325}
\goldrow{diag}{thiele\_CO\_alumina\_prater}{T}{573.15}
\goldrow{diag}{thiele\_CO\_alumina\_prater}{T\_centre\_K}{588.722}
\goldrow{diag}{thiele\_CO\_alumina\_prater}{c\_centre}{0.289874}
\goldrow{diag}{thiele\_CO\_alumina\_prater}{c\_centre\_closedForm}{0.289873}
\goldrow{diag}{thiele\_CO\_alumina\_prater}{deltaT\_centre\_K}{15.5718}
\goldrow{diag}{thiele\_CO\_alumina\_prater}{dimension}{0.0015}
\goldrow{diag}{thiele\_CO\_alumina\_prater}{epsilon\_p}{0.45}
\goldrow{diag}{thiele\_CO\_alumina\_prater}{eta}{0.665742}
\goldrow{diag}{thiele\_CO\_alumina\_prater}{eta\_closedForm}{0.665742}
\goldrow{diag}{thiele\_CO\_alumina\_prater}{eta\_flux}{0.665731}
\goldmore{23}
\end{tutgolden}

\subsection{reconciliation}
\grouptheory{reconciliation}
\clearpage
\tutentry{reconcile01\_recycle\_loop}{tutorials/props/reconciliation/reconcile01\_recycle\_loop}
\begin{tutscheme}
\begin{tikzpicture}[>={Stealth[length=2mm]}, every node/.style={font=\scriptsize}]
  \node[draw, rounded corners=2pt, align=center, minimum width=3.1cm, minimum height=0.8cm, fill=black!4] at (0.0,-0.0) {\textbf{honest}\\[-1pt]{\tiny dataReconciliation}};
  \node[draw, rounded corners=2pt, align=center, minimum width=3.1cm, minimum height=0.8cm, fill=black!4] at (3.6,-0.0) {\textbf{drifted}\\[-1pt]{\tiny dataReconciliation}};
  \node[draw, rounded corners=2pt, align=center, minimum width=3.1cm, minimum height=0.8cm, fill=black!4] at (7.2,-0.0) {\textbf{driftedWithout}\\[-1pt]{\tiny dataReconciliation}};
  \node[draw, rounded corners=2pt, align=center, minimum width=3.1cm, minimum height=0.8cm, fill=black!4] at (0.0,-1.3) {\textbf{yesterday}\\[-1pt]{\tiny dataReconciliation}};
\end{tikzpicture}
\end{tutscheme}
\tutwhat{Seven flow meters on an invented recycle loop, reconciled under four node balances: the residuals of the raw readings, the weighted least squares that is also the Bayesian posterior mode, the chi-square and measurement tests that catch a meter planted to read 15 t/h high, and a prior from yesterday -- the witness of the EduTool on data reconciliation}
\begin{tutnarrative}
reconcile01 -- SEVEN METERS THAT DO NOT CLOSE FOUR BALANCES (DEV.md 4c C121).

An invented recycle loop (constant/plant/network: a mixer, a reactor, a separator and a purge splitter, seven streams, one flow meter on each) read once (constant/plant/readings). The readings carry noise, so not one node balance closes; the residuals are of the size the meters' own sigmas predict. The witness of the EduTool "Data reconciliation, and how Bayes helps", which runs this case in the browser.

WHAT IS SYNTHETIC, said plainly. The plant, its true flows, the meter sigmas, the noise, the planted gross error and yesterday's prior are all invented. Nothing here is a measurement of anything. The point is the method, and an invented plant lets the page show the TRUTH beside the readings and the reconciled flows -- which no real plant can.

THE FOUR OPERATIONS, one per question the page asks.

honest the readings as they are. Weighted least squares under the four balances: the posterior mode for Gaussian meters and a flat prior. The objective is small against chi-square(4), and no meter's test exceeds the critical value.

drifted the same snapshot with F3 planted to read 15 t/h high (syntheticDrift). The global test rejects, the measurement test points at F3, and the reconciled flows show the error SMEARED over F3's neighbours.

driftedWithout the drifted snapshot with F3's meter taken out of service (exclude). The balances are projected onto the six remaining meters (the reactor and the separator merge into one envelope), the tests pass again, and F3 is read back from the balances -- with how far the excluded meter sits from it.

yesterday the honest readings with a prior per flow from yesterday's reconciled values (constant/plant/ yesterday): each reading and its prior combine by precision weighting, exactly, before the same reconciliation; the reconciled sigmas shrink again.

The significance level is DECLARED (alpha 0.05), never defaulted: the probability of a false alarm a plant accepts is a decision, not a constant.

\end{tutnarrative}
\begin{tutgolden}
\goldrow{csv}{prior\_streams.csv}{0:combined}{100.5}
\goldrow{csv}{prior\_streams.csv}{0:combinedSigma}{1.41421356237}
\goldrow{csv}{prior\_streams.csv}{1:combined}{147.7}
\goldrow{csv}{prior\_streams.csv}{1:combinedSigma}{2.12132034356}
\goldrow{csv}{prior\_streams.csv}{2:combined}{150.25}
\goldrow{csv}{prior\_streams.csv}{2:combinedSigma}{2.12132034356}
\goldrow{csv}{prior\_streams.csv}{3:combined}{89.4}
\goldrow{csv}{prior\_streams.csv}{3:combinedSigma}{1.27279220614}
\goldrow{csv}{prior\_streams.csv}{4:combined}{60.3195121951}
\goldrow{csv}{prior\_streams.csv}{4:combinedSigma}{0.937042571332}
\goldrow{csv}{prior\_streams.csv}{5:combined}{10.04}
\goldrow{csv}{prior\_streams.csv}{5:combinedSigma}{0.1788854382}
\goldrow{csv}{prior\_streams.csv}{6:combined}{50.0967213115}
\goldrow{csv}{prior\_streams.csv}{6:combinedSigma}{0.768221279597}
\goldmore{264}
\end{tutgolden}

\subsection{safety}
\grouptheory{safety}
\clearpage
\tutentry{relief01\_ammonia\_loop\_gas}{tutorials/props/safety/relief01\_ammonia\_loop\_gas}
\begin{tutscheme}
\begin{tikzpicture}[>={Stealth[length=2mm]}, every node/.style={font=\scriptsize}]
  \node[draw, rounded corners=2pt, align=center, minimum width=3.1cm, minimum height=0.8cm, fill=black!4] at (0.0,-0.0) {\textbf{loop150}\\[-1pt]{\tiny propertyPoint}};
  \node[draw, rounded corners=2pt, align=center, minimum width=3.1cm, minimum height=0.8cm, fill=black!4] at (3.6,-0.0) {\textbf{loop50}\\[-1pt]{\tiny propertyPoint}};
\end{tikzpicture}
\end{tutscheme}
\tutwhat{The ammonia synthesis loop's converter feed at the relieving pressure of a blocked outlet, at 150 bar and at 50 bar: the compressibility factor and Cp/Cv the gas relief-valve equation needs, computed with the SRK the loop runs on}
\begin{tutnarrative}
THE GAS SIDE OF A BLOCKED-OUTLET RELIEF VALVE ON THE AMMONIA SYNTHESIS LOOP (DEV.md 4c C61, the "Designing for overpressure" EduTool).

THE QUESTION. The converter's outlet is blocked while the recycle compressor goes on delivering the converter feed. The relief valve on the feed side must pass that whole flow at the relieving pressure, and the gas relief equation needs two properties of the gas there besides its molar mass: the compressibility factor Z and the ratio of specific heats k = Cp/Cv. This case computes both, for the same loop at 150 bar and at 50 bar -- the two green-ammonia feasibility cases the Process-safety EduTool compares (tutorials/plant/greenAmmonia02\_feasibility150bar and greenAmmonia02\_feasibility50bar).

WHERE THE STATES COME FROM. composition -- each plant's converged converter feed (`PreheatedFeed`), copied as the component flows in kmol/h exactly as the plant's converged/PreheatedFeed prints them. propertyPoint normalises the composition, so the flows ARE the composition; the EduTool reads them again as the relieving rate. gui/tests/overpressure.test.ts holds every number to the plant's own converged file. T -- the same stream's temperature (653.15 K and 593.15 K): a blocked outlet relieves the gas as it arrives. P -- the RELIEVING pressure, worked from the stream's operating pressure in the EduTool and written here to the digit: operating (gauge) = P\_stream - 1.01325 bar design (gauge) = max(1.1 x operating, operating + 1.7 bar) (the Design Guide's rule of thumb for a PRELIMINARY design pressure, docs/designGuide.tex, "Pressure relief, flare \& safety basics") MAWP (gauge) = design (gauge) -- an ASSUMPTION, the lowest MAWP the vessel can have; a real MAWP is a property of the built vessel, never of a rule of thumb set pressure = MAWP (a single device) relieving (gauge) = max(1.10 x MAWP, MAWP + 3 psi) relieving (abs) = relieving (gauge) + 1.01325 bar The design pressure is a rule of thumb applied to a feasibility case that sizes no vessel, and the MAWP is assumed equal to it; no number here is the design pressure or the MAWP of any real loop.

WHAT IS COMPUTED. propertyPoint's base set: Z on the SRK vapour root, and gamma = Cp\_ig/(Cp\_ig - R) from the records' ideal-gas heat capacities -- the ideal-gas ratio the gas relief equation is derived with.

\end{tutnarrative}
\tutreadme{The gas side of a blocked-outlet relief valve on the green-ammonia synthesis loop: the converter's outlet is shut while the recycle compressor goes on delivering the converter feed, so the valve on the feed side must pass that whole flow at the relieving pressure.}
\tutreadme{Two propertyPoint operations, one per loop pressure of the two C57 feasibility cases ([greenAmmonia02\_feasibility150bar](../../../plant/greenAmmonia02\_feasibility150bar) at 150 bar and [greenAmmonia02\_feasibility50bar](../../../plant/greenAmmonia02\_feasibility50bar)):}
\tutreadme{$\ast$ the composition is each plant's converged converter feed (PreheatedFeed), copied as component flows in kmol/h --- propertyPoint normalises it, and the EduTool reads the same numbers back as the relieving rate; $\ast$ the temperature is that stream's; $\ast$ the pressure is the relieving pressure of a single device set at an ASSUMED MAWP, taken equal to the preliminary design pressure the Design Guide's rule of thumb gives, max(1.1 $\times$ operating, operating + 1.7 bar), plus the 10 \% accumulation --- worked out in the propsDict header. A rule of thumb gives a preliminary design pressure, never a MAWP.}
\tutreadme{What the run gives the relief equation is Z (SRK vapour root, kij = 0, announced) and k = Cp/Cv (the ideal-gas ratio). The sizing itself is the EduTool's arithmetic: $\ast$Designing for overpressure$\ast$ (DEV.md 4c C61), where the 50 bar loop turns out to need a relief valve about four and a half times larger than the 150 bar one in the sizing arithmetic, on the same declared blocked-outlet flow hypothesis --- not a comparison of two real relief systems.}
\tutreadme{Nothing here is a design pressure or a relief load for a real loop.}
\begin{tutgolden}
\goldrow{diag}{loop150}{Cp\_ig}{29.9418368462}
\goldrow{diag}{loop150}{H\_R}{256.573139636}
\goldrow{diag}{loop150}{H\_ig}{9665.59492918}
\goldrow{diag}{loop150}{H\_real}{9922.16806882}
\goldrow{diag}{loop150}{P}{18216642.165}
\goldrow{diag}{loop150}{S\_R}{-0.125249214427}
\goldrow{diag}{loop150}{S\_ig}{129.797638756}
\goldrow{diag}{loop150}{S\_real}{129.672389542}
\goldrow{diag}{loop150}{T}{653.15}
\goldrow{diag}{loop150}{Z}{1.06279185274}
\goldrow{diag}{loop150}{gamma}{1.3844416123}
\goldrow{diag}{loop150}{v\_molar}{0.000316830516475}
\goldrow{diag}{loop50}{Cp\_ig}{30.1534149462}
\goldrow{diag}{loop50}{H\_R}{64.0995153664}
\goldmore{10}
\end{tutgolden}

\clearpage
\tutentry{relief02\_acetone\_fire\_case}{tutorials/props/safety/relief02\_acetone\_fire\_case}
\begin{tutscheme}
\begin{tikzpicture}[>={Stealth[length=2mm]}, every node/.style={font=\scriptsize}]
  \node[draw, rounded corners=2pt, align=center, minimum width=3.1cm, minimum height=0.8cm, fill=black!4] at (0.0,-0.0) {\textbf{relieving}\\[-1pt]{\tiny propertyPoint}};
\end{tikzpicture}
\end{tutscheme}
\tutwhat{An acetone drum in a pool fire: the boiling temperature at the relieving pressure, the latent heat there, and the vapour's Z and Cp/Cv -- the fluid properties the fire-case relief load and the gas sizing equation need, computed by the engine rather than typed in}
\begin{tutnarrative}
THE FLUID SIDE OF A FIRE-CASE RELIEF VALVE (DEV.md 4c C61, the "Designing for overpressure" EduTool).

THE QUESTION. A drum of acetone, a flammable solvent, sits in a pool fire. Its relief valve opens at the set pressure and the pressure climbs to the relieving pressure P1 = set + allowed accumulation; there the liquid BOILS, and the fire's heat leaves as vapour. The relief load is the heat absorbed divided by the latent heat, and the valve passes that vapour. So four properties of the fluid decide the valve, and all four are computed here:

T\_bubble the boiling temperature at P1 (the relieving temperature) Hvap\_acetone the latent heat of vaporisation at that temperature Z the vapour's compressibility factor at (T, P1) gamma Cp/Cv of the vapour (the ideal-gas ratio, Cp\_ig/(Cp\_ig-R))

THE NUMBERS ARE AUTHOR-SET, and the EduTool moves them. P is the relieving pressure of a drum whose MAWP is 5 barg, protected for the fire case with the set pressure at the MAWP: 1.21 x 5 barg + 1.01325 bar = 7.06325 bar absolute. The 5 barg is a round number chosen for the lesson, not a property of any real drum. T is the boiling temperature the SAME run reports at that pressure (T\_bubble), written back by hand so that Z, gamma and the latent heat are evaluated AT the relieving state; the EduTool does the same thing live, in two runs, whenever the reader moves the MAWP, and its test holds this file's T to its own T\_bubble.

WHERE EACH NUMBER COMES FROM. T\_bubble: the bubble point of pure acetone on the package in constant/thermoPhysPropDict (ideal liquid, SRK vapour). Hvap\_acetone: the record's HvapTb carried to T by the Watson correlation, Component::Hvap\_latent -- the same term the engine subtracts from the ideal-gas enthalpy to price every liquid stream. Z: the SRK vapour root. gamma: the record's ideal-gas heat capacity.

WHY ACETONE AND NOT TOLUENE. The first draft of this case was a toluene drum, and its run announced that toluene's Antoine record is fitted only from 273 to 380 K -- below the normal boiling point, so every boiling temperature of a pressurised toluene drum is an EXTRAPOLATED vapour pressure. A relief temperature resting on an extrapolation is a weak first lesson. Acetone's record covers 259.16 to 507.60 K, which holds the relieving state of any drum this EduTool lets the reader set.

WHAT THIS CASE DOES NOT SAY. Nothing here is a relief-valve design: the sizing is the EduTool's arithmetic on these numbers, and no number in either is a design for a real vessel.

\end{tutnarrative}
\tutreadme{The fluid side of a fire-case relief valve. A drum of acetone sits in a pool fire; its valve opens at the set pressure (here the MAWP, 5 barg --- a round number for the lesson) and the pressure climbs to the relieving pressure with the 21 \% fire accumulation, 7.06325 bar absolute. There the liquid boils, and the fire's heat leaves as vapour.}
\tutreadme{One propertyPoint operation, at the relieving pressure and at the boiling temperature the same run reports there, asks for the four numbers the relief load and the gas sizing equation need:}
\tutreadme{Hvap\_<comp> is a property key added for this case (DEV.md 4c C61): it is the same latent heat the engine subtracts from the ideal-gas enthalpy to price every liquid, and it refuses at or above the critical temperature.}
\tutreadme{Why acetone and not toluene: toluene's Antoine record is fitted only from 273 to 380 K, below its own normal boiling point, so every boiling temperature of a pressurised toluene drum would be an extrapolated vapour pressure. Acetone's covers 259.16 to 507.60 K.}
\tutreadme{The sizing is the EduTool's arithmetic: $\ast$Designing for overpressure$\ast$. Nothing here is a relief design for a real vessel.}
\begin{tutgolden}
\goldrow{diag}{relieving}{Cp\_ig}{91.7054083639}
\goldrow{diag}{relieving}{H\_R}{-1250.60447353}
\goldrow{diag}{relieving}{H\_ig}{-208718.53434}
\goldrow{diag}{relieving}{H\_real}{-209969.138813}
\goldrow{diag}{relieving}{Hvap\_acetone}{24076.6281528}
\goldrow{diag}{relieving}{P}{706325}
\goldrow{diag}{relieving}{S\_R}{-2.16115539516}
\goldrow{diag}{relieving}{S\_ig}{303.115727719}
\goldrow{diag}{relieving}{S\_real}{300.954572324}
\goldrow{diag}{relieving}{T}{399.459968689}
\goldrow{diag}{relieving}{T\_bubble}{399.459968689}
\goldrow{diag}{relieving}{Z}{0.876603897883}
\goldrow{diag}{relieving}{Z\_liquid}{0.0242807497275}
\goldrow{diag}{relieving}{gamma}{1.09970462073}
\goldmore{2}
\end{tutgolden}

\subsection{scan}
\grouptheory{scan}
\clearpage
\tutentry{binary01\_lle\_water\_nbutanol}{tutorials/props/scan/binary01\_lle\_water\_nbutanol}
\begin{tutscheme}
\begin{tikzpicture}[>={Stealth[length=2mm]}, every node/.style={font=\scriptsize}]
  \node[draw, rounded corners=2pt, align=center, minimum width=3.1cm, minimum height=0.8cm, fill=black!4] at (0.0,-0.0) {\textbf{lle}\\[-1pt]{\tiny propertyScanBinary}};
\end{tikzpicture}
\end{tutscheme}
\tutwhat{Binary LLE g\_mix + common-tangent (water/1-butanol) at 298.15 K via propertyScanBinary, cited Winkelman 2009 UNIQUAC}
\begin{tutnarrative}
Binary liquid-liquid teaching map -- water / 1-butanol at 25 C.

The molar Gibbs energy of mixing g\_mix(x1) swept across composition PLUS the two coexisting liquid compositions from ONE liquid-liquid IsothermalFlash (PhaseSet::LL). The two binodal points lie on the g\_mix curve and share its common tangent -- the student reads the split straight off the diagram.

gamma comes from the CITED Winkelman et al. (2009) UNIQUAC LLE set (Table 2; structural r,q from Table 1), so the g\_mix curve goes NON-CONVEX and the flash emits the real water/1-butanol binodal: \textasciitilde{}1.9 mol\% butanol in the water-rich phase and \textasciitilde{}49 mol\% butanol in the butanol-rich phase (matching their Figs 1-2 at 298 K). See README.md.

The flash feed is set to z(water) = 0.80 (20 mol\% butanol), INSIDE the asymmetric gap; an equimolar feed would lie outside it (both binodal branches sit below 50 mol\% butanol). Run-only smoke (the CSV is a teaching curve, not pinned KPIs).

\end{tutnarrative}
\tutreadme{Water + 1-butanol (n-butanol) is $\ast$the$\ast$ classic partially-miscible binary. At 25 $^\circ$C it splits into a water-rich liquid (\textasciitilde{}2 mol \% butanol) and a butanol-rich liquid (\textasciitilde{}50 mol \% butanol). A g\_mix(x\_water) scan should therefore show a clearly non-convex region: two binodal points sharing a common tangent, which a student reads the split straight off.}
\tutreadme{and binaryLle.csv is a g\_mix(x1) curve that goes non-convex, plus the two binodal points the liquid-liquid flash locates on it:}
\tutreadme{The agreement is essentially exact: the cited UNIQUAC fit reproduces the published mutual solubilities.}
\tutreadme{- J.G.M. Winkelman, G.N. Kraai, H.J. Heeres, $\ast$"Binary, ternary and quaternary liquid-liquid equilibria in 1-butanol, oleic acid, water and n-heptane mixtures"$\ast$, Fluid Phase Equilibria 284 (2009) 71--79, doi:10.1016/j.fluid.2009.06.013.}
\tutreadme{- Interaction parameters --- their Table 2 (p.73), binary 1-butanol(1)/water(3), valid 273--363 K, with the quadratic temperature correlation A\_ij(T) = a + b$\cdot$T + c$\cdot$T$^2$ (Eq. 10), $\tau$\_ij = exp(-A\_ij/T) (Eq. 7): - A\_1,3 (butanol$\to$water): a = 155.31, b = 1.0822, 10?$\cdot$c = -43.711 - A\_3,1 (water$\to$butanol): a = -579.36, b = 2.7517, 10?$\cdot$c = -6.7700 - Structural r, q --- their Table 1 (p.72): 1-butanol r = 3.9243, q = 3.668; water r = 0.9200, q = 1.400.}
\tutreadme{The numbers are inline here, fully provenanced above; the record is STAGED CASE-LOCALLY, not curated: it lives at tutorials/steady/flash/marcilla02\_lls\_ternary\_invariant/constant/parameters/UNIQUAC/nButanol-water.dat with reviewStatus interim, and promotion into data/standards/ is the maintainer's batch review. This README used to claim the numbers already lived under data/standards/ -- asserting a curation status that had been deliberately withheld, which is a worse error than saying nothing.}
\begin{tutgolden}
\goldrow{diag}{lle}{T}{298.15}
\goldrow{diag}{lle}{split}{1}
\end{tutgolden}

\clearpage
\tutentry{hxy01\_ethanol\_water\_1atm}{tutorials/props/scan/hxy01\_ethanol\_water\_1atm}
\begin{tutscheme}
\begin{tikzpicture}[>={Stealth[length=2mm]}, every node/.style={font=\scriptsize}]
  \node[draw, rounded corners=2pt, align=center, minimum width=3.1cm, minimum height=0.8cm, fill=black!4] at (0.0,-0.0) {\textbf{hxy}\\[-1pt]{\tiny enthalpyConcentration}};
\end{tikzpicture}
\end{tutscheme}
\tutwhat{Enthalpy-concentration locus (h-x-y) of ethanol/water at 1 atm on the elements datum -- the data a Ponchon-Savarit construction is drawn on}
\begin{tutnarrative}
The enthalpy-concentration locus of ethanol/water at 1 atm.

McCabe-Thiele carries equilibrium and material balance under CONSTANT MOLAR OVERFLOW -- one mole condensing releases exactly what one mole vaporising absorbs, whatever the composition. This case publishes the quantity that assumption is about: the saturated-liquid enthalpy h(x) on its bubble line, the saturated-vapour enthalpy H(y) on its dew line, and the equilibrium tie between them, all on Choupo's ONE enthalpy datum (the elements at 298.15 K).

lambda(x) = H(y$\ast$) - h(x) is the molar heat of vaporisation ALONG the locus. CMO says it is a constant. Read the two lambda diagnostics and decide for yourself.

Run-only + KPI golden; see README.md for the provenance of every number and for what is NOT validated.

\end{tutnarrative}
\tutreadme{The saturated-liquid enthalpy on its bubble line, the saturated-vapour enthalpy on its dew line, and the equilibrium tie between them --- the data an enthalpy-concentration (Ponchon--Savarit) construction is drawn on, and the first thing the engine publishes that carries enthalpy against composition.}
\tutreadme{enthalpyConcentration.csv --- a regular table, one header, homogeneous rows:}
\tutreadme{Every row is one equilibrium. It carries $\ast$both$\ast$ ends of a tie line --- the liquid (x1, h\_liquid) and the vapour (y1, H\_vapour) that are in equilibrium with each other --- at the one temperature T\_K they share. So nothing downstream pairs rows by index and nothing interpolates one end of a tie. role says only which composition was the uniform grid node:}
\tutreadme{The saturated-liquid curve is (x1, h\_liquid) over all rows; the saturated-vapour curve is (y1, H\_vapour) over all rows. status is ok or names a refusal; a refused node keeps its grid coordinate and leaves the solved fields empty, so a gap is visible rather than silently bridged.}
\tutreadme{enthalpyConcentration.meta --- a key,value sidecar, written on every run. It states the two components, the pressure, which enthalpy datum the two columns are on, which kernel computed each, the tie contract, whether the heat of mixing is included, and the reason for every refused node. The construction's lever arms are enthalpy $\ast$differences$\ast$, so a constant offset cancels in the geometry --- which is exactly why the file has to say what its axis means instead of letting the next reader assume.}
\tutreadme{Both enthalpy columns are absolute on Choupo's one enthalpy base, the elements at 298.15 K, from the same two kernels every energy balance and duty in the engine is priced on:}
\begin{tutgolden}
\goldrow{diag}{hxy}{H\_vapour\_pure1\_J\_per\_mol}{-231233.18018}
\goldrow{diag}{hxy}{H\_vapour\_pure2\_J\_per\_mol}{-239303.954931}
\goldrow{diag}{hxy}{P\_Pa}{101325}
\goldrow{diag}{hxy}{T\_bubble\_pure1\_K}{351.452745505}
\goldrow{diag}{hxy}{T\_bubble\_pure2\_K}{372.453622433}
\goldrow{diag}{hxy}{h\_liquid\_pure1\_J\_per\_mol}{-271912.766061}
\goldrow{diag}{hxy}{h\_liquid\_pure2\_J\_per\_mol}{-280791.66824}
\goldrow{diag}{hxy}{lambda\_max\_J\_per\_mol}{45052.6104262}
\goldrow{diag}{hxy}{lambda\_min\_J\_per\_mol}{40679.5858814}
\goldrow{diag}{hxy}{n\_nodes}{82}
\goldrow{diag}{hxy}{n\_refused}{0}
\end{tutgolden}

\clearpage
\tutentry{phase01\_water\_full}{tutorials/props/scan/phase01\_water\_full}
\begin{tutscheme}
\begin{tikzpicture}[>={Stealth[length=2mm]}, every node/.style={font=\scriptsize}]
  \node[draw, rounded corners=2pt, align=center, minimum width=3.1cm, minimum height=0.8cm, fill=black!4] at (0.0,-0.0) {\textbf{pd}\\[-1pt]{\tiny purePhaseDiagram}};
\end{tikzpicture}
\end{tutscheme}
\tutwhat{Complete P-T phase diagram of water (purePhaseDiagram, solid+liquid+vapour)}
\begin{tutnarrative}
Complete pure-compound P-T phase diagram of water: liquid-vapour saturation (to the critical point) + solid-vapour sublimation + solid-liquid fusion (the negative-slope ice line) + triple and critical points. Clapeyron physics in the purePhaseDiagram engine op; cited solid data in the solid\{\} block (IAPWS/CRC, see solidPhaseData.ts). Run-only smoke.

\end{tutnarrative}
\begin{tutgolden}
\goldrow{diag}{pd}{Pc\_bar}{220.64}
\goldrow{diag}{pd}{Tc}{647.14}
\goldrow{diag}{pd}{has\_solid}{1}
\end{tutgolden}

\clearpage
\tutentry{scan2d01\_co2\_compressibility}{tutorials/props/scan/scan2d01\_co2\_compressibility}
\begin{tutscheme}
\begin{tikzpicture}[>={Stealth[length=2mm]}, every node/.style={font=\scriptsize}]
  \node[draw, rounded corners=2pt, align=center, minimum width=3.1cm, minimum height=0.8cm, fill=black!4] at (0.0,-0.0) {\textbf{Z\_co2\_TP}\\[-1pt]{\tiny propertyScan2D}};
\end{tikzpicture}
\end{tutscheme}
\tutwhat{Z(T,P) for CO2 on Peng-Robinson over a 15 x 20 grid that crosses the critical point: the compressibility surface, one row per node}
\begin{tutnarrative}
THE COMPRESSIBILITY SURFACE, NOT A CURVE.

`propertyScan2D` sweeps TWO state variables over a rectangle and reports the property at every node -- long-form CSV, one row per (T, P), ready to pivot into a heat map. Here it is Z(T, P) for CO2 on Peng-Robinson, 280-420 K by 1-200 bar: a 15 x 20 grid that walks straight over the critical point (304.13 K, 73.77 bar).

What to read: along an isotherm below Tc the surface DROPS steeply and then turns up -- the liquid-like branch; above Tc the same isotherm is smooth: the discontinuity is a coexistence line, and the grid crosses where it ends; Z $\to$ 1 in the whole low-P edge, which is the ideal-gas limit falling out of a cubic rather than being assumed.

The op reports `failures` -- nodes where the property could not be evaluated -- as a diagnostic, so a grid that quietly returned fewer points than it promised is visible in the payload rather than only in the CSV's line count. The expected file pins it at zero: on this rectangle every node must solve.

Why the case exists: `propertyScan2D` had NO case and no gate, though `propertyScan1D`, `propertyScanBinary` and `propertyScanTernary` all have one. Its sibling vaporPressureFit was in the same position, and writing the first case for THAT one segfaulted choupoProps twice.

\end{tutnarrative}
\begin{tutgolden}
\goldrow{diag}{Z\_co2\_TP}{failures}{0}
\goldrow{diag}{Z\_co2\_TP}{n\_points}{300}
\goldrow{diag}{Z\_co2\_TP}{n\_x}{15}
\goldrow{diag}{Z\_co2\_TP}{n\_y}{20}
\end{tutgolden}

\clearpage
\tutentry{ternary01\_audit\_vlle}{tutorials/props/scan/ternary01\_audit\_vlle}
\begin{tutscheme}
\begin{tikzpicture}[>={Stealth[length=2mm]}, every node/.style={font=\scriptsize}]
  \node[draw, rounded corners=2pt, align=center, minimum width=3.1cm, minimum height=0.8cm, fill=black!4] at (0.0,-0.0) {\textbf{lle}\\[-1pt]{\tiny propertyScanTernary}};
\end{tikzpicture}
\end{tutscheme}
\tutwhat{Ternary VLLE phase map (compA/compB/compC) at 350 K, 1 bar via propertyScanTernary}
\begin{tutnarrative}
ILLUSTRATIVE / SYNTHETIC DATA. compA / compB / compC are NOT real species: they are artificial components (case-local constant/components/ + inline NRTL pairs in thermoPhysPropDict) deliberately tuned to guarantee a ternary VLLE region, so this case exercises the propertyScanTernary classifier on a known-3-phase map. Architecture (see thermoPhysPropDict): A and B mutually immiscible (LL gap on the A-B axis), C a volatile bridge soluble in both liquids. At 350 K, 1 bar the equimolar feed gives vapour rich in C, L\_alpha rich in A, L\_beta rich in B.

Expected map (14 intervals/edge, 78 nodes): VLLE 36, LL 29, VL 0, 1-phase 13, 22 tie-lines, 0 failures. The point to SEE: a healthy spread of 3-phase (VLLE) and LL nodes with zero solver failures --- the classifier resolves the engineered gap node-by-node.

\end{tutnarrative}
\begin{tutgolden}
\goldrow{diag}{lle}{failures}{0}
\goldrow{diag}{lle}{n\_ll}{32}
\goldrow{diag}{lle}{n\_nodes}{78}
\goldrow{diag}{lle}{n\_onephase}{10}
\goldrow{diag}{lle}{n\_tielines}{23}
\goldrow{diag}{lle}{n\_vl}{0}
\goldrow{diag}{lle}{n\_vlle}{36}
\end{tutgolden}

\clearpage
\tutentry{ternary02\_bubbleT\_benzene\_toluene\_water}{tutorials/props/scan/ternary02\_bubbleT\_benzene\_toluene\_water}
\begin{tutscheme}
\begin{tikzpicture}[>={Stealth[length=2mm]}, every node/.style={font=\scriptsize}]
  \node[draw, rounded corners=2pt, align=center, minimum width=3.1cm, minimum height=0.8cm, fill=black!4] at (0.0,-0.0) {\textbf{bt}\\[-1pt]{\tiny propertyScanTernary}};
\end{tikzpicture}
\end{tutscheme}
\tutwhat{Ternary T\_bubble surface (benzene/toluene/water) at 1 bar via propertyScanTernary mode bubbleT}
\begin{tutnarrative}
Ternary boiling-temperature surface (T\_bubble over the composition triangle) via propertyScanTernary mode bubbleT --- reuses BubblePoint per simplex node. Works with ideal binaries (only the 3 Psat are needed). Run-only smoke.

\end{tutnarrative}
\begin{tutgolden}
\goldrow{diag}{bt}{failures}{0}
\goldrow{diag}{bt}{n\_nodes}{150}
\end{tutgolden}

\clearpage
\tutentry{ternary03\_lle\_water\_ethanol\_benzene}{tutorials/props/scan/ternary03\_lle\_water\_ethanol\_benzene}
\begin{tutscheme}
\begin{tikzpicture}[>={Stealth[length=2mm]}, every node/.style={font=\scriptsize}]
  \node[draw, rounded corners=2pt, align=center, minimum width=3.1cm, minimum height=0.8cm, fill=black!4] at (0.0,-0.0) {\textbf{lle}\\[-1pt]{\tiny propertyScanTernary}};
\end{tikzpicture}
\end{tutscheme}
\tutwhat{Ternary LLE solubility (water/ethanol/benzene) at 298 K via UNIFAC + propertyScanTernary mode lle}
\begin{tutnarrative}
Ternary LIQUID-LIQUID solubility map (water / ethanol / benzene) --- the classic extraction triangle. Activity from UNIFAC (group contribution): NO fitted binary pairs needed --- gamma is predicted from molecular groups, so a real system shows a real miscibility gap with no curated pair data. mode lle uses the liquid-liquid flash (PhaseSet::LL): clean binodal + tie-lines, no spurious vapour. Run-only smoke (a phase map, not KPIs).

\end{tutnarrative}
\begin{tutgolden}
\goldrow{diag}{lle}{failures}{0}
\goldrow{diag}{lle}{n\_ll}{84}
\goldrow{diag}{lle}{n\_nodes}{105}
\goldrow{diag}{lle}{n\_onephase}{21}
\goldrow{diag}{lle}{n\_tielines}{21}
\goldrow{diag}{lle}{n\_vl}{0}
\goldrow{diag}{lle}{n\_vlle}{0}
\end{tutgolden}

\subsection{solvents}
\grouptheory{solvents}
\clearpage
\tutentry{solvent01\_synthetic\_mixture}{tutorials/props/solvents/solvent01\_synthetic\_mixture}
\begin{tutscheme}
\begin{tikzpicture}[>={Stealth[length=2mm]}, every node/.style={font=\scriptsize}]
  \node[draw, rounded corners=2pt, align=center, minimum width=3.1cm, minimum height=0.8cm, fill=black!4] at (0.0,-0.0) {\textbf{blend}\\[-1pt]{\tiny hansenMixture}};
\end{tikzpicture}
\end{tutscheme}
\tutwhat{DIDACTIC EXAMPLE WITH SYNTHETIC PARAMETERS: two fictitious solvents A and B, a fictitious polymer P between them -- the Hansen distance along the A/B line has its minimum at 50:50 by volume, hand-checkable}
\begin{tutnarrative}
solvent01 -- DIDACTIC EXAMPLE WITH SYNTHETIC PARAMETERS (DEV.md 4c C128, adjustment 2: start with a synthetic case, then a real one, then the confrontation with measurements).

THE ARITHMETIC, by hand (MPa\textasciicircum{}0.5):

solventA dD 16 dP 2 dH 4 solventB dD 16 dP 10 dH 12 polymerP dD 17 dP 6 dH 8 R0 4

Along the A/B line at volume fraction s of B the mixture rule gives (16, 2 + 8 s, 4 + 8 s), so

Ra\textasciicircum{}2 = 4 (16 - 17)\textasciicircum{}2 + (8 s - 4)\textasciicircum{}2 + (8 s - 4)\textasciicircum{}2 = 4 + 2 (8 s - 4)\textasciicircum{}2

which is 36 at either pure end (Ra = 6.000, RED = 1.5) and 4 at s = 0.5 (Ra = 2.000, RED = 0.5): the minimum sits at 50:50 BY CONSTRUCTION, and the dispersion offset of 1 keeps it from being a degenerate zero. The two solvents have equal Vliq, so the volume, mole and mass bases coincide here -- the basis conversion is exercised on the real case instead.

THE PHASE WORD IS `notVerified`, honestly: the package declares an ideal liquid (gamma = 1), which has no excess Gibbs energy and can never demix, so a stability test on it would say "homogeneous" about every pair ever asked. The op marks every blend point rather than claim a verification the model cannot give.

\end{tutnarrative}
\tutreadme{Two fictitious solvents A and B and a fictitious polymer P between them, so the Hansen distance along the A/B line has its minimum at 50:50 by volume --- by construction, and checkable by hand. Nothing here is a real substance; the engine says so on every run ([synthetic] component 'solventA' is NOT a real substance), and the records carry provenance \{ source synthetic; \}.}
\tutreadme{All values in MPa\textasciicircum{}0.5 (the records store Pa\textasciicircum{}0.5, 1000 times larger):}
\tutreadme{Hansen's distance, with the factor 4 on the dispersion term (his 1967 thesis, Eq. 8.2), along the line at volume fraction s of B:}
\tutreadme{so Ra = 6.000 at either pure end (RED = 1.5, outside the sphere) and Ra = 2.000 at s = 0.5 (RED = 0.5, inside). The dispersion offset of 1 keeps the minimum from being a degenerate zero. The two solvents have the same molar volume, so the volume, mole and mass bases coincide here; the conversion is exercised on the real case (solvent02).}
\tutreadme{The package declares an ideal liquid, and gamma = 1 has no excess Gibbs energy: it can never demix, so a stability test on it would say "one liquid" about every pair ever asked. The op marks every blend point notVerified rather than claim a verification the model cannot give --- the honest third outcome of the phase check (the other two, homogeneous and twoPhase, are live on solvent02 under UNIFAC). The pure ends carry pure: one component, nothing to split.}
\tutreadme{$\ast$ the mixture rule is the volume-fraction average of each parameter --- an approximation (no excess volume; each solvent keeps its own split); $\ast$ the parameters are held constant at their reference temperature (Tref 298.15 K on every record); no correction exists in the engine, and the declared T decides the phase check alone; $\ast$ Ra is an affinity, not a solubility: no saturation concentration is published.}
\begin{tutgolden}
\goldrow{diag}{blend}{R0\_MPa05}{4}
\goldrow{diag}{blend}{RED\_min}{0.5}
\goldrow{diag}{blend}{Ra\_min\_MPa05}{2}
\goldrow{diag}{blend}{Ra\_pure\_solventA}{6}
\goldrow{diag}{blend}{Ra\_pure\_solventB}{6}
\goldrow{diag}{blend}{T\_K}{298.15}
\goldrow{diag}{blend}{Tref\_target\_K}{298.15}
\goldrow{diag}{blend}{n\_homogeneous}{0}
\goldrow{diag}{blend}{n\_notVerified}{19}
\goldrow{diag}{blend}{n\_points}{21}
\goldrow{diag}{blend}{n\_pure}{2}
\goldrow{diag}{blend}{n\_twoPhase}{0}
\goldrow{diag}{blend}{phiB\_at\_min}{0.5}
\goldrow{diag}{blend}{s\_at\_min}{0.5}
\goldmore{35}
\end{tutgolden}

\clearpage
\tutentry{solvent02\_polystyrene\_hexane\_blends}{tutorials/props/solvents/solvent02\_polystyrene\_hexane\_blends}
\begin{tutscheme}
\begin{tikzpicture}[>={Stealth[length=2mm]}, every node/.style={font=\scriptsize}]
  \node[draw, rounded corners=2pt, align=center, minimum width=3.1cm, minimum height=0.8cm, fill=black!4] at (0.0,-0.0) {\textbf{hexaneAcetone}\\[-1pt]{\tiny hansenMixture}};
  \node[draw, rounded corners=2pt, align=center, minimum width=3.1cm, minimum height=0.8cm, fill=black!4] at (3.6,-0.0) {\textbf{hexaneNitromethane}\\[-1pt]{\tiny hansenMixture}};
\end{tikzpicture}
\end{tutscheme}
\tutwhat{Polystyrene (Hansen 1967, solute G) against two hexane blends: hexane + acetone, two solvents each short of the polymer that meet it between them, and hexane + nitromethane, whose blend the mixture rule likes exactly where UNIFAC finds two liquids}
\begin{tutnarrative}
solvent02 -- THE REAL CASE (DEV.md 4c C128, adjustment 2, step two).

The target is polystyrene as Hansen placed it in his 1967 thesis (solute G, Table 4.1: centre (8.6, 3.0, 2.0) (cal/cm3)\textasciicircum{}0.5, radius 3.5); the solvents are three of his Table 2.1 placements, read from the page images and converted to Pa\textasciicircum{}0.5 in constant/components/$\ast$.dat, one citation per value. Two blends, two lessons:

1. hexane + acetone, by VOLUME. In the engine's MPa\textasciicircum{}0.5 (the thesis's (cal/cm3)\textasciicircum{}0.5 times 2.0455) polystyrene's radius is R0 = 7.159; pure hexane lies outside the sphere (Ra 9.238, RED 1.290) and pure acetone just inside its edge (Ra 6.638, RED 0.927 -- and rated 3, "strongly swollen, slight solubility", in the thesis's own test, Table 2.3). The blend dips well inside: the minimum is Ra 4.724 (RED 0.660) at 65 \% acetone by volume, and 50:50 gives 4.992 (RED 0.697). The mixture rule lets two solvents on opposite sides of the target meet it between them, which is the observation Hansen checked his placements against (Section 2.5: "mixtures of non-solvents ... should ... dissolve the solute. This has always been found to be the case"). UNIFAC is asked whether the pair is one liquid at every point.

2. hexane + nitromethane, by MOLE (the basis conversion exercised: mole fractions become volume fractions through the records' Vliq, and the run prints both). Both ends are far outside the sphere (RED 1.290 and 1.850); the mixture rule puts the middle of the line inside (Ra 5.399, RED 0.754 at x = 0.6, which is 38 \% nitromethane by volume). But hexane and nitromethane are an immiscible pair at 25 C, and the case's UNIFAC is asked: where the Michelsen test finds a second liquid the op REFUSES to represent the pair as one solvent -- no averaged triple, no Ra, the point marked `twoPhase` -- and the minimum is taken over the points that survive. Measured: 17 of 19 blends (x = 0.05 to 0.85) are two-phase, every one where the rule would have put the blend inside the sphere among them, and the survivors all lie outside it -- the minimum over survivors is pure hexane, RED 1.290. The averaged number would have been the best on the line, and it describes a solvent that does not exist. (A two-phase pair is not useless -- it may serve an extraction; it is only not one solvent.)

T = 298.15 K is also every record's Tref, so no parameter is used away from the temperature it was placed at; the run says so either way.

\end{tutnarrative}
\tutreadme{The real case of the solvent-selection lesson (DEV.md 4c C128). The target is polystyrene as Hansen placed it in his 1967 thesis (solute G, Table 4.1: centre (8.6, 3.0, 2.0), radius 3.5 in (cal/cm$^3$)\textasciicircum{}0.5 --- R0 = 7.159 MPa\textasciicircum{}0.5 in the engine's unit); the solvents are three of his Table 2.1 placements, transcribed from the page images into case-local overlay records (constant/components/$\ast$.dat) with the thesis page and row beside each value, converted to Pa\textasciicircum{}0.5 by sqrt(4.184e6) = 2045.48, and marked interim. The reference temperature is adopted as 25 $^\circ$C (the thesis prints none beside its tables); the run is at 25 $^\circ$C so nothing is moved away from it.}
\tutreadme{Pure hexane lies outside the sphere and pure acetone just inside its edge (in the thesis's own test both were rated 3, "strongly swollen, slight solubility"); the blend dips well inside near 65 \% acetone. That is the observation Hansen checked his placements against: $\ast$two non-solvents on opposite sides of a solute's sphere give a blend that dissolves it$\ast$ (thesis, ?2.5, "this has always been found to be the case"). UNIFAC --- the case's activity model, predictive from the records' groups --- finds the pair one liquid at every point, so the word is homogeneous all along, which is what makes the blend a candidate at all.}
\tutreadme{The composition is declared by MOLE, so the run converts it to volume fractions through the records' Vliq (54.1 cm$^3$/mol for nitromethane against 131.6 for hexane: equal moles is 29 \% nitromethane by volume) and prints both. Both ends are far outside the sphere (RED 1.290 and 1.850), and the mixture rule puts the middle of the line inside it --- on paper, a better blend than the first (Ra 5.399, RED 0.754 at x = 0.6).}
\tutreadme{But hexane and nitromethane are an immiscible pair at room temperature, and the case's UNIFAC says so: where Michelsen's tangent-plane test finds a second liquid the op refuses to represent the pair as one solvent --- no averaged triple, no Ra, no RED for that point, the row marked twoPhase with its negative tm --- and the minimum is taken over the points that survive. The averaged number would have been the best on the line, and it describes a solvent that does not exist. A two-phase pair is not useless (it may serve an extraction); it is only not one solvent.}
\tutreadme{What UNIFAC answered, measured: 17 of the 19 blends are two-phase, from x = 0.05 to 0.85 (tm between -0.10 and -0.40), and only the two nitromethane-rich blends (x = 0.90 and 0.95) are one liquid --- both outside the sphere (RED 1.340 and 1.572). Every composition where the mixture rule would have put the blend inside the sphere (roughly x = 0.3 to 0.8, RED down to 0.754) is refused, so the minimum over the points that survive is pure hexane, RED 1.290: no blend of these two is a solvent for polystyrene. The counts are pinned in expected (n\_twoPhase, n\_homogeneous).}
\tutreadme{This case also found an engine defect, fixed in the same slice. Its first run reported hexane + acetone two-phase at EVERY composition (tm down to -0.43), for a pair whose UNIFAC activity coefficients at 50:50 are 1.52 and 1.56 and which is miscible. The Michelsen tangent-plane distance in src/solver/StabilityTest.cpp omitted the ln $\gamma$$_i$(y) term, so at the trivial point it read -G\textasciicircum{}E/RT instead of 0 and called every positive-deviation mixture unstable. With the term restored the pair is one liquid (tm at round-off) and the 38 corpus cases that reach the test kept every golden.}
\begin{tutgolden}
\goldrow{diag}{hexaneAcetone}{R0\_MPa05}{7.1592}
\goldrow{diag}{hexaneAcetone}{RED\_min}{0.659789277374}
\goldrow{diag}{hexaneAcetone}{Ra\_min\_MPa05}{4.72356339458}
\goldrow{diag}{hexaneAcetone}{Ra\_pure\_acetone}{6.63812052919}
\goldrow{diag}{hexaneAcetone}{Ra\_pure\_nHexane}{9.23837412102}
\goldrow{diag}{hexaneAcetone}{T\_K}{298.15}
\goldrow{diag}{hexaneAcetone}{Tref\_target\_K}{298.15}
\goldrow{diag}{hexaneAcetone}{n\_homogeneous}{19}
\goldrow{diag}{hexaneAcetone}{n\_notVerified}{0}
\goldrow{diag}{hexaneAcetone}{n\_points}{21}
\goldrow{diag}{hexaneAcetone}{n\_pure}{2}
\goldrow{diag}{hexaneAcetone}{n\_twoPhase}{0}
\goldrow{diag}{hexaneAcetone}{phiB\_at\_min}{0.65}
\goldrow{diag}{hexaneAcetone}{s\_at\_min}{0.65}
\goldmore{67}
\end{tutgolden}

\clearpage
\tutentry{solvent03\_polystyrene\_observed}{tutorials/props/solvents/solvent03\_polystyrene\_observed}
\begin{tutscheme}
\begin{tikzpicture}[>={Stealth[length=2mm]}, every node/.style={font=\scriptsize}]
  \node[draw, rounded corners=2pt, align=center, minimum width=3.1cm, minimum height=0.8cm, fill=black!4] at (0.0,-0.0) {\textbf{screen}\\[-1pt]{\tiny hansenMixture}};
\end{tikzpicture}
\end{tutscheme}
\tutwhat{Polystyrene (Hansen 1967, solute G) screened against thirteen pure solvents of the same thesis, beside the thesis's own dissolution ratings (Table 2.3) kept in constant/experimental/ -- the confrontation of Ra and RED with what was observed}
\begin{tutnarrative}
solvent03 -- THE CONFRONTATION WITH EXPERIMENT (DEV.md 4c C128, adjustment 2, step three).

No `composition \{\}`: the op SCREENS each declared solvent pure -- one row each, Ra and RED against polystyrene's sphere (Hansen 1967, Table 4.1, solute G). Beside the run, constant/experimental/ hansen1967\_table23\_polystyrene.csv holds what the same thesis OBSERVED for each of these thirteen solvents (Table 2.3, column G, the 1-6 scale of page 20: 1 soluble ... 6 no visible effect), with the citation on every row. The op does not read that file: the prediction and the observation stay two tables, joined by the reader (the README does it by hand, the EduTool draws them side by side, the gate counts the agreement).

WHAT THE CONFRONTATION IS, said plainly. Hansen's sphere for this polymer was DRAWN THROUGH these very observations (Section 4.4: "the centers of the volumes of solubility ... have been determined" from the solvents that dissolve each solute), so agreement here shows that the transcription and the distance arithmetic reproduce his placement -- it is a check of the method on its own training set, not an independent prediction. The thesis itself states the residual: about 2.5 \% of the solvents beyond R0 + 0.5 dissolve, and about 2.5 \% within R0 - 0.5 do not (page 39). Nitroethane (RED 1.39, rated 2 "almost soluble") and cyclohexane (RED 1.06, rated 1 "soluble") are that residual, in this set.

The package is `ideal` because no mixture is asked for: every row is one pure solvent and the phase word is `pure`.

\end{tutnarrative}
\tutreadme{The confrontation with experiment of the solvent-selection lesson (DEV.md 4c C128). No composition \{\}: the hansenMixture op screens each declared solvent PURE against polystyrene (Hansen 1967, solute G) --- one row each, the phase word pure --- and constant/experimental/hansen1967\_table23\_polystyrene.csv holds what the same thesis observed for each solvent (Table 2.3, column G, the 1--6 scale of its page 20: 1 soluble, 2 almost soluble, 3 strongly swollen with slight solubility, 4 swollen, 5 slight effect, 6 no visible effect), with the citation on every row. The op does not read that file; the two tables are joined here, by the EduTool, and by the gate.}
\tutreadme{Ra in MPa\textasciicircum{}0.5 against R0 = 7.159 (the thesis's 3.5 (cal/cm$^3$)\textasciicircum{}0.5); RED = Ra/R0.}
\tutreadme{Every solvent the thesis rated 6 lies outside the sphere; every solvent it rated 1 lies inside or --- cyclohexane, RED 1.058 --- within the $\pm$0.5 (cal/cm$^3$)\textasciicircum{}0.5 band the thesis states for its own radii. The residual, read across the line RED = 1: nitroethane (RED 1.39, rated 2) and cyclohexane (RED 1.06, rated 1) dissolve from outside; acetone (RED 0.93, rated 3) does not quite from inside. Three of thirteen, where the thesis itself expects about one solvent in forty on each side beyond $\pm$0.5 of the radius (page 39) --- this set is small and its residual sits on the boundary, as the thesis says the deviations do.}
\tutreadme{Hansen's sphere for this polymer was drawn through these very observations (thesis ?4.4: the centres were determined from the solvents that dissolve each solute). Agreement here therefore shows that the transcription and the distance arithmetic reproduce his placement --- a check of the method on its own training set, not an independent prediction of a solvent he never tried. No blend was dissolved by anyone: the thesis tabulates pure solvents only, and a measured point along a mixture line is what this corpus does not have.}
\tutreadme{The package is ideal because no mixture is asked for. Every parameter is a case-local overlay record with the thesis page and row beside the value, reviewStatus interim (transcribed from the page images, not the OCR text, whose columns are scrambled; a second reading pending).}
\tutreadme{expected pins each solvent's Ra\_pure\_<name> and the CSV cells of the screen; check\_hansen\_mixture rebuilds them from the records and prints the residual above as a measurement, failing only on a rated-6 solvent inside the sphere or a rated-1 solvent beyond the thesis's own band.}
\begin{tutgolden}
\goldrow{diag}{screen}{R0\_MPa05}{7.1592}
\goldrow{diag}{screen}{RED\_min}{0.636685199779}
\goldrow{diag}{screen}{Ra\_min\_MPa05}{4.55815668226}
\goldrow{diag}{screen}{Ra\_pure\_acetone}{6.63812052919}
\goldrow{diag}{screen}{Ra\_pure\_benzene}{5.6906984791}
\goldrow{diag}{screen}{Ra\_pure\_cyclohexane}{7.57262071677}
\goldrow{diag}{screen}{Ra\_pure\_diethylEther}{7.20865879204}
\goldrow{diag}{screen}{Ra\_pure\_ethanol}{15.9714933081}
\goldrow{diag}{screen}{Ra\_pure\_ethylAcetate}{7.0242390805}
\goldrow{diag}{screen}{Ra\_pure\_methanol}{19.8084619809}
\goldrow{diag}{screen}{Ra\_pure\_methylEthylKetone}{4.55815668226}
\goldrow{diag}{screen}{Ra\_pure\_nButanol}{12.1057857919}
\goldrow{diag}{screen}{Ra\_pure\_nHexane}{9.23837412102}
\goldrow{diag}{screen}{Ra\_pure\_nitroethane}{9.97063233401}
\goldmore{48}
\end{tutgolden}

\subsection{steam}
\grouptheory{steam}
\clearpage
\tutentry{steam01\_if97\_verification}{tutorials/props/steam/steam01\_if97\_verification}
\begin{tutscheme}
\begin{tikzpicture}[>={Stealth[length=2mm]}, every node/.style={font=\scriptsize}]
  \node[draw, rounded corners=2pt, align=center, minimum width=3.1cm, minimum height=0.8cm, fill=black!4] at (0.0,-0.0) {\textbf{r1\_T300\_p3MPa}\\[-1pt]{\tiny steamTables}};
  \node[draw, rounded corners=2pt, align=center, minimum width=3.1cm, minimum height=0.8cm, fill=black!4] at (3.6,-0.0) {\textbf{r1\_T300\_p80MPa}\\[-1pt]{\tiny steamTables}};
  \node[draw, rounded corners=2pt, align=center, minimum width=3.1cm, minimum height=0.8cm, fill=black!4] at (7.2,-0.0) {\textbf{r1\_T500\_p3MPa}\\[-1pt]{\tiny steamTables}};
  \node[draw, rounded corners=2pt, align=center, minimum width=3.1cm, minimum height=0.8cm, fill=black!4] at (0.0,-1.3) {\textbf{r2\_T300\_p3p5kPa}\\[-1pt]{\tiny steamTables}};
  \node[draw, rounded corners=2pt, align=center, minimum width=3.1cm, minimum height=0.8cm, fill=black!4] at (3.6,-1.3) {\textbf{r2\_T700\_p3p5kPa}\\[-1pt]{\tiny steamTables}};
  \node[draw, rounded corners=2pt, align=center, minimum width=3.1cm, minimum height=0.8cm, fill=black!4] at (7.2,-1.3) {\textbf{r2\_T700\_p30MPa}\\[-1pt]{\tiny steamTables}};
  \node[draw, rounded corners=2pt, align=center, minimum width=3.1cm, minimum height=0.8cm, fill=black!4] at (0.0,-2.6) {\textbf{satline}\\[-1pt]{\tiny steamTables}};
  \node[draw, rounded corners=2pt, align=center, minimum width=3.1cm, minimum height=0.8cm, fill=black!4] at (3.6,-2.6) {\textbf{isobar\_1bar}\\[-1pt]{\tiny steamTables}};
  \node[draw, rounded corners=2pt, align=center, minimum width=3.1cm, minimum height=0.8cm, fill=black!4] at (7.2,-2.6) {\textbf{isobar\_1MPa}\\[-1pt]{\tiny steamTables}};
  \node[draw, rounded corners=2pt, align=center, minimum width=3.1cm, minimum height=0.8cm, fill=black!4] at (0.0,-3.9) {\textbf{isobar\_10MPa}\\[-1pt]{\tiny steamTables}};
\end{tikzpicture}
\end{tutscheme}
\tutwhat{IAPWS-IF97 steam tables (R7-97(2012)) -- the verification case. The six computer-program verification states of the release (Table 5 region 1, Table 15 region 2) evaluated as points, the saturation line 0.01-350 degC, and three isobars whose Tsat crossings pin the release's Table 36 backward values. The published numbers ARE the golden master (reltol 1e-8).}
\begin{tutnarrative}
the IAPWS-IF97 weapon case.

The release (IAPWS R7-97(2012)) publishes computer-program verification tables; this case evaluates EVERY one of them this slice covers and locks the published numbers as golden KPIs at reltol 1e-8 -- a single wrong digit in a transcribed coefficient fails the regression:

Table 5 -- region 1 at (300 K, 3 MPa), (300 K, 80 MPa), (500 K, 3 MPa) Table 15 -- region 2 at (300 K, 3.5 kPa), (700 K, 3.5 kPa), (700 K, 30 MPa) Table 35 -- psat at 300 / 500 K (the psat\_MPa diag of the point ops) Table 36 -- Tsat at 0.1 / 1 / 10 MPa (the Tsat\_K diag of the isobars)

Plus the working artefacts: the saturated-steam table 0.01-350 degC (saturation.csv) and three isobars (h, s, v, cp vs T) crossing the saturation line -- the crossing is ANNOUNCED in the log.

\end{tutnarrative}
\begin{tutgolden}
\goldrow{diag}{r1\_T300\_p3MPa}{region}{1}
\goldrow{diag}{r1\_T300\_p3MPa}{v}{0.100215168e-2}
\goldrow{diag}{r1\_T300\_p3MPa}{h}{0.115331273e3}
\goldrow{diag}{r1\_T300\_p3MPa}{u}{0.112324818e3}
\goldrow{diag}{r1\_T300\_p3MPa}{s}{0.392294792}
\goldrow{diag}{r1\_T300\_p3MPa}{cp}{0.417301218e1}
\goldrow{diag}{r1\_T300\_p3MPa}{w}{0.150773921e4}
\goldrow{diag}{r1\_T300\_p3MPa}{psat\_MPa}{0.353658941e-2}
\goldrow{diag}{r1\_T300\_p80MPa}{region}{1}
\goldrow{diag}{r1\_T300\_p80MPa}{v}{0.971180894e-3}
\goldrow{diag}{r1\_T300\_p80MPa}{h}{0.184142828e3}
\goldrow{diag}{r1\_T300\_p80MPa}{u}{0.106448356e3}
\goldrow{diag}{r1\_T300\_p80MPa}{s}{0.368563852}
\goldrow{diag}{r1\_T300\_p80MPa}{cp}{0.401008987e1}
\goldmore{47}
\end{tutgolden}

\clearpage
\tutentry{steam02\_water\_transport}{tutorials/props/steam/steam02\_water\_transport}
\begin{tutscheme}
\begin{tikzpicture}[>={Stealth[length=2mm]}, every node/.style={font=\scriptsize}]
  \node[draw, rounded corners=2pt, align=center, minimum width=3.1cm, minimum height=0.8cm, fill=black!4] at (0.0,-0.0) {\textbf{pin\_25C}\\[-1pt]{\tiny propertyPoint}};
  \node[draw, rounded corners=2pt, align=center, minimum width=3.1cm, minimum height=0.8cm, fill=black!4] at (3.6,-0.0) {\textbf{pin\_50C}\\[-1pt]{\tiny propertyPoint}};
  \node[draw, rounded corners=2pt, align=center, minimum width=3.1cm, minimum height=0.8cm, fill=black!4] at (7.2,-0.0) {\textbf{pin\_75C}\\[-1pt]{\tiny propertyPoint}};
  \node[draw, rounded corners=2pt, align=center, minimum width=3.1cm, minimum height=0.8cm, fill=black!4] at (0.0,-1.3) {\textbf{pin\_100C}\\[-1pt]{\tiny propertyPoint}};
  \node[draw, rounded corners=2pt, align=center, minimum width=3.1cm, minimum height=0.8cm, fill=black!4] at (3.6,-1.3) {\textbf{transport\_scan}\\[-1pt]{\tiny propertyScan1D}};
\end{tikzpicture}
\end{tutscheme}
\tutwhat{IAPWS water transport properties -- viscosity (R12-08), thermal conductivity (R15-11) and surface tension (R1-76). Water flagged as a pure fluid (method IF97) routes the transport accessors to the matching IAPWS releases: the density comes from the IF97 thermodynamic kernel at the saturation state, fed to the (rho,T) transport kernels; surface tension is the R1-76 saturation-line form. A coarse 4-point scan (25/50/75/100 degC) pinned against the published IAPWS values, plus a fine sweep for the plot.}
\begin{tutnarrative}
the IAPWS water transport route.

Water flagged `pureFluids \{ water \{ method IF97; \} \}` routes the transport accessors to the matching IAPWS releases:

viscosity $\to$ R12-08 (industrial, critical enhancement off) thermal\_conductivity $\to$ R15-11 (background, critical enhancement off) surface\_tension $\to$ R1-76 (saturation line)

The (rho,T) viscosity / conductivity kernels get their density from the IF97 thermodynamic kernel at the SATURATION state (the legacy transport accessors carry no pressure, so p = p\_sat(T) is the physical convention -- saturated liquid / saturated vapour). Surface tension is a saturation-line property and takes only T.

`pins` evaluates the three properties at 25/50/75/100 degC; the golden master locks them at reltol 1e-5 (the formulation is verified to 1e-8 by IAPWSTransport::verify(), which fires on every run). `scan` writes a fine sweep for the plot.

\end{tutnarrative}
\begin{tutgolden}
\goldrow{diag}{pin\_25C}{surface\_tension}{0.0719722052302}
\goldrow{diag}{pin\_25C}{thermal\_conductivity}{0.606460966251}
\goldrow{diag}{pin\_25C}{viscosity}{0.000890036037708}
\goldrow{diag}{pin\_50C}{surface\_tension}{0.0679439059449}
\goldrow{diag}{pin\_50C}{thermal\_conductivity}{0.640589598777}
\goldrow{diag}{pin\_75C}{surface\_tension}{0.0635830195097}
\goldrow{diag}{pin\_75C}{thermal\_conductivity}{0.66354437705}
\goldrow{diag}{pin\_100C}{surface\_tension}{0.0589118685877}
\goldrow{diag}{pin\_100C}{thermal\_conductivity}{0.677216843797}
\goldrow{diag}{transport\_scan}{n\_points}{46}
\goldrow{diag}{transport\_scan}{n\_failures}{0}
\goldrow{diag}{pin\_25C}{Cp\_ig}{33.6596038227}
\goldrow{diag}{pin\_25C}{H\_R}{0}
\goldrow{diag}{pin\_25C}{H\_ig}{-241826}
\goldmore{49}
\end{tutgolden}

\clearpage
\tutentry{steam03\_if97\_region3}{tutorials/props/steam/steam03\_if97\_region3}
\begin{tutscheme}
\begin{tikzpicture}[>={Stealth[length=2mm]}, every node/.style={font=\scriptsize}]
  \node[draw, rounded corners=2pt, align=center, minimum width=3.1cm, minimum height=0.8cm, fill=black!4] at (0.0,-0.0) {\textbf{r3\_T650\_rho500}\\[-1pt]{\tiny steamTables}};
  \node[draw, rounded corners=2pt, align=center, minimum width=3.1cm, minimum height=0.8cm, fill=black!4] at (3.6,-0.0) {\textbf{r3\_T650\_rho200}\\[-1pt]{\tiny steamTables}};
  \node[draw, rounded corners=2pt, align=center, minimum width=3.1cm, minimum height=0.8cm, fill=black!4] at (7.2,-0.0) {\textbf{r3\_T750\_rho500}\\[-1pt]{\tiny steamTables}};
  \node[draw, rounded corners=2pt, align=center, minimum width=3.1cm, minimum height=0.8cm, fill=black!4] at (0.0,-1.3) {\textbf{r3\_T650\_p25p58MPa}\\[-1pt]{\tiny steamTables}};
  \node[draw, rounded corners=2pt, align=center, minimum width=3.1cm, minimum height=0.8cm, fill=black!4] at (3.6,-1.3) {\textbf{drum\_water}\\[-1pt]{\tiny steamTables}};
  \node[draw, rounded corners=2pt, align=center, minimum width=3.1cm, minimum height=0.8cm, fill=black!4] at (7.2,-1.3) {\textbf{drum\_steam}\\[-1pt]{\tiny steamTables}};
  \node[draw, rounded corners=2pt, align=center, minimum width=3.1cm, minimum height=0.8cm, fill=black!4] at (0.0,-2.6) {\textbf{oncethrough\_25MPa\_650K}\\[-1pt]{\tiny steamTables}};
  \node[draw, rounded corners=2pt, align=center, minimum width=3.1cm, minimum height=0.8cm, fill=black!4] at (3.6,-2.6) {\textbf{satline\_near\_critical}\\[-1pt]{\tiny steamTables}};
  \node[draw, rounded corners=2pt, align=center, minimum width=3.1cm, minimum height=0.8cm, fill=black!4] at (7.2,-2.6) {\textbf{isobar\_25MPa}\\[-1pt]{\tiny steamTables}};
\end{tikzpicture}
\end{tutscheme}
\tutwhat{IAPWS-IF97 region 3 (near-critical and supercritical water, the Helmholtz form f(rho,T)). The three verification states printed in Wagner et al. 2000 Table A10 evaluated at (T, rho) as printed, one of them again from its printed pressure through the backward density, a 20 MPa drum boiler's water and steam on either side of the saturation line, a 25 MPa supercritical once-through point, the saturation line up to the critical point and a 25 MPa isobar crossing regions 1, 3 and 2.}
\begin{tutnarrative}
IAPWS-IF97 region 3, where the boiler drum and the supercritical once-through boiler live.

Region 3 is IF97's near-critical region: 623.15 K < T <= T\_B23(p) and p\_B23(T) < p <= 100 MPa. Its equation is a Helmholtz energy f(rho, T) (Wagner et al., J. Eng. Gas Turbines Power 122 (2000) 150-182, Eq. 25, 40 coefficients in Table A9), so the engine reaches it from (p, T) by solving p(rho, T) = p for the density on ONE branch of the isotherm -- below Tc the liquid branch when p >= psat(T), the vapour branch below.

VALIDATED (the golden rows of the four Table A10 operations ARE the paper's printed values, not a recording, at reltol 1e-8 -- the printed precision; the engine's own 1e-8 tripwire is IF97::verify(), which every steamTables run executes and which refuses the run above it): r3\_T650\_rho500, r3\_T650\_rho200, r3\_T750\_rho500 -- Table A10 at the printed (T, rho): p, h, u, s, cp, w; r3\_T650\_p25p58MPa -- the first state again, from its PRINTED pressure: the backward density must return rho = 500 kg/m3. NOT VALIDATED against any printed value (self-recorded goldens; they pin that the answer has not moved, never that it is right): drum\_water / drum\_steam -- a 20 MPa drum boiler (Tsat = 638.9 K): water 4 K subcooled, steam 6 K superheated, both region 3; oncethrough\_25MPa\_650K -- supercritical pressure, just below the pseudo-critical temperature, where cp peaks; satline\_near\_critical -- the saturation line 600 K .. Tc: regions 1/2 on the line up to 623.15 K, region 3's two branches above it, h\_fg falling to zero at Tc; isobar\_25MPa -- 600 .. 900 K: region 1, then region 3 (the cp peak), then region 2 above T\_B23(25 MPa). IF97's own stated uncertainty against IAPWS-95 near the critical point is large (cp up to -53 \% near 648 K and 22.3 MPa; the paper's Fig. 16b) -- a near-critical cp from this case is the industrial formulation's number, not the scientific one.

\end{tutnarrative}
\begin{tutgolden}
\goldrow{diag}{r3\_T650\_rho500}{region}{3}
\goldrow{diag}{r3\_T650\_rho500}{p\_MPa}{0.255837018e2}
\goldrow{diag}{r3\_T650\_rho500}{h}{0.186343019e4}
\goldrow{diag}{r3\_T650\_rho500}{u}{0.181226279e4}
\goldrow{diag}{r3\_T650\_rho500}{s}{0.405427273e1}
\goldrow{diag}{r3\_T650\_rho500}{cp}{0.138935717e2}
\goldrow{diag}{r3\_T650\_rho500}{w}{0.502005554e3}
\goldrow{diag}{r3\_T650\_rho200}{region}{3}
\goldrow{diag}{r3\_T650\_rho200}{p\_MPa}{0.222930643e2}
\goldrow{diag}{r3\_T650\_rho200}{h}{0.237512401e4}
\goldrow{diag}{r3\_T650\_rho200}{u}{0.226365868e4}
\goldrow{diag}{r3\_T650\_rho200}{s}{0.485438792e1}
\goldrow{diag}{r3\_T650\_rho200}{cp}{0.446579342e2}
\goldrow{diag}{r3\_T650\_rho200}{w}{0.383444594e3}
\goldmore{34}
\end{tutgolden}

\subsection{thermo}
\grouptheory{thermo}
\clearpage
\tutentry{bubble01\_reaction\_mass\_boiling\_point}{tutorials/props/thermo/bubble01\_reaction\_mass\_boiling\_point}
\begin{tutscheme}
\begin{tikzpicture}[>={Stealth[length=2mm]}, every node/.style={font=\scriptsize}]
  \node[draw, rounded corners=2pt, align=center, minimum width=3.1cm, minimum height=0.8cm, fill=black!4] at (0.0,-0.0) {\textbf{reactionMassBubblePoint}\\[-1pt]{\tiny propertyPoint}};
\end{tikzpicture}
\end{tutscheme}
\tutwhat{The boiling point of a reaction mass: the bubble point of the cooling-failure batch's contents (fictitious A + B $\to$ P in solvent S, ideal), the technical limit (MTT) of the scenario}
\begin{tutnarrative}
the BOILING POINT OF A REACTION MASS, which is what the thermal-safety methodology takes as the technical limit MTT of an open vessel at atmospheric pressure.

The mass is the one in tutorials/batch/reactor/batch26\_cooling\_failure (and batch27/batch28, the same reaction dosed): the FICTITIOUS reactants A and B, their product P and the inert solvent S. constant/ is a COPY of batch26's sealed constant/ (the same four authored records under the same manifest, the same thermoPhysPropDict; only the reactor's own constant/reactions is left out, because no property operation reads it), so the boiling point below is computed on exactly the thermodynamics the reactor case declares: an IDEAL liquid (Raoult, gamma = 1) over an ideal gas, at the reactor's pressure, 1.013 bar.

The composition is the reactor's at batch26's shipped cooling failure (t\_f = 600 s; mole fractions of the trajectory row at 600 s, 39.3 \% of B converted). The process-safety EduTool rewrites the four lines with the reactor contents at the failure instant the reader chooses, read from the engine's own trajectory, and reads T\_bubble back. ONE KEY PER LINE: the GUI's override anchors on a key at the start of a line, so a composition written on one line is a knob that silently misses.

WHAT THIS IS AND IS NOT. The mixture is ideal BY DEFINITION (the substances are fictitious and the case declares it so), so it has no azeotrope: its bubble point lies between the normal boiling points of the components present -- here between B's (345.01 K) and P's (479.99 K), and much nearer the light end, because the volatile components dominate the vapour pressure. It is a bubble point at the failure instant: as a runaway converts the mass, its composition, and so its boiling point, moves -- upward here, as the light reactants become the heavy product.

\end{tutnarrative}
\begin{tutgolden}
\goldrow{diag}{reactionMassBubblePoint}{Cp\_ig}{113.71926727}
\goldrow{diag}{reactionMassBubblePoint}{H\_R}{0}
\goldrow{diag}{reactionMassBubblePoint}{H\_ig}{-304682.453301}
\goldrow{diag}{reactionMassBubblePoint}{H\_real}{-304682.453301}
\goldrow{diag}{reactionMassBubblePoint}{P}{101300}
\goldrow{diag}{reactionMassBubblePoint}{S\_R}{0}
\goldrow{diag}{reactionMassBubblePoint}{S\_ig}{401.345731157}
\goldrow{diag}{reactionMassBubblePoint}{S\_real}{401.345731157}
\goldrow{diag}{reactionMassBubblePoint}{T}{330}
\goldrow{diag}{reactionMassBubblePoint}{T\_bubble}{376.000570165}
\goldrow{diag}{reactionMassBubblePoint}{Z}{1}
\goldrow{diag}{reactionMassBubblePoint}{gamma}{1.07888124878}
\goldrow{diag}{reactionMassBubblePoint}{v\_molar}{0.0270856136618}
\end{tutgolden}

\clearpage
\tutentry{flashSpec01\_wet\_steam\_expansion}{tutorials/props/thermo/flashSpec01\_wet\_steam\_expansion}
\begin{tutscheme}
\begin{tikzpicture}[>={Stealth[length=2mm]}, every node/.style={font=\scriptsize}]
  \node[draw, rounded corners=2pt, align=center, minimum width=3.1cm, minimum height=0.8cm, fill=black!4] at (0.0,-0.0) {\textbf{expand}\\[-1pt]{\tiny flashSpec}};
  \node[draw, rounded corners=2pt, align=center, minimum width=3.1cm, minimum height=0.8cm, fill=black!4] at (3.6,-0.0) {\textbf{throttle}\\[-1pt]{\tiny flashSpec}};
  \node[draw, rounded corners=2pt, align=center, minimum width=3.1cm, minimum height=0.8cm, fill=black!4] at (7.2,-0.0) {\textbf{halfBoil}\\[-1pt]{\tiny flashSpec}};
  \node[draw, rounded corners=2pt, align=center, minimum width=3.1cm, minimum height=0.8cm, fill=black!4] at (0.0,-1.3) {\textbf{dewP}\\[-1pt]{\tiny flashSpec}};
\end{tikzpicture}
\end{tutscheme}
\tutwhat{The general flash specification (DEV.md 4c C214 W1): steam at 80 bar and 500 C expanded isentropically to 0.1 bar on IAPWS-IF97 ends WET (vapour fraction about 0.81, checkable against the steam table); a subcooled liquid throttled at constant enthalpy flashes; a vapour-fraction specification on a pure fluid lands on its saturation curve -- the (P,S), (P,H), (P,V/F) and (T,V/F) pairs through one home}
\begin{tutnarrative}
A FLASH SPECIFIED BY A PAIR OTHER THAN (T, P) (DEV.md 4c C214 W1; docs/design/a-general-flash-specification.md).

A state of a fluid of known composition has two degrees of freedom. The isothermal flash fixes them as (T, P); a turbine knows its outlet PRESSURE and, ideally, its ENTROPY; a valve knows its pressure and its ENTHALPY; a condenser drum knows its pressure and how much of the stream is VAPOUR. Each op below asks one of those pairs of ONE home (flashState::specifiedFlash), and every answer is the state the isothermal flash publishes there -- except a pure fluid ON its saturation curve, where no flash can decide the split and the saturated legs price it.

THE HAND CHECK (IAPWS-IF97 steam table, per kg): expand 80 bar, 500 C: s = 6.7262 kJ/(kg K), h = 3399.4 kJ/kg 0.1 bar saturated: Tsat = 45.81 C, s\_f = 0.6492, s\_g = 8.1501, h\_f = 191.81, h\_g = 2583.9 quality x = (6.7262 - 0.6492)/(8.1501 - 0.6492) = 0.8102 h\_out = 191.81 + 0.8102 (2583.9 - 191.81) = 2129.8 kJ/kg isentropic work = 3399.4 - 2129.8 = 1269.6 kJ/kg throttle 10 bar, 450 K (subcooled liquid): h = 749.3 kJ/kg 1 bar saturated: h\_f = 417.4, h\_g = 2674.9 flash fraction x = (749.3 - 417.4)/(2674.9 - 417.4) = 0.147 halfBoil 1 atm, half vapour: T = Tsat(1 atm) = 373.124 K dewP 373.15 K, all vapour: P = Psat(100 C) = 101.418 kPa

WHAT THE STUDENT WATCHES. `expand` lands on the saturation PLATEAU: for a pure fluid T and P are one condition there, so the state's coordinate is the vapour fraction, found by the lever rule on the saturated legs (no search). An isentropic search over a single-phase VAPOUR of fixed composition -- what Choupo's turbine does today -- has no root here: below s\_g at 0.1 bar no vapour carries that entropy (the vapour entropy jumps by the whole latent entropy at Tsat), so that search cannot land on this answer. That gap is what this op makes visible.

NOT MODELLED, said plainly: any non-equilibrium (a real turbine's supersaturated "Wilson line" expansion is a kinetic effect); the irreversibility of a real expansion (the isentropic state is the ideal reference an efficiency is measured against). No number here is fitted: every value is an IF97 evaluation, a STRUCTURAL witness checkable by hand.

\end{tutnarrative}
\begin{tutgolden}
\goldrow{diag}{expand}{H}{38374.9393008}
\goldrow{diag}{expand}{H\_per\_mass}{2130.16593399}
\goldrow{diag}{expand}{P}{10000}
\goldrow{diag}{expand}{S}{121.175719105}
\goldrow{diag}{expand}{S\_per\_mass}{6.72637907882}
\goldrow{diag}{expand}{T}{318.957548207}
\goldrow{diag}{expand}{dH\_fromReference}{-22864.7579963}
\goldrow{diag}{expand}{flashes}{0}
\goldrow{diag}{expand}{iterations}{0}
\goldrow{diag}{expand}{onSaturationPlateau}{1}
\goldrow{diag}{expand}{ref\_H}{61239.6972971}
\goldrow{diag}{expand}{ref\_P}{8000000}
\goldrow{diag}{expand}{ref\_S}{121.175719105}
\goldrow{diag}{expand}{ref\_T}{773.15}
\goldmore{53}
\end{tutgolden}

\clearpage
\tutentry{heavyWater01\_isotope\_effect}{tutorials/props/thermo/heavyWater01\_isotope\_effect}
\begin{tutscheme}
\begin{tikzpicture}[>={Stealth[length=2mm]}, every node/.style={font=\scriptsize}]
  \node[draw, rounded corners=2pt, align=center, minimum width=3.1cm, minimum height=0.8cm, fill=black!4] at (0.0,-0.0) {\textbf{psat\_isotope}\\[-1pt]{\tiny propertyScan1D}};
  \node[draw, rounded corners=2pt, align=center, minimum width=3.1cm, minimum height=0.8cm, fill=black!4] at (3.6,-0.0) {\textbf{boiling\_water}\\[-1pt]{\tiny propertyScan1D}};
  \node[draw, rounded corners=2pt, align=center, minimum width=3.1cm, minimum height=0.8cm, fill=black!4] at (7.2,-0.0) {\textbf{boiling\_heavyWater}\\[-1pt]{\tiny propertyScan1D}};
\end{tikzpicture}
\end{tutscheme}
\tutwhat{Heavy water against ordinary water: the vapour-pressure isotope effect and the boiling points, from the records}
\begin{tutnarrative}
D2O against H2O, from the catalogue records.

WHAT THIS CASE EVALUATES. The saturation pressure of heavy water (record `heavyWater`: the Wagner equation of Herrig, Thol, Harvey \& Lemmon, J. Phys. Chem. Ref. Data 47, 043102 (2018), Eq. (1)) and of ordinary water (IAPWS-IF97, via `pureFluids`) at 25/50/75/100 degC, and the normal boiling point of each pure liquid as its bubble temperature at 101325 Pa. Nothing here is stated from memory: bin/curate/check\_heavy\_water.py recomputes the heavy-water side from the IAPWS R16-17(2018) Helmholtz equation (Maxwell criterion), reads the ordinary-water side from the IF97 saturation equation, and holds this case's published numbers (CSV cells, pinned by `csv` rows) to both.

WHAT IS VERIFIED. That the engine reproduces the primary equation: the Wagner arithmetic to 1e-6, the equation against the full D2O EOS to 8e-5, the boiling point of D2O against the EOS's. WHAT IS NOT. No measured datum is compared here; the ordinary-water side is IF97 (itself a correlation of IAPWS-95); `Hvap` is the engine's Watson estimate from the record's Tc, Tb and HvapTb (exact at Tb, an extrapolation elsewhere -- NOT the EOS latent heat). The record carries no formation datum, so `propertyPoint` (which prints H\_ig) refuses heavyWater by name: this case uses the scan operation, which evaluates only what it is asked for. No mixture property of D2O + H2O is claimed (the isotope-exchange equilibrium H2O + D2O = 2 HDO is not modelled).

\end{tutnarrative}
\begin{tutgolden}
\goldrow{csv}{psat\_isotope.csv}{0:Psat\_water}{3.1697468500e+03}
\goldrow{csv}{psat\_isotope.csv}{0:Psat\_heavyWater}{2.7362295800e+03}
\goldrow{csv}{psat\_isotope.csv}{0:Hvap\_heavyWater}{4.5626456900e+04}
\goldrow{csv}{psat\_isotope.csv}{1:Psat\_water}{1.2351270400e+04}
\goldrow{csv}{psat\_isotope.csv}{1:Psat\_heavyWater}{1.1117923400e+04}
\goldrow{csv}{psat\_isotope.csv}{1:Hvap\_heavyWater}{4.4343347400e+04}
\goldrow{csv}{psat\_isotope.csv}{2:Psat\_water}{3.8595362700e+04}
\goldrow{csv}{psat\_isotope.csv}{2:Psat\_heavyWater}{3.5826392400e+04}
\goldrow{csv}{psat\_isotope.csv}{2:Hvap\_heavyWater}{4.2996612600e+04}
\goldrow{csv}{psat\_isotope.csv}{3:Psat\_water}{1.0141797800e+05}
\goldrow{csv}{psat\_isotope.csv}{3:Psat\_heavyWater}{9.6302716000e+04}
\goldrow{csv}{psat\_isotope.csv}{3:Hvap\_heavyWater}{4.1577284500e+04}
\goldrow{csv}{boiling\_water.csv}{0:T\_bubble}{3.7312430000e+02}
\goldrow{csv}{boiling\_heavyWater.csv}{0:T\_bubble}{3.7455020600e+02}
\goldmore{9}
\end{tutgolden}

\clearpage
\tutentry{phaseSplit01\_freeze\_water\_ethanol}{tutorials/props/thermo/phaseSplit01\_freeze\_water\_ethanol}
\begin{tutscheme}
\begin{tikzpicture}[>={Stealth[length=2mm]}, every node/.style={font=\scriptsize}]
  \node[draw, rounded corners=2pt, align=center, minimum width=3.1cm, minimum height=0.8cm, fill=black!4] at (0.0,-0.0) {\textbf{freeze}\\[-1pt]{\tiny phaseSplitScan}};
\end{tikzpicture}
\end{tutscheme}
\tutwhat{Freeze-concentration scan: a 10 mol\% ethanol liquor cooled from 0 to -20 C, one SLE flash per kelvin -- the ice point, the liquor concentrating along the freezing-point-depression identity}
\begin{tutnarrative}
THE PHASE SPLIT ALONG T (DEV.md 4c C118).

The same liquor flash21 freezes at ONE temperature (-15 C), cooled here from 0 to -20 C one kelvin at a time: at every grid point ONE isothermal flash, the flash21 unit's own solver (IsothermalFlash::solveCore), under flash21's own sealed records -- water, ethanol, the curated NRTL pair, the C3 `phases ( liquor liquid; ice solid crystallizing )` declaration. Nothing is re-derived; the point at 258.15 K IS flash21's answer (the gate holds it to flash21's golden).

WHAT A STUDENT WATCHES. Above the liquor's freezing point the feed is one liquid and the declared ice stays subsaturated: the water activity gamma\_w x\_w (0.917) sits BELOW the saturation value exp(-dG\_fus/RT) the crystal would need (1 at the triple point, falling as T falls). Between -8 and -9 C the two cross, ice APPEARS (announced, between the two grid values it lies between), and from there down the liquor concentrates exactly along the freezing-point-depression identity

gamma\_w $\ast$ x\_w = exp(-dHfus (1 - T/Tt) / (R T))

with every quantity from the record's own data (Hfus 6008 J/mol, triple point 273.16 K) -- published at every point as `a\_water` beside `aEq\_water`, equal to round-off wherever ice is present, and recomputed independently by check\_phase\_split\_scan. By -20 C nearly 60 \% of the feed is ice and the liquor is about a quarter ethanol.

NOT MODELLED, said plainly (the run says it too): nucleation and supercooling -- every point is the equilibrium ceiling; an ethanol solid or a eutectic -- no record declares one, so the liquor never solidifies completely in this model; solute inclusion in the ice; the liquid/solid heat-capacity term of dG\_fus (announced by the record, |T - Tt| up to 20 K here). Every sub-273 K vapour pressure is an announced extrapolation of the Antoine fit (it cancels in the SLE identity: both phases share one Psat), and the NRTL pair is used below its fitted window (announced). No number here is anchored to a measured freezing-point datum: a STRUCTURAL witness, like flash21.

\end{tutnarrative}
\tutreadme{What this witnesses (DEV.md 4c C118 stage 1, 2026-10-20): the phaseSplitScan props operation. At a FIXED overall composition (90/10 water/ethanol, moles) it sweeps temperature at fixed pressure and, at every grid point, solves the full phase split through the ONE flash home --- IsothermalFlash::solveCore, the isothermalFlash unit's own solver, with the options the unit's own reader parses --- under the case's own sealed records. The constant/ tree is flash21's, copied verbatim: the same water, the same ethanol, the same curated NRTL pair, the same C3 phases ( liquor liquid; ice solid crystallizing ) declaration. Nothing is re-derived, which is why the point at 258.15 K is flash21's answer (check\_phase\_split\_scan arm (a) holds it to flash21's golden; its arms also hold single points on copies of flash16 and flash19 to those units' pH and V/F).}
\tutreadme{What a student watches. One isothermal flash per kelvin, 273.15 $\to$ 253.15 K:}
\tutreadme{Above the liquor's freezing point the feed is one liquid and the declared crystal stays subsaturated: the water activity $\gamma$\_w$\cdot$x\_w (0.917) lies BELOW the saturation value exp(-$\Delta$G\_fus/RT) the crystal would need (1 at the triple point, falling as T falls). Between 265.15 and 264.15 K the two cross and ice APPEARS --- announced by the run between the two grid values it lies between, recorded as an appeared event in the result block, pinned as diag first\_ice. From there down the liquor concentrates exactly along the freezing-point-depression identity}
\tutreadme{$\gamma$\_w $\cdot$ x\_w = exp(-$\Delta$H\_fus (1 - T/T\_t) / (R T)), $\Delta$H\_fus = 6008 J/mol, T\_t = 273.16 K}
\tutreadme{with every quantity from water.dat's own data --- published at every point as a\_water beside aEq\_water, equal to 1e-10 wherever ice is present, and recomputed independently by the gate from the record. By -20 $^\circ$C 63 \% of the feed is ice and the liquor is 27 mol\% ethanol.}
\tutreadme{What is published. freeze.csv (T first, then solved, nPhases, the amount of each phase --- L1, ice --- their compositions, the flash's own a\_water / aEq\_water / solidFraction) and the top-level result block phaseSplitScans (points, the event, the record-level announcements at the first point that raised each, what is not modelled). Every number in the block is a CSV cell or a diagnostic, so the csv and diag golden kinds pin all of it: the golden carries the diagnostics and ten authored CSV cells (rows 9, 15, 20).}
\begin{tutgolden}
\goldrow{diag}{freeze}{first\_L1}{273.15}
\goldrow{diag}{freeze}{first\_ice}{264.15}
\goldrow{diag}{freeze}{last\_L1}{253.15}
\goldrow{diag}{freeze}{last\_ice}{253.15}
\goldrow{diag}{freeze}{n\_announcements}{6}
\goldrow{diag}{freeze}{n\_events}{1}
\goldrow{diag}{freeze}{n\_points}{21}
\goldrow{diag}{freeze}{n\_refused}{0}
\goldrow{diag}{freeze}{n\_solved}{21}
\goldrow{diag}{freeze}{points\_L1}{21}
\goldrow{diag}{freeze}{points\_ice}{12}
\goldrow{csv}{freeze.csv}{9:ice}{3.2042311316e-02}
\goldrow{csv}{freeze.csv}{9:x\_L1\_ethanol}{1.0331030082e-01}
\goldrow{csv}{freeze.csv}{15:ice}{4.7599462325e-01}
\goldmore{7}
\end{tutgolden}

\clearpage
\tutentry{standardState01\_hydrogen\_enthalpy\_TP}{tutorials/props/thermo/standardState01\_hydrogen\_enthalpy\_TP}
\begin{tutscheme}
\begin{tikzpicture}[>={Stealth[length=2mm]}, every node/.style={font=\scriptsize}]
  \node[draw, rounded corners=2pt, align=center, minimum width=3.1cm, minimum height=0.8cm, fill=black!4] at (0.0,-0.0) {\textbf{h\_T300\_P1}\\[-1pt]{\tiny propertyPoint}};
  \node[draw, rounded corners=2pt, align=center, minimum width=3.1cm, minimum height=0.8cm, fill=black!4] at (3.6,-0.0) {\textbf{h\_T300\_P100}\\[-1pt]{\tiny propertyPoint}};
  \node[draw, rounded corners=2pt, align=center, minimum width=3.1cm, minimum height=0.8cm, fill=black!4] at (7.2,-0.0) {\textbf{h\_T300\_P300}\\[-1pt]{\tiny propertyPoint}};
  \node[draw, rounded corners=2pt, align=center, minimum width=3.1cm, minimum height=0.8cm, fill=black!4] at (0.0,-1.3) {\textbf{h\_T300\_P600}\\[-1pt]{\tiny propertyPoint}};
  \node[draw, rounded corners=2pt, align=center, minimum width=3.1cm, minimum height=0.8cm, fill=black!4] at (3.6,-1.3) {\textbf{h\_T300\_P1000}\\[-1pt]{\tiny propertyPoint}};
  \node[draw, rounded corners=2pt, align=center, minimum width=3.1cm, minimum height=0.8cm, fill=black!4] at (7.2,-1.3) {\textbf{h\_T400\_P1}\\[-1pt]{\tiny propertyPoint}};
  \node[draw, rounded corners=2pt, align=center, minimum width=3.1cm, minimum height=0.8cm, fill=black!4] at (0.0,-2.6) {\textbf{h\_T400\_P100}\\[-1pt]{\tiny propertyPoint}};
  \node[draw, rounded corners=2pt, align=center, minimum width=3.1cm, minimum height=0.8cm, fill=black!4] at (3.6,-2.6) {\textbf{h\_T400\_P300}\\[-1pt]{\tiny propertyPoint}};
  \node[draw, rounded corners=2pt, align=center, minimum width=3.1cm, minimum height=0.8cm, fill=black!4] at (7.2,-2.6) {\textbf{h\_T400\_P600}\\[-1pt]{\tiny propertyPoint}};
  \node[draw, rounded corners=2pt, align=center, minimum width=3.1cm, minimum height=0.8cm, fill=black!4] at (0.0,-3.9) {\textbf{h\_T400\_P1000}\\[-1pt]{\tiny propertyPoint}};
  \node[draw, rounded corners=2pt, align=center, minimum width=3.1cm, minimum height=0.8cm, fill=black!4] at (3.6,-3.9) {\textbf{h\_T500\_P1}\\[-1pt]{\tiny propertyPoint}};
  \node[draw, rounded corners=2pt, align=center, minimum width=3.1cm, minimum height=0.8cm, fill=black!4] at (7.2,-3.9) {\textbf{h\_T500\_P100}\\[-1pt]{\tiny propertyPoint}};
\end{tikzpicture}
\end{tutscheme}
\tutwhat{h(T,P) of hydrogen on SRK, 300-1000 K by 1-1000 bar, 42 points: the ideal-gas part against the residual, for the standard-state EduTool}
\begin{tutnarrative}
One propertyPoint per node of the page's grid. Names are the keys the GUI reads: h\_T<K>\_P<bar>; h\_converter is the green-ammonia converter's (T, P), h\_converter\_1bar its standard-pressure twin.

\end{tutnarrative}
\tutreadme{Witness of the standard-state EduTool ($\ast$Enthalpy, the standard state and equilibrium: what does the $^\circ$ mean?$\ast$). Forty-two propertyPoint operations on hydrogen, one per node of the page's grid --- 300 to 1000 K by 100 K, at 1, 100, 300, 600 and 1000 bar --- plus the green-ammonia converter's state (733.15 K, 150 bar) and its 1 bar twin. Each publishes the three numbers the page separates: H\_ig, the ideal-gas enthalpy (a function of temperature alone, on the elements datum --- hydrogen is an element, so it is zero at 298.15 K); H\_R, the SRK residual at the real temperature and pressure; and their sum H\_real.}
\tutreadme{A table gives a standard enthalpy with a degree sign. Does the sign mean 1 bar? Does it also mean 25 $^\circ$C? If the process runs at 150 bar, is using a standard property not assuming the gas is ideal? The page answers in the order a student can follow --- behaviour first, then the reference, only then the symbol --- and this case supplies the behaviour.}
\tutreadme{$\ast$ Along 1 bar, heating from 300 to 1000 K raises the enthalpy by 20.64 kJ/mol. At 300 K, compressing from 1 to 1000 bar at constant temperature raises it by 1.29 kJ/mol; at 1000 K by 1.69. The pressure effect is an order of magnitude smaller than the temperature effect in this range, and it is not zero. $\ast$ At the converter's state the ideal-gas rise from 298.15 K is 12.73 kJ/mol and the residual adds 0.23 kJ/mol --- 1.8 \% of the rise. At 1 bar and the same temperature the residual is 1.5 J/mol: the real gas at the standard pressure is a hair from the ideal gas, and still not the same definition.}
\tutreadme{constant/experimental/hydrogen\_Hmolar\_leachman2009\_coolprop.csv carries the same 42 states from a fundamental equation of state fitted to measurements --- Leachman, Jacobsen, Penoncello \& Lemmon, $\ast$J. Phys. Chem. Ref. Data$\ast$ 38 (2009) 721--748, as implemented in CoolProp 8.0.0 (Bell et al., $\ast$Ind. Eng. Chem. Res.$\ast$ 53 (2014) 2498; MIT licence). It is RE-COMPUTED here by bin/curate/reference\_h2\_enthalpy.py (--check re-derives every row), not pasted from anywhere, and it sits in constant/experimental/ because that is the one location the GUI bundle and .gitignore admit for a CSV that is a case input; its header says it is a model's output, not a measurement.}
\tutreadme{The two agree on the ideal-gas part to a few joules per mole over the whole range. They do not agree on the residual: SRK is 6 \% below the reference at 300 K and 1000 bar, a third below at 300 K and 100 bar, 11 \% below at the converter. Measured before the model was chosen: Peng-Robinson is worse here --- a third below at 1000 bar and the wrong sign at 100 bar --- so the case runs SRK and the page draws both sets and states the gap. Nothing was tuned; hydrogen's anomalous acentric factor (-0.220) is what a cubic equation handles poorly, and the page says so rather than hiding it.}
\tutreadme{A propertyScan2D would do this in one operation and write a CSV --- which the browser cannot read back. The GUI reads each operation's diagnostics, so the grid is 42 point evaluations whose names (h\_T<K>\_P<bar>) are the keys the page looks up. The golden pins every one of them.}
\begin{tutgolden}
\goldrow{diag}{h\_T300\_P1}{Cp\_ig}{28.885715}
\goldrow{diag}{h\_T300\_P1}{H\_R}{0.477247091364}
\goldrow{diag}{h\_T300\_P1}{H\_ig}{53.4333351109}
\goldrow{diag}{h\_T300\_P1}{H\_real}{53.9105822023}
\goldrow{diag}{h\_T300\_P1}{P}{100000}
\goldrow{diag}{h\_T300\_P1}{S\_R}{-0.00272994535243}
\goldrow{diag}{h\_T300\_P1}{S\_ig}{130.858662543}
\goldrow{diag}{h\_T300\_P1}{S\_real}{130.855932598}
\goldrow{diag}{h\_T300\_P1}{T}{300}
\goldrow{diag}{h\_T300\_P1}{Z}{1.00051987687}
\goldrow{diag}{h\_T300\_P1}{gamma}{1.40417872785}
\goldrow{diag}{h\_T300\_P1}{v\_molar}{0.0249563553443}
\goldrow{diag}{h\_T300\_P100}{Cp\_ig}{28.885715}
\goldrow{diag}{h\_T300\_P100}{H\_R}{63.6322761263}
\goldmore{490}
\end{tutgolden}

\clearpage
\tutentry{temperature01\_gas\_thermometer}{tutorials/props/thermo/temperature01\_gas\_thermometer}
\begin{tutscheme}
\begin{tikzpicture}[>={Stealth[length=2mm]}, every node/.style={font=\scriptsize}]
  \node[draw, rounded corners=2pt, align=center, minimum width=3.1cm, minimum height=0.8cm, fill=black!4] at (0.0,-0.0) {\textbf{Z\_of\_P}\\[-1pt]{\tiny propertyScan1D}};
  \node[draw, rounded corners=2pt, align=center, minimum width=3.1cm, minimum height=0.8cm, fill=black!4] at (3.6,-0.0) {\textbf{nearly\_ideal}\\[-1pt]{\tiny propertyPoint}};
  \node[draw, rounded corners=2pt, align=center, minimum width=3.1cm, minimum height=0.8cm, fill=black!4] at (7.2,-0.0) {\textbf{laboratory\_pressure}\\[-1pt]{\tiny propertyPoint}};
\end{tikzpicture}
\end{tutscheme}
\tutwhat{The gas thermometer: why the ideal gas is a limit, not a substance}
\begin{tutnarrative}
A constant-volume gas thermometer reads a temperature off PV = nRT. A real gas does not obey that, so the reading depends on the PRESSURE it was taken at -- and only in the limit P $\to$ 0 does it become the thermodynamic temperature. This case measures that departure, at the temperature the accompanying notes are about.

500.012 K is not a special temperature. It is the number the notes use to ask what the last digits of a temperature actually assert, and the scan is run there so the two are talking about the same state.

\end{tutnarrative}
\tutreadme{$\ast$"Thermodynamics is a funny subject. The first time you go through it, you don't understand it at all. The second time you go through it, you think you understand it, except for one or two small points. The third time you go through it, you know you don't understand it, but by that time you are so used to it, it doesn't bother you any more."$\ast$}
\tutreadme{--- attributed to Arnold Sommerfeld. PROVENANCE UNVERIFIED: Sommerfeld published no such sentence; it reaches us by report, and this repository does not quote a source it has not read. It is here as an epigraph, marked, and must be either sourced or dropped before it appears in any published guide.}
\tutreadme{What does the 012 assert? Not "roughly five hundred kelvin" --- three digits past the decimal point is a specific claim about the world. Whose claim, and resting on what?}
\tutreadme{The kelvin is no longer defined by a substance. It is defined by fixing the Boltzmann constant:}
\tutreadme{Not measured. Fixed. So a temperature is a conversion factor between energy and degrees, and nothing else. The triple point of water --- which $\ast$defined$\ast$ the kelvin for sixty-five years --- became, overnight, a quantity that is measured, with an uncertainty.}
\tutreadme{That is the first dogma, and it is worth sitting with: the definition tells you what a kelvin $\ast$means$\ast$ and gives you no way whatsoever to $\ast$measure$\ast$ one.}
\begin{tutgolden}
\goldrow{diag}{Z\_of\_P}{n\_emptyCurves}{0}
\goldrow{diag}{Z\_of\_P}{n\_failures}{0}
\goldrow{diag}{Z\_of\_P}{n\_points}{25}
\goldrow{diag}{nearly\_ideal}{Cp\_ig}{29.6050536413}
\goldrow{diag}{nearly\_ideal}{H\_R}{-0.00396437322822}
\goldrow{diag}{nearly\_ideal}{H\_ig}{5915.34999676}
\goldrow{diag}{nearly\_ideal}{H\_real}{5915.34603239}
\goldrow{diag}{nearly\_ideal}{P}{1000}
\goldrow{diag}{nearly\_ideal}{S\_R}{-4.73470040529e-05}
\goldrow{diag}{nearly\_ideal}{S\_ig}{245.041374669}
\goldrow{diag}{nearly\_ideal}{S\_real}{245.041327322}
\goldrow{diag}{nearly\_ideal}{T}{500.012}
\goldrow{diag}{nearly\_ideal}{Z}{1.00000474096}
\goldrow{diag}{nearly\_ideal}{gamma}{1.39052286566}
\goldmore{13}
\end{tutgolden}

\clearpage
\tutentry{temperature02\_engineers\_ladder}{tutorials/props/thermo/temperature02\_engineers\_ladder}
\begin{tutscheme}
\begin{tikzpicture}[>={Stealth[length=2mm]}, every node/.style={font=\scriptsize}]
  \node[draw, rounded corners=2pt, align=center, minimum width=3.1cm, minimum height=0.8cm, fill=black!4] at (0.0,-0.0) {\textbf{rung\_H2}\\[-1pt]{\tiny propertyPoint}};
  \node[draw, rounded corners=2pt, align=center, minimum width=3.1cm, minimum height=0.8cm, fill=black!4] at (3.6,-0.0) {\textbf{rung\_neon}\\[-1pt]{\tiny propertyPoint}};
  \node[draw, rounded corners=2pt, align=center, minimum width=3.1cm, minimum height=0.8cm, fill=black!4] at (7.2,-0.0) {\textbf{rung\_N2}\\[-1pt]{\tiny propertyPoint}};
  \node[draw, rounded corners=2pt, align=center, minimum width=3.1cm, minimum height=0.8cm, fill=black!4] at (0.0,-1.3) {\textbf{rung\_Ar}\\[-1pt]{\tiny propertyPoint}};
  \node[draw, rounded corners=2pt, align=center, minimum width=3.1cm, minimum height=0.8cm, fill=black!4] at (3.6,-1.3) {\textbf{rung\_O2}\\[-1pt]{\tiny propertyPoint}};
  \node[draw, rounded corners=2pt, align=center, minimum width=3.1cm, minimum height=0.8cm, fill=black!4] at (7.2,-1.3) {\textbf{rung\_CH4}\\[-1pt]{\tiny propertyPoint}};
  \node[draw, rounded corners=2pt, align=center, minimum width=3.1cm, minimum height=0.8cm, fill=black!4] at (0.0,-2.6) {\textbf{rung\_C2H4}\\[-1pt]{\tiny propertyPoint}};
  \node[draw, rounded corners=2pt, align=center, minimum width=3.1cm, minimum height=0.8cm, fill=black!4] at (3.6,-2.6) {\textbf{rung\_benzene}\\[-1pt]{\tiny propertyPoint}};
  \node[draw, rounded corners=2pt, align=center, minimum width=3.1cm, minimum height=0.8cm, fill=black!4] at (7.2,-2.6) {\textbf{rung\_water}\\[-1pt]{\tiny propertyPoint}};
  \node[draw, rounded corners=2pt, align=center, minimum width=3.1cm, minimum height=0.8cm, fill=black!4] at (0.0,-3.9) {\textbf{rung\_ethyleneGlycol}\\[-1pt]{\tiny propertyPoint}};
  \node[draw, rounded corners=2pt, align=center, minimum width=3.1cm, minimum height=0.8cm, fill=black!4] at (3.6,-3.9) {\textbf{rung\_naphthalene}\\[-1pt]{\tiny propertyPoint}};
  \node[draw, rounded corners=2pt, align=center, minimum width=3.1cm, minimum height=0.8cm, fill=black!4] at (7.2,-3.9) {\textbf{rung\_biphenyl}\\[-1pt]{\tiny propertyPoint}};
\end{tikzpicture}
\end{tutscheme}
\tutwhat{The range a chemical engineer actually works in}
\begin{tutnarrative}
THE LADDER A CHEMICAL ENGINEER ACTUALLY CLIMBS, and the engine CHECKS it rather than quoting it.

Each rung asks one question: at the temperature this substance's record calls its normal boiling point, what does its VAPOUR-PRESSURE CORRELATION say the pressure is? The answer should be one atmosphere -- that is what a normal boiling point means -- and how close it lands is a measure of the record, published rather than assumed.

THE OTHER WAY ROUND WOULD HAVE BEEN NICER AND DOES NOT WORK. Asking for `T\_bubble` at 1 atm solves for the temperature, which is the more direct question; `BubblePoint::compute` carries its own initial guess and does not converge at 20 K, a gap this repository already has on record (see docs/design/held-out-pressure.md). Verifying the pressure at a declared temperature asks the same physics of the same correlation and is robust, so that is what runs here -- said plainly rather than papered over.

WHY THE LADDER STOPS AT GLYCEROL, and it is not because the list got boring. Above roughly 600 K an ordinary organic does not have a normal boiling point to check: it CRACKS first, so there is no one-atmosphere equilibrium to ask about. The scale itself changes character there -- ITS-90 stops interpolating on vapour pressure and gas thermometry altogether and realises its upper range on the FREEZING POINTS OF METALS, then on Planck radiation above the silver point. Those fixed points are declared values of the scale, not something this engine computes, and the page says so where it lists them.

DOWTHERM A IS ABSENT AND THAT IS A FINDING, not an oversight. The hot-oil fluid an engineer actually runs a 530 K loop on is curated here with `role nonvolatile` and carries NO vapour-pressure correlation, so this case cannot ask it the question it asks the others. Named rather than quietly skipped: an omission a reader cannot see is the kind of gap this project refuses to leave.

\end{tutnarrative}
\tutreadme{Eight rungs, from liquid hydrogen to boiling water, checked by the engine rather than quoted at you.}
\tutreadme{At the temperature each substance's record calls its normal boiling point, what does its vapour-pressure correlation say the pressure is?}
\tutreadme{It should be one atmosphere --- that is what a normal boiling point $\ast$means$\ast$. How close it lands measures the record, and the answer is published rather than assumed.}
\tutreadme{The substance every engineer trusts most is the worst rung on the ladder.}
\tutreadme{Water's Antoine fit declares Trange (273 373). Its normal boiling point is 373.15 K. The fit stops 0.15 K before the temperature everybody uses it at, and evaluated there it returns 1.026 atm instead of 1.000 --- 2.6 \% out, on the one substance whose boiling point every reader could recite.}
\tutreadme{Seven cryogens the reader has no intuition for land inside 1.4 \%. The familiar one does not. Familiarity is not accuracy, and a declared validity range is worth more than a feeling.}
\begin{tutgolden}
\goldrow{diag}{rung\_H2}{Cp\_ig}{27.3234201316}
\goldrow{diag}{rung\_H2}{H\_R}{0}
\goldrow{diag}{rung\_H2}{H\_ig}{-7841.81404035}
\goldrow{diag}{rung\_H2}{H\_real}{-7841.81404035}
\goldrow{diag}{rung\_H2}{P}{101325}
\goldrow{diag}{rung\_H2}{Psat\_H2}{101850.765938}
\goldrow{diag}{rung\_H2}{S\_R}{0}
\goldrow{diag}{rung\_H2}{S\_ig}{55.7335657417}
\goldrow{diag}{rung\_H2}{S\_real}{55.7335657417}
\goldrow{diag}{rung\_H2}{T}{20.39}
\goldrow{diag}{rung\_H2}{Z}{1}
\goldrow{diag}{rung\_H2}{gamma}{1.43739708567}
\goldrow{diag}{rung\_H2}{v\_molar}{0.00167314969436}
\goldrow{diag}{rung\_neon}{Cp\_ig}{20.7861}
\goldmore{155}
\end{tutgolden}

\subsection{thermoTest}
\grouptheory{thermoTest}
\clearpage
\tutentry{atoms01\_formula\_parser}{tutorials/props/thermoTest/atoms01\_formula\_parser}
\begin{tutscheme}
\begin{tikzpicture}[>={Stealth[length=2mm]}, every node/.style={font=\scriptsize}]
  \node[draw, rounded corners=2pt, align=center, minimum width=3.1cm, minimum height=0.8cm, fill=black!4] at (0.0,-0.0) {\textbf{atoms}\\[-1pt]{\tiny elementalComposition}};
\end{tikzpicture}
\end{tutscheme}
\tutwhat{The shared formula parser, opened up: nested groups, both hydrate spellings, an ion charge that is not an element}
\begin{tutnarrative}
THE FORMULA PARSER, WITH ITS LID OFF.

The element balance is only as good as the thing that reads a formula, so that reader is ONE piece of code (src/thermo/ElementComposition) and `elementalComposition` is its glass-box surface: give it a formula, see the atoms it counted.

The formulas here are chosen to be the awkward ones, because those are where a hand-rolled parser goes wrong:

C2H6O plain, the control; Ca(OH)2 a nested group with a multiplier -- 2 O and 2 H, not 1; CaSO4:2H2O a HYDRATE, the colon form: 6 O in total, 4 from the sulfate and 2 from the two waters; MgSO4:7H2O the same salt, colon form: Mg 1, S 1, O 11, H 14; MgSO4.7H2O THE SAME SALT AS A CHEMIST WRITES IT, and the one formula here that is REFUSED. The ASCII '.' is ambiguous, and Choupo resolved it in favour of the decimal subscript -- see sepiolite below -- so this spelling would silently count the seven waters into the oxygen: O 5.7 and H 2 where the salt has O 11 and H 14. A wrong element balance from a spelling nobody would flag by eye. It refuses and names both unambiguous forms; Mg2Si3O7.5OH:3H2O why the dot cannot simply BE the hydrate separator: this is curated sepiolite, and its 7.5 is a real half-integer; SO4-2 an ION: the charge is not an element and must not become one; (NH4)2SO4 a nested group of a nested group's worth of hydrogen -- 8 H.

A parser that silently mis-reads any of these produces an element balance that closes on the wrong atoms, which is worse than one that refuses.

\end{tutnarrative}
\begin{tutgolden}
\goldrow{diag}{atoms}{C2H6O\_C}{2}
\goldrow{diag}{atoms}{C2H6O\_H}{6}
\goldrow{diag}{atoms}{C2H6O\_O}{1}
\goldrow{diag}{atoms}{CO2\_C}{1}
\goldrow{diag}{atoms}{CO2\_O}{2}
\goldrow{diag}{atoms}{ethanol\_C}{2}
\goldrow{diag}{atoms}{ethanol\_H}{6}
\goldrow{diag}{atoms}{ethanol\_O}{1}
\goldrow{diag}{atoms}{formulas\_available}{10}
\goldrow{diag}{atoms}{formulas\_total}{11}
\goldrow{diag}{atoms}{water\_H}{2}
\goldrow{diag}{atoms}{water\_O}{1}
\end{tutgolden}

\clearpage
\tutentry{h\_surface\_identity}{tutorials/props/thermoTest/h\_surface\_identity}
\begin{tutscheme}
\begin{tikzpicture}[>={Stealth[length=2mm]}, every node/.style={font=\scriptsize}]
  \node[draw, rounded corners=2pt, align=center, minimum width=3.1cm, minimum height=0.8cm, fill=black!4] at (0.0,-0.0) {\textbf{h\_identity}\\[-1pt]{\tiny hConsistency}};
\end{tikzpicture}
\end{tutscheme}
\tutwhat{The state identities of the canonical enthalpy surface: dh/dT == Cp\_phase, the 298 K vaporisation anchor, and Kirchhoff on Hvap\_state -- pinned to round-off}
\begin{tutnarrative}
thermodynamics as a REGRESSION TEST.

A surface that anchors each phase leg on the 298.15 K formation datum and integrates the TARGET phase's own Cp obeys, by construction:

(1) dh\_liq/dT = Cp\_liquid(T) (2) dh\_gas/dT = Cp\_idealGas(T) (3) h\_g(298.15) - h\_liq(298.15) = Hvap(298.15) (the anchor) (4) d[h\_g - h\_liq]/dT = Cp\_g - Cp\_liq (Kirchhoff)

This op checks all four by central finite difference for benzene, toluene, water and ethanol at 350 K and REFUSES above 1e-6 relative error -- so no future "convenient" leg (a Watson-slope liquid, a mixed-datum gas) can re-enter the balances unnoticed. The golden pins every per-species error.

\end{tutnarrative}
\begin{tutgolden}
\goldrow{diag}{h\_identity}{benzene\_anchor298\_rel}{0.000e+00}
\goldrow{diag}{h\_identity}{benzene\_dhgas\_vs\_cpig\_rel}{7.289e-11}
\goldrow{diag}{h\_identity}{benzene\_dhliq\_vs\_cpliq\_rel}{4.729e-13}
\goldrow{diag}{h\_identity}{benzene\_kirchhoff\_rel}{2.058e-10}
\goldrow{diag}{h\_identity}{ethanol\_anchor298\_rel}{1.694e-16}
\goldrow{diag}{h\_identity}{ethanol\_dhgas\_vs\_cpig\_rel}{3.826e-11}
\goldrow{diag}{h\_identity}{ethanol\_dhliq\_vs\_cpliq\_rel}{1.907e-11}
\goldrow{diag}{h\_identity}{ethanol\_kirchhoff\_rel}{1.287e-10}
\goldrow{diag}{h\_identity}{toluene\_anchor298\_rel}{0.000e+00}
\goldrow{diag}{h\_identity}{toluene\_dhgas\_vs\_cpig\_rel}{6.085e-11}
\goldrow{diag}{h\_identity}{toluene\_dhliq\_vs\_cpliq\_rel}{1.854e-13}
\goldrow{diag}{h\_identity}{toluene\_kirchhoff\_rel}{2.214e-10}
\goldrow{diag}{h\_identity}{water\_anchor298\_rel}{6.529e-16}
\goldrow{diag}{h\_identity}{water\_dhgas\_vs\_cpig\_rel}{3.627e-11}
\goldmore{3}
\end{tutgolden}

\clearpage
\tutentry{model4\_speciation\_scaling}{tutorials/props/thermoTest/model4\_speciation\_scaling}
\begin{tutscheme}
\begin{tikzpicture}[>={Stealth[length=2mm]}, every node/.style={font=\scriptsize}]
  \node[draw, rounded corners=2pt, align=center, minimum width=3.1cm, minimum height=0.8cm, fill=black!4] at (0.0,-0.0) {\textbf{brine}\\[-1pt]{\tiny speciate}};
  \node[draw, rounded corners=2pt, align=center, minimum width=3.1cm, minimum height=0.8cm, fill=black!4] at (3.6,-0.0) {\textbf{concentrate}\\[-1pt]{\tiny scalingScan}};
\end{tikzpicture}
\end{tutscheme}
\tutwhat{RO membrane-scaling audit of a brackish groundwater: aqueous speciation (Davies) with pH SOLVED from electroneutrality (the feed charge imbalance is announced first) + saturation indices SI = log10(IAP/K), scanned over water recovery to r = 0.85 (concentrate = feed/(1-r)) TWICE: closed system (DIC conserved and concentrating) vs open to the atmosphere (a(CO2aq) pinned by Henry at pCO2 = 4.2e-4 atm, DIC a solved outcome -- this CO2-rich water degasses, pH and calcite SI jump above the closed curve) -- the calcite SI curves differ; that contrast is the lesson. Three analytic pins lock the physics: pure water at given pH 7 reproduces Kw (m\_OH \textasciitilde{} 1e-7), a back-computed gypsum-saturated CaSO4 solution gives SI\_gypsum \textasciitilde{} 0.00, and pure water with pH solved lands at exactly 7.000.}
\begin{tutnarrative}
thermoTest / model4\_speciation\_scaling

The SAME system \{NaCl, CaSO4, water\} in its THIRD role: a MULTI-ION broth. Here we no longer think of "NaCl" or "CaSO4" -- the aqueous speciation engine reasons in the shared ions Na+, Cl-, Ca2+, SO4-2 (plus the neutral ion pair CaSO4aq), and reports SATURATION INDICES SI = log10(IAP/Ksp) for every solid the ions can form: gypsum (CaSO4.2H2O) and halite (NaCl).

Na+ comes from NaCl, Cl- from NaCl, Ca2+/SO4-2 from CaSO4 -- but once dissolved there are no salt boundaries, only ions. This is why a "salt" cannot carry its own activity model: the ion pool is shared (the species/ tier of the tree).

MODEL: Pitzer-HMW is MANDATORY here. The ionic strength is I \textasciitilde{} 3.3 mol/kg (NaCl-dominated); the Davies law is only trustworthy to I \textasciitilde{} 0.5 m, so the `activityModel` is declared explicitly (no silent default) -- with Davies the saturation indices at this I are quantitatively meaningless.

The brine is NaCl-dominated (3 mol/kg, near seawater) with CaSO4 (0.10 mol/kg) right at gypsum saturation (SI \textasciitilde{} 0). Charge balances exactly (Na + 2 Ca = Cl + 2 SO4), so pH is pinned at 7. The scan concentrates the brine (evaporation; r = fraction of water removed): gypsum scales first, at r \textasciitilde{} 0.035 (it starts saturated), halite only near r \textasciitilde{} 0.50. Note m\_CaSO4aq \textasciitilde{} 0.005: the neutral ion pair is a real fraction of the calcium -- "dissociation" is partial.

\end{tutnarrative}
\begin{tutgolden}
\goldrow{diag}{brine}{I}{3.26782}
\goldrow{diag}{brine}{SI\_anhydrite}{-0.22134}
\goldrow{diag}{brine}{SI\_gypsum}{-0.0221674}
\goldrow{diag}{brine}{SI\_halite}{-0.910511}
\goldrow{diag}{brine}{SI\_mirabilite}{-2.13799}
\goldrow{diag}{brine}{SI\_thenardite}{-2.98985}
\goldrow{diag}{brine}{aw}{0.890402}
\goldrow{diag}{brine}{gamma\_CaSO4aq}{1}
\goldrow{diag}{brine}{m\_CaSO4aq}{0.00466294}
\goldrow{diag}{brine}{pH}{7}
\goldrow{diag}{concentrate}{I\_r0}{3.26782}
\goldrow{diag}{concentrate}{I\_rmax}{15.6285}
\goldrow{diag}{concentrate}{SI\_anhydrite\_r0}{-0.22134}
\goldrow{diag}{concentrate}{SI\_anhydrite\_rmax}{1.39807}
\goldmore{26}
\end{tutgolden}

\subsection{transport}
\grouptheory{transport}
\clearpage
\tutentry{transport01\_gas\_bench}{tutorials/props/transport/transport01\_gas\_bench}
\begin{tutscheme}
\begin{tikzpicture}[>={Stealth[length=2mm]}, every node/.style={font=\scriptsize}]
  \node[draw, rounded corners=2pt, align=center, minimum width=3.1cm, minimum height=0.8cm, fill=black!4] at (0.0,-0.0) {\textbf{bench}\\[-1pt]{\tiny transportBench}};
\end{tikzpicture}
\end{tutscheme}
\tutwhat{Every transport correlation proves itself against its own anchor -- theory or arithmetic, said which -- and then all are asked the same question for four gases at 373 K; NOTHING here is measured}
\begin{tutnarrative}
transport01 -- THE TRANSPORT CORRELATIONS, AS OBJECTS THAT PROVE THEMSELVES.

NOTHING IN THIS CASE IS MEASURED. Read that first. Choupo's catalogue carries no measured gas viscosity, thermal conductivity or diffusivity, so no number below is compared with an experiment. What the case shows is what CAN be shown without one, and it says which kind each check is:

THEORY an identity kinetic theory or the correlation's own author supplies -- Eucken on a monatomic gas MUST give the Chapman-Enskog lambda = (15/4) R mu / M exactly; Wilke's phi\_ii MUST be 1 for identical species. A wrong constant fails here. ARITHMETIC the correlation's own closed form at a stated point, written as a literal -- Chung's Neufeld Omega at T$\ast$ = 1, Fuller at round numbers, Sato-Riedel at T = Tb where its bracket ratio is 1 by construction, the fitted liquid forms reducing to Andrade. This checks the TRANSCRIPTION and the units, and the output says so rather than letting 0.0000 \% read as a validation.

WHAT THE RUN DOES, IN TWO PARTS. (1) VERIFY. Every registered model in five families, plus the shared gas mixing rule, reproduces its own anchor; each prints its window, its PRIMARY citation and the anchor's kind. (2) COMPARE. N2, O2, CO2 and water vapour at 373.15 K, 1 bar: every registered gas-viscosity and gas-conductivity model for every component, each with its computable window verdict (Chung's T$\ast$ against the Neufeld domain; Eucken's "polyatomic" by the record's own Cp), then the case's mixing rules -- Wilke for mu, Wassiljewa / Mason-Saxena for k -- at the equimolar composition (announced).

WHAT TO READ IN THE OUTPUT. Three things. The spread across models is 0.00 \% -- for BOTH gas families -- and the bench says why: ONE model is registered per family. A spread over one is not agreement; Chapman-Enskog (Lennard-Jones sigma and eps/k, which no record carries yet) is the named next member. Chung ANNOUNCES, for each of the four components, that Vc was estimated from Zc(omega) and that the polar and association terms were dropped. Water is polar; the announcement is the only thing in the run that says the non-polar form was applied to it anyway. Eucken flags CO2 and water (and, by their Cp, N2 and O2 -- diatomics are polyatomic) as APPROXIMATION HERE, because plain Eucken is exact only for a monatomic gas and this case has none.

WHAT THIS CASE DOES NOT DO. It does not say which model is right, and it does not compare with a measurement. The golden pins that the anchors, the values and the (absent) spread have not MOVED -- nothing more. Record: docs/design/transport-correlations-as-objects.md.

\end{tutnarrative}
\begin{tutgolden}
\goldrow{diag}{bench}{allPass}{1}
\goldrow{diag}{bench}{dev\_Andrade}{5.35123586857e-08}
\goldrow{diag}{bench}{dev\_Chung}{2.53636147631e-07}
\goldrow{diag}{bench}{dev\_Eucken}{0}
\goldrow{diag}{bench}{dev\_Fuller}{5.617814941e-08}
\goldrow{diag}{bench}{dev\_SatoRiedel}{3.4111766641e-06}
\goldrow{diag}{bench}{dev\_Vogel}{0}
\goldrow{diag}{bench}{dev\_chemsepEq101}{0}
\goldrow{diag}{bench}{dev\_chemsepEq16}{0}
\goldrow{diag}{bench}{dev\_wilkePhi}{0}
\goldrow{diag}{bench}{k\_Eucken\_CO2}{0.021231358218}
\goldrow{diag}{bench}{k\_Eucken\_N2}{0.0297082727854}
\goldrow{diag}{bench}{k\_Eucken\_O2}{0.0306615645526}
\goldrow{diag}{bench}{k\_Eucken\_water}{0.0250221542205}
\goldmore{21}
\end{tutgolden}

\clearpage
\tutentry{transport02\_chapman\_enskog}{tutorials/props/transport/transport02\_chapman\_enskog}
\begin{tutscheme}
\begin{tikzpicture}[>={Stealth[length=2mm]}, every node/.style={font=\scriptsize}]
  \node[draw, rounded corners=2pt, align=center, minimum width=3.1cm, minimum height=0.8cm, fill=black!4] at (0.0,-0.0) {\textbf{bench}\\[-1pt]{\tiny transportBench}};
\end{tikzpicture}
\end{tutscheme}
\tutwhat{Chapman-Enskog from Svehla's own force constants: seven gases at 300 K priced by two viscosity and two conductivity models, the spread a real disagreement for the first time; nothing measured here, the sigma/eps-k are Svehla's fits to pre-1962 viscosities, cited by page}
\begin{tutnarrative}
transport02 -- CHAPMAN-ENSKOG FROM SVEHLA'S CONSTANTS, AND THE FIRST REAL SPREAD.

NOTHING IN THIS CASE IS MEASURED, except in the one sense the anchors state. The seven Lennard-Jones pairs (sigma, eps/k) on the records in constant/components/ are Svehla's LEAST-SQUARES FITS to viscosities measured before 1962 (his Table I(a) method 1, one row per gas, cited by report page in each record's own provenance block); Choupo holds none of those measurements and compares with none. What the run shows:

TWO models per gas family, so the spread column is a DISAGREEMENT and not the 0.00 \% of transport01: Chung (corresponding states from Tc, Pc, omega -- non-polar truncation, Vc estimated, both announced) beside Chapman-Enskog (kinetic theory over the fitted sigma, eps/k); plain Eucken beside the modified Eucken (Svehla eq. (2), the internal modes weighted 1.32). Water is where the two viscosities part furthest: Chung's non-polar form against a pair FITTED to steam data. Argon is the control: monatomic, so Eucken and the modified Eucken coincide to the extent the record's Cp polynomial returns 5R/2 (it does, to four figures at 300 K; lambda'' is the residual) and the polyatomic flag stays silent. Every Chapman-Enskog component is ANNOUNCED with the ORIGIN of both constants (FITTED to measured transport data, Svehla method 1) and the page they were read from; an estimated pair would say ESTIMATED in the same place. The window verdict is T$\ast$ = kT/eps against the Neufeld fit's 0.3 <= T$\ast$ <= 100; water at 300 K sits at T$\ast$ = 0.37, inside, near the low edge. The comparison point is 300 K, 1 bar -- the 300 K row of Svehla's Table III -- so a reader can set the printed viscosities beside his (report pages 31 Ar, 56 CH4, 58 CO2, 74 H2, 75 H2O, 90 N2, 98 O2) and see the residual the anchor row measures: the Neufeld 1972 fit of the collision integral against the Hirschfelder 1954 tables he interpolated, 0.1-0.8 \% across these gases, 1.8 \% for water.

THE ANCHORS. ChapmanEnskog: ARITHMETIC -- with Svehla's own M, sigma, eps/k for N2 the engine returns his printed eta at 300 K (report page 90) within a DECLARED 1 \% (0.40 \% observed); it checks the constant 26.693, the unit chain and the shared collision-integral home, nothing about any gas. modifiedEucken: THEORY -- a monatomic gas must give exactly (15/4) R mu / M -- after reproducing Svehla's printed N2 conductivity from his own eta and Cp/R through eq. (2). The bench prints the kind.

WHAT THIS CASE DOES NOT DO. It does not say which viscosity is right for water or for anything else, and it does not compare with a measurement. The records are CASE-LOCAL (adopted, on the edwards02 precedent): the public catalogue carries no lennardJones block until Vitor promotes one; bin/curate/propose\_lennard\_jones.py drafts them into data/local/ from the committed transcription bin/curate/svehla1962/. Binary diffusivity by Chapman-Enskog is NOT here: the Omega(1,1)$\ast$ fit is not citable from any source in the tree. Record: docs/design/transport-correlations-as-objects.md, section 8.

\end{tutnarrative}
\begin{tutgolden}
\goldrow{diag}{bench}{allPass}{1}
\goldrow{diag}{bench}{dev\_Andrade}{5.35123586857e-08}
\goldrow{diag}{bench}{dev\_ChapmanEnskog}{0.00400952242196}
\goldrow{diag}{bench}{dev\_Chung}{2.53636147631e-07}
\goldrow{diag}{bench}{dev\_Eucken}{0}
\goldrow{diag}{bench}{dev\_Fuller}{5.617814941e-08}
\goldrow{diag}{bench}{dev\_SatoRiedel}{3.4111766641e-06}
\goldrow{diag}{bench}{dev\_Vogel}{0}
\goldrow{diag}{bench}{dev\_chemsepEq101}{0}
\goldrow{diag}{bench}{dev\_chemsepEq16}{0}
\goldrow{diag}{bench}{dev\_modifiedEucken}{0}
\goldrow{diag}{bench}{dev\_wilkePhi}{0}
\goldrow{diag}{bench}{k\_Eucken\_Ar}{0.0178116749867}
\goldrow{diag}{bench}{k\_Eucken\_CH4}{0.0321131371617}
\goldmore{37}
\end{tutgolden}

